% ############################################### PREAMBLE AND CONTROL PANEL #

% ########################################################### DOCUMENT CLASS #
\documentclass[  %   draft
             % , demo  % black boxes instead of figures
              , 8pt, twoside, openright]{extbook}

% ########################################################## AND CONTROL PANEL #

% true - hmarginration 1:1, urlcolor - darkgreen, version - for searching
\newif\ifscreen 	   \screentrue
% true - no apendix, no author of grammer, no backmatter ISBN, 
\newif\ifPDF               \PDFtrue
% true - rotate pages of morph verb tables
\newif\ifPDFrotate         \PDFrotatetrue
% true - adds white searchable value !
\newif\ifPDFsearchvalue    \PDFsearchvaluetrue
\newif\ifmakecovers        \makecoverstrue
\newif\ifmakefrontmatter   \makefrontmattertrue
\newif\ifinputintroduction \inputintroductiontrue
\newif\ifinputletters      \inputletterstrue
\newif\ifinputphon         \inputphontrue
\newif\ifinputmorpho       \inputmorphotrue
\newif\ifmakebackmatter    \makebackmattertrue
\newif\ifshowcrop           \showcroptrue
% corrections
\newif\ifC                  \Ctrue

\newif\ifPRINTHOUSE 	   \PRINTHOUSEfalse
\newif\ifPDFSEARCH  	   \PDFSEARCHfalse
\newif\ifPDFPRINT 	   \PDFPRINTfalse
\newif\ifPDFtest	   \PDFtesttrue

%versions
%PRINTHOUSE
\ifPRINTHOUSE
\makecoversfalse
\screenfalse
\PDFfalse
\PDFrotatefalse
\PDFsearchvaluefalse
\showcroptrue
\Ctrue
\fi

%PDFSEARCH
\ifPDFSEARCH
\screentrue
\PDFtrue
\PDFrotatetrue
\PDFsearchvaluetrue
\showcropfalse
\Ctrue
\fi

%PDFPRINT
\ifPDFPRINT
\screenfalse
\PDFtrue
\PDFrotatefalse
\PDFsearchvaluefalse
\showcropfalse
\Ctrue
\fi

%PDFtest
\ifPDFtest
\makecoversfalse
\makefrontmatterfalse
\inputintroductionfalse
\inputletterstrue
\inputphonfalse
\inputmorphofalse
\makebackmatterfalse
\screentrue
\PDFtrue
\PDFrotatefalse
\PDFsearchvaluefalse
\showcropfalse
\Cfalse
\fi

% 1. PDF (for printing)  
% a) searchable value, b) no apendix, c) version= PDF d) no ISBN and no ISBN bar
% e) automatic rotation of landscape pages, f) no author of Icelandic Grammer
% 2. PDF (for viewing offline and searching) + screen (horizontal ratio 1:1) 
% 3. PRINTHOUSE
% a) no searchable value, b) apendix, c) version ISBN d) ISBN bar
% e) no rotation of landscape pages, f)urllinkcolor black
% f) A4 paper dimentions, C5 layout, shows cropping marks

% ########################################################## PAGE GEOMETRIES #

\usepackage%[showframe]
           {geometry}
           

\ifshowcrop
% Basic geometry of document
\geometry
  { headsep    =   \baselineskip
  , textwidth  = 42\baselineskip
  , textheight = 60\baselineskip
  , hmarginratio = \ifscreen 1:1\else 2:3\fi
  , vmarginratio = 2:3
  , bindingoffset = 0cm
  , onecolumn
  , a4paper 
  , layout=c5paper,
  , layouthoffset=\dimexpr(\paperwidth-\csname Gm@layoutwidth\endcsname)/2\relax
   , layoutvoffset=\dimexpr(\paperheight-\csname Gm@layoutheight\endcsname)/2\relax
  , showcrop
  }
\else
\geometry
  { headsep    =   \baselineskip
  , textwidth  = 42\baselineskip
  , textheight = 60\baselineskip
  , hmarginratio = \ifscreen 1:1\else 2:3\fi
  , vmarginratio = 2:3
  , bindingoffset = 0cm
  , onecolumn,
  , c5paper
  }
\fi

\ifshowcrop
% Geometry of the dictionary section
\newcommand\dictionarygeometry{\newgeometry
  { textwidth  = 42\baselineskip
  , textheight = 60\baselineskip
  , hmarginratio = \ifscreen 1:1\else 2:3\fi
  , vmarginratio = 3:2
  , bindingoffset = 0cm
  , twocolumn
  , a4paper 
  , layout=c5paper
  , layouthoffset=\dimexpr(\paperwidth-\csname Gm@layoutwidth\endcsname)/2\relax
   , layoutvoffset=\dimexpr(\paperheight-\csname Gm@layoutheight\endcsname)/2\relax
  , showcrop
  }}
   % Geometry of the cover pages
\newcommand\covergeometry{\newgeometry
  { textwidth  = 42\baselineskip
  , textheight = 60\baselineskip
  , hmarginratio = 1:1
  , vmarginratio = 2:3
  , bindingoffset = 0cm
  , onecolumn
  , layout=c5paper
 , layouthoffset=\dimexpr(\paperwidth-\csname Gm@layoutwidth\endcsname)/2\relax
   , layoutvoffset=\dimexpr(\paperheight-\csname Gm@layoutheight\endcsname)/2\relax
  , showcrop
  }} 
\else
% Geometry of the dictionary section
\newcommand\dictionarygeometry{\newgeometry
  { textwidth  = 42\baselineskip
  , textheight = 60\baselineskip
  , hmarginratio = \ifscreen 1:1\else 2:3\fi
  , vmarginratio = 3:2
  , bindingoffset = 0cm
  , twocolumn,
  , c5paper
  }}
   % Geometry of the cover pages
\newcommand\covergeometry{\newgeometry
  { textwidth  = 42\baselineskip
  , textheight = 60\baselineskip
  , hmarginratio = 1:1
  , vmarginratio = 2:3
  , bindingoffset = 0cm
  , onecolumn
  , c5paper
  }} 
\fi

% ################################################################# LANGUAGES #

\usepackage[icelandic, latin, czech]{babel}
\usepackage{csquotes}


% ################################################################## ENCODING #

\usepackage[utf8]{inputenc}

% Smashing uppercase letters with diacritics
% ------------------ THIS IS EVIL AND SHOULD NEVER BE DONE -------------------
% Don't smash anything if we are on the first run to produce the index
\IfFileExists{figures.idx}{
  % Czech uppercase letters with diacritics:
  \DeclareUnicodeCharacter{00C1}{\strut\smash{\' A}} % Á
  \DeclareUnicodeCharacter{010C}{\strut\smash{\v C}} % Č
  \DeclareUnicodeCharacter{010E}{\strut\smash{\v D}} % Ď
  \DeclareUnicodeCharacter{00C9}{\strut\smash{\' E}} % É
  \DeclareUnicodeCharacter{011A}{\strut\smash{\v E}} % Ě
  \DeclareUnicodeCharacter{00CD}{\strut\smash{\' I}} % Í
  \DeclareUnicodeCharacter{0147}{\strut\smash{\v N}} % Ň
  \DeclareUnicodeCharacter{00D3}{\strut\smash{\' O}} % Ó
  \DeclareUnicodeCharacter{0158}{\strut\smash{\v R}} % Ř
  \DeclareUnicodeCharacter{0160}{\strut\smash{\v S}} % Š
  \DeclareUnicodeCharacter{0164}{\strut\smash{\v T}} % Ť
  \DeclareUnicodeCharacter{00DA}{\strut\smash{\' U}} % Ú
  \DeclareUnicodeCharacter{016E}{\strut\smash{\r U}} % Ů
  \DeclareUnicodeCharacter{00DD}{\strut\smash{\' Y}} % Ý
  \DeclareUnicodeCharacter{017D}{\strut\smash{\v Z}} % Ž
  % Icelandic uppercase letters with diacritics: ÁÉÍÓÚÝ (ÐÞÆ) and
  \DeclareUnicodeCharacter{00D6}{\strut\smash{\" O}} % Ö
}{}
% ------------------ THAT WAS EVIL AND SHOULD NEVER BE DONE ------------------

% Non-breakable space
\DeclareUnicodeCharacter{00A0}{~}

% Non-breakable hyphen
\newcommand*{\nobreakhyphen}{\mbox{-}}
\DeclareUnicodeCharacter{2011}{\nobreakhyphen}


% ############################################################ FONTS, SYMBOLS #

\usepackage[T1]{fontenc}

\usepackage{tgpagella}
\usepackage[scaled=0.88]{helvet}      % relative scale of the two fonts

\def\phvfamily{\fontfamily{phv}\selectfont}

%\usepackage{stmaryrd}
% Let's keep this minimal instead of using the package:
\DeclareSymbolFont{stmry}{U}{stmry}{m}{n}
%\SetSymbolFont{stmry}{bold}{U}{stmry}{b}{n}
\DeclareMathSymbol\shortrightarrow\mathrel{stmry}{"01}
\DeclareMathSymbol\shortuparrow\mathrel{stmry}{"02}

% Methinks these are not necessary
%\usepackage{dingbat}
%\usepackage{manfnt}
%\usepackage{latexsym}

\usepackage{pifont}




% #################################################################### LAYOUT #

% \normalsize should be {8pt}{9.6pt}
\def\HUGE {\fontsize{23.887872pt}{3\baselineskip}\selectfont}
\def\Huge {\fontsize{19.906560pt}{3\baselineskip}\selectfont}
\def\huge {\fontsize{16.588800pt}{3\baselineskip}\selectfont}
\def\LARGE{\fontsize{13.824000pt}{2\baselineskip}\selectfont}
\def\Large{\fontsize{11.520000pt}{2\baselineskip}\selectfont}
\def\large{\fontsize{ 9.600000pt}{2\baselineskip}\selectfont}

% Two columns layout ruler
\setlength\columnsep    {2\baselineskip}
\setlength\columnseprule{0.4pt}

% Necessary for baseline alignment
\topskip=\baselineskip
\raggedbottom
\setlength\parskip{0pt} % it's better to avoid glue

% Temporarily suppress warnings
\hbadness=5000
\vbadness=5000

\setlength\emergencystretch{17pt}

% Allow smaller emergencystretch in several cases
\newenvironment{xtolerant}[2]{%
  \par
  \ifx\empty#1\empty\else\tolerance=#1\relax\fi
  \ifx\empty#2\empty\else\emergencystretch=#2\relax\fi
}{%
  \par
}

% ######################################################### ASSORTED PACKAGES #

\usepackage{tikz}
\usetikzlibrary{calc}

\usepackage{tipa}                      % IPA phonetics

\usepackage{fourier-orns}              % ornaments
\usepackage{amsmath}                   % non-breakable dash
\usepackage{hyphenat}                  % no hyphen in abbreviations

\usepackage{hanging}
\usepackage{fix2col}
\usepackage{fixltx2e}                  % subscript
\usepackage{pagecolor}
\usepackage{afterpage}                 % Used in coloring trick for covers
\usepackage{pdfpages}

\usepackage{multirow}
\usepackage{multicol}
\usepackage{rotating}

\usepackage{etoolbox} % http://ctan.org/pkg/etoolbox
%\usepackage{newunicodechar} 		% new definition of middle dot
\usepackage[ISBN=978-80-260-2385-2,SC0]{ean13isbn}

\usepackage{pgffor}

\usepackage{newunicodechar}

\newunicodechar{·}{\chejnikcdot}
\makeatletter
\newcommand{\chejnikcdot}{%
  \ifnum\lastpenalty=10042
    \kern-0.06667em\relax % dots should always be kerned
  \fi
  \textperiodcentered
  \@ifnextchar|{\kern\chejnikkern\relax}{\penalty10042 }%
}

\newcommand{\closeupdotbar}{%
  \renewcommand\chejnikkern{-0.06667em}%
}
\newcommand\chejnikkern{0pt} % default

% Middle dot + | little bit closer, nicer
%\newunicodechar{·}{\advancedcdot}
%\makeatletter
%\newcommand*\advancedcdot{\textperiodcentered
%  \@ifnextchar|{\hspace{-0.03667em}}{\checkifcdotnext}}
%\newcommand*\checkifcdotnext[2]
%  {\def\tmpa{·}\def\tmpb{#1#2}\ifx\tmpa\tmpb\hspace{-0.06667em}\fi#1#2}
%\makeatother

% ############################################################ BASELINE GRID #

\usepackage{eso-pic}
\usepackage{tikzpagenodes}

% Command to draw a baseline grid
\newcommand\drawbaselinegrid{%
  \begin{tikzpicture}[overlay,remember picture]
    \draw [red!30!white, ultra thin, dashed]
      (0,0) grid [ ystep  = \baselineskip, xstep = \textwidth
                 , shift  = (current page text area.north west)
                 , yshift = -\dp\strutbox
                 ] ++(\textwidth,-\textheight);
    \draw [red!30!white, thin]
      (0,0) grid [ step = \baselineskip, xstep = \textwidth
                 , shift=(current page text area.north)
                 ] ++(0.5\textwidth,-\textheight)
      (0,0) grid [ step = \baselineskip, xstep = \textwidth
                 , shift=(current page text area.north)
                 ] ++(-0.5\textwidth,-\textheight);
    \draw[black!10!white]
      (current page text area.north west)
        rectangle (current page text area.south east);
  \end{tikzpicture}}

% Draw a baseline grid on every page
% \AddToShipoutPicture{\drawbaselinegrid}


% ################################################################## METADATA #

\title{%
  \textbf{Islandsko-český studijní slovník}%
  \thanks{Tato kniha byla vytvořena v \LaTeX{}u pod Ubuntu 14.04.%
    Poděkování patří všem autorům, kteří publikují pod svobodnými licencemi.}}
\author{Aleš Chejn, Ján Zaťko, Jón Gíslason}
\date{říjen 2011}


% ############################################################## BIBLIOGRAPHY #

\usepackage[style=authoryear]{biblatex}
\addbibresource{slovnik.bib}


% ##################################################################### INDEX #

\usepackage{imakeidx}

% Index of authors of photographs
\makeindex
  [ name = figures
  , title = {Seznam autorů fotografií a ilustrací}
  , intoc
  , columnseprule
  , columnsep = \columnsep
  , noautomatic                       % We handle this in make.sh
  ]

% icelandic rules start
\usepackage{filecontents}
\def\mygroup#1{{\bfseries#1}}
\begin{filecontents*}{style.xdy}
;; style.xdy
(markup-letter-group :open-head "~n\mygroup{" :close-head "}")
\end{filecontents*}
% icelandic rules end


% ############################################################### PAGE STYLES #

\usepackage{fancyhdr}

% -------------------------------------- BASIC PAGE STYLE
\fancypagestyle{plain}{
  \fancyhf\relax                       % Clear header/footer
  \renewcommand \headrulewidth {0.0pt} % No header rule
  \renewcommand \footrulewidth {0.0pt} % No footer rule
  \fancyfoot[C]{\thepage}              % Page number in footer, centred
}

% -------------------------------------- INDEX PAGE STYLE  
%\fancypagestyle{indexstyle}{
%  \fancyhf\relax                       % Clear header/footer
%  \renewcommand \headrulewidth {0.4pt} % Thin header rule
%  \renewcommand \footrulewidth {0.0pt} % No footer rule
%  \fancyhead[C]{\thepage}              % Page number in footer, centred
%}

% -------------------------------------- DICTIONARY PAGE STYLE
\fancypagestyle{myheadings}{
  \fancyhf\relax                       % Clear header/footer
  \renewcommand \headrulewidth {0.4pt} % Thin header rule
  \fancyhead[CO,CE]{\thepage}          % Page number in header, centred
  % NOTE: the page numbers will be printed when the dictionary is ready
  \fancyhead[LE,LO]{\phvfamily\bfseries\rightmark}
  \fancyhead[RE,RO]{\phvfamily\bfseries\leftmark}
}

\pagestyle{plain}  

% ##################################################################### LISTS #

\usepackage{expdlist}
\usepackage{enumitem}

\setlist{nosep, topsep=\baselineskip}
\setlist[itemize]{leftmargin=*,label=\scalebox{.6}{\textbullet}}


% #################################################################### FLOATS #

\usepackage{caption}
\captionsetup{labelformat=empty}

% Alter some LaTeX defaults for better treatment of figures:
% See p.105 of "TeX Unbound" for suggested values. 
% See pp. 199-200 of Lamport's "LaTeX" book for details.
% ------------------------------------- Parameters for ALL pages:
\renewcommand \topfraction    {0.9}   %   max fraction of floats at top
\renewcommand \bottomfraction {0.8}   %   max fraction of floats at bottom
% ------------------------------------- Parameters for TEXT (non float) pages:
\setcounter{topnumber}    {1}
\setcounter{bottomnumber} {1}
\setcounter{totalnumber}  {1}         %   2 may work better
\setcounter{dbltopnumber} {2}         %   for 2-column pages
\renewcommand \dbltopfraction   {0.9} %   fit big float above 2-col. text
\renewcommand \textfraction     {0.07}%   allow minimal text w. figs
% ------------------------------------- Parameters for FLOAT (non text) pages:
\renewcommand \floatpagefraction{0.7} %   require fuller float pages
% N.B.: floatpagefraction MUST be less than topfraction !!
\renewcommand{\dblfloatpagefraction}{0.7} % require fuller float pages
% remember to use [htp] or [htpb] for placement

\makeatletter
\setlength{\@fptop}{0pt}
\setlength{\@fpbot}{0pt plus 1fil}
\makeatother

% This kills spacing around floats
\setlength \intextsep        {0pt}
\setlength \floatsep         {0pt}
\setlength \textfloatsep     {0pt}
\setlength \dbltextfloatsep  {0pt}
\setlength \dblfloatsep      {0pt}

% This kills spacing around captions
\setlength \abovecaptionskip {0pt}
\setlength \belowcaptionskip {0pt}


% #################################################################### TABLES #

\usepackage{booktabs}
\usepackage{tabularx}
\usepackage{longtable}
\usepackage{ltxtable}
\usepackage{lscape}

\ifPDFrotate
\else
\makeatletter			       
\renewcommand\PLS@Rotate[1]{}          % this command prevents pdflscape 
\makeatother                           % to rotate pages
\fi

%\LetLtxMacro\oldtabularx\tabularx
%\LetLtxMacro\oldendtabularx\endtabularx
%\renewenvironment{tabularx}[1]
%  {\oldtabularx{\columnwidth}{#2}}
%  {\oldendtabularx}

\newcommand\LTXfw[1]{\LTXtable{\columnwidth}{#1}}

%%% \LetLtxMacro\oldtabularx\tabularx
%%% \LetLtxMacro\oldendtabularx\endtabularx
%%% \newenvironment{longtabularx}[2]
%%%   {\oldtabularx{#1}{#2}}
%%%   {\oldendtabularx}

% ######################################################### czech babel vs cmidrule
% ###################################### http://tex.stackexchange.com/a/112001/8380
\begingroup
  \makeatletter
  \catcode`\-=\active
  \AtBeginDocument{
  \def\@@@cmidrule[#1-#2]#3#4{\global\@cmidla#1\relax
    \global\advance\@cmidla\m@ne
    \ifnum\@cmidla>0\global\let\@gtempa\@cmidrulea\else
    \global\let\@gtempa\@cmidruleb\fi
    \global\@cmidlb#2\relax
    \global\advance\@cmidlb-\@cmidla
    \global\@thisrulewidth=#3
    \@setrulekerning{#4}
    \ifnum\@lastruleclass=\z@\vskip \aboverulesep\fi
    \ifnum0=`{\fi}\@gtempa
    \noalign{\ifnum0=`}\fi\futurenonspacelet\@tempa\@xcmidrule}
  }
\endgroup

% macro that supress cmidrule at page brake
% http://tex.stackexchange.com/questions/110548/using-longtable-with-cmidrule-page-breaking-issues
\makeatletter
\def\LT@output{%
  \ifnum\outputpenalty <-\@Mi
    \ifnum\outputpenalty > -\LT@end@pen
      \LT@err{floats and marginpars not allowed in a longtable}\@ehc
    \else
      \setbox\z@\vbox{\unvbox\@cclv}%
      \ifdim \ht\LT@lastfoot>\ht\LT@foot
        \dimen@\pagegoal
        \advance\dimen@-\ht\LT@lastfoot
        \ifdim\dimen@<\ht\z@
          \setbox\@cclv\vbox{\unvbox\z@
\copy\LT@foot\vss}%
          \@makecol
          \@outputpage
          \setbox\z@\vbox{\box\LT@head}%
        \fi
      \fi
      \global\@colroom\@colht
      \global\vsize\@colht
      \vbox
        {\unvbox\z@
\box\ifvoid\LT@lastfoot\LT@foot\else\LT@lastfoot\fi}%
    \fi
  \else
    \setbox\@cclv\vbox{\unvbox\@cclv

\setbox0\lastbox
\ifdim\ht0=0.29999pt \unskip % hope that was a cline we threw away
\else
\nointerlineskip
\box0     % put it back, whatever it was
\fi

\copy\LT@foot\vss}%
    \@makecol
    \@outputpage
      \global\vsize\@colroom
    \copy\LT@head\nobreak
  \fi}

\makeatother

% These values ensure that for top, mid and bottom rules
% space above + rule width + space below = baselineskip
\setlength \abovetopsep    {0.43\baselineskip}
\setlength \heavyrulewidth {0.10\baselineskip}
\setlength \belowbottomsep {0.43\baselineskip}
\setlength \aboverulesep   {0.47\baselineskip}
\setlength \lightrulewidth {0.06\baselineskip}
\setlength \belowrulesep   {0.47\baselineskip}

%\usepackage{letltxmacro}

%%% \LetLtxMacro\oldtabular\tabular
%%% \LetLtxMacro\oldendtabular\endtabular
%%% \renewenvironment{tabular}[2][t]
%%%   {\oldtabular[#1]{#2}}
%%%   {\oldendtabular}

%%% \LetLtxMacro\oldtable\table
%%% \LetLtxMacro\oldendtable\endtable
%%% \renewenvironment{table}[1][h]
%%%   {\leavevmode\oldtable[#1]\vspace{\dp\strutbox}}
%%%   {\vspace{-\dp\strutbox}\oldendtable}
%%%   % This aligns a [t] aligned tab-envo to baseline grid


%%% % declention and conjugation tables
%%% \usepackage{floatrow}
%%% \DeclareFloatFont{footnotesize}{\footnotesize}
%%% % "scriptsize" is defined by floatrow, "tiny" not
%%% \floatsetup[table]{font=footnotesize}

% ################################################################### FIGURES #

\usepackage{adjustbox} % loads also graphicx

\graphicspath{%
  {/home/chejnik/Dokumenty/hvalur.org/images/biolib/full/}%
  {/home/chejnik/Dokumenty/hvalur.org/images/uploaded_files/}}

\newlength\dicfigheight
\newlength\dicfigskip

\newcommand\dicFigure[3]{ % {filename}{caption}{index entry}
  \setlength\fboxsep{0pt}\setlength\fboxrule{0.5pt}
  \def\dicfigbox{\fbox{%
    \adjincludegraphics [ max height = 0.8\columnwidth
                        , max width  = 0.8\columnwidth ] {#1}}}
  \setlength\abovecaptionskip {\dp\strutbox}
  \setlength\belowcaptionskip {\ht\strutbox} 
  \setlength\dicfigheight {\heightof\dicfigbox+\fboxrule}
  \setlength\dicfigskip
    {\baselineskip*\numexpr1+\dicfigheight/\baselineskip\relax-\dicfigheight}
  \begin{figure}[ht]
    \centering\vspace*{\dicfigskip}\dicfigbox\caption{#2}\index[figures]{#3}
  \end{figure}
}


% #################################################################### COLORS #

\usepackage{color}

% http://www.colorschemer.com/online.html
% main color - grammar color
% #66023C
\definecolor {darkgreen}          {rgb} {0.40, 0.01, 0.24}

\definecolor {royalazure}         {rgb} {0.00, 0.22, 0.66}
\definecolor {brown}              {rgb} {0.40, 0.01, 0.24}

% COLORS FOR THUMB INDEXES
\definecolor {babyblueeyes}       {rgb} {0.63, 0.79, 0.95}
\definecolor {unitednationsblue}  {rgb} {0.36, 0.57, 0.90}
\definecolor {blue(ryb)}          {rgb} {0.01, 0.28, 1.00}
\definecolor {darkblue}           {rgb} {0.00, 0.00, 0.55}
\definecolor {screamingreen}      {rgb} {0.46, 1.00, 0.44}
\definecolor {limegreen}          {rgb} {0.20, 0.80, 0.20}
\definecolor {islamicgreen}       {rgb} {0.00, 0.56, 0.00}
\definecolor {upforestgreen}      {rgb} {0.00, 0.27, 0.13}
\definecolor {icterine}           {rgb} {0.99, 0.97, 0.37}
\definecolor {orange(colorwheel)} {rgb} {1.00, 0.50, 0.00}
\definecolor {orange-red}         {rgb} {1.00, 0.27, 0.00}
\definecolor {oucrimsonred}       {rgb} {0.60, 0.00, 0.00}
\definecolor {cottoncandy}        {rgb} {1.00, 0.74, 0.85}
\definecolor {orchid}             {rgb} {0.85, 0.44, 0.84}
\definecolor {vividcerise}        {rgb} {0.85, 0.11, 0.51}
\definecolor {patriarch}          {rgb} {0.50, 0.00, 0.50}

\definecolor {title}              {RGB} {  16,   13,   32}
\definecolor {golden}             {RGB} { 241,  184,   45}
\definecolor {freqstar}           {RGB} {  91,    3,   99}
\definecolor {newdev}             {RGB} {   3,   11,   99}

\definecolor {color1}             {RGB} { 182,   86,    0}
\definecolor {color2}             {RGB} { 143,    9,    6}
\definecolor {color3}             {RGB} {   3,   23,  118}
\definecolor {color4}             {RGB} {   0,   82,  168}
\definecolor {color5}             {RGB} {   0,   85,  142}
\definecolor {color6}             {RGB} {   0,  115,  162}
\definecolor {color7}             {RGB} {  34,  146,  186}
\definecolor {color8}             {RGB} {  40,  159,  153}
\definecolor {color9}             {RGB} {   0,  125,  111}
\definecolor {color10}            {RGB} {   4,  107,   60}
\definecolor {color11}            {RGB} {  71,  134,   81}
\definecolor {color12}            {RGB} { 109,  134,   42}
\definecolor {color13}            {RGB} { 205,  194,   18}
\definecolor {color14}            {RGB} { 204,  162,   24}

%\definecolor {color1}             {RGB} { 239,  222,  205}
%\definecolor {color2}             {RGB} { 251,  206,  177}
%\definecolor {color3}             {RGB} { 251,  185,  142}
%\definecolor {color4}             {RGB} { 248,  166,  112}
%\definecolor {color5}             {RGB} { 216,  231,  246}
%\definecolor {color6}             {RGB} { 186,  215,  244}
%\definecolor {color7}             {RGB} { 152,  197,  243}
%\definecolor {color8}             {RGB} { 118,  181,  245}
%\definecolor {color9}             {RGB} { 181,  249,  185}
%\definecolor {color10}            {RGB} { 160,  248,  167}
%\definecolor {color11}            {RGB} { 128,  247,  136}
%\definecolor {color12}            {RGB} {  87,  248,   98}
%\definecolor {color13}            {RGB} { 250,  241,  138}
%\definecolor {color14}            {RGB} { 250,  240,  111}
%\definecolor {color15}            {RGB} { 244,  239,   85}
%\definecolor {color16}            {RGB} { 245,  239,   57}

\definecolor {lightgrey}          {RGB} { 105,  105,  105}
\definecolor {darkgreen_real}     {rgb} {0.00, 0.50, 0.00}


% ############################################################# THUMB INDEXES #

% Thumb indexes' colors
%\newcommand\BoxColor{%
%  \ifcase\theletternum darkgreen!30\or babyblueeyes\or unitednationsblue\or blue(ryb)\or screamingreen\or limegreen\or islamicgreen\or upforestgreen\or icterine\or orange(colorwheel)\or orange-red%
%  \or oucrimsonred\or cottoncandy\or orchid\or vividcerise\or patriarch\or babyblueeyes\or unitednationsblue\or blue(ryb)\or screamingreen\or limegreen\or islamicgreen\or upforestgreen\or icterine\or orange(colorwheel)\or orange-red%
%  \or oucrimsonred\or cottoncandy\or orchid\or vividcerise\or patriarch\else darkgreen!30\fi}
  
% Thumb indexes' colors
%\newcommand\BoxColor{%
%  \ifcase\theletternum darkgreen!30\or color1\or color2\or color3\or color4\or color5\or color6\or color7\or color8\or color9\or color10%
%  \or color11\or color12\or color13\or color14\or color1\or color2\or color3\or color4\or color5\or color6\or color7\or color8\or color9%
%  \or color10\or color11\or color12\or color13\or color14\or color1\or color2\or color3\or color4\or color5\or color6\else color7\or color8\or color9\or color10\or color11\or color12\fi}

%\def\BoxColor#1{darkgreen!\the\numexpr103-#1-#1-#1\relax!blue}

\newcommand\BoxColor[1]{%
\ifcase#1 darkgreen!30\or color13\or color14\or color1\or color2\or color3\or color4\or color5\or color6\or color7\or color8\or color9\or color10%
\or color11\or color12\or color13\or color14\or color1\or color2\or color3\or color4\or color5\or color6\or color7\or color8\or color9%
\or color10\or color11\or color12\or color13\or color14\or color1\or color2\or color3\or color4\or color5\or color6\else color7\or color8\or color9\or color10\or color11\or color12\fi}

\ifshowcrop
%%%%%%%%%%%%%%%%%%%%%%%%%%%%%%%%%%%%%%%%%%%%%%%%%%%%%%%%%%%%%%%%%%%%%%%%%%%%%%%%

\newcounter{letternum}
\newcounter{lettersum}                 \setcounter{lettersum}{34}
% these margins separate top thumb from top of page and last thumb from
% bottom of page. BEWARE, it's from top of the (cut out) page, not the folio!
\newlength{\thumbtopmargin}            \setlength{\thumbtopmargin}{1cm}
\newlength{\thumbbottommargin}         \setlength{\thumbbottommargin}{1cm}
\newlength{\thumbwidth}                \setlength{\thumbwidth}{2.20cm}
\newlength{\thumbheight}
\pgfmathsetlength{\thumbheight}%
  {(\csname Gm@layoutheight\endcsname-\thumbtopmargin-\thumbbottommargin)/\value{lettersum}}
% inner xsep of letter from box margin
\newlength{\thumbsep}                  \setlength{\thumbsep}{.30cm}
% overlap over cutting line
\newlength{\thumboverlap}              \setlength{\thumboverlap}{3cm}

\tikzset{
  thumb/.style={
    text=white,
    outer sep=0pt,
    inner xsep=\thumbsep,
    minimum height=\thumbheight,
    text width=\thumbwidth-2\thumbsep,
    font=\sffamily\bfseries,},
  thumb east/.style={thumb,
    xshift=\dimexpr(+\thumboverlap-\paperwidth+\csname Gm@layoutwidth\endcsname)/2\relax,
    align=left,anchor=north east,},
  thumb west/.style={thumb,
    xshift=\dimexpr(-\thumboverlap+\paperwidth-\csname Gm@layoutwidth\endcsname)/2\relax,
    align=right,anchor=north west,},
}

\def\thumbnew{}
\def\thumbold{}
\usepackage{everypage}
\AddEverypageHook{\if\relax\thumbnew\relax\xdef\thumbnew{\thumbold}\fi}

\def\ethumbs#1,#2\relax%
  {\if\relax#1\relax\else\eventhumb{#1}\fi%
   \if\relax#2\relax\else\ethumbs#2\relax\fi%
   \gdef\thumbnew{}%
   \gdef\thumbold{#1,}}

\def\othumbs#1,#2\relax%
  {\if\relax#1\relax\else\oddthumb{#1}\fi%
   \if\relax#2\relax\else\othumbs#2\relax\fi%
   \gdef\thumbold{#1,}%
   \gdef\thumbnew{}}

\newcommand{\drawthumb}[2]{%
  \begin{tikzpicture}[remember picture, overlay]
    \node [thumb #2, fill=\BoxColor{#1}]
       at ($(current page.north #2)-%
            (0,\csname Gm@layoutvoffset\endcsname+\thumbtopmargin-\thumbheight+%
               #1*\thumbheight)$) {\csname Let#1\endcsname};
   \end{tikzpicture}}



%%%%%%%%%%%%%%%%%%%%%%%%%%%%%%%%%%%%%%%%%%%%%%%%%%%%%%%%%%%%%%%%%%%%%%%%%%%%%%%%
\else
% THUMB INDEXES
% new counter to hold the current number of the letter to determine the vertical position
\newcounter{letternum}
% newcounter for the sum of all letters to get the right height of a box
\newcounter{lettersum}
\setcounter{lettersum}{34}
% some margin settings
\newlength{\thumbtopmargin}
\setlength{\thumbtopmargin}{\ifshowcrop 3cm\else 1.5cm\fi} %1cm
\newlength{\thumbbottommargin}
\setlength{\thumbbottommargin}{\ifshowcrop 6cm\else 2cm\fi} %2.5cm
% calculate the box height by dividing the page height
\newlength{\thumbheight}
\pgfmathsetlength{\thumbheight}{%
(\paperheight-\thumbtopmargin-\thumbbottommargin)%
/%
\value{lettersum}
}
% box width
\newlength{\thumbwidth}
\setlength{\thumbwidth}{\ifshowcrop 2cm\else 0.5cm\fi} %0.5cm
% style the boxes
\tikzset{
thumb/.style={
   text=white,
   minimum height=\thumbheight,
   text width=\thumbwidth,
   outer sep=0pt,
   font=\sffamily\bfseries,
 }
 }

\def\thumbnew{}
\def\thumbold{}
\usepackage{everypage}
\AddEverypageHook{\if\relax\thumbnew\relax\xdef\thumbnew{\thumbold}\fi}

\def\ethumbs#1,#2\relax{\if\relax#1\relax\else\eventhumb{#1}\fi%
                        \if\relax#2\relax\else\ethumbs#2\relax\fi%
                        \gdef\thumbnew{}%
                        \gdef\thumbold{#1,}%
}

\def\othumbs#1,#2\relax{\if\relax#1\relax\else\oddthumb{#1}\fi%
                        \if\relax#2\relax\else\othumbs#2\relax\fi%
                        \gdef\thumbold{#1,}%
                        \gdef\thumbnew{}%
}

\newcommand{\drawthumb}[2]{%
  % see pgfmanual.pdf for more information about this part
  \begin{tikzpicture}[remember picture, overlay]
    \node [thumb, fill=\BoxColor{#1}, text centered, anchor=north #2]
       at ($(current page.north #2)-%
%            (0,\thumbtopmargin+\value{letternum}*\thumbheight)%
            (0,\thumbtopmargin+#1*\thumbheight)$) {\csname Let#1\endcsname};
   \end{tikzpicture}}
\fi

\newcommand{\oddthumb} [1]{\drawthumb{#1}{west}}
\newcommand{\eventhumb}[1]{\drawthumb{#1}{east}}

\newcommand{\lettergroup}[1]%
  {\refstepcounter{letternum}%
   \expandafter\gdef\csname Let\theletternum\endcsname{#1}%
   \xdef\thumbnew{\theletternum,\thumbnew}%
   \fancyhead[LO]{\phvfamily\bfseries\rightmark%
     \expandafter\ethumbs\thumbnew\relax\relax\relax}%
   \fancyhead[RE]{\phvfamily\bfseries\leftmark%
     \expandafter\othumbs\thumbnew\relax\relax\relax}}
     
% ############################################################## LOCALIZATION #

% renames the index name
%\addto\captionsczech{%
%  \renewcommand\indexname{{Seznam autorů fotografií a ilustrací}}}

% renames the contents name
\addto\captionsczech{%
  \renewcommand\contentsname{Obsah}}

\renewcommand \contentsname {Obsah}
\renewcommand \chaptername  {Kapitola}


% ################################################################ SECTIONING #
\usepackage[explicit]{titlesec}

\setlength{\beforetitleunit}{\baselineskip}
\setlength{\aftertitleunit} {\baselineskip}

\titlespacing*{\chapter}   {0em}{*0}{*3} % Remembere there's a topskip too
\titlespacing*{\section}   {0em}{*1}{*2}
\titlespacing*{\subsection}{0em}{*1}{*2}

\newcommand\ghostdrop[2]{%
  \raisebox{-#1\baselineskip}[0pt][0pt]{#2}}

\titleformat {\chapter} {} {} {0pt} % {#1}
  {\ghostdrop{1.0}{\parbox[t]{\columnwidth}
     {\LARGE\bfseries\centering\MakeUppercase{#1}}}}

\def\longchapterskip{\vspace{2\baselineskip}}
\def\longsectionskip{\vspace{2\baselineskip}}

\titleformat {\section} {\bfseries} {} {0pt}
  {\ghostdrop{1.0}{\parbox[t]{\columnwidth}
     {\Large\bfseries\centering#1}}}

\titleformat {\subsection} {\bfseries} {} {0pt}
  {\ghostdrop{1.0}{\hspace{1em}\parbox[t]{\columnwidth-1em}
     {\large\bfseries#1}}}


\setlength\multicolsep{0pt}
%\BeforeBeginEnvironment{multicols}{\vspace{-\dp\strutbox}}


% ######################################################## TABLE OF CONTENTS #
\usepackage{titletoc}

\titlecontents {chapter} [0pt] {\addvspace\baselineskip\bfseries}
  {} {} {\titlerule*[2\baselineskip]{}\contentspage}

\titlecontents {section} [\baselineskip] {} {} {}
  {\titlerule*[1\baselineskip]{.}\contentspage}


% #################################################### ASSORTED CUSTOM MACROS #

\def\textIS#1{\foreignlanguage{icelandic}{#1}}
\def\textCS#1{\foreignlanguage{czech}{#1}}
\def\textLA#1{\foreignlanguage{latin}{#1}}

\usepackage{accsupp}

\makeatletter                                                                                                          
  \def\ACCSUPP@bdc{\special{pdf:code \ACCSUPP@span BDC}}                                                               
  \def\ACCSUPP@emc{\special{pdf:code EMC}}                                                                             
\makeatother

\usepackage{pgffor}

\newcommand\blspace[1][1]{\vspace{#1\baselineskip}}

\usepackage{xparse}


\newcommand \dicsymFrequent  {\raisebox{-.2ex}{\color{darkgreen}\ding{167}}}
\newcommand \dicsymSee       {$\shortrightarrow$}
\newcommand \dicsymCompare   {$\shortuparrow$}
\newcommand \dicsymIdiom     {$\diamondsuit$}
\newcommand \dicsymExampleIS {$\triangleright$}
\newcommand \dicsymExampleCS {\guilsinglright}
\newcommand \dicsymProverb   {{\color{newdev}\ding{96}}}
\newcommand \dicsymPhoto
  {\includegraphics[keepaspectratio,height=0.65em]{photo_icon}}

%\NewDocumentCommand\dicTerm{mo}
 % {{\phvfamily\bfseries\textIS{#1}}\IfValueT{#2}{, \phvfamily\bfseries\textIS{#2}}}

\NewDocumentCommand\dicEntry{o}
  {\par\hangpara{0.25\baselineskip}{1}%
   \IfValueT{#1}{\markboth{#1}{#1}}%
    \ifPDFsearchvalue
    \settowidth{\eyja}{#1}%
  \makebox[\eyja]{\color{white}#1}\hspace{-\eyja}%
  \fi
   \ignorespaces}
   
   \NewDocumentCommand\dicEntryGuide{o}
  {\par\hangpara{0.29\baselineskip}{1}%
   \IfValueT{#1}{\markboth{#1}{#1}}%
    \ifPDFsearchvalue
    \settowidth{\eyja}{#1}%
  \makebox[\eyja]{\color{white}#1}\hspace{-\eyja}%
  \fi
   \ignorespaces}
   
\newcommand\dicSubEntry
  {\\\hspace*{-\hangindent}}

\NewDocumentCommand\textAlt{mm}
  {\BeginAccSupp{method=pdfstringdef,ActualText={#2}}#1\EndAccSupp{}}

\newcounter{terms}

\NewDocumentCommand\dicStyleIntTerm{m}{\textIS{{\phvfamily\bfseries#1}}}
\NewDocumentCommand\dicStyleExtTerm{m}{\textIS{\textit{#1}}}

   \NewDocumentCommand\dicTerm{mo}
  {{\closeupdotbar \phvfamily\bfseries\textIS{#1}}\IfValueT{#2}{, \phvfamily\bfseries\textIS{#2}}}

  \newcommand\dicTermList[1]
  {\foreach \t [count=\n] in {#1}{\ifnum\n=1\else, \fi\dicTerm{\t}}}
  
\def\dicIPA#1{\textipa{[#1]}}

\NewDocumentCommand\dicPos{smo}
  {{\color{darkgreen}\textbf{\small#2}%
   \IfValueT{#3}{\nobreak\textsubscript{#3}}}}

%\NewDocumentCommand\dicFlx{smo}
 % {{\color{darkgreen}\footnotesize#1\IfValueT{#2}{\textsubscript{#2}}}}

  \NewDocumentCommand\dicFlx{smo}
  {{\color{darkgreen}\footnotesize#2%
   \IfValueT{#3}{\nobreak\textsubscript{#3}}}}
  
\newcommand\dicExampleCS[1]{\dicsymExampleCS~\textit{#1}}
\newcommand\dicExampleIS[1]{\dicsymExampleIS~\textIS{\textit{#1}}}

\NewDocumentCommand\dicLink{sm}
  {\IfBooleanF#1{\dicsymSee\nobreak\hspace{.08667em plus .03333em}}\nobreak\hspace{.03667em}\dicTerm{#2}}
\NewDocumentCommand\dicSynonym{sm}
  {(\IfBooleanF#1{\dicsymSee\nobreak\hspace{.08667em plus .03333em}}\nobreak\hspace{.03667em}\textIS{\textit{#2}})}
\NewDocumentCommand\dicAntonym{sm}
  {(\IfBooleanF#1{\dicsymCompare\nobreak\hspace{.08667em plus .03333em}}\nobreak\hspace{.03667em}\textIS{\textit{#2}})}
\NewDocumentCommand\dicCompare{sm}
  {(\IfBooleanF#1{\dicsymCompare\nobreak\hspace{.08667em plus .03333em}}\nobreak\hspace{.03667em}\textIS{\textit{#2}})}


\newcommand\dicPhraseIS[1]{\textIS{\textbf{#1}}}

\NewDocumentCommand\dicIdiom{mo}
  {\dicSubEntry{\color{newdev}\dicStyleIntTerm{#1} %
     \IfValueTF{#2}{+ \dicStyleIntTerm{#2} \markboth{#1 + #2}{#1 + #2}}{\markboth{#1}{#1}}}\dicsymIdiom}
     
\newcommand\dicProverb
  {\dicSubEntry\dicsymProverb}



\newcommand\dicLangCat[1]{{\footnotesize#1}}
\newcommand\dicFieldCat[1]{{\footnotesize#1}}


\newcommand\dicDirectTranslationCS[1]{#1}
\newcommand\dicIndirectTranslationCS[1]{{\footnotesize#1}}

\usepackage{placeins}

\newcommand*{\dicLetter}[2]{%
  \FloatBarrier\newpage\phantomsection\pdfbookmark[1]{#1}{#2}%
  \noindent\parbox[b][9\baselineskip][c]{\columnwidth}
    {\centering\HUGE\strut\smash{\MakeUppercase{#1}~\MakeLowercase{#1}}}%
      %\par\xdef\tpd{\the\prevdepth}
      %\par\prevdepth\tpd%
  \expandafter\lettergroup{\MakeLowercase{#1}}}

  
%tabulky
\newcommand{\specialcell}[2][l]{%
\begin{tabular}[#1]{@{}l@{}}#2\end{tabular}}


%\newcommand{\devisionguide}[1]{\hspace*{0.3cm}{{{{\foreignlanguage{icelandic}{\color{newdev}{\phvfamily{\textbf{#1}}}}}}}}}
%\newcommand{\devisionguideindent}[1]{\hspace*{-0.1cm}{{{{\foreignlanguage{icelandic}{\color{newdev}{\phvfamily{\textbf{#1}}}}}}}}}

%\newcommand\n{$n$\nobreakdash-\hspace{0pt}} %nonbreakable dash in grammar endings
%\def \nobreakseq {\nobreak \hskip 0pt \hbox} %nolinebreak in case like (-s, -)

% rule line
\newcommand{\HRule}{\rule{\linewidth}{0.1mm}}

% width of pdf found box
\newlength\eyja

%\newcommand\strutsmash[1]{\strut\smash{#1}}


% Dictionary symbols


% Command synonyms
%\let\freqstar\dicsymFrequent
%\let\photoicon\dicsymPhoto

\newcommand{\vstretch}[1]{\vspace*{\stretch{#1}}}

\newcommand*{\addthin}{\hskip0.16667em\relax}
\newcommand*{\addthinS}{\hskip0.06667em\relax}
\newcommand*{\addthinM}{\hskip0.10007em\relax}
\newcommand*{\addthinSS}{\hskip0.00007em\relax}
% #################################################################### COVERS #

\def\dictnameCZ{islandsko-český studijní slovník}
\def\dictnameIS{íslensk-tékknesk stúdentaorðabók}

\newcommand*{\frontcoverimages}{%
  ds_image_blaberjalyng_0_2.jpg,
  ds_image_lundi_0_1.jpg,
  ds_image_hestur_0_1.jpg,
  ds_image_hrafnafifa_0_2.jpg,
  ds_image_baldursbra_0_1.jpg,
  20948.jpg,
  ds_image_hvalur_0_1.jpg,
  ds_image_blodberg_0_1.jpg}

\newcommand*{\backcoverimages}{%
  ds_image_spoi_0_1.jpg,
  ds_image_mosalyng_0_1.jpg,
  ds_image_blaklukkulyng_0_1.jpg,
  20787.jpg,
  ds_image_ufsi_0_2.jpg,
  %ds_image_gullbra_0_1.jpg,
  ds_image_vatnsberi_0_1.jpg,
  ds_image_islandsfifill_0_2.jpg,
  21137.jpg}

\newcommand\makecoverwith[1]{
 \pagecolor{title}\afterpage\nopagecolor
 \begin{center}
 \vspace*{0.3cm}%{\fill}
 {\Huge\scshape\color{white}%
   \dictnameCZ\\\dictnameIS}
 \vspace*{1em}                        % I think this is necessary
 \begin{figure}[ht]\centering%
   \setlength\fboxsep{0pt}\setlength\fboxrule{0.5pt}%
   \foreach \file in #1{
     \fbox{\includegraphics[width=5.5cm]{\file}}}
 \end{figure}
 \vspace*{\fill}
 \end{center}
}


% ###################################################### HYPERREF & PDF INFO #

% unicode neccessary so that national characters in hypersetup appear ok
\usepackage[pdftex, unicode, hyperfootnotes=false]{hyperref}

% hyperlinks in black
\makeatletter
\let\Hy@linktoc\Hy@linktoc@none
\makeatother

\hypersetup
  { pdftitle={Islandsko-český studijní slovník}
  , pdfauthor={Aleš Chejn, Jón Gíslason, Ján Zaťko}
  , pdfsubject={Islandsko-český studijní slovník}
  , pdfkeywords={Islandsko-český studijní slovník%
               , islandština, čeština, slovník}
  , bookmarks=true
  , colorlinks=true
  , citecolor=black
  , urlcolor=\ifscreen darkgreen\else black\fi
  , linkcolor=black}


% ################################################################# MICROTYPE #

\usepackage[final]{microtype}

% ############################################################### HYPHENATION #

% this allows to include brackets inside hyphenation rules f.e (pro)-dis-ku-to-vat
\lccode`\(`\(
\lccode`\)`\)
\lccode`\[`\[
\lccode`\]`\]

\begin{hyphenrules}{icelandic}
\hyphenation{
  frum-semja
  burm-neskur
  ó-breytan-le-gur
  pla-tó-nskt
  skatta-fram-tal
  skrúf-gang-ur
  strand-rauð-viður
  prent-sver-ta
  grun-nur
  uppi-skrop-pa
  ger-vi-mennska
  sö-lu-stað-ur
  sý-na-gó-ga
  mark-mið
  trau-st-ur
  ís-kyggi-leg-ur
  brjó-sta-höld
  kné-krjú-pa
  nel-li-ka
  kaffi-kvörn
  margs-konar
  skrúf-gang-ur
  uppi-skrop-pa
  veð-rað-ur
  vest-ur-átt
  mis-mun-andi
  Bar-ba-do-si
  tung-u-mál
  au-ga-brún
  há-sker-pu-sjón-varp
  rúnn-stykki
  ger-vi-tungl
  yfir-borðs-menn-ska
  gæ-fu-mun-ur
  hlí-fi-skil-di
  á-vís-un
  ó-um-flýj-an-leg-ur
  san-mar-ín-sk-ur
  gæ-ða-stað-all
}
\end{hyphenrules}
\begin{hyphenrules}{czech}
\hyphenation{
  teo-rie
  mi-liardy
  (z)-od-po-věd-nost
  tan-tié-ma
  alian-ce
  pe-dia-tr(ička)
  pa-cienty
  zá-chra-nář(ka)
  dia-man-to-vá
  poe-zie
  ma-te-riál-ní
  egoi-s-ta
  egoist-ka
  nai-vi-ta
  ne-upřím-nost
  pě-ti-úhel-ník
  še-s-ti-úhel-ník
  za-šmo-dr-chat
  ar-cheo-lo-gic-ký
  ne-omlu-vi-tel-ný
  (vy)-pro-du-ko-vat
  dia-gno-s-ti-ko-vat
  teo-lož-ka
  ateist-ka
  oceá-nu
  im-pe-ria-li-s-ta
  (vy)-zdvih-nout
  ne-úměr-ný
  kon-ti-nuál-ní
  neu-t-rum
  ne-opráv-ně-ně
  vig-vam
  pro-zaic-ký
  geo-lož-ka
  seiz-mic-ká
  teo-rém
  olean-dr
  poe-ta
  (vy)-fo-to-gra-fo-vat
  (vy)-ven-ti-lo-vat
  pneu-mo-nie
  Lu-cem-bur-čan-(ka)
  ma-te-riály
  am-bi-cióz-ní
  Myan-ma(r)
  Myan-ma-řan-(ka)
  mol-dav-šti-na
  si-tua-cí
  ne-ukáz-ně-ný
  he-te-ro-trof-ní
  ne-omlu-vi-tel-ná
  ne-upřím-ný
  dvou-bliz-ný
  ne-úmy-sl-ný
  ne-ustá-lý-mi
  ne-ohle-du-pl-ný
  ne-osob-ní
  ne-oso-le-ný
  ne-uvě-ři-tel-ný
  ne-ohro-že-ný
  ne-oby-čej-ný
  ne-útul-ný
  ne-únav-ný
  ne-ože-nit
  po-ukáz-ka
  za-úto-čit
  aero-dy-na-mic-ký
  hru-bián-ství
  (pro)-di-s-ku-to-vat
  (vy)-ča-lou-nit
  (na)-oč-ko-vá-ní
  Bhú-tán-ka
  (za)-šně-ro-vat
  soud-ní
  re-vo-lu-cio-nář(ka)
  á-zer-báj-d-žán
  čtyř-hran-ný
  ze-m-dle-lý
  klou-zač-ka
  Vest-mann-ské
  tří-stran-ný
  (za)-bu-rá-ce-ní
  os-tud-ný
  pro-du-cent-(ka)
  chvá-s-tal-(ka)
  apar-theid
  asis-ten-ce
  asis-tent-(ka)
  (u)-či-nit
  (u)-za-mknout
  past-vi-na
  (za)-zá-lo-ho-vat
  Af-ri-čan-(ka)
  (za)-an-ga-žo-vá-ní
  se-rióz-ní
  opo-nent-(ka)
  pro-tiv-ník
  zvlášt-ní
  (po)-div-ný
  ho-mo-se-xuál-ní
  ko-alič-ní
  si-tua-ce
  si-tua-cí
  egoi-s-mus
  spe-ciál-ní
  spe-cia-li-zo-va-ný
  spe-cia-li-zo-va-né
  sobě-stač-ný
  sa-mo-zvo-le-ný
  sa-mo-čin-ně
 chvá-s-ta-vost
 klient-(ka)
 ku-rióz-ní
 (na)-šrou-bo-vat
 (za)-šrou-bo-vat
 krea-ti-vi-ta
 pneu-ma-ti-ka
 toa-le-ta
 so-cia-list-ka
 (vy)-dez-in-fi-ko-va-ný
 vše-uměl
 vše-uměl-(kyně)
 (pře)-škrt-nout
 (vy)-svlék-nout
 vy-učo-va-cí
 vy-učo-va-cích
 (vy)-chla-s-tat
 (vy)-te-to-vat
 pa-cient-(ka)
 ne-opod-s-tat-ně-ný
 de-se-ti-úhel-ník
 pe-rio-di-kum
 (za)-mum-lá-ní
 ne-smě-lá
 ne-dů-vě-ři-vá
 tee-na-ger-(ka)
 (z)-de-ge-ne-ro-vá-ní
 mo-ci-chti-vý
 bez-úhon-nost
 zne-uctít
 vo-dě-odol-ný
 we-bo-vý
 web-de-sign
 in-že-nýr-(ka)
 zá-pa-do-is-land-ský
 klient-(ka)
 za-uja-tý
 vy-uží-va-né
 vy-uží-vat
 (za)-re-gi-s-t-ro-vat
 (z)-kom-pri-mo-va-ný
 na-cio-na-li-s-mus
 ma-k-ro-eko-no-mic-ká
 bez-bo-lest-ný
 viol-ka
 troj-úhel-ník
 (za)-bu-rá-ce-ní
 ne-ústup-nost
 (vy)-cvi-če-ný
 zprá-va
 ge-ria-tr-(ička)
 osud-ný
 ná-pra-vy
 vstup-ní
 obec-ně
 re-pub-li-ka
 (po)-div-ný
 pa-slun-ce
 vliv-ný
 ba-ham-ský
 pe-dia-tr-(ička)
 ha-sič-(ka)
 vnuč-ka
 bro-s-kvoň
 lo-di-vod-(ka)
 mi-s-tr-(yně)
 in-že-nýr-(ka)
 dr-ži-tel-(ka)
 to-pič-(ka)
 pro-mi-nent-ní
 takt-ní
 frag-men-tár-ní
 po-zo-ru-hod-ný
 úspě-š-ný
 trans-fu-ze
 bur-žoa-zie
 pře-stro-je-ní
 oh-ni-vzdor-ný
 dis-tan-ční
 de-frau-da-ce
 fi-s-kál-ní
 zběh-lý
 kvap-ný
 střed-ní
 sa-mo-hlá-s-ka
 pře-vo-di-tel-ný
 zban-k-ro-to-va-ný
 zá-clo-na
 zná-mý
 zá-klad-ní
 na-uka
 ro-š-ťa-č-ka
 lí-bán-ky
 sou-du
 zkou-š-ka
 mlý-nek
 prá-šiv-ka
 kul-ma
 la-hod-ný
 zá-meč-ni-ce
 ne-vý-znam-ný
 ovzdu-ší
 mno-ho-úhel-ní-ko-vý
 žra-vý
 zcest-ný
 esen-ce
 (při)-nu-ce-ný
 skr-b-lic-ký
 skle-s-lý
 bez-bož-ný
 ne-opráv-ně-ný
 ošu-mě-lý
 pru-der-ní
 chvá-s-ta-vý
 sou-drž-ný
 van-dal-(ka)
 sa-bo-tér-(ka)
 (po)-div-ně
 mo-k-řad
 plyt-ký
 na-spo-du
 izo-la-ce
 le-toun
 za-han-be-ný
 ří-ze-ní
 kom-plex-ní
 pr-kén-ko
 ba-vič-(ka)
 plyt-ký
 ne-mo-tor-ně
 šéf-(ka)
 přes-čas
 zji-s-ti-tel-ný
 ta-ko-vý
 kňou-ral-(ka)
 se-me-s-tr
 po-pra-s-ka-ný
 han-ba
 vzpří-me-ný
 chys-ta-ný
 lik-vi-da-ce
 or-ga-ni-za-ce
 ri-s-kant-ní
 (za)-pís-kat
 (za)-hvíz-dat
 vrou-cí
  }
\end{hyphenrules}


% ############################################################# CZECH TYPOGRAPHY #
% https://math.feld.cvut.cz/tkadlec/ftp/text/textypo.pdf



% ################################################################# DOCUMENT #

\begin{document}
\pagenumbering{gobble}

% -------------------------------------------------------------- Front cover -

\ifmakecovers
\covergeometry
  \makecoverwith\frontcoverimages
\restoregeometry  
\fi


% ============================================================= FRONT MATTER =

\ifmakefrontmatter

% ------------------------------------------------------------------ Patitul -

\cleardoublepage

\begin{center}
  %\vspace*{2cm}{\fontsize{20}{20}\selectfont\scshape\dictnameCZ\\\dictnameIS}
  \vspace*{\stretch 1}
  {\Huge\scshape\dictnameCZ\\\dictnameIS}
  \vspace{\stretch 4}
\end{center}

% -------------------------------------------------------------------- Titul -

\cleardoublepage

\begin{center}
  %\vspace*{2cm}{\fontsize{24}{24}\selectfont\scshape\dictnameCZ\\\dictnameIS}
  \vspace*{\stretch 1}
  {\HUGE\scshape\dictnameCZ\\\dictnameIS}
  \vspace{\stretch 2}
\end{center}

% ----------------------------------------------------------------- Impresum -

\clearpage

\noindent\begin{minipage}[t][\textheight][t]{\textwidth}
Autoři \textit{Islandsko-českého studijního slovníku}:\\
  \textbf{Aleš Chejn, Ján Zaťko, Jón Gíslason}\\[\baselineskip]
\ifPDF
\else
Autor \textit{Islandské gramatiky}:\\
  \textbf{Vojtěch Kupča}\\[\baselineskip]
\fi
Poděkování:\\
  \textbf{Dorotě Nierychlewské-Chejn, Jiřině Hanzlové, Jiřině Chejnové,
          Simoně Petrásové,\\ Petru Mikešovi, Vojtěchu Kupčovi, 
          Renatě Peškové Emilsson, Amiru Mulahumici, \\ Kristýně Antonové, 
          Martině Kašparové, Lucii Brabcové, Samoriele, Adamu Kožouškovi}\\[\baselineskip]
Finanční příspěvek na výtisk exemplářů určených knihovnám\\ v České republice,
Slovenské republice a na Islandu poskytli:\\
  \textbf{Petr Mikeš, Dorota Nierychlewska-Chejn, Jiřina Chejnová, Otta Chejn}
  \ifPDF
  \\[39\baselineskip]
  \else
  \\[36\baselineskip]
  \fi
\textbf{Hvalur.org - 2016}\\
Licence: \textbf{GNU Free Documentation License}\\
Sazba: \textbf{Aleš Chejn} v prostředí  \textbf\LaTeX\\
  Zvláštní poděkování patří \textbf{Paolu al-Imk\=an\=\i\ Brasolin} a uživatelům  \textbf{\TeX\ - \LaTeX\ Stack Exchange}\\[\baselineskip]
\ifPDF
Verze: \textbf{PDF verze} \textit{Islandsko-českého studijního slovníku} vhodná pro \textbf{\ifscreen offline prohlížení a vyhledávání\else vlastní tisk\fi}\\ volně dostupná na \url{www.hvalur.org}\\[\baselineskip]
\else
\textbf\ISBN\\[\baselineskip]
\fi
\ifPDF
\else
Tisk: \textbf{Tiskárna Melmen, Pardubice}\\
\fi
Heslová část vysázena patkovým fontem \TeX\ Gyre Pagella,
hesla bezpatkovým fontem {\phvfamily Helvetica}.
\end{minipage}

% -------------------------------------------------------- Table of contents -

\cleardoublepage 

\pdfbookmark\contentsname{toc}    
\tableofcontents
\thispagestyle{empty}
\cleardoublepage 

\fi

\pagenumbering{arabic}
\setcounter{page}{3}
% ============================================================= INTRODUCTION =

\ifinputintroduction

% ------------------------------------- The study Icelandic-Czech dictionary -

\chapter{O Islandsko-českém studijním slovníku}
Cílem vytvoření\textit{ Islandsko-českého studijního slovníku} je zpřístupnění islandského jazyka všem zájemcům o~tento jazyk. 

\textit{Islandsko-český studijní slovník} je určen především českým a~slovenským studentům islandštiny, začátečníkům i~pokročilým, kteří zde naleznou ucelené fonetické, morfologické a~syntaktické informace, jakož
i~slovní zásobu potřebnou k~aktivnímu i~pasivnímu užití s~rozsahem vhodným ke studiu islandštiny na vysoké škole.
Slovník je rovněž určen všem zájemcům o~islandštinu a~Island, jimž slovník nabízí encyklopedické informace, úvod do fonetiky a~gramatiky a~velké množství fotografií a~ilustrací. 
Třebaže se jedná o~slovník primárně zaměřený na potřeby českých a~slovenských uživatelů, obsahuje také řadu informací, které usnadní práci islandskému uživateli (např. použití více jak 10 000 synonym a~antonym, překlady příkladů ap.).

Slovní zásoba původně vychází z~\textit{Concise Icelandic-English Dictionary} (\cite {ic_en}), značně však byla rozšířena díky přístupu k~webovému projektu \textit{ISLEX} (\cite {int1}). Velké množství hesel bylo
přidáno na podnět uživatelů online verze \textit{Islandsko-českého studijního slovníku}.

Slovník obsahuje 30~575 hesel s~46~725 významy nebo slovními spojeními. Z~celkového počtu 30~575 je 28~097 svébytných hesel obsahujících alespoň jeden význam nebo slovní spojení. 
Zbylá hesla (v~počtu 2~478) jsou varianty, dublety nebo nepravidelné tvary.

Slovník je primárně online elektronický slovník, volně přístupný na internetových stránkách \url{www.hvalur.org}  (\cite {int14}). 
Tištěná verze slovníku je publikována na vlastní náklady v~počtu 30 kusů, přičemž většina exemplářů je určena knihovnám a~institucím v~České republice, Slovenské republice a~na Islandu
a~několik exemplářů je určeno autorům nebo zájemcům z~řad veřejnosti. Dvě PDF verze (vyhledávání\,/\addthin prohlížení offline a~vlastní tisk) jsou dostupné zdarma na internetových stránkách slovníku (viz výše).

\textit{Islandsko-český studijní slovník} je otevřený svobodný projekt, jehož se v~průběhu tvorby zúčastnilo velké množství spoluautorů (viz Historie tvorby IČSS), kterým tímto jménem autorů velmi děkuji.
Děkuji také všem uživatelům online verze slovníku za trpělivost, dobrou vůli a~veškeré připomínky a~návrhy, jmenovitě především Petru Mikešovi, Adamu Kožouškovi a~Samoriele za neutuchající víru v~náš projekt.
Poděkování patří také našim rodinám a~blízkým, díky kterým jsme se mohli soustavně věnovat tvorbě slovníku.
Srdečné poděkování patří mé ženě Dorotce za podporu a~nadšení potřebné k~tvorbě slovníku.

\blspace[5]

{\centering Aleš Chejn, Brwinów, 24. prosince 2015\par}


% ------------ The history of the formation Icelandic-Czech dictionary study -

\chapter{Historie tvorby Islandsko-českého studijního slovníku}
\longchapterskip
Projekt \textit{Islandsko-českého studijního slovníku} (IČSS) vznikl v~roce 2001 na nepřímý podnět doc. PhDr. Heleny Kadečkové, CSc. z~Univerzity Karlovy,
kdy byl vypracován seznam slov. 
Tento seznam obsahoval islandská slova a~jejich překlady, později byly přidány základní morfologické informace o~slovním druhu a~deklinační\,/\addthin konjugační koncovky podle vzoru z~\textit{Concise Icelandic-English Dictionary} (\cite {ic_en}). 
Seznam v~této době obsahoval okolo 1 200 slov. Byl zveřejněn pomocí statických webových stránek a~později byl umístěn na webu Univerzity Karlovy. 

V~roce 2006 byl slovník rozšířen o~nových 1 500 heslových slov a~stávající hesla byla rozšířena o~nové významy, zvláště významy frázových sloves. 
Ke konci roku 2006  byl slovník převeden z programu Excel do databáze MySQL a~byly zahájeny práce na vytvoření aplikace \textit{Dictionary System} (DS) (\cite {int19}), která by podporovala tvorbu IČSS. 
Obsah slovníku byl zveřejněn pod svobodnou licencí na webových stránkách \url{www.hvalur.org} (\cite {int14}), kde se nachází dodnes. 

V~roce 2007 byla vytvořena první verze aplikace DS. V~této době se k~projektu přidala Renata Pešková Emilsson, která se nyní věnuje podpoře výuky češtiny na Islandu. 
Ta slovník obohatila rozsáhlým seznamem slov z~praktického života na Islandu. Dále se slovník rozšiřoval o~heslová slova podle \textit{Concise Icelandic-English Dictionary} (\cite {ic_en}). 
V průběhu následujících třech let bylo doplněno přibližně 15~000 slov. V této době byla hesla rozšířena o~synonyma, příklady, frekvenci, oborovou a~stylovou charakteristiku. K projektu se na Islandu připojila na krátkou dobu Kristýna Antonová, studentka islandštiny na Islandské univerzitě, která 
nás upozornila na nesrovnalosti v~uživatelském rozhraní aplikace DS.

Během léta 2009 byl vytvořen skript na generování skloňování a~časování heslových slov. Výsledky generování byly porovnány s~\textit{Beygingarlýsing íslensks nútímamáls} (\cite {DalvikVM}) a~během následujících let konzultovány
s~Jónem Gíslasonem. Kristín Bjarnadóttir, vedoucí projektu \textit{Beygingarlýsing íslensks nútímamáls}, zodpověděla řadu našich dotazů týkajících se deklinací a~konjugací.

V~roce 2009 začala spolupráce s~\textit{Biolib.cz} (\cite {int5}), jmenovitě s~Ondřejem Zichou. Obdrželi jsme  databázi s~náhledy obrázků, autory a~licencemi a~latinskými názvy botanických a~zoologických druhů.

V~zimním semestru 2009 jsme požádali s~Amirem Mulahumicem, bosenským studentem islandštiny na Islandské univerzitě, Studentský fond Islandské univerzity o~grant na tvorbu zvukové databáze islandských slov. 
Studentský fond návrh přijal. Peníze z~grantu jsme určili pro islandského rodilého mluvčího, Jóna Gíslasona, lektora islandského jazyka na Islandské univerzitě v~Reykjavíku, který nahrál více jak 22~000 islandských slov do elektronické formy. 
S~Dorotou Nierychlewskou-Chejn jsme připravili dva seznamy heslových slov (jeden pro nahrávajícího a~druhý pro zpracování nahrávek). Na podporu nahrávek jsme vytvořili webové stránky nazvané \textit{Databáze s~výslovností islandských slov} (\cite {int13}) a~tyto stránky jsme přeložili do 6~jazyků (Jón Gíslason do islandštiny, Ján Zaťko do slovenštiny, francouzštiny a~angličtiny, Amir Mulamuhic do srbochorvatštiny a~Dorota Nierychlewska-Chejn do polštiny).

V~roce 2010 jsme přidali do databáze pravidla islandské výslovnosti, která jsou uvedena v~knize \textit{Handbók um íslenskan framburð} (\cite {is2}). Z~těchto pravidel jsme automaticky vygenerovali fonetický zápis výslovnosti všech hesel ve slovníku. 
Jón Gíslason ověřil správnost fonetického zápisu u~všech heslových slov.

Na konci roku 2011 jsme spojili webové stránky slovníku se stránkami aplikace a~sjednotili jsme a~vylepšili vzhled webových stránek. Cennými připomínkami se na vytvoření aplikace podílela
RNDr. Ingrid Nagyová, PhD., vysokoškolská pedagožka Ostravské Univerzity. Ján Zaťko přeložil nové webové stránky, Průvodce po slovníku do slovenštiny, francouzštiny a~angličtiny a~celý manuál aplikace do angličtiny. 
Dorota Nierychlewska-Chejn přeložila stránky do polštiny, Jón Gíslason do islandštiny a~Marcos Helena do portugalštiny.

Od prosince 2011 do března 2014 jsme realizovali první čtení, ve kterém jsme doplnili chybějící informace, přidali nová hesla a~opravili pravopisné chyby. 
Obsah slovníku jsme porovnávali především s~\textit{Íslensk orðabók} (\mbox{\cite {is}}) a~\textit{Islandsko-českým slovníkem} (\cite {is7}). S~Vojtěchem Kupčou jsme konzultovali mnohé nesnáze při překladech.

V~letech 2013 až 2015 byl vytvořen morfologický klíč, který rozšiřuje klasický systém zobrazení morfologické informace v~islandských slovnících o~přesné zařazení daného hesla do určité deklinační\,/\addthin konjugační skupiny, jejíž
skloňování\,/\addthin časování je ukázáno na příkladu. V~roce 2014 jsme obdrželi roční grant od Ostravské univerzity v~Ostravě na tvorbu morfologického klíče. S~realizací grantu nám velmi pomohla doc. PhDr. Diana Svobodová, Ph.D. a~Ing. Zuzana Havrlantová.

Od března 2014 do srpna 2015 jsme realizovali druhé finální čtení, ve kterém jsme porovnávali informace s~webovým slovníkem \textit{ISLEX} (\cite {int1}). 
Na požádání jsme obdrželi od Þórdís Úlfarsdóttir (vedoucí projektu \textit{ISLEX}) seznam víceslovných příslovečných frází. Takto poskytnutý seznam nám usnadnil opravy slovníku. 
S~Þórdís Úlfarsdóttir jsme konzultovali různé gramatické jevy a~záležitosti týkající se deklinací\,/\addthin konjugací.
Martina Kašparová, studentka islandštiny na Islandské univerzitě, prošla celý slovník, přičemž se soustředila mimo jiné na stylistické a~pravopisné korektury překladů. 
Lucie Brabcová udělala jazykové a~typografické korektury celého slovníku.

Vojtěch Kupča vydal v~roce 2014 \textit{Islandskou gramatiku} (\cite {is77}) a~souhlasil s~připojením této publikace ke~knižní verzi IČSS.
V~roce 2015 jsme s~přispěním Paola al-Imk\=an\=\i\ Brasolin vysázeli slovník v~typografickém prostředí {\LaTeX}. V~únoru 2016 byl slovník vytištěn v~Tiskárně MELMEN, s.r.o v~Pardubicích a svázán 
v~Knihařství Vojtíšková Hrubá v~Pardubicích. 

\blspace[5]

{\centering Všem spoluautorům děkujeme za spolupráci.\par}


% -------------------------------------------------  Guide to the dictionary -

\chapter{Průvodce po slovníku}
\begin{multicols}{2}
\def\tableA{\noindent\begin{tabularx}{\columnwidth}{>{\bfseries}lll}
  \toprule
  \textit{přípona} & \textit{přípona s~koncovkou} & \textit{příklad} \\
  \midrule
  -ing-   & -ingur   & rugl··ingur \\
          & -ing     & hreyf··ing \\
          & -ingi    & erf··ingi \\
          & -ingar   & lækn··ingar \\
  -ling-  & -lingur  & vett··lingur \\
  -ning-  & -ningur  & flut··ningur \\
          & -ning    & set··ning \\
          & -ningar  & kos··ningar \\
  -að-    & -aður    & mark··aður \\
  -nað-   & -naður   & skil··naður \\
  -ar-    & -ari     & dóm··ari \\
  -skap-  & -skapur  & kveð··skapur \\
  -dóm-   & -dómur   & lær··dómur \\
  -leik-  & -leikur  & sann··leikur \\
          & -leiki   & mögu··leiki \\
  -lát-   & -látur   & laus··látur \\
  -átt-   & -átta    & víð··átta \\
  -sam-   & -samur   & starf··samur \\
  -sem-   & -semi    & hamingju··semi \\
  -nesk-  & -neskur  & tékk··neskur \\
          & -neska   & tékk··neska \\
  -leg-   & -legur   & blíð··legur \\
          & -lega    & blíð··lega \\
  -ung-   & -ungur   & gler··ungur \\
          & -unga    & snopp··unga \\
          & -ungi    & sveit··ungi \\
          & -ung     & nýj··ung \\
  -ótt-   & -óttur   & blett··óttur \\
  -ræn-   & -rænn    & suð··rænn \\
          & -ræna    & nor··ræna \\
  -ern-   & -erni    & fað··erni \\
  -ind-   & -indi    & vís··indi \\
  -and-   & -andi    & gef··andi \\
  -ug-    & -ugur    & göf··ugur \\
          & -ugt     & stöð··ugt \\
  -ust-   & -usta    & þjón··usta \\
  -ald-   & -aldur   & far··aldur \\
  -uð-    & -uður    & söfn··uður \\
  -úð-    & -úð      & dul··úð \\
  -versk- & -verskur & ung··verskur \\
          & -verska  & ung··verska \\
  -ist-   & -isti    & hass··isti \\
  -ism-   & -ismi    & húman··ismi \\
  -ang-   & -angur   & far··angur \\
  \bottomrule
  \end{tabularx}
  \captionof{table}{Seznam přípon označených ve slovníku}
  \label{table1}}

\def\tableB{\noindent\begin{tabularx}{\columnwidth}{>{\bfseries}lX}
  \toprule
  \textbf{symbol} & \textbf{vysvětlivka} \\
  \midrule
  m                   & označuje slovní druh \\
  \textsubscript{1}   & označuje třídu \\
  (-a, -ar)           & označuje deklinační / konjugační koncovky hesla \\
  \textsubscript{8}   & označuje podtřídu\\
  \bottomrule
  \end{tabularx}
  \captionof{table}{Morfologická informace}
  \label{table2}}

\def\tableC{\noindent\begin{tabularx}{\columnwidth}
    {>{\bfseries}l>{\footnotesize\itshape}lllll}
  \toprule
  \multirow{3}{*}{{{\textbf{f{\textsubscript{1}}} \Large{\textbf{7}}}}}
    & nom & kisa & kisan    & kisur & kisurnar \\
    & acc & kisu & kisuna   & kisur & kisurnar \\
    & dat & kisu & kisunni  & kisum & kisunum \\
    & gen & kisu & kisunnar & kisa  & kisanna \\
  \bottomrule
  \end{tabularx}
  \captionof{table}{Skloňování hesla \textit{kisa}}
  \label{table3}}

\def\tableD{\noindent\begin{tabularx}{\columnwidth}{>{\bfseries}lllll}
  \toprule
  \textit{zkratka} & \textit{slovo} & \textit{pád} & \textit{číslo}
    & \specialcell{\textit{předmět}\\\textit{podmět}} \\
  \midrule
  e-að  & eitthvað   & 1. & sg & neživotný \\
  e-ð   & eitthvað   & 4. & sg & neživotný \\
  e-ir  & einhverjir & 1. & pl & životný \\
  e-ja  & einhverja  & 4. & pl & životný \\
  e-jum & einhverjum & 3. & pl & životný \\
  e-m   & einhverjum & 3. & sg & životný \\
  e-n   & einhvern   & 4. & sg & životný \\
  e-r   & einhver    & 1. & sg & životný \\
  e-rra & einhverra  & 2. & pl & životný \\
  e-rs  & einhvers   & 2. & sg & životný \\
  e-s   & einhvers  & 2. & sg & neživotný \\
  e-u   & einhverju  & 3. & sg & neživotný \\
  \bottomrule
  \end{tabularx}
  \captionof{table}{Seznam použitých zkratek neurčitých zájmen}
  \label{table4}}

%#############################################################################

\section{Heslo}

\subsection*{Řazení slov}

Heslová slova jsou zobrazena tučně a~jsou seřazena podle islandské abecedy:

\blspace
{\bfseries\centering
  a~á b c d ð e é f g h i~í j k~l m n\\
  o~ó p q r s~t u~ú v~w x y ý z~þ æ ö\par}
\blspace

K~snadnější orientaci v~heslové části je použit u~každého písmene tzv. palcový index, který se nachází na hřbetě slovníku. Každé písmeno je v~palcovém indexu barevně odlišeno.

\subsection*{Homonyma}

Homonyma (heslová slova stejného reprezentativního tvaru) jsou označena povýšenými arabskými číslicemi uváděnými za heslovým slovem.

\blspace
  \dicEntry \dicTerm{vor\textsuperscript{1}} \dicPos{n} \dicFlx{(-s, -)}
  \dicEntry \dicTerm{vor\textsuperscript{2}} \dicPos{pron} \dicFlx{pers gen pl}
\blspace

Homonyma jsou seřazena podle českého ustáleného řazení slovních druhů. Poté následují předpony a~zkratky.
V~případě dvou podstatných jmen je pořadí mužský rod, ženský rod, střední rod. V~případě dvou podstatných jmen stejného rodu má přednost slovo se slabým skloňováním před slovem se silným skloňováním. V~případě dvou sloves stejného reprezentativního tvaru má přednost sloveso se slabým časováním před slovesem se silným časováním.

\subsection*{Varianty a~dublety}

U~některých heslových slov jsou v~záhlaví hesla uvedeny a) varianty (obměny)

\blspace
  \dicEntry \dicTerm{ærsla··fenginn}, \dicTerm{ærslafullur} \dicPos{adj}
\blspace

nebo b) dublety (dvojtvary).

\blspace
  \dicEntry \dicTerm{{ævin··týri}, \dicTerm{ævintýr}} \dicPos{n}
\blspace

Varianty a~dublety jsou méně frekventované než samotné heslo. Lze je vyhledat ve slovníku jako hesla obsahující fonetické a~morfologické informace s~odkazem na původní heslové slovo.

\blspace
  \dicEntry \dicTerm{ærsla··|fullur} \dicIPA{{a}{i}{\textsubring{r}}{s}{\textsubring{d}}{l}{a}{f}{\textscy}{\textsubring{d}}{l}{\textscy}{\textsubring{r}}} \dicPos{adj} \dicFlx{(comp -fyllri, sup -fyllstur)} \dicLink{ærslafenginn}
  \dicEntry \dicTerm{ævin··týr} \dicIPA{{a}{i}{\textlengthmark}{v}{\textsci}{n}{t\textsuperscript{h}}{i}{\textsubring{r}}} \dicPos{n} \dicFlx{(-s, -)} \dicLink{ævintýri}
\blspace

\subsection*{Vybrané tvary jako hesla}

Ve slovníku jsou uvedeny vybrané tvary heslových slov, u~kterých by mohl mít uživatel potíže s~identifikací. Jedná se o~a) genitiv jednotného čísla ({\textbf{sg}}) a~nominativ množného čísla ({\textbf{pl}}) podstatných jmen, b) ženský a~střední rod přídavných jmen a~2. a~3. stupeň rodu mužského přídavných jmen, c) oznamovací způsob ({\textbf{ind}}) přítomného času ({\textbf{praes}}) a~minulého času ({\textbf{pf}}) 1. osoby ({\textbf{pers}}) sg a~pl a~spojovací způsob ({\textbf{con}}) minulého času 1. osoby sg a~jiné nepravidelné tvary.

\blspace
  \dicEntry \dicTerm{öttum} \dicPos{v} \dicFlx{ind pf pl 1 pers} \dicLink{etja}
\blspace

\subsection*{Dělení heslových slov}

K~dělení heslových slov je použita střední tečka \textbf{'·'}, která odděluje části složeného slova, přípony a~předpony. Platí následující pravidla:

\renewcommand{\labelenumi}{\alph{enumi})}
\begin{enumerate}
\item jediné dělení heslového slova je označeno dvěma středními tečkami \textbf{'··'}. Může se jednat o~dělení částí složeného slova
\blspace
\dicEntry \dicTerm{ráð··hús}
\blspace

nebo o~dělení mezi kořenem a~předponou / příponou.

\blspace
\dicEntry \dicTerm{þjón··usta}
\blspace
\item V~případě, že se v~hesle vyskytuje více dělení, jsou použity dvě střední tečky \textbf{'··'} pro významově poslední dělení a~jedna střední tečka \textbf{'·'} pro zbylé / zbylá dělení.
\blspace
\dicEntry \dicTerm{við·bótar··líf·eyris·sparn·aður}
\blspace
\item V~hesle s~předponou je označena jedna nebo i~více předpon (každá zvlášť).
\blspace
\dicEntry \dicTerm{ó··að·finnan·legur}
\blspace
\item V~hesle s~příponou je označena jedna nebo i~více přípon (každá zvlášť).
\blspace
\dicEntry \dicTerm{létt·úð··ugur}
\blspace
\item V~hesle s~příponou je označena pouze přípona / přípony poslední části složeného slova. Důvodem pro neoznačení přípon jiných částí složených slov je lepší čitelnost hesla. Příklad označení pouze přípony poslední části složeniny:
\blspace
\dicEntry \dicTerm{flutninga··miðl·ari}
\end{enumerate}

Následuje seznam přípon, které jsou ve slovníku označeny:

\tableA

\section{Frekvence slov}

Symbolem \dicsymFrequent\ je označeno 2~825 hesel patřících mezi velmi frekventovaná islandská slova. Tento seznam vychází z~korpusu \textit{Íslenskur Orðasjóður} (\cite{int7}) a~knihy \textit{Íslensk orðtíðnibók} (\cite{is2}).

\blspace
  \dicEntry \dicTerm{koma}\textsuperscript{2} \dicsymFrequent
\blspace

\section{Fonetický zápis}

V~hranatých závorkách v~záhlaví hesla je uveden fonetický zápis výslovnosti hesla v~mezinárodní fonetické abecedě (IPA, angl. International Phonetic Alphabet).

\blspace
  \dicEntry \dicTerm{öldu··rót} \dicIPA{{\oe}{l}{\textsubring{d}}{\textscy}{r}{ou}{\textsubring{d}}}
\blspace

Znak '{\textlengthmark}' označuje dlouhou samohlásku nebo souhlásku.

\blspace
  \dicEntry \dicTerm{ör}\textsuperscript{2} \dicIPA{{\oe}{\textlengthmark}{\textsubring{r}}}
\blspace

\section{Slovní druhy, předpony a~zkratky}

V~záhlaví hesla je uvedena morfologická informace popisující a) slovní druh hesla, b) předponu nebo c) zkratku. Pokud se heslo řadí do více slovních druhů, jsou jednotlivé zkratky pro slovní druhy odděleny lomítkem.

\blspace
  \dicEntry \dicTerm{af} \dicPos{prep / adv}
\blspace

Informace o~slovním druhu, předponě či zkratce je vysázena barevně a~tučně.

\subsection*{Podstatná jména}

Podstatná jména jsou uvedena v~1. pádě (\textbf{nom}) jednotného čísla (\textbf{sg}) nebo v~1. pádě množného čísla (\textbf{pl}), pokud se heslo v~jednotném čísle nevyskytuje. Podstatná jména jsou označena následovně: \textbf{m} (l. masculinum) - podstatné jméno rodu mužského,
\textbf{f} (l. femininum) - podstatné jméno rodu ženského, \textbf{n} (l. neutrum) - podstatné jméno rodu středního, \textbf{subs} (l. substantivum) - podstatné jméno bez jednoznačného mluvnického rodu.

\blspace
  \dicEntry \dicTerm{hestur} \dicPos{m}
  \dicEntry \dicTerm{Hanoí} \dicPos{subs}
\blspace

\subsection*{Přídavná jména}

Přídavná jména jsou uvedena v~1. pádě (\textbf{nom}) rodu mužského (\textbf{m}) jednotného čísla (\textbf{sg}) 1. stupně (\textbf{pos}) (l. positivus) a~jsou označena zkratkou \textbf{adj} (l. adjectivum).

\blspace
  \dicEntry \dicTerm{fallegur} \dicPos{adj}
\blspace

V~případě, že u~přídavného jména dochází v~rodě ženském (\textbf{f}) nebo středním (\textbf{n}) k~přehlásce anebo je tvar druhého stupně (\textbf{comp} (l. comparativus)) či třetího stupně (\textbf{sup} (l. superlativus)) nepravidelný, pak jsou tyto tvary uvedeny v~nominativu jako heslová slova s~odkazem na hlavní heslový tvar

\blspace
  \dicEntry \dicTerm{yngri} \dicPos{adj} \dicFlx{comp m} \dicLink{ungur}
\blspace

nebo jako svébytná hesla s~významy.

\blspace
  \dicEntry \dicTerm{fremri} \dicPos{adj} \dicFlx{comp (sup fremstur)}
\blspace

\subsection*{Zájmena}

Zájmena jsou označena zkratkou \textbf{pron} (l. pronomen) a~jsou tříděna na zájmena ukazovací \textbf{dem} (angl. demonstrative), přivlastňovací \textbf{poss} (angl. possessive), osobní \textbf{pers} (angl. personal), neurčitá \textbf{indef} (angl. indefinite), zvratná \textbf{refl} (angl. reflexive) a~vztažná \textbf{rel} (angl. relative).

\blspace
  \dicEntry \dicTerm{þessi} \dicPos{pron} \dicFlx{dem}
\blspace

U~zájmen přivlastňovacích je v~3. osobě označen rod.

\blspace
  \dicEntry \dicTerm{hans} \dicPos{pron} \dicFlx{poss m}
\blspace

U~zájmen osobních je označen pád a~číslo.

\blspace
  \dicEntry \dicTerm{mig} \dicPos{pron} \dicFlx{pers acc sg}
\blspace

\subsection*{Číslovky}

Číslovky jsou označeny zkratkou \textbf{num} (l. numerale). V~případě vybraných číslovek (1 až 4) jsou ve slovníku uvedeny jako heslová slova také tvary ženského i~středního rodu. V~takovém případě jsou označeny zkratkou \textbf{m} pro mužský rod, zkratkou \textbf{f} pro ženský rod a~zkratkou \textbf{n} pro střední rod. Zkratkou \textbf{ord} (angl. ordinal number) jsou označeny číslovky řadové.

\blspace
  \dicEntry \dicTerm{tut··tugu} \dicPos{num}
  \dicEntry \dicTerm{tveir} \dicPos{num} \dicFlx{pl m}
  \dicEntry \dicTerm{fyrsti} \dicPos{num} \dicFlx{ord}
\blspace

\subsection*{Slovesa}

Slovesa jsou uvedena v~infinitivu a~jsou označena zkratkou \textbf{v} (l. verbum).

\blspace
  \dicEntry \dicTerm{fara} \dicPos{v}
\blspace

Slovesa ve středním rodě jsou označena jako \textbf{refl} (angl. reflexive).

\blspace
  \dicEntry \dicTerm{nálgast} \dicPos{v} \dicFlx{refl}
\blspace

Slovesa, která jsou vždy neosobní, jsou označena \textbf{impers} (angl. impersonal).

\blspace
  \dicEntry \dicTerm{svima} \dicPos{v} \dicFlx{impers}
\blspace

U~sloves je dále uvedeno, s~jakým pádem se pojí. Zkratka \textbf{acc} označuje 4. pád, \textbf{dat} 3. pád, \textbf{gen} 2. pád a~\textbf{nom} 1. pád. V~případě, že se sloveso pojí s~různými pády, je použito lomítko.

\blspace
  \dicEntry \dicTerm{klóra} \dicPos{v} \dicFlx{acc / dat}
\blspace

V~případě, že se sloveso pojí s~více pády, je použito znaménko plus.

\blspace
  \dicEntry \dicTerm{gefa} \dicPos{v} \dicFlx{dat + acc}
\blspace

\subsection*{Příslovce}

Příslovce jsou označena zkratkou \textbf{adv} (l. adverbium). Pokud chybí informace o~stupni, jedná se o~1. stupeň stupňování (\textbf{pos}). 2. stupeň je vždy označen zkratkou \textbf{comp} a~3. stupeň zkratkou \textbf{sup}.

\blspace
  \dicEntry \dicTerm{þar} \dicPos{adv}
  \dicEntry \dicTerm{ofar} \dicPos{adv} \dicFlx{comp (pos uppi, sup efst)}
  \dicEntry \dicTerm{síðast} \dicPos{adv} \dicFlx{sup (comp síðar)}
\blspace

\subsection*{Předložky}

Předložky jsou označeny zkratkou \textbf{prep} (l. praepositio). Pokud se předložka pojí výhradně s~jedním pádem, jsou použity zkratky pro 4. pád (\textbf{acc}), 3. pád (\textbf{dat}), 2. pád (\textbf{gen}).

\blspace
  \dicEntry \dicTerm{af} \dicPos{prep / adv} \dicFlx{dat}
\blspace

V~případě, že se předložka pojí s~různými pády, je použito lomítko.

\blspace
  \dicEntry \dicTerm{fyrir} \dicPos{prep / adv} \dicFlx{acc / dat}
\blspace

\subsection*{Spojky}

Spojky jsou označeny zkratkou \textbf{conj} (l. conjunctio).

\subsection*{Částice}

Částice jsou označeny zkratkou \textbf{part} (l. particula).

\subsection*{Citoslovce}

Citoslovce jsou označeny zkratkou \textbf{inter} (l. interjectio).

\subsection*{Předpony}

Předpony jsou označeny zkratkou \textbf{predp}.

\subsection*{Zkratky}

Zkratky jsou označeny zkratkou \textbf{zkr}.

\section{Skloňování, časování a~stupňování}

Každé heslo obsahuje zápis morfologické informace, která přesně informuje, jak se heslo skloňuje / časuje nebo stupňuje.

V~slovníku jsou použity dva odlišné systémy zápisu morfologické informace, které se vzájemně doplňují:

\begin{enumerate}
\item tradiční systém - do závorky za slovním druhem se zapisují deklinační / konjugační koncovky. Počet koncovek závisí na slovním druhu (viz dále).
\item numerický systém (vychází ze systému Helgi Haraldssona (\cite {is6})) - číslice odkazující na morfologický klíč, kde jsou podrobně vypsány deklinační / konjugační třídy.
Další informace o~numerickém systému lze nalézt na straně \pageref{sec:dec}.
\end{enumerate}

Deklinační / konjugační tvary nebo koncovky podstatných jmen, přídavných jmen, zájmen, číslovek, sloves a~příslovcí jsou uvedeny v~závorkách za slovním druhem.

Hnízdovací šev \textbf{'|'} označuje místo ve slově, kde se deklinační koncovka připojuje ke slovu.

\blspace
  \dicEntry \dicTerm{heils|a{\textsuperscript{1}}} \dicPos{f} \dicFlx{(-u)}
\blspace

Pokud není hnízdovací šev \textbf{'|'} uveden, tak se koncovka připojuje k~samotnému heslu.

\blspace
  \dicEntry \dicTerm{mynd} \dicPos{f} \dicFlx{(-ar, -ir)}
\blspace

Pokud dochází v~ohýbání k~přehlásce nebo je ohýbání nepravidelné, pak jsou uvedeny celé tvary hesla.

\blspace
  \dicEntry \dicTerm{maður} \dicPos{m} \dicFlx{(manns, menn)}
\blspace

V~případě, že heslové slovo má více tvarů, jsou varianty koncovky odděleny lomítkem.

\blspace
  \dicEntry \dicTerm{beð|ur} \dicPos{m} \dicFlx{(-s / -jar, -ir)}
\blspace

\subsection*{Podstatná jména}

U~podstatných jmen je uveden tvar hesla pro 2.~pád jednotného čísla a~1.~pád množného čísla.

\blspace
  \dicEntry \dicTerm{hval|ur} \dicPos{m} \dicFlx{(-s, -ir)}
\blspace

kde '\textit{hvals}' je 2.~pád jednotného čísla a~'\textit{hvalir}' 1.~pád množného čísla.
Pokud je uveden pouze jeden tvar, jedná se o~2.~pád jednotného čísla a~znamená to rovněž, že heslové slovo se neohýbá v~množném čísle.

\blspace
  \dicEntry \dicTerm{heils|a{\textsuperscript{1}}} \dicPos{f} \dicFlx{(-u)}
\blspace

Pokud je použita zkratka pro množné číslo (\textbf{pl}), znamená to, že se heslo ohýbá pouze v~množném čísle.

\blspace
  \dicEntry \dicTerm{afar··kostir} \dicPos{m} \dicFlx{pl}
\blspace

Pokud je použita zkratka \textbf{indecl}, znamená to, že podstatné jméno je nesklonné.

\blspace
  \dicEntry \dicTerm{október} \dicPos{m} \dicFlx{indecl}
\blspace

\subsection*{Přídavná jména}

U~přídavných jmen jsou uvedeny v~závorkách tvary, které poukazují na nepravidelnosti ve skloňování.

Pokud přídavné jméno v~ženském rodě neztrácí koncové '\textit{-r}', je uveden tvar hesla v~ženském rodě.

\blspace
  \dicEntry \dicTerm{hýr} \dicPos{adj} \dicFlx{(f -)}
\blspace

Dále jsou uvedeny tvary ženského rodu, ve~kterých dochází k~přehlásce:

\blspace
  \dicEntry \dicTerm{reglu··|samur} \dicPos{adj} \dicFlx{(f -söm)}
\blspace

Pokud je přídavné jméno nesklonné, je uvedena zkratka \textbf{indecl}.

\blspace
  \dicEntry \dicTerm{tví··mála} \dicPos{adj} \dicFlx{indecl}
\blspace

Pokud má přídavné jméno nepravidelné tvary 2.~a~3.~stupně stupňování, jsou tyto tvary uvedeny v~závorce spolu se zkratkou pro 2.~stupeň (\textbf{comp}) a~3.~stupeň (\textbf{sup}).

\blspace
  \dicEntry \dicTerm{þunnur} \dicPos{adj} \dicFlx{(comp þynnri, sup þynnstur)}
\blspace

V~případě, že je uveden pouze jeden ze dvou zbývajících stupňů, znamená to, že heslo se stupňuje pouze v~uvedených stupních.

\blspace
  \dicEntry \dicTerm{syðri} \dicPos{adj} \dicFlx{comp (sup syðstur)}
\blspace

\subsection*{Slovesa}

Islandská slovesa se dělí na slabá, silná a~nepravidelná. Vzhledem k~tomuto dělení se liší počet slovních tvarů uvedených v~závorce za zkratkou \textbf{v}. Princip uvádění slovních tvarů spočívá v~tom, že uživatel
je schopen (více méně) z~těchto tvarů odvodit časování slovesa. V~tomto slovníku jsme
záměrně umístili některá slabá slovesa mezi silná (místo dvou slovesných tvarů je uvedeno pět tvarů) kvůli tomu, aby byl uživatel přesněji informován o~časování slovesa.

Slabá slovesa, která tvoří 1.~osobu jednotného čísla minulého času oznamovacího způsobu koncovkou '\textit{-aði}', mají v~závorce uvedenu jen jednu koncovku, jmenovitě '\textit{(-aði)}'.

\blspace
  \dicEntry \dicTerm{ætl|a} \dicPos{v} \dicFlx{(-aði)}
\blspace

U~zbývajících skupin slabých sloves se v~závorce nachází koncovka 1.~osoby jednotného čísla minulého času oznamovacího způsobu a~supinum.

\blspace
  \dicEntry \dicTerm{kenn|a} \dicPos{v} \dicFlx{(-di, -t)}
\blspace

Silná nebo nepravidelná slovesa mají v~závorce vždy pět tvarů, jmenovitě 1.~osobu jednotného čísla přítomného času oznamovacího způsobu, 1.~osobu jednotného čísla minulého času oznamovacího způsobu, 1.~osobu množného čísla minulého času oznamovacího způsobu, 1.~osobu jednotného čísla minulého času spojovacího způsobu a~supinum.

\blspace
  \dicEntry \dicTerm{grípa} \dicPos{v} \dicFlx{(gríp, greip, gripum, gripi, gripið)}
\blspace

Silná nebo nepravidelná slovesa vyskytující se pouze v~neosobním užití mají v~závorce 4 tvary, jmenovitě 3.~osobu jednotného čísla přítomného času oznamovacího způsobu, 3.~osobu jednotného čísla minulého času oznamovacího způsobu, 3.~osobu jednotného čísla minulého času spojovacího způsobu a~supinum.

\blspace
  \dicEntry \dicTerm{kross··|bregða} \dicPos{v} \dicFlx{(-bregður, -brá, -brygði, -brugðið)}
\blspace

U~některých významů v~heslovém slově se vyskytují gramatické informace, které popisují chování daného významu.

Zkratka \textbf{impers} označuje, že se daný význam slovesa vyskytuje v~neosobním užití. Pomocí slovního spojení je vyjádřeno užití slovesa.

\blspace
  \dicEntry \dicTerm{batn|a} \dicPos{v} \dicFlx{(-aði)} \dicTerm{e-m batnar} \dicFlx{impers} (kdo) se uzdravuje, (kdo) se zotavuje, (komu) je lépe
\blspace

Zkratka \textbf{refl} označuje mediopasivní tvary slovesa.

\blspace
  \dicEntry \dicTerm{að··lag|a} \dicPos{v} \dicFlx{(-aði)} \dicIdiom{aðlagast} {\dicPhraseIS{aðlagast e-u}} \dicFlx{refl} {přizpůsobit se (čemu)}
\blspace

\subsection*{Příslovce}

U~příslovcí se v~záhlaví nachází tvary stupňování. Tyto tvary jsou uvedeny v~závorce spolu se zkratkou pro 2.~stupeň (\textbf{comp}) a~3.~stupeň (\textbf{sup}).

\blspace
  \dicEntry \dicTerm{vel\textsuperscript{2}} \dicPos{adv} \dicFlx{(comp betur, sup best)}
\blspace

Nepravidelné tvary 2.~a~3.~stupně jsou uvedeny buď jako hesla odkazující na základní tvar v~1.~stupni (\textbf{pos}), anebo jako svébytná hesla. V~obou případech jsou uvedeny v~závorce zbývající stupně.

\blspace
  \dicEntry \dicTerm{síst} \dicPos{adv} \dicFlx{sup (pos varla, comp síður)}
\blspace

Pokud některý ze stupňů uveden není, znamená to, že pro dané slovo neexistuje.

\blspace
  \dicEntry \dicTerm{síðar} \dicPos{adv} \dicFlx{comp (sup síðast)}
\blspace

\section{Morfologická informace} \label{sec:dec}

V~záhlaví hesla se kromě deklinačních / konjugačních koncovek nachází také morfologické numerické informace, které odkazují na morfologický klíč (viz strana \pageref{sec:morpho}), kde jsou podrobně vypsány kompletní deklinační / konjugační tabulky.

\blspace
  \dicEntry \dicTerm{ap|i} \dicPos{m}[1] \dicFlx{(-a, -ar)}[8]
\blspace

Morfologická informace je zapsána podle následujícího schématu: \textbf{m{\textsubscript{1}} (-a, -ar){\textsubscript{8}}}

\blspace
\tableB
\blspace

Deklinační / konjugační třídy a~jejich počet nejsou v~islandštině ustálené a~slouží pouze pro rozlišení odlišných skloňování / časování.

Názvy podtříd slouží pro rozlišení odlišných skloňování / časování.

\subsection*{Práce s~morfologickým klíčem}

Pro vyhledání kompletních deklinačních / konjugačních tabulek slouží následující postup, který demonstrujeme na ukázce:

\blspace
  \dicEntry \dicTerm{kis|a} \dicPos{f}[1] \dicFlx{(-u, -ur)}[7]
\blspace

Nejdříve určíme slovní druh (v~případě podstatných jmen také rod), v~tomto případě podstatné jméno rodu ženského, a~navštívíme Morfologický klíč (str. \pageref{sec:morpho}) pro podstatná jména rodu ženského (str. \pageref{sec:morpho_f}).
Následně dohledáme třídu (v~tomto případě číslo {\textsubscript{\textbf{1}}}) a~následně dohledáme dodatečné numerické informace za závorkou (v~tomto případě číslo {\textsubscript{\textbf{7}}}).

\tableC

\section{Definice}

\subsection*{Základní tvar}

Základním tvarem definice je, že jednomu islandskému slovu odpovídá jeden nebo více českých významů.

\blspace
  \dicEntry \dicTerm{landa··fræð|i} \dicPos{f} \dicFlx{(-i)} \textCS{zeměpis, geografie}
\blspace

V~některých případech je v~české části použita závorka. V~prvním případě je závorka užívaná ke~zkrácení zápisu dvou slov ('\textit{pomoc}', '\textit{výpomoc}').

\blspace
  \dicEntry \dicTerm{að··stoð} \dicPos{f} \dicFlx{(-ar)} \textCS{(vý)pomoc, asistence}
\blspace

V~druhém případě se v~závorce nachází upřesňující informace, která zabraňuje dvojznačnosti.

\blspace
  \dicEntry \dicTerm{fíkj|a} \dicPos{f} \dicFlx{(-u, -ur)} \textCS{fík} {\footnotesize {\textCS{(plod)}}}
\blspace

V~mnoha případech je závorka v~české části použita pro příklad užití českého slova a~tím vymezení jeho významu.

\blspace
  \dicEntry \dicTerm{að·gengi··legur} \dicPos{adj} {přístupný, dostupný (informace ap.)}
\blspace

Lomítko \textbf{'/'} slouží v~české části významu k~oddělení více významů pojících se s~jedním nebo více slovy a~je použito v~české části následujícím způsobem:

\blspace
  \dicEntry \dicTerm{aðal··|gata} \dicPos{f} \dicFlx{(-götu, -götur)} \textCS{hlavní třída / ulice}
\blspace

Význam přečteme jako '\textit{hlavní třída, hlavní ulice}'.

\subsection*{Opis heslového slova}

V~případech, kdy islandské slovo nemá ekvivalent v~českém jazyce, jsou významy islandských slov opsány delším opisem a~zobrazeny menším písmem. Příkladem jsou slova z~lidových pověstí, kulinářské speciality nebo botanické pojmy.

\blspace
  \dicEntry \dicTerm{til··ber|i} \dicPos{m} \dicFlx{(-a, -ar)} {\footnotesize{pov.}} {\footnotesize{stvoření, které dojí krávy a~ovce hospodářů}}
\blspace

\subsection*{Slovní spojení}

V~definici hesla jsou kromě samotných významů uváděna také slovní spojení, která jsou sázena tučně. Po slovním spojení následuje překlad slovního spojení.

\blspace
  \dicEntry \dicTerm{lær|a} \dicPos{v} \dicFlx{(-ði, -t)} \textbf{læra e-ð utanbókar} {učit se (co) zpaměti, memorovat (co)}
\blspace

Závorky jsou použity na tu část slovního spojení, kterou je možné vypustit.

\blspace
  \dicEntry \dicTerm{ár\textsuperscript{1}} \dicPos{f} \dicFlx{(-ar, -ar)} \textbf{taka (of) djúpt í árinni} {\footnotesize{přen.}} {přehnat (co v~tvrzení)}
\blspace

Příklad můžeme přečíst jako '\textit{taka djúpt í árinni}' nebo '\textit{taka of djúpt í árinni}'.

Mezi slovní spojení v~tomto slovníku řadíme a) tvary heslového slova, b) příkladová spojení, c) ustálená slovní spojení a~d) rčení nebo přísloví.

\begin{enumerate}
\item Tvary heslového slova\\
Do této skupiny řadíme tvar množného čísla hesla (zkratka \textbf{pl}).
\blspace
\dicEntry \dicTerm{bót} \dicPos{f} \dicFlx{(-ar, bætur)} \textbf{bætur} \dicFlx{pl} {dávky, přídavky, příspěvky}\\

Dále do této skupiny řadíme tvar slovesa v~středním slovesném rodě (zkratka \textbf{refl}).\\

\dicEntry \dicTerm{finna} \dicPos{v} \dicFlx{(finn, fann, fundum, fyndi, fundið)} \dicIdiom{finnast} {\dicPhraseIS{finnast}} \dicFlx{refl} {být nalezen, najít se, nalézt se, nacházet se}
\blspace
\item Příkladová spojení\\
Příkladová spojení ilustrují komunikační chování heslových slov v~příslušných významech a~v~různých syntagmatech (např. sloveso s~určitou předložkou ap.). Jsou označena tučně modrou barvou a~odsazená pro snadnou orientaci jako svébytné heslo.
\blspace
\dicEntry \dicTerm{koma\textsuperscript{2}} \dicPos{v} \dicFlx{(kem, kom, komum, kæmi, komið)} \dicIdiom{koma}[að] {\dicPhraseIS{koma að e-m}} {přiblížit se ke (komu), přijít ke (komu)}
\blspace
\item Ustálená slovní spojení\\
Tento druh slovních spojení je charakterizován většinou přeneseným významem. Tato spojení jsou uváděna za posledním významem nebo příkladovým spojením a~jsou označená tučně modrou barvou a~odsazená pro snadnou orientaci jako svébytné heslo.
\blspace
\dicEntry \dicTerm{kast} \dicPos{n} \dicFlx{(-s, köst)} \dicIdiom{kast} {\dicPhraseIS{þegar til kastanna kemur}} {\footnotesize{přen.}} {když jde do tuhého}
\blspace
\item Rčení a~přísloví\\
Rčení a~přísloví se nacházejí vždy na konci definice a~jsou odlišená symbolem \dicsymProverb\ a~odsazená pro snadnou orientaci.
\blspace
\dicEntry \dicTerm{æfing} \dicPos{f} \dicFlx{(-ar, -ar)} \dicProverb{} \dicPhraseIS{Æfing skapar meistara.} {\footnotesize{přís.}} \textCS{Cvičení dělá mistra.}
\end{enumerate}

\section{Vícevýznamovost}

\subsection*{Vícevýznamová hesla}

Značné množství hesel má více než jeden význam. Významy jsou odlišeny arabskými číslicemi a~v~některých případech i~islandskými synonymy.

\blspace
  \dicEntry \dicTerm{land} \dicPos{n} \dicFlx{(-s, lönd)} {\textbf{1.}} \dicSynonym{þurrlendi} {souš, pevnina, země} {\textbf{2.}} \dicSynonym{árbakki} {břeh} {\textbf{3.}} \dicSynonym{ríki} {země, stát, vlast} {\textbf{4.}} \dicSynonym{landareign} {země, půda, pozemek}
\blspace

\subsection*{Vícevýznamová slovní spojení}

V~případě, že slovní spojení uvnitř definice má více významů, jsou jednotlivé významy odlišeny písmeny {\textbf{a., b., c.}} atd.

\blspace
  \dicEntry \dicTerm{drag} {\dicPos{n} \dicFlx{(-s, drög)} {\textbf{1.}} {mokřina, podmoklý terén;} {\textbf{2.}} \textbf{drög} {\footnotesize{pl}} {\textbf{a.}} \dicSynonym{uppsprettur} {prameny (řeky ap.)} {\textbf{b.}} \dicSynonym{undirbúningur} {náčrt, návrh.}}
\blspace

\section{Řazení významů a~slovních spojení v~definici}
\longsectionskip

V~definici jsou nejdříve seřazeny významy slova a~odlišeny arabskými číslicemi.

\subsection*{V rámci významu}

Slovní spojení jsou uváděna v~rámci jednotlivého významu, pokud se s~ním významově vážou.

\blspace
  \dicEntry \dicTerm{tím|i} \dicPos{m} \dicFlx{(-a, -ar)} {\textbf{1.}} \dicSynonym{tíð} {čas, doba}; \textbf{í þann tíma} \dicFlx{adv} \dicSynonym{þá\textsuperscript{2}} {pak} {\textbf{2.}} \dicSynonym{rétt stund} {(správný) čas / moment}
\blspace

\subsection*{Za posledním očíslovaným významem}

Pokud se slovní spojení neváže k~žádnému významu, je řazeno za posledním očíslovaným významem. Ustálená slovní spojení jsou uvedena symbolem \dicsymIdiom.

\blspace
  \dicEntry \dicTerm{finna} \dicPos{v} \dicFlx{(finn, fann, fundum, fyndi, fundið)} {\textbf{3.}} \dicSynonym{skynja} {cítit, vnímat} \dicIdiom{finna}[að] {\dicPhraseIS{finna að e-u}} {kritizovat (co), nacházet chyby na (čem)}
\blspace

\subsection*{Způsob řazení slovních spojení}

Slovní spojení, která se nachází za posledním očíslovaným významem, jsou řazena podle předložky / příslovce / zájmena, se kterou se pojí. Pro snadnou orientaci v~rozsáhlých definicích je uvedeno heslové slovo nebo tvar heslového slova plus předložka / příslovce / zájmeno a~barevně odlišeno modrou barvou.

\blspace
  \dicEntry \dicTerm{koma\textsuperscript{2}} \dicPos{v} \dicFlx{(kem, kom, komum, kæmi, komið)} \dicIdiom{koma}[út] {\dicPhraseIS{koma út}} {vyjít (kniha ap.), objevit se (na trhu ap.)}
\blspace

Poté následují slovní spojení v~slovesném středním rodě řazená opět podle předložky / příslovce / zájmena, se kterou se pojí:

\blspace
  \dicEntry \dicTerm{koma\textsuperscript{2}} \dicPos{v} \dicFlx{(kem, kom, komum, kæmi, komið)} \dicIdiom{komast} {\dicPhraseIS{komast}} \dicFlx{refl} {dostat se, dostávat se}
\blspace

Poté následují ustálená slovní spojení, která nebylo možné zařadit do předchozích položek:

\blspace
  \dicEntry \dicTerm{koma\textsuperscript{2}} \dicPos{v} \dicFlx{(kem, kom, komum, kæmi, komið)} \dicIdiom{koma} {\dicPhraseIS{(komdu) sæll (og blessaður)}} {ahoj, nazdar, dobrý den}
\blspace

Nakonec jsou uvedena přísloví a~rčení označená symbolem \dicsymProverb\:

\blspace
  \dicEntry \dicTerm{eik} \dicPos{n} \dicFlx{(-ar / -ur, -ur)} \dicProverb{} \dicPhraseIS{Ekki / eigi fellur eik við fyrsta högg.} {\footnotesize{přís.}} {Jednou ranou dub nespadne.}
\blspace

\section{Synonyma a~antonyma / hesla ke~srovnání}
\longsectionskip

\subsection*{Synonyma}

Synonyma nebo synonymní slovní spojení jsou uvedeny před opisem významu hesla. Jsou umístěna v~závorkách. Islandská synonyma slouží a) islandskému uživateli k~rozlišení českých významů b) českému uživateli jako dodatečná informace o~významu slova.

\blspace
  \dicEntry \dicTerm{staða} \dicPos{f} \dicFlx{(stöðu, stöður)} {\textbf{1.}} \dicSynonym{ástand} {situace, stav, poměry} {\textbf{2.}} \dicSynonym{embætti} {post, (pracovní) místo / pozice (ve firmě ap.)}
\blspace

V~případě, že synonymum lze nalézt jako heslo ve slovníku, je použita šipka '\dicsymSee\ '.

\blspace
  \dicEntry \dicTerm{bygging} {\dicPos{f} \dicFlx{(-ar, -ar)} {\textbf{1.}} \dicSynonym{hús} {budova, stavba, stavení} {\textbf{2.}} \textit{(það að byggja)} {(vý)stavba} }
\blspace

\subsection*{Antonyma / Hesla ke~srovnání}

Antonyma či hesla ke~srovnání jsou uvedena za opisem významu heslového slova a~jsou označena šipkou '$\shortuparrow$'.

\blspace
  \dicEntry \dicTerm{kaldur} \dicPos{adj} \dicFlx{(f köld)} {\textbf{1.}} {studený, chladný} \textit{($\shortuparrow$ heitur)}
\blspace

\section{Oborové a~stylové charakteristiky}

K~upřesnění významu heslového slova slouží oborové a~stylové charakteristiky.

\subsection*{Oborové charakteristiky}

Oborové charakteristiky zařazují islandské slovo do určitého oboru a~naznačují použití daného slova v~jazyce

\blspace
  \dicEntry \dicTerm{berg} \dicPos{n} \dicFlx{(-s, -)} {\textbf{1.}} {\footnotesize{geol.}} {hornina, kámen} {\textbf{2.}} \dicSynonym{klöpp\textsuperscript{1}} {skalní stěna, skalisko}
\blspace

a~vymezují význam českého slova.

\blspace
  \dicEntry \dicTerm{mús} \dicPos{f} \dicFlx{(-ar, mýs)} {\textbf{1.}} {\footnotesize{zool.}} {myš} (\textit{l. Mus}) {\textbf{2.}} {\footnotesize{poč.}} {myš}
\blspace

\subsection*{Stylové charakteristiky}

Stylové charakteristiky jsou uváděny v~případě, kdy význam heslového slova je možné použít pouze při určité příležitosti nebo je význam jazykově zabarven.

\blspace
  \dicEntry \dicTerm{gaur} \dicPos{m} \dicFlx{(-s, -ar)} {\textbf{1.}} {\footnotesize{hovor.}} {chlápek, maník, člověk}
\blspace

\section{Syntax}

Syntaktické informace popisují:

\begin{enumerate}
\item zdali je sloveso tranzitivní nebo intranzitivní a~s~jakým(i) pádem (pády) se pojí, pokud je tranzitivní,
\item zdali je podmět / předmět životný či neživotný,
\item zdali je sloveso osobní či neosobní a~v~jakém pádu je podmět, pokud je sloveso neosobní,
\item zdali je podmět gramatický,
\item zdali je sloveso použito ve středním rodě,
\item nejčastější příklady, které se pojí se slovesem,
\item se kterými příslovci nebo předložkami se sloveso pojí.
\end{enumerate}

Syntaktické informace se ve slovníku vyskytují ve třech formách: a) informace uvedené zkratkou, b) teoretické příklady a~c) praktické příklady.

\subsection*{Syntaktické informace uvedené zkratkou}

Zde máme na mysli informace o~pádě (\textbf{nom, acc, dat, gen}), slovesu v~středním slovesném rodě (\textbf{refl}), neosobním slovesu (\textbf{impers}), čísle (\textbf{sg, pl}), rozkazovacím způsobu (\textbf{imper}) a~spojovacím způsobu (\textbf{con}).

\subsection*{Teoretické příklady}

Teoretickým příkladem nazýváme příklad, kdy podmět nebo předmět je zastoupen zájmenem neurčitým v~pádě, se~kterým se slovo pojí a~sloveso je zpravidla v~infinitivu. V~případě, že podmět není v~1.~pádě (např. u~sloves neosobních), je sloveso uvedeno v~3.~osobě jednotného čísla přítomného času.
Ve slovních spojeních používáme zkratky pro tvary neurčitého zájmena \textbf{einhver} (někdo) a~\textbf{eitthvað} (něco). Tyto zkratky označují a) pád, b) číslo a~c) životnost nebo neživotnost podmětu / předmětu.\\
Následuje seznam zkratek neurčitých zájmen použitých ve slovníku: \\

\tableD
\blspace
Příklad použití zkratek s~neurčitými zájmeny:

\blspace
  \dicEntry \dicTerm{leit|a} \dicPos{v} \dicFlx{(-aði)} \textbf{leita að e-u / e-m} {hledat (co / koho)} \dicExampleIS{leita að lyklunum} \dicExampleCS{hledat klíče}
\blspace

Příklad tedy přečteme jako '\textit{leita að einhverju}' (\textit{hledat co}) a~'\textit{leita að einhverjum}' (\textit{hledat koho}). V~tomto případě se sloveso pojí s~3.~pádem a~předmět může být životný i~neživotný.

\subsection*{Praktické příklady}

Ve slovníku se vyskytuje velké množství praktických příkladů a~jejich překladů. Praktický příklad se liší od příkladu teoretického tím, že místo zájmen neurčitých jsou uvedena frekventovaná slova z~běžného projevu. Praktické příklady slouží k~ilustraci použití
heslového slova.\\
Praktickým příkladem může být a) (složené) slovo, b) slovní spojení nebo c) věta.

\begin{enumerate}
\item Slovo jako praktický příklad:
\blspace
\dicEntry \dicTerm{hett|a} \dicPos{f} \dicFlx{(-u, -ur)} \textbf{1.} {kapuce, kápě} \dicExampleIS{hettuúlpa} \dicExampleCS{bunda s~kapucí}
\blspace
\item Slovní spojení jako praktický příklad:
\blspace
\dicEntry \dicTerm{tím|i} \dicPos{m} \dicFlx{(-a, -ar)} \textbf{3.} \dicSynonym{klukkustund} {hodina} {\footnotesize{(šedesát minut)}} \dicExampleIS{tveggja tíma gangur} \dicExampleCS{dvouhodinový pochod}
\blspace
\item Věta jako praktický příklad:
\blspace
\dicEntry \dicTerm{mæt|a} \dicPos{v} \dicFlx{(-ti, -t)} \dicIdiom{mætast} \dicPhraseIS{mætast} \dicSynonym*{ná saman} \dicFlx{refl} {setkat se, potkat se} \dicExampleIS{Þau mættust á gangstéttinni.} \dicExampleCS{Potkali se na chodníku.}
\end{enumerate}

\section{Fotografie a~ilustrace}

Ve slovníku je publikováno 630 fotografií rostlin, živočichů a~fotografií, která dokumentují místopisně heslová slova (např. kulinářské speciality).
Fotografie jsou publikovány pod veřejnými licencemi. Jména autorů a~názvy licencí lze nalézt v~Seznamu autorů fotografií a~ilustrací na straně \pageref{sec:photo}. U~hesla a~konkrétně u~daného významu, ke kterému se fotografie vztahuje, je použit symbol fotoaparátu \dicsymPhoto. Fotografie je umístěna buď pod heslem, anebo v~horní části stránky, kde se heslo nachází, anebo na stránce (stránkách) následující(ch).

%\par\begin{center}\setlength\fboxsep{0pt}\setlength\fboxrule{0.5pt}
%\fbox{\includegraphics[width=6cm]{ds_image_posthus_0_2.jpg}}\end{center}
%\par\begin{center}\footnotesize {Pósthús}\end{center}

\end{multicols}

% ---------------------------------------- List of abbreviations and symbols -

\chapter{Seznam použitých zkratek a symbolů}
\begin{multicols}{2}
\begin{description}[leftmargin=!,labelwidth=1.5cm]
%[\breaklabel\setleftmargin{40pt}\setlabelstyle{\bfseries}
%\addtolength{\itemsep}{-0.5\baselineskip}]
\item[{abb}] {zkrácený}
\item[{acc}] {4. pád, akuzativ}
\item[{adj}] {přídavné jméno}
\item[{adv}] {příslovce}
\item[{akt}] {rod činný, aktivum}
\item[{anat.}] {anatomie}
\item[{angl.}] {anglicky}
\item[{astro.}] {astronomie, výzkum vesmíru}
\item[{ap.}] {a podobně}
\item[{básn.}] {básnický, poetický výraz}
\item[{biol.}] {biologie}
\item[{bot.}] {botanika}

\item[{comp}] {2. stupeň stupňování, komparativ}
\item[{con}] {spojovací způsob, konjunktiv}
\item[{conj}] {spojka}
\item[{dat}] {3. pád, dativ}
\item[{def}] {určitý tvar}
\item[{dem}] {zájmeno ukazovací}
\item[{dět.}] {dětsky}

\item[{e-að}] {co (sg nom neživotný) (\textit{eitthvað})}
\item[{e-ð}] {co (sg acc neživotný) (\textit{eitthvað})}
\item[{e-ir}] {kdo (pl nom životný) (\textit{einhverjir})}
\item[{e-ja}] {koho (pl acc životný) (\textit{einhverja})}
\item[{e-jum}] {komu (pl dat životný) (\textit{einhverjum})}
\item[{e-m}] {komu (sg dat životný) (\textit{einhverjum})}
\item[{e-n}] {koho (sg acc životný) (\textit{einhvern})}
\item[{e-r}] {kdo (sg nom životný) (\textit{einhver})}
\item[{e-rra}] {koho (pl gen životný) (\textit{einhverra})}
\item[{e-rs}] {koho (sg gen životný) (\textit{einhvers})}
\item[{e-s}] {čeho (sg gen neživotný) (\textit{einhvers})}
\item[{e-u}] {čemu (sg dat neživotný) (\textit{einhverju})}

\item[{ekon.}] {ekonomika, obchod}
\item[{elek.}] {elektřina}
\item[{f}] {rod ženský}
\item[{filos.}] {filosofie, logika}
\item[{form.}] {formálně}
\item[{fyz.}] {fyzika}
\item[{gen}] {2. pád, genitiv}
\item[{geog.}] {zeměpis, geografie}
\item[{geol.}] {geologie}
\item[{han.}] {hanlivě, pejorativně}
\item[{hist.}] {historie}
\item[{hovor.}] {hovorově}
\item[{hrub.}] {hrubě}
\item[{hud.}] {hudba}
\item[{imper}] {rozkazovací způsob, imperativ}
\item[{impers}] {sloveso neosobní}
\item[{ind}] {oznamovací způsob, indikativ}
\item[{indecl}] {nesklonný}
\item[{indef}] {neurčitý tvar, (pron +) zájmeno neurčité}
\item[{inf}] {infinitiv}
\item[{int}] {tázací zájmeno}
\item[{inter}] {citoslovce}
\item[{chem.}] {chemie}
\item[{jaz.}] {jazykověda}
\item[{kulin.}] {kulinářství, vaření}
\item[{l.}] {latinsky}
\item[{let.}] {letectví}
\item[{lit.}] {literatura, vydavatelství}
\item[{m}] {rod mužský}
\item[{mat.}] {matematika}
\item[{med. }] {lékařství}
\item[{meteo.}] {meteorologie}
\item[{myt.}] {mytologie}
\item[{n}] {rod střední}
\item[{náb.}] {náboženství}
\item[{nám.}] {námořnictví, rybolov}
\item[{nom}] {1. pád, nominativ}
\item[{num}] {číslovka}
\item[{ord}] {řadová číslovka}
\item[{p}] {osoba}
\item[{part}] {částice}
\item[{pers}] {(v +) osoba, (pron +) zájmeno osobní}
\item[{pf}] {minulý čas}
\item[{pl}] {množné číslo, plurál}
\item[{poč.}] {informatika}
\item[{pol.}] {politika, politologie}
\item[{poss}] {zájmeno přivlastňovací}
\item[{pos}] {1. stupeň stupňování, pozitiv}
\item[{pov.}] {lidové pověsti, folkloristika}
\item[{pp}] {příčestí minulé}
\item[{praes}] {přítomný čas}
\item[{práv.}] {právnictví, soudnictví}
\item[{predp}] {předpona}
\item[{prep}] {předložka}
\item[{presp}] {příčestí přítomné}
\item[{pron}] {zájmeno}
\item[{prop}] {vlastní jméno, proprium}
\item[{přen.}] {přeneseně}
\item[{přís.}] {přísloví}
\item[{psych.}] {psychologie}
\item[{refl}] {(v +) médium, (pron +) zvratné zájmeno}
\item[{rel}] {vztažné zájmeno}
\item[{sg}] {jednotné číslo, singulár}
\item[{slang.}] {slang}
\item[{sport.}] {sport}
\item[{stav.}] {stavebnictví, architektura}
\item[{subs}] {podstatné jméno (bez rodu)}
\item[{sup}] {3. stupeň stupňování, superlativ}
\item[{supin}] {supinum}
\item[{škol.}] {školství}
\item[{techn.}] {technika, mechanika}
\item[{v}] {sloveso}
\item[{voj.}] {vojenství}
\item[{zast.}] {zastarale}
\item[{zkr}] {zkratka}
\item[{zool.}] {zoologie}

\item[{;}] {středník uprostřed definice odděluje věty; středník na konci definice odděluje významy}

\item[{(...)}] {v kulatých závorkách jsou uvedeny doplňkové informace}
\item[{[...]}] {v hranatých závorkách je uveden fonetický zápis výslovnosti}\footnote{Vysvětlení fonetických symbolů je uvedeno v kapitole Fonetika \ref{sec:phon_phonems} na straně \pageref{sec:phon_phonems}}
\item[{|}] {naznačuje místo, kde se deklinační koncovka připojuje ke slovnímu základu}
\item[{·}] {vedlejší dělení složeného slova}
\item[{··}] {hlavní dělení složeného slova}
\item[{+}] {plus}
\item[{/}] {lomítko označuje více možností}
\item[{\dicsymSee}] {šipka odkazuje na heslo ve slovníku}
\item[{\dicsymCompare}] {šipka odkazuje na heslo ve slovníku ke srovnání}
\item[{\dicsymExampleIS}] {uvozuje islandský příklad}
\item[{\dicsymExampleCS}] {uvozuje český překlad příkladu}
\item[{\dicsymIdiom}] {uvozuje slovní spojení}
\item[{\dicsymFrequent}] {heslo patří mezi frekventované výrazy}

\item[{\dicsymPhoto}] {fotografie, ilustrace}
\item[{\dicsymProverb}] {přísloví, rčení}

\end{description}
\end{multicols}

\chapter{Skrá yfir skammstafanir og tákn}
\begin{multicols}{2}
\begin{description}[leftmargin=!,labelwidth=1.5cm]
%[\breaklabel\setleftmargin{40pt}\setlabelstyle{\bfseries}
%\addtolength{\itemsep}{-0.5\baselineskip}]
\item[{abb}] {stýfður}
\item[{acc}] {þolfall}
\item[{adj}] {lýsingarorð}
\item[{adv}] {atviksorð}
\item[{akt}] {germynd}
\item[{anat.}] {líffærafræði}
\item[{angl.}] {enska}
\item[{astro.}] {geimvísindi, stjörnufræði}
\item[{ap.}] {e.þ.h.}
\item[{básn.}] {ljóðrænn}
\item[{biol.}] {líffræði}
\item[{bot.}] {grasafræði}

\item[{comp}] {miðstig}
\item[{con}] {viðtengingarháttur}
\item[{conj}] {samtenging}
\item[{dat}] {þágufall}
\item[{def}] {með greini}
\item[{dem}] {ábendingarfornafn}
\item[{dět.}] {barnamál}

\item[{e-að}] {(\textit{eitthvað})}
\item[{e-ð}] {(\textit{eitthvað})}
\item[{e-ir}] {(\textit{einhverjir})}
\item[{e-ja}] {(\textit{einhverja})}
\item[{e-jum}] {(\textit{einhverjum})}
\item[{e-m}] {(\textit{einhverjum})}
\item[{e-n}] {(\textit{einhvern})}
\item[{e-r}] { (\textit{einhver})}
\item[{e-rra}] {(\textit{einhverra})}
\item[{e-rs}] {(\textit{einhvers})}
\item[{e-s}] {(\textit{einhvers})}
\item[{e-u}] { (\textit{einhverju})}

\item[{ekon.}] {hagfræði, viðskipti}
\item[{elek.}] {rafmagn}
\item[{f}] {kvenkyn}
\item[{filos.}] {heimspeki, rökfræði}
\item[{form.}] {formlegt mál}
\item[{fyz.}] {eðlisfræði}
\item[{gen}] {eignarfall}
\item[{geog.}] {landafræði}
\item[{geol.}] {jarðfræði}
\item[{han.}] {niðrandi}
\item[{hist.}] {sagnfræði}
\item[{hovor.}] {talmál}
\item[{hrub.}] {gróft}
\item[{hud.}] {tónlist}
\item[{imper}] {boðháttur}
\item[{impers}] {ópersónuleg sögn}
\item[{ind}] {framsöguháttur}
\item[{indecl}] {óbeygjanlegur}
\item[{indef}] {óákveðinn, (pron +) óákveðið fornafn}
\item[{inf}] {nafnháttur}
\item[{int}] {spurnigarfornafn}
\item[{inter}] {upphrópun}
\item[{chem.}] {efnafræði}
\item[{jaz.}] {málfræði, málvísindi}
\item[{kulin.}] {eldhús}
\item[{l.}] {latína}
\item[{let.}] {flug, flugmál}
\item[{lit.}] {skáldskapur}
\item[{m}] {nafnorð, karlkyn}
\item[{mat.}] {stærðfræði}
\item[{med. }] {læknisfræði}
\item[{meteo.}] {veðurfræði}
\item[{myt.}] {goðafræði}
\item[{n}] {nafnorð, hvorugkyn}
\item[{náb.}] {trúfræði}
\item[{nám.}] {sjómennska}
\item[{nom}] {nefnifall}
\item[{num}] {töluorð}
\item[{ord}] {raðtala}
\item[{p}] {persóna}
\item[{part}] {ögn}
\item[{pers}] {(v +) persóna, (pron +) persónufornafn}
\item[{pf}] {þátíð}
\item[{pl}] {fleirtala}
\item[{poč.}] {tölvufræði, tölvur}
\item[{pol.}] {stjórnmál, stjórnmálafræði}
\item[{poss}] {eignarfornafn}
\item[{pos}] {frumstig}
\item[{pov.}] {þjóðtrú}
\item[{pp}] {lýsingarháttur þátíðar}
\item[{praes}] {nútíð}
\item[{práv.}] {lögfræði, dómsmál}
\item[{predp}] {forskeyti}
\item[{prep}] {forsetning}
\item[{presp}] {lýsingarháttur nútíðar}
\item[{pron}] {fornafn}
\item[{prop}] {sérnafn}
\item[{přen.}] {afleidd merking}
\item[{přís.}] {málsháttur}
\item[{psych.}] {sálfræði}
\item[{refl}] {(v +) miðmynd, (pron +) afturbeygt fornafn}
\item[{rel}] {tilvísunarfornafn}
\item[{sg}] {eintala}
\item[{slang.}] {slangur}
\item[{sport.}] {íþróttir}
\item[{stav.}] {húsasmíði, byggingarlist}
\item[{subs}] {nafnorð (án kyns)}
\item[{sup}] {efstastig}
\item[{supin}] {sagnbót}
\item[{škol.}] {skólamál}
\item[{techn.}] {tækni, verkfræði}
\item[{v}] {sagnorð}
\item[{voj.}] {her}
\item[{zast.}] {úrelt mál}
\item[{zkr}] {skammstöfun}
\item[{zool.}] {dýrafræði}

\item[{;}] {semíkomma í miðri skýringu aðgreinir setningu; semíkomma í enda skýringar aðgreinir mismunandi merkingarsvið}

\item[{(...)}] {innan sviga eru viðbótarupplýsingar}
\item[{[...]}] {innan hornklofa er hljóðritaður framburður}
\item[{|}] {merkir stað, þar sem beygingarending bætist við stofn}
\item[{·}] {aukasamskeyti samsetts orðs}
\item[{··}] {aðalsamskeyti samsetts orðs}
\item[{+}] {plús}
\item[{/}] {skástrik táknar fleiri möguleika}
\item[{\dicsymSee}] {ör vísar til uppflettiorðs í orðabókinni}
\item[{\dicsymCompare}] {ör vísar til uppflettiorðs í bókinni til samanburðar}
\item[{\dicsymExampleIS}] {kemur á undan dæmi á íslensku}
\item[{\dicsymExampleCS}] {kemur á undan tékkneskri þýðingu á dæminu}
\item[{\dicsymIdiom}] {orð með háa tíðni}

\item[{\dicsymPhoto}] {ljósmynd}
\item[{\dicsymProverb}] {málsháttur}

\end{description}
\end{multicols}

% ------------------------------------------------- Headword part and alphabet

\chapter{Heslová část}
{\renewcommand{\arraystretch}{2}%
%\topskip0pt
\vstretch{1}
{\setlength{\tabcolsep}{1.5em}
\begin{table}[htpb]
\centering

\begin{tabular}{ccccccc} \hline
A a &  & G g &  & O o &  & V v \\
Á á  &  & H h &  & Ó ó &  & W w \\
B b &  & I i &  & P p &  & X x \\
C c &  & Í í &  & Q q &  & Y y \\
D d &  & J j &  & R r &  & Ý ý \\
Ð ð &  & K k &  & S s &  & Z z \\
E e &  & L l &  & T t &  & Þ þ \\
É é  &  & M m &  & U u &  & Æ æ \\
F f &  & N n &  & Ú ú &  & Ö ö \\
\hline
\end{tabular}

%\caption{\textbf{Islandská abeceda}}
\label{my-islandska-abeceda}
\end{table}
}
\vstretch{1}
{\renewcommand{\arraystretch}{1}%

\clearpage

\fi

% =========================================================== THE DICTIONARY =

\ifinputletters

\cleardoublepage

\dictionarygeometry
\pagestyle{myheadings}

\newcommand*{\alphabetF}{%
  a,aa,b,c,d,e,ee,f,g,h,i,ii,j,k,l,m,n,o,oo,p,r,s1,s2,t,u,uu,v,w,x,y,yy,th,ae,oe}
\ifPDFtest
\dicLetter{f}{letter8}
\dicEntry[f.h.] \dicTerm{f.h.} \dicPos{zkr} \dicPhraseIS{fyrir hádegi} \dicFlx{adv} \dicDirectTranslationCS{dopoledne} \dicAntonym{e.h.}
\dicEntry[f.hl.] \dicTerm{f.hl.} \dicPos{zkr} \dicPhraseIS{fyrri hluti} \dicDirectTranslationCS{první polovina\,/\addthin část} \dicExampleIS{fyrri hluti dags} \dicExampleCS{první polovina dne}
\dicEntry[f.Kr.] \dicTerm{f.Kr.} \dicPos{zkr} \dicPhraseIS{fyrir Krist, fyrir Krists burð} \dicFlx{adv} \dicDirectTranslationCS{před Kristem, před naším letopočtem} \dicAntonym{e.Kr.}
\dicEntry[faðerni] \dicTerm{fað··erni} \dicIPA{{f}{a}{\textlengthmark}{ð}{\textepsilon}{r}{\textsubring{d}}{n}{\textsci}} \dicPos{n}[2] \dicFlx{(‑s)}[20] \dicDirectTranslationCS{otcovství}
\dicEntry[faðir] \dicTerm{faðir} \dicsymFrequent\  \dicIPA{{f}{a}{\textlengthmark}{ð}{\textsci}{\textsubring{r}}} \dicPos{m}[13] \dicFlx{(föður, feður)}[3] \textbf{1.} \dicDirectTranslationCS{otec, táta} \dicExampleIS{Þau dvelja hjá föðurnum aðra hverja helgi.} \dicExampleCS{Zůstávají u~otce každý druhý víkend.}  \textbf{2.} \dicDirectTranslationCS{Otec} \dicIndirectTranslationCS{(Bůh)}
\dicEntry[faðirvor] \dicTerm{faðir··vor} \dicIPA{{f}{a}{\textlengthmark}{ð}{\textsci}{r}{v}{\textopeno}{\textsubring{r}}} \dicPos{n}[2] \dicFlx{(‑s, ‑)}[5] \dicFieldCat{náb.} \dicSynonym*{drottinleg bæn} \dicDirectTranslationCS{Otčenáš}
\dicEntry[faðma] \dicTerm{faðm|a} \dicsymFrequent\  \dicIPA{{f}{a}{ð}{m}{a}} \dicPos{v}[1] \dicFlx{(‑aði)}[13] \dicFlx{acc} \dicDirectTranslationCS{obejmout, objímat, vzít do náručí} \dicExampleIS{Mamma faðmaði hana.} \dicExampleCS{Máma ji objala.};  \dicPhraseIS{faðma e‑n að sér} \dicDirectTranslationCS{obejmout (koho)} \dicExampleIS{Hann faðmaði mig að sér.} \dicExampleCS{Objal mě.};  \dicIdiom{faðmast}{ \dicPhraseIS{faðmast}} \dicFlx{refl} \dicDirectTranslationCS{obejmout se, objímat se}
\dicEntry[faðmlag] \dicTerm{faðm··|lag} \dicIPA{{f}{a}{ð}{m}{l}{a}{x}} \dicPos{n}[2] \dicFlx{(‑lags, ‑lög)}[8] \dicDirectTranslationCS{objetí, sevření} \dicExampleIS{liggja í faðmlögum} \dicExampleCS{ležet v~objetí}
\dicEntry[faðmur] \dicTerm{faðm|ur} \dicsymFrequent\  \dicIPA{{f}{a}{ð}{m}{\textscy}{\textsubring{r}}} \dicPos{m}[6] \dicFlx{(‑s, ‑ar)}[5] \textbf{1.} \dicSynonym*{útbreiddir armar} \dicDirectTranslationCS{objetí, (otevřená) náruč} \dicExampleIS{hvíla í faðmi hennar} \dicExampleCS{odpočívat v~jejím objetí}  \textbf{2.} \dicLangCat{zast.} \dicSynonym{mælieining} \dicDirectTranslationCS{sáh} \dicIndirectTranslationCS{(167 cm)}
\dicEntry[fafla] \dicTerm{fafla} \dicIPA{{f}{a}{\textsubring{b}}{l}{a}} \dicPos{f}[1] \dicFlx{(föflu, föflur)}[20] \dicFieldCat{lit.} \dicSynonym{söguþráður} \dicDirectTranslationCS{fabule}
\dicEntry[fag] \dicTerm{fag} \dicIPA{{f}{a}{\textlengthmark}{x}} \dicPos{n}[2] \dicFlx{(‑s, fög)}[8] \textbf{1.} \dicSynonym{grein} \dicDirectTranslationCS{obor, odborná oblast} \dicExampleIS{kunna sitt fag} \dicExampleCS{vyznat se ve svém oboru}  \textbf{2.} \dicSynonym*{gluggarammi} \dicDirectTranslationCS{okenní rám}
\dicEntry[faglegur] \dicTerm{fag··legur} \dicIPA{{f}{a}{\textbabygamma}{l}{\textepsilon}{\textbabygamma}{\textscy}{\textsubring{r}}} \dicPos{adj}[1]\dicFlx{}[-8] \dicDirectTranslationCS{odborný, profesionální}
\dicEntry[faglærður] \dicTerm{fag··lærður} \dicIPA{{f}{a}{\textbabygamma}{l}{a}{i}{r}{ð}{\textscy}{\textsubring{r}}} \dicPos{adj}[2]\dicFlx{}[-1] \dicDirectTranslationCS{odborně školený, vyškolený}
\dicEntry[fagmaður] \dicTerm{fag··|maður} \dicIPA{{f}{a}{\textbabygamma}{m}{a}{ð}{\textscy}{\textsubring{r}}} \dicPos{m}[13] \dicFlx{(‑manns, ‑menn)}[2] \dicDirectTranslationCS{odborník, odbornice, profesionál(ka)}
\dicEntry[fagmál] \dicTerm{fag··mál} \dicIPA{{f}{a}{\textbabygamma}{m}{au}{\textsubring{l}}} \dicPos{n}[2] \dicFlx{(‑s)}[2] \dicDirectTranslationCS{odborný jazyk, hantýrka, žargon}
\dicEntry[fagna] \dicTerm{fagn|a} \dicsymFrequent\  \dicIPA{{f}{a}{\r{g}}{n}{a}} \dicPos{v}[1] \dicFlx{(‑aði)}[50] \dicFlx{dat} \dicPhraseIS{fagna e‑m} \dicDirectTranslationCS{uvítat (koho), přivítat (koho)} \dicExampleIS{fagna e‑m blíðlega} \dicExampleCS{přivítat (koho) srdečně};  \dicPhraseIS{fagna e‑u} \dicDirectTranslationCS{uvítat (co), oslavit (co), těšit se z~(čeho)} \dicExampleIS{fagna sumri} \dicExampleCS{uvítat léto}
\dicEntry[fagnaðarboðskapur] \dicTerm{fagnaðar··boð·skap|ur} \dicIPA{{f}\-{a}\-{\r{g}}\-{n}\-{a}\-{ð}\-{a}\-{r}\-{\textsubring{b}}\-{\textopeno}\-{ð}\-{s}\-{\r{g}}\-{a}\-{\textsubring{b}}\-{\textscy}\-{\textsubring{r}}\-} \dicPos{m}[10] \dicFlx{(‑ar)}[15] \textbf{1.} \dicSynonym{gleðitíðindi} \dicDirectTranslationCS{radostné zvěsti\,/\addthin zprávy}  \textbf{2.} \dicFieldCat{náb.} \dicSynonym{guðspjall} \dicDirectTranslationCS{evangelium}
\dicEntry[fagnaðarerindi] \dicTerm{fagnaðar··erindi} \dicIPA{{f}{a}{\r{g}}{n}{a}{ð}{a}{r}{\textepsilon}{r}{\textsci}{n}{\textsubring{d}}{\textsci}} \dicPos{n}[2] \dicFlx{(‑s, ‑)}[14] \textbf{1.} \dicSynonym{gleðitíðindi} \dicDirectTranslationCS{radostné zvěsti\,/\addthin zprávy}  \textbf{2.} \dicFieldCat{náb.} \dicSynonym{guðspjall} \dicDirectTranslationCS{evangelium}
\dicEntry[fagnaðarfundur] \dicTerm{fagnaðar··fund|ur} \dicIPA{{f}{a}{\r{g}}{n}{a}{ð}{a}{\textsubring{r}}{f}{\textscy}{n}{\textsubring{d}}{\textscy}{\textsubring{r}}} \dicPos{m}[10] \dicFlx{(‑ar, ‑ir)}[4] \dicDirectTranslationCS{oslava (opětovného) setkání\,/\addthin shledání}
\dicEntry[fagnaður] \dicTerm{fag··nað|ur} \dicIPA{{f}{a}{\r{g}}{n}{a}{ð}{\textscy}{\textsubring{r}}} \dicPos{m}[10] \dicFlx{(‑ar, ‑ir)}[8] \dicDirectTranslationCS{oslava, veselí}
\dicEntry[fagorð] \dicTerm{fag··orð} \dicIPA{{f}{a}{\textlengthmark}{\textbabygamma}{\textopeno}{r}{\texttheta}} \dicPos{n}[2] \dicFlx{(‑s, ‑)}[5] \dicDirectTranslationCS{odborný název, termín}
\dicEntry[fagott] \dicTerm{fagott} \dicIPA{{f}{a}{\textlengthmark}{\r{g}}{\textopeno}{h}{\textsubring{d}}} \dicPos{n}[2] \dicFlx{(‑s, ‑)}[5] \dicFieldCat{hud.} \dicDirectTranslationCS{fagot}
\dicEntry[fagtímarit] \dicTerm{fag··tíma·rit} \dicIPA{{f}{a}{x}{t\smash{\textsuperscript{h}}}{i}{m}{a}{r}{\textsci}{\textsubring{d}}} \dicPos{n}[2] \dicFlx{(‑s, ‑)}[5] \dicDirectTranslationCS{odborný časopis}
\dicEntry[fagur] \dicTerm{fagur} \dicsymFrequent\  \dicIPA{{f}{a}{\textlengthmark}{\textbabygamma}{\textscy}{\textsubring{r}}} \dicPos{adj}[9] \dicFlx{(f fögur)}[6] \dicDirectTranslationCS{krásný, úchvatný, nádherný} \dicExampleIS{fögur kona} \dicExampleCS{krásná žena};  \dicPhraseIS{fara fögrum orðum um e‑n} \dicDirectTranslationCS{mluvit o~(kom) v~superlativech};  \dicPhraseIS{lofa öllu fögru} \dicLangCat{přen.} \dicDirectTranslationCS{slibovat hory doly}
\dicEntry[fagurblár] \dicTerm{fagur··blár} \dicIPA{{f}{a}{\textlengthmark}{\textbabygamma}{\textscy}{r}{\textsubring{b}}{l}{au}{\textsubring{r}}} \dicPos{adj}[4]\dicFlx{}[-3] \dicDirectTranslationCS{azurový}
\dicEntry[fagurbókmenntir] \dicTerm{fagur··bók·menntir} \dicIPA{{f}{a}{\textlengthmark}{\textbabygamma}{\textscy}{r}{\textsubring{b}}{ou}{\r{g}}{m}{\textepsilon}{\textsubring{n}}{\textsubring{d}}{\textsci}{\textsubring{r}}} \dicPos{f}[7] \dicFlx{pl}[2] \dicDirectTranslationCS{beletrie, krásná literatura}
\dicEntry[fagurfífill] \dicTerm{fagur··fíf|ill} \dicIPA{{f}{a}{\textlengthmark}{\textbabygamma}{\textscy}{\textsubring{r}}{f}{i}{v}{\textsci}{\textsubring{d}}{\textsubring{l}}} \dicPos{m}[6] \dicFlx{(‑ils, ‑lar)}[35] \dicFieldCat{bot.} \dicDirectTranslationCS{sedmikráska, sedmikráska chudobka} \textit{(l.~{\textLA{Bellis perennis}})}  \dicsymPhoto\ 
\dicFigure{744.jpg}{Fagurfífill}{Fagurfífill - Zicha Ondřej, Biolib, Copyright/CC-BY-NC}
\dicEntry[fagurfræði] \dicTerm{fagur··fræð|i} \dicIPA{{f}{a}{\textlengthmark}{\textbabygamma}{\textscy}{\textsubring{r}}{f}{r}{a}{i}{ð}{\textsci}} \dicPos{f}[3] \dicFlx{(‑i)}[3] \dicDirectTranslationCS{estetika}
\dicEntry[fagurfræðilegur] \dicTerm{fagur·fræði··legur} \dicIPA{{f}{a}{\textlengthmark}{\textbabygamma}{\textscy}{\textsubring{r}}{f}{r}{a}{i}{ð}{\textsci}{l}{\textepsilon}{\textbabygamma}{\textscy}{\textsubring{r}}} \dicPos{adj}[1]\dicFlx{}[-8] \dicDirectTranslationCS{estetický}
\dicEntry[fagurkeri] \dicTerm{fagur··ker|i} \dicIPA{{f}{a}{\textlengthmark}{\textbabygamma}{\textscy}{\textsubring{r}}{c\smash{\textsuperscript{h}}}{\textepsilon}{r}{\textsci}} \dicPos{m}[1] \dicFlx{(‑a, ‑ar)}[1] \dicDirectTranslationCS{estét(ka)}
\dicEntry[fagurrauður] \dicTerm{fagur··rauður} \dicIPA{{f}{a}{\textlengthmark}{\textbabygamma}{\textscy}{r}{\oe i}{ð}{\textscy}{\textsubring{r}}} \dicPos{adj}[2]\dicFlx{}[-6] \dicDirectTranslationCS{karmínový, sytě červený}
\dicEntry[faktúra] \dicTerm{faktúr|a} \dicIPA{{f}{a}{x}{\textsubring{d}}{u}{r}{a}} \dicPos{f}[1] \dicFlx{(‑u, ‑ur)}[7] \dicFieldCat{ekon.} \dicDirectTranslationCS{faktura}
\dicEntry[fala] \dicTerm{fal|a} \dicIPA{{f}{a}{\textlengthmark}{l}{a}} \dicPos{v}[1] \dicFlx{(‑aði)}[13] \dicFlx{acc} \dicLangCat{zast.} \dicDirectTranslationCS{požádat};  \dicIdiom{falast}[eftir]{ \dicPhraseIS{falast eftir e‑u}} \dicFlx{refl} \dicDirectTranslationCS{požádat o~(co) (práci ap.)}
\dicEntry[falda] \dicTerm{fald|a} \dicIPA{{f}{a}{l}{\textsubring{d}}{a}} \dicPos{v}[1] \dicFlx{(‑aði)}[13] \dicFlx{acc} \dicDirectTranslationCS{olemovat, obroubit} \dicExampleIS{falda dúkinn} \dicExampleCS{olemovat ubrus}
\dicEntry[Faldafeykir] \dicTerm{Falda··feyk|ir} \dicIPA{{f}{a}{l}{\textsubring{d}}{a}{f}{ei}{\r{\textObardotlessj}}{\textsci}{\textsubring{r}}} \dicPos{m}[7] \dicFlx{(‑is)}[3] \dicFlx{prop} \dicLink{Hurðaskellir}
\dicEntry[faldi] \dicTerm{faldi} \dicIPA{{f}{a}{l}{\textsubring{d}}{\textsci}} \dicPos{v} \dicFlx{ind pf sg 1 pers} \dicLink{fela\smash{\textsuperscript{1}}}
\dicEntry[faldur] \dicTerm{fald|ur} \dicIPA{{f}{a}{l}{\textsubring{d}}{\textscy}{\textsubring{r}}} \dicPos{m}[6] \dicFlx{(‑s, ‑ar)}[5] \textbf{1.} \dicSynonym*{innafbrot} \dicDirectTranslationCS{obruba, lem} \dicExampleIS{setja fald á pilsið} \dicExampleCS{obroubit sukni}  \textbf{2.} \dicSynonym*{höfuðbúnaður} \dicDirectTranslationCS{diadém, ozdoba hlavy}
\dicEntry[falið] \dicTerm{falið} \dicIPA{{f}{a}{\textlengthmark}{l}{\textsci}{\texttheta}} \dicPos{v} \dicFlx{supin} \textbf{1.} \dicLink{fela\smash{\textsuperscript{1}}}  \textbf{2.} \dicLink{fela\smash{\textsuperscript{2}}}
\dicEntry[Falklandseyjar] \dicTerm{Falk·lands··eyjar} \dicIPA{{f}{a}{\textsubring{l}}{\r{g}}{l}{a}{n}{\textsubring{d}}{s}{ei}{j}{a}{\textsubring{r}}} \dicPos{f}[4] \dicFlx{pl}[15] \dicFieldCat{geog.} \dicDirectTranslationCS{Falklandy, Malvíny} \dicIndirectTranslationCS{(ostrovy v~jižním Atlantiku)}
\dicEntry[fall] \dicTerm{fall} \dicsymFrequent\  \dicIPA{{f}{a}{\textsubring{d}}{\textsubring{l}}} \dicPos{n}[2] \dicFlx{(‑s, föll)}[8] \textbf{1.} \dicSynonym{hrun} \dicDirectTranslationCS{pád, upadnutí}  \textbf{2.} \dicFieldCat{jaz.} \dicDirectTranslationCS{pád} \dicExampleIS{eignarfall} \dicExampleCS{genitiv}  \textbf{3.} \dicFieldCat{mat.} \dicDirectTranslationCS{funkce};  \dicIdiom{fall}{ \dicPhraseIS{e‑að riðar til falls}} \dicDirectTranslationCS{(co) se žene do záhuby};  \dicPhraseIS{e‑að varð e‑m að falli} \dicLangCat{přen.} \dicDirectTranslationCS{(co komu) přivodilo pád}; { \dicPhraseIS{í besta\,/\addthin versta falli}} \dicFlx{adv} \dicDirectTranslationCS{v~nejlepším\,/\addthin nejhorším případě};  \dicPhraseIS{með fyrra fallinu} \dicFlx{adv} \dicDirectTranslationCS{ve správný čas}
\dicEntry[falla] \dicTerm{falla} \dicsymFrequent\  \dicIPA{{f}{a}{\textsubring{d}}{l}{a}} \dicPos{v}[6] \dicFlx{(fell, féll, féllum, félli, fallið)}[54] \textbf{1.} \dicSynonym{detta} \dicDirectTranslationCS{(s)padnout, upadnout, padat} \dicExampleIS{falla á gólfið} \dicExampleCS{upadnout na podlahu}  \textbf{2.} \dicSynonym*{vera drepinn} \dicDirectTranslationCS{padnout, být zabit} \dicExampleIS{falla í stríði} \dicExampleCS{padnout ve válce}  \textbf{3.} \dicSynonym{streyma} \dicDirectTranslationCS{vtékat, vlévat se} \dicExampleIS{Áin fellur til sjávar.} \dicExampleCS{Řeka vtéká do moře.}  \textbf{4.} \dicPhraseIS{e‑m fellur e‑að vel} \dicFlx{impers} \dicSynonym*{líka} \dicDirectTranslationCS{(co) se (komu) líbí, (co) má (kdo) rád} \dicExampleIS{Mér fellur þetta vel.} \dicExampleCS{Líbí se mi to.}  \textbf{5.} \dicSynonym{hæfa\smash{\textsuperscript{2}}} \dicDirectTranslationCS{padnout (šaty ap.), sedět, pasovat} \dicExampleIS{Stykkið fellur vel.} \dicExampleCS{Ten kus padne dobře.}  \textbf{6.} \dicDirectTranslationCS{padnout, padat (stín ap.)}  \textbf{7.} \dicDirectTranslationCS{propadnout, neuspět (u~zkoušky ap.)} \dicExampleIS{falla í málfræði} \dicExampleCS{neuspět v~gramatice};  \dicIdiom{falla}[að]{ \dicPhraseIS{e‑að fellur (vel) að e‑u}} \dicDirectTranslationCS{(co) padne (dobře) k~(čemu)}; { \dicPhraseIS{það fellur að}} \dicFlx{impers} \dicDirectTranslationCS{nastává příliv};  \dicIdiom{falla}[á]{ \dicPhraseIS{falla á prófi}} \dicSynonym*{standast ekki próf} \dicDirectTranslationCS{propadnout\,/\addthin neuspět u~zkoušky (na VŠ ap.)} \dicExampleIS{Nemandinn féll á prófi.} \dicExampleCS{Student neuspěl u~zkoušky.}; { \dicPhraseIS{það fellur á e‑ð}} \dicFlx{impers} \dicDirectTranslationCS{(co) oxiduje};  \dicIdiom{falla}[frá]{ \dicPhraseIS{falla frá}} \dicDirectTranslationCS{zemřít, skonat}; { \dicPhraseIS{falla frá e‑u}} \dicDirectTranslationCS{upustit od (čeho)} \dicExampleIS{falla frá áformum sínum} \dicExampleCS{upustit od svých záměrů};  \dicIdiom{falla}[fyrir]{ \dicPhraseIS{falla fyrir e‑u\,/\addthin e‑m}} \dicDirectTranslationCS{oblíbit si (ihned) (co\,/\addthin koho)};  \dicIdiom{falla}[niður]{ \dicPhraseIS{e‑að fellur niður}} \dicDirectTranslationCS{(co) je zrušeno, (co) odpadá} \dicExampleIS{Æfing féll niður.} \dicExampleCS{Cvičení odpadlo.};  \dicIdiom{falla}[saman]{ \dicPhraseIS{e‑að fellur saman}} \dicDirectTranslationCS{(co) se hroutí, (co) se řítí};  \dicIdiom{falla}[til]{ \dicPhraseIS{það fellur e‑að til}} \dicFlx{impers} \dicDirectTranslationCS{(co) se stane, (co) se přihodí};  \dicIdiom{falla}[undir]{ \dicPhraseIS{e‑að fellur undir e‑n}} \dicDirectTranslationCS{(co) spadá pod (koho) (pravomoc ap.)};  \dicIdiom{falla}[út]{ \dicPhraseIS{það fellur út}} \dicFlx{impers} \dicDirectTranslationCS{nastává odliv};  \dicIdiom{falla}[við]{ \dicPhraseIS{e‑m fellur (vel) við e‑n}} \dicFlx{impers} \dicDirectTranslationCS{(kdo) má (koho) rád, (kdo) s~(kým) (dobře) vychází} \dicExampleIS{Mér fellur vel við nágrannann.} \dicExampleCS{Vycházím dobře se sousedem.};  \dicIdiom{fallast}{ \dicPhraseIS{e‑m fallast hendur}} \dicFlx{impers refl} \dicDirectTranslationCS{(kdo) je ztracen, (kdo) si neví rady};  \dicIdiom{fallast}[á]{ \dicPhraseIS{fallast á e‑ð}} \dicFlx{refl} \dicSynonym{samþykkja} \dicDirectTranslationCS{přistoupit na (co), souhlasit s~(čím)};  \dicIdiom{fallast}[í]{ \dicPhraseIS{fallast í faðma}} \dicFlx{refl} \dicDirectTranslationCS{obejmout se, padnout si do náručí}
\dicEntry[fallbeyging] \dicTerm{fall··beyg·ing} \dicIPA{{f}{a}{\textsubring{d}}{\textsubring{l}}{\textsubring{b}}{ei}{j}{i}{\ng}{\r{g}}} \dicPos{f}[4] \dicFlx{(‑ar)}[7] \dicFieldCat{jaz.} \dicDirectTranslationCS{skloňování, deklinace}
\dicEntry[fallbeygja] \dicTerm{fall··beyg|ja} \dicIPA{{f}{a}{\textsubring{d}}{\textsubring{l}}{\textsubring{b}}{ei}{j}{a}} \dicPos{v}[2] \dicFlx{(‑ði, ‑t)}[89] \dicFlx{acc} \dicFieldCat{jaz.} \dicDirectTranslationCS{skloňovat}
\dicEntry[fallbyssa] \dicTerm{fall··byss|a} \dicIPA{{f}{a}{\textsubring{d}}{\textsubring{l}}{\textsubring{b}}{\textsci}{s}{a}} \dicPos{f}[1] \dicFlx{(‑u, ‑ur)}[19] \dicDirectTranslationCS{kanón, dělo}
\dicEntry[fallega] \dicTerm{falleg|a} \dicsymFrequent\  \dicIPA{{f}{a}{\textsubring{d}}{l}{\textepsilon}{\textbabygamma}{a}} \dicPos{adv} \dicFlx{(comp ‑ar, sup ‑ast)} \dicDirectTranslationCS{hezky, pěkně} \dicExampleIS{Rósarunninn blómstrar svo fallega.} \dicExampleCS{Růžový keř tak krásně kvete.}
\dicEntry[fallegur] \dicTerm{fallegur} \dicsymFrequent\  \dicIPA{{f}{a}{\textsubring{d}}{l}{\textepsilon}{\textbabygamma}{\textscy}{\textsubring{r}}} \dicPos{adj}[1]\dicFlx{}[-8] \dicSynonym{fagur} \dicDirectTranslationCS{hezký, pěkný, krásný} \dicExampleIS{falleg stúlka} \dicExampleCS{hezká dívka}
\dicEntry[fallending] \dicTerm{fall··end·ing} \dicIPA{{f}{a}{\textsubring{d}}{l}{\textepsilon}{n}{\textsubring{d}}{i}{\ng}{\r{g}}} \dicPos{f}[4] \dicFlx{(‑ar, ‑ar)}[5] \dicFieldCat{jaz.} \dicDirectTranslationCS{pádová\,/\addthin deklinační koncovka}
\dicEntry[fallháttur] \dicTerm{fall··|háttur} \dicIPA{{f}{a}{\textsubring{d}}{\textsubring{l}}{h}{au}{h}{\textsubring{d}}{\textscy}{\textsubring{r}}} \dicPos{m}[12] \dicFlx{(‑háttar, ‑hættir)}[7] \dicFieldCat{jaz.} \dicDirectTranslationCS{neurčitý způsob}
\dicEntry[fallhlíf] \dicTerm{fall··hlíf} \dicIPA{{f}{a}{\textsubring{d}}{\textsubring{l}}{i}{f}} \dicPos{f}[4] \dicFlx{(‑ar, ‑ar)}[1] \dicFieldCat{let.} \dicDirectTranslationCS{padák}
\dicEntry[fallhlífarstökkvari] \dicTerm{fall·hlífar··stökkv·ar|i} \dicIPA{{f}\-{a}\-{\textsubring{d}}\-{\textsubring{l}}\-{i}\-{v}\-{a}\-{\textsubring{r}}\-{s}\-{\textsubring{d}}\-{\oe}\-{h}\-{\r{g}}\-{v}\-{a}\-{r}\-{\textsci}\-} \dicPos{m}[1] \dicFlx{(‑a, ‑ar)}[13] \dicDirectTranslationCS{parašutista, parašutistka}
\dicEntry[fallið] \dicTerm{fallið} \dicIPA{{f}{a}{\textsubring{d}}{l}{\textsci}{\texttheta}} \dicPos{v} \dicFlx{supin} \dicLink{falla}
\dicEntry[fallinn] \dicTerm{fallinn} \dicsymFrequent\  \dicIPA{{f}{a}{\textsubring{d}}{l}{\textsci}{\textsubring{n}}} \dicPos{adj}[6]\dicFlx{}[-4] \textbf{1.} \dicPhraseIS{vera (vel) til þess fallinn að (gera e‑ð)} \dicDirectTranslationCS{být vhodný k~(dělání (čeho))}  \textbf{2.} \dicDirectTranslationCS{propadlý (u~zkoušky), neúspěšný} \dicExampleIS{vera fallinn í stærðfræði} \dicExampleCS{neuspět u~zkoušky z~matematiky}  \textbf{3.} \dicSynonym*{drepinn} \dicDirectTranslationCS{padlý, zabitý} \dicExampleIS{fallinn hermaður} \dicExampleCS{padlý voják}
\dicEntry[fallorð] \dicTerm{fall··orð} \dicIPA{{f}{a}{\textsubring{d}}{l}{\textopeno}{r}{\texttheta}} \dicPos{n}[2] \dicFlx{(‑s, ‑)}[5] \dicFieldCat{jaz.} \dicDirectTranslationCS{ohebné slovo}
\dicEntry[fals] \dicTerm{fals} \dicIPA{{f}{a}{l}{s}} \dicPos{n}[2] \dicFlx{(‑, föls)}[23] \textbf{1.} \dicSynonym*{óhreinlyndi} \dicDirectTranslationCS{faleš, falešnost, neupřímnost}  \textbf{2.} \dicSynonym{fölsun} \dicDirectTranslationCS{padělání, falšování}
\dicEntry[falsa] \dicTerm{fals|a} \dicIPA{{f}{a}{l}{s}{a}} \dicPos{v}[1] \dicFlx{(‑aði)}[13] \dicFlx{acc} \dicDirectTranslationCS{padělat, (z)falšovat} \dicExampleIS{falsa undirskriftina} \dicExampleCS{zfalšovat podpis}
\dicEntry[falsaður] \dicTerm{fals··aður} \dicIPA{{f}{a}{l}{s}{a}{ð}{\textscy}{\textsubring{r}}} \dicPos{adj}[3] \dicFlx{(f fölsuð)}[1] \dicDirectTranslationCS{falešný, (z)falšovaný} \dicExampleIS{falsaðir peningar} \dicExampleCS{falešné peníze}
\dicEntry[falsari] \dicTerm{fals··ar|i} \dicIPA{{f}{a}{l}{s}{a}{r}{\textsci}} \dicPos{m}[1] \dicFlx{(‑a, ‑ar)}[10] \dicSynonym{svikari} \dicDirectTranslationCS{padělatel(ka)} \dicExampleIS{peningafalsari} \dicExampleCS{padělatel peněz}
\dicEntry[falsetta] \dicTerm{falsett|a} \dicIPA{{f}{a}{l}{s}{\textepsilon}{h}{\textsubring{d}}{a}} \dicPos{f}[1] \dicFlx{(‑u, ‑ur)}[7] \dicFieldCat{hud.} \dicDirectTranslationCS{fistule, falzet}
\dicEntry[falskur] \dicTerm{falskur} \dicIPA{{f}{a}{l}{s}{\r{g}}{\textscy}{\textsubring{r}}} \dicPos{adj}[1] \dicFlx{(f fölsk)}[2] \textbf{1.} \dicSynonym{svikull} \dicDirectTranslationCS{falešný, neupřímný, pokrytecký} \dicExampleIS{falskur svipur} \dicExampleCS{falešný výraz}  \textbf{2.} \dicDirectTranslationCS{padělaný, zfalšovaný, falešný} \dicExampleIS{falskur peningur} \dicExampleCS{padělané peníze}  \textbf{3.} \dicFieldCat{hud.} \dicDirectTranslationCS{falešný, nečistý (zpěv ap.)} \dicExampleIS{falskur tónn} \dicExampleCS{falešný tón};  \dicPhraseIS{falskar tennur} \dicDirectTranslationCS{falešné zuby, umělý chrup}
\dicEntry[falur] \dicTerm{falur} \dicIPA{{f}{a}{\textlengthmark}{l}{\textscy}{\textsubring{r}}} \dicPos{adj}[1] \dicFlx{(f föl)}[2] \dicSynonym*{til sölu} \dicDirectTranslationCS{(jsoucí) na prodej} \dicExampleIS{Hesturinn er falur.} \dicExampleCS{Kůň je na prodej.}
\dicEntry[fanatískur] \dicTerm{fanatískur} \dicIPA{{f}{a}{\textlengthmark}{n}{a}{t\smash{\textsuperscript{h}}}{i}{s}{\r{g}}{\textscy}{\textsubring{r}}} \dicPos{adj}[1]\dicFlx{}[-1] \dicDirectTranslationCS{fanatický}
\dicEntry[fang] \dicTerm{fang} \dicsymFrequent\  \dicIPA{{f}{au}{\ng}{\r{g}}} \dicPos{n}[2] \dicFlx{(‑s, föng)}[8] \textbf{1.} \dicSynonym{faðmur} \dicDirectTranslationCS{náruč};  \dicPhraseIS{taka e‑n í fangið} \dicDirectTranslationCS{vzít (koho) do náruče}  \textbf{2.} \dicFieldCat{biol.} \dicSynonym{fóstur} \dicDirectTranslationCS{zárodek, plod};  \dicPhraseIS{fá fang} \dicDirectTranslationCS{otelit se}  \textbf{3.} \dicPhraseIS{föng} \dicFlx{pl} \dicSynonym{aðföng} \dicDirectTranslationCS{strava, potrava}  \textbf{4.} \dicPhraseIS{föng} \dicFlx{pl} \dicSynonym*{möguleikar} \dicDirectTranslationCS{možnost};  \dicPhraseIS{eftir því sem föng eru á} \dicDirectTranslationCS{jak nejlépe to jde};  \dicPhraseIS{eftir föngum} \dicFlx{adv} \dicDirectTranslationCS{podle možností};  \dicIdiom{fang}{ \dicPhraseIS{eiga fullt í fangi með e‑ð}} \dicDirectTranslationCS{mít plné ruce práce s~(čím)}; { \dicPhraseIS{færast e‑ð í fang}} \dicFlx{refl} \dicLangCat{přen.} \dicDirectTranslationCS{brát si (co) k~srdci};  \dicPhraseIS{hafa storminn í fangið} \dicDirectTranslationCS{jít proti větru, mít protivítr}
\dicEntry[fanga] \dicTerm{fang|a} \dicIPA{{f}{au}{\ng}{\r{g}}{a}} \dicPos{v}[1] \dicFlx{(‑aði)}[13] \dicFlx{acc} \dicSynonym{grípa} \dicDirectTranslationCS{uchopit, zachytit} \dicExampleIS{fanga atburðinn á ljósmynd} \dicExampleCS{zachytit událost na fotografii}
\dicEntry[fangaklefi] \dicTerm{fanga··klef|i} \dicIPA{{f}{au}{\ng}{\r{g}}{a}{k\smash{\textsuperscript{h}}}{l}{\textepsilon}{v}{\textsci}} \dicPos{m}[1] \dicFlx{(‑a, ‑ar)}[1] \dicDirectTranslationCS{vězeňská cela}
\dicEntry[fangamark] \dicTerm{fanga··|mark} \dicIPA{{f}{au}{\ng}{\r{g}}{a}{m}{a}{\textsubring{r}}{\r{g}}} \dicPos{n}[2] \dicFlx{(‑marks, ‑mörk)}[8] \dicDirectTranslationCS{monogram, iniciála}
\dicEntry[fangavist] \dicTerm{fanga··vist} \dicIPA{{f}{au}{\ng}{\r{g}}{a}{v}{\textsci}{s}{\textsubring{d}}} \dicPos{f}[7] \dicFlx{(‑ar)}[3] \dicDirectTranslationCS{pobyt ve vězení, vězeňský pobyt}
\dicEntry[fangavörður] \dicTerm{fanga··|vörður} \dicIPA{{f}{au}{\ng}{\r{g}}{a}{v}{\oe}{r}{ð}{\textscy}{\textsubring{r}}} \dicPos{m}[11] \dicFlx{(‑varðar, ‑verðir)}[5] \dicDirectTranslationCS{vězeňský dozorce, vězeňská dozorkyně, bachař(ka), žalářník, žalářnice}
\dicEntry[fangelsa] \dicTerm{fang··els|a} \dicIPA{{f}{au}{\textltailn}{\r{\textObardotlessj}}{\textepsilon}{l}{s}{a}} \dicPos{v}[1] \dicFlx{(‑aði)}[1] \dicFlx{acc} \dicDirectTranslationCS{(u)věznit, vsadit do vězení}
\dicEntry[fangelsi] \dicTerm{fang··elsi} \dicIPA{{f}{au}{\textltailn}{\r{\textObardotlessj}}{\textepsilon}{l}{s}{\textsci}} \dicPos{n}[2] \dicFlx{(‑s, ‑)}[14] \dicDirectTranslationCS{vězení, věznice} \dicExampleIS{sitja lengi í fangelsi} \dicExampleCS{sedět dlouho ve vězení}
\dicEntry[fangelsisvist] \dicTerm{fangelsis··vist} \dicIPA{{f}{au}{\textltailn}{\r{\textObardotlessj}}{\textepsilon}{l}{s}{\textsci}{s}{v}{\textsci}{s}{\textsubring{d}}} \dicPos{f}[7] \dicFlx{(‑ar)}[3] \dicDirectTranslationCS{pobyt ve vězení}
\dicEntry[fangi] \dicTerm{fang|i} \dicsymFrequent\  \dicIPA{{f}{au}{\textltailn}{\r{\textObardotlessj}}{\textsci}} \dicPos{m}[1] \dicFlx{(‑a, ‑ar)}[8] \dicDirectTranslationCS{vězeň(kyně), trestanec, trestankyně} \dicExampleIS{pólitískur fangi} \dicExampleCS{politický vězeň};  \dicPhraseIS{taka e‑n til fanga} \dicDirectTranslationCS{uvěznit (koho)}
\dicEntry[fann] \dicTerm{fann} \dicIPA{{f}{a}{n}{\textlengthmark}} \dicPos{v} \dicFlx{ind pf sg 1 pers} \dicLink{finna}
\dicEntry[fannar] \dicTerm{fannar} \dicIPA{{f}{a}{n}{\textlengthmark}{a}{\textsubring{r}}} \dicPos{f} \dicFlx{sg gen} \dicLink{fönn}
\dicEntry[fannhvítur] \dicTerm{fann··hvítur} \dicIPA{{f}{a}{n}{k\smash{\textsuperscript{h}}}{v}{i}{\textsubring{d}}{\textscy}{\textsubring{r}}} \dicPos{adj}[1]\dicFlx{}[-6] \dicDirectTranslationCS{sněhobílý}
\dicEntry[fantasía] \dicTerm{fantasí|a} \dicIPA{{f}{a}{\textsubring{n}}{\textsubring{d}}{a}{s}{i}{j}{a}} \dicPos{f}[1] \dicFlx{(‑u, ‑ur)}[7] \textbf{1.} \dicSynonym{hugarflug} \dicDirectTranslationCS{fantazie, obrazotvornost}  \textbf{2.} \dicSynonym{ímyndun} \dicDirectTranslationCS{výtvor obrazotvornosti, smyšlenka}  \textbf{3.} \dicFieldCat{lit.} \dicSynonym*{furðusaga} \dicDirectTranslationCS{fantasy}  \textbf{4.} \dicFieldCat{hud.} \dicDirectTranslationCS{fantazie}
\dicEntry[fantaskapur] \dicTerm{fanta··skap|ur} \dicIPA{{f}{a}{\textsubring{n}}{\textsubring{d}}{a}{s}{\r{g}}{a}{\textsubring{b}}{\textscy}{\textsubring{r}}} \dicPos{m}[10] \dicFlx{(‑ar)}[15] \dicDirectTranslationCS{surovost, zlotřilost}
\dicEntry[fantur] \dicTerm{fant|ur} \dicIPA{{f}{a}{\textsubring{n}}{\textsubring{d}}{\textscy}{\textsubring{r}}} \dicPos{m}[6] \dicFlx{(‑s, ‑ar)}[5] \textbf{1.} \dicSynonym{þorpari} \dicDirectTranslationCS{násilník, surovec}  \textbf{2.} \dicSynonym{krús} \dicDirectTranslationCS{hrnek}
\dicEntry[far] \dicTerm{far} \dicsymFrequent\  \dicIPA{{f}{a}{\textlengthmark}{\textsubring{r}}} \dicPos{n}[2] \dicFlx{(‑s, för)}[8] \textbf{1.} \dicSynonym{flutningur} \dicDirectTranslationCS{jízda, svezení, projížďka} \dicExampleIS{fá far með bílnum} \dicExampleCS{projet se autem}  \textbf{2.} \dicSynonym{spor} \dicDirectTranslationCS{stopa, otisk} \dicExampleIS{far eftir hamarinn} \dicExampleCS{otisk po kladivu}  \textbf{3.} \dicSynonym{bátur} \dicDirectTranslationCS{plavidlo, loď}  \textbf{4.} \dicSynonym{atferli} \dicDirectTranslationCS{způsob, chování};  \dicPhraseIS{gera sér far um e‑ð} \dicDirectTranslationCS{předsevzít si (co)};  \dicPhraseIS{í fari e‑rs} \dicDirectTranslationCS{v~chování (koho)};  \dicPhraseIS{frá fornu fari} \dicFlx{adv} \dicDirectTranslationCS{odpradávna, od dávných časů};  \dicPhraseIS{það sækir í sama farið} \dicLangCat{přen.} \dicDirectTranslationCS{dostává se to zpět do zajetých kolejí, vše je opět při starém}
\dicEntry[fara] \dicTerm{fara} \dicsymFrequent\  \dicIPA{{f}{a}{\textlengthmark}{r}{a}} \dicPos{v}[8] \dicFlx{(fer, fór, fórum, færi, farið)}[3] \textbf{1.} \dicDirectTranslationCS{jet, jít} \dicExampleIS{fara í bíó} \dicExampleCS{jít do kina};  \dicPhraseIS{fara leiðar sinnar} \dicDirectTranslationCS{(vy)jít, (vy)jet}  \textbf{2.} \dicDirectTranslationCS{(za)chovat se, jednat} \dicExampleIS{farðu varlega!} \dicExampleCS{buď opatrný!};  \dicPhraseIS{fara huldu höfði} \dicDirectTranslationCS{chovat se nenápadně};  \dicPhraseIS{fara sér hægt} \dicDirectTranslationCS{brát to s~klidem, nespěchat}  \textbf{3.} \dicSynonym{enda\smash{\textsuperscript{1}}} \dicDirectTranslationCS{skončit, dopadnout} \dicExampleIS{Allt fór vel að lokum.} \dicExampleCS{Nakonec všechno dobře dopadlo.}  \textbf{4.} \dicSynonym*{eiga við} \dicDirectTranslationCS{hodit se, sedět (šaty ap.), pasovat} \dicExampleIS{Kjóllinn fer illa.} \dicExampleCS{Šaty nepasují.}  \textbf{5.} \dicSynonym{hefja} \dicDirectTranslationCS{začít, začínat} \dicIndirectTranslationCS{(užívané jako pomocné sloveso s~infinitivem -- vyjadřuje začátek akce)} \dicExampleIS{fara að lesa} \dicExampleCS{začít si číst}  \textbf{6.} \dicIndirectTranslationCS{(užívané jako pomocné sloveso s~příčestím přítomným -- vyjadřuje vývoj\,/\addthin pokrok)} \dicExampleIS{Honum fór batnandi.} \dicExampleCS{Udělalo se mu líp.}  \textbf{7.} \dicDirectTranslationCS{jít, poslat} \dicExampleIS{fari e‑r norður og niður} \dicExampleCS{ať už (kdo) zmizí z~očí};  \dicPhraseIS{fari e‑r til fjandans} \dicDirectTranslationCS{ať jde (kdo) k~čertu};  \dicIdiom{fara}[að]{ \dicPhraseIS{fara (svona) að e‑u}} \dicDirectTranslationCS{přistoupit k~(čemu) (určitým způsobem), provést (co) (určitým způsobem)}; { \dicPhraseIS{fara að e‑u}} \dicDirectTranslationCS{řídit se (čím)} \dicExampleIS{fara að reglunum} \dicExampleCS{řídit se pravidly}; { \dicPhraseIS{fara vel að e‑m}} \dicDirectTranslationCS{přistupovat ke (komu) s~úctou\,/\addthin respektem};  \dicIdiom{fara}[af]{ \dicPhraseIS{fara af stað}} \dicDirectTranslationCS{vyjít, vyjet, vyrazit (na cestu ap.)};  \dicIdiom{fara}[aftur]{ \dicPhraseIS{e‑m fer aftur}} \dicFlx{impers} \dicDirectTranslationCS{(komu) se přitěžuje}; { \dicPhraseIS{e‑u fer aftur}} \dicFlx{impers} \dicDirectTranslationCS{(co) upadá, (co) se zhoršuje};  \dicIdiom{fara}[á]{ \dicPhraseIS{fara á hausinn\,/\addthin höfuðið}} \dicDirectTranslationCS{zbankrotovat, zkrachovat}; { \dicPhraseIS{fara á skauta}} \dicDirectTranslationCS{bruslit}; { \dicPhraseIS{fara á skíði}} \dicDirectTranslationCS{lyžovat}; { \dicPhraseIS{það fer vel á e‑u}} \dicFlx{impers} \dicDirectTranslationCS{(co) se hodí, (co) pasuje};  \dicIdiom{fara}[á bak við]{ \dicPhraseIS{fara á bak við e‑n\,/\addthin e‑ð}} \dicLangCat{přen.} \dicDirectTranslationCS{obejít (koho\,/\addthin co), obcházet (koho\,/\addthin co)};  \dicIdiom{fara}[á eftir]{ \dicPhraseIS{fara á eftir e‑m}} \dicDirectTranslationCS{jít\,/\addthin jet za (kým)};  \dicIdiom{fara}[á mis við]{ \dicPhraseIS{fara á mis við e‑ð}} \dicDirectTranslationCS{ztratit šanci na (co)};  \dicIdiom{fara}[eftir]{ \dicPhraseIS{fara eftir e‑u}} {\textbf{a.}} \dicSynonym{sækja} \dicDirectTranslationCS{donést (co), přinést (co), (do)jít pro (co)} \dicExampleIS{fara eftir sykri út í búð} \dicExampleCS{dojít pro cukr do obchodu};  {\textbf{b.}} \dicDirectTranslationCS{záležet na (čem), záviset na (čem)} \dicExampleIS{Það fer eftir merkingu sagnarinnar.} \dicExampleCS{To záleží na významu slovesa.};  {\textbf{c.}} \dicDirectTranslationCS{řídit se (čím), jednat podle (čeho)};  \dicIdiom{fara}[fjarri]{ \dicPhraseIS{e‑u fer fjarri}} \dicFlx{impers} \dicDirectTranslationCS{(co) je daleko od pravdy};  \dicIdiom{fara}[fram]{ \dicPhraseIS{e‑að fer fram}} \dicDirectTranslationCS{(co) probíhá, (co) se uskutečňuje}; { \dicPhraseIS{e‑m fer fram}} \dicFlx{impers} \dicDirectTranslationCS{(kdo) postupuje kupředu, (kdo) dělá krok kupředu, (kdo) se zlepšuje};  \dicIdiom{fara}[fram á]{ \dicPhraseIS{fara fram á e‑ð}} \dicDirectTranslationCS{požadovat (co), dožadovat se (čeho)};  \dicIdiom{fara}[fram hjá]{ \dicPhraseIS{e‑að fer fram hjá e‑m}} \dicDirectTranslationCS{(co) jde mimo (koho), (co komu) unikne (informace ap.)};  \dicIdiom{fara}[fram úr]{ \dicPhraseIS{fara fram úr (bílnum)}} \dicDirectTranslationCS{předjet (auto)}; { \dicPhraseIS{e‑að fer fram úr (mínum) björtustu vonum}} \dicDirectTranslationCS{(co) předčí (moje) největší očekávání};  \dicIdiom{fara}[frá]{ \dicPhraseIS{fara frá e‑m}} \dicDirectTranslationCS{nechat (koho), opustit (koho)}; { \dicPhraseIS{fara frá}} \dicDirectTranslationCS{pospíšit si, hnout sebou};  \dicIdiom{fara}[fyrir]{ \dicPhraseIS{fara fyrir e‑u}} \dicDirectTranslationCS{vést (co), mít na starost (co)}; { \dicPhraseIS{það fer mikið\,/\addthin lítið fyrir e‑u}} \dicFlx{impers} \dicDirectTranslationCS{(co) zabere hodně\,/\addthin málo místa};  \dicIdiom{fara}[hjá]{ \dicPhraseIS{fara hjá sér}} \dicDirectTranslationCS{přijít do rozpaků}; { \dicPhraseIS{það fer ekki hjá e‑u}} \dicFlx{impers} \dicDirectTranslationCS{(co) je nevyhnutelné};  \dicIdiom{fara}[í]{ \dicPhraseIS{fara í e‑ð}} \dicDirectTranslationCS{obléknout si (co), nandat si (co), nazout\,/\addthin obout si (co)} \dicExampleIS{fara í skóna} \dicExampleCS{nazout si boty}; { \dicPhraseIS{fara í vinnuna}} \dicDirectTranslationCS{jít do práce}; { \dicPhraseIS{e‑að fer í mig}} \dicDirectTranslationCS{(co) mi leze na nervy}; { \dicPhraseIS{e‑að fer (vel) í mig}} \dicDirectTranslationCS{(co) mi chutná (jídlo ap.)};  \dicIdiom{fara}[í kringum]{ \dicPhraseIS{fara í kringum e‑ð}} \dicDirectTranslationCS{obejít (co), obcházet (co)} \dicExampleIS{fara í kringum lögin} \dicExampleCS{obcházet zákon};  \dicIdiom{fara}[í sundur]{ \dicPhraseIS{e‑að fer í sundur}} \dicDirectTranslationCS{(co) se roztrhne\,/\addthin přetrhne};  \dicIdiom{fara}[með]{ \dicPhraseIS{fara vel með e‑ð}} \dicDirectTranslationCS{dobře zacházet\,/\addthin nakládat s~(čím)}; { \dicPhraseIS{fara illa með e‑ð}} \dicDirectTranslationCS{špatně zacházet s~(čím)}; { \dicPhraseIS{fara með e‑ð}} \dicDirectTranslationCS{odříkat (co), (za)recitovat (co)}; { \dicPhraseIS{fara með e‑n (þangað)}} \dicDirectTranslationCS{vzít (koho) (tam)}; { \dicPhraseIS{fara með e‑m (þangað)}} \dicDirectTranslationCS{jít s~(kým) (tam)}; { \dicPhraseIS{fara með heilsuna}} \dicDirectTranslationCS{(z)ničit si zdraví};  \dicIdiom{fara}[ofan af]{ \dicPhraseIS{fara ekki ofan af e‑u}} \dicDirectTranslationCS{nezměnit (co) (názor ap.)};  \dicIdiom{fara}[saman]{ \dicPhraseIS{e‑að fer saman}} \dicLangCat{přen.} \dicDirectTranslationCS{(co) jde ruku v~ruce}; { \dicPhraseIS{e‑að fer (vel) saman}} \dicDirectTranslationCS{(co) jde (dobře) k~sobě};  \dicIdiom{fara}[um]{ \dicPhraseIS{það fer (vel) um e‑n}} \dicFlx{impers} \dicDirectTranslationCS{(kdo) je v~pohodě, (kdo) je na dobrém místě}; { \dicPhraseIS{fara um koll}} \dicDirectTranslationCS{převrátit se}; { \dicPhraseIS{fara um e‑ð}} \dicDirectTranslationCS{navštívit (co)} \dicExampleIS{fara um þjóðgarðinn} \dicExampleCS{navštívit národní park}; { \dicPhraseIS{það fer um e‑n}} \dicFlx{impers} \dicDirectTranslationCS{(kdo) se začíná bát};  \dicIdiom{fara}[utan]{ \dicPhraseIS{fara utan}} \dicDirectTranslationCS{jet do zahraničí\,/\addthin ciziny};  \dicIdiom{fara}[úr]{ \dicPhraseIS{fara úr e‑u}} \dicFlx{dat} \dicDirectTranslationCS{svléknout si (co), sundat si (co) (kalhoty ap.), vyzout si (co) (boty ap.)} \dicExampleIS{fara úr buxunum} \dicExampleCS{sundat si kalhoty}; { \dicPhraseIS{fara úr}} \dicDirectTranslationCS{svléknout se, vyzout se};  \dicIdiom{fara}[út]{ \dicPhraseIS{fara út}} \dicDirectTranslationCS{jet do zahraničí\,/\addthin ciziny}; { \dicPhraseIS{fara út að borða}} \dicDirectTranslationCS{jít se najíst do restaurace}; { \dicPhraseIS{fara út að skemmta sér}} \dicDirectTranslationCS{jít (ven) se bavit (do města ap.)};  \dicIdiom{fara}[út í]{ \dicPhraseIS{fara út í e‑ð}} \dicDirectTranslationCS{zajít do (čeho), zacházet do (čeho) (do detailů ap.)} \dicExampleIS{fara út í smáatriði} \dicExampleCS{zajít do detailů};  \dicIdiom{fara}[út úr]{ \dicPhraseIS{fara út úr e‑u}} \dicDirectTranslationCS{vyjít ven z~(čeho) (se ziskem ap.)};  \dicIdiom{fara}[við]{ \dicPhraseIS{e‑að fer (vel) við e‑ð}} \dicDirectTranslationCS{(co) se hodí k~(čemu)};  \dicIdiom{fara}[yfir]{ \dicPhraseIS{fara yfir e‑ð}} \dicDirectTranslationCS{projít (co), proběhnout si (co)} \dicExampleIS{fara yfir heimaverkefnin} \dicExampleCS{projít si domácí úkol};  \dicIdiom{fara}[yfir um]{ \dicPhraseIS{fara yfir um}} \dicDirectTranslationCS{zbláznit se};  \dicIdiom{farast}{ \dicPhraseIS{farast}} \dicFlx{refl} \dicSynonym{deyja} \dicDirectTranslationCS{zemřít, zahynout}; { \dicPhraseIS{e‑m ferst verkið (vel)}} \dicFlx{refl impers} \dicDirectTranslationCS{(komu) jde práce (dobře)}; { \dicPhraseIS{e‑m ferst}} \dicFlx{refl impers} \dicDirectTranslationCS{(kdo) není ten pravý (na posouzení situace ap.)};  \dicIdiom{farast}[á]{ \dicPhraseIS{farast á mis}} \dicFlx{refl} \dicDirectTranslationCS{minout se};  \dicIdiom{farast}[fyrir]{ \dicPhraseIS{farast fyrir}} \dicFlx{refl} \dicDirectTranslationCS{neuskutečnit se};  \dicIdiom{fara}{ \dicPhraseIS{sem betur fer}} \dicFlx{adv} \dicDirectTranslationCS{naštěstí}
\dicEntry[faraldsfræði] \dicTerm{faralds··fræð|i} \dicIPA{{f}{a}{\textlengthmark}{r}{a}{l}{\textsubring{d}}{s}{f}{r}{a}{i}{ð}{\textsci}} \dicPos{f}[3] \dicFlx{(‑i)}[3] \dicDirectTranslationCS{epidemiologie}
\dicEntry[faraldsfræðingur] \dicTerm{faralds·fræð··ing|ur} \dicIPA{{f}{a}{\textlengthmark}{r}{a}{l}{\textsubring{d}}{s}{f}{r}{a}{i}{ð}{i}{\ng}{\r{g}}{\textscy}{\textsubring{r}}} \dicPos{m}[6] \dicFlx{(‑s, ‑ar)}[8] \dicDirectTranslationCS{epidemiolog, epidemioložka}
\dicEntry[faraldur] \dicTerm{far··ald|ur} \dicIPA{{f}{a}{\textlengthmark}{r}{a}{l}{\textsubring{d}}{\textscy}{\textsubring{r}}} \dicPos{m}[5] \dicFlx{(‑urs, ‑rar)}[3] \textbf{1.} \dicSynonym{farsótt} \dicDirectTranslationCS{epidemie} \dicIndirectTranslationCS{(náhlé rozšíření nemoci)}  \textbf{2.} \dicDirectTranslationCS{epidemie} \dicIndirectTranslationCS{(silná vlna např. krádeží ap.)}
\dicEntry[farand] \dicTerm{far··and-} \dicIPA{{f}{a}{\textlengthmark}{r}{a}{n}{\textsubring{d}}} \dicPos{predp} \dicDirectTranslationCS{putovní, pocestný, pouliční}
\dicEntry[farandsali] \dicTerm{farand··sal|i} \dicIPA{{f}{a}{\textlengthmark}{r}{a}{n}{\textsubring{d}}{s}{a}{l}{\textsci}} \dicPos{m}[1] \dicFlx{(‑a, ‑ar)}[8] \dicDirectTranslationCS{pouliční prodejce\,/\addthin prodejkyně}
\dicEntry[farandsýning] \dicTerm{farand··sýn·ing} \dicIPA{{f}{a}{\textlengthmark}{r}{a}{n}{\textsubring{d}}{s}{i}{n}{i}{\ng}{\r{g}}} \dicPos{f}[4] \dicFlx{(‑ar, ‑ar)}[5] \dicDirectTranslationCS{putovní výstava}
\dicEntry[farandverkamaður] \dicTerm{farand··verka·|maður} \dicIPA{{f}{a}{r}{a}{n}{\textsubring{d}}{v}{\textepsilon}{\textsubring{r}}{\r{g}}{a}{m}{a}{ð}{\textscy}{\textsubring{r}}} \dicPos{m}[13] \dicFlx{(‑manns, ‑menn)}[2] \dicDirectTranslationCS{sezónní pracovník\,/\addthin pracovnice}
\dicEntry[farangur] \dicTerm{far··ang|ur} \dicIPA{{f}{a}{\textlengthmark}{r}{au}{\ng}{\r{g}}{\textscy}{\textsubring{r}}} \dicPos{m}[5] \dicFlx{(‑urs)}[2] \dicSynonym{dót} \dicDirectTranslationCS{zavazadlo}
\dicEntry[farangursgeymsla] \dicTerm{farangurs··geymsl|a} \dicIPA{{f}{a}{\textlengthmark}{r}{au}{\ng}{\r{g}}{\textscy}{\textsubring{r}}{s}{\r{\textObardotlessj}}{ei}{m}{s}{\textsubring{d}}{l}{a}} \dicPos{f}[1] \dicFlx{(‑u, ‑ur)}[13] \textbf{1.} \dicDirectTranslationCS{kufr, zavazadlový prostor (v~autě ap.)}  \textbf{2.} \dicDirectTranslationCS{úschovna zavazadel}
\dicEntry[farangursgrind] \dicTerm{farangurs··grind} \dicIPA{{f}{a}{\textlengthmark}{r}{au}{\ng}{\r{g}}{\textscy}{\textsubring{r}}{s}{\r{g}}{r}{\textsci}{n}{\textsubring{d}}} \dicPos{f}[8] \dicFlx{(‑ar, ‑ur)}[1] \dicDirectTranslationCS{střešní nosič, zahrádka (na střeše auta ap.)}
\dicEntry[farangursmiði] \dicTerm{farangurs··mið|i} \dicIPA{{f}{a}{\textlengthmark}{r}{au}{\ng}{\r{g}}{\textscy}{\textsubring{r}}{s}{m}{\textsci}{ð}{\textsci}} \dicPos{m}[1] \dicFlx{(‑a, ‑ar)}[1] \dicDirectTranslationCS{jízdenka pro zavazadlo}
\dicEntry[farangursrými] \dicTerm{farangurs··rými} \dicIPA{{f}{a}{\textlengthmark}{r}{au}{\ng}{\r{g}}{\textscy}{\textsubring{r}}{s}{r}{i}{m}{\textsci}} \dicPos{n}[2] \dicFlx{(‑s, ‑)}[14] \dicDirectTranslationCS{zavazadlový prostor}
\dicEntry[farar] \dicTerm{farar} \dicIPA{{f}{a}{\textlengthmark}{r}{a}{\textsubring{r}}} \dicPos{f} \dicFlx{sg gen} \dicLink{för\smash{\textsuperscript{1}}}
\dicEntry[fararbroddur] \dicTerm{farar··brodd|ur} \dicIPA{{f}{a}{\textlengthmark}{r}{a}{r}{\textsubring{b}}{r}{\textopeno}{\textsubring{d}}{\textscy}{\textsubring{r}}} \dicPos{m}[6] \dicFlx{(‑s)}[7] \dicPhraseIS{vera í fararbroddi} \dicLangCat{přen.} \dicDirectTranslationCS{být ve vedení, být na špici}
\dicEntry[fararstjóri] \dicTerm{farar··stjór|i} \dicIPA{{f}{a}{\textlengthmark}{r}{a}{\textsubring{r}}{s}{\textsubring{d}}{j}{ou}{r}{\textsci}} \dicPos{m}[1] \dicFlx{(‑a, ‑ar)}[1] \dicDirectTranslationCS{(turistický) průvodce, (turistická) průvodkyně}
\dicEntry[farartæki] \dicTerm{farar··tæki} \dicIPA{{f}{a}{\textlengthmark}{r}{a}{\textsubring{r}}{t\smash{\textsuperscript{h}}}{a}{i}{\r{\textObardotlessj}}{\textsci}} \dicPos{n}[2] \dicFlx{(‑s, ‑)}[16] \dicDirectTranslationCS{vůz, vozidlo, dopravní prostředek}
\dicEntry[farði] \dicTerm{farð|i} \dicIPA{{f}{a}{r}{ð}{\textsci}} \dicPos{m}[1] \dicFlx{(‑a)}[3] \dicDirectTranslationCS{líčidlo}
\dicEntry[farfugl] \dicTerm{far··fugl} \dicIPA{{f}{a}{\textsubring{r}}{f}{\textscy}{\r{g}}{\textsubring{l}}} \dicPos{m}[4] \dicFlx{(‑s, ‑ar)}[3] \dicDirectTranslationCS{stěhovavý\,/\addthin tažný pták} \dicAntonym{staðfugl}
\dicEntry[farfuglaheimili] \dicTerm{far·fugla··heimili} \dicIPA{{f}{a}{\textsubring{r}}{f}{\textscy}{\r{g}}{l}{a}{h}{ei}{m}{\textsci}{l}{\textsci}} \dicPos{n}[2] \dicFlx{(‑s, ‑)}[14] \dicDirectTranslationCS{hostel, ubytovna}
\dicEntry[farg] \dicTerm{farg} \dicIPA{{f}{a}{r}{\r{g}}} \dicPos{n}[2] \dicFlx{(‑s, förg)}[8] \dicSynonym{þungi} \dicDirectTranslationCS{břemeno, zátěž} \dicExampleIS{Það liggur þungt farg á mér.} \dicExampleCS{Leží na mně těžké břemeno.}
\dicEntry[farga] \dicTerm{farg|a} \dicIPA{{f}{a}{r}{\r{g}}{a}} \dicPos{v}[1] \dicFlx{(‑aði)}[50] \dicFlx{dat} \textbf{1.} \dicSynonym*{láta e‑ð af hendi} \dicDirectTranslationCS{zbavit se, zlikvidovat}  \textbf{2.} \dicSynonym{drepa} \dicDirectTranslationCS{zabít, porazit} \dicExampleIS{farga sauðfénu} \dicExampleCS{porazit ovce};  \dicPhraseIS{farga sér} \dicSynonym*{fyrirfara sér} \dicDirectTranslationCS{zabít se, spáchat sebevraždu}
\dicEntry[fargjald] \dicTerm{far··|gjald} \dicIPA{{f}{a}{r}{\r{\textObardotlessj}}{a}{l}{\textsubring{d}}} \dicPos{n}[2] \dicFlx{(‑gjalds, ‑gjöld)}[8] \dicDirectTranslationCS{jízdné}
\dicEntry[farið] \dicTerm{farið} \dicIPA{{f}{a}{\textlengthmark}{r}{\textsci}{\texttheta}} \dicPos{v} \dicFlx{supin} \dicLink{fara}
\dicEntry[farinn] \dicTerm{farinn} \dicIPA{{f}{a}{\textlengthmark}{r}{\textsci}{\textsubring{n}}} \dicPos{adj}[6]\dicFlx{}[-4] \dicPhraseIS{vera farinn að heilsu} \dicDirectTranslationCS{přijít o~zdraví};  \dicPhraseIS{vera farinn að kröftum} \dicDirectTranslationCS{ztrácet síly}
\dicEntry[farir] \dicTerm{farir} \dicIPA{{f}{a}{\textlengthmark}{r}{\textsci}{\textsubring{r}}} \dicPos{f} \dicFlx{pl nom} \dicLink{för\smash{\textsuperscript{1}}}
\dicEntry[farlama] \dicTerm{far··lama} \dicIPA{{f}{a}{r}{l}{a}{m}{a}} \dicPos{adj}[13] \dicFlx{indecl}[1] \dicSynonym{örvasa} \dicDirectTranslationCS{obtížně se pohybující, (jsoucí) se sníženou pohyblivostí};  \dicPhraseIS{farlama gamalmenni} \dicDirectTranslationCS{senior se sníženou pohyblivostí}
\dicEntry[farmaður] \dicTerm{far··|maður} \dicIPA{{f}{a}{r}{m}{a}{ð}{\textscy}{\textsubring{r}}} \dicPos{m}[13] \dicFlx{(‑manns, ‑menn)}[2] \dicDirectTranslationCS{námořník, námořnice}
\dicEntry[farmflytjandi] \dicTerm{farm··flytj·|andi} \dicIPA{{f}{a}{r}{m}{f}{l}{\textsci}{\textsubring{d}}{j}{a}{n}{\textsubring{d}}{\textsci}} \dicPos{m}[2] \dicFlx{(‑anda, ‑endur)}[1] \dicDirectTranslationCS{náklad\-ní dopravce}
\dicEntry[farmgjald] \dicTerm{farm··|gjald} \dicIPA{{f}{a}{r}{m}{\r{\textObardotlessj}}{a}{l}{\textsubring{d}}} \dicPos{n}[2] \dicFlx{(‑gjalds, ‑gjöld)}[8] \dicDirectTranslationCS{dopravné}
\dicEntry[farmiði] \dicTerm{far··mið|i} \dicIPA{{f}{a}{r}{m}{\textsci}{ð}{\textsci}} \dicPos{m}[1] \dicFlx{(‑a, ‑ar)}[1] \dicSynonym{farseðill} \dicDirectTranslationCS{jízdenka, lístek}
\dicEntry[farmrúm] \dicTerm{farm··rúm} \dicIPA{{f}{a}{r}{m}{r}{u}{\textsubring{m}}} \dicPos{n}[2] \dicFlx{(‑s)}[2] \dicFieldCat{nám.} \dicDirectTranslationCS{nákladní prostor}
\dicEntry[farmskip] \dicTerm{farm··skip} \dicIPA{{f}{a}{r}{m}{s}{\r{\textObardotlessj}}{\textsci}{\textsubring{b}}} \dicPos{n}[2] \dicFlx{(‑s, ‑)}[5] \dicDirectTranslationCS{nákladní loď}
\dicEntry[farmskírteini] \dicTerm{farm··skír·teini} \dicIPA{{f}{a}{r}{m}{s}{\r{\textObardotlessj}}{i}{\textsubring{r}}{t\smash{\textsuperscript{h}}}{ei}{n}{\textsci}} \dicPos{n}[2] \dicFlx{(‑s, ‑)}[14] \dicDirectTranslationCS{nákladní list, konosament}
\dicEntry[farmskrá] \dicTerm{farm··skrá} \dicIPA{{f}{a}{r}{m}{s}{\r{g}}{r}{au}} \dicPos{f}[4] \dicFlx{(‑r\,/\addthin ‑ar, ‑r)}[21] \dicDirectTranslationCS{konosament, nákladní list}
\dicEntry[farmtrygging] \dicTerm{farm··trygg·ing} \dicIPA{{f}{a}{r}{m}{t\smash{\textsuperscript{h}}}{r}{\textsci}{\r{\textObardotlessj}}{i}{\ng}{\r{g}}} \dicPos{f}[4] \dicFlx{(‑ar, ‑ar)}[5] \dicFieldCat{ekon.} \dicDirectTranslationCS{pojištění nákladu}
\dicEntry[farmur] \dicTerm{farm|ur} \dicIPA{{f}{a}{r}{m}{\textscy}{\textsubring{r}}} \dicPos{m}[6] \dicFlx{(‑s, ‑ar)}[5] \dicDirectTranslationCS{náklad, zásilka}
\dicEntry[farnaður] \dicTerm{far··nað|ur} \dicIPA{{f}{a}{r}{\textsubring{d}}{n}{a}{ð}{\textscy}{\textsubring{r}}} \dicPos{m}[10] \dicFlx{(‑ar)}[9] \dicSynonym{gengi\smash{\textsuperscript{1}}} \dicDirectTranslationCS{úspěch, zdar}
\dicEntry[farnast] \dicTerm{farn|ast} \dicIPA{{f}{a}{r}{\textsubring{d}}{n}{a}{s}{\textsubring{d}}} \dicPos{v}[1] \dicFlx{(‑aðist)}[102] \dicFlx{refl impers} \dicPhraseIS{e‑m farnast vel\,/\addthin illa} \dicDirectTranslationCS{(komu) se daří\,/\addthin nedaří, (kdo) prospívá\,/\addthin neprospívá}
\dicEntry[farrými] \dicTerm{far··rými} \dicIPA{{f}{a}{r}{\textlengthmark}{i}{m}{\textsci}} \dicPos{n}[2] \dicFlx{(‑s, ‑)}[14] \dicDirectTranslationCS{třída (ve vlaku ap.)} \dicExampleIS{fyrsta farrými} \dicExampleCS{první třída}
\dicEntry[fars] \dicTerm{fars} \dicIPA{{f}{a}{\textsubring{r}}{s}} \dicPos{n}[2] \dicFlx{(‑)}[24] \dicDirectTranslationCS{sekané maso}
\dicEntry[farseðill] \dicTerm{far··seð|ill} \dicIPA{{f}{a}{\textsubring{r}}{s}{\textepsilon}{ð}{\textsci}{\textsubring{d}}{\textsubring{l}}} \dicPos{m}[6] \dicFlx{(‑ils, ‑lar)}[35] \dicDirectTranslationCS{jízdenka, lístek}
\dicEntry[farsi] \dicTerm{fars|i} \dicIPA{{f}{a}{\textsubring{r}}{s}{\textsci}} \dicPos{m}[1] \dicFlx{(‑a, ‑ar)}[8] \textbf{1.} \dicFieldCat{lit.} \dicDirectTranslationCS{fraška} \dicIndirectTranslationCS{(veselohra s~obhroublou komikou)}  \textbf{2.} \dicDirectTranslationCS{fraška} \dicIndirectTranslationCS{(nesmyslné jednání\,/\addthin dění)}
\dicEntry[farsímanúmer] \dicTerm{far·síma··númer} \dicIPA{{f}{a}{\textsubring{r}}{s}{i}{m}{a}{n}{u}{m}{\textepsilon}{\textsubring{r}}} \dicPos{n}[2] \dicFlx{(‑s, ‑)}[5] \dicDirectTranslationCS{mobilní číslo, číslo mobilního telefonu}
\dicEntry[farsími] \dicTerm{far··sím|i} \dicIPA{{f}{a}{\textsubring{r}}{s}{i}{m}{\textsci}} \dicPos{m}[1] \dicFlx{(‑a, ‑ar)}[1] \dicDirectTranslationCS{mobilní telefon, mobil}
\dicEntry[farskrá] \dicTerm{far··skrá} \dicIPA{{f}{a}{\textsubring{r}}{s}{\r{g}}{r}{au}} \dicPos{f}[4] \dicFlx{(‑r\,/\addthin ‑ar, ‑r)}[21] \dicDirectTranslationCS{seznam pasažérů\,/\addthin cestujících}
\dicEntry[farsótt] \dicTerm{far··sótt} \dicIPA{{f}{a}{\textsubring{r}}{s}{ou}{h}{\textsubring{d}}} \dicPos{f}[7] \dicFlx{(‑ar, ‑ir)}[1] \dicSynonym{faraldur} \dicDirectTranslationCS{epidemie}
\dicEntry[farsæld] \dicTerm{far··sæld} \dicIPA{{f}{a}{\textsubring{r}}{s}{a}{i}{l}{\textsubring{d}}} \dicPos{f}[7] \dicFlx{(‑ar)}[3] \dicDirectTranslationCS{prosperita, úspěch, zdar};  \dicPhraseIS{óska e‑m farsældar} \dicDirectTranslationCS{popřát (komu) úspěch}
\dicEntry[farsæll] \dicTerm{far··sæll} \dicIPA{{f}{a}{\textsubring{r}}{s}{a}{i}{\textsubring{d}}{\textsubring{l}}} \dicPos{adj}[8]\dicFlx{}[-1] \dicDirectTranslationCS{úspěšný, zdárný};  \dicPhraseIS{farsælt hjó\-na\-band} \dicDirectTranslationCS{šťastné manželství}
\dicEntry[fartölva] \dicTerm{far··tölv|a} \dicIPA{{f}{a}{\textsubring{r}}{t\smash{\textsuperscript{h}}}{\oe}{l}{v}{a}} \dicPos{f}[1] \dicFlx{(‑u, ‑ur)}[7] \dicFieldCat{poč.} \dicDirectTranslationCS{notebook, laptop}
\dicEntry[farvegur] \dicTerm{far··veg|ur} \dicIPA{{f}{a}{r}{v}{\textepsilon}{\textbabygamma}{\textscy}{\textsubring{r}}} \dicPos{m}[10] \dicFlx{(‑s\,/\addthin ‑ar, ‑ir)}[29] \dicDirectTranslationCS{řečiště, koryto řeky};  \dicPhraseIS{e‑að er í (réttum) farvegi} \dicDirectTranslationCS{(co) je na (správné) cestě}
\dicEntry[farþegaskip] \dicTerm{far·þega··skip} \dicIPA{{f}{a}{\textsubring{r}}{\texttheta}{\textepsilon}{\textbabygamma}{a}{s}{\r{\textObardotlessj}}{\textsci}{\textsubring{b}}} \dicPos{n}[2] \dicFlx{(‑s, ‑)}[5] \dicFieldCat{nám.} \dicDirectTranslationCS{trajekt, přívoz}
\dicEntry[farþegaþota] \textls[15]{\dicTerm{far·þega··þot|a} \dicIPA{{f}{a}{\textsubring{r}}{\texttheta}{\textepsilon}{\textbabygamma}{a}{\texttheta}{\textopeno}{\textsubring{d}}{a}}\addthinM\dicPos{f}[1]\addthinM\dicFlx{(‑u, ‑ur)}[19] \dicDirectTranslationCS{dopravní le\-toun\,/\addthin letadlo}}
\dicEntry[farþegi] \dicTerm{far··þeg|i} \dicsymFrequent\  \dicIPA{{f}{a}{\textsubring{r}}{\texttheta}{ei}{\textsci}} \dicPos{m}[1] \dicFlx{(‑a, ‑ar)}[1] \dicDirectTranslationCS{cestující, pasažér(ka)} \dicExampleIS{Skipstjóri hefur neitað farþegum um far vegna slæms veðurs.} \dicExampleCS{Kapitán odepřel cestujícím plavbu kvůli špatnému počasí.}
\dicEntry[fas] \dicTerm{fas} \dicIPA{{f}{a}{\textlengthmark}{s}} \dicPos{n}[2] \dicFlx{(‑s)}[2] \dicSynonym{framkoma} \dicDirectTranslationCS{vystupování, chování} \dicExampleIS{vera vingjarnlegur í fasi} \dicExampleCS{mít přátelské vystupování}
\dicEntry[fasani] \dicTerm{fasan|i} \dicIPA{{f}{a}{\textlengthmark}{s}{a}{n}{\textsci}} \dicPos{m}[1] \dicFlx{(‑a, ‑ar)}[11] \dicFieldCat{zool.} \dicDirectTranslationCS{bažant, bažant obecný} \textit{(l.~{\textLA{Phasianus colchicus}})}  \dicsymPhoto\ 
\dicFigure{113961.jpg}{Fasani}{Fasani - Šinko Jan, Biolib, Copyright/PD}
\dicEntry[fasismi] \dicTerm{fas··ism|i} \dicIPA{{f}{a}{\textlengthmark}{s}{\textsci}{s}{m}{\textsci}} \dicPos{m}[1] \dicFlx{(‑a)}[3] \dicDirectTranslationCS{fašismus}
\dicEntry[fasisti] \dicTerm{fas··ist|i} \dicIPA{{f}{a}{\textlengthmark}{s}{\textsci}{s}{\textsubring{d}}{\textsci}} \dicPos{m}[1] \dicFlx{(‑a, ‑ar)}[1] \dicDirectTranslationCS{fašista, fašistka}
\dicEntry[fasískur] \dicTerm{fasískur} \dicIPA{{f}{a}{\textlengthmark}{s}{i}{s}{\r{g}}{\textscy}{\textsubring{r}}} \dicPos{adj}[1]\dicFlx{}[-1] \dicDirectTranslationCS{fašistický}
\dicEntry[fast] \dicTerm{fast} \dicsymFrequent\  \dicIPA{{f}{a}{s}{\textsubring{d}}} \dicPos{adv} \dicFlx{(comp ‑ar, sup ‑ast)} \dicSynonym{sterklega} \dicDirectTranslationCS{pevně, napevno} \dicExampleIS{haltu þér fast!} \dicExampleCS{drž se pevně!};  \dicPhraseIS{fyrir fullt og fast} \dicFlx{adv} \dicDirectTranslationCS{definitivně, napořád};  \dicPhraseIS{halda fast við rétt sinn} \dicDirectTranslationCS{stát si pevně za svým právem};  \dicPhraseIS{sofa fast} \dicLangCat{přen.} \dicDirectTranslationCS{spát jako zabitý, spát jako špalek};  \dicPhraseIS{standa fast á sínu} \dicLangCat{přen.} \dicDirectTranslationCS{stát si za svým};  \dicPhraseIS{sækja þetta fast} \dicDirectTranslationCS{vydobýt si to, vybojovat si to}
\dicEntry[fasta] \dicTerm{fasta\smash{\textsuperscript{1}}} \dicIPA{{f}{a}{s}{\textsubring{d}}{a}} \dicPos{f}[1] \dicFlx{(föstu, föstur)}[20] \dicDirectTranslationCS{půst, abstinence};  \dicPhraseIS{halda föstu} \dicFieldCat{náb.} \dicDirectTranslationCS{postit se}
\dicEntry[fasta] \dicTerm{fast|a\smash{\textsuperscript{2}}} \dicIPA{{f}{a}{s}{\textsubring{d}}{a}} \dicPos{v}[1] \dicFlx{(‑aði)}[50] \dicFieldCat{náb.} \dicSynonym*{halda föstu} \dicDirectTranslationCS{postit se}
\dicEntry[fastafjármunir] \dicTerm{fasta··fjár·munir} \dicIPA{{f}{a}{s}{\textsubring{d}}{a}{f}{j}{au}{r}{m}{\textscy}{n}{\textsci}{\textsubring{r}}} \dicPos{m}[10] \dicFlx{pl}[1] \dicFieldCat{ekon.} \dicDirectTranslationCS{fixní aktiva\,/\addthin majetek}
\dicEntry[fastagestur] \dicTerm{fasta··gest|ur} \dicIPA{{f}{a}{s}{\textsubring{d}}{a}{\r{\textObardotlessj}}{\textepsilon}{s}{\textsubring{d}}{\textscy}{\textsubring{r}}} \dicPos{m}[9] \dicFlx{(‑s, ‑ir)}[4] \dicDirectTranslationCS{stálý host\,/\addthin zákazník\,/\addthin návštěvník, stálá zákaznice\,/\addthin návštěvnice}
\begin{xtolerant}{}{1pt}
\dicEntry[fastakostnaður] \dicTerm{fasta··kost·nað|ur}\addthinS\dicIPA{{f}{a}{s}{\textsubring{d}}{a}{k\smash{\textsuperscript{h}}}{\textopeno}{s}{\textsubring{d}}{n}{a}{ð}{\textscy}{\textsubring{r}}}\addthinS\dicPos{m}[10]\addthinS\dicFlx{(‑ar)}[9]\addthinS\dicFieldCat{ekon.}\addthinS\dicDirectTranslationCS{fixní\hspace{-0.10667em}\,/\addthin režijní náklady}
\end{xtolerant}
\dicEntry[fastakúnni] \dicTerm{fasta··kúnn|i} \dicIPA{{f}{a}{s}{\textsubring{d}}{a}{k\smash{\textsuperscript{h}}}{u}{n}{\textlengthmark}{\textsci}} \dicPos{m}[1] \dicFlx{(‑a, ‑ar)}[1] \dicDirectTranslationCS{stálý host\,/\addthin zákazník, stálá zákaznice}
\dicEntry[fastaland] \dicTerm{fasta··|land} \dicIPA{{f}{a}{s}{\textsubring{d}}{a}{l}{a}{n}{\textsubring{d}}} \dicPos{n}[2] \dicFlx{(‑lands, ‑lönd)}[8] \dicSynonym{meginland} \dicDirectTranslationCS{pevnina}
\dicEntry[fasteign] \dicTerm{fast··eign} \dicIPA{{f}{a}{s}{\textsubring{d}}{ei}{\r{g}}{\textsubring{n}}} \dicPos{f}[7] \dicFlx{(‑ar, ‑ir)}[1] \dicDirectTranslationCS{nemovitost, nemovitý majetek}
\dicEntry[fasteignasala] \dicTerm{fast·eigna··|sala} \dicIPA{{f}{a}{s}{\textsubring{d}}{ei}{\r{g}}{n}{a}{s}{a}{l}{a}} \dicPos{f}[1] \dicFlx{(‑sölu, ‑sölur)}[20] \dicDirectTranslationCS{realitní agentura\,/\addthin kancelář}
\dicEntry[fasteignasali] \dicTerm{fast·eigna··sal|i} \dicIPA{{f}{a}{s}{\textsubring{d}}{ei}{\r{g}}{n}{a}{s}{a}{l}{\textsci}} \dicPos{m}[1] \dicFlx{(‑a, ‑ar)}[8] \dicDirectTranslationCS{realitní agent(ka)\,/\addthin makléř(ka), obchodník\,/\addthin obchodnice s~nemovitostmi}
\dicEntry[fasteignaskattur] \dicTerm{fast·eigna··skatt|ur} \dicIPA{{f}{a}{s}{\textsubring{d}}{ei}{\r{g}}{n}{a}{s}{\r{g}}{a}{h}{\textsubring{d}}{\textscy}{\textsubring{r}}} \dicPos{m}[6] \dicFlx{(‑s, ‑ar)}[5] \dicFieldCat{ekon.} \dicDirectTranslationCS{daň z~nemovitosti}
\dicEntry[fasteignaveð] \dicTerm{fast·eigna··veð} \dicIPA{{f}{a}{s}{\textsubring{d}}{ei}{\r{g}}{n}{a}{v}{\textepsilon}{\texttheta}} \dicPos{n}[2] \dicFlx{(‑s, ‑)}[5] \dicFieldCat{ekon.} \dicDirectTranslationCS{zástava nemovitosti}
\dicEntry[fastheldinn] \dicTerm{fast··heldinn} \dicIPA{{f}{a}{s}{t\smash{\textsuperscript{h}}}{\textepsilon}{l}{\textsubring{d}}{\textsci}{\textsubring{n}}} \dicPos{adj}[6]\dicFlx{}[-2] \textbf{1.} \dicSynonym{íhaldssamur} \dicDirectTranslationCS{konzervativní}  \textbf{2.} \dicSynonym{nískur} \dicDirectTranslationCS{lakotný, skoupý}
\dicEntry[fastheldni] \dicTerm{fast··heldn|i} \dicIPA{{f}{a}{s}{t\smash{\textsuperscript{h}}}{\textepsilon}{l}{\textsubring{d}}{n}{\textsci}} \dicPos{f}[3] \dicFlx{(‑i)}[3] \textbf{1.} \dicSynonym*{vanafesta} \dicDirectTranslationCS{lpění, ulpívání, konzervatismus}  \textbf{2.} \dicSynonym{níska} \dicDirectTranslationCS{lakota, skoupost}
\dicEntry[fasti] \dicTerm{fast|i} \dicIPA{{f}{a}{s}{\textsubring{d}}{\textsci}} \dicPos{m}[1] \dicFlx{(‑a, ‑ar)}[8] \dicFieldCat{mat.} \dicDirectTranslationCS{konstanta}
\dicEntry[fastlega] \dicTerm{fast··lega} \dicIPA{{f}{a}{s}{\textsubring{d}}{l}{\textepsilon}{\textbabygamma}{a}} \dicPos{adv} \dicDirectTranslationCS{pevně, jistě}
\dicEntry[fastmótaður] \dicTerm{fast··mót·|aður} \dicIPA{{f}{a}{s}{\textsubring{d}}{m}{ou}{\textsubring{d}}{a}{ð}{\textscy}{\textsubring{r}}} \dicPos{adj}[3] \dicFlx{(f ‑uð)}[3] \dicDirectTranslationCS{pevný, pev\-ně stanovený, fixní}
\dicEntry[fastráða] \dicTerm{fast··|ráða} \dicIPA{{f}{a}{s}{\textsubring{d}}{r}{au}{ð}{a}} \dicPos{v}[6] \dicFlx{(‑ræð, ‑réð, ‑réðum, ‑réði, ‑ráðið)}[88] \dicFlx{acc} \textbf{1.} \dicDirectTranslationCS{(rychle) rozhodnout}  \textbf{2.} \dicDirectTranslationCS{zaměstnat na trvalý pracovní poměr}
\dicEntry[fastráðið] \dicTerm{fast··ráðið} \dicIPA{{f}{a}{s}{\textsubring{d}}{r}{au}{ð}{\textsci}{\texttheta}} \dicPos{v} \dicFlx{supin} \dicLink{fastráða}
\dicEntry[fastréð] \dicTerm{fast··réð} \dicIPA{{f}{a}{s}{\textsubring{d}}{r}{j}{\textepsilon}{\texttheta}} \dicPos{v} \dicFlx{ind pf sg 1 pers} \dicLink{fastráða}
\dicEntry[fastréði] \dicTerm{fast··réði} \dicIPA{{f}{a}{s}{\textsubring{d}}{r}{j}{\textepsilon}{ð}{\textsci}} \dicPos{v} \dicFlx{con pf sg 1 pers} \dicLink{fastráða}
\dicEntry[fastréðum] \dicTerm{fast··réðum} \dicIPA{{f}{a}{s}{\textsubring{d}}{r}{j}{\textepsilon}{ð}{\textscy}{\textsubring{m}}} \dicPos{v} \dicFlx{ind pf pl 1 pers} \dicLink{fastráða}
\dicEntry[fastræð] \dicTerm{fast··ræð} \dicIPA{{f}{a}{s}{\textsubring{d}}{r}{a}{i}{\texttheta}} \dicPos{v} \dicFlx{ind praes sg 1 pers} \dicLink{fastráða}
\dicEntry[fastur] \dicTerm{fastur} \dicsymFrequent\  \dicIPA{{f}{a}{s}{\textsubring{d}}{\textscy}{\textsubring{r}}} \dicPos{adj}[1] \dicFlx{(f föst)}[11] \textbf{1.} \dicSynonym*{sem hreyfist ekki} \dicDirectTranslationCS{pevný, upevněný, uchycený}  \textbf{2.} \dicSynonym{reglulegur} \dicDirectTranslationCS{stálý, trvalý}  \textbf{3.} \dicDirectTranslationCS{zaklesnutý, uvízlý, zaklíněný};  \dicIdiom{fastur}{ \dicPhraseIS{fast efni}} \dicFieldCat{fyz.} \dicDirectTranslationCS{pevná látka};  \dicPhraseIS{taka e‑n fastan} \dicDirectTranslationCS{zatknout (koho), zadržet (koho)};  \dicPhraseIS{taka e‑ð föstum tökum} \dicLangCat{přen.} \dicDirectTranslationCS{vzít (co) pevně do rukou};  \dicPhraseIS{vera fastur á fé} \dicDirectTranslationCS{být skoupý na peníze};  \dicPhraseIS{vera fastur í sessi} \dicLangCat{přen.} \dicDirectTranslationCS{mít jisté místo (v~bance ap.)}
\dicEntry[fat] \dicTerm{fat} \dicsymFrequent\  \dicIPA{{f}{a}{\textlengthmark}{\textsubring{d}}} \dicPos{n}[2] \dicFlx{(‑s, föt)}[8] \textbf{1.} \dicSynonym{ílát} \dicDirectTranslationCS{tác, podnos, mísa} \dicExampleIS{fat undir fisk} \dicExampleCS{mísa na rybu}  \textbf{2.} \dicSynonym{flík} \dicDirectTranslationCS{oděv, kus oblečení}
\dicEntry[fata] \dicTerm{fata\smash{\textsuperscript{1}}} \dicIPA{{f}{a}{\textlengthmark}{\textsubring{d}}{a}} \dicPos{f}[1] \dicFlx{(fötu, fötur)}[20] \dicDirectTranslationCS{kbelík, kyblík}
\dicEntry[fata] \dicTerm{fat|a\smash{\textsuperscript{2}}} \dicIPA{{f}{a}{\textlengthmark}{\textsubring{d}}{a}} \dicPos{v}[1] \dicFlx{(‑aði)}[43] \dicPhraseIS{fata sig upp} \dicDirectTranslationCS{nakoupit oblečení, obnovit šatník}
\dicEntry[fatabursti] \dicTerm{fata··burst|i} \dicIPA{{f}{a}{\textlengthmark}{\textsubring{d}}{a}{\textsubring{b}}{\textscy}{\textsubring{r}}{s}{\textsubring{d}}{\textsci}} \dicPos{m}[1] \dicFlx{(‑a, ‑ar)}[1] \dicDirectTranslationCS{kartáč na šaty}
\dicEntry[fatageymsla] \dicTerm{fata··geymsl|a} \dicIPA{{f}{a}{\textlengthmark}{\textsubring{d}}{a}{\r{\textObardotlessj}}{ei}{m}{s}{\textsubring{d}}{l}{a}} \dicPos{f}[1] \dicFlx{(‑u, ‑ur)}[13] \dicDirectTranslationCS{šatna, garderoba}
\dicEntry[fatahengi] \dicTerm{fata··hengi} \dicIPA{{f}{a}{\textlengthmark}{\textsubring{d}}{a}{h}{ei}{\textltailn}{\r{\textObardotlessj}}{\textsci}} \dicPos{n}[2] \dicFlx{(‑s, ‑)}[16] \textbf{1.} \dicDirectTranslationCS{věšák (na šaty)}  \textbf{2.} \dicDirectTranslationCS{šatna, šatník, garderoba}
\dicEntry[fatahreinsun] \dicTerm{fata··hreins|un} \dicIPA{{f}{a}{\textlengthmark}{\textsubring{d}}{a}{\textsubring{r}}{ei}{n}{s}{\textscy}{\textsubring{n}}} \dicPos{f}[7] \dicFlx{(‑unar, ‑anir)}[8] \dicDirectTranslationCS{čistírna (oděvů)}
\dicEntry[fatahönnuður] \dicTerm{fata··hönn·uð|ur} \dicIPA{{f}{a}{\textlengthmark}{\textsubring{d}}{a}{h}{\oe}{n}{\textscy}{ð}{\textscy}{\textsubring{r}}} \dicPos{m}[10] \dicFlx{(‑ar, ‑ir)}[4] \dicDirectTranslationCS{módní návrhář(ka)} \dicExampleIS{heimsfrægur fatahönnuður} \dicExampleCS{světoznámý módní návrhář}
\dicEntry[fatahönnun] \dicTerm{fata··hönn|un} \dicIPA{{f}{a}{\textlengthmark}{\textsubring{d}}{a}{h}{\oe}{n}{\textscy}{\textsubring{n}}} \dicPos{f}[7] \dicFlx{(‑unar)}[12] \dicDirectTranslationCS{módní návrhářství}
\dicEntry[fataskápur] \dicTerm{fata··skáp|ur} \dicIPA{{f}{a}{\textlengthmark}{\textsubring{d}}{a}{s}{\r{g}}{au}{\textsubring{b}}{\textscy}{\textsubring{r}}} \dicPos{m}[6] \dicFlx{(‑s, ‑ar)}[15] \dicDirectTranslationCS{šatník, skříň na šaty}
\dicEntry[fataverslun] \dicTerm{fata··versl|un} \dicIPA{{f}{a}{\textlengthmark}{\textsubring{d}}{a}{v}{\textepsilon}{\textsubring{r}}{s}{\textsubring{d}}{l}{\textscy}{\textsubring{n}}} \dicPos{f}[7] \dicFlx{(‑unar, ‑anir)}[8] \dicDirectTranslationCS{prodejna\,/\addthin obchod s~oděvy}
\dicEntry[fatlaður] \dicTerm{fatl··aður} \dicIPA{{f}{a}{h}{\textsubring{d}}{l}{a}{ð}{\textscy}{\textsubring{r}}} \dicPos{adj}[3] \dicFlx{(f fötluð)}[1] \dicDirectTranslationCS{(tělesně) postižený, handicapovaný} \dicExampleIS{Hún er fötluð eftir slysið.} \dicExampleCS{Je postižená po nehodě.}
\dicEntry[fatlast] \dicTerm{fatl|ast} \dicIPA{{f}{a}{h}{\textsubring{d}}{l}{a}{s}{\textsubring{d}}} \dicPos{v}[1] \dicFlx{(‑aðist)}[101] \dicFlx{refl} \dicDirectTranslationCS{být postiženým} \dicExampleIS{fatlast í bílslysi} \dicExampleCS{stát se postiženým při automobilové nehodě}
\dicEntry[fatli] \dicTerm{fatl|i} \dicIPA{{f}{a}{h}{\textsubring{d}}{l}{\textsci}} \dicPos{m}[1] \dicFlx{(‑a, ‑ar)}[8] \dicSynonym{fetill} \dicDirectTranslationCS{zdravotnický\,/\addthin trojcípý šátek};  \dicPhraseIS{bera hönd í fatla} \dicDirectTranslationCS{mít ruku zavázanou v~šátku (po zlomenině ap.)}
\dicEntry[fatnaður] \dicTerm{fat··nað|ur} \dicIPA{{f}{a}{h}{\textsubring{d}}{n}{a}{ð}{\textscy}{\textsubring{r}}} \dicPos{m}[10] \dicFlx{(‑ar)}[9] \dicDirectTranslationCS{oblečení, oděv, šat}
\dicEntry[fatta] \dicTerm{fatt|a} \dicsymFrequent\  \dicIPA{{f}{a}{h}{\textsubring{d}}{a}} \dicPos{v}[1] \dicFlx{(‑aði)}[28] \dicFlx{acc} \dicLangCat{hovor.} \dicDirectTranslationCS{(po)rozumět, pochopit, chápat} \dicExampleIS{Fattarðu þetta ekki?} \dicExampleCS{Ty to nechápeš?}
\dicEntry[fauk] \dicTerm{fauk} \dicIPA{{f}{\oe i}{\textlengthmark}{\r{g}}} \dicPos{v} \dicFlx{ind pf sg 1 pers} \dicLink{fjúka}
\dicEntry[fauti] \dicTerm{faut|i} \dicIPA{{f}{\oe i}{\textlengthmark}{\textsubring{d}}{\textsci}} \dicPos{m}[1] \dicFlx{(‑a, ‑ar)}[1] \dicSynonym{fól\smash{\textsuperscript{1}}} \dicDirectTranslationCS{surovec, surovkyně, hrubián(ka)}
\dicEntry[fax] \dicTerm{fax\smash{\textsuperscript{1}}} \dicIPA{{f}{a}{x}{s}} \dicPos{n}[2] \dicFlx{(‑, föx)}[23] \dicDirectTranslationCS{hříva}
\dicEntry[fax] \dicTerm{fax\smash{\textsuperscript{2}}} \dicIPA{{f}{a}{x}{s}} \dicPos{n}[2] \dicFlx{(‑, föx)}[23] \textbf{1.} \dicSynonym{faxtæki} \dicDirectTranslationCS{(tele)fax} \dicIndirectTranslationCS{(zařízení)}  \textbf{2.} \dicDirectTranslationCS{(tele)fax} \dicIndirectTranslationCS{(dokument poslaný faxem)}
\dicEntry[faxtæki] \dicTerm{fax··tæki} \dicIPA{{f}{a}{x}{s}{t\smash{\textsuperscript{h}}}{a}{i}{\r{\textObardotlessj}}{\textsci}} \dicPos{n}[2] \dicFlx{(‑s, ‑)}[16] \dicDirectTranslationCS{(tele)fax, faxovací zařízení}
\dicEntry[fá] \dicTerm{fá} \dicsymFrequent\  \dicIPA{{f}{au}{\textlengthmark}} \dicPos{v}[7] \dicFlx{(fæ, fékk, fengum, fengi, fengið)}[6] \dicFlx{acc} \textbf{1.} \dicSynonym*{taka við} \dicDirectTranslationCS{dostat, dostávat, obdržet} \dicExampleIS{fá bréf} \dicExampleCS{dostat dopis};  \dicPhraseIS{fá e‑ð lánað} \dicDirectTranslationCS{(vy)půjčit si (co)} \dicExampleIS{Má ég fá bílinn lánaðan?} \dicExampleCS{Můžu si půjčit auto?}  \textbf{2.} \dicSynonym*{fá leyfi} \dicDirectTranslationCS{moct, moci, smět, dostat svolení}  \textbf{3.} \dicDirectTranslationCS{vzít si, pořídit si};  \dicPhraseIS{fá sér sæti} \dicDirectTranslationCS{sednout si, posadit se};  \dicPhraseIS{fá sér e‑ð} {\textbf{a.}} \dicDirectTranslationCS{vzít si (co) k~jídlu\,/\addthin pití};  {\textbf{b.}} \dicDirectTranslationCS{získat (co), pořídit si (co)} \dicExampleIS{fá sér nýja tölvu} \dicExampleCS{pořídit si nový počítač}  \textbf{4.} \dicDirectTranslationCS{dát, dávat, podat, podávat};  \dicPhraseIS{fá e‑m e‑ð} \dicSynonym{afhenda} \dicDirectTranslationCS{podat (komu co)} \dicExampleIS{fá henni peningana} \dicExampleCS{podat jí peníze};  \dicIdiom{fá}[að]{ \dicPhraseIS{fá að láni}} \dicDirectTranslationCS{(vy)půjčit si} \dicExampleIS{Ég fékk að láni bók hjá Ellu.} \dicExampleCS{Půjčil jsem si knížku od Elly.};  \dicIdiom{fá}[af]{ \dicPhraseIS{fá e‑ð af sér}} \dicDirectTranslationCS{dovolit si (co)} \dicIndirectTranslationCS{(používané v~otázce nebo záporu)} \dicExampleIS{Ég fæ ekki af mér að láta barnið vera heima.} \dicExampleCS{Nedovolím si nechat dítě doma.};  \dicIdiom{fá}[á]{ \dicPhraseIS{e‑að fær á e‑n}} \dicSynonym*{hafa neikvæð áhrif} \dicDirectTranslationCS{(co) se (koho) hluboce dotýká, (co) má na (koho) (špatný) dopad} \dicExampleIS{Slysið fékk á hann.} \dicExampleCS{Ta nehoda na něj měla hluboký dopad.}; { \dicPhraseIS{fá á baukinn}} \dicLangCat{přen.} \dicDirectTranslationCS{dostat přes prsty};  \dicIdiom{fá}[fyrir]{ \dicPhraseIS{fá fyrir ferðina}} \dicLangCat{přen.} \dicDirectTranslationCS{dostat zaplaceno, dostat zaslouženou odměnu (za zlobení ap.)}; { \dicPhraseIS{fá e‑ð fyrir e‑ð}} \dicDirectTranslationCS{dostat (co) za (co)};  \dicIdiom{fá}[inni]{ \dicPhraseIS{fá inni}} \dicDirectTranslationCS{sehnat ubytování} \dicExampleIS{fá inni á gistihúsi} \dicExampleCS{sehnat ubytování v~penzionu};  \dicIdiom{fá}[í]{ \dicPhraseIS{fá e‑n í e‑ð}} \dicDirectTranslationCS{přimět (koho) k~(čemu)};  \dicIdiom{fá}[ofan af]{ \dicPhraseIS{fá e‑n ofan af e‑u}} \dicDirectTranslationCS{přesvědčit\,/\addthin přimět (koho) k~(čemu)};  \dicIdiom{fá}[til]{ \dicPhraseIS{fá e‑n til e‑s}} \dicDirectTranslationCS{přimět\,/\addthin přemluvit (koho) k~(čemu)};  \dicIdiom{fá}[til baka]{ \dicPhraseIS{fá e‑ð til baka}} \dicDirectTranslationCS{získat (co) nazpátek (drobné ap.)};  \dicIdiom{fá}[upp úr]{ \dicPhraseIS{fá e‑ð upp úr e‑m}} \dicSynonym*{afla vitneskju} \dicDirectTranslationCS{dostat (co) z~(koho) (přiznání ap.), vytáhnout (co) z~(koho)} \dicExampleIS{Lögreglan fékk ekkert upp úr honum við yfirheyrsluna.} \dicExampleCS{Policie z~něj při výslechu nic nedostala.};  \dicIdiom{fá}[út]{ \dicPhraseIS{fá e‑ð út}} \dicDirectTranslationCS{spočítat (co), vypočítat (co)};  \dicIdiom{fá}[út úr]{ \dicPhraseIS{fá e‑ð út úr e‑u}} \dicDirectTranslationCS{získat (co) z~(čeho), mít (co) z~(čeho)};  \dicIdiom{fást}{ \dicPhraseIS{fást}} \dicFlx{refl} \dicSynonym*{vera fáanlegt} \dicDirectTranslationCS{dostat se, být k~dostání} \dicExampleIS{Þeir fást bara ekki.} \dicExampleCS{Prostě nejsou k~dostání.};  \dicIdiom{fást}[til]{ \dicPhraseIS{fást til e‑u}} \dicFlx{refl} \dicDirectTranslationCS{chtít (co), být ochotný k~(čemu)};  \dicIdiom{fást}[um]{ \dicPhraseIS{fást ekki um e‑ð}} \dicFlx{refl} \dicDirectTranslationCS{nelámat si s~(čím) hlavu, pustit (co) z~hlavy} \dicExampleIS{Fástu ekki um það!} \dicExampleCS{Nelam si s~tím hlavu!};  \dicIdiom{fást}[við]{ \dicPhraseIS{fást við e‑ð}} \dicFlx{refl} \dicDirectTranslationCS{věnovat se (čemu), pracovat na (čem), zabývat se (čím)} \dicExampleIS{fást við húsbyggingar} \dicExampleCS{zabývat se výstavbou domů}
\dicEntry[fáanlegur] \dicTerm{fáan··legur} \dicIPA{{f}{au}{\textlengthmark}{a}{n}{l}{\textepsilon}{\textbabygamma}{\textscy}{\textsubring{r}}} \dicPos{adj}[1]\dicFlx{}[-8] \textbf{1.} \dicDirectTranslationCS{(jsoucí) k~dostání, dostupný}  \textbf{2.} \dicDirectTranslationCS{ochotný, vstřícný}
\dicEntry[fábjáni] \dicTerm{fá··bján|i} \dicIPA{{f}{au}{\textlengthmark}{\textsubring{b}}{j}{au}{n}{\textsci}} \dicPos{m}[1] \dicFlx{(‑a, ‑ar)}[1] \dicSynonym{fáviti} \dicDirectTranslationCS{idiot, blbec, pitomec}
\dicEntry[fábreytilegur] \dicTerm{fá··breyti·legur} \dicIPA{{f}{au}{\textlengthmark}{\textsubring{b}}{r}{ei}{\textsubring{d}}{\textsci}{l}{\textepsilon}{\textbabygamma}{\textscy}{\textsubring{r}}} \dicPos{adj}[1]\dicFlx{}[-8] \dicDirectTranslationCS{jednotvárný, monotónní}
\dicEntry[fábreytni] \dicTerm{fá··breytn|i} \dicIPA{{f}{au}{\textlengthmark}{\textsubring{b}}{r}{ei}{h}{\textsubring{d}}{n}{\textsci}} \dicPos{f}[3] \dicFlx{(‑i)}[3] \dicDirectTranslationCS{jednotvárnost, monotónnost}
\dicEntry[fábrotinn] \dicTerm{fá··brotinn} \dicIPA{{f}{au}{\textlengthmark}{\textsubring{b}}{r}{\textopeno}{\textsubring{d}}{\textsci}{\textsubring{n}}} \dicPos{adj}[6]\dicFlx{}[-2] \dicDirectTranslationCS{jednoduchý, nepříliš komplikovaný, obyčejný, všední} \dicExampleIS{fábrotinn matur} \dicExampleCS{jednoduché jídlo}
\dicEntry[fádæmi] \dicTerm{fá··dæmi} \dicIPA{{f}{au}{\textlengthmark}{\textsubring{d}}{a}{i}{m}{\textsci}} \dicPos{n}[2] \dicFlx{pl}[19] \dicDirectTranslationCS{výjimečnost, nezvyklost};  \dicPhraseIS{e‑að er með fádæmum} \dicDirectTranslationCS{(co) je neslýchané}
\dicEntry[fáeinir] \dicTerm{fá··einir} \dicsymFrequent\  \dicIPA{{f}{au}{\textlengthmark}{ei}{n}{\textsci}{\textsubring{r}}} \dicPos{pron}[10] \dicFlx{indef pl} \dicDirectTranslationCS{několik, pár} \dicExampleIS{bíða í fáeina tíma} \dicExampleCS{počkat několik hodin}
\dicEntry[fáfróður] \dicTerm{fá··fróður} \dicIPA{{f}{au}{\textlengthmark}{f}{r}{ou}{ð}{\textscy}{\textsubring{r}}} \dicPos{adj}[2]\dicFlx{}[-6] \dicSynonym{fákunnandi} \dicDirectTranslationCS{nevzdělaný, nevědomý, neinformovaný, neznalý}
\dicEntry[fáfræði] \dicTerm{fá··fræð|i} \dicIPA{{f}{au}{\textlengthmark}{f}{r}{a}{i}{ð}{\textsci}} \dicPos{f}[3] \dicFlx{(‑i)}[3] \dicSynonym{fákunnátta} \dicDirectTranslationCS{neznalost, nevědomost, nevzdělanost, neinformovanost}
\dicEntry[fága] \dicTerm{fág|a} \dicIPA{{f}{au}{\textlengthmark}{a}} \dicPos{v}[1] \dicFlx{(‑aði)}[1] \dicFlx{acc} \dicDirectTranslationCS{(vy)čistit, (vy)leštit} \dicExampleIS{fága yfirborðið} \dicExampleCS{leštit povrch}
\dicEntry[fágaður] \dicTerm{fág··|aður} \dicIPA{{f}{au}{\textlengthmark}{a}{ð}{\textscy}{\textsubring{r}}} \dicPos{adj}[3] \dicFlx{(f ‑uð)}[3] \textbf{1.} \dicDirectTranslationCS{uhlazený, kultivovaný, vytříbený} \dicExampleIS{fágaðar hreyfingar} \dicExampleCS{uhlazené pohyby}  \textbf{2.} \dicSynonym*{fægður} \dicDirectTranslationCS{nablýskaný, vyčištěný}
\dicEntry[fágun] \dicTerm{fág|un} \dicIPA{{f}{au}{\textlengthmark}{\textscy}{\textsubring{n}}} \dicPos{f}[7] \dicFlx{(‑unar)}[9] \textbf{1.} \dicDirectTranslationCS{kultivovanost, vytříbenost}  \textbf{2.} \dicDirectTranslationCS{čištění, leštění}
\dicEntry[fágæti] \dicTerm{fá··gæti} \dicIPA{{f}{au}{\textlengthmark}{\r{\textObardotlessj}}{a}{i}{\textsubring{d}}{\textsci}} \dicPos{n}[2] \dicFlx{(‑s)}[20] \dicDirectTranslationCS{vzácnost, jedinečnost, unikát}
\dicEntry[fágætur] \dicTerm{fá··gætur} \dicIPA{{f}{au}{\textlengthmark}{\r{\textObardotlessj}}{a}{i}{\textsubring{d}}{\textscy}{\textsubring{r}}} \dicPos{adj}[1]\dicFlx{}[-1] \dicDirectTranslationCS{vzácný, unikátní, jedinečný} \dicExampleIS{fágætt safn} \dicExampleCS{unikátní sbírka}
\dicEntry[fáheyrður] \dicTerm{fá··heyrður} \dicIPA{{f}{au}{\textlengthmark}{h}{ei}{r}{ð}{\textscy}{\textsubring{r}}} \dicPos{adj}[2]\dicFlx{}[-1] \dicDirectTranslationCS{neslýchaný} \dicExampleIS{fáheyrður dónaskapur} \dicExampleCS{neslýchaná hrubost}
\dicEntry[fáklæddur] \dicTerm{fá··klæddur} \dicIPA{{f}{au}{\textlengthmark}{k\smash{\textsuperscript{h}}}{l}{a}{i}{\textsubring{d}}{\textscy}{\textsubring{r}}} \dicPos{adj}[2]\dicFlx{}[-18] \dicDirectTranslationCS{spoře oděný\,/\addthin oblečený}
\dicEntry[fákunnandi] \dicTerm{fá··kunn·andi} \dicIPA{{f}{au}{\textlengthmark}{k\smash{\textsuperscript{h}}}{\textscy}{n}{a}{n}{\textsubring{d}}{\textsci}} \dicPos{adj}[13] \dicFlx{indecl}[1] \dicSynonym{fáfróður} \dicDirectTranslationCS{nevědomý, nevzdělaný}
\dicEntry[fákunnátta] \dicTerm{fá··kunn·átt|a} \dicIPA{{f}{au}{\textlengthmark}{k\smash{\textsuperscript{h}}}{\textscy}{n}{au}{h}{\textsubring{d}}{a}} \dicPos{f}[1] \dicFlx{(‑u)}[5] \dicSynonym{fáfræði} \dicDirectTranslationCS{nevědomost, neznalost}
\dicEntry[fákur] \dicTerm{fák|ur} \dicIPA{{f}{au}{\textlengthmark}{\r{g}}{\textscy}{\textsubring{r}}} \dicPos{m}[6] \dicFlx{(‑s, ‑ar)}[24] \dicLangCat{básn.} \dicDirectTranslationCS{oř, kůň}
\dicEntry[fálátur] \dicTerm{fá··látur} \dicIPA{{f}{au}{\textlengthmark}{l}{au}{\textsubring{d}}{\textscy}{\textsubring{r}}} \dicPos{adj}[1]\dicFlx{}[-1] \dicSynonym{fámáll} \dicDirectTranslationCS{odměřený, nepřístupný, rezervovaný}
\dicEntry[fálki] \dicTerm{fálk|i} \dicIPA{{f}{au}{\textsubring{l}}{\r{\textObardotlessj}}{\textsci}} \dicPos{m}[1] \dicFlx{(‑a, ‑ar)}[1] \dicFieldCat{zool.} \dicDirectTranslationCS{raroh, raroh lovecký} \textit{(l.~{\textLA{Falco rusticolus}})}  \dicsymPhoto\ 
\dicFigure{29595.jpg}{Fálki}{Fálki - Velechovský Marek, Biolib, Copyright/CC-BY-NC}
\dicEntry[fálma] \dicTerm{fálm|a} \dicIPA{{f}{au}{l}{m}{a}} \dicPos{v}[1] \dicFlx{(‑aði)}[1] \dicDirectTranslationCS{šmátrat, tápat} \dicExampleIS{fálma í kringum sig} \dicExampleCS{tápat kolem sebe};  \dicPhraseIS{fálma eftir e‑u} \dicDirectTranslationCS{šmátrat po (čem)}
\dicEntry[fálmari] \dicTerm{fálm··ar|i} \dicIPA{{f}{au}{l}{m}{a}{r}{\textsci}} \dicPos{m}[1] \dicFlx{(‑a, ‑ar)}[13] \dicFieldCat{biol.} \dicDirectTranslationCS{tykadlo}
\dicEntry[fálmkenndur] \dicTerm{fálm··kenndur} \dicIPA{{f}{au}{l}{m}{c\smash{\textsuperscript{h}}}{\textepsilon}{n}{\textsubring{d}}{\textscy}{\textsubring{r}}} \dicPos{adj}[2]\dicFlx{}[-14] \dicDirectTranslationCS{tápající, šmátrající}
\dicEntry[fálæti] \dicTerm{fá··læti} \dicIPA{{f}{au}{\textlengthmark}{l}{a}{i}{\textsubring{d}}{\textsci}} \dicPos{n}[2] \dicFlx{(‑s)}[20] \dicDirectTranslationCS{nezájem, neúčast, lhostejnost} \dicExampleIS{sýna e‑m fálæti} \dicExampleCS{projevovat o~(koho) nezájem}
\dicEntry[fámáll] \dicTerm{fá··máll} \dicIPA{{f}{au}{\textlengthmark}{m}{au}{\textsubring{d}}{\textsubring{l}}} \dicPos{adj}[8]\dicFlx{}[-1] \dicSynonym{fátalaður} \dicDirectTranslationCS{zamlklý, tichý, málomluvný, mlčenlivý}
\dicEntry[fámenni] \dicTerm{fá··menni} \dicIPA{{f}{au}{\textlengthmark}{m}{\textepsilon}{n}{\textsci}} \dicPos{n}[2] \dicFlx{(‑s)}[20] \dicDirectTranslationCS{pár\,/\addthin málo lidí}
\dicEntry[fámennur] \dicTerm{fá··mennur} \dicIPA{{f}{au}{\textlengthmark}{m}{\textepsilon}{n}{\textscy}{\textsubring{r}}} \dicPos{adj}[1]\dicFlx{}[-1] \dicDirectTranslationCS{nepočetný, málo početný (rodina ap.)} \dicExampleIS{fámennt heimili} \dicExampleCS{nepočetná domácnost}
\dicEntry[fána] \dicTerm{fán|a} \dicIPA{{f}{au}{\textlengthmark}{n}{a}} \dicPos{f}[1] \dicFlx{(‑u, ‑ur)}[7] \dicDirectTranslationCS{fauna, zvířena, zvířectvo}
\dicEntry[fánaberi] \dicTerm{fána··ber|i} \dicIPA{{f}{au}{\textlengthmark}{n}{a}{\textsubring{b}}{\textepsilon}{r}{\textsci}} \dicPos{m}[1] \dicFlx{(‑a, ‑ar)}[1] \dicDirectTranslationCS{vlajkonoš(ka)}
\dicEntry[fáni] \dicTerm{fán|i} \dicIPA{{f}{au}{\textlengthmark}{n}{\textsci}} \dicPos{m}[1] \dicFlx{(‑a, ‑ar)}[1] \dicDirectTranslationCS{vlajka, prapor} \dicExampleIS{Fáninn blaktir í vindinum.} \dicExampleCS{Vlajka vlaje ve větru.}
\dicEntry[fánýti] \dicTerm{fá··nýti} \dicIPA{{f}{au}{\textlengthmark}{n}{i}{\textsubring{d}}{\textsci}} \dicPos{n}[2] \dicFlx{(‑s)}[20] \dicSynonym{hjóm} \dicDirectTranslationCS{bezvýslednost, marnost}
\dicEntry[fánýtur] \dicTerm{fá··nýtur} \dicIPA{{f}{au}{\textlengthmark}{n}{i}{\textsubring{d}}{\textscy}{\textsubring{r}}} \dicPos{adj}[1]\dicFlx{}[-1] \dicDirectTranslationCS{marný, bezcenný, (jsoucí) k~ničemu, neužitečný} \dicExampleIS{fánýtt starf} \dicExampleCS{práce k~ničemu}
\dicEntry[fáorður] \dicTerm{fá··orður} \dicIPA{{f}{au}{\textlengthmark}{\textopeno}{r}{ð}{\textscy}{\textsubring{r}}} \dicPos{adj}[2]\dicFlx{}[-1] \dicDirectTranslationCS{málomluvný, mlčenlivý}
\dicEntry[fár] \dicTerm{fár\smash{\textsuperscript{1}}} \dicIPA{{f}{au}{\textlengthmark}{\textsubring{r}}} \dicPos{n}[2] \dicFlx{(‑s)}[2] \dicDirectTranslationCS{povyk, rozruch, pozdvižení}
\dicEntry[fár] \dicTerm{fár\smash{\textsuperscript{2}}} \dicsymFrequent\  \dicIPA{{f}{au}{\textlengthmark}{\textsubring{r}}} \dicPos{adj}[10] \dicFlx{(comp færri, sup fæstur)}[14] \textbf{1.} \dicDirectTranslationCS{nemnohý, nepočetný};  \dicPhraseIS{það er fátt (að frétta)} \dicDirectTranslationCS{není toho mnoho (k~vyprávění)}  \textbf{2.} \dicDirectTranslationCS{odměřený, rezervovaný};  \dicPhraseIS{verða fár við} \dicDirectTranslationCS{zamlknout}
\dicEntry[fáráðlingur] \dicTerm{fá··ráð·ling|ur} \dicIPA{{f}{au}{\textlengthmark}{r}{au}{ð}{l}{i}{\ng}{\r{g}}{\textscy}{\textsubring{r}}} \dicPos{m}[6] \dicFlx{(‑s, ‑ar)}[8] \dicSynonym{kjáni} \dicDirectTranslationCS{ťulpas, trouba, hňup}
\dicEntry[fáránlegur] \dicTerm{fárán··legur} \dicsymFrequent\  \dicIPA{{f}{au}{\textlengthmark}{r}{au}{n}{l}{\textepsilon}{\textbabygamma}{\textscy}{\textsubring{r}}} \dicPos{adj}[1]\dicFlx{}[-8] \dicDirectTranslationCS{absurdní, nesmyslný} \dicExampleIS{Þetta er fáránlegt viðhorf.} \dicExampleCS{To je absurdní názor.}
\dicEntry[fáránleiki] \dicTerm{fárán··leik|i} \dicIPA{{f}{au}{r}{au}{n}{l}{ei}{\r{\textObardotlessj}}{\textsci}} \dicPos{m}[1] \dicFlx{(‑a)}[3] \dicDirectTranslationCS{absurdita, nesmyslnost}
\dicEntry[fárveikur] \dicTerm{fár··veikur} \dicIPA{{f}{au}{r}{v}{ei}{\r{g}}{\textscy}{\textsubring{r}}} \dicPos{adj}[1]\dicFlx{}[-6] \dicDirectTranslationCS{vážně nemocný}
\dicEntry[fárviðri] \dicTerm{fár··viðri} \dicIPA{{f}{au}{r}{v}{\textsci}{ð}{r}{\textsci}} \dicPos{n}[2] \dicFlx{(‑s, ‑)}[14] \dicDirectTranslationCS{orkán} \dicIndirectTranslationCS{(Beaufort 12) (od 32,7~m/s)}
\dicEntry[fáséður] \dicTerm{fá··séður} \dicIPA{{f}{au}{\textlengthmark}{s}{j}{\textepsilon}{ð}{\textscy}{\textsubring{r}}} \dicPos{adj}[2]\dicFlx{}[-12] \dicDirectTranslationCS{neobvyklý, nezvyklý, málo vídaný, vzácný}
\dicEntry[fásinna] \dicTerm{fá··sinn|a} \dicIPA{{f}{au}{\textlengthmark}{s}{\textsci}{n}{a}} \dicPos{f}[1] \dicFlx{(‑u)}[5] \dicSynonym{heimska} \dicDirectTranslationCS{nerozum, pošetilost}
\dicEntry[fásinni] \dicTerm{fá··sinni} \dicIPA{{f}{au}{\textlengthmark}{s}{\textsci}{n}{\textsci}} \dicPos{n}[2] \dicFlx{(‑s)}[20] \dicDirectTranslationCS{osamění, opuštěnost, samota} \dicExampleIS{í fásinni smábæjarins} \dicExampleCS{v~opuštěném malém městečku}
\dicEntry[fáskiptinn] \dicTerm{fá··skiptinn} \dicIPA{{f}{au}{\textlengthmark}{s}{\r{\textObardotlessj}}{\textsci}{f}{\textsubring{d}}{\textsci}{\textsubring{n}}} \dicPos{adj}[6]\dicFlx{}[-2] \dicSynonym{fálátur} \dicDirectTranslationCS{odměřený, rezervovaný, stranící se} \dicAntonym{afskiptasamur}
\dicEntry[fáskrúðugur] \textls[15]{\dicTerm{fá··skrúð·ugur} \dicIPA{{f}{au}{\textlengthmark}{s}{\r{g}}{r}{u}{ð}{\textscy}{\textbabygamma}{\textscy}{\textsubring{r}}} \dicPos{adj}[1]\dicFlx{}[-8] \dicDirectTranslationCS{fádní, (jsoucí) beze změn}} \dicAntonym{fjölskrúðugur}
\dicEntry[fát] \dicTerm{fát} \dicIPA{{f}{au}{\textlengthmark}{\textsubring{d}}} \dicPos{n}[2] \dicFlx{(‑s)}[2] \dicDirectTranslationCS{zmatek, zmatenost, popletenost};  \dicPhraseIS{það kemur fát á e‑n} \dicFlx{impers} \dicDirectTranslationCS{(kdo) zmatkuje}
\dicEntry[fátalaður] \dicTerm{fá··|tal·aður} \dicIPA{{f}{au}{\textlengthmark}{t\smash{\textsuperscript{h}}}{a}{l}{a}{ð}{\textscy}{\textsubring{r}}} \dicPos{adj}[3] \dicFlx{(f ‑töluð)}[1] \dicSynonym{fámáll} \dicDirectTranslationCS{tichý, mlčenlivý, zamlklý}
\dicEntry[fátíður] \dicTerm{fá··tíður} \dicIPA{{f}{au}{\textlengthmark}{t\smash{\textsuperscript{h}}}{i}{ð}{\textscy}{\textsubring{r}}} \dicPos{adj}[2]\dicFlx{}[-6] \dicSynonym{sjaldgæfur} \dicDirectTranslationCS{vzácný, ojedinělý}
\dicEntry[fátæklegur] \dicTerm{fá·tæk··legur} \dicIPA{{f}{au}{\textlengthmark}{t\smash{\textsuperscript{h}}}{a}{i}{h}{\r{g}}{l}{\textepsilon}{\textbabygamma}{\textscy}{\textsubring{r}}} \dicPos{adj}[1]\dicFlx{}[-8] \dicSynonym{aumur} \dicDirectTranslationCS{ubohý, mizerný, nevalný}
\dicEntry[fátæklingur] \dicTerm{fá·tæk··ling|ur} \dicIPA{{f}{au}{\textlengthmark}{t\smash{\textsuperscript{h}}}{a}{i}{h}{\r{g}}{l}{i}{\ng}{\r{g}}{\textscy}{\textsubring{r}}} \dicPos{m}[6] \dicFlx{(‑s, ‑ar)}[8] \dicSynonym{öreigi} \dicDirectTranslationCS{chudák, chudačka, ubožák, ubožačka}
\dicEntry[fátækt] \dicTerm{fá··tækt} \dicIPA{{f}{au}{\textlengthmark}{t\smash{\textsuperscript{h}}}{a}{i}{x}{\textsubring{d}}} \dicPos{f}[4] \dicFlx{(‑ar)}[3] \dicSynonym{eymd} \dicDirectTranslationCS{chudoba, bída}
\dicEntry[fátæktarmörk] \dicTerm{fá·tæktar··mörk} \dicIPA{{f}{au}{\textlengthmark}{t\smash{\textsuperscript{h}}}{a}{i}{x}{\textsubring{d}}{a}{r}{m}{\oe}{\textsubring{r}}{\r{g}}} \dicPos{n}[2] \dicFlx{pl}[9] \dicDirectTranslationCS{hranice chudoby}
\dicEntry[fátækur] \dicTerm{fá··tækur} \dicsymFrequent\  \dicIPA{{f}{au}{\textlengthmark}{t\smash{\textsuperscript{h}}}{a}{i}{\r{g}}{\textscy}{\textsubring{r}}} \dicPos{adj}[1]\dicFlx{}[-1] \dicDirectTranslationCS{chudý, ubohý, bídný, mizerný} \dicExampleIS{Þeir voru fátækastir allra þjóða.} \dicExampleCS{Byli nejchudší z~národů.}
\dicEntry[fáviska] \dicTerm{fá··visk|a} \dicIPA{{f}{au}{\textlengthmark}{v}{\textsci}{s}{\r{g}}{a}} \dicPos{f}[1] \dicFlx{(‑u)}[5] \dicSynonym{fáfræði} \dicDirectTranslationCS{neznalost, nevědomost} \dicExampleIS{Fyrirgefðu mér fáviskuna.} \dicExampleCS{Omluvte moji nevědomost.}
\dicEntry[fáviti] \dicTerm{fá··vit|i} \dicIPA{{f}{au}{\textlengthmark}{v}{\textsci}{\textsubring{d}}{\textsci}} \dicPos{m}[1] \dicFlx{(‑a, ‑ar)}[1] \textbf{1.} \dicSynonym{bjáni} \dicDirectTranslationCS{idiot, imbecil, blbec}  \textbf{2.} \dicDirectTranslationCS{člověk s~mentální retardací}
\dicEntry[fávís] \dicTerm{fá··vís} \dicIPA{{f}{au}{\textlengthmark}{v}{i}{s}} \dicPos{adj}[5]\dicFlx{}[-1] \dicSynonym{fáfróður} \dicDirectTranslationCS{nevědomý, neinformovaný, neznalý}
\dicEntry[feb.] \dicTerm{feb.} \dicPos{zkr} \dicPhraseIS{febrúar} \dicDirectTranslationCS{únor}
\dicEntry[febrúar] \dicTerm{febrúar} \dicsymFrequent\  \dicIPA{{f}{\textepsilon}{\textlengthmark}{\textsubring{b}}{r}{u}{a}{\textsubring{r}}} \dicPos{m}[15] \dicFlx{indecl}[2] \dicDirectTranslationCS{únor} \dicExampleIS{í fyrstu viku febrúar} \dicExampleCS{v~prvním únorovém týdnu}
\dicEntry[febrúarmánuður] \dicTerm{febrúar··mán·|uður} \dicIPA{{f}{\textepsilon}{\textlengthmark}{\textsubring{b}}{r}{u}{a}{r}{m}{au}{n}{\textscy}{ð}{\textscy}{\textsubring{r}}} \dicPos{m}[10] \dicFlx{(‑aðar, ‑uð\-ir)}[42] \dicDirectTranslationCS{(měsíc) únor}
\dicEntry[feðgar] \dicTerm{feðgar} \dicIPA{{f}{\textepsilon}{ð}{\r{g}}{a}{\textsubring{r}}} \dicPos{m}[14] \dicFlx{pl}[1] \dicDirectTranslationCS{otec a~syn(ové)}
\dicEntry[feðgin] \dicTerm{feðgin} \dicIPA{{f}{\textepsilon}{ð}{\r{\textObardotlessj}}{\textsci}{\textsubring{n}}} \dicPos{n}[2] \dicFlx{pl}[1] \dicDirectTranslationCS{otec a~dcera\,/\addthin dcery}
\dicEntry[feðra] \dicTerm{feðr|a} \dicIPA{{f}{\textepsilon}{ð}{r}{a}} \dicPos{v}[1] \dicFlx{(‑aði)}[1] \dicFlx{acc} \dicDirectTranslationCS{určit otce} \dicExampleIS{feðra barnið} \dicExampleCS{určit otce dítěte}
\dicEntry[feðralaun] \dicTerm{feðra··laun} \dicIPA{{f}{\textepsilon}{ð}{r}{a}{l}{\oe i}{\textsubring{n}}} \dicPos{n}[2] \dicFlx{pl}[1] \dicDirectTranslationCS{příspěvky na dítě} \dicIndirectTranslationCS{(určené otcům samoživitelům)}
\dicEntry[feðraveldi] \dicTerm{feðra··veldi} \dicIPA{{f}{\textepsilon}{ð}{r}{a}{v}{\textepsilon}{l}{\textsubring{d}}{\textsci}} \dicPos{n}[2] \dicFlx{(‑s, ‑)}[14] \dicDirectTranslationCS{patriarchát}
\dicEntry[feður] \dicTerm{feður} \dicIPA{{f}{\textepsilon}{\textlengthmark}{ð}{\textscy}{\textsubring{r}}} \dicPos{m} \dicFlx{pl nom} \dicLink{faðir}
\dicEntry[feginn] \dicTerm{feginn} \dicsymFrequent\  \dicIPA{{f}{ei}{\textsci}{\textsubring{n}}} \dicPos{adj}[6]\dicFlx{}[-2] \dicDirectTranslationCS{rád, potěšený} \dicExampleIS{Ég er svo feginn að sjá þig.} \dicExampleCS{Moc rád tě vidím.};  \dicPhraseIS{taka e‑u fegins hendi} \dicLangCat{přen.} \dicDirectTranslationCS{přivítat (co) s~otevřenou náručí}
\dicEntry[fegra] \dicTerm{fegr|a} \dicIPA{{f}{\textepsilon}{\textbabygamma}{r}{a}} \dicPos{v}[1] \dicFlx{(‑aði)}[1] \dicFlx{acc} \dicDirectTranslationCS{(z)krášlit, (vy)zdobit} \dicExampleIS{fegra umhverfið} \dicExampleCS{zkrášlit okolí};  \dicPhraseIS{fegra e‑ð fyrir e‑m} \dicLangCat{přen.} \dicDirectTranslationCS{přikrášlit (co) před (kým)} \dicExampleIS{fegra dvölina fyrir sér} \dicExampleCS{přikrášlit si pobyt}
\dicEntry[fegrun] \dicTerm{fegr|un} \dicIPA{{f}{\textepsilon}{\textbabygamma}{r}{\textscy}{\textsubring{n}}} \dicPos{f}[7] \dicFlx{(‑unar)}[9] \dicDirectTranslationCS{zkrášlení, vyzdobení, výzdoba}
\dicEntry[fegrunaraðgerð] \dicTerm{fegrunar··að·gerð} \dicIPA{{f}{\textepsilon}{\textbabygamma}{r}{\textscy}{n}{a}{r}{a}{ð}{\r{\textObardotlessj}}{\textepsilon}{r}{\texttheta}} \dicPos{f}[7] \dicFlx{(‑ar, ‑ir)}[1] \dicDirectTranslationCS{plastická operace\,/\addthin chirurgie} \dicExampleIS{gangast undir fegrunaraðgerð} \dicExampleCS{podstoupit plastickou operaci}
\dicEntry[fegurð] \dicTerm{fegurð} \dicsymFrequent\  \dicIPA{{f}{\textepsilon}{\textlengthmark}{\textbabygamma}{\textscy}{r}{\texttheta}} \dicPos{f}[4] \dicFlx{(‑ar)}[3] \dicDirectTranslationCS{krása, půvab} \dicExampleIS{fegurð náttúrunnar} \dicExampleCS{krása přírody}
\dicEntry[fegurðarsamkeppni] \dicTerm{fegurðar··sam·keppn|i} \dicIPA{{f}\-{\textepsilon}\-{\textlengthmark}\-{\textbabygamma}\-{\textscy}\-{r}\-{ð}\-{a}\-{\textsubring{r}}\-{s}\-{a}\-{m}\-{c\smash{\textsuperscript{h}}}\-{\textepsilon}\-{h}\-{\textsubring{b}}\-{n}\-{\textsci}\-} \dicPos{f}[3] \dicFlx{(‑i, ‑ir)}[1] \dicDirectTranslationCS{soutěž krásy}
\dicEntry[fegurstur] \dicTerm{fegurstur} \dicIPA{{f}{\textepsilon}{\textlengthmark}{\textbabygamma}{\textscy}{\textsubring{r}}{s}{\textsubring{d}}{\textscy}{\textsubring{r}}} \dicPos{adj} \dicFlx{m sg nom sup} \dicLink{fagur}
\dicEntry[feigð] \dicTerm{feigð} \dicIPA{{f}{ei}{\textbabygamma}{\texttheta}} \dicPos{f}[4] \dicFlx{(‑ar)}[3] \dicDirectTranslationCS{(bezprostřední) smrt, blízká smrt}
\dicEntry[feigur] \dicTerm{feigur} \dicIPA{{f}{ei}{\textlengthmark}{\textbabygamma}{\textscy}{\textsubring{r}}} \dicPos{adj}[1]\dicFlx{}[-1] \dicDirectTranslationCS{blízký smrti, umírající}
\dicEntry[feiki] \dicTerm{feiki-} \dicIPA{{f}{ei}{\textlengthmark}{\r{\textObardotlessj}}{\textsci}} \dicPos{predp} \dicDirectTranslationCS{velmi, mega-}
\dicEntry[feikilegur] \dicTerm{feiki··legur} \dicIPA{{f}{ei}{\textlengthmark}{\r{\textObardotlessj}}{\textsci}{l}{\textepsilon}{\textbabygamma}{\textscy}{\textsubring{r}}} \dicPos{adj}[1]\dicFlx{}[-6] \dicDirectTranslationCS{enormní, kolosální}
\dicEntry[feikn] \dicTerm{feikn} \dicIPA{{f}{ei}{h}{\r{g}}{\textsubring{n}}} \dicPos{f}[4] \dicFlx{(‑ar)}[3] \dicDirectTranslationCS{(neuvěřitelně) mnoho, spousta} \dicExampleIS{feikn af e‑u} \dicExampleCS{spousta (čeho)}
\dicEntry[feila] \dicTerm{feil|a} \dicIPA{{f}{ei}{\textlengthmark}{l}{a}} \dicPos{v}[1] \dicFlx{(‑aði)}[73] \dicLangCat{hovor.} \dicDirectTranslationCS{selhat, neuspět}
\dicEntry[feill] \dicTerm{feil|l} \dicIPA{{f}{ei}{\textsubring{d}}{\textsubring{l}}} \dicPos{m}[6] \dicFlx{(‑s, ‑ar)}[48] \dicLangCat{hovor.} \dicSynonym{mistök} \dicDirectTranslationCS{chyba, omyl, selhání};  \dicPhraseIS{taka feil á e‑m} \dicDirectTranslationCS{splést si (koho)}
\dicEntry[feiminn] \dicTerm{feiminn} \dicIPA{{f}{ei}{\textlengthmark}{m}{\textsci}{\textsubring{n}}} \dicPos{adj}[6]\dicFlx{}[-2] \dicSynonym{kindarlegur} \dicDirectTranslationCS{plachý, nesmělý, ostýchavý};  \dicPhraseIS{vera feiminn við e‑n} \dicDirectTranslationCS{ostýchat se před (kým)}
\dicEntry[feimni] \dicTerm{feimn|i} \dicIPA{{f}{ei}{m}{n}{\textsci}} \dicPos{f}[3] \dicFlx{(‑i)}[3] \dicSynonym*{óframfærni} \dicDirectTranslationCS{plachost, nesmělost, ostýchavost} \dicExampleIS{feimni við ókunnuga} \dicExampleCS{plachost před cizími}
\dicEntry[feimnismál] \dicTerm{feimnis··mál} \dicIPA{{f}{ei}{m}{n}{\textsci}{s}{m}{au}{\textsubring{l}}} \dicPos{n}[2] \dicFlx{(‑s, ‑)}[5] \dicDirectTranslationCS{delikátní\,/\addthin trapná záležitost} \dicIndirectTranslationCS{(téma, o~kterém se nehovoří)}
\dicEntry[feiti] \dicTerm{feit|i} \dicIPA{{f}{ei}{\textlengthmark}{\textsubring{d}}{\textsci}} \dicPos{f}[3] \dicFlx{(‑i)}[3] \dicDirectTranslationCS{tuk}
\dicEntry[feitlaginn] \dicTerm{feit··laginn} \dicIPA{{f}{ei}{\textlengthmark}{\textsubring{d}}{l}{a}{i}{j}{\textsci}{\textsubring{n}}} \dicPos{adj}[6]\dicFlx{}[-3] \dicSynonym{feitur} \dicDirectTranslationCS{korpulentní, tělnatý, tlustý}
\dicEntry[feitletra] \dicTerm{feit··letr|a} \dicIPA{{f}{ei}{\textlengthmark}{\textsubring{d}}{l}{\textepsilon}{\textsubring{d}}{r}{a}} \dicPos{v}[1] \dicFlx{(‑aði)}[1] \dicFlx{acc} \dicDirectTranslationCS{napsat tučně} \dicIndirectTranslationCS{(tučným písmem)}
\dicEntry[feitletraður] \dicTerm{feit··letr·|aður} \dicIPA{{f}{ei}{\textlengthmark}{\textsubring{d}}{l}{\textepsilon}{\textsubring{d}}{r}{a}{ð}{\textscy}{\textsubring{r}}} \dicPos{adj}[3] \dicFlx{(f ‑uð)}[3] \dicDirectTranslationCS{(psaný) tučným písmem}
\dicEntry[feitletur] \dicTerm{feit··letur} \dicIPA{{f}{ei}{\textlengthmark}{\textsubring{d}}{l}{\textepsilon}{\textsubring{d}}{\textscy}{\textsubring{r}}} \dicPos{n}[2] \dicFlx{(‑s, ‑)}[25] \dicDirectTranslationCS{tučné písmo}
\dicEntry[feitmeti] \dicTerm{feit··meti} \dicIPA{{f}{ei}{\textlengthmark}{\textsubring{d}}{m}{\textepsilon}{\textsubring{d}}{\textsci}} \dicPos{n}[2] \dicFlx{(‑s)}[20] \dicSynonym*{feitur matur} \dicDirectTranslationCS{tučné jídlo}
\dicEntry[feitur] \dicTerm{feitur} \dicsymFrequent\  \dicIPA{{f}{ei}{\textlengthmark}{\textsubring{d}}{\textscy}{\textsubring{r}}} \dicPos{adj}[1]\dicFlx{}[-1] \textbf{1.} \dicSynonym{feitlaginn} \dicDirectTranslationCS{tlustý, korpulentní}  \textbf{2.} \dicDirectTranslationCS{tučný (jídlo ap.)} \dicExampleIS{feitur matur} \dicExampleCS{tučné jídlo}  \textbf{3.} \dicDirectTranslationCS{tučný, ziskový}  \textbf{4.} \dicDirectTranslationCS{tučný (písmo ap.)};  \dicPhraseIS{feitt letur} \dicDirectTranslationCS{tučné písmo}
\dicEntry[fel] \dicTerm{fel} \dicIPA{{f}{\textepsilon}{\textlengthmark}{\textsubring{l}}} \dicPos{v} \dicFlx{ind praes sg 1 pers} \textbf{1.} \dicLink{fela\smash{\textsuperscript{1}}}  \textbf{2.} \dicLink{fela\smash{\textsuperscript{2}}}
\dicEntry[fela] \dicTerm{fela\smash{\textsuperscript{1}}} \dicsymFrequent\  \dicIPA{{f}{\textepsilon}{\textlengthmark}{l}{a}} \dicPos{v}[4] \dicFlx{(fel, faldi, földum, feldi, falið)}[1] \dicFlx{acc} \dicDirectTranslationCS{schovat, schovávat, ukrýt, ukrývat} \dicExampleIS{fela hlutinn inni í skáp} \dicExampleCS{schovat věc do skříně};  \dicPhraseIS{fela sig} \dicDirectTranslationCS{schovat se} \dicExampleIS{fela sig á bak við hurð} \dicExampleCS{schovat se za dveřmi}
\dicEntry[fela] \dicTerm{fela\smash{\textsuperscript{2}}} \dicIPA{{f}{\textepsilon}{\textlengthmark}{l}{a}} \dicPos{v}[6] \dicFlx{(fel, fól, fólum, fæli, falið)}[113] \dicFlx{dat + acc} \dicDirectTranslationCS{svěřit, svěřovat, pověřit, pověřovat} \dicIndirectTranslationCS{(dát na starost)};  \dicPhraseIS{fela e‑m e‑ð} \dicDirectTranslationCS{svěřit (komu co)};  \dicPhraseIS{fela e‑m e‑ð á vald} \dicDirectTranslationCS{svěřit (komu co) na milost} \dicExampleIS{fela sig guði á vald} \dicExampleCS{svěřit se do božích rukou};  \dicIdiom{fela}[í]{ \dicPhraseIS{e‑að felur e‑ð í sér}} \dicDirectTranslationCS{(co) v~sobě skrývá (co), (co) v~sobě zahrnuje (co)};  \dicIdiom{felast}{ \dicPhraseIS{felast}} \dicFlx{refl} {\textbf{a.}} \dicDirectTranslationCS{schovat se, schovávat se, ukrýt se, ukrývat se};  {\textbf{b.}} \dicDirectTranslationCS{být obsaženo\,/\addthin zahrnuto}
\dicEntry[feldi] \dicTerm{feldi} \dicIPA{{f}{\textepsilon}{l}{\textsubring{d}}{\textsci}} \dicPos{v} \dicFlx{con pf sg 1 pers} \dicLink{fela\smash{\textsuperscript{1}}}
\dicEntry[feldur] \dicTerm{feld|ur} \dicIPA{{f}{\textepsilon}{l}{\textsubring{d}}{\textscy}{\textsubring{r}}} \dicPos{m}[10] \dicFlx{(‑ar, ‑ir)}[4] \textbf{1.} \dicDirectTranslationCS{srst, kožešina} \dicExampleIS{feldur kattarins} \dicExampleCS{kočičí srst}  \textbf{2.} \dicSynonym{grunn\-ur\smash{\textsuperscript{1}}} \dicDirectTranslationCS{pozadí} \dicExampleIS{hvítur kross á bláum feldi} \dicExampleCS{bílý kříž na modrém pozadí}
\dicEntry[felga] \dicTerm{felg|a} \dicIPA{{f}{\textepsilon}{l}{\r{g}}{a}} \dicPos{f}[1] \dicFlx{(‑u, ‑ur)}[19] \dicDirectTranslationCS{ráfek} \dicIndirectTranslationCS{(obvodová část kola)}
\dicEntry[fell] \dicTerm{fell\smash{\textsuperscript{1}}} \dicIPA{{f}{\textepsilon}{\textsubring{d}}{\textsubring{l}}} \dicPos{n}[2] \dicFlx{(‑s, ‑)}[5] \dicDirectTranslationCS{osamělá hora\,/\addthin kopec}
\dicEntry[fell] \dicTerm{fell\smash{\textsuperscript{2}}} \dicIPA{{f}{\textepsilon}{\textsubring{d}}{\textsubring{l}}} \dicPos{v} \dicFlx{ind praes sg 1 pers} \dicLink{falla}
\dicEntry[fella] \dicTerm{fell|a} \dicsymFrequent\  \dicIPA{{f}{\textepsilon}{\textsubring{d}}{l}{a}} \dicPos{v}[2] \dicFlx{(‑di, ‑t)}[155] \dicFlx{acc} \textbf{1.} \dicSynonym*{láta falla} \dicDirectTranslationCS{povalit, sklátit, (po)kácet, skácet} \dicExampleIS{fella tré} \dicExampleCS{pokácet strom}  \textbf{2.} \dicSynonym{drepa} \dicDirectTranslationCS{zabít} \dicExampleIS{fella marga hermenn í bardaga} \dicExampleCS{zabít v~bitvě mnoho vojáků}  \textbf{3.} \dicDirectTranslationCS{zamítnout (v~hlasování ap.), přehlasovat} \dicExampleIS{fella tillöguna} \dicExampleCS{zamítnout návrh}  \textbf{4.} \dicDirectTranslationCS{nechat propadnout (studenta u~zkoušky ap.)};  \dicIdiom{fella}[brott]{ \dicPhraseIS{fella e‑ð brott}} \dicDirectTranslationCS{vypustit (co), vyškrtnout (co), vymazat (co)};  \dicIdiom{fella}[niður]{ \dicPhraseIS{fella e‑ð niður}} \dicSynonym*{nema burt} \dicDirectTranslationCS{vymazat (co), anulovat (co), vyškrtnout (co)} \dicExampleIS{fella niður skuldir} \dicExampleCS{vymazat dluhy};  \dicIdiom{fella}[saman]{ \dicPhraseIS{fella e‑ð saman}} \dicDirectTranslationCS{spojit (co), skloubit (co)};  \dicIdiom{fella}[við]{ \dicPhraseIS{fella sig ekki við e‑ð}} \dicDirectTranslationCS{nemoct přijmout (co), nemoct se smířit s~(čím)}
\dicEntry[fellibylur] \dicTerm{felli··byl|ur} \dicIPA{{f}{\textepsilon}{\textsubring{d}}{l}{\textsci}{\textsubring{b}}{\textsci}{l}{\textscy}{\textsubring{r}}} \dicPos{m}[9] \dicFlx{(‑s\,/\addthin ‑jar, ‑jir)}[32] \dicFieldCat{meteo.} \dicDirectTranslationCS{hurikán, (tropický) cyklón, tajfun, uragán}
\dicEntry[felling] \dicTerm{fell··ing} \dicIPA{{f}{\textepsilon}{\textsubring{d}}{l}{i}{\ng}{\r{g}}} \dicPos{f}[4] \dicFlx{(‑ar, ‑ar)}[5] \textbf{1.} \dicDirectTranslationCS{záhyb, zahnutí, puk (na šatech ap.)}  \textbf{2.} \dicFieldCat{geol.} \dicDirectTranslationCS{vrása, zvrásnění}  \textbf{3.} \dicFieldCat{chem.} \dicDirectTranslationCS{sedimentace}
\dicEntry[felmtur] \dicTerm{felmtur} \dicIPA{{f}{\textepsilon}{\textsubring{l}}{\textsubring{m}}{\textsubring{d}}{\textscy}{\textsubring{r}}} \dicPos{n}[2] \dicFlx{(‑s)}[28] \dicPhraseIS{vera\,/\addthin verða felmtri sleginn} \dicDirectTranslationCS{vyděsit se, zpanikařit}
\dicEntry[feluleikur] \dicTerm{felu··leik|ur} \dicIPA{{f}{\textepsilon}{\textlengthmark}{l}{\textscy}{l}{ei}{\r{g}}{\textscy}{\textsubring{r}}} \dicPos{m}[6] \dicFlx{(‑s)}[17] \textbf{1.} \dicDirectTranslationCS{hra na schovávanou}  \textbf{2.} \dicSynonym{baktjaldamakk} \dicDirectTranslationCS{utajování}
\dicEntry[felur] \dicTerm{felur} \dicIPA{{f}{\textepsilon}{\textlengthmark}{l}{\textscy}{\textsubring{r}}} \dicPos{f}[1] \dicFlx{pl}[24] \dicDirectTranslationCS{skrývání, ukrývání};  \dicPhraseIS{fara í felur} \dicDirectTranslationCS{skrýt se, schovat se};  \dicPhraseIS{vera í felum} \dicDirectTranslationCS{ukrývat se, schovávat se}
\dicEntry[felustaður] \dicTerm{felu··stað|ur} \dicIPA{{f}{\textepsilon}{\textlengthmark}{l}{\textscy}{s}{\textsubring{d}}{a}{ð}{\textscy}{\textsubring{r}}} \dicPos{m}[10] \dicFlx{(‑ar, ‑ir)}[14] \dicSynonym{fylgsni} \dicDirectTranslationCS{úkryt, skrýš, útočiště}
\dicEntry[femínismi] \dicTerm{femín··ism|i} \dicIPA{{f}{\textepsilon}{\textlengthmark}{m}{i}{n}{\textsci}{s}{m}{\textsci}} \dicPos{m}[1] \dicFlx{(‑a)}[3] \dicDirectTranslationCS{feminismus}
\dicEntry[femínisti] \dicTerm{femín··ist|i} \dicIPA{{f}{\textepsilon}{\textlengthmark}{m}{i}{n}{\textsci}{s}{\textsubring{d}}{\textsci}} \dicPos{m}[1] \dicFlx{(‑a, ‑ar)}[1] \dicDirectTranslationCS{feminista, feministka}
\dicEntry[femínískur] \dicTerm{femínískur} \dicIPA{{f}{\textepsilon}{\textlengthmark}{m}{i}{n}{i}{s}{\r{g}}{\textscy}{\textsubring{r}}} \dicPos{adj}[1]\dicFlx{}[-1] \dicDirectTranslationCS{feministický}
\dicEntry[fen] \dicTerm{fen} \dicIPA{{f}{\textepsilon}{\textlengthmark}{\textsubring{n}}} \dicPos{n}[2] \dicFlx{(‑s, ‑)}[13] \dicSynonym{kviksyndi} \dicDirectTranslationCS{močál, bažina}
\dicEntry[Feneyjar] \dicTerm{Fen··eyjar} \dicIPA{{f}{\textepsilon}{\textlengthmark}{n}{ei}{j}{a}{\textsubring{r}}} \dicPos{f}[4] \dicFlx{pl}[15] \dicFieldCat{geog.} \dicDirectTranslationCS{Benátky} \dicIndirectTranslationCS{(město v~Itálii)}
\dicEntry[fengi] \dicTerm{fengi} \dicIPA{{f}{ei}{\textltailn}{\r{\textObardotlessj}}{\textsci}} \dicPos{v} \dicFlx{con pf sg 1 pers} \dicLink{fá}
\dicEntry[fengið] \dicTerm{fengið} \dicIPA{{f}{ei}{\textltailn}{\r{\textObardotlessj}}{\textsci}{\texttheta}} \dicPos{v} \dicFlx{supin} \dicLink{fá}
\dicEntry[fengrani] \dicTerm{fen··gran|i} \dicIPA{{f}{\textepsilon}{n}{\r{g}}{r}{a}{n}{\textsci}} \dicPos{m}[1] \dicFlx{(‑a, ‑ar)}[8] \dicFieldCat{zool.} \dicDirectTranslationCS{sumec, sumec velký} \textit{(l.~{\textLA{Silurus glanis}})}  \dicsymPhoto\ 
\dicFigure{75323.jpg}{Fengrani}{Fengrani - Fiala Miroslav, Biolib, Copyright/CC-BY-NC}
\dicEntry[fengum] \dicTerm{fengum} \dicIPA{{f}{ei}{\ng}{\r{g}}{\textscy}{\textsubring{m}}} \dicPos{v} \dicFlx{ind pf pl 1 pers} \dicLink{fá}
\dicEntry[fengur] \dicTerm{feng|ur} \dicIPA{{f}{ei}{\ng}{\r{g}}{\textscy}{\textsubring{r}}} \dicPos{m}[9] \dicFlx{(‑s\,/\addthin ‑jar, ‑ir)}[26] \textbf{1.} \dicSynonym{veiði} \dicDirectTranslationCS{úlovek, kořist}  \textbf{2.} \dicSynonym{ávinningur} \dicDirectTranslationCS{prospěch, kořist} \dicExampleIS{fengur í ritinu} \dicExampleCS{prospěch z~díla}
\dicEntry[Fenrisúlfur] \dicTerm{Fenris··úlf|ur} \dicIPA{{f}{\textepsilon}{n}{r}{\textsci}{s}{u}{l}{v}{\textscy}{\textsubring{r}}} \dicPos{m}[6] \dicFlx{(‑s)}[3] \dicFlx{prop} \dicFieldCat{myt.} \dicDirectTranslationCS{Fenrir} \dicIndirectTranslationCS{(obrovský vlk, nejnebezpečnější nepřítel Ásů, Lokiho syn)}
\dicEntry[fer] \dicTerm{fer} \dicIPA{{f}{\textepsilon}{\textlengthmark}{\textsubring{r}}} \dicPos{v} \dicFlx{ind praes sg 1 pers} \dicLink{fara}
\dicEntry[ferð] \dicTerm{ferð} \dicsymFrequent\  \dicIPA{{f}{\textepsilon}{r}{\texttheta}} \dicPos{f}[7] \dicFlx{(‑ar, ‑ir)}[1] \textbf{1.} \dicSynonym{ferðalag} \dicDirectTranslationCS{cesta, jízda};  \dicPhraseIS{góða ferð!} \dicDirectTranslationCS{šťastnou cestu!};  \dicPhraseIS{hvert er ferðinni heitið} \dicDirectTranslationCS{kam směřuje cesta};  \dicPhraseIS{vera frjáls ferða sinna} \dicDirectTranslationCS{být volný};  \dicPhraseIS{vera á ferð og flugi} \dicDirectTranslationCS{mít hektické období};  \dicPhraseIS{hér er á ferðinni e‑að} \dicLangCat{přen.} \dicDirectTranslationCS{teď je na pořadu (co)};  \dicPhraseIS{það\,/\addthin nú er alvara á ferðum} \dicLangCat{přen.} \dicDirectTranslationCS{teď jde do tuhého};  \dicPhraseIS{taka sér ferð á hendur} \dicDirectTranslationCS{vydat se na cestu}  \textbf{2.} \dicDirectTranslationCS{výlet, výprava, túra}  \textbf{3.} \dicDirectTranslationCS{zájezd, turné}  \textbf{4.} \dicSynonym{hraði} \dicDirectTranslationCS{rychlost}
\dicEntry[ferðaáætlun] \dicTerm{ferða··á·ætl|un} \dicIPA{{f}{\textepsilon}{r}{ð}{a}{au}{a}{i}{h}{\textsubring{d}}{l}{\textscy}{\textsubring{n}}} \dicPos{f}[7] \dicFlx{(‑unar, ‑anir)}[8] \dicDirectTranslationCS{jízdní řád\,/\addthin plán}
\dicEntry[ferðabók] \dicTerm{ferða··|bók} \dicIPA{{f}{\textepsilon}{r}{ð}{a}{\textsubring{b}}{ou}{\r{g}}} \dicPos{f}[8] \dicFlx{(‑bókar, ‑bækur)}[5] \dicFieldCat{lit.} \dicDirectTranslationCS{cestopis, reportáž}
\dicEntry[ferðabúnaður] \dicTerm{ferða··bú·nað|ur} \dicIPA{{f}{\textepsilon}{r}{ð}{a}{\textsubring{b}}{u}{n}{a}{ð}{\textscy}{\textsubring{r}}} \dicPos{m}[10] \dicFlx{(‑ar)}[9] \dicDirectTranslationCS{cestovní vybavení}
\dicEntry[ferðafélag] \dicTerm{ferða··fé·|lag} \dicIPA{{f}{\textepsilon}{r}{ð}{a}{f}{j}{\textepsilon}{l}{a}{x}} \dicPos{n}[2] \dicFlx{(‑lags, ‑lög)}[8] \dicDirectTranslationCS{cestovní kancelář}
\dicEntry[ferðafélagi] \dicTerm{ferða··fé·lag|i} \dicIPA{{f}{\textepsilon}{r}{ð}{a}{f}{j}{\textepsilon}{l}{a}{i}{j}{\textsci}} \dicPos{m}[1] \dicFlx{(‑a, ‑ar)}[8] \dicDirectTranslationCS{společník\,/\addthin společnice na cestách}
\dicEntry[ferðafólk] \dicTerm{ferða··fólk} \dicIPA{{f}{\textepsilon}{r}{ð}{a}{f}{ou}{\textsubring{l}}{\r{g}}} \dicPos{n}[2] \dicFlx{(‑s)}[2] \dicDirectTranslationCS{cestující, turisté}
\dicEntry[ferðahandbók] \dicTerm{ferða··hand·|bók} \dicIPA{{f}{\textepsilon}{r}{ð}{a}{h}{a}{n}{\textsubring{d}}{\textsubring{b}}{ou}{\r{g}}} \dicPos{f}[8] \dicFlx{(‑bókar, ‑bækur)}[5] \dicDirectTranslationCS{cestovní průvodce (kniha), cestovní příručka}
\dicEntry[ferðakostnaður] \dicTerm{ferða··kost·nað|ur} \dicIPA{{f}{\textepsilon}{r}{ð}{a}{k\smash{\textsuperscript{h}}}{\textopeno}{s}{\textsubring{d}}{n}{a}{ð}{\textscy}{\textsubring{r}}} \dicPos{m}[10] \dicFlx{(‑ar)}[9] \dicDirectTranslationCS{cestovní náklady}
\dicEntry[ferðalag] \dicTerm{ferða··|lag} \dicsymFrequent\  \dicIPA{{f}{\textepsilon}{r}{ð}{a}{l}{a}{x}} \dicPos{n}[2] \dicFlx{(‑lags, ‑lög)}[8] \dicDirectTranslationCS{cesta, výlet, výprava, túra};  \dicPhraseIS{fara í ferðalag} \dicDirectTranslationCS{jet na cesty, cestovat};  \dicPhraseIS{vera á ferðalagi} \dicDirectTranslationCS{být na cestách, cestovat, putovat}
\dicEntry[ferðaleyfi] \dicTerm{ferða··leyfi} \dicIPA{{f}{\textepsilon}{r}{ð}{a}{l}{ei}{v}{\textsci}} \dicPos{n}[2] \dicFlx{(‑s, ‑)}[14] \dicDirectTranslationCS{cestovní povolení, vízum}
\dicEntry[ferðamaður] \dicTerm{ferða··|maður} \dicIPA{{f}{\textepsilon}{r}{ð}{a}{m}{a}{ð}{\textscy}{\textsubring{r}}} \dicPos{m}[13] \dicFlx{(‑manns, ‑menn)}[2] \dicDirectTranslationCS{cestovatel(ka), turista, turistka, výletník, výletnice}
\dicEntry[ferðamál] \dicTerm{ferða··mál} \dicIPA{{f}{\textepsilon}{r}{ð}{a}{m}{au}{\textsubring{l}}} \dicPos{n}[2] \dicFlx{pl}[1] \dicDirectTranslationCS{turismus}
\dicEntry[ferðasaga] \dicTerm{ferða··|saga} \dicIPA{{f}{\textepsilon}{r}{ð}{a}{s}{a}{\textbabygamma}{a}} \dicPos{f}[1] \dicFlx{(‑sögu, ‑sögur)}[14] \dicDirectTranslationCS{cestopis}
\dicEntry[ferðaskrifstofa] \dicTerm{ferða··skrif·stof|a} \dicIPA{{f}{\textepsilon}{r}{ð}{a}{s}{\r{g}}{r}{\textsci}{f}{s}{\textsubring{d}}{\textopeno}{v}{a}} \dicPos{f}[1] \dicFlx{(‑u, ‑ur)}[7] \dicDirectTranslationCS{cestovní kancelář}
\dicEntry[ferðast] \dicTerm{ferð|ast} \dicsymFrequent\  \dicIPA{{f}{\textepsilon}{r}{ð}{a}{s}{\textsubring{d}}} \dicPos{v}[1] \dicFlx{(‑aðist)}[91] \dicFlx{refl} \dicDirectTranslationCS{cestovat, putovat} \dicExampleIS{ferðast til útlanda} \dicExampleCS{cestovat do ciziny}
\dicEntry[ferðataska] \dicTerm{ferða··|taska} \dicIPA{{f}{\textepsilon}{r}{ð}{a}{t\smash{\textsuperscript{h}}}{a}{s}{\r{g}}{a}} \dicPos{f}[1] \dicFlx{(‑tösku, ‑töskur)}[8] \dicDirectTranslationCS{cestovní taška, kufr}
\dicEntry[ferðatékki] \dicTerm{ferða··tékk|i} \dicIPA{{f}{\textepsilon}{r}{ð}{a}{t\smash{\textsuperscript{h}}}{j}{\textepsilon}{h}{\r{\textObardotlessj}}{\textsci}} \dicPos{m}[1] \dicFlx{(‑a, ‑ar)}[1] \dicDirectTranslationCS{cestovní šek}
\dicEntry[ferðatrygging] \dicTerm{ferða··trygg·ing} \dicIPA{{f}{\textepsilon}{r}{ð}{a}{t\smash{\textsuperscript{h}}}{r}{\textsci}{\r{\textObardotlessj}}{i}{\ng}{\r{g}}} \dicPos{f}[4] \dicFlx{(‑ar, ‑ar)}[5] \dicDirectTranslationCS{cestovní pojištění}
\dicEntry[ferðatölva] \dicTerm{ferða··tölv|a} \dicIPA{{f}{\textepsilon}{r}{ð}{a}{t\smash{\textsuperscript{h}}}{\oe}{l}{v}{a}} \dicPos{f}[1] \dicFlx{(‑u, ‑ur)}[7] \dicFieldCat{poč.} \dicSynonym{fartölva} \dicDirectTranslationCS{notebook, přenosný počítač}
\dicEntry[ferðaþjónusta] \dicTerm{ferða··þjón·ust|a} \dicIPA{{f}{\textepsilon}{r}{ð}{a}{\texttheta}{j}{ou}{n}{\textscy}{s}{\textsubring{d}}{a}} \dicPos{f}[1] \dicFlx{(‑u)}[5] \dicDirectTranslationCS{cestovní ruch, turismus}
\dicEntry[ferfaldur] \dicTerm{fer··|faldur} \dicIPA{{f}{\textepsilon}{\textsubring{r}}{f}{a}{l}{\textsubring{d}}{\textscy}{\textsubring{r}}} \dicPos{adj}[2] \dicFlx{(f ‑föld)}[16] \textbf{1.} \dicDirectTranslationCS{čtyřnásobný}  \textbf{2.} \dicDirectTranslationCS{čtyřdílný}
\dicEntry[ferfættur] \dicTerm{fer··fættur} \dicIPA{{f}{\textepsilon}{\textsubring{r}}{f}{a}{i}{h}{\textsubring{d}}{\textscy}{\textsubring{r}}} \dicPos{adj}[1]\dicFlx{}[-13] \dicDirectTranslationCS{čtyřnohý} \dicExampleIS{ferfætt dýr} \dicExampleCS{čtyřnohé zvíře}
\dicEntry[fergin] \dicTerm{fergin} \dicIPA{{f}{\textepsilon}{r}{\r{\textObardotlessj}}{\textsci}{\textsubring{n}}} \dicPos{n}[2] \dicFlx{(‑s)}[2] \dicFieldCat{bot.} \dicDirectTranslationCS{přeslička poříční} \textit{(l.~{\textLA{Equisetum fluviatile}})}  \dicsymPhoto\ 
\dicFigure{16069.jpg}{Fergin}{Fergin - Kesl Michael, Biolib, Copyright/CC-BY-NC}
\dicEntry[ferhenda] \dicTerm{fer··hend|a} \dicIPA{{f}{\textepsilon}{\textlengthmark}{\textsubring{r}}{h}{\textepsilon}{n}{\textsubring{d}}{a}} \dicPos{f}[1] \dicFlx{(‑u, ‑ur)}[19] \dicFieldCat{lit.} \dicDirectTranslationCS{čtyřveršová strofa, čtyřverší}
\dicEntry[ferhyrndur] \textls[15]{\dicTerm{fer··hyrndur} \dicIPA{{f}{\textepsilon}{\textlengthmark}{\textsubring{r}}{h}{\textsci}{r}{n}{\textsubring{d}}{\textscy}{\textsubring{r}}} \dicPos{adj}[2]\dicFlx{}[-17]} \textbf{1.} \dicDirectTranslationCS{čtyřúhelníkový, čtver\-hran\-ný, čtyřhranný, čtyřboký} \dicExampleIS{ferhyrndur garður} \dicExampleCS{čtyřhranný park}  \textbf{2.} \dicDirectTranslationCS{čtyřrohý}
\dicEntry[ferhyrningur] \dicTerm{fer··hyrn·ing|ur} \dicIPA{{f}{\textepsilon}{\textlengthmark}{\textsubring{r}}{h}{\textsci}{r}{\textsubring{d}}{n}{i}{\ng}{\r{g}}{\textscy}{\textsubring{r}}} \dicPos{m}[6] \dicFlx{(‑s, ‑ar)}[8] \dicFieldCat{mat.} \dicDirectTranslationCS{čtyřúhelník}
\dicEntry[ferill] \dicTerm{fer|ill} \dicsymFrequent\  \dicIPA{{f}{\textepsilon}{\textlengthmark}{r}{\textsci}{\textsubring{d}}{\textsubring{l}}} \dicPos{m}[6] \dicFlx{(‑ils, ‑lar)}[35] \textbf{1.} \dicSynonym{braut\smash{\textsuperscript{1}}} \dicDirectTranslationCS{dráha, stezka}  \textbf{2.} \dicSynonym{æviferill} \dicDirectTranslationCS{profesní dráha, kariéra} \dicExampleIS{hefja nýjan feril} \dicExampleCS{začít novou kariéru};  \dicPhraseIS{vera á ferli} \dicDirectTranslationCS{být v~pohybu, být na nohou}  \textbf{3.} \dicFieldCat{mat.} \dicDirectTranslationCS{křivka};  \dicPhraseIS{kúptur ferill} \dicFieldCat{mat.} \dicDirectTranslationCS{konvexní křivka};  \dicPhraseIS{lokaður ferill} \dicFieldCat{mat.} \dicDirectTranslationCS{uzavřená křivka}
\dicEntry[ferja] \dicTerm{ferj|a\smash{\textsuperscript{1}}} \dicIPA{{f}{\textepsilon}{r}{j}{a}} \dicPos{f}[1] \dicFlx{(‑u, ‑ur)}[7] \dicDirectTranslationCS{trajekt, přívoz} \dicExampleIS{Ferjan siglir milli landanna.} \dicExampleCS{Trajekt se plaví mezi zeměmi.}
\dicEntry[ferja] \dicTerm{ferj|a\smash{\textsuperscript{2}}} \dicIPA{{f}{\textepsilon}{r}{j}{a}} \dicPos{v}[1] \dicFlx{(‑aði)}[1] \dicFlx{dat} \dicDirectTranslationCS{převést, přepravit (přes řeku ap.)}
\dicEntry[ferjumaður] \dicTerm{ferju··|maður} \dicIPA{{f}{\textepsilon}{r}{j}{\textscy}{m}{a}{ð}{\textscy}{\textsubring{r}}} \dicPos{m}[13] \dicFlx{(‑manns, ‑menn)}[2] \dicLangCat{zast.} \dicDirectTranslationCS{převozník, převoznice}
\dicEntry[ferkílómetri] \dicTerm{fer··kíló·metr|i} \dicIPA{{f}{\textepsilon}{\textsubring{r}}{c\smash{\textsuperscript{h}}}{i}{l}{ou}{m}{\textepsilon}{\textsubring{d}}{r}{\textsci}} \dicPos{m}[1] \dicFlx{(‑a, ‑ar)}[1] \dicDirectTranslationCS{čtvereční kilometr}
\dicEntry[ferlaufungur] \dicTerm{fer··lauf·ung|ur} \dicIPA{{f}{\textepsilon}{r}{l}{\oe i}{v}{u}{\ng}{\r{g}}{\textscy}{\textsubring{r}}} \dicPos{m}[6] \dicFlx{(‑s, ‑ar)}[8] \dicFieldCat{bot.} \dicDirectTranslationCS{vraní oko čtyřlisté} \textit{(l.~{\textLA{Paris quadrifolia}})}  \dicsymPhoto\ 
\dicFigure{ds_image_ferlaufungur_0_2.jpg}{Ferlaufungur}{Ferlaufungur - Rosser1954, PD}
\dicEntry[ferlegur] \dicTerm{fer··legur} \dicIPA{{f}{\textepsilon}{r}{l}{\textepsilon}{\textbabygamma}{\textscy}{\textsubring{r}}} \dicPos{adj}[1]\dicFlx{}[-8] \textbf{1.} \dicSynonym*{ógnvænlegur} \dicDirectTranslationCS{škaredý, monstrózní}  \textbf{2.} \dicLangCat{hovor.} \dicSynonym{slæmur} \dicDirectTranslationCS{strašný, příšerný} \dicExampleIS{vera ferlegur asni} \dicExampleCS{být příšerný hlupák}
\dicEntry[ferli] \dicTerm{ferli} \dicIPA{{f}{\textepsilon}{r}{\textsubring{d}}{l}{\textsci}} \dicPos{n}[2] \dicFlx{(‑s, ‑)}[14] \dicSynonym{framvinda} \dicDirectTranslationCS{průběh, chod, proces, procedura}
\dicEntry[ferlíki] \dicTerm{fer··líki} \dicIPA{{f}{\textepsilon}{r}{l}{i}{\r{\textObardotlessj}}{\textsci}} \dicPos{n}[2] \dicFlx{(‑s, ‑)}[16] \dicDirectTranslationCS{kolos, monstrum}
\dicEntry[ferm.] \dicTerm{ferm.} \dicPos{zkr} \dicPhraseIS{fermetri} \dicDirectTranslationCS{čtvereční metr}
\dicEntry[ferma] \dicTerm{ferm|a} \dicIPA{{f}{\textepsilon}{r}{m}{a}} \dicPos{v}[2] \dicFlx{(‑di, ‑t)}[141] \dicFlx{acc} \textbf{1.} \dicSynonym{hlaða\smash{\textsuperscript{2}}} \dicDirectTranslationCS{naložit, nakládat} \dicExampleIS{ferma skipið} \dicExampleCS{naložit loď} \dicAntonym{afferma}  \textbf{2.} \dicFieldCat{náb.} \dicDirectTranslationCS{konfirmovat, biřmovat} \dicExampleIS{Strákurinn neitaði að láta ferma sig.} \dicExampleCS{Hoch odmítal dát se konfirmovat.};  \dicIdiom{fermast}{ \dicPhraseIS{fermast}} \dicFlx{refl} \dicDirectTranslationCS{být biřmován}
\dicEntry[fermdur] \dicTerm{fermdur} \dicIPA{{f}{\textepsilon}{r}{m}{\textsubring{d}}{\textscy}{\textsubring{r}}} \dicPos{adj}[2]\dicFlx{}[-17] \dicDirectTranslationCS{konfirmovaný, biřmovaný}
\dicEntry[fermetri] \dicTerm{fer··metr|i} \dicIPA{{f}{\textepsilon}{r}{m}{\textepsilon}{\textsubring{d}}{r}{\textsci}} \dicPos{m}[1] \dicFlx{(‑a, ‑ar)}[1] \dicDirectTranslationCS{čtvereční metr}
\dicEntry[ferming] \dicTerm{ferm··ing} \dicIPA{{f}{\textepsilon}{r}{m}{i}{\ng}{\r{g}}} \dicPos{f}[4] \dicFlx{(‑ar, ‑ar)}[5] \textbf{1.} \dicDirectTranslationCS{nakládání (zboží ap.)}  \textbf{2.} \dicFieldCat{náb.} \dicDirectTranslationCS{konfirmace, biřmování}
\dicEntry[fermingarbarn] \dicTerm{fermingar··|barn} \dicIPA{{f}{\textepsilon}{r}{m}{i}{\ng}{\r{g}}{a}{r}{\textsubring{b}}{a}{r}{\textsubring{d}}{\textsubring{n}}} \dicPos{n}[2] \dicFlx{(‑barns, ‑börn)}[8] \dicFieldCat{náb.} \dicDirectTranslationCS{konfirmand(ka), biřmovanec, biřmovanka}
\dicEntry[fern] \dicTerm{fern} \dicIPA{{f}{\textepsilon}{r}{\textsubring{d}}{\textsubring{n}}} \dicPos{adj}[5]\dicFlx{}[-7] \dicDirectTranslationCS{čtverý} \dicExampleIS{fernar buxur} \dicExampleCS{čtvery kalhoty}
\dicEntry[ferna] \dicTerm{fern|a} \dicIPA{{f}{\textepsilon}{r}{\textsubring{d}}{n}{a}} \dicPos{f}[1] \dicFlx{(‑u, ‑ur)}[7] \dicDirectTranslationCS{krabice, kartón} \dicExampleIS{mjólkurferna} \dicExampleCS{krabice mléka}
\dicEntry[ferningur] \dicTerm{fern··ing|ur} \dicIPA{{f}{\textepsilon}{r}{\textsubring{d}}{n}{i}{\ng}{\r{g}}{\textscy}{\textsubring{r}}} \dicPos{m}[6] \dicFlx{(‑s, ‑ar)}[8] \dicFieldCat{mat.} \dicDirectTranslationCS{čtverec}
\dicEntry[ferskeytla] \dicTerm{fer··skeytl|a} \dicIPA{{f}{\textepsilon}{\textsubring{r}}{s}{\r{\textObardotlessj}}{ei}{h}{\textsubring{d}}{l}{a}} \dicPos{f}[1] \dicFlx{(‑u, ‑ur)}[19] \dicFieldCat{lit.} \dicDirectTranslationCS{kvartet} \dicIndirectTranslationCS{(čtyřveršová strofa)}
\dicEntry[ferskja] \dicTerm{ferskj|a} \dicIPA{{f}{\textepsilon}{\textsubring{r}}{s}{\r{\textObardotlessj}}{a}} \dicPos{f}[1] \dicFlx{(‑u, ‑ur)}[25] \dicDirectTranslationCS{broskev}
\dicEntry[ferskjutré] \dicTerm{ferskju··tré} \dicIPA{{f}{\textepsilon}{\textsubring{r}}{s}{\r{\textObardotlessj}}{\textscy}{t\smash{\textsuperscript{h}}}{r}{j}{\textepsilon}} \dicPos{n}[2] \dicFlx{(‑s, ‑)}[36] \dicFieldCat{bot.} \dicDirectTranslationCS{broskvoň, broskvoň obecná} \textit{(l.~{\textLA{Prunus persica}})}
\dicEntry[ferskur] \dicTerm{ferskur} \dicsymFrequent\  \dicIPA{{f}{\textepsilon}{\textsubring{r}}{s}{\r{g}}{\textscy}{\textsubring{r}}} \dicPos{adj}[1]\dicFlx{}[-1] \textbf{1.} \dicSynonym{nýr} \dicDirectTranslationCS{čerstvý, svěží} \dicExampleIS{ferskt loft} \dicExampleCS{čerstvý vzduch}  \textbf{2.} \dicSynonym{hress} \dicDirectTranslationCS{svěží, plný energie}
\dicEntry[ferskvatn] \dicTerm{fersk··vatn} \dicIPA{{f}{\textepsilon}{\textsubring{r}}{s}{\r{g}}{v}{a}{h}{\textsubring{d}}{\textsubring{n}}} \dicPos{n}[2] \dicFlx{(‑s)}[2] \dicDirectTranslationCS{sladká\,/\addthin říční voda}
\dicEntry[ferskvatnsfiskur] \dicTerm{fersk·vatns··fisk|ur} \dicIPA{{f}{\textepsilon}{\textsubring{r}}{s}{\r{g}}{v}{a}{s}{f}{\textsci}{s}{\r{g}}{\textscy}{\textsubring{r}}} \dicPos{m}[6] \dicFlx{(‑s, ‑ar)}[10] \dicDirectTranslationCS{sladkovodní ryba}
\dicEntry[ferstrendur] \dicTerm{fer··strendur} \dicIPA{{f}{\textepsilon}{\textsubring{r}}{s}{\textsubring{d}}{r}{\textepsilon}{n}{\textsubring{d}}{\textscy}{\textsubring{r}}} \dicPos{adj}[2]\dicFlx{}[-17] \dicDirectTranslationCS{čtyřstranný, čtyřhran\-ný, čtyřboký}
\dicEntry[fertugasti] \dicTerm{fer·tug··asti} \dicIPA{{f}{\textepsilon}{\textsubring{r}}{t\smash{\textsuperscript{h}}}{\textscy}{\textbabygamma}{a}{s}{\textsubring{d}}{\textsci}} \dicPos{num}[6] \dicFlx{ord} \dicDirectTranslationCS{čtyřicátý}
\dicEntry[fertugfaldur] \dicTerm{fer·tug··|faldur} \dicIPA{{f}{\textepsilon}{\textsubring{r}}{t\smash{\textsuperscript{h}}}{\textscy}{x}{f}{a}{l}{\textsubring{d}}{\textscy}{\textsubring{r}}} \dicPos{adj}[2] \dicFlx{(f ‑föld)}[16] \textbf{1.} \dicDirectTranslationCS{čtyřicetinásobný}  \textbf{2.} \dicDirectTranslationCS{čtyřicetidílný}
\dicEntry[fertugsaldur] \dicTerm{fer·tugs··ald|ur} \dicIPA{{f}{\textepsilon}{\textsubring{r}}{t\smash{\textsuperscript{h}}}{\textscy}{x}{s}{a}{l}{\textsubring{d}}{\textscy}{\textsubring{r}}} \dicPos{m}[5] \dicFlx{(‑urs, ‑rar)}[3] \dicDirectTranslationCS{věk po třicítce, věk mezi 30. a~39. rokem};  \dicPhraseIS{e‑n er á fertugsaldri} \dicDirectTranslationCS{(kdo) je třicátník}
\dicEntry[fertugur] \dicTerm{fer··tugur} \dicIPA{{f}{\textepsilon}{\textsubring{r}}{t\smash{\textsuperscript{h}}}{\textscy}{\textbabygamma}{\textscy}{\textsubring{r}}} \dicPos{adj}[1]\dicFlx{}[-6] \textbf{1.} \dicDirectTranslationCS{čtyřicetiletý}  \textbf{2.} \dicDirectTranslationCS{měřící čtyřicet sáhů}
\dicEntry[festa] \dicTerm{fest|a\smash{\textsuperscript{1}}} \dicIPA{{f}{\textepsilon}{s}{\textsubring{d}}{a}} \dicPos{f}[1] \dicFlx{(‑u)}[5] \textbf{1.} \dicSynonym{hald} \dicDirectTranslationCS{chyt, úchop}  \textbf{2.} \dicSynonym{stað\-festa\smash{\textsuperscript{1}}} \dicDirectTranslationCS{rozhodnost, pevnost, cílevědomost}
\dicEntry[festa] \dicTerm{fest|a\smash{\textsuperscript{2}}} \dicsymFrequent\  \dicIPA{{f}{\textepsilon}{s}{\textsubring{d}}{a}} \dicPos{v}[2] \dicFlx{(‑i, ‑)}[12] \dicFlx{acc} \dicSynonym*{setja fast} \dicDirectTranslationCS{připevnit, připevňovat, upevnit, upevňovat, upnout, upínat} \dicExampleIS{festa myndina á vegginn með nagla} \dicExampleCS{připevnit obraz na zeď hřebíkem};  \dicPhraseIS{festa blund\,/\addthin svefn} \dicDirectTranslationCS{usnout};  \dicPhraseIS{festa hugann í e‑ð} \dicDirectTranslationCS{upnout mysl na (co)};  \dicPhraseIS{festa kaup á e‑u} \dicDirectTranslationCS{koupit (co)};  \dicPhraseIS{festa e‑ð á filmu} \dicDirectTranslationCS{nafilmovat (co), zachytit (co) na film};  \dicIdiom{festast}{ \dicPhraseIS{festast}} \dicFlx{refl} \dicDirectTranslationCS{připevnit se, uchytit se}
\dicEntry[festi] \dicTerm{fest|i} \dicIPA{{f}{\textepsilon}{s}{\textsubring{d}}{\textsci}} \dicPos{f}[5] \dicFlx{(‑ar, ‑ar)}[3] \textbf{1.} \dicSynonym{keðja} \dicDirectTranslationCS{řetěz}  \textbf{2.} \dicSynonym{kaðall} \dicDirectTranslationCS{lano, provaz}  \textbf{3.} \dicSynonym{hálsfesti} \dicDirectTranslationCS{řetízek}  \textbf{4.} \dicSynonym{sjálfhelda} \dicDirectTranslationCS{slepá ulička} \dicExampleIS{Kindurnar komust í festi.} \dicExampleCS{Ovce se dostaly do slepé uličky.}
\dicEntry[festing] \dicTerm{fest··ing} \dicIPA{{f}{\textepsilon}{s}{\textsubring{d}}{i}{\ng}{\r{g}}} \dicPos{f}[4] \dicFlx{(‑ar, ‑ar)}[5] \textbf{1.} \dicDirectTranslationCS{upevnění, připevnění}  \textbf{2.} \dicLangCat{kniž.} \dicSynonym*{himinhvel} \dicDirectTranslationCS{nebeská klenba, obloha, nebe}
\dicEntry[festulaus] \dicTerm{festu··laus} \dicIPA{{f}{\textepsilon}{s}{\textsubring{d}}{\textscy}{l}{\oe i}{s}} \dicPos{adj}[5]\dicFlx{}[-1] \dicDirectTranslationCS{vrtkavý, nestálý}
\dicEntry[fet] \dicTerm{fet} \dicsymFrequent\  \dicIPA{{f}{\textepsilon}{\textlengthmark}{\textsubring{d}}} \dicPos{n}[2] \dicFlx{(‑s, ‑)}[5] \dicDirectTranslationCS{stopa} \dicIndirectTranslationCS{(stará délková míra)} \dicExampleIS{Fet var notað sem mælieining á Íslandi.} \dicExampleCS{Stopa se používala na Islandě jako měrná jednotka.};  \dicPhraseIS{fet fyrir fet} \dicFlx{adv} \dicDirectTranslationCS{krok za krokem}
\dicEntry[feta] \dicTerm{fet|a} \dicIPA{{f}{\textepsilon}{\textlengthmark}{\textsubring{d}}{a}} \dicPos{v}[1] \dicFlx{(‑aði)}[44] \dicDirectTranslationCS{kráčet, jít, šlapat};  \dicPhraseIS{feta sig áfram} \dicDirectTranslationCS{jít postupně kupředu};  \dicPhraseIS{feta í fótspor e‑rs} \dicLangCat{přen.} \dicDirectTranslationCS{jít v~(čích) stopách\,/\addthin šlépějích} \dicExampleIS{Börn sem feta í fótspor foreldranna.} \dicExampleCS{Děti, které kráčejí ve stopách rodičů.}
\dicEntry[fetill] \dicTerm{fet|ill} \dicIPA{{f}{\textepsilon}{\textlengthmark}{\textsubring{d}}{\textsci}{\textsubring{d}}{\textsubring{l}}} \dicPos{m}[6] \dicFlx{(‑ils, ‑lar)}[35] \dicSynonym{fatli} \dicDirectTranslationCS{zdravotnický\,/\addthin trojcípý šátek}
\dicEntry[fetta] \dicTerm{fett|a} \dicIPA{{f}{\textepsilon}{h}{\textsubring{d}}{a}} \dicPos{v}[2] \dicFlx{(‑i, ‑)}[1] \dicFlx{acc} \dicPhraseIS{fetta sig} \dicSynonym*{reigja sig} \dicDirectTranslationCS{protáhnout se, protahovat se};  \dicPhraseIS{fetta sig og bretta} \dicDirectTranslationCS{šklebit se, dělat obličeje};  \dicPhraseIS{fetta fingur út í e‑ð} \dicLangCat{přen.} \dicDirectTranslationCS{kritizovat (co), ukazovat prstem na (co)}
\dicEntry[feykja] \dicTerm{feyk|ja} \dicIPA{{f}{ei}{\textlengthmark}{\r{\textObardotlessj}}{a}} \dicPos{v}[2] \dicFlx{(‑ti, ‑t)}[25] \dicFlx{dat} \dicDirectTranslationCS{odfouknout, odvát} \dicExampleIS{feykja laufinu burt} \dicExampleCS{odvát lístek pryč}
\dicEntry[fé] \dicTerm{fé} \dicsymFrequent\  \dicIPA{{f}{j}{\textepsilon}{\textlengthmark}} \dicPos{n}[3] \dicFlx{(fjár)}[5] \textbf{1.} \dicSynonym{sauðfé} \dicDirectTranslationCS{ovce}  \textbf{2.} \dicSynonym*{peningar} \dicDirectTranslationCS{peníze};  \dicPhraseIS{bera fé á e‑n} \dicDirectTranslationCS{podplatit (koho), podplácet (koho)};  \dicPhraseIS{leggja fé í e‑ð} \dicDirectTranslationCS{vložit do (čeho) peníze};  \dicPhraseIS{farið hefur fé betra} \dicLangCat{přen.} \dicDirectTranslationCS{mohlo to být ještě horší};  \dicPhraseIS{þetta er bitamunur en ekki fjár} \dicDirectTranslationCS{to je zanedbatelný rozdíl}
\dicEntry[fébætur] \dicTerm{fé··bætur} \dicIPA{{f}{j}{\textepsilon}{\textlengthmark}{\textsubring{b}}{a}{i}{\textsubring{d}}{\textscy}{\textsubring{r}}} \dicPos{f}[8] \dicFlx{pl}[6] \dicDirectTranslationCS{náhrada, odškodnění}
\dicEntry[féfletta] \dicTerm{fé··flett|a} \dicIPA{{f}{j}{\textepsilon}{\textlengthmark}{f}{l}{\textepsilon}{h}{\textsubring{d}}{a}} \dicPos{v}[2] \dicFlx{(‑i, ‑)}[6] \dicFlx{acc} \dicDirectTranslationCS{ošidit, obrat, oškubat (člověka ap.)}
\dicEntry[fégirnd] \dicTerm{fé··girnd} \dicIPA{{f}{j}{\textepsilon}{\textlengthmark}{\r{\textObardotlessj}}{\textsci}{r}{n}{\textsubring{d}}} \dicPos{f}[7] \dicFlx{(‑ar)}[3] \dicDirectTranslationCS{lakota, chamtivost, touha po penězích}
\dicEntry[fékk] \dicTerm{fékk} \dicIPA{{f}{j}{\textepsilon}{h}{\r{g}}} \dicPos{v} \dicFlx{ind pf sg 1 pers} \dicLink{fá}
\dicEntry[fél.] \dicTerm{fél.} \dicPos{zkr} \dicPhraseIS{félag} \dicDirectTranslationCS{společnost}
\dicEntry[félag] \dicTerm{fé··|lag} \dicsymFrequent\  \dicIPA{{f}{j}{\textepsilon}{\textlengthmark}{l}{a}{x}} \dicPos{n}[2] \dicFlx{(‑lags, ‑lög)}[8] \textbf{1.} \dicDirectTranslationCS{organizace, společnost, spolek, klub} \dicExampleIS{íþróttafélag} \dicExampleCS{sportovní klub}  \textbf{2.} \dicSynonym*{samfélag fólks} \dicDirectTranslationCS{společnost, přítomnost (lidí ap.)} \dicExampleIS{í félagi við systur sína} \dicExampleCS{ve společnosti své sestry}
\dicEntry[félagasamtök] \dicTerm{fé·laga··sam·tök} \textls[-10]{\dicIPA{{f}{j}{\textepsilon}{\textlengthmark}{l}{a}{\textbabygamma}{a}{s}{a}{m}{t\smash{\textsuperscript{h}}}{\oe}{\r{g}}}} \dicPos{n}[2] \dicFlx{pl}[9] \dicDirectTranslationCS{asociace, svaz, unie}
\dicEntry[félagi] \dicTerm{fé··lag|i} \dicsymFrequent\  \dicIPA{{f}{j}{\textepsilon}{\textlengthmark}{l}{a}{i}{j}{\textsci}} \dicPos{m}[1] \dicFlx{(‑a, ‑ar)}[8] \textbf{1.} \dicSynonym{samherji} \dicDirectTranslationCS{společník, společnice, partner(ka), druh, družka} \dicExampleIS{deila mikilli gleði með félögum sínum} \dicExampleCS{sdílet velkou radost se svými společníky}  \textbf{2.} \dicSynonym*{meðlimur í félagi} \dicDirectTranslationCS{člen(ka)}
\dicEntry[félagsdómur] \dicTerm{fé·lags··dóm|ur} \dicIPA{{f}{j}{\textepsilon}{\textlengthmark}{l}{a}{x}{s}{\textsubring{d}}{ou}{m}{\textscy}{\textsubring{r}}} \dicPos{m}[6] \dicFlx{(‑s)}[12] \dicFieldCat{práv.} \dicDirectTranslationCS{pracovněprávní soud}
\dicEntry[félagseign] \dicTerm{fé·lags··eign} \dicIPA{{f}{j}{\textepsilon}{\textlengthmark}{l}{a}{x}{s}{ei}{\r{g}}{\textsubring{n}}} \dicPos{f}[7] \dicFlx{(‑ar, ‑ir)}[1] \dicDirectTranslationCS{společné vlastnictví}
\dicEntry[félagsfræði] \dicTerm{fé·lags··fræð|i} \dicIPA{{f}{j}{\textepsilon}{\textlengthmark}{l}{a}{x}{s}{f}{r}{a}{i}{ð}{\textsci}} \dicPos{f}[3] \dicFlx{(‑i)}[3] \dicDirectTranslationCS{sociologie}
\dicEntry[félagsfræðingur] \dicTerm{fé·lags·fræð··ing|ur} \dicIPA{{f}{j}{\textepsilon}{\textlengthmark}{l}{a}{x}{s}{f}{r}{a}{i}{ð}{i}{\ng}{\r{g}}{\textscy}{\textsubring{r}}} \dicPos{m}[6] \dicFlx{(‑s, ‑ar)}[8] \dicDirectTranslationCS{sociolog, socioložka}
\dicEntry[félagsgjald] \dicTerm{fé·lags··|gjald} \dicIPA{{f}{j}{\textepsilon}{\textlengthmark}{l}{a}{x}{s}{\r{\textObardotlessj}}{a}{l}{\textsubring{d}}} \dicPos{n}[2] \dicFlx{(‑gjalds, ‑gjöld)}[8] \dicDirectTranslationCS{členský poplatek\,/\addthin příspěvek}
\dicEntry[félagsheimili] \dicTerm{fé·lags··heimili} \dicIPA{{f}{j}{\textepsilon}{\textlengthmark}{l}{a}{x}{s}{h}{ei}{m}{\textsci}{l}{\textsci}} \dicPos{n}[2] \dicFlx{(‑s, ‑)}[14] \dicDirectTranslationCS{komunitní\,/\addthin kulturní centrum, společenská místnost}
\dicEntry[félagshyggja] \dicTerm{fé·lags··hyggj|a} \dicIPA{{f}{j}{\textepsilon}{\textlengthmark}{l}{a}{x}{s}{h}{\textsci}{\r{\textObardotlessj}}{a}} \dicPos{f}[1] \dicFlx{(‑u)}[5] \textbf{1.} \dicFieldCat{pol.} \dicDirectTranslationCS{socialismus}  \textbf{2.} \dicDirectTranslationCS{sociální cítění}
\dicEntry[félagshyggjumaður] \dicTerm{fé·lags·hyggju··|maður} \dicIPA{{f}\-{j}\-{\textepsilon}\-{\textlengthmark}\-{l}\-{a}\-{x}\-{s}\-{h}\-{\textsci}\-{\r{\textObardotlessj}}\-{\textscy}\-{m}\-{a}\-{ð}\-{\textscy}\-{\textsubring{r}}\-} \dicPos{m}[13] \dicFlx{(‑manns, ‑menn)}[2] \dicDirectTranslationCS{sociálně smýšlející člověk, socialista, socialistka}
\dicEntry[félagslegur] \dicTerm{fé·lags··legur} \dicsymFrequent\  \dicIPA{{f}{j}{\textepsilon}{\textlengthmark}{l}{a}{x}{s}{l}{\textepsilon}{\textbabygamma}{\textscy}{\textsubring{r}}} \dicPos{adj}[1]\dicFlx{}[-8] \dicDirectTranslationCS{sociální, společenský};  \dicPhraseIS{félagsleg réttindi} \dicDirectTranslationCS{sociální spravedlnost};  \dicPhraseIS{félagsleg vandamál} \dicDirectTranslationCS{sociální problémy}
\dicEntry[félagslyndur] \dicTerm{fé·lags··lyndur} \dicIPA{{f}{j}{\textepsilon}{\textlengthmark}{l}{a}{x}{s}{l}{\textsci}{n}{\textsubring{d}}{\textscy}{\textsubring{r}}} \dicPos{adj}[2]\dicFlx{}[-14] \dicDirectTranslationCS{družný, společenský, pospolitý}
\dicEntry[félagsmaður] \dicTerm{fé·lags··|maður} \dicIPA{{f}{j}{\textepsilon}{\textlengthmark}{l}{a}{x}{s}{m}{a}{ð}{\textscy}{\textsubring{r}}} \dicPos{m}[13] \dicFlx{(‑manns, ‑menn)}[2] \dicDirectTranslationCS{člen(ka) (určité společnosti ap.)}
\dicEntry[félagsmál] \dicTerm{fé·lags··mál} \dicIPA{{f}{j}{\textepsilon}{\textlengthmark}{l}{a}{x}{s}{m}{au}{\textsubring{l}}} \dicPos{n}[2] \dicFlx{pl}[1] \dicDirectTranslationCS{sociální otázky\,/\addthin politika}
\dicEntry[félagsmálanefnd] \dicTerm{fé·lags·mála··nefnd} \dicIPA{{f}{j}{\textepsilon}{\textlengthmark}{l}{a}{x}{s}{m}{au}{l}{a}{n}{\textepsilon}{m}{\textsubring{d}}} \dicPos{f}[7] \dicFlx{(‑ar, ‑ir)}[1] \dicFieldCat{pol.} \dicIndirectTranslationCS{parlamentní komise zabývající se sociálními otázkami}
\dicEntry[félagsmálaráðherra] \dicTerm{fé·lags·mála··ráð·herr|a} \dicIPA{{f}\-{j}\-{\textepsilon}\-{\textlengthmark}\-{l}\-{a}\-{x}\-{s}\-{m}\-{au}\-{l}\-{a}\-{r}\-{au}\-{\texttheta}\-{h}\-{\textepsilon}\-{r}\-{a}\-} \dicPos{m}[1] \dicFlx{(‑a, ‑ar)}[17] \dicDirectTranslationCS{ministr(yně) sociálních věcí}
\dicEntry[félagsmálaráðuneyti] \dicTerm{fé·lags·mála··ráðu·neyti} \dicIPA{{f}\-{j}\-{\textepsilon}\-{\textlengthmark}\-{l}\-{a}\-{x}\-{s}\-{m}\-{au}\-{l}\-{a}\-{r}\-{au}\-{ð}\-{\textscy}\-{n}\-{ei}\-{\textsubring{d}}\-{\textsci}\-} \dicPos{n}[2] \dicFlx{(‑s, ‑)}[14] \dicDirectTranslationCS{ministerstvo sociálních věcí}
\dicEntry[félagsráðgjafi] \dicTerm{fé·lags··ráð·gjaf|i} \dicIPA{{f}{j}{\textepsilon}{\textlengthmark}{l}{a}{x}{s}{r}{au}{ð}{\r{\textObardotlessj}}{a}{v}{\textsci}} \dicPos{m}[1] \dicFlx{(‑a, ‑ar)}[8] \dicDirectTranslationCS{sociální pracovník\,/\addthin pracovnice}
\dicEntry[félagsráðgjöf] \dicTerm{fé·lags··ráð·|gjöf} \dicIPA{{f}{j}{\textepsilon}{\textlengthmark}{l}{a}{x}{s}{r}{au}{ð}{\r{\textObardotlessj}}{\oe}{f}} \dicPos{f}[7] \dicFlx{(‑gjafar)}[19] \dicDirectTranslationCS{sociální poradenství}
\dicEntry[félagsskapur] \dicTerm{fé·lags··skap|ur} \dicsymFrequent\  \dicIPA{{f}{j}{\textepsilon}{\textlengthmark}{l}{a}{x}{s}{\r{g}}{a}{\textsubring{b}}{\textscy}{\textsubring{r}}} \dicPos{m}[10] \dicFlx{(‑ar)}[15] \textbf{1.} \dicSynonym{félag} \dicDirectTranslationCS{sdružení, spolek} \dicExampleIS{félagsskapur fyrir einhleypa} \dicExampleCS{spolek pro nezadané}  \textbf{2.} \dicDirectTranslationCS{společenství, partnerství, družba}
\dicEntry[félagsvísindi] \dicTerm{fé·lags··vís·indi}\addthinS\textls[-13]{\dicIPA{{f}{j}{\textepsilon}{\textlengthmark}{l}{a}{x}{s}{v}{i}{s}{\textsci}{n}{\textsubring{d}}{\textsci}}} \dicPos{n}[2] \dicFlx{pl}[19] \dicDirectTranslationCS{společenské vědy}
\dicEntry[félaus] \dicTerm{fé··laus} \dicIPA{{f}{j}{\textepsilon}{\textlengthmark}{l}{\oe i}{s}} \dicPos{adj}[5]\dicFlx{}[-1] \dicSynonym{peningalaus} \dicDirectTranslationCS{(jsoucí) bez peněz, (jsoucí) švorc, švorcový}
\dicEntry[féleysi] \dicTerm{fé··leysi} \dicIPA{{f}{j}{\textepsilon}{\textlengthmark}{l}{ei}{s}{\textsci}} \dicPos{n}[2] \dicFlx{(‑s)}[20] \dicDirectTranslationCS{nedostatek peněz}
\dicEntry[félítill] \dicTerm{fé··|lítill} \dicIPA{{f}{j}{\textepsilon}{\textlengthmark}{l}{i}{\textsubring{d}}{\textsci}{\textsubring{d}}{\textsubring{l}}} \dicPos{adj}[11] \dicFlx{(comp ‑minni, sup ‑minn\-st\-ur)}[10] \dicDirectTranslationCS{nemajetný, chudý, mající málo peněz}
\dicEntry[féll] \dicTerm{féll} \dicIPA{{f}{j}{\textepsilon}{\textsubring{d}}{\textsubring{l}}} \dicPos{v} \dicFlx{ind pf sg 1 pers} \dicLink{falla}
\dicEntry[félli] \dicTerm{félli} \dicIPA{{f}{j}{\textepsilon}{\textsubring{d}}{l}{\textsci}} \dicPos{v} \dicFlx{con pf sg 1 pers} \dicLink{falla}
\dicEntry[féllum] \dicTerm{féllum} \dicIPA{{f}{j}{\textepsilon}{\textsubring{d}}{l}{\textscy}{\textsubring{m}}} \dicPos{v} \dicFlx{ind pf pl 1 pers} \dicLink{falla}
\dicEntry[félög] \dicTerm{fé··lög} \dicIPA{{f}{j}{\textepsilon}{\textlengthmark}{l}{\oe}{x}} \dicPos{n} \dicFlx{pl nom} \dicLink{félag}
\dicEntry[fénaður] \dicTerm{fé··nað|ur} \dicIPA{{f}{j}{\textepsilon}{\textlengthmark}{n}{a}{ð}{\textscy}{\textsubring{r}}} \dicPos{m}[10] \dicFlx{(‑ar)}[7] \textbf{1.} \dicSynonym{búfé} \dicDirectTranslationCS{hospodářská zvířata (zejména ovce)}  \textbf{2.} \dicLangCat{han.} \dicDirectTranslationCS{ovce} \dicIndirectTranslationCS{(o~lidech)}
\dicEntry[fés] \dicTerm{fés} \dicIPA{{f}{j}{\textepsilon}{\textlengthmark}{s}} \dicPos{n}[2] \dicFlx{(‑s, ‑)}[5] \dicLangCat{han.} \dicDirectTranslationCS{ksicht, držka}
\dicEntry[fésbók] \dicTerm{fés··bók} \dicIPA{{f}{j}{\textepsilon}{\textlengthmark}{s}{\textsubring{b}}{ou}{\r{g}}} \dicPos{f}[8] \dicFlx{(‑ar)}[7] \dicFieldCat{poč.} \dicDirectTranslationCS{Facebook}
\dicEntry[fésekt] \dicTerm{fé··sekt} \dicIPA{{f}{j}{\textepsilon}{\textlengthmark}{s}{\textepsilon}{x}{\textsubring{d}}} \dicPos{f}[7] \dicFlx{(‑ar, ‑ir)}[1] \dicDirectTranslationCS{peněžitá pokuta}
\dicEntry[féþúfa] \dicTerm{fé··þúf|a} \dicIPA{{f}{j}{\textepsilon}{\textlengthmark}{\texttheta}{u}{a}} \dicPos{f}[1] \dicFlx{(‑u)}[5] \dicPhraseIS{hafa e‑n að féþúfu} \dicLangCat{přen.} \dicDirectTranslationCS{ždímat (koho) (finančně ap.), mít (koho) jako dojnou krávu}
\dicEntry[fh.] \dicTerm{fh.} \dicPos{zkr} \dicPhraseIS{framsöguháttur} \dicFieldCat{jaz.} \dicDirectTranslationCS{oznamovací způsob}
\dicEntry[fid.] \dicTerm{fid.} \dicPos{zkr} \dicPhraseIS{fimmtudagur} \dicDirectTranslationCS{čtvrtek}
\dicEntry[fiðla] \dicTerm{fiðl|a} \dicIPA{{f}{\textsci}{ð}{l}{a}} \dicPos{f}[1] \dicFlx{(‑u, ‑ur)}[7] \dicFieldCat{hud.} \dicDirectTranslationCS{housle}
\dicEntry[fiðluleikari] \dicTerm{fiðlu··leik·ar|i} \dicIPA{{f}{\textsci}{ð}{l}{\textscy}{l}{ei}{\r{g}}{a}{r}{\textsci}} \dicPos{m}[1] \dicFlx{(‑a, ‑ar)}[13] \dicDirectTranslationCS{houslista, houslistka}
\dicEntry[fiðraður] \dicTerm{fiðr··|aður} \dicIPA{{f}{\textsci}{ð}{r}{a}{ð}{\textscy}{\textsubring{r}}} \dicPos{adj}[3] \dicFlx{(f ‑uð)}[3] \dicDirectTranslationCS{opeřený}
\dicEntry[fiðrast] \dicTerm{fiðr|ast} \dicIPA{{f}{\textsci}{ð}{r}{a}{s}{\textsubring{d}}} \dicPos{v}[1] \dicFlx{(‑aðist)}[95] \dicFlx{refl} \dicDirectTranslationCS{opeřit se}
\dicEntry[fiðrildi] \dicTerm{fiðrildi} \dicIPA{{f}{\textsci}{ð}{r}{\textsci}{l}{\textsubring{d}}{\textsci}} \dicPos{n}[2] \dicFlx{(‑s, ‑)}[14] \dicFieldCat{zool.} \dicDirectTranslationCS{motýl} \textit{(l.~{\textLA{Lepidoptera}})}  \dicsymPhoto\ 
\dicFigure{77.jpg}{Fiðrildi}{Fiðrildi - Zicha Ondřej, Biolib, Copyright/CC-BY-NC}
\dicEntry[fiðrildislirfa] \dicTerm{fiðrildis··lirf|a} \dicIPA{{f}{\textsci}{ð}{r}{\textsci}{l}{\textsubring{d}}{\textsci}{s}{l}{\textsci}{r}{v}{a}} \dicPos{f}[1] \dicFlx{(‑u, ‑ur)}[19] \dicSynonym{lirfa} \dicDirectTranslationCS{housenka}
\dicEntry[fiðringur] \dicTerm{fiðr··ing|ur} \dicIPA{{f}{\textsci}{ð}{r}{i}{\ng}{\r{g}}{\textscy}{\textsubring{r}}} \dicPos{m}[6] \dicFlx{(‑s)}[9] \dicDirectTranslationCS{napětí, chvění, nervozita};  \dicPhraseIS{grái fiðringurinn} \dicDirectTranslationCS{starý kanec} \dicIndirectTranslationCS{(postarší pán jevící zájem o~mladé dívky)}
\dicEntry[fiður] \dicTerm{fiður} \dicIPA{{f}{\textsci}{\textlengthmark}{ð}{\textscy}{\textsubring{r}}} \dicPos{n}[2] \dicFlx{(‑s)}[28] \dicDirectTranslationCS{(prachové) peří}
\dicEntry[fiff] \dicTerm{fiff} \dicIPA{{f}{\textsci}{f}{\textlengthmark}} \dicPos{n}[2] \dicFlx{(‑s)}[2] \dicLangCat{hovor.} \dicDirectTranslationCS{trik, obratné řešení}
\dicEntry[fikra] \dicTerm{fikr|a} \dicIPA{{f}{\textsci}{\textlengthmark}{\r{g}}{r}{a}} \dicPos{v}[1] \dicFlx{(‑aði)}[22] \dicPhraseIS{fikra sig áfram} \dicDirectTranslationCS{šinout se kupředu}
\dicEntry[fikt] \dicTerm{fikt} \dicIPA{{f}{\textsci}{x}{\textsubring{d}}} \dicPos{n}[2] \dicFlx{(‑s)}[2] \dicSynonym*{rælni} \dicDirectTranslationCS{hraní si, manipulace} \dicExampleIS{fikt við eldspýtur} \dicExampleCS{hraní si se sirkami}
\dicEntry[fikta] \dicTerm{fikt|a} \dicIPA{{f}{\textsci}{x}{\textsubring{d}}{a}} \dicPos{v}[1] \dicFlx{(‑aði)}[44] \textbf{1.} \dicDirectTranslationCS{hrát si, pohrávat si, zahrávat si};  \dicPhraseIS{fikta við e‑ð} \dicDirectTranslationCS{pohrávat si s~(čím)} \dicExampleIS{fikta við tækið} \dicExampleCS{pohrávat si s~přístrojem}  \textbf{2.} \dicDirectTranslationCS{experimentovat} \dicExampleIS{fikta við vímuefni} \dicExampleCS{experimentovat s~drogami}
\dicEntry[Filippseyingur] \dicTerm{Filipps·ey··ing|ur} \dicIPA{{f}{\textsci}{\textlengthmark}{l}{\textsci}{\textsubring{b}}{s}{ei}{\textlengthmark}{\ng}{\r{g}}{\textscy}{\textsubring{r}}} \dicPos{m}[6] \dicFlx{(‑s, ‑ar)}[8] \dicDirectTranslationCS{Filipínec, Filipínka}
\dicEntry[Filippseyjar] \dicTerm{Filipps··eyjar} \dicIPA{{f}{\textsci}{\textlengthmark}{l}{\textsci}{\textsubring{b}}{s}{ei}{j}{a}{\textsubring{r}}} \dicPos{f}[4] \dicFlx{pl}[15] \dicFieldCat{geog.} \dicDirectTranslationCS{Filipíny}
\dicEntry[filippseyskur] \dicTerm{filipps··ey·skur} \dicIPA{{f}{\textsci}{\textlengthmark}{l}{\textsci}{\textsubring{b}}{s}{ei}{s}{\r{g}}{\textscy}{\textsubring{r}}} \dicPos{adj}[1]\dicFlx{}[-6] \dicDirectTranslationCS{filipínský}
\dicEntry[filma] \dicTerm{film|a\smash{\textsuperscript{1}}} \dicIPA{{f}{\textsci}{l}{m}{a}} \dicPos{f}[1] \dicFlx{(‑u, ‑ur)}[19] \textbf{1.} \dicDirectTranslationCS{film (materiál na záznam fotografií, videa ap.)}  \textbf{2.} \dicSynonym{kvikmynd} \dicDirectTranslationCS{film (dokumentární ap.)}  \textbf{3.} \dicDirectTranslationCS{povlak, tenká blána, film}
\dicEntry[filma] \dicTerm{film|a\smash{\textsuperscript{2}}} \dicIPA{{f}{\textsci}{l}{m}{a}} \dicPos{v}[1] \dicFlx{(‑aði)}[1] \dicFlx{acc} \dicLangCat{hovor.} \dicSynonym*{taka kvikmynd} \dicDirectTranslationCS{filmovat, točit\,/\addthin natáčet film}
\dicEntry[fimbulfamb] \dicTerm{fimbul··famb} \dicIPA{{f}{\textsci}{m}{\textsubring{b}}{\textscy}{l}{f}{a}{m}{\textsubring{b}}} \dicPos{n}[2] \dicFlx{(‑s)}[2] \dicDirectTranslationCS{nesmysl, žvásty}
\dicEntry[fimbulvetur] \dicTerm{fimbul··vet|ur} \dicIPA{{f}{\textsci}{m}{\textsubring{b}}{\textscy}{l}{v}{\textepsilon}{\textsubring{d}}{\textscy}{\textsubring{r}}} \dicPos{m}[13] \dicFlx{(‑rar, ‑ur)}[6] \dicDirectTranslationCS{tuhá zima}
\dicEntry[fimi] \dicTerm{fim|i} \dicIPA{{f}{\textsci}{\textlengthmark}{m}{\textsci}} \dicPos{f}[3] \dicFlx{(‑i)}[3] \dicDirectTranslationCS{obratnost, zručnost}
\dicEntry[fimleikur] \dicTerm{fim··leik|ur} \dicIPA{{f}{\textsci}{m}{l}{ei}{\r{g}}{\textscy}{\textsubring{r}}} \dicPos{m}[6] \dicFlx{(‑s, ‑ar)}[15] \textbf{1.} \dicSynonym{lipurð} \dicDirectTranslationCS{obratnost, zručnost}  \textbf{2.} \dicPhraseIS{fimleikar} \dicFlx{pl} \dicFieldCat{sport.} \dicDirectTranslationCS{gymnastika}
\dicEntry[fimm] \dicTerm{fimm} \dicsymFrequent\  \dicIPA{{f}{\textsci}{m}{\textlengthmark}} \dicPos{num} \dicFlx{indecl} \dicDirectTranslationCS{pět} \dicExampleIS{borga fimm þúsund fyrir bók} \dicExampleCS{zaplatit pět tisíc za knížku}
\dicEntry[fimma] \dicTerm{fimm|a} \dicIPA{{f}{\textsci}{m}{\textlengthmark}{a}} \dicPos{f}[1] \dicFlx{(‑u, ‑ur)}[7] \textbf{1.} \dicDirectTranslationCS{pětka} \dicIndirectTranslationCS{(v~kartách)} \dicExampleIS{tígulfimma} \dicExampleCS{kárová pětka}  \textbf{2.} \dicDirectTranslationCS{pětka (autobus číslo 5 ap.)}
\dicEntry[fimmfaldur] \dicTerm{fimm··|faldur} \dicIPA{{f}{\textsci}{m}{f}{a}{l}{\textsubring{d}}{\textscy}{\textsubring{r}}} \dicPos{adj}[2] \dicFlx{(f ‑föld)}[16] \textbf{1.} \dicDirectTranslationCS{pětinásobný}  \textbf{2.} \dicDirectTranslationCS{pětidílný}
\dicEntry[fimmhundruðkall] \dicTerm{fimm·hundruð··kall} \dicIPA{{f}{\textsci}{m}{h}{\textscy}{n}{\textsubring{d}}{r}{\textscy}{ð}{k\smash{\textsuperscript{h}}}{a}{\textsubring{d}}{\textsubring{l}}} \dicPos{m}[4] \dicFlx{(‑s, ‑ar)}[4] \textbf{1.} \dicLangCat{hovor.} \dicDirectTranslationCS{pět set (korun)} \dicIndirectTranslationCS{(částka)}  \textbf{2.} \dicLangCat{hovor.} \dicDirectTranslationCS{pětistovka} \dicIndirectTranslationCS{(bankovka)}
\dicEntry[fimmhyrndur] \dicTerm{fimm··hyrndur} \dicIPA{{f}{\textsci}{m}{h}{\textsci}{r}{n}{\textsubring{d}}{\textscy}{\textsubring{r}}} \dicPos{adj}[2]\dicFlx{}[-17] \textbf{1.} \dicDirectTranslationCS{pětiúhelníkový, pětihranný}  \textbf{2.} \dicDirectTranslationCS{pětirohý}
\dicEntry[fimmhyrningur] \dicTerm{fimm··hyrn·ing|ur} \dicIPA{{f}{\textsci}{m}{h}{\textsci}{r}{\textsubring{d}}{n}{i}{\ng}{\r{g}}{\textscy}{\textsubring{r}}} \dicPos{m}[6] \dicFlx{(‑s, ‑ar)}[8] \dicFieldCat{mat.} \dicDirectTranslationCS{pětiúhelník}
\dicEntry[fimmkall] \dicTerm{fimm··kall} \dicIPA{{f}{\textsci}{m}{k\smash{\textsuperscript{h}}}{a}{\textsubring{d}}{\textsubring{l}}} \dicPos{m}[4] \dicFlx{(‑s, ‑ar)}[4] \textbf{1.} \dicLangCat{hovor.} \dicDirectTranslationCS{pět (korun)} \dicIndirectTranslationCS{(částka)}  \textbf{2.} \dicLangCat{hovor.} \dicDirectTranslationCS{pětikoruna, pětikačka} \dicIndirectTranslationCS{(mince)}
\dicEntry[fimmstrendur] \dicTerm{fimm··strendur} \dicIPA{{f}{\textsci}{m}{s}{\textsubring{d}}{r}{\textepsilon}{n}{\textsubring{d}}{\textscy}{\textsubring{r}}} \dicPos{adj}[2]\dicFlx{}[-17] \dicDirectTranslationCS{pětistranný, pětihranný, pětiboký}
\dicEntry[fimmtardómur] \dicTerm{fimmtar··dóm|ur} \dicIPA{{f}{\textsci}{\textsubring{m}}{\textsubring{d}}{a}{r}{\textsubring{d}}{ou}{m}{\textscy}{\textsubring{r}}} \dicPos{m}[6] \dicFlx{(‑s, ‑ar)}[10] \dicFieldCat{hist.} \dicDirectTranslationCS{nejvyšší soud původního Altingu}
\dicEntry[fimmtarþraut] \dicTerm{fimmtar··þraut} \dicIPA{{f}{\textsci}{\textsubring{m}}{\textsubring{d}}{a}{\textsubring{r}}{\texttheta}{r}{\oe i}{\textsubring{d}}} \dicPos{f}[7] \dicFlx{(‑ar)}[3] \dicFieldCat{sport.} \dicDirectTranslationCS{pětiboj}
\dicEntry[fimmtán] \dicTerm{fimm··tán} \dicsymFrequent\  \dicIPA{{f}{\textsci}{\textsubring{m}}{\textsubring{d}}{au}{\textsubring{n}}} \dicPos{num} \dicFlx{indecl} \dicDirectTranslationCS{patnáct} \dicExampleIS{fimmtán ára hjónaband} \dicExampleCS{patnáctileté manželství}
\dicEntry[fimmtándi] \dicTerm{fimm··tándi} \dicIPA{{f}{\textsci}{\textsubring{m}}{\textsubring{d}}{au}{n}{\textsubring{d}}{\textsci}} \dicPos{num}[5] \dicFlx{ord} \dicDirectTranslationCS{patnáctý}
\dicEntry[fimmti] \dicTerm{fimmti} \dicsymFrequent\  \dicIPA{{f}{\textsci}{\textsubring{m}}{\textsubring{d}}{\textsci}} \dicPos{num}[5] \dicFlx{ord} \dicDirectTranslationCS{pátý} \dicExampleIS{í fimmta hluta bókar} \dicExampleCS{v~páté části knihy}
\dicEntry[fimmtíu] \dicTerm{fimm··tíu} \dicsymFrequent\  \dicIPA{{f}{\textsci}{\textsubring{m}}{t\smash{\textsuperscript{h}}}{i}{j}{\textscy}} \dicPos{num} \dicFlx{indecl} \dicDirectTranslationCS{padesát} \dicExampleIS{Við vorum fimmtíu strákarnir í skólanum.} \dicExampleCS{Ve škole nás bylo padesát chlapců.}
\dicEntry[fimmtíukall] \dicTerm{fimm·tíu··kall} \dicIPA{{f}{\textsci}{\textsubring{m}}{t\smash{\textsuperscript{h}}}{i}{j}{\textscy}{k\smash{\textsuperscript{h}}}{a}{\textsubring{d}}{\textsubring{l}}} \dicPos{m}[4] \dicFlx{(‑s, ‑ar)}[4] \textbf{1.} \dicLangCat{hovor.} \dicDirectTranslationCS{padesát (korun)} \dicIndirectTranslationCS{(částka)}  \textbf{2.} \dicLangCat{hovor.} \dicDirectTranslationCS{padesátikoruna} \dicIndirectTranslationCS{(mince)}
\dicEntry[fimmtudagur] \dicTerm{fimmtu··dag|ur} \dicIPA{{f}{\textsci}{\textsubring{m}}{\textsubring{d}}{\textscy}{\textsubring{d}}{a}{\textbabygamma}{\textscy}{\textsubring{r}}} \dicPos{m}[6] \dicFlx{(‑s, ‑ar)}[62] \dicDirectTranslationCS{čtvrtek}
\dicEntry[fimmtugasti] \dicTerm{fimm·tug··asti} \dicIPA{{f}{\textsci}{\textsubring{m}}{\textsubring{d}}{\textscy}{\textbabygamma}{a}{s}{\textsubring{d}}{\textsci}} \dicPos{num}[6] \dicFlx{ord} \dicDirectTranslationCS{padesátý}
\dicEntry[fimmtugfaldur] \dicTerm{fimm·tug··|faldur} \dicIPA{{f}{\textsci}{\textsubring{m}}{\textsubring{d}}{\textscy}{x}{f}{a}{l}{\textsubring{d}}{\textscy}{\textsubring{r}}} \dicPos{adj}[2] \dicFlx{(f ‑föld)}[16] \textbf{1.} \dicDirectTranslationCS{padesátinásobný}  \textbf{2.} \dicDirectTranslationCS{padesátidílný}
\dicEntry[fimmtugsaldur] \dicTerm{fimm·tugs··ald|ur} \dicIPA{{f}{\textsci}{\textsubring{m}}{\textsubring{d}}{\textscy}{x}{s}{a}{l}{\textsubring{d}}{\textscy}{\textsubring{r}}} \dicPos{m}[5] \dicFlx{(‑urs, ‑rar)}[3] \dicDirectTranslationCS{věk po čtyřicítce, věk mezi 40. a~49. rokem};  \dicPhraseIS{e‑r er á fimmtugsaldri} \dicDirectTranslationCS{(kdo) je čtyřicátník}
\dicEntry[fimmtugur] \dicTerm{fimm··tugur} \dicIPA{{f}{\textsci}{\textsubring{m}}{\textsubring{d}}{\textscy}{\textbabygamma}{\textscy}{\textsubring{r}}} \dicPos{adj}[1]\dicFlx{}[-6] \textbf{1.} \dicDirectTranslationCS{padesátiletý}  \textbf{2.} \dicDirectTranslationCS{měřící padesát sáhů}
\dicEntry[fimmund] \dicTerm{fimm··und} \dicIPA{{f}{\textsci}{m}{\textlengthmark}{\textscy}{n}{\textsubring{d}}} \dicPos{f}[7] \dicFlx{(‑ar, ‑ir)}[1] \dicFieldCat{hud.} \dicDirectTranslationCS{kvinta}
\dicEntry[fimmþúsundkall] \dicTerm{fimm·þúsund··kall} \dicIPA{{f}{\textsci}{m}{\texttheta}{u}{s}{\textscy}{n}{\textsubring{d}}{k\smash{\textsuperscript{h}}}{a}{\textsubring{d}}{\textsubring{l}}} \dicPos{m}[4] \dicFlx{(‑s, ‑ar)}[4] \textbf{1.} \dicLangCat{hovor.} \dicDirectTranslationCS{pět tisíc (korun)} \dicIndirectTranslationCS{(částka)}  \textbf{2.} \dicLangCat{hovor.} \dicDirectTranslationCS{pětitisícikoruna, pětitisícovka} \dicIndirectTranslationCS{(bankovka)}
\dicEntry[fimur] \dicTerm{fimur} \dicIPA{{f}{\textsci}{\textlengthmark}{m}{\textscy}{\textsubring{r}}} \dicPos{adj}[1]\dicFlx{}[-1] \dicSynonym{liðugur} \dicDirectTranslationCS{hbitý, obratný} \dicExampleIS{vera fimur í hreyfingum} \dicExampleCS{mít hbité pohyby}
\dicEntry[fingramál] \dicTerm{fingra··mál} \dicIPA{{f}{i}{\ng}{\r{g}}{r}{a}{m}{au}{\textsubring{l}}} \dicPos{n}[2] \dicFlx{(‑s, ‑)}[5] \dicDirectTranslationCS{znaková řeč}
\dicEntry[fingrasetning] \dicTerm{fingra··set·ning} \dicIPA{{f}{i}{\ng}{\r{g}}{r}{a}{s}{\textepsilon}{h}{\textsubring{d}}{n}{i}{\ng}{\r{g}}} \dicPos{f}[4] \dicFlx{(‑ar)}[7] \textbf{1.} \dicSynonym*{stelling við vélritun} \dicDirectTranslationCS{pozice\,/\addthin postavení prstů}  \textbf{2.} \dicDirectTranslationCS{prstoklad}
\dicEntry[fingur] \dicTerm{fing|ur} \dicsymFrequent\  \dicIPA{{f}{i}{\ng}{\r{g}}{\textscy}{\textsubring{r}}} \dicPos{m}[5] \dicFlx{(‑urs, ‑ur)}[13] \dicDirectTranslationCS{prst};  \dicPhraseIS{leika við hvern sinn fingur} \dicLangCat{přen.} \dicDirectTranslationCS{být dobře naladěn, mít dobrou náladu, být v~dobrém rozmaru};  \dicPhraseIS{sjá í gegnum fingur við e‑n} \dicLangCat{přen.} \dicDirectTranslationCS{dívat se na (koho) skrz prsty};  \dicPhraseIS{vefja e‑m um fingur sér} \dicLangCat{přen.} \dicDirectTranslationCS{omotat si (koho) kolem prstu}
\dicEntry[fingurbjörg] \dicTerm{fingur··|björg} \dicIPA{{f}{i}{\ng}{\r{g}}{\textscy}{r}{\textsubring{b}}{j}{\oe}{r}{\r{g}}} \dicPos{f}[7] \dicFlx{(‑bjargar, ‑bjargir)}[16] \dicDirectTranslationCS{náprstek}
\dicEntry[fingurgómur] \dicTerm{fingur··góm|ur} \dicIPA{{f}{i}{\ng}{\r{g}}{\textscy}{r}{\r{g}}{ou}{m}{\textscy}{\textsubring{r}}} \dicPos{m}[6] \dicFlx{(‑s, ‑ar)}[10] \dicDirectTranslationCS{bříško prstu}
\dicEntry[finka] \dicTerm{fink|a} \dicIPA{{f}{i}{\r{\ng}}{\r{g}}{a}} \dicPos{f}[1] \dicFlx{(‑u, ‑ur)}[19] \dicFieldCat{zool.} \dicDirectTranslationCS{pěnkava} \textit{(l.~{\textLA{Fringilla}})}  \dicsymPhoto\ 
\dicFigure{862.jpg}{Finka}{Finka - Zicha Ondřej, Biolib, Copyright/CC-BY-NC}
\dicEntry[finn] \dicTerm{finn} \dicIPA{{f}{\textsci}{n}{\textlengthmark}} \dicPos{v} \dicFlx{ind praes sg 1 pers} \dicLink{finna}
\dicEntry[finna] \dicTerm{finna} \dicsymFrequent\  \dicIPA{{f}{\textsci}{n}{\textlengthmark}{a}} \dicPos{v}[6] \dicFlx{(finn, fann, fundum, fyndi, fundið)}[32] \dicFlx{acc} \textbf{1.} \dicSynonym{uppgötva} \dicDirectTranslationCS{najít, nalézt, nalézat, nacházet} \dicExampleIS{finna úrið} \dicExampleCS{najít hodinky}  \textbf{2.} \dicSynonym{hitta} \dicDirectTranslationCS{potkat, potkávat} \dicExampleIS{Hann bað mig að finna sig eftir fundinn.} \dicExampleCS{Požádal mě, abychom se setkali po schůzi.}  \textbf{3.} \dicSynonym{skynja} \dicDirectTranslationCS{cítit, vnímat} \dicExampleIS{Finnurðu nokkra reykjarlykt?} \dicExampleCS{Cítíš nějaký kouř?};  \dicIdiom{finna}[að]{ \dicPhraseIS{finna að e‑u}} \dicSynonym{gagnrýna} \dicDirectTranslationCS{kritizovat (co), nacházet chyby na (čem)} \dicExampleIS{finna að hegðun hans} \dicExampleCS{kritizovat jeho chování};  \dicIdiom{finna}[á]{ \dicPhraseIS{finna e‑ð á sér}} \dicDirectTranslationCS{cítit (co) (na sobě)}; { \dicPhraseIS{finna á sér}} \dicDirectTranslationCS{cítit to (vliv alkoholu ap.)};  \dicIdiom{finna}[fyrir]{ \dicPhraseIS{finna fyrir e‑u}} \dicSynonym*{verða var við e‑ð} \dicDirectTranslationCS{cítit (co), pociťovat (co)} \dicExampleIS{Ég finn fyrir naglanum.} \dicExampleCS{Cítím hřebík.};  \dicIdiom{finna}[til]{ \dicPhraseIS{finna til}} \dicSynonym*{kenna sársauka} \dicDirectTranslationCS{cítit bolest, bolet, mít bolesti} \dicExampleIS{Ég finn mikið til fætinum.} \dicExampleCS{Hodně mě bolí noha.}; { \dicPhraseIS{finna til e‑s}} \dicDirectTranslationCS{cítit (co), pociťovat (co) (únavu ap.)} \dicExampleIS{finna til þreytu} \dicExampleCS{cítit únavu}; { \dicPhraseIS{finna til e‑ð}} \dicDirectTranslationCS{připravit\,/\addthin připravovat (co) (svačinu ap.)} \dicExampleIS{finna til nesti handa e‑m} \dicExampleCS{připravit pro (koho) svačinu na cestu}; { \dicPhraseIS{e‑að er vel til fundið}} \dicDirectTranslationCS{(co) je dobrý nápad, (co) je dobře vymyšleno};  \dicIdiom{finna}[til með]{ \dicPhraseIS{finna til með e‑m}} \dicSynonym{vorkenna} \dicDirectTranslationCS{mít soucit s~(kým), soucítit s~(kým)};  \dicIdiom{finna}[upp]{ \dicPhraseIS{finna upp e‑ð}} \dicSynonym*{upphugsa e‑ð} \dicDirectTranslationCS{vymyslet (co), vymýšlet (co), vynalézt (co), vynalézat (co)} \dicExampleIS{finna upp vél} \dicExampleCS{vynalézt přístroj};  \dicIdiom{finna}[upp á]{ \dicPhraseIS{finna upp á e‑u}} \dicDirectTranslationCS{přijít na (co)};  \dicIdiom{finna}[út]{ \dicPhraseIS{finna e‑ð út}} \dicDirectTranslationCS{zjistit (co), přijít na (co)};  \dicIdiom{finna}[út úr]{ \dicPhraseIS{finna út úr e‑u}} \dicDirectTranslationCS{přijít na řešení (čeho), nalézt cestu z~(čeho)};  \dicIdiom{finnast}{ \dicPhraseIS{finnast}} \dicFlx{refl} {\textbf{a.}} \dicSynonym*{vera fundinn} \dicDirectTranslationCS{být nalezen, najít se, nalézt se, nacházet se} \dicExampleIS{Orðið finnst ekki í orðabókinni.} \dicExampleCS{Slovo se nenachází ve slovníku.};  {\textbf{b.}} \dicLangCat{zast.} \dicDirectTranslationCS{potkat se, setkat se, shledat se}; { \dicPhraseIS{e‑m finnst e‑að (fallegt)}} \dicFlx{refl impers} \dicDirectTranslationCS{(komu) připadá (co (hezké)), (kdo) shledává (co (hezkým))} \dicExampleIS{Mér finnst sagan skemmtileg.} \dicExampleCS{Myslím si, že příběh je zajímavý.}; { \dicPhraseIS{láta sér fátt um finnast}} \dicDirectTranslationCS{být lhostejný\,/\addthin nezaujatý}
\dicEntry[finnandi] \dicTerm{finn··|andi} \dicIPA{{f}{\textsci}{n}{\textlengthmark}{a}{n}{\textsubring{d}}{\textsci}} \dicPos{m}[2] \dicFlx{(‑anda, ‑endur)}[1] \dicDirectTranslationCS{nálezce, nálezkyně}
\dicEntry[finnanlegur] \dicTerm{finnan··legur} \dicIPA{{f}{\textsci}{n}{\textlengthmark}{a}{n}{l}{\textepsilon}{\textbabygamma}{\textscy}{\textsubring{r}}} \dicPos{adj}[1]\dicFlx{}[-8] \dicDirectTranslationCS{(jsoucí) k~nalezení}
\dicEntry[Finni] \dicTerm{Finn|i} \dicIPA{{f}{\textsci}{n}{\textlengthmark}{\textsci}} \dicPos{m}[1] \dicFlx{(‑a, ‑ar)}[1] \dicDirectTranslationCS{Fin(ka)}
\dicEntry[Finnland] \dicTerm{Finn··land} \dicIPA{{f}{\textsci}{n}{l}{a}{n}{\textsubring{d}}} \dicPos{n}[2] \dicFlx{(‑s)}[4] \dicFieldCat{geog.} \dicDirectTranslationCS{Finsko}
\dicEntry[Finnmörk] \dicTerm{Finn··|mörk} \dicIPA{{f}{\textsci}{n}{m}{\oe}{\textsubring{r}}{\r{g}}} \dicPos{f}[10] \dicFlx{(‑markar\,/\addthin ‑merkur)}[6] \dicFieldCat{geog.} \dicDirectTranslationCS{Finnmark} \dicIndirectTranslationCS{(kraj v~Norsku)}
\dicEntry[finnska] \dicTerm{finnsk|a} \dicIPA{{f}{\textsci}{n}{s}{\r{g}}{a}} \dicPos{f}[1] \dicFlx{(‑u)}[5] \dicDirectTranslationCS{finština}
\dicEntry[finnskur] \dicTerm{finnskur} \dicIPA{{f}{\textsci}{n}{s}{\r{g}}{\textscy}{\textsubring{r}}} \dicPos{adj}[1]\dicFlx{}[-1] \dicDirectTranslationCS{finský}
\dicEntry[finnungur] \dicTerm{finn··ung|ur} \dicIPA{{f}{\textsci}{n}{\textlengthmark}{u}{\ng}{\r{g}}{\textscy}{\textsubring{r}}} \dicPos{m}[6] \dicFlx{(‑s, ‑ar)}[8] \dicFieldCat{bot.} \dicDirectTranslationCS{smilka tuhá} \textit{(l.~{\textLA{Nardus stricta}})}
\dicEntry[fipa] \dicTerm{fip|a} \dicIPA{{f}{\textsci}{\textlengthmark}{\textsubring{b}}{a}} \dicPos{v}[1] \dicFlx{(‑aði)}[52] \dicFlx{acc} \dicSynonym{trufla} \dicDirectTranslationCS{(vy)rušit, poplést, zmást};  \dicIdiom{fipast}{ \dicPhraseIS{e‑m fipast}} \dicFlx{refl impers} \dicDirectTranslationCS{(kdo) je vyveden z~míry, (kdo) ztrácí nit (vyprávění ap.)}
\dicEntry[firðir] \dicTerm{firðir} \dicIPA{{f}{\textsci}{r}{ð}{\textsci}{\textsubring{r}}} \dicPos{m} \dicFlx{pl nom} \dicLink{fjörður}
\dicEntry[firma] \dicTerm{firm|a} \dicIPA{{f}{\textsci}{r}{m}{a}} \dicPos{n}[1] \dicFlx{(‑a, ‑u)}[1] \dicSynonym{fyrirtæki} \dicDirectTranslationCS{firma}
\dicEntry[firmaskrá] \dicTerm{firma··skrá} \dicIPA{{f}{\textsci}{r}{m}{a}{s}{\r{g}}{r}{au}} \dicPos{f}[4] \dicFlx{(‑r\,/\addthin ‑ar, ‑r)}[21] \dicDirectTranslationCS{seznam firem}
\dicEntry[firn] \dicTerm{firn} \dicIPA{{f}{\textsci}{r}{\textsubring{d}}{\textsubring{n}}} \dicPos{n}[2] \dicFlx{(‑s, ‑)}[5] \dicSynonym{býsn} \dicDirectTranslationCS{spousty, velké množství};  \dicPhraseIS{firnin öll af e‑u} \dicDirectTranslationCS{hory (čeho)}
\dicEntry[firra] \dicTerm{firr|a\smash{\textsuperscript{1}}} \dicIPA{{f}{\textsci}{r}{\textlengthmark}{a}} \dicPos{f}[1] \dicFlx{(‑u, ‑ur)}[7] \dicSynonym{vitleysa} \dicDirectTranslationCS{hloupost, nesmyslnost, absurdita}
\dicEntry[firra] \dicTerm{firr|a\smash{\textsuperscript{2}}} \dicIPA{{f}{\textsci}{r}{\textlengthmark}{a}} \dicPos{v}[2] \dicFlx{(‑ti, ‑t)}[63] \dicFlx{acc + dat} \dicDirectTranslationCS{vyhnout, vyhýbat, uniknout, unikat (zodpovědnosti ap.)} \dicExampleIS{firra sig vandræðum} \dicExampleCS{vyhnout se problémům}
\dicEntry[firring] \dicTerm{firr··ing} \dicIPA{{f}{\textsci}{r}{\textlengthmark}{i}{\ng}{\r{g}}} \dicPos{f}[4] \dicFlx{(‑ar)}[7] \dicSynonym*{ókynni} \dicDirectTranslationCS{odcizení} \dicIndirectTranslationCS{(citové ap.)}
\dicEntry[firrtur] \dicTerm{firrtur} \dicIPA{{f}{\textsci}{\textsubring{r}}{\textsubring{d}}{\textscy}{\textsubring{r}}} \dicPos{adj}[1]\dicFlx{}[-10] \dicDirectTranslationCS{odcizený} \dicIndirectTranslationCS{(člověk od společnosti ap.)}
\dicEntry[fis] \dicTerm{fis} \dicIPA{{f}{\textsci}{\textlengthmark}{s}} \dicPos{n}[2] \dicFlx{(‑s, ‑)}[5] \dicSynonym{ögn\smash{\textsuperscript{1}}} \dicDirectTranslationCS{trocha, troška};  \dicPhraseIS{létt eins og fis} \dicDirectTranslationCS{lehký jako peříčko\,/\addthin pírko}
\dicEntry[fiska] \dicTerm{fisk|a} \dicIPA{{f}{\textsci}{s}{\r{g}}{a}} \dicPos{v}[1] \dicFlx{(‑aði)}[22] \dicFlx{acc} \dicDirectTranslationCS{rybařit, chytat\,/\addthin lovit ryby};  \dicIdiom{fiska}[eftir]{ \dicPhraseIS{fiska eftir e‑u}} \dicDirectTranslationCS{slídit po (čem), pást po (čem)} \dicExampleIS{fiska eftir upplýsingum} \dicExampleCS{slídit po informacích};  \dicIdiom{fiska}[upp]{ \dicPhraseIS{fiska e‑ð upp}} \dicDirectTranslationCS{vylovit (co), vyhrabat (co) (z~kapsy ap.)};  \dicIdiom{fiskast}{ \dicPhraseIS{fiskast}} \dicFlx{refl} \dicDirectTranslationCS{být uloven\,/\addthin chycen (rybolovem)}
\dicEntry[fiskabúr] \dicTerm{fiska··búr} \dicIPA{{f}{\textsci}{s}{\r{g}}{a}{\textsubring{b}}{u}{\textsubring{r}}} \dicPos{n}[2] \dicFlx{(‑s, ‑)}[5] \dicDirectTranslationCS{akvárium}
\dicEntry[fiskbolla] \dicTerm{fisk··boll|a}\dicTerm{, fiskibolla} \dicIPA{{f}\-{\textsci}\-{s}\-{\r{g}}\-{\textsubring{b}}\-{\textopeno}\-{l}\-{\textlengthmark}\-{a}\-} \dicPos{f}[1] \dicFlx{(‑u, ‑ur)}[7] \dicFieldCat{kulin.} \dicDirectTranslationCS{rybí karbanátek}
\dicEntry[fiskbúð] \dicTerm{fisk··búð} \dicIPA{{f}{\textsci}{s}{\r{g}}{\textsubring{b}}{u}{\texttheta}} \dicPos{f}[7] \dicFlx{(‑ar, ‑ir)}[1] \dicDirectTranslationCS{prodejna ryb, rybárna}
\dicEntry[fiskeldi] \dicTerm{fisk··eldi} \dicIPA{{f}{\textsci}{s}{\r{g}}{\textepsilon}{l}{\textsubring{d}}{\textsci}} \dicPos{n}[2] \dicFlx{(‑s)}[20] \dicDirectTranslationCS{chov ryb}
\dicEntry[fiskflak] \dicTerm{fisk··|flak} \dicIPA{{f}{\textsci}{s}{\r{g}}{f}{l}{a}{\r{g}}} \dicPos{n}[2] \dicFlx{(‑flaks, ‑flök)}[8] \dicDirectTranslationCS{rybí filé}
\dicEntry[fiski] \dicTerm{fisk|i} \dicIPA{{f}{\textsci}{s}{\r{\textObardotlessj}}{\textsci}} \dicPos{f}[5] \dicFlx{(‑jar)}[7] \dicSynonym{fiskveiði} \dicDirectTranslationCS{rybaření};  \dicPhraseIS{róa til fiskjar} \dicDirectTranslationCS{jet na ryby}
\dicEntry[fiskibátur] \dicTerm{fiski··bát|ur} \dicIPA{{f}{\textsci}{s}{\r{\textObardotlessj}}{\textsci}{\textsubring{b}}{au}{\textsubring{d}}{\textscy}{\textsubring{r}}} \dicPos{m}[6] \dicFlx{(‑s, ‑ar)}[22] \dicDirectTranslationCS{rybářský člun, rybářská loď}
\dicEntry[fiskibolla] \dicTerm{fiski··boll|a} \dicIPA{{f}{\textsci}{s}{\r{\textObardotlessj}}{\textsci}{\textsubring{b}}{\textopeno}{l}{\textlengthmark}{a}} \dicPos{f}[1] \dicFlx{(‑u, ‑ur)}[7] \dicLink{fiskbolla}
\dicEntry[fiskiðnaður] \dicTerm{fisk··iðn·að|ur} \dicIPA{{f}{\textsci}{s}{\r{g}}{\textsci}{ð}{n}{a}{ð}{\textscy}{\textsubring{r}}} \dicPos{m}[10] \dicFlx{(‑ar)}[9] \dicDirectTranslationCS{rybářský průmysl}
\dicEntry[fiskifluga] \dicTerm{fiski··flug|a} \dicIPA{{f}{\textsci}{s}{\r{\textObardotlessj}}{\textsci}{f}{l}{\textscy}{\textbabygamma}{a}} \dicPos{f}[1] \dicFlx{(‑u, ‑ur)}[13] \dicFieldCat{zool.} \dicDirectTranslationCS{bzučivka} \textit{(l.~{\textLA{Calliphoridae}})}
\dicEntry[fiskifræði] \dicTerm{fiski··fræð|i} \dicIPA{{f}{\textsci}{s}{\r{\textObardotlessj}}{\textsci}{f}{r}{a}{i}{ð}{\textsci}} \dicPos{f}[3] \dicFlx{(‑i)}[3] \dicDirectTranslationCS{ichtyologie}
\dicEntry[fiskifræðingur] \dicTerm{fiski·fræð··ing|ur} \dicIPA{{f}{\textsci}{s}{\r{\textObardotlessj}}{\textsci}{f}{r}{a}{i}{ð}{i}{\ng}{\r{g}}{\textscy}{\textsubring{r}}} \dicPos{m}[6] \dicFlx{(‑s, ‑ar)}[8] \dicDirectTranslationCS{ichtyolog, ichtyoložka}
\dicEntry[fiskilína] \dicTerm{fiski··lín|a} \dicIPA{{f}{\textsci}{s}{\r{\textObardotlessj}}{\textsci}{l}{i}{n}{a}} \dicPos{f}[1] \dicFlx{(‑u, ‑ur)}[7] \dicFieldCat{nám.} \dicSynonym{lóð\smash{\textsuperscript{2}}} \dicDirectTranslationCS{(rybářský) vlasec}
\dicEntry[fiskimaður] \dicTerm{fiski··|maður} \dicIPA{{f}{\textsci}{s}{\r{\textObardotlessj}}{\textsci}{m}{a}{ð}{\textscy}{\textsubring{r}}} \dicPos{m}[13] \dicFlx{(‑manns, ‑menn)}[2] \dicDirectTranslationCS{rybář(ka)}
\dicEntry[fiskimið] \dicTerm{fiski··mið} \dicIPA{{f}{\textsci}{s}{\r{\textObardotlessj}}{\textsci}{m}{\textsci}{\texttheta}} \dicPos{n}[2] \dicFlx{pl}[1] \dicDirectTranslationCS{loviště ryb, oblast rybolovu}
\dicEntry[fiskimjöl] \dicTerm{fiski··mjöl} \dicIPA{{f}{\textsci}{s}{\r{\textObardotlessj}}{\textsci}{m}{j}{\oe}{\textsubring{l}}} \dicPos{n}[2] \dicFlx{(‑s)}[2] \dicDirectTranslationCS{rybí moučka}
\dicEntry[fiskirækt] \dicTerm{fiski··rækt} \dicIPA{{f}{\textsci}{s}{\r{\textObardotlessj}}{\textsci}{r}{a}{i}{x}{\textsubring{d}}} \dicPos{f}[4] \dicFlx{(‑ar)}[3] \dicDirectTranslationCS{chov ryb}
\dicEntry[fiskiskip] \dicTerm{fiski··skip} \dicIPA{{f}{\textsci}{s}{\r{\textObardotlessj}}{\textsci}{s}{\r{\textObardotlessj}}{\textsci}{\textsubring{b}}} \dicPos{n}[2] \dicFlx{(‑s, ‑)}[5] \dicDirectTranslationCS{rybářská loď, rybářský člun}
\dicEntry[Fiskistofa] \dicTerm{Fiski··stof|a} \dicIPA{{f}{\textsci}{s}{\r{\textObardotlessj}}{\textsci}{s}{\textsubring{d}}{\textopeno}{v}{a}} \dicPos{f}[1] \dicFlx{(‑u)}[6] \dicDirectTranslationCS{Ředitelství rybolovu} \dicIndirectTranslationCS{(islandská vládní instituce spadající pod Ministerstvo rybolovu)}
\dicEntry[fiskistofn] \dicTerm{fiski··stofn} \dicIPA{{f}{\textsci}{s}{\r{\textObardotlessj}}{\textsci}{s}{\textsubring{d}}{\textopeno}{\textsubring{b}}{\textsubring{n}}} \dicPos{m}[4] \dicFlx{(‑s, ‑ar)}[3] \dicDirectTranslationCS{rybí kolonie}
\dicEntry[fiskisúpa] \dicTerm{fiski··súp|a} \dicIPA{{f}{\textsci}{s}{\r{\textObardotlessj}}{\textsci}{s}{u}{\textsubring{b}}{a}} \dicPos{f}[1] \dicFlx{(‑u, ‑ur)}[19] \dicFieldCat{kulin.} \dicDirectTranslationCS{rybí polévka}
\dicEntry[fiskiþorp] \dicTerm{fiski··þorp} \dicIPA{{f}{\textsci}{s}{\r{\textObardotlessj}}{\textsci}{\texttheta}{\textopeno}{\textsubring{r}}{\textsubring{b}}} \dicPos{n}[2] \dicFlx{(‑s, ‑)}[5] \dicDirectTranslationCS{rybářská vesnice}
\dicEntry[fiskmarkaður] \dicTerm{fisk··mark·að|ur} \dicIPA{{f}{\textsci}{s}{\r{g}}{m}{a}{\textsubring{r}}{\r{g}}{a}{ð}{\textscy}{\textsubring{r}}} \dicPos{m}[10] \dicFlx{(‑ar\,/\addthin ‑s, ‑ir)}[31] \dicDirectTranslationCS{trh s~rybami, rybí trh}
\dicEntry[fisksali] \dicTerm{fisk··sal|i} \dicIPA{{f}{\textsci}{s}{\r{g}}{s}{a}{l}{\textsci}} \dicPos{m}[1] \dicFlx{(‑a, ‑ar)}[8] \dicDirectTranslationCS{prodavač(ka) ryb, obchodník\,/\addthin obchodnice s~rybami}
\dicEntry[fiskstautur] \dicTerm{fisk··staut|ur} \dicIPA{{f}{\textsci}{s}{\r{g}}{s}{\textsubring{d}}{\oe i}{\textsubring{d}}{\textscy}{\textsubring{r}}} \dicPos{m}[6] \dicFlx{(‑s, ‑ar)}[24] \dicFieldCat{kulin.} \dicDirectTranslationCS{rybí prst}
\dicEntry[fiskur] \dicTerm{fisk|ur} \dicsymFrequent\  \dicIPA{{f}{\textsci}{s}{\r{g}}{\textscy}{\textsubring{r}}} \dicPos{m}[6] \dicFlx{(‑s, ‑ar)}[10] \textbf{1.} \dicFieldCat{zool.} \dicDirectTranslationCS{ryba} \textit{(l.~{\textLA{Pisces}})} \textbf{2.} \dicPhraseIS{fiskar} \dicDirectTranslationCS{Ryby} \dicIndirectTranslationCS{(znamení zvěrokruhu)} \textit{(l.~{\textLA{Pisces}})};  \dicPhraseIS{e‑m hefur vaxið fiskur um hrygg} \dicFlx{impers} \dicDirectTranslationCS{(kdo) vyrostl jako z~vody};  \dicPhraseIS{e‑að er ekki upp á marga fiska} \dicLangCat{přen.} \dicDirectTranslationCS{(co) nestojí za zlámanou grešli}
\dicEntry[fiskveiði] \dicTerm{fisk··veið|i} \dicIPA{{f}{\textsci}{s}{\r{g}}{v}{ei}{ð}{\textsci}} \dicPos{f}[2] \dicFlx{(‑i\,/\addthin ‑ar, ‑ar)}[3] \dicSynonym{fiski} \dicDirectTranslationCS{rybaření, rybolov}
\dicEntry[fiskveiðikvóti] \dicTerm{fisk·veiði··kvót|i} \dicIPA{{f}{\textsci}{s}{\r{g}}{v}{ei}{ð}{\textsci}{k\smash{\textsuperscript{h}}}{v}{ou}{\textsubring{d}}{\textsci}} \dicPos{m}[1] \dicFlx{(‑a, ‑ar)}[1] \dicDirectTranslationCS{rybolovná kvóta}
\dicEntry[fiskverð] \dicTerm{fisk··verð} \dicsymFrequent\  \dicIPA{{f}{\textsci}{s}{\r{g}}{v}{\textepsilon}{r}{\texttheta}} \dicPos{n}[2] \dicFlx{(‑s)}[2] \dicDirectTranslationCS{cena ryb} \dicExampleIS{Breytingarnar leiða til hækkunar fiskverðs.} \dicExampleCS{Změny vedou ke zvýšení cen ryb.}
\dicEntry[fiskverkun] \dicTerm{fisk··verk|un} \dicIPA{{f}{\textsci}{s}{\r{g}}{v}{\textepsilon}{\textsubring{r}}{\r{g}}{\textscy}{\textsubring{n}}} \dicPos{f}[7] \dicFlx{(‑unar, ‑anir)}[8] \dicDirectTranslationCS{zpracování ryb}
\dicEntry[fiskvinnsla] \dicTerm{fisk··vinnsl|a} \dicIPA{{f}{\textsci}{s}{\r{g}}{v}{\textsci}{n}{s}{\textsubring{d}}{l}{a}} \dicPos{f}[1] \dicFlx{(‑u, ‑ur)}[19] \dicDirectTranslationCS{zpracování ryb}
\dicEntry[fisléttur] \dicTerm{fis··léttur} \dicIPA{{f}{\textsci}{\textlengthmark}{s}{l}{j}{\textepsilon}{h}{\textsubring{d}}{\textscy}{\textsubring{r}}} \dicPos{adj}[1]\dicFlx{}[-13] \dicDirectTranslationCS{lehoučký, lehký jako pírko}
\dicEntry[fit] \dicTerm{fit} \dicIPA{{f}{\textsci}{\textlengthmark}{\textsubring{d}}} \dicPos{f}[4] \dicFlx{(‑jar, ‑jar)}[9] \dicDirectTranslationCS{blána, blanka}
\dicEntry[fita] \dicTerm{fit|a\smash{\textsuperscript{1}}} \dicIPA{{f}{\textsci}{\textlengthmark}{\textsubring{d}}{a}} \dicPos{f}[1] \dicFlx{(‑u)}[5] \dicDirectTranslationCS{tuk, mastnota};  \dicPhraseIS{mettuð fita} \dicDirectTranslationCS{nasycený tuk}
\dicEntry[fita] \dicTerm{fit|a\smash{\textsuperscript{2}}} \dicIPA{{f}{\textsci}{\textlengthmark}{\textsubring{d}}{a}} \dicPos{v}[1] \dicFlx{(‑aði)}[1] \dicFlx{acc} \dicDirectTranslationCS{vykrmit, vykrmovat} \dicExampleIS{fita sig} \dicExampleCS{vykrmit se}
\dicEntry[fitja] \dicTerm{fitj|a} \dicIPA{{f}{\textsci}{\textlengthmark}{\textsubring{d}}{j}{a}} \dicPos{v}[1] \dicFlx{(‑aði)}[1] \textbf{1.} \dicPhraseIS{fitja upp (lykkjur)} \dicDirectTranslationCS{nahodit (pletení)}  \textbf{2.} \dicPhraseIS{fitja upp á e‑u} \dicDirectTranslationCS{nadhodit (co) (téma ap.), začít hovořit o~(čem)} \dicExampleIS{fitja upp á nýju umræðuefni} \dicExampleCS{nadhodit nové téma}  \textbf{3.} \dicPhraseIS{fitja upp á trýnið\,/\addthin nefið} \dicSynonym*{hrukka nefið} \dicDirectTranslationCS{ohrnovat nos, dělat grimasy, šklebit se}
\dicEntry[fitl] \dicTerm{fitl} \dicIPA{{f}{\textsci}{h}{\textsubring{d}}{\textsubring{l}}} \dicPos{n}[2] \dicFlx{(‑s)}[2] \dicDirectTranslationCS{hraní si, pohrávání si (nervózní ap.)}
\dicEntry[fitla] \dicTerm{fitl|a} \dicIPA{{f}{\textsci}{h}{\textsubring{d}}{l}{a}} \dicPos{v}[1] \dicFlx{(‑aði)}[44] \dicDirectTranslationCS{hrát si, pohrávat si (bezmyšlenkovitě s~prstýnkem ap.)};  \dicPhraseIS{fitla við e‑ð} \dicDirectTranslationCS{hrát si s~(čím), pohrávat si s~(čím) (nervózně)} \dicExampleIS{fitla við skeggið á sér} \dicExampleCS{pohrávat si se svými vousy}
\dicEntry[fitna] \dicTerm{fitn|a} \dicIPA{{f}{\textsci}{h}{\textsubring{d}}{n}{a}} \dicPos{v}[1] \dicFlx{(‑aði)}[44] \dicDirectTranslationCS{(z)tloustnout, přibrat}
\dicEntry[fituefni] \dicTerm{fitu··efni} \dicIPA{{f}{\textsci}{\textlengthmark}{\textsubring{d}}{\textscy}{\textepsilon}{\textsubring{b}}{n}{\textsci}} \dicPos{n}[2] \dicFlx{(‑s, ‑)}[14] \dicFieldCat{chem.} \dicDirectTranslationCS{tuk, lipid}
\dicEntry[fitugur] \dicTerm{fit··ugur} \dicIPA{{f}{\textsci}{\textlengthmark}{\textsubring{d}}{\textscy}{\textbabygamma}{\textscy}{\textsubring{r}}} \dicPos{adj}[1]\dicFlx{}[-8] \dicDirectTranslationCS{tučný, mastný}
\dicEntry[fitukirtill] \dicTerm{fitu··kirt|ill} \dicIPA{{f}{\textsci}{\textlengthmark}{\textsubring{d}}{\textscy}{c\smash{\textsuperscript{h}}}{\textsci}{\textsubring{r}}{\textsubring{d}}{\textsci}{\textsubring{d}}{\textsubring{l}}} \dicPos{m}[6] \dicFlx{(‑ils, ‑lar)}[35] \dicFieldCat{anat.} \dicDirectTranslationCS{tuková žláza}
\dicEntry[fituskertur] \dicTerm{fitu··skertur} \dicIPA{{f}{\textsci}{\textlengthmark}{\textsubring{d}}{\textscy}{s}{\r{\textObardotlessj}}{\textepsilon}{\textsubring{r}}{\textsubring{d}}{\textscy}{\textsubring{r}}} \dicPos{adj}[1]\dicFlx{}[-10] \dicDirectTranslationCS{odtučněný, (jsoucí) se sníženým obsahem tuku} \dicExampleIS{fituskert mjólk} \dicExampleCS{odtučněné mléko}
\dicEntry[fíaskó] \dicTerm{fíaskó} \dicIPA{{f}{i}{j}{a}{s}{\r{g}}{ou}} \dicPos{n}[2] \dicFlx{(‑s, ‑)}[34] \dicDirectTranslationCS{fiasko}
\dicEntry[Fídjeyingur] \dicTerm{Fídj·ey··ing|ur} \dicIPA{{f}{i}{\textlengthmark}{\textsubring{d}}{j}{ei}{\textlengthmark}{\ng}{\r{g}}{\textscy}{\textsubring{r}}} \dicPos{m}[6] \dicFlx{(‑s, ‑ar)}[8] \dicDirectTranslationCS{Fidžijec, Fidžijka}
\dicEntry[fídjeyskur] \dicTerm{fídj··ey·skur} \dicIPA{{f}{i}{\textlengthmark}{\textsubring{d}}{j}{ei}{s}{\r{g}}{\textscy}{\textsubring{r}}} \dicPos{adj}[1]\dicFlx{}[-6] \dicDirectTranslationCS{fidžijský}
\dicEntry[Fídjíeyjar] \dicTerm{Fídjí··eyjar} \dicIPA{{f}{i}{\textlengthmark}{\textsubring{d}}{j}{i}{ei}{j}{a}{\textsubring{r}}} \dicPos{f}[4] \dicFlx{pl}[15] \dicFieldCat{geog.} \dicDirectTranslationCS{Fidži} \dicIndirectTranslationCS{(souostroví v~Tichém oceánu)}
\dicEntry[fífa] \dicTerm{fíf|a} \dicIPA{{f}{i}{\textlengthmark}{v}{a}} \dicPos{f}[1] \dicFlx{(‑u, ‑ur)}[19] \dicFieldCat{bot.} \dicDirectTranslationCS{suchopýr} \textit{(l.~{\textLA{Eriophorum}})}  \dicsymPhoto\ 
\dicFigure{ds_image_fifa_0_1.jpg}{Fífa}{Fífa - Kim Hansen, CC BY-SA 3.0}
\dicEntry[fífill] \dicTerm{fíf|ill} \dicIPA{{f}{i}{\textlengthmark}{v}{\textsci}{\textsubring{d}}{\textsubring{l}}} \dicPos{m}[6] \dicFlx{(‑ils, ‑lar)}[35] \dicFieldCat{bot.} \dicDirectTranslationCS{pampeliška, smetánka} \textit{(l.~{\textLA{Taraxacum}})};  \dicPhraseIS{mega muna sinn fífil fegri} \dicLangCat{přen.} \dicDirectTranslationCS{pamatovat lepší časy}
\dicEntry[fífl] \dicTerm{fífl} \dicIPA{{f}{i}{\textsubring{b}}{\textsubring{l}}} \dicPos{n}[2] \dicFlx{(‑s, ‑)}[5] \dicDirectTranslationCS{hlupák, šašek, idiot} \dicExampleIS{haga sér eins og fífl} \dicExampleCS{chovat se jako hlupák};  \dicPhraseIS{hafa e‑n að fífli} \dicDirectTranslationCS{mít (koho) za blázna\,/\addthin hlupáka}
\dicEntry[fífla] \dicTerm{fífl|a} \dicIPA{{f}{i}{\textsubring{b}}{l}{a}} \dicPos{v}[1] \dicFlx{(‑aði)}[1] \dicFlx{acc} \dicSynonym{fleka} \dicDirectTranslationCS{pobláznit, svést, svádět} \dicExampleIS{fífla stúlkuna} \dicExampleCS{svést dívku};  \dicIdiom{fíflast}{ \dicPhraseIS{fíflast}} \dicFlx{refl} \dicDirectTranslationCS{šaškovat, blbnout, vyvádět}
\dicEntry[fíflalæti] \dicTerm{fífla··læti} \dicIPA{{f}{i}{\textsubring{b}}{l}{a}{l}{a}{i}{\textsubring{d}}{\textsci}} \dicPos{n}[2] \dicFlx{pl}[19] \dicDirectTranslationCS{šaškárny, šaškování, bláznoviny}
\dicEntry[fífldirfska] \dicTerm{fífl··dirfsk|a} \dicIPA{{f}{i}{\textsubring{b}}{\textsubring{l}}{\textsubring{d}}{\textsci}{\textsubring{r}}{s}{\r{g}}{a}} \dicPos{f}[1] \dicFlx{(‑u)}[5] \dicDirectTranslationCS{potřeštěnost, ztřeštěnost, nerozvážnost}
\dicEntry[fífldjarfur] \dicTerm{fífl··|djarfur} \dicIPA{{f}{i}{\textsubring{b}}{\textsubring{l}}{\textsubring{d}}{j}{a}{r}{v}{\textscy}{\textsubring{r}}} \dicPos{adj}[1] \dicFlx{(f ‑djörf)}[2] \dicDirectTranslationCS{ztřeštěný, šíleně odvážný, nerozvážný}
\dicEntry[fígúra] \dicTerm{fígúr|a} \dicIPA{{f}{i}{\textlengthmark}{\r{g}}{u}{r}{a}} \dicPos{f}[1] \dicFlx{(‑u, ‑ur)}[7] \dicDirectTranslationCS{směšný\,/\addthin afektovaný člověk}
\dicEntry[fíkill] \dicTerm{fík|ill} \dicIPA{{f}{i}{\textlengthmark}{\r{\textObardotlessj}}{\textsci}{\textsubring{d}}{\textsubring{l}}} \dicPos{m}[6] \dicFlx{(‑ils, ‑lar)}[35] \dicDirectTranslationCS{závislák} \dicIndirectTranslationCS{(na drogách, hazardu ap.)}
\dicEntry[fíkinn] \dicTerm{fíkinn} \dicIPA{{f}{i}{\textlengthmark}{\r{\textObardotlessj}}{\textsci}{\textsubring{n}}} \dicPos{adj}[6]\dicFlx{}[-2] \dicDirectTranslationCS{chtivý, lačný};  \dicPhraseIS{vera fíkinn í e‑ð} \dicSynonym{ágjarn} \dicDirectTranslationCS{být (dy)chtivý po (čem), žíznit po (čem)} \dicExampleIS{vera fíkinn í sælgæti} \dicExampleCS{dychtit po sladkostech}
\dicEntry[fíkja] \dicTerm{fíkj|a} \dicIPA{{f}{i}{\textlengthmark}{\r{\textObardotlessj}}{a}} \dicPos{f}[1] \dicFlx{(‑u, ‑ur)}[17] \dicDirectTranslationCS{fík} \dicIndirectTranslationCS{(plod)}
\dicEntry[fíkjutré] \dicTerm{fíkju··tré} \dicIPA{{f}{i}{\textlengthmark}{\r{\textObardotlessj}}{\textscy}{t\smash{\textsuperscript{h}}}{r}{j}{\textepsilon}} \dicPos{n}[2] \dicFlx{(‑s, ‑)}[36] \dicFieldCat{bot.} \dicDirectTranslationCS{fíkovník, fíkovník smokvoň} \textit{(l.~{\textLA{Ficus carica}})}  \dicsymPhoto\ 
\dicFigure{ds_image_fikjutre_0_1.jpg}{Fíkjutré}{Fíkjutré - AxelRohdeElias, CC BY 2.0}
\dicEntry[fíkn] \dicTerm{fíkn} \dicIPA{{f}{i}{h}{\r{g}}{\textsubring{n}}} \dicPos{f}[4] \dicFlx{(‑ar)}[3] \dicSynonym{fýsn} \dicDirectTranslationCS{závislost, (chorobná) touha}
\dicEntry[fíkniefnaneytandi] \dicTerm{fíkni·efna··neyt·|andi}\addthinS\textls[-10]{\dicIPA{{f}\-{i}\-{h}\-{\r{g}}\-{n}\-{\textsci}\-{\textepsilon}\-{\textsubring{b}}\-{n}\-{a}\-{n}\-{ei}\-{\textsubring{d}}\-{a}\-{n}\-{\textsubring{d}}\-{\textsci}}} \dicPos{m}[2] \dicFlx{(‑anda, ‑end\-ur)}[1] \dicDirectTranslationCS{narkoman(ka)}
\dicEntry[fíkniefnasali] \dicTerm{fíkni·efna··sal|i} \dicIPA{{f}{i}{h}{\r{g}}{n}{\textsci}{\textepsilon}{\textsubring{b}}{n}{a}{s}{a}{l}{\textsci}} \dicPos{m}[1] \dicFlx{(‑a, ‑ar)}[8] \dicDirectTranslationCS{prodejce\,/\addthin prodejkyně drog, drogový dealer, drogová dealerka}
\dicEntry[fíkniefni] \dicTerm{fíkni··efni} \dicIPA{{f}{i}{h}{\r{g}}{n}{\textsci}{\textepsilon}{\textsubring{b}}{n}{\textsci}} \dicPos{n}[2] \dicFlx{(‑s, ‑)}[14] \dicDirectTranslationCS{narkotikum, droga}
\dicEntry[fíla] \dicTerm{fíl|a} \dicIPA{{f}{i}{\textlengthmark}{l}{a}} \dicPos{v}[1] \dicFlx{(‑aði)}[44] \dicFlx{acc} \dicLangCat{slang.} \dicDirectTranslationCS{cítit} \dicExampleIS{fíla sig vel} \dicExampleCS{cítit se dobře}
\dicEntry[fílabein] \dicTerm{fíla··bein} \dicIPA{{f}{i}{\textlengthmark}{l}{a}{\textsubring{b}}{ei}{\textsubring{n}}} \dicPos{n}[2] \dicFlx{(‑s)}[2] \dicDirectTranslationCS{sloní kel, slonovina}
\dicEntry[Fílabeinsstrendingur] \dicTerm{Fíla·beins··strend·ing|ur} \dicIPA{{f}\-{i}\-{\textlengthmark}\-{l}\-{a}\-{\textsubring{b}}\-{ei}\-{n}\-{s}\-{\textsubring{d}}\-{r}\-{\textepsilon}\-{n}\-{\textsubring{d}}\-{i}\-{\ng}\-{\r{g}}\-{\textscy}\-{\textsubring{r}}\-} \dicPos{m}[6] \dicFlx{(‑s, ‑ar)}[8] \dicDirectTranslationCS{obyvatel(ka) Pobřeží slonoviny}
\dicEntry[Fílabeinsströndin] \dicTerm{Fíla·beins··|ströndin} \dicIPA{{f}\-{i}\-{\textlengthmark}\-{l}\-{a}\-{\textsubring{b}}\-{ei}\-{n}\-{s}\-{\textsubring{d}}\-{r}\-{\oe}\-{n}\-{\textsubring{d}}\-{\textsci}\-{\textsubring{n}}\-} \dicPos{f}[8] \dicFlx{(‑strand\-ar\-in\-nar)}[11] \dicFlx{def} \dicFieldCat{geog.} \dicDirectTranslationCS{Pobřeží slonoviny}
\dicEntry[fílapensill] \dicTerm{fíla··pens|ill} \dicIPA{{f}{i}{\textlengthmark}{l}{a}{p\smash{\textsuperscript{h}}}{\textepsilon}{n}{s}{\textsci}{\textsubring{d}}{\textsubring{l}}} \dicPos{m}[6] \dicFlx{(‑ils, ‑lar)}[35] \dicDirectTranslationCS{uher, pupínek}
\dicEntry[fílefldur] \dicTerm{fíl··efldur} \dicIPA{{f}{i}{\textlengthmark}{l}{\textepsilon}{l}{\textsubring{d}}{\textscy}{\textsubring{r}}} \dicPos{adj}[2]\dicFlx{}[-14] \dicDirectTranslationCS{udělaný, statný}
\dicEntry[fíll] \dicTerm{fíl|l} \dicIPA{{f}{i}{\textsubring{d}}{\textsubring{l}}} \dicPos{m}[6] \dicFlx{(‑s, ‑ar)}[48] \dicFieldCat{zool.} \dicDirectTranslationCS{slon} \textit{(l.~{\textLA{Elephas}})}
\dicEntry[fíngerður] \dicTerm{fín··gerður} \dicsymFrequent\  \dicIPA{{f}{i}{\textltailn}{\r{\textObardotlessj}}{\textepsilon}{r}{ð}{\textscy}{\textsubring{r}}} \dicPos{adj}[2]\dicFlx{}[-1] \dicDirectTranslationCS{jemný, drobný} \dicExampleIS{fíngerð kona} \dicExampleCS{drobná žena}
\dicEntry[fínlegur] \dicTerm{fín··legur} \dicIPA{{f}{i}{n}{l}{\textepsilon}{\textbabygamma}{\textscy}{\textsubring{r}}} \dicPos{adj}[1]\dicFlx{}[-8] \dicDirectTranslationCS{jemný, delikátní} \dicExampleIS{fínlegt smíði} \dicExampleCS{jemný výrobek}
\dicEntry[fínn] \dicTerm{fínn} \dicsymFrequent\  \dicIPA{{f}{i}{\textsubring{d}}{\textsubring{n}}} \dicPos{adj}[7]\dicFlx{}[-1] \textbf{1.} \dicSynonym{fíngerður} \dicDirectTranslationCS{jemný, delikátní, drobný} \dicExampleIS{fínn þráður} \dicExampleCS{jemné vlákno}  \textbf{2.} \dicSynonym*{í hárri stöðu} \dicDirectTranslationCS{panský, (jsoucí) z~lepší společnosti, noblesní} \dicExampleIS{fínn maður} \dicExampleCS{člověk z~lepší společnosti}  \textbf{3.} \dicSynonym{vandaður} \dicDirectTranslationCS{kvalitní} \dicExampleIS{fínar buxur} \dicExampleCS{kvalitní kalhoty}  \textbf{4.} \dicSynonym{ágætur} \dicDirectTranslationCS{dobrý, fajn, prima} \dicExampleIS{fínt veður} \dicExampleCS{dobré počasí}  \textbf{5.} \dicSynonym*{í betri fötunum} \dicDirectTranslationCS{hezky\,/\addthin dobře oblečený};  \dicPhraseIS{allt í fína lagi} \dicDirectTranslationCS{všechno v~pořádku, pohoda}
\dicEntry[fínpússa] \dicTerm{fín··púss|a} \dicIPA{{f}{i}{n}{p\smash{\textsuperscript{h}}}{u}{s}{a}} \dicPos{v}[1] \dicFlx{(‑aði)}[1] \dicFlx{acc} \textbf{1.} \dicSynonym*{fínhúða} \dicDirectTranslationCS{omítnout, nahodit} \dicExampleIS{fínpússa vegg} \dicExampleCS{omítnout zeď}  \textbf{2.} \dicDirectTranslationCS{doladit} \dicExampleIS{fínpússa ritgerðina} \dicExampleCS{doladit slohovou práci}
\dicEntry[físibelgur] \dicTerm{físi··belg|ur} \dicIPA{{f}{i}{\textlengthmark}{s}{\textsci}{\textsubring{b}}{\textepsilon}{l}{\r{g}}{\textscy}{\textsubring{r}}} \dicPos{m}[9] \dicFlx{(‑s\,/\addthin ‑jar, ‑ir)}[26] \dicDirectTranslationCS{měch} \dicIndirectTranslationCS{(zařízení na vhánění vzduchu do kovářské výhně ap.)}
\dicEntry[físisveppur] \dicTerm{físi··svepp|ur} \dicIPA{{f}{i}{\textlengthmark}{s}{\textsci}{s}{v}{\textepsilon}{h}{\textsubring{b}}{\textscy}{\textsubring{r}}} \dicPos{m}[6] \dicFlx{(‑s, ‑ir\,/\addthin ‑ar)}[73] \dicFieldCat{bot.} \dicDirectTranslationCS{pýchavka, pýchavka obecná} \textit{(l.~{\textLA{Lycoperdon}})}
\dicEntry[fjaðra] \dicTerm{fjaðr|a} \dicIPA{{f}{j}{a}{ð}{r}{a}} \dicPos{v}[1] \dicFlx{(‑aði)}[13] \dicDirectTranslationCS{(od)pružit} \dicExampleIS{Sætið fjaðrar.} \dicExampleCS{Sedadlo pruží.}
\dicEntry[fjaðradýna] \dicTerm{fjaðra··dýn|a} \dicIPA{{f}{j}{a}{ð}{r}{a}{\textsubring{d}}{i}{n}{a}} \dicPos{f}[1] \dicFlx{(‑u, ‑ur)}[7] \dicFieldCat{sport.} \dicDirectTranslationCS{trampolína}
\dicEntry[fjaðrar] \dicTerm{fjaðrar} \dicIPA{{f}{j}{a}{ð}{r}{a}{\textsubring{r}}} \dicPos{f} \dicFlx{sg gen} \dicLink{fjöður}
\dicEntry[fjaðrir] \dicTerm{fjaðrir} \dicIPA{{f}{j}{a}{ð}{r}{\textsci}{\textsubring{r}}} \dicPos{f} \dicFlx{pl nom} \dicLink{fjöður}
\dicEntry[fjaðurmagnaður] \dicTerm{fjaður··|magn·aður} \dicIPA{{f}{j}{a}{ð}{\textscy}{r}{m}{a}{\r{g}}{n}{a}{ð}{\textscy}{\textsubring{r}}} \dicPos{adj}[3] \dicFlx{(f ‑mögnuð)}[1] \dicDirectTranslationCS{ohebný, pružný, elastický}
\dicEntry[fjalar] \dicTerm{fjalar} \dicIPA{{f}{j}{a}{\textlengthmark}{l}{a}{\textsubring{r}}} \dicPos{f} \dicFlx{sg gen} \dicLink{fjöl\smash{\textsuperscript{1}}}
\dicEntry[fjalir] \dicTerm{fjalir} \dicIPA{{f}{j}{a}{\textlengthmark}{l}{\textsci}{\textsubring{r}}} \dicPos{f} \dicFlx{pl nom} \dicLink{fjöl\smash{\textsuperscript{1}}}
\dicEntry[fjall] \dicTerm{fjall} \dicsymFrequent\  \dicIPA{{f}{j}{a}{\textsubring{d}}{\textsubring{l}}} \dicPos{n}[2] \dicFlx{(‑s, fjöll)}[8] \dicDirectTranslationCS{hora} \dicExampleIS{í hlíðum stærsta fjalls heims} \dicExampleCS{na svazích největší hory světa};  \dicPhraseIS{koma af fjöllum} \dicLangCat{přen.} \dicDirectTranslationCS{být velmi překvapen\,/\addthin zaskočen} \dicExampleIS{Hann virðist koma algjörlega af fjöllum.} \dicExampleCS{Vypadá, jako by spadl z~višně.};  \dicPhraseIS{það gengur fjöllunum hærra að} \dicLangCat{přen.} \dicDirectTranslationCS{proslýchá se, že};  \dicProverb\  \dicPhraseIS{Fjarlægðin gerir fjöllin blá og mennina mikla.} \dicLangCat{přís.} \dicDirectTranslationCS{Sousedova tráva je vždycky zelenější.}
\dicEntry[fjalla] \dicTerm{fjall|a} \dicsymFrequent\  \dicIPA{{f}{j}{a}{\textsubring{d}}{l}{a}} \dicPos{v}[1] \dicFlx{(‑aði)}[53] \dicPhraseIS{fjalla um e‑ð} \dicSynonym{ræða\smash{\textsuperscript{2}}} \dicDirectTranslationCS{pojednávat o~(čem), vyprávět o~(čem), zabývat se (čím)} \dicExampleIS{Bókin fjallar um íslenska þjóðhætti.} \dicExampleCS{Knížka vypráví o~islandských lidových zvycích.}
\dicEntry[fjallafinka] \dicTerm{fjalla··fink|a} \dicIPA{{f}{j}{a}{\textsubring{d}}{l}{a}{f}{i}{\r{\ng}}{\r{g}}{a}} \dicPos{f}[1] \dicFlx{(‑u, ‑ur)}[19] \dicFieldCat{zool.} \dicDirectTranslationCS{pěnkava jikavec} \textit{(l.~{\textLA{Fringilla montifringilla}})}  \dicsymPhoto\ 
\dicFigure{ds_image_fjallafinka_0_1.jpg}{Fjallafinka}{Fjallafinka - mmlolek, CC BY 2.0}
\dicEntry[fjallafoxgras] \dicTerm{fjalla··fox·|gras} \dicIPA{{f}{j}{a}{\textsubring{d}}{l}{a}{f}{\textopeno}{x}{s}{\r{g}}{r}{a}{s}} \dicPos{n}[2] \dicFlx{(‑grass, ‑grös)}[8] \dicFieldCat{bot.} \dicDirectTranslationCS{bojínek alpinský} \textit{(l.~{\textLA{Phleum alpinum}})}  \dicsymPhoto\ 
\dicFigure{45152.jpg}{Fjallafoxgras}{Fjallafoxgras - Kesl Michael, Biolib, Copyright/CC-BY-NC}
\dicEntry[fjallagrös] \dicTerm{fjalla··grös} \dicIPA{{f}{j}{a}{\textsubring{d}}{l}{a}{\r{g}}{r}{\oe}{s}} \dicPos{n}[2] \dicFlx{pl}[9] \dicFieldCat{bot.} \dicDirectTranslationCS{pukléřka islandská, lišejník islandský} \textit{(l.~{\textLA{Cetraria islandica}})}  \dicsymPhoto\ 
\dicFigure{5301.jpg}{Fjallagrös}{Fjallagrös - Bouda František, Biolib, Copyright/PD}
\dicEntry[fjallahjól] \dicTerm{fjalla··hjól} \dicIPA{{f}{j}{a}{\textsubring{d}}{l}{a}{\c{c}}{ou}{\textsubring{l}}} \dicPos{n}[2] \dicFlx{(‑s, ‑)}[5] \dicDirectTranslationCS{horské kolo}
\dicEntry[fjallakornblóm] \dicTerm{fjalla··korn·blóm} \dicIPA{{f}{j}{a}{\textsubring{d}}{l}{a}{k\smash{\textsuperscript{h}}}{\textopeno}{r}{\textsubring{d}}{\textsubring{n}}{\textsubring{b}}{l}{ou}{\textsubring{m}}} \dicPos{n}[2] \dicFlx{(‑s, ‑)}[5] \dicFieldCat{bot.} \dicDirectTranslationCS{chrpa horská} \textit{(l.~{\textLA{Centaurea montana}})}  \dicsymPhoto\ 
\dicFigure{87146.jpg}{Fjallakornblóm}{Fjallakornblóm - Pleva František, Biolib, PD}
\dicEntry[fjallaljón] \dicTerm{fjalla··ljón} \dicIPA{{f}{j}{a}{\textsubring{d}}{l}{a}{l}{j}{ou}{\textsubring{n}}} \dicPos{n}[2] \dicFlx{(‑s, ‑)}[5] \dicFieldCat{zool.} \dicDirectTranslationCS{puma} \textit{(l.~{\textLA{Puma}})}
\dicEntry[fjallamaður] \dicTerm{fjalla··|maður} \dicIPA{{f}{j}{a}{\textsubring{d}}{l}{a}{m}{a}{ð}{\textscy}{\textsubring{r}}} \dicPos{m}[13] \dicFlx{(‑manns, ‑menn)}[2] \dicDirectTranslationCS{horolezec, horolezkyně}
\dicEntry[fjallarefur] \dicTerm{fjalla··ref|ur} \dicIPA{{f}{j}{a}{\textsubring{d}}{l}{a}{r}{\textepsilon}{v}{\textscy}{\textsubring{r}}} \dicPos{m}[9] \dicFlx{(‑s, ‑ir)}[8] \dicFieldCat{zool.} \dicDirectTranslationCS{liška polární} \textit{(l.~{\textLA{Alopex lagopus}})}
\dicEntry[fjallasmári] \dicTerm{fjalla··smár|i} \dicIPA{{f}{j}{a}{\textsubring{d}}{l}{a}{s}{m}{au}{r}{\textsci}} \dicPos{m}[1] \dicFlx{(‑a, ‑ar)}[1] \dicFieldCat{bot.} \dicDirectTranslationCS{zubatka poléhavá} \textit{(l.~{\textLA{Sibbaldia procumbens}})}  \dicsymPhoto\ 
\dicFigure{ds_image_fjallasmari_0_2.jpg}{Fjallasmári}{Fjallasmári - pellaea, CC BY 2.0}
\dicEntry[fjalldalafífill] \dicTerm{fjall·dala··fíf|ill} \dicIPA{{f}{j}{a}{\textsubring{d}}{\textsubring{l}}{\textsubring{d}}{a}{l}{a}{f}{i}{v}{\textsci}{\textsubring{d}}{\textsubring{l}}} \dicPos{m}[6] \dicFlx{(‑ils, ‑lar)}[35] \dicFieldCat{bot.} \dicDirectTranslationCS{kuklík potoční} \textit{(l.~{\textLA{Geum rivale}})}  \dicsymPhoto\ 
\dicFigure{5143.jpg}{Fjalldalafífill}{Fjalldalafífill - Zicha Ondřej, Biolib, Copyright/CC-BY-NC}
\dicEntry[fjalldrapi] \dicTerm{fjall··drap|i} \dicIPA{{f}{j}{a}{\textsubring{d}}{\textsubring{l}}{\textsubring{d}}{r}{a}{\textsubring{b}}{\textsci}} \dicPos{m}[1] \dicFlx{(‑a)}[3] \dicFieldCat{bot.} \dicDirectTranslationCS{bříza zakrslá} \textit{(l.~{\textLA{Betula nana}})}  \dicsymPhoto\ 
\dicFigure{ds_image_fjalldrapi_0_1.jpg}{Fjalldrapi}{Fjalldrapi - Sten Porse, CC BY-SA 3.0}
\dicEntry[fjallganga] \dicTerm{fjall··|ganga} \dicIPA{{f}{j}{a}{\textsubring{d}}{\textsubring{l}}{\r{g}}{au}{\ng}{\r{g}}{a}} \dicPos{f}[1] \dicFlx{(‑göngu, ‑göngur)}[14] \dicDirectTranslationCS{(vysokohorská) turistika, (vysokohorský) pochod} \dicExampleIS{fara í fjallgöngu} \dicExampleCS{vydat se na vysokohorský pochod}
\dicEntry[fjallgarður] \dicTerm{fjall··garð|ur} \dicIPA{{f}{j}{a}{\textsubring{d}}{\textsubring{l}}{\r{g}}{a}{r}{ð}{\textscy}{\textsubring{r}}} \dicPos{m}[6] \dicFlx{(‑s, ‑ar)}[5] \dicDirectTranslationCS{pohoří, horské pásmo}
\dicEntry[fjallgöngumaður] \dicTerm{fjall·göngu··|maður} \textls[18]{\dicIPA{{f}{j}{a}{\textsubring{d}}{\textsubring{l}}{\r{g}}{\oe i}{\ng}{\r{g}}{\textscy}{m}{a}{ð}{\textscy}{\textsubring{r}}} \dicPos{m}[13] \dicFlx{(‑manns, ‑menn)}[2]} \dicDirectTranslationCS{horolezec, horolezkyně, alpinista, alpinistka}
\dicEntry[fjallkjói] \dicTerm{fjall··kjó|i} \dicIPA{{f}{j}{a}{\textsubring{d}}{\textsubring{l}}{c\smash{\textsuperscript{h}}}{ou}{\textsci}} \dicPos{m}[1] \dicFlx{(‑a, ‑ar)}[1] \dicFieldCat{zool.} \dicDirectTranslationCS{chaluha malá} \textit{(l.~{\textLA{Stercorarius longicaudus}})}  \dicsymPhoto\ 
\dicFigure{21093.jpg}{Fjallkjói}{Fjallkjói - Bowman, Tim, Biolib, PD}
\dicEntry[Fjallkona] \dicTerm{Fjall··|kona} \dicIPA{{f}{j}{a}{\textsubring{d}}{\textsubring{l}}{k\smash{\textsuperscript{h}}}{\textopeno}{n}{a}} \dicPos{f}[1] \dicFlx{(‑kvenu, ‑kvenur)}[28] \dicIndirectTranslationCS{královna hor -- zosobnění Islandu}
\dicEntry[fjalllendi] \dicTerm{fjall··lendi} \dicIPA{{f}{j}{a}{\textsubring{d}}{l}{\textepsilon}{n}{\textsubring{d}}{\textsci}} \dicPos{n}[2] \dicFlx{(‑s)}[20] \dicDirectTranslationCS{(horský) masiv, horská oblast}
\dicEntry[fjallmyndarlegur] \dicTerm{fjall··myndar·legur} \dicIPA{{f}{j}{a}{\textsubring{d}}{l}{m}{\textsci}{n}{\textsubring{d}}{a}{r}{l}{\textepsilon}{\textbabygamma}{\textscy}{\textsubring{r}}} \dicPos{adj}[1]\dicFlx{}[-6] \dicDirectTranslationCS{velmi elegantní\,/\addthin pohledný}
\dicEntry[fjallshlíð] \dicTerm{fjalls··hlíð} \dicIPA{{f}{j}{a}{l}{s}{\textsubring{l}}{i}{\texttheta}} \dicPos{f}[4] \dicFlx{(‑ar, ‑ar\,/\addthin ‑ir)}[30] \dicDirectTranslationCS{svah, úbočí, stráň (kopce ap.)}
\dicEntry[fjallvegur] \dicTerm{fjall··veg|ur} \dicIPA{{f}{j}{a}{\textsubring{d}}{l}{v}{\textepsilon}{\textbabygamma}{\textscy}{\textsubring{r}}} \dicPos{m}[10] \dicFlx{(‑ar, ‑ir)}[4] \dicDirectTranslationCS{horská cesta}
\dicEntry[fjandi] \dicTerm{fjand|i\smash{\textsuperscript{1}}} \dicsymFrequent\  \dicIPA{{f}{j}{a}{n}{\textsubring{d}}{\textsci}} \dicPos{m}[1] \dicFlx{(‑a, ‑ar\,/\addthin fjendur)}[18] \dicDirectTranslationCS{ďábel, čert};  \dicPhraseIS{fara með e‑ð til fjandans} \dicDirectTranslationCS{jít s~(čím) k~čertu};  \dicPhraseIS{farðu til fjandans!} \dicDirectTranslationCS{jdi k~čertu!};  \dicPhraseIS{fjandinn hafi það} \dicDirectTranslationCS{čert to vem};  \dicPhraseIS{hver fjandinn!} \dicDirectTranslationCS{hergot!, sakramentsky!}
\dicEntry[fjandi] \dicTerm{fjandi\smash{\textsuperscript{2}}} \dicIPA{{f}{j}{a}{n}{\textsubring{d}}{\textsci}} \dicPos{adv} \dicDirectTranslationCS{zatraceně, sakramentsky} \dicExampleIS{fjandi erfitt verkefni} \dicExampleCS{zatraceně těžký úkol}
\dicEntry[fjandmaður] \dicTerm{fjand··|maður} \dicIPA{{f}{j}{a}{n}{\textsubring{d}}{m}{a}{ð}{\textscy}{\textsubring{r}}} \dicPos{m}[13] \dicFlx{(‑manns, ‑menn)}[2] \dicSynonym{óvinur} \dicDirectTranslationCS{nepřítel(kyně), protivník, protivnice}
\dicEntry[fjandsamlegur] \dicTerm{fjand·sam··legur} \dicIPA{{f}{j}{a}{n}{\textsubring{d}}{s}{a}{m}{l}{\textepsilon}{\textbabygamma}{\textscy}{\textsubring{r}}} \dicPos{adj}[1]\dicFlx{}[-8] \dicSynonym{óvinveittur} \dicDirectTranslationCS{nepřátelský}
\dicEntry[fjandskapur] \dicTerm{fjand··skap|ur} \dicIPA{{f}{j}{a}{n}{\textsubring{d}}{s}{\r{g}}{a}{\textsubring{b}}{\textscy}{\textsubring{r}}} \dicPos{m}[10] \dicFlx{(‑ar)}[15] \dicDirectTranslationCS{nepřátelství, zášť} \dicExampleIS{sýna e‑m fullan fjandskap} \dicExampleCS{chovat ke (komu) naprostou zášť}
\dicEntry[fjara] \dicTerm{fjara\smash{\textsuperscript{1}}} \dicsymFrequent\  \dicIPA{{f}{j}{a}{\textlengthmark}{r}{a}} \dicPos{f}[1] \dicFlx{(fjöru, fjörur)}[8] \textbf{1.} \dicSynonym{útfall} \dicDirectTranslationCS{odliv} \dicAntonym{flóð}  \textbf{2.} \dicSynonym{strönd\smash{\textsuperscript{1}}} \dicDirectTranslationCS{pobřeží, (mořský) břeh} \dicExampleIS{Það dreymdi mig að ég stóð í fjöru og horfði til hafs.} \dicExampleCS{Zdálo se mi, že stojím na pobřeží a~dívám se na moře.};  \dicPhraseIS{finna e‑n í fjöru} \dicLangCat{přen.} \dicSynonym*{hefnast á e‑m} \dicDirectTranslationCS{pomstít se (komu)}
\dicEntry[fjara] \dicTerm{fjar|a\smash{\textsuperscript{2}}} \dicIPA{{f}{j}{a}{\textlengthmark}{r}{a}} \dicPos{v}[1] \dicFlx{(‑aði)}[36] \textbf{1.} \dicPhraseIS{það fjarar út} \dicFlx{impers} \dicDirectTranslationCS{je odliv}  \textbf{2.} \dicPhraseIS{e‑að fjarar út} \dicDirectTranslationCS{(co) doznívá}
\dicEntry[fjarðar] \dicTerm{fjarðar} \dicIPA{{f}{j}{a}{r}{ð}{a}{\textsubring{r}}} \dicPos{m} \dicFlx{sg gen} \dicLink{fjörður}
\dicEntry[fjargviðrast] \dicTerm{fjarg··viðr|ast} \dicIPA{{f}{j}{a}{r}{\r{g}}{v}{\textsci}{ð}{r}{a}{s}{\textsubring{d}}} \dicPos{v}[1] \dicFlx{(‑aðist)}[91] \dicFlx{refl} \dicDirectTranslationCS{slovíčkařit, reptat, láteřit} \dicExampleIS{fjargviðrast yfir e‑u} \dicExampleCS{láteřit nad (čím)}
\dicEntry[fjarhrif] \dicTerm{fjar··hrif} \dicIPA{{f}{j}{a}{\textsubring{r}}{\textlengthmark}{\textsci}{f}} \dicPos{n}[2] \dicFlx{pl}[1] \dicDirectTranslationCS{telepatie, jasnovidnost}
\dicEntry[fjarkennsla] \dicTerm{fjar··kennsl|a} \dicIPA{{f}{j}{a}{\textsubring{r}}{c\smash{\textsuperscript{h}}}{\textepsilon}{n}{s}{\textsubring{d}}{l}{a}} \dicPos{f}[1] \dicFlx{(‑u)}[5] \dicFieldCat{škol.} \dicDirectTranslationCS{dálková\,/\addthin distanční výuka}
\dicEntry[fjarki] \dicTerm{fjark|i} \dicIPA{{f}{j}{a}{\textsubring{r}}{\r{\textObardotlessj}}{\textsci}} \dicPos{m}[1] \dicFlx{(‑a, ‑ar)}[8] \textbf{1.} \dicDirectTranslationCS{čtyřka} \dicIndirectTranslationCS{(v~kartách)} \dicExampleIS{spaðafjarki} \dicExampleCS{piková čtyřka}  \textbf{2.} \dicDirectTranslationCS{čtyřka (autobus číslo 4 ap.)}
\dicEntry[fjarlægð] \dicTerm{fjar··lægð} \dicsymFrequent\  \dicIPA{{f}{j}{a}{r}{l}{a}{i}{\textbabygamma}{\texttheta}} \dicPos{f}[7] \dicFlx{(‑ar, ‑ir)}[1] \textbf{1.} \dicSynonym{vegalengd} \dicDirectTranslationCS{vzdálenost, odlehlost} \dicExampleIS{fjarlægðin milli þorpanna} \dicExampleCS{vzdálenost mezi vesnicemi}  \textbf{2.} \dicSynonym{fjarski} \dicDirectTranslationCS{dálka, odlehlost} \dicExampleIS{úr fjarlægð} \dicExampleCS{z~dálky}
\dicEntry[fjarlægja] \dicTerm{fjar··læg|ja} \dicsymFrequent\  \dicIPA{{f}{j}{a}{r}{l}{a}{i}{j}{a}} \dicPos{v}[2] \dicFlx{(‑ði, ‑t)}[89] \dicFlx{acc} \dicDirectTranslationCS{odstranit, odstraňovat, odklidit, odklízet, odebrat, odebírat} \dicExampleIS{fjarlægja bókina af borðinu} \dicExampleCS{odstranit knížku ze stolu};  \dicIdiom{fjarlægjast}{ \dicPhraseIS{fjarlægjast}} \dicFlx{refl} \dicDirectTranslationCS{vzdálit se, vzdalovat se}
\dicEntry[fjarlægur] \dicTerm{fjar··lægur} \dicsymFrequent\  \dicIPA{{f}{j}{a}{r}{l}{a}{i}{\textbabygamma}{\textscy}{\textsubring{r}}} \dicPos{adj}[1]\dicFlx{}[-1] \textbf{1.} \dicDirectTranslationCS{vzdálený, odlehlý, daleký} \dicExampleIS{fjarlægt fjall} \dicExampleCS{vzdálená hora} \dicAntonym{nálægur}  \textbf{2.} \dicFieldCat{jaz.} \dicDirectTranslationCS{otevřený (samohláska ap.)}
\dicEntry[fjarnám] \dicTerm{fjar··nám} \dicIPA{{f}{j}{a}{r}{n}{au}{\textsubring{m}}} \dicPos{n}[2] \dicFlx{(‑s)}[2] \dicFieldCat{škol.} \dicDirectTranslationCS{dálkové\,/\addthin distanční studium}
\dicEntry[fjarri] \dicTerm{fjarri} \dicsymFrequent\  \dicIPA{{f}{j}{a}{r}{\textlengthmark}{\textsci}} \dicPos{adv} \dicFlx{(comp fjær, sup fjærst)} \dicSynonym*{langt frá} \dicDirectTranslationCS{daleko};  \dicPhraseIS{halda sig fjarri} \dicDirectTranslationCS{držet se stranou};  \dicPhraseIS{vera fjarri sanni} \dicLangCat{přen.} \dicDirectTranslationCS{být daleko od pravdy};  \dicPhraseIS{það er e‑m fjarri skapi} \dicDirectTranslationCS{to je (komu) proti mysli};  \dicPhraseIS{því fer fjarri} \dicDirectTranslationCS{tak to není}
\dicEntry[fjarriti] \dicTerm{fjar··rit|i} \dicIPA{{f}{j}{a}{r}{\textlengthmark}{\textsci}{\textsubring{d}}{\textsci}} \dicPos{m}[1] \dicFlx{(‑a, ‑ar)}[1] \dicFieldCat{hist.} \dicDirectTranslationCS{dálnopis}
\dicEntry[fjarrænn] \dicTerm{fjar··rænn} \dicIPA{{f}{j}{a}{r}{\textlengthmark}{a}{i}{\textsubring{d}}{\textsubring{n}}} \dicPos{adj}[7]\dicFlx{}[-1] \dicDirectTranslationCS{nepřítomný, netečný}
\dicEntry[fjarski] \dicTerm{fjarsk|i} \dicsymFrequent\  \dicIPA{{f}{j}{a}{\textsubring{r}}{s}{\r{\textObardotlessj}}{\textsci}} \dicPos{m}[1] \dicFlx{(‑a)}[3] \dicDirectTranslationCS{dálka, dál} \dicExampleIS{í fjarska} \dicExampleCS{v~dálce}
\dicEntry[fjarskiptakerfi] \dicTerm{fjar·skipta··kerfi} \dicIPA{{f}{j}{a}{\textsubring{r}}{s}{\r{\textObardotlessj}}{\textsci}{f}{\textsubring{d}}{a}{c\smash{\textsuperscript{h}}}{\textepsilon}{r}{v}{\textsci}} \dicPos{n}[2] \dicFlx{(‑s, ‑)}[14] \dicDirectTranslationCS{telekomunikační systém}
\dicEntry[fjarskiptatæki] \dicTerm{fjar·skipta··tæki} \dicIPA{{f}{j}{a}{\textsubring{r}}{s}{\r{\textObardotlessj}}{\textsci}{f}{\textsubring{d}}{a}{t\smash{\textsuperscript{h}}}{a}{i}{\r{\textObardotlessj}}{\textsci}} \dicPos{n}[2] \dicFlx{(‑s, ‑)}[16] \dicDirectTranslationCS{telekomunikační zařízení}
\dicEntry[fjarskiptatækni] \dicTerm{fjar·skipta··tækn|i} \dicIPA{{f}{j}{a}{\textsubring{r}}{s}{\r{\textObardotlessj}}{\textsci}{f}{\textsubring{d}}{a}{t\smash{\textsuperscript{h}}}{a}{i}{h}{\r{g}}{n}{\textsci}} \dicPos{f}[3] \dicFlx{(‑i)}[3] \dicDirectTranslationCS{telekomunikační technologie}
\dicEntry[fjarskipti] \dicTerm{fjar··skipti} \dicIPA{{f}{j}{a}{\textsubring{r}}{s}{\r{\textObardotlessj}}{\textsci}{f}{\textsubring{d}}{\textsci}} \dicPos{n}[2] \dicFlx{pl}[19] \dicDirectTranslationCS{telekomunikace}
\dicEntry[fjarskyggn] \dicTerm{fjar··skyggn} \dicIPA{{f}{j}{a}{\textsubring{r}}{s}{\r{\textObardotlessj}}{\textsci}{\r{g}}{\textsubring{n}}} \dicPos{adj}[5]\dicFlx{}[-5] \dicDirectTranslationCS{jasnovidný}
\dicEntry[fjarstaddur] \dicTerm{fjar··|staddur} \dicIPA{{f}{j}{a}{\textsubring{r}}{s}{\textsubring{d}}{a}{\textsubring{d}}{\textscy}{\textsubring{r}}} \dicPos{adj}[2] \dicFlx{(f ‑stödd)}[20] \dicDirectTranslationCS{nepřítomný, (jsoucí) pryč}
\dicEntry[fjarstýra] \dicTerm{fjar··stýr|a} \dicIPA{{f}{j}{a}{\textsubring{r}}{s}{\textsubring{d}}{i}{r}{a}} \dicPos{v}[2] \dicFlx{(‑ði, ‑t)}[99] \dicFlx{dat} \dicDirectTranslationCS{dálkově ovládat}
\dicEntry[fjarstýrður] \dicTerm{fjar··stýrður} \dicIPA{{f}{j}{a}{\textsubring{r}}{s}{\textsubring{d}}{i}{r}{ð}{\textscy}{\textsubring{r}}} \dicPos{adj}[2]\dicFlx{}[-4] \dicDirectTranslationCS{dálkově ovládaný}
\dicEntry[fjarstýring] \dicTerm{fjar··stýr·ing} \dicIPA{{f}{j}{a}{\textsubring{r}}{s}{\textsubring{d}}{i}{r}{i}{\ng}{\r{g}}} \dicPos{f}[4] \dicFlx{(‑ar, ‑ar)}[5] \textbf{1.} \dicDirectTranslationCS{dálkový ovladač (k~televizi ap.)}  \textbf{2.} \dicDirectTranslationCS{dálkové ovládání}
\dicEntry[fjarstæða] \dicTerm{fjar··stæð|a} \dicIPA{{f}{j}{a}{\textsubring{r}}{s}{\textsubring{d}}{a}{i}{ð}{a}} \dicPos{f}[1] \dicFlx{(‑u, ‑ur)}[13] \dicSynonym{firra\smash{\textsuperscript{1}}} \dicDirectTranslationCS{absurdita, nesmyslnost}
\dicEntry[fjarstæður] \dicTerm{fjar··stæður} \dicIPA{{f}{j}{a}{\textsubring{r}}{s}{\textsubring{d}}{a}{i}{ð}{\textscy}{\textsubring{r}}} \dicPos{adj}[2]\dicFlx{}[-6] \dicSynonym{fráleitur} \dicDirectTranslationCS{absurdní, nereálný, nedávající smysl};  \dicPhraseIS{það er fjarstætt að (gera e‑ð)} \dicDirectTranslationCS{to je absurdní (udělat (co))}
\dicEntry[fjarsýni] \dicTerm{fjar··sýn|i} \dicIPA{{f}{j}{a}{\textsubring{r}}{s}{i}{n}{\textsci}} \dicPos{f}[3] \dicFlx{(‑i)}[3] \textbf{1.} \dicFieldCat{med.} \dicDirectTranslationCS{dalekozrakost} \dicAntonym{nærsýni}  \textbf{2.} \dicDirectTranslationCS{dalekozrakost, předvídavost}
\dicEntry[fjarsýnn] \dicTerm{fjar··sýnn} \dicIPA{{f}{j}{a}{\textsubring{r}}{s}{i}{\textsubring{d}}{\textsubring{n}}} \dicPos{adj}[7]\dicFlx{}[-1] \dicDirectTranslationCS{dalekozraký} \dicAntonym{nærsýnn}
\dicEntry[fjarvera] \dicTerm{fjar··ver|a} \textls[15]{\dicIPA{{f}{j}{a}{r}{v}{\textepsilon}{r}{a}} \dicPos{f}[1] \dicFlx{(‑u)}[5] \dicDirectTranslationCS{nepřítomnost, absence}} \dicExampleIS{koma aftur til vinnu eftir langa fjarveru} \dicExampleCS{vrátit se zpátky do práce po dlouhé nepřítomnosti}
\dicEntry[fjarverandi] \dicTerm{fjar··ver·andi} \dicIPA{{f}{j}{a}{r}{v}{\textepsilon}{r}{a}{n}{\textsubring{d}}{\textsci}} \dicPos{adj}[13] \dicFlx{indecl}[1] \dicSynonym{fjarstaddur} \dicDirectTranslationCS{nepřítomný} \dicIndirectTranslationCS{(nevyskytující se na určitém místě)}
\dicEntry[fjarverufall] \dicTerm{fjar·veru··|fall} \dicIPA{{f}{j}{a}{r}{v}{\textepsilon}{r}{\textscy}{f}{a}{\textsubring{d}}{\textsubring{l}}} \dicPos{n}[2] \dicFlx{(‑falls, ‑föll)}[8] \dicFieldCat{jaz.} \dicDirectTranslationCS{abessiv}
\dicEntry[fjarvist] \dicTerm{fjar··vist} \dicIPA{{f}{j}{a}{r}{v}{\textsci}{s}{\textsubring{d}}} \dicPos{f}[7] \dicFlx{(‑ar, ‑ir)}[1] \dicSynonym{frávera} \dicDirectTranslationCS{nepřítomnost, absence} \dicExampleIS{langar fjarvistir frá fjölskyldunni} \dicExampleCS{dlouhá nepřítomnost v~rodině}
\dicEntry[fjarvistarsönnun] \dicTerm{fjar·vistar··|sönnun} \dicIPA{{f}{j}{a}{r}{v}{\textsci}{s}{\textsubring{d}}{a}{\textsubring{r}}{s}{\oe}{n}{\textscy}{\textsubring{n}}} \dicPos{f}[7] \dicFlx{(‑sönnunar, ‑sann\-an\-ir)}[11] \dicDirectTranslationCS{alibi}
\dicEntry[fjas] \dicTerm{fjas} \dicIPA{{f}{j}{a}{\textlengthmark}{s}} \dicPos{n}[2] \dicFlx{(‑s)}[2] \dicSynonym{mælgi} \dicDirectTranslationCS{brebentění, klábosení, tlachání}
\dicEntry[fjasa] \dicTerm{fjas|a} \dicIPA{{f}{j}{a}{\textlengthmark}{s}{a}} \dicPos{v}[1] \dicFlx{(‑aði)}[28] \dicSynonym{þusa} \dicDirectTranslationCS{vyvádět, remcat, povykovat} \dicExampleIS{fjasa um e‑ð} \dicExampleCS{vyvádět kvůli (čemu)}
\dicEntry[fjáður] \dicTerm{fjáður} \dicIPA{{f}{j}{au}{\textlengthmark}{ð}{\textscy}{\textsubring{r}}} \dicPos{adj}[2]\dicFlx{}[-12] \dicSynonym{efnaður} \dicDirectTranslationCS{zámožný, majetný}
\dicEntry[fjárdráttur] \dicTerm{fjár··drátt|ur} \dicIPA{{f}{j}{au}{r}{\textsubring{d}}{r}{au}{h}{\textsubring{d}}{\textscy}{\textsubring{r}}} \dicPos{m}[12] \dicFlx{(‑ar)}[8] \dicDirectTranslationCS{zpronevěra, defraudace}
\dicEntry[fjárfesta] \dicTerm{fjár··fest|a} \dicIPA{{f}{j}{au}{\textsubring{r}}{f}{\textepsilon}{s}{\textsubring{d}}{a}} \dicPos{v}[2] \dicFlx{(‑i, ‑)}[17] \dicDirectTranslationCS{investovat (peníze ap.)} \dicExampleIS{fjárfesta eina milljón í hlutabréfum} \dicExampleCS{investovat jeden milión do akcií}
\dicEntry[fjárfesting] \dicTerm{fjár··fest·ing} \dicIPA{{f}{j}{au}{\textsubring{r}}{f}{\textepsilon}{s}{\textsubring{d}}{i}{\ng}{\r{g}}} \dicPos{f}[4] \dicFlx{(‑ar, ‑ar)}[5] \dicFieldCat{ekon.} \dicDirectTranslationCS{investování, investice}
\dicEntry[fjárfestingarfélag] \dicTerm{fjár·festingar··fé·|lag} \dicIPA{{f}\-{j}\-{au}\-{\textsubring{r}}\-{f}\-{\textepsilon}\-{s}\-{\textsubring{d}}\-{i}\-{\ng}\-{\r{g}}\-{a}\-{\textsubring{r}}\-{f}\-{j}\-{\textepsilon}\-{l}\-{a}\-{x}\-} \dicPos{n}[2] \dicFlx{(‑lags, ‑lög)}[8] \dicFieldCat{ekon.} \dicDirectTranslationCS{investiční společnost}
\dicEntry[fjárfestingarsjóður] \dicTerm{fjár·festingar··sjóð|ur} \dicIPA{{f}\-{j}\-{au}\-{\textsubring{r}}\-{f}\-{\textepsilon}\-{s}\-{\textsubring{d}}\-{i}\-{\ng}\-{\r{g}}\-{a}\-{\textsubring{r}}\-{s}\-{j}\-{ou}\-{ð}\-{\textscy}\-{\textsubring{r}}\-} \dicPos{m}[9] \dicFlx{(‑s, ‑ir)}[6] \dicFieldCat{ekon.} \dicDirectTranslationCS{investiční fond}
\dicEntry[fjárglæfrar] \dicTerm{fjár··glæfrar} \dicIPA{{f}{j}{au}{r}{\r{g}}{l}{a}{i}{v}{r}{a}{\textsubring{r}}} \dicPos{m}[14] \dicFlx{pl}[1] \dicDirectTranslationCS{hospodářská kriminalita}
\dicEntry[fjárhagsaðstoð] \dicTerm{fjár·hags··að·stoð} \dicIPA{{f}{j}{au}{\textlengthmark}{\textsubring{r}}{h}{a}{x}{s}{a}{ð}{s}{\textsubring{d}}{\textopeno}{\texttheta}} \dicPos{f}[7] \dicFlx{(‑ar)}[3] \dicDirectTranslationCS{finanční pomoc\,/\addthin podpora}
\dicEntry[fjárhagsár] \dicTerm{fjár·hags··ár} \dicIPA{{f}{j}{au}{\textlengthmark}{\textsubring{r}}{h}{a}{x}{s}{au}{\textsubring{r}}} \dicPos{n}[2] \dicFlx{(‑s, ‑)}[5] \dicDirectTranslationCS{fiskální rok}
\dicEntry[fjárhagsáætlun] \dicTerm{fjár·hags··á·ætl|un} \dicIPA{{f}{j}{au}{\textlengthmark}{\textsubring{r}}{h}{a}{x}{s}{au}{a}{i}{h}{\textsubring{d}}{l}{\textscy}{\textsubring{n}}} \dicPos{f}[7] \dicFlx{(‑unar, ‑anir)}[8] \dicFieldCat{ekon.} \dicDirectTranslationCS{rozpočet}
\dicEntry[fjárhagslegur] \dicTerm{fjár·hags··legur} \dicIPA{{f}{j}{au}{\textlengthmark}{\textsubring{r}}{h}{a}{x}{s}{l}{\textepsilon}{\textbabygamma}{\textscy}{\textsubring{r}}} \dicPos{adj}[1]\dicFlx{}[-8] \dicDirectTranslationCS{finanční, peněžní}
\dicEntry[fjárhagsstaða] \dicTerm{fjár·hags··|staða} \dicIPA{{f}{j}{au}{\textlengthmark}{\textsubring{r}}{h}{a}{x}{s}{\textsubring{d}}{a}{ð}{a}} \dicPos{f}[1] \dicFlx{(‑stöðu)}[2] \dicDirectTranslationCS{finanční situace}
\dicEntry[fjárhagstímabil] \dicTerm{fjár·hags··tíma·bil} \dicIPA{{f}{j}{au}{\textlengthmark}{\textsubring{r}}{h}{a}{x}{s}{t\smash{\textsuperscript{h}}}{i}{m}{a}{\textsubring{b}}{\textsci}{\textsubring{l}}} \dicPos{n}[2] \dicFlx{(‑s, ‑)}[5] \dicFieldCat{ekon.} \dicDirectTranslationCS{rozpočtové období}
\dicEntry[fjárhagsyfirlit] \dicTerm{fjár·hags··yfir·lit} \dicIPA{{f}{j}{au}{\textlengthmark}{\textsubring{r}}{h}{a}{x}{s}{\textsci}{v}{\textsci}{r}{l}{\textsci}{\textsubring{d}}} \dicPos{n}[2] \dicFlx{(‑s, ‑)}[5] \dicDirectTranslationCS{finanční přehled\,/\addthin zpráva}
\dicEntry[fjárhagur] \dicTerm{fjár··hag|ur} \dicIPA{{f}{j}{au}{\textlengthmark}{\textsubring{r}}{h}{a}{\textbabygamma}{\textscy}{\textsubring{r}}} \dicPos{m}[9] \dicFlx{(‑s, ‑ir)}[9] \dicSynonym{fjármál} \dicDirectTranslationCS{stav financí, finance}
\begin{xtolerant}{}{1pt}
\dicEntry[fjárhaldsmaður] \dicTerm{fjár·halds··|maður}\addthinS\dicIPA{{f}{j}{au}{\textlengthmark}{\textsubring{r}}{h}{a}{l}{\textsubring{d}}{s}{m}{a}{ð}{\textscy}{\textsubring{r}}}\addthinS\dicPos{m}[13]\addthinS\dicFlx{(‑manns, ‑menn)}[2] \dicDirectTranslationCS{správce financí}
\end{xtolerant}
\dicEntry[fjárhundur] \dicTerm{fjár··hund|ur} \dicIPA{{f}{j}{au}{\textlengthmark}{\textsubring{r}}{h}{\textscy}{n}{\textsubring{d}}{\textscy}{\textsubring{r}}} \dicPos{m}[6] \dicFlx{(‑s, ‑ar)}[4] \dicDirectTranslationCS{ovčák (pes), ovčácký pes};  \dicPhraseIS{þýskur fjárhundur} \dicDirectTranslationCS{německý ovčák}
\dicEntry[fjárhús] \dicTerm{fjár··hús} \dicIPA{{f}{j}{au}{\textlengthmark}{\textsubring{r}}{h}{u}{s}} \dicPos{n}[2] \dicFlx{(‑s, ‑)}[5] \dicDirectTranslationCS{ovčín, salaš}
\dicEntry[fjárhæð] \dicTerm{fjár··hæð} \dicIPA{{f}{j}{au}{\textlengthmark}{\textsubring{r}}{h}{a}{i}{\texttheta}} \dicPos{f}[7] \dicFlx{(‑ar, ‑ir)}[1] \dicDirectTranslationCS{peněžní částka}
\dicEntry[fjári] \dicTerm{fjár|i\smash{\textsuperscript{1}}} \dicIPA{{f}{j}{au}{\textlengthmark}{r}{\textsci}} \dicPos{m}[1] \dicFlx{(‑a, ‑ar)}[1] \dicSynonym{fjandi\smash{\textsuperscript{1}}} \dicDirectTranslationCS{čert, ďábel};  \dicPhraseIS{fjárinn hafi það} \dicDirectTranslationCS{vem to čert};  \dicPhraseIS{hver fjárinn} \dicDirectTranslationCS{sakra, zatraceně}
\dicEntry[fjári] \dicTerm{fjári\smash{\textsuperscript{2}}} \dicIPA{{f}{j}{au}{\textlengthmark}{r}{\textsci}} \dicPos{adv} \dicDirectTranslationCS{zatraceně, sakramentsky} \dicExampleIS{fjári góður kennari} \dicExampleCS{zatraceně dobrý učitel}
\dicEntry[fjárlagahalli] \dicTerm{fjár·laga··hall|i} \dicIPA{{f}{j}{au}{r}{l}{a}{\textbabygamma}{a}{h}{a}{\textsubring{d}}{l}{\textsci}} \dicPos{m}[1] \dicFlx{(‑a)}[3] \dicFieldCat{ekon.} \dicDirectTranslationCS{rozpočtový schodek}
\dicEntry[fjárlög] \dicTerm{fjár··lög} \dicIPA{{f}{j}{au}{r}{l}{\oe}{x}} \dicPos{n}[2] \dicFlx{pl}[9] \dicFieldCat{ekon.} \dicDirectTranslationCS{státní rozpočet}
\dicEntry[fjármagn] \dicTerm{fjár··magn} \dicIPA{{f}{j}{au}{r}{m}{a}{\r{g}}{\textsubring{n}}} \dicPos{n}[2] \dicFlx{(‑s)}[2] \dicFieldCat{ekon.} \dicDirectTranslationCS{kapitál, finanční zdroje}
\dicEntry[fjármagna] \dicTerm{fjár··magn|a} \dicIPA{{f}{j}{au}{r}{m}{a}{\r{g}}{n}{a}} \dicPos{v}[1] \dicFlx{(‑aði)}[13] \dicFlx{acc} \dicDirectTranslationCS{financovat} \dicExampleIS{fjármagna framkvæmdina} \dicExampleCS{financovat realizaci}
\begin{xtolerant}{}{1pt}
\dicEntry[fjármagnsfrekur] \dicTerm{fjár·magns··frekur}\addthinS\textls[-10]{\dicIPA{{f}{j}{au}{r}{m}{a}{\r{g}}{\textsubring{n}}{s}{f}{r}{\textepsilon}{\r{g}}{\textscy}{\textsubring{r}}}} \dicPos{adj}[1]\dicFlx{}[-1]\addthinS\dicDirectTranslationCS{kapitálově\hspace{-0.10667em}\,/\addthin finančně náročný}
\end{xtolerant}
\dicEntry[fjármagnshreyfing] \dicTerm{fjár·magns··hreyf·ing} \dicIPA{{f}\-{j}\-{au}\-{r}\-{m}\-{a}\-{\r{g}}\-{\textsubring{n}}\-{s}\-{\textsubring{r}}\-{ei}\-{v}\-{i}\-{\ng}\-{\r{g}}\-} \dicPos{f}[4] \dicFlx{(‑ar, ‑ar)}[5] \dicFieldCat{ekon.} \dicDirectTranslationCS{kapitálová transakce}
\dicEntry[fjármagnsjöfnuður] \dicTerm{fjár·magns··|jöfn·uður} \textls[-10]{\dicIPA{{f}\-{j}\-{au}\-{r}\-{m}\-{a}\-{\r{g}}\-{\textsubring{n}}\-{s}\-{j}\-{\oe}\-{\textsubring{b}}\-{n}\-{\textscy}\-{ð}\-{\textscy}\-{\textsubring{r}}}} \dicPos{m}[10] \dicFlx{(‑jafn\-aðar\,/\addthin ‑jöfnuðar)}[43] \dicFieldCat{ekon.} \dicDirectTranslationCS{kapitálový účet}
\dicEntry[fjármagnskostnaður] \dicTerm{fjár·magns··kost·nað|ur} \dicIPA{{f}\-{j}\-{au}\-{r}\-{m}\-{a}\-{\r{g}}\-{\textsubring{n}}\-{s}\-{k\smash{\textsuperscript{h}}}\-{\textopeno}\-{s}\-{\textsubring{d}}\-{n}\-{a}\-{ð}\-{\textscy}\-{\textsubring{r}}\-} \dicPos{m}[10] \dicFlx{(‑ar)}[9] \dicFieldCat{ekon.} \dicDirectTranslationCS{kapitálové\,/\addthin finanční náklady}
\dicEntry[fjármagnsmarkaður] \dicTerm{fjár·magns··mark·að|ur} \dicIPA{{f}\-{j}\-{au}\-{r}\-{m}\-{a}\-{\r{g}}\-{\textsubring{n}}\-{s}\-{m}\-{a}\-{\textsubring{r}}\-{\r{g}}\-{a}\-{ð}\-{\textscy}\-{\textsubring{r}}\-} \dicPos{m}[10] \dicFlx{(‑ar\,/\addthin ‑s, ‑ir)}[31] \dicFieldCat{ekon.} \dicDirectTranslationCS{kapitálový\,/\addthin finanční trh}
\dicEntry[fjármagnsstreymi] \dicTerm{fjár·magns··streymi} \dicIPA{{f}{j}{au}{r}{m}{a}{\r{g}}{\textsubring{n}}{s}{\textsubring{d}}{r}{ei}{m}{\textsci}} \dicPos{n}[2] \dicFlx{(‑s, ‑)}[14] \dicFieldCat{ekon.} \dicDirectTranslationCS{tok kapitálu, kapitálový tok}
\dicEntry[fjármagnsþörf] \dicTerm{fjár·magns··|þörf} \dicIPA{{f}{j}{au}{r}{m}{a}{\r{g}}{\textsubring{n}}{s}{\texttheta}{\oe}{r}{f}} \dicPos{f}[7] \dicFlx{(‑þarfar, ‑þarfir)}[16] \dicFieldCat{ekon.} \dicDirectTranslationCS{finanční potřeby}
\dicEntry[fjármál] \dicTerm{fjár··mál} \dicIPA{{f}{j}{au}{r}{m}{au}{\textsubring{l}}} \dicPos{n}[2] \dicFlx{pl}[1] \dicFieldCat{ekon.} \dicSynonym*{peningamál} \dicDirectTranslationCS{finance, finanční záležitosti}
\dicEntry[fjármálafyrirtæki] \dicTerm{fjár·mála··fyrir·tæki} \dicIPA{{f}\-{j}\-{au}\-{r}\-{m}\-{au}\-{l}\-{a}\-{f}\-{\textsci}\-{r}\-{\textsci}\-{\textsubring{r}}\-{t\smash{\textsuperscript{h}}}\-{a}\-{i}\-{\r{\textObardotlessj}}\-{\textsci}\-} \dicPos{n}[2] \dicFlx{(‑s, ‑)}[16] \dicDirectTranslationCS{finanční instituce}
\dicEntry[fjármálamaður] \dicTerm{fjár·mála··|maður} \dicIPA{{f}{j}{au}{r}{m}{au}{l}{a}{m}{a}{ð}{\textscy}{\textsubring{r}}} \dicPos{m}[13] \dicFlx{(‑manns, ‑menn)}[2] \dicDirectTranslationCS{finančník, finančnice}
\dicEntry[fjármálaráðherra] \dicTerm{fjár·mála··ráð·herr|a} \dicIPA{{f}\-{j}\-{au}\-{r}\-{m}\-{au}\-{l}\-{a}\-{r}\-{au}\-{\texttheta}\-{h}\-{\textepsilon}\-{r}\-{a}\-} \dicPos{m}[1] \dicFlx{(‑a, ‑ar)}[17] \dicDirectTranslationCS{ministr(yně) financí}
\dicEntry[fjármálaráðuneyti] \dicTerm{fjár·mála··ráðu·neyti} \dicIPA{{f}\-{j}\-{au}\-{r}\-{m}\-{au}\-{l}\-{a}\-{r}\-{au}\-{ð}\-{\textscy}\-{n}\-{ei}\-{\textsubring{d}}\-{\textsci}\-} \dicPos{n}[2] \dicFlx{(‑s, ‑)}[14] \dicDirectTranslationCS{ministerstvo financí}
\dicEntry[fjármálastarfsemi] \dicTerm{fjár·mála··starf·sem|i} \dicIPA{{f}\-{j}\-{au}\-{r}\-{m}\-{au}\-{l}\-{a}\-{s}\-{\textsubring{d}}\-{a}\-{r}\-{v}\-{s}\-{\textepsilon}\-{m}\-{\textsci}\-} \dicPos{f}[2] \dicFlx{(‑i)}[2] \dicFieldCat{ekon.} \dicDirectTranslationCS{finančnictví}
\dicEntry[fjármálastefna] \dicTerm{fjár·mála··stefn|a} \dicIPA{{f}{j}{au}{r}{m}{au}{l}{a}{s}{\textsubring{d}}{\textepsilon}{\textsubring{b}}{n}{a}} \dicPos{f}[1] \dicFlx{(‑u)}[5] \dicFieldCat{ekon.} \dicDirectTranslationCS{fiskální politika}
\dicEntry[fjármálastjóri] \dicTerm{fjár·mála··stjór|i} \dicIPA{{f}{j}{au}{r}{m}{au}{l}{a}{s}{\textsubring{d}}{j}{ou}{r}{\textsci}} \dicPos{m}[1] \dicFlx{(‑a, ‑ar)}[1] \dicDirectTranslationCS{hlavní ekonom(ka)}
\dicEntry[fjármunamyndun] \dicTerm{fjár·muna··mynd|un} \dicIPA{{f}{j}{au}{r}{m}{\textscy}{n}{a}{m}{\textsci}{n}{\textsubring{d}}{\textscy}{\textsubring{n}}} \dicPos{f}[7] \dicFlx{(‑unar)}[9] \dicFieldCat{ekon.} \dicDirectTranslationCS{tvorba kapitálu\,/\addthin aktiv}
\dicEntry[fjármunir] \dicTerm{fjár··munir} \dicsymFrequent\  \dicIPA{{f}{j}{au}{r}{m}{\textscy}{n}{\textsci}{\textsubring{r}}} \dicPos{m}[10] \dicFlx{pl}[1] \dicFieldCat{ekon.} \dicDirectTranslationCS{finanční prostředky, finance} \dicExampleIS{Hún gaf honum mikla fjármuni.} \dicExampleCS{Věnovala mu velké finanční prostředky.}
\dicEntry[fjármögnun] \dicTerm{fjár··mögn|un} \dicIPA{{f}{j}{au}{r}{m}{\oe}{\r{g}}{n}{\textscy}{\textsubring{n}}} \dicPos{f}[7] \dicFlx{(‑unar)}[12] \dicFieldCat{ekon.} \dicDirectTranslationCS{financování}
\dicEntry[fjármögnunarfélag] \dicTerm{fjár·mögnunar··fé·|lag} \dicIPA{{f}\-{j}\-{au}\-{r}\-{m}\-{\oe}\-{\r{g}}\-{n}\-{\textscy}\-{n}\-{a}\-{\textsubring{r}}\-{f}\-{j}\-{\textepsilon}\-{l}\-{a}\-{x}\-} \dicPos{n}[2] \dicFlx{(‑lags, ‑lög)}[8] \dicDirectTranslationCS{finanční společnost\,/\addthin dům}
\dicEntry[fjárnám] \dicTerm{fjár··nám} \dicIPA{{f}{j}{au}{r}{n}{au}{\textsubring{m}}} \dicPos{n}[2] \dicFlx{(‑s)}[2] \dicFieldCat{práv.} \dicDirectTranslationCS{exekuce, zabavení}
\dicEntry[fjárráð] \dicTerm{fjár··ráð} \dicIPA{{f}{j}{au}{r}{au}{\texttheta}} \dicPos{n}[2] \dicFlx{pl}[1] \textbf{1.} \dicSynonym*{forræði yfir eignum} \dicDirectTranslationCS{disponování financemi}  \textbf{2.} \dicSynonym*{peningaeign} \dicDirectTranslationCS{finanční prostředky}
\dicEntry[fjársjóður] \dicTerm{fjár··sjóð|ur} \dicIPA{{f}{j}{au}{\textsubring{r}}{s}{j}{ou}{ð}{\textscy}{\textsubring{r}}} \dicPos{m}[9] \dicFlx{(‑s, ‑ir)}[6] \dicDirectTranslationCS{poklad} \dicExampleIS{finna fjársjóð} \dicExampleCS{nalézt poklad}
\dicEntry[fjárstyrkur] \dicTerm{fjár··styrk|ur} \dicIPA{{f}{j}{au}{\textsubring{r}}{s}{\textsubring{d}}{\textsci}{\textsubring{r}}{\r{g}}{\textscy}{\textsubring{r}}} \dicPos{m}[9] \dicFlx{(‑s, ‑ir)}[15] \dicDirectTranslationCS{finanční podpora}
\dicEntry[fjársvik] \dicTerm{fjár··svik} \dicIPA{{f}{j}{au}{\textsubring{r}}{s}{v}{\textsci}{\r{g}}} \dicPos{n}[2] \dicFlx{pl}[1] \dicFieldCat{práv.} \dicDirectTranslationCS{finanční podvod}
\dicEntry[fjárveiting] \dicTerm{fjár··veit·ing} \dicIPA{{f}{j}{au}{r}{v}{ei}{\textsubring{d}}{i}{\ng}{\r{g}}} \dicPos{f}[4] \dicFlx{(‑ar, ‑ar)}[5] \dicDirectTranslationCS{dotace, grant}
\dicEntry[fjáröflun] \dicTerm{fjár··öfl|un} \dicIPA{{f}{j}{au}{\textlengthmark}{r}{\oe}{\textsubring{b}}{l}{\textscy}{\textsubring{n}}} \dicPos{f}[7] \dicFlx{(‑unar)}[12] \textbf{1.} \dicSynonym*{útvegun fjár} \dicDirectTranslationCS{financování}  \textbf{2.} \dicSynonym*{fjársöfnun} \dicDirectTranslationCS{finanční sbírka}
\dicEntry[fjóla] \dicTerm{fjól|a} \dicIPA{{f}{j}{ou}{\textlengthmark}{l}{a}} \dicPos{f}[1] \dicFlx{(‑u, ‑ur)}[19] \dicFieldCat{bot.} \dicDirectTranslationCS{violka, fialka} \textit{(l.~{\textLA{Viola}})}  \dicsymPhoto\ 
\dicFigure{1234.jpg}{Fjóla}{Fjóla - Zicha Ondřej, Biolib, Copyright/CC-BY-NC}
\dicEntry[fjólublár] \dicTerm{fjólu··blár} \dicIPA{{f}{j}{ou}{\textlengthmark}{l}{\textscy}{\textsubring{b}}{l}{au}{\textsubring{r}}} \dicPos{adj}[4]\dicFlx{}[-1] \dicDirectTranslationCS{fialový}
\dicEntry[Fjón] \dicTerm{Fjón} \dicIPA{{f}{j}{ou}{\textlengthmark}{\textsubring{n}}} \dicPos{n}[2] \dicFlx{(‑s)}[4] \dicFieldCat{geog.} \dicDirectTranslationCS{Fyn} \dicIndirectTranslationCS{(ostrov v~Dánsku)}
\dicEntry[fjórar] \dicTerm{fjórar} \dicIPA{{f}{j}{ou}{\textlengthmark}{r}{a}{\textsubring{r}}} \dicPos{num} \dicFlx{pl} \dicFlx{f nom\,/\addthin acc} \dicLink{fjórir}
\dicEntry[fjórði] \dicTerm{fjórði} \dicsymFrequent\  \dicIPA{{f}{j}{ou}{r}{ð}{\textsci}} \dicPos{num}[5] \dicFlx{ord} \dicDirectTranslationCS{čtvrtý} \dicExampleIS{á fjórðu hæðinni} \dicExampleCS{ve čtvrtém patře}
\dicEntry[fjórðungur] \dicTerm{fjórð··ung|ur} \dicsymFrequent\  \dicIPA{{f}{j}{ou}{r}{ð}{u}{\ng}{\r{g}}{\textscy}{\textsubring{r}}} \dicPos{m}[6] \dicFlx{(‑s, ‑ar)}[8] \textbf{1.} \dicDirectTranslationCS{čtvrtina, čtvrtka} \dicExampleIS{fækka um fjórðung} \dicExampleCS{snížit o~čtvrtinu}  \textbf{2.} \dicLangCat{kniž.} \dicSynonym{korter} \dicDirectTranslationCS{čtvrt}
\dicEntry[fjórfaldur] \dicTerm{fjór··|faldur} \dicIPA{{f}{j}{ou}{\textsubring{r}}{f}{a}{l}{\textsubring{d}}{\textscy}{\textsubring{r}}} \dicPos{adj}[2] \dicFlx{(f ‑föld)}[16] \textbf{1.} \dicDirectTranslationCS{čtyřnásobný}  \textbf{2.} \dicDirectTranslationCS{čtyřdílný}
\dicEntry[fjórhjól] \dicTerm{fjór··hjól} \dicIPA{{f}{j}{ou}{\textsubring{r}}{\c{c}}{ou}{\textsubring{l}}} \dicPos{n}[2] \dicFlx{(‑s, ‑)}[5] \dicDirectTranslationCS{čtyřkolka}
\dicEntry[fjórir] \dicTerm{fjórir} \dicsymFrequent\  \dicIPA{{f}{j}{ou}{\textlengthmark}{r}{\textsci}{\textsubring{r}}} \dicPos{num}[4] \dicFlx{pl} \dicFlx{m} \dicDirectTranslationCS{čtyři} \dicExampleIS{eftir fjóra daga} \dicExampleCS{po čtyřech dnech}
\dicEntry[fjórtán] \dicTerm{fjór··tán} \dicsymFrequent\  \dicIPA{{f}{j}{ou}{\textsubring{r}}{\textsubring{d}}{au}{\textsubring{n}}} \dicPos{num} \dicFlx{indecl} \dicDirectTranslationCS{čtrnáct} \dicExampleIS{Tólf til fjórtán manns bjuggu saman í litlu húsi.} \dicExampleCS{V~malém domě bydlelo spolu dvanáct až čtrnáct lidí.}
\dicEntry[fjórtándi] \dicTerm{fjór··tándi} \dicIPA{{f}{j}{ou}{\textsubring{r}}{\textsubring{d}}{au}{n}{\textsubring{d}}{\textsci}} \dicPos{num}[5] \dicFlx{ord} \dicDirectTranslationCS{čtrnáctý}
\dicEntry[fjós] \dicTerm{fjós} \dicsymFrequent\  \dicIPA{{f}{j}{ou}{\textlengthmark}{s}} \dicPos{n}[2] \dicFlx{(‑s, ‑)}[5] \dicDirectTranslationCS{kravín, chlév} \dicExampleIS{Á veturna eru kýrnar í fjósinu.} \dicExampleCS{Během zimy jsou krávy v~kravíně.}
\dicEntry[fjúka] \dicTerm{fjúka} \dicsymFrequent\  \dicIPA{{f}{j}{u}{\textlengthmark}{\r{g}}{a}} \dicPos{v}[6] \dicFlx{(fýk, fauk, fukum, fyki, fokið)}[101] \dicSynonym{blása} \dicDirectTranslationCS{být unášen (větrem ap.), být odfouknut (klobouk ap.)} \dicExampleIS{Laufið fýkur.} \dicExampleCS{List je unášen větrem.};  \dicPhraseIS{láta e‑ð fjúka} \dicLangCat{přen.} \dicDirectTranslationCS{říct (co) bez obalu, vysypat (co) ze sebe};  \dicPhraseIS{láta e‑n fjúka} \dicDirectTranslationCS{vykopnout (koho), vyhodit (koho) (z~práce ap.)};  \dicPhraseIS{það fýkur} \dicFlx{impers} \dicSynonym*{það snjóar} \dicDirectTranslationCS{sněží};  \dicPhraseIS{það fýkur í e‑n} \dicFlx{impers} \dicDirectTranslationCS{(kdo) se zlobí, (kdo) se rozčiluje} \dicExampleIS{Það fauk í hana.} \dicExampleCS{Rozzlobila se.}
\dicEntry[fjær] \dicTerm{fjær} \dicIPA{{f}{j}{a}{i}{\textlengthmark}{\textsubring{r}}} \dicPos{adv} \dicFlx{comp} \dicLink{fjarri}
\dicEntry[fjærst] \dicTerm{fjærst} \dicIPA{{f}{j}{a}{i}{\textsubring{r}}{s}{\textsubring{d}}} \dicPos{adv} \dicFlx{sup} \dicLink{fjarri}
\dicEntry[fjærstur] \dicTerm{fjærstur} \dicIPA{{f}{j}{a}{i}{\textsubring{r}}{s}{\textsubring{d}}{\textscy}{\textsubring{r}}} \dicPos{adj}[12]\dicFlx{}[-12] \dicFlx{sup} \dicDirectTranslationCS{nejvzdálenější, nejodlehlejší}
\dicEntry[fjöður] \dicTerm{fjöður} \dicsymFrequent\  \dicIPA{{f}{j}{\oe}{\textlengthmark}{ð}{\textscy}{\textsubring{r}}} \dicPos{f}[7] \dicFlx{(fjaðrar, fjaðrir)}[16] \textbf{1.} \dicDirectTranslationCS{pero, pírko}  \textbf{2.} \dicSynonym{gormur} \dicDirectTranslationCS{pružina, péro};  \dicPhraseIS{draga fjöður yfir e‑ð} \dicLangCat{přen.} \dicDirectTranslationCS{přivřít oči nad (čím)};  \dicPhraseIS{skreyta sig með stolnum fjöðrum} \dicLangCat{přen.} \dicDirectTranslationCS{chlubit se cizím peřím};  \dicPhraseIS{vera\,/\addthin verða fjöðrum fenginn} \dicLangCat{přen.} \dicDirectTranslationCS{být štěstím bez sebe}
\dicEntry[fjögur] \dicTerm{fjögur} \dicIPA{{f}{j}{\oe}{\textlengthmark}{\textbabygamma}{\textscy}{\textsubring{r}}} \dicPos{num} \dicFlx{pl} \dicFlx{n nom\,/\addthin acc} \dicLink{fjórir}
\dicEntry[fjöl] \dicTerm{fjöl\smash{\textsuperscript{1}}} \dicIPA{{f}{j}{\oe}{\textlengthmark}{\textsubring{l}}} \dicPos{f}[7] \dicFlx{(fjalar, fjalir)}[16] \dicDirectTranslationCS{prkno, deska};  \dicPhraseIS{vera ekki við eina fjölina felldur} \dicLangCat{přen.} \dicDirectTranslationCS{mít více želízek v~ohni}
\dicEntry[fjöl] \dicTerm{fjöl-\smash{\textsuperscript{2}}} \dicIPA{{f}{j}{\oe}{l}} \dicPos{predp} \dicDirectTranslationCS{více-, poly-, multi-, různo-}
\dicEntry[fjölbrautaskóli] \textls[10]{\dicTerm{fjöl·brauta··skól|i} \dicIPA{{f}{j}{\oe}{l}{\textsubring{b}}{r}{\oe i}{\textsubring{d}}{a}{s}{\r{g}}{ou}{l}{\textsci}} \dicPos{m}[1] \dicFlx{(‑a, ‑ar)}[1] \dicFieldCat{škol.} \dicDirectTranslationCS{střed\-ní odborná škola}}
\dicEntry[fjölbreytilegur] \dicTerm{fjöl·breyti··legur} \dicIPA{{f}{j}{\oe}{l}{\textsubring{b}}{r}{ei}{\textsubring{d}}{\textsci}{l}{\textepsilon}{\textbabygamma}{\textscy}{\textsubring{r}}} \dicPos{adj}[1]\dicFlx{}[-8] \dicDirectTranslationCS{rozmanitý, různorodý, pestrý}
\dicEntry[fjölbreytni] \dicTerm{fjöl··breytn|i} \dicIPA{{f}{j}{\oe}{l}{\textsubring{b}}{r}{ei}{h}{\textsubring{d}}{n}{\textsci}} \dicPos{f}[3] \dicFlx{(‑i)}[3] \dicDirectTranslationCS{rozmanitost, různorodost, rozličnost}
\dicEntry[fjölbreyttur] \dicTerm{fjöl··breyttur} \dicIPA{{f}{j}{\oe}{l}{\textsubring{b}}{r}{ei}{h}{\textsubring{d}}{\textscy}{\textsubring{r}}} \dicPos{adj}[1]\dicFlx{}[-10] \dicDirectTranslationCS{různorodý, rozmanitý}
\dicEntry[fjölbýlishús] \dicTerm{fjöl·býlis··hús} \dicIPA{{f}{j}{\oe}{l}{\textsubring{b}}{i}{l}{\textsci}{s}{h}{u}{s}} \dicPos{n}[2] \dicFlx{(‑s, ‑)}[5] \dicDirectTranslationCS{činžovní\,/\addthin panelový dům, činžák, panelák}
\dicEntry[fjöldaframleiðsla] \dicTerm{fjölda··fram·leiðsl|a} \dicIPA{{f}\-{j}\-{\oe}\-{l}\-{\textsubring{d}}\-{a}\-{f}\-{r}\-{a}\-{m}\-{l}\-{ei}\-{ð}\-{s}\-{\textsubring{d}}\-{l}\-{a}\-} \dicPos{f}[1] \dicFlx{(‑u)}[5] \dicDirectTranslationCS{velkovýroba, masová produkce}
\dicEntry[fjöldamorðingi] \dicTerm{fjölda··morð·ing|i} \dicIPA{{f}{j}{\oe}{l}{\textsubring{d}}{a}{m}{\textopeno}{r}{ð}{i}{\textltailn}{\r{\textObardotlessj}}{\textsci}} \dicPos{m}[1] \dicFlx{(‑ja, ‑jar)}[14] \dicDirectTranslationCS{masový vrah, masová vražedkyně}
\dicEntry[fjöldi] \dicTerm{fjöld|i} \dicsymFrequent\  \dicIPA{{f}{j}{\oe}{l}{\textsubring{d}}{\textsci}} \dicPos{m}[1] \dicFlx{(‑a)}[3] \dicDirectTranslationCS{počet, množství (lidí ap.)} \dicExampleIS{fjöldi fólks} \dicExampleCS{množství lidí}
\dicEntry[fjölfalda] \dicTerm{fjöl··fald|a} \dicIPA{{f}{j}{\oe}{l}{f}{a}{l}{\textsubring{d}}{a}} \dicPos{v}[1] \dicFlx{(‑aði)}[13] \dicFlx{acc} \dicDirectTranslationCS{(roz)množit, (z)kopírovat} \dicExampleIS{fjölfalda bréfið} \dicExampleCS{rozmnožit dopis}
\dicEntry[fjölfarinn] \dicTerm{fjöl··farinn} \dicIPA{{f}{j}{\oe}{l}{f}{a}{r}{\textsci}{\textsubring{n}}} \dicPos{adj}[6]\dicFlx{}[-3] \dicDirectTranslationCS{rušný, frekventovaný (ulice ap.)} \dicExampleIS{fjölfarin gata} \dicExampleCS{rušná ulice}
\dicEntry[fjölfatlaður] \dicTerm{fjöl··|fatl·aður} \dicIPA{{f}{j}{\oe}{l}{f}{a}{h}{\textsubring{d}}{l}{a}{ð}{\textscy}{\textsubring{r}}} \dicPos{adj}[3] \dicFlx{(f ‑fötluð)}[2] \dicDirectTranslationCS{vícenásobně postižený\,/\addthin handicapovaný}
\dicEntry[fjölfræðibók] \dicTerm{fjöl·fræði··|bók} \dicIPA{{f}{j}{\oe}{l}{f}{r}{a}{i}{ð}{\textsci}{\textsubring{b}}{ou}{\r{g}}} \dicPos{f}[8] \dicFlx{(‑bókar, ‑bækur)}[5] \dicDirectTranslationCS{encyklopedie, lexikon}
\dicEntry[fjölföldun] \dicTerm{fjöl··föld|un} \dicIPA{{f}{j}{\oe}{l}{f}{\oe}{l}{\textsubring{d}}{\textscy}{\textsubring{n}}} \dicPos{f}[7] \dicFlx{(‑unar)}[12] \dicDirectTranslationCS{kopírování, rozmnožování} \dicExampleIS{Öll fjölföldun er óheimil.} \dicExampleCS{Veškeré kopírování je nelegální.}
\dicEntry[fjölga] \dicTerm{fjölg|a} \dicsymFrequent\  \dicIPA{{f}{j}{\oe}{l}{\r{g}}{a}} \dicPos{v}[1] \dicFlx{(‑aði)}[44] \dicFlx{dat} \textbf{1.} \dicDirectTranslationCS{(roz)množit, zvýšit počet};  \dicPhraseIS{fjölga sér} \dicDirectTranslationCS{množit se, rozmnožovat se}  \textbf{2.} \dicPhraseIS{e‑u fjölgar} \dicFlx{impers} \dicDirectTranslationCS{(co) se rozmnožuje, (čeho) přibývá}
\dicEntry[fjölgervi] \dicTerm{fjöl··gervi} \dicIPA{{f}{j}{\oe}{l}{\r{\textObardotlessj}}{\textepsilon}{r}{v}{\textsci}} \dicPos{n}[2] \dicFlx{(‑s)}[20] \dicFieldCat{chem.} \dicDirectTranslationCS{alotropie}
\dicEntry[fjölgun] \dicTerm{fjölg|un} \dicIPA{{f}{j}{\oe}{l}{\r{g}}{\textscy}{\textsubring{n}}} \dicPos{f}[7] \dicFlx{(‑unar, ‑anir)}[8] \dicDirectTranslationCS{nárůst, přírůstek, zvýšení počtu} \dicExampleIS{fjölgun á e‑u} \dicExampleCS{přírůstek (čeho)}
\dicEntry[fjölgyðistrú] \dicTerm{fjöl·gyðis··trú} \dicIPA{{f}{j}{\oe}{l}{\r{\textObardotlessj}}{\textsci}{ð}{\textsci}{s}{t\smash{\textsuperscript{h}}}{r}{u}} \dicPos{f}[4] \dicFlx{(‑ar)}[27] \dicFieldCat{náb.} \dicDirectTranslationCS{mnohobožství, polyteismus}
\dicEntry[fjölhyggja] \dicTerm{fjöl··hyggj|a} \dicIPA{{f}{j}{\oe}{\textlengthmark}{\textsubring{l}}{h}{\textsci}{\r{\textObardotlessj}}{a}} \dicPos{f}[1] \dicFlx{(‑u)}[5] \dicFieldCat{filos.} \dicDirectTranslationCS{pluralismus}
\dicEntry[fjölhæfni] \dicTerm{fjöl··hæfn|i} \dicIPA{{f}{j}{\oe}{\textlengthmark}{\textsubring{l}}{h}{a}{i}{\textsubring{b}}{n}{\textsci}} \dicPos{f}[3] \dicFlx{(‑i)}[3] \dicDirectTranslationCS{všestrannost}
\dicEntry[fjölhæfur] \dicTerm{fjöl··hæfur} \dicIPA{{f}{j}{\oe}{\textlengthmark}{\textsubring{l}}{h}{a}{i}{v}{\textscy}{\textsubring{r}}} \dicPos{adj}[1]\dicFlx{}[-1] \dicSynonym{alhliða} \dicDirectTranslationCS{všestranný, všestranně nadaný}
\begin{xtolerant}{}{1pt}
\dicEntry[fjölkunnugur] \dicTerm{fjöl··kunn·ugur}\addthinS\dicIPA{{f}{j}{\oe}{l}{k\smash{\textsuperscript{h}}}{\textscy}{n}{\textscy}{\textbabygamma}{\textscy}{\textsubring{r}}}\addthinS\dicPos{adj}[1]\dicFlx{}[-8]\addthinS\dicSynonym{göldróttur}\addthinS\dicDirectTranslationCS{zběh\-lý v~magii, znalý magie\,/\addthin kouzel}
\end{xtolerant}
\dicEntry[fjölkvæni] \dicTerm{fjöl··kvæni} \dicIPA{{f}{j}{\oe}{l}{k\smash{\textsuperscript{h}}}{v}{a}{i}{n}{\textsci}} \dicPos{n}[2] \dicFlx{(‑s)}[20] \dicDirectTranslationCS{mnohoženství, polygynie}
\dicEntry[fjölkynngi] \dicTerm{fjöl··kynng|i} \dicIPA{{f}{j}{\oe}{l}{c\smash{\textsuperscript{h}}}{\textsci}{\ng}{\r{\textObardotlessj}}{\textsci}} \dicPos{f}[3] \dicFlx{(‑i)}[3] \dicSynonym*{galdrar} \dicDirectTranslationCS{čarodějnictví, čarování}
\dicEntry[fjöll] \dicTerm{fjöll} \dicIPA{{f}{j}{\oe}{\textsubring{d}}{\textsubring{l}}} \dicPos{n} \dicFlx{pl nom} \dicLink{fjall}
\dicEntry[fjölleikahús] \dicTerm{fjöl·leika··hús} \dicIPA{{f}{j}{\oe}{l}{\textlengthmark}{ei}{\r{g}}{a}{h}{u}{s}} \dicPos{n}[2] \dicFlx{(‑s, ‑)}[5] \dicSynonym{sirkus} \dicDirectTranslationCS{cirkus}
\dicEntry[fjölliða] \dicTerm{fjöl··lið|a} \dicIPA{{f}{j}{\oe}{l}{\textlengthmark}{\textsci}{ð}{a}} \dicPos{f}[1] \dicFlx{(‑u, ‑ur)}[19] \dicFieldCat{chem.} \dicDirectTranslationCS{polymer}
\dicEntry[fjöllyndi] \dicTerm{fjöl··lyndi} \dicIPA{{f}{j}{\oe}{l}{\textlengthmark}{\textsci}{n}{\textsubring{d}}{\textsci}} \dicPos{n}[2] \dicFlx{(‑s)}[20] \dicDirectTranslationCS{promiskuita}
\dicEntry[fjöllyndur] \dicTerm{fjöl··lyndur} \dicIPA{{f}{j}{\oe}{l}{\textlengthmark}{\textsci}{n}{\textsubring{d}}{\textscy}{\textsubring{r}}} \dicPos{adj}[2]\dicFlx{}[-14] \dicDirectTranslationCS{promiskuitní}
\dicEntry[fjölmargur] \dicTerm{fjöl··|margur} \dicsymFrequent\  \dicIPA{{f}{j}{\oe}{l}{m}{a}{r}{\r{g}}{\textscy}{\textsubring{r}}} \dicPos{adj}[11] \dicFlx{(f ‑mörg)}[9] \dicDirectTranslationCS{početný, četný} \dicExampleIS{fjölmörg dæmi um notkun orðsins} \dicExampleCS{četné příklady použití slova}
\dicEntry[fjölmenna] \dicTerm{fjöl··menn|a} \dicIPA{{f}{j}{\oe}{l}{m}{\textepsilon}{n}{a}} \dicPos{v}[2] \dicFlx{(‑ti, ‑t)}[78] \dicDirectTranslationCS{shromáždit se (v~hojném počtu)}
\dicEntry[fjölmenni] \dicTerm{fjöl··menni} \dicIPA{{f}{j}{\oe}{l}{m}{\textepsilon}{n}{\textsci}} \dicPos{n}[2] \dicFlx{(‑s)}[20] \dicDirectTranslationCS{dav, zástup, množství lidí}
\dicEntry[fjölmenning] \dicTerm{fjöl··menn·ing} \dicIPA{{f}{j}{\oe}{l}{m}{\textepsilon}{n}{i}{\ng}{\r{g}}} \dicPos{f}[4] \dicFlx{(‑ar)}[7] \dicDirectTranslationCS{multikulturalismus}
\dicEntry[fjölmennur] \dicTerm{fjöl··mennur} \dicIPA{{f}{j}{\oe}{l}{m}{\textepsilon}{n}{\textscy}{\textsubring{r}}} \dicPos{adj}[1]\dicFlx{}[-1] \textbf{1.} \dicDirectTranslationCS{slušně\,/\addthin hojně navštívený (schůze ap.)} \dicExampleIS{fjölmennur fundur} \dicExampleCS{hojně navštívená schůze}  \textbf{2.} \dicDirectTranslationCS{početný} \dicExampleIS{fjölmennur hópur} \dicExampleCS{početná skupina}
\dicEntry[fjölmiðill] \dicTerm{fjöl··mið|ill} \dicIPA{{f}{j}{\oe}{l}{m}{\textsci}{ð}{\textsci}{\textsubring{d}}{\textsubring{l}}} \dicPos{m}[6] \dicFlx{(‑ils, ‑lar)}[35] \dicDirectTranslationCS{(masové) médium, masmédium}
\dicEntry[fjölmiðlafræði] \dicTerm{fjöl·miðla··fræð|i} \dicIPA{{f}{j}{\oe}{l}{m}{\textsci}{ð}{l}{a}{f}{r}{a}{i}{ð}{\textsci}} \dicPos{f}[3] \dicFlx{(‑i)}[3] \dicDirectTranslationCS{mediální studia}
\dicEntry[fjölmiðlafræðingur] \dicTerm{fjöl·miðla·fræð··ing|ur} \dicIPA{{f}\-{j}\-{\oe}\-{l}\-{m}\-{\textsci}\-{ð}\-{l}\-{a}\-{f}\-{r}\-{a}\-{i}\-{ð}\-{i}\-{\ng}\-{\r{g}}\-{\textscy}\-{\textsubring{r}}\-} \dicPos{m}[6] \dicFlx{(‑s, ‑ar)}[8] \dicDirectTranslationCS{odborník\,/\addthin odbornice v~oboru mediálních studií}
\dicEntry[fjölmiðlun] \dicTerm{fjöl··miðl|un} \dicIPA{{f}{j}{\oe}{l}{m}{\textsci}{ð}{l}{\textscy}{\textsubring{n}}} \dicPos{f}[7] \dicFlx{(‑unar, ‑anir)}[8] \dicDirectTranslationCS{masová komunikace}
\dicEntry[fjölnota] \dicTerm{fjöl··nota} \dicIPA{{f}{j}{\oe}{l}{n}{\textopeno}{\textsubring{d}}{a}} \dicPos{adj}[13] \dicFlx{indecl}[1] \dicDirectTranslationCS{víceúčelový, mnohoúčelový}
\dicEntry[fjölraddaður] \dicTerm{fjöl··|radd·aður} \dicIPA{{f}{j}{\oe}{l}{r}{a}{\textsubring{d}}{a}{ð}{\textscy}{\textsubring{r}}} \dicPos{adj}[3] \dicFlx{(f ‑rödduð)}[2] \dicFieldCat{hud.} \dicDirectTranslationCS{polyfonní}
\dicEntry[fjölrit] \dicTerm{fjöl··rit} \dicIPA{{f}{j}{\oe}{l}{r}{\textsci}{\textsubring{d}}} \dicPos{n}[2] \dicFlx{(‑s, ‑)}[5] \dicDirectTranslationCS{kopie, duplikát}
\dicEntry[fjölrita] \dicTerm{fjöl··rit|a} \dicIPA{{f}{j}{\oe}{l}{r}{\textsci}{\textsubring{d}}{a}} \dicPos{v}[1] \dicFlx{(‑aði)}[1] \dicDirectTranslationCS{(o)kopírovat, (roz)množit} \dicExampleIS{fjölrita blaðið} \dicExampleCS{kopírovat stránku}
\dicEntry[fjölritari] \dicTerm{fjöl··rit·ar|i} \dicIPA{{f}{j}{\oe}{l}{r}{\textsci}{\textsubring{d}}{a}{r}{\textsci}} \dicPos{m}[1] \dicFlx{(‑a, ‑ar)}[13] \dicDirectTranslationCS{kopírka}
\dicEntry[fjölskrúðugur] \dicTerm{fjöl··skrúð·ugur} \dicIPA{{f}{j}{\oe}{l}{s}{\r{g}}{r}{u}{ð}{\textscy}{\textbabygamma}{\textscy}{\textsubring{r}}} \dicPos{adj}[1]\dicFlx{}[-8] \dicDirectTranslationCS{rozmanitý, různorodý, pestrý}
\dicEntry[fjölskylda] \dicTerm{fjöl··skyld|a} \dicsymFrequent\  \dicIPA{{f}{j}{\oe}{l}{s}{\r{\textObardotlessj}}{\textsci}{l}{\textsubring{d}}{a}} \dicPos{f}[1] \dicFlx{(‑u, ‑ur)}[13] \dicDirectTranslationCS{rodina} \dicExampleIS{fjögurra manna fjölskylda} \dicExampleCS{čtyřčlenná rodina}
\dicEntry[fjölskyldubætur] \dicTerm{fjöl·skyldu··bætur} \dicIPA{{f}{j}{\oe}{l}{s}{\r{\textObardotlessj}}{\textsci}{l}{\textsubring{d}}{\textscy}{\textsubring{b}}{a}{i}{\textsubring{d}}{\textscy}{\textsubring{r}}} \dicPos{f}[8] \dicFlx{pl}[6] \dicDirectTranslationCS{rodinné dáv\-ky\,/\addthin příspěvky}
\dicEntry[fjöltengi] \dicTerm{fjöl··tengi} \dicIPA{{f}{j}{\oe}{l}{t\smash{\textsuperscript{h}}}{ei}{\textltailn}{\r{\textObardotlessj}}{\textsci}} \dicPos{n}[2] \dicFlx{(‑s, ‑)}[16] \dicDirectTranslationCS{prodlužovací kabel s~vícenásobnou zásuvkou}
\dicEntry[fjölveri] \dicTerm{fjöl··veri} \dicIPA{{f}{j}{\oe}{l}{v}{\textepsilon}{r}{\textsci}} \dicPos{n}[2] \dicFlx{(‑s)}[20] \dicDirectTranslationCS{mnohomužství, polyandrie}
\dicEntry[fjölyrða] \dicTerm{fjöl··yr|ða} \dicIPA{{f}{j}{\oe}{\textlengthmark}{l}{\textsci}{r}{ð}{a}} \dicPos{v}[2] \dicFlx{(‑ti, ‑t)}[49] \dicPhraseIS{fjölyrða um e‑ð} \dicDirectTranslationCS{rozpovídat se o~(čem), vykládat o~(čem) zeširoka}
\dicEntry[fjölþjóðafyrirtæki] \dicTerm{fjöl·þjóða··fyrir·tæki} \textls[-10]{\dicIPA{{f}\-{j}\-{\oe}\-{l}\-{\texttheta}\-{j}\-{ou}\-{ð}\-{a}\-{f}\-{\textsci}\-{r}\-{\textsci}\-{\textsubring{r}}\-{t\smash{\textsuperscript{h}}}\-{a}\-{i}\-{\r{\textObardotlessj}}\-{\textsci}}} \dicPos{n}[2] \dicFlx{(‑s, ‑)}[16] \dicDirectTranslationCS{nad\-ná\-rod\-ní společnost\,/\addthin firma}
\dicEntry[fjölþjóðlegur] \dicTerm{fjöl··þjóð·legur} \textls[15]{\dicIPA{{f}{j}{\oe}{l}{\texttheta}{j}{ou}{ð}{l}{\textepsilon}{\textbabygamma}{\textscy}{\textsubring{r}}} \dicPos{adj}[1]\dicFlx{}[-8] \dicDirectTranslationCS{mezinárodní, nadnárodní, mnohonárodnostní}}
\dicEntry[fjör] \dicTerm{fjör} \dicsymFrequent\  \dicIPA{{f}{j}{\oe}{\textlengthmark}{\textsubring{r}}} \dicPos{n}[2] \dicFlx{(‑s)}[2] \textbf{1.} \dicSynonym*{lífsafl} \dicDirectTranslationCS{vitalita, životní energie};  \dicPhraseIS{vera í fullu fjöri} \dicDirectTranslationCS{být v~plné síle};  \dicPhraseIS{það er mikið fjör í e‑m} \dicDirectTranslationCS{(kdo) je plný života}  \textbf{2.} \dicSynonym{gleðskapur} \dicDirectTranslationCS{veselí, zábava};  \dicPhraseIS{það er líf og fjör e‑s staðar} \dicDirectTranslationCS{(kde) je živo a~veselo}
\dicEntry[fjörður] \dicTerm{fjörður} \dicsymFrequent\  \dicIPA{{f}{j}{\oe}{r}{ð}{\textscy}{\textsubring{r}}} \dicPos{m}[12] \dicFlx{(fjarðar, firðir)}[3] \dicDirectTranslationCS{záliv, fjord} \dicExampleIS{Austurströnd Grænlands er sundurskorin af fjörðum og sundum.} \dicExampleCS{Východní pobřeží Grónska je poseté zálivy a~průlivy.}
\dicEntry[fjörefni] \dicTerm{fjör··efni} \dicIPA{{f}{j}{\oe}{\textlengthmark}{r}{\textepsilon}{\textsubring{b}}{n}{\textsci}} \dicPos{n}[2] \dicFlx{(‑s, ‑)}[14] \dicSynonym{bætiefni} \dicDirectTranslationCS{vitamín}
\dicEntry[fjörga] \dicTerm{fjörg|a} \dicIPA{{f}{j}{\oe}{r}{\r{g}}{a}} \dicPos{v}[1] \dicFlx{(‑aði)}[1] \dicFlx{acc} \dicSynonym{lífga} \dicDirectTranslationCS{oživit, oživovat, stimulovat, povzbudit, povzbuzovat}
\dicEntry[fjörgamall] \dicTerm{fjör··|gamall} \dicIPA{{f}{j}{\oe}{r}{\r{g}}{a}{m}{a}{\textsubring{d}}{\textsubring{l}}} \dicPos{adj}[11] \dicFlx{(f ‑gömul)}[15] \dicDirectTranslationCS{letitý, velmi starý}
\dicEntry[fjörlegur] \dicTerm{fjör··legur} \dicIPA{{f}{j}{\oe}{r}{l}{\textepsilon}{\textbabygamma}{\textscy}{\textsubring{r}}} \dicPos{adj}[1]\dicFlx{}[-8] \dicSynonym{fjörugur} \dicDirectTranslationCS{živý, čilý, intenzivní}
\dicEntry[fjöru] \dicTerm{fjöru} \dicIPA{{f}{j}{\oe}{\textlengthmark}{r}{\textscy}} \dicPos{f} \dicFlx{sg gen} \dicLink{fjara\smash{\textsuperscript{1}}}
\dicEntry[fjörugrös] \dicTerm{fjöru··grös} \dicIPA{{f}{j}{\oe}{\textlengthmark}{r}{\textscy}{\r{g}}{r}{\oe}{s}} \dicPos{n}[2] \dicFlx{pl}[9] \dicFieldCat{bot.} \dicDirectTranslationCS{puchratka kadeřavá} \textit{(l.~{\textLA{Chondrus crispus}})}
\dicEntry[fjörugur] \dicTerm{fjör··ugur} \dicIPA{{f}{j}{\oe}{\textlengthmark}{r}{\textscy}{\textbabygamma}{\textscy}{\textsubring{r}}} \dicPos{adj}[1]\dicFlx{}[-8] \dicSynonym{frískur} \dicDirectTranslationCS{živý, bujný, bujarý, jarý, svižný}
\dicEntry[fjörukál] \dicTerm{fjöru··kál} \dicIPA{{f}{j}{\oe}{\textlengthmark}{r}{\textscy}{k\smash{\textsuperscript{h}}}{au}{\textsubring{l}}} \dicPos{n}[2] \dicFlx{(‑s)}[2] \dicFieldCat{bot.} \dicDirectTranslationCS{pomořanka přímořská} \dicIndirectTranslationCS{(poddruh pomořanky rostoucí na Islandu, Faerských ostrovech, Špicberkách, v~severní Skandinávii a~severním Rusku)} \textit{(l.~{\textLA{Cakile edentula subsp. Islandica}})}  \dicsymPhoto\ 
\dicFigure{ds_image_fjorukal_0_2.jpg}{Fjörukál}{Fjörukál - Jurgen Howaldt, CC BY-SA 2.0 DE}
\dicEntry[fjöruspói] \dicTerm{fjöru··spó|i} \dicIPA{{f}{j}{\oe}{\textlengthmark}{r}{\textscy}{s}{\textsubring{b}}{ou}{\textsci}} \dicPos{m}[1] \dicFlx{(‑a, ‑ar)}[1] \dicFieldCat{zool.} \dicDirectTranslationCS{koliha velká} \textit{(l.~{\textLA{Numenius arquata}})}  \dicsymPhoto\ 
\dicFigure{ds_image_fjoruspoi_0_1.jpg}{Fjöruspói}{Fjöruspói - Andreas Trepte, CC BY-SA 2.5}
\dicEntry[fjörutíu] \dicTerm{fjöru··tíu} \dicsymFrequent\  \dicIPA{{f}{j}{\oe}{\textlengthmark}{r}{\textscy}{t\smash{\textsuperscript{h}}}{i}{j}{\textscy}} \dicPos{num} \dicFlx{indecl} \dicDirectTranslationCS{čtyřicet} \dicExampleIS{Hitastigið úti var tæplega fjörutíu gráður.} \dicExampleCS{Teplota venku byla téměř čtyřicet stupňů.}
\dicEntry[fjörutíukall] \dicTerm{fjöru·tíu··kall} \dicIPA{{f}{j}{\oe}{\textlengthmark}{r}{\textscy}{t\smash{\textsuperscript{h}}}{i}{j}{\textscy}{k\smash{\textsuperscript{h}}}{a}{\textsubring{d}}{\textsubring{l}}} \dicPos{m}[4] \dicFlx{(‑s, ‑ar)}[4] \dicLangCat{hovor.} \dicDirectTranslationCS{čtyřicet (korun)} \dicIndirectTranslationCS{(částka)}
\dicEntry[fjötra] \dicTerm{fjötr|a} \dicIPA{{f}{j}{\oe}{\textlengthmark}{\textsubring{d}}{r}{a}} \dicPos{v}[1] \dicFlx{(‑aði)}[1] \dicFlx{acc} \dicDirectTranslationCS{spoutat, svázat} \dicExampleIS{fjötra e‑n á höndum og fótum} \dicExampleCS{spoutat (komu) ruce a~nohy}
\dicEntry[fjötur] \dicTerm{fjöt|ur} \dicIPA{{f}{j}{\oe}{\textlengthmark}{\textsubring{d}}{\textscy}{\textsubring{r}}} \dicPos{m}[5] \dicFlx{(‑urs, ‑rar)}[1] \dicDirectTranslationCS{pouto, okov} \dicExampleIS{setja fjötur á fætur e‑rs} \dicExampleCS{spoutat (komu) nohy};  \dicPhraseIS{e‑að er e‑m fjötur um fót} \dicLangCat{přen.} \dicDirectTranslationCS{(co komu) svazuje ruce, (co koho) brzdí}
\dicEntry[fl.] \dicTerm{fl.} \dicPos{zkr} \dicPhraseIS{fleira} \dicDirectTranslationCS{více}
\dicEntry[flaðra] \dicTerm{flaðr|a} \dicIPA{{f}{l}{a}{ð}{r}{a}} \dicPos{v}[1] \dicFlx{(‑aði)}[50] \dicPhraseIS{flaðra upp um e‑n} \dicSynonym*{skríða fyrir} \dicDirectTranslationCS{podlézat (komu)}
\dicEntry[flaður] \dicTerm{flaður} \dicIPA{{f}{l}{a}{\textlengthmark}{ð}{\textscy}{\textsubring{r}}} \dicPos{n}[2] \dicFlx{(‑s)}[28] \dicDirectTranslationCS{podlézání}
\dicEntry[flag] \dicTerm{flag} \dicIPA{{f}{l}{a}{\textlengthmark}{x}} \dicPos{n}[2] \dicFlx{(‑s, flög)}[8] \dicDirectTranslationCS{nezarostlý úhor}
\dicEntry[flaga] \dicTerm{flaga} \dicIPA{{f}{l}{a}{\textlengthmark}{\textbabygamma}{a}} \dicPos{f}[1] \dicFlx{(flögu, flögur)}[20] \textbf{1.} \dicSynonym{hella\smash{\textsuperscript{1}}} \dicDirectTranslationCS{oblázek} \dicExampleIS{steinflaga} \dicExampleCS{oblázek}  \textbf{2.} \dicDirectTranslationCS{plátek} \dicExampleIS{skera ostinn í þunnar flögur} \dicExampleCS{pokrájet sýr na tenké plátky}  \textbf{3.} \dicDirectTranslationCS{destička, lamela} \dicExampleIS{málmflaga} \dicExampleCS{kovová destička}
\dicEntry[flagahnoðri] \dicTerm{flaga··hnoðr|i} \dicIPA{{f}{l}{a}{\textlengthmark}{\textbabygamma}{a}{\textsubring{n}}{\textopeno}{ð}{r}{\textsci}} \dicPos{m}[1] \dicFlx{(‑a, ‑ar)}[1] \dicFieldCat{bot.} \dicDirectTranslationCS{rozchodník huňatý} \textit{(l.~{\textLA{Sedum villosum}})}  \dicsymPhoto\ 
\dicFigure{ds_image_flagahnodri_0_2.jpg}{Flagahnoðri}{Flagahnoðri - Thommybe, GFDL}
\dicEntry[flagari] \dicTerm{flag··ar|i} \dicIPA{{f}{l}{a}{\textlengthmark}{\textbabygamma}{a}{r}{\textsci}} \dicPos{m}[1] \dicFlx{(‑a, ‑ar)}[10] \dicSynonym*{kvennabósi} \dicDirectTranslationCS{svůdce, svůdník, lamač dámských srdcí}
\dicEntry[flagð] \dicTerm{flagð} \dicIPA{{f}{l}{a}{\textbabygamma}{\texttheta}} \dicPos{n}[2] \dicFlx{(‑s, flögð)}[8] \dicLangCat{han.} \dicDirectTranslationCS{čarodějnice, babizna}
\dicEntry[flagg] \dicTerm{flagg} \dicIPA{{f}{l}{a}{\r{g}}{\textlengthmark}} \dicPos{n}[2] \dicFlx{(‑s, flögg)}[8] \dicDirectTranslationCS{vlajka, prapor} \dicExampleIS{draga upp flagg} \dicExampleCS{vytáhnout vlajku nahoru};  \dicPhraseIS{sigla undir fölsku flaggi} \dicLangCat{přen.} \dicDirectTranslationCS{plout pod falešnou\,/\addthin cizí vlajkou} \dicIndirectTranslationCS{(vydávat se za někoho jiného, než kým člověk opravdu je)}
\dicEntry[flagga] \dicTerm{flagg|a} \dicIPA{{f}{l}{a}{\r{g}}{\textlengthmark}{a}} \dicPos{v}[1] \dicFlx{(‑aði)}[50] \dicFlx{dat} \textbf{1.} \dicDirectTranslationCS{vyvěsit vlajku} \dicExampleIS{Hann flaggaði íslenska fánanum í gær.} \dicExampleCS{Včera vyvěsil islandskou vlajku.}  \textbf{2.} \dicDirectTranslationCS{dávat na obdiv}
\dicEntry[flaggskip] \dicTerm{flagg··skip} \dicIPA{{f}{l}{a}{\r{g}}{s}{\r{\textObardotlessj}}{\textsci}{\textsubring{b}}} \dicPos{n}[2] \dicFlx{(‑s, ‑)}[5] \dicFieldCat{nám.} \dicDirectTranslationCS{vlajková loď}
\dicEntry[flaggstöng] \dicTerm{flagg··|stöng} \dicIPA{{f}{l}{a}{\r{g}}{s}{\textsubring{d}}{\oe i}{\ng}{\r{g}}} \dicPos{f}[8] \dicFlx{(‑stangar, ‑stengur\,/\addthin ‑stangir)}[9] \dicDirectTranslationCS{stožár}
\dicEntry[flagna] \dicTerm{flagn|a} \dicIPA{{f}{l}{a}{\r{g}}{n}{a}} \dicPos{v}[1] \dicFlx{(‑aði)}[13] \dicDirectTranslationCS{loupat se, sloupnout se, sloupávat se} \dicExampleIS{Skinnið er farið að flagna af.} \dicExampleCS{Kůže se začíná sloupávat.}
\dicEntry[flak] \dicTerm{flak} \dicIPA{{f}{l}{a}{\textlengthmark}{\r{g}}} \dicPos{n}[2] \dicFlx{(‑s, flök)}[8] \textbf{1.} \dicSynonym{rekald} \dicDirectTranslationCS{vrak, troska (lodi ap.)} \dicExampleIS{skipsflak} \dicExampleCS{vrak lodi}  \textbf{2.} \dicDirectTranslationCS{(rybí) filé, filet} \dicExampleIS{fiskflak} \dicExampleCS{rybí filé}
\dicEntry[flaka] \dicTerm{flak|a\smash{\textsuperscript{1}}} \dicIPA{{f}{l}{a}{\textlengthmark}{\r{g}}{a}} \dicPos{v}[1] \dicFlx{(‑aði)}[13] \dicFlx{acc} \dicDirectTranslationCS{rozřezat na filety, vykostit (rybu ap.)}
\dicEntry[flaka] \dicTerm{flak|a\smash{\textsuperscript{2}}} \dicIPA{{f}{l}{a}{\textlengthmark}{\r{g}}{a}} \dicPos{v}[2] \dicFlx{(‑ti, ‑að)}[86] \dicSynonym*{vera opinn} \dicDirectTranslationCS{zet, rozvírat se, otvírat se} \dicExampleIS{Sárið flakir opið.} \dicExampleCS{Rána je otevřená.}
\dicEntry[flakk] \dicTerm{flakk} \dicIPA{{f}{l}{a}{h}{\r{g}}} \dicPos{n}[2] \dicFlx{(‑s)}[2] \dicSynonym*{fótavist} \dicDirectTranslationCS{toulání, potulka, vandr} \dicExampleIS{vera á flakki um landið} \dicExampleCS{toulat se po zemi}
\dicEntry[flakka] \dicTerm{flakk|a} \dicIPA{{f}{l}{a}{h}{\r{g}}{a}} \dicPos{v}[1] \dicFlx{(‑aði)}[50] \dicDirectTranslationCS{toulat se, putovat, vandrovat} \dicExampleIS{flakka um sveitina} \dicExampleCS{putovat po venkově}
\dicEntry[flakkari] \dicTerm{flakk··ar|i} \dicIPA{{f}{l}{a}{h}{\r{g}}{a}{r}{\textsci}} \dicPos{m}[1] \dicFlx{(‑a, ‑ar)}[10] \textbf{1.} \dicSynonym{flækingur} \dicDirectTranslationCS{tulák, tulačka, vandrák, vandračka}  \textbf{2.} \dicFieldCat{poč.} \dicDirectTranslationCS{externí disk}
\dicEntry[flan] \dicTerm{flan} \dicIPA{{f}{l}{a}{\textlengthmark}{\textsubring{n}}} \dicPos{n}[2] \dicFlx{(‑s)}[2] \dicDirectTranslationCS{ukvapenost, nerozvážnost}
\dicEntry[flana] \dicTerm{flan|a} \dicIPA{{f}{l}{a}{\textlengthmark}{n}{a}} \dicPos{v}[1] \dicFlx{(‑aði)}[28] \dicDirectTranslationCS{ukvapit se, unáhlit se, pustit se (nerozvážně)} \dicExampleIS{flana út í e‑ð} \dicExampleCS{nerozvážně se pustit do (čeho)}
\dicEntry[flas] \dicTerm{flas} \dicIPA{{f}{l}{a}{\textlengthmark}{s}} \dicPos{n}[2] \dicFlx{(‑s)}[2] \dicSynonym{asi} \dicDirectTranslationCS{spěch, kvap};  \dicPhraseIS{ekki er flas  til fagnaðar} \dicLangCat{přen.} \dicDirectTranslationCS{spěchat se nevyplácí};  \dicPhraseIS{koma í flasið á e‑m} \dicDirectTranslationCS{narazit\,/\addthin natrefit na (koho) náhodou}
\dicEntry[flasa] \dicTerm{flasa\smash{\textsuperscript{1}}} \dicIPA{{f}{l}{a}{\textlengthmark}{s}{a}} \dicPos{f}[1] \dicFlx{(flösu)}[2] \dicDirectTranslationCS{lupy}
\dicEntry[flasa] \dicTerm{flas|a\smash{\textsuperscript{2}}} \dicIPA{{f}{l}{a}{\textlengthmark}{s}{a}} \dicPos{v}[1] \dicFlx{(‑aði)}[50] \dicSynonym{ana} \dicDirectTranslationCS{spěchat, chvátat, hnát se}
\dicEntry[flaska] \dicTerm{flaska} \dicsymFrequent\  \dicIPA{{f}{l}{a}{s}{\r{g}}{a}} \dicPos{f}[1] \dicFlx{(flösku, flöskur)}[14] \dicDirectTranslationCS{láhev, flaška} \dicExampleIS{flöskur með glæru innihaldi} \dicExampleCS{láhve s~průzračným obsahem}
\dicEntry[flatar] \dicTerm{flatar\smash{\textsuperscript{1}}} \dicIPA{{f}{l}{a}{\textlengthmark}{\textsubring{d}}{a}{\textsubring{r}}} \dicPos{m} \dicFlx{sg gen} \dicLink{flötur}
\dicEntry[flatar] \dicTerm{flatar\smash{\textsuperscript{2}}} \dicIPA{{f}{l}{a}{\textlengthmark}{\textsubring{d}}{a}{\textsubring{r}}} \dicPos{f} \dicFlx{sg gen} \dicLink{flöt}
\dicEntry[flatarmál] \dicTerm{flatar··mál} \dicIPA{{f}{l}{a}{\textlengthmark}{\textsubring{d}}{a}{r}{m}{au}{\textsubring{l}}} \dicPos{n}[2] \dicFlx{(‑s)}[2] \dicDirectTranslationCS{povrch, plocha, obsah} \dicExampleIS{flatarmál þríhyrnings} \dicExampleCS{obsah trojúhelníka}
\dicEntry[flatbrauð] \dicTerm{flat··brauð}\dicTerm{, flatkaka} \dicIPA{{f}\-{l}\-{a}\-{\textlengthmark}\-{\textsubring{d}}\-{\textsubring{b}}\-{r}\-{\oe i}\-{\texttheta}\-} \dicPos{n}[2] \dicFlx{(‑s, ‑)}[5] \dicFieldCat{kulin.} \dicIndirectTranslationCS{tradiční islandský chléb kruhového tvaru pečený na uhlících nebo na plotně}
\dicEntry[flatfiskur] \dicTerm{flat··fisk|ur} \dicIPA{{f}{l}{a}{\textlengthmark}{\textsubring{d}}{f}{\textsci}{s}{\r{g}}{\textscy}{\textsubring{r}}} \dicPos{m}[6] \dicFlx{(‑s, ‑ar)}[10] \dicFieldCat{zool.} \dicDirectTranslationCS{platýs} \textit{(l.~{\textLA{Pleuronectiformes}})}
\dicEntry[flatir] \dicTerm{flatir} \dicIPA{{f}{l}{a}{\textlengthmark}{\textsubring{d}}{\textsci}{\textsubring{r}}} \dicPos{f} \dicFlx{pl nom} \dicLink{flöt}
\dicEntry[flatkaka] \dicTerm{flat··|kaka} \dicIPA{{f}{l}{a}{\textlengthmark}{\textsubring{d}}{k\smash{\textsuperscript{h}}}{a}{\r{g}}{a}} \dicPos{f}[1] \dicFlx{(‑köku, ‑kökur)}[20] \dicLink{flatbrauð}
\dicEntry[flatmaga] \dicTerm{flat··mag|a} \dicIPA{{f}{l}{a}{\textlengthmark}{\textsubring{d}}{m}{a}{\textbabygamma}{a}} \dicPos{v}[1] \dicFlx{(‑aði)}[50] \dicDirectTranslationCS{polehávat, vyvalovat se, hovět si} \dicExampleIS{flatmaga á ströndinni} \dicExampleCS{polehávat na pláži}
\dicEntry[flatsæng] \dicTerm{flat··sæng} \dicIPA{{f}{l}{a}{\textlengthmark}{\textsubring{d}}{s}{a}{i}{\ng}{\r{g}}} \dicPos{f}[10] \dicFlx{(‑ur, ‑ur)}[1] \dicDirectTranslationCS{provizorní lůžko, lůžko na zemi, slamník}
\dicEntry[flatur] \dicTerm{flatur} \dicsymFrequent\  \dicIPA{{f}{l}{a}{\textlengthmark}{\textsubring{d}}{\textscy}{\textsubring{r}}} \dicPos{adj}[1] \dicFlx{(f flöt)}[2] \dicDirectTranslationCS{plochý, rovný} \dicExampleIS{flöt þök} \dicExampleCS{rovné střechy};  \dicPhraseIS{fara flatt fyrir e‑m} \dicSynonym*{fara halloka} \dicDirectTranslationCS{utrpět porážku od (koho)};  \dicPhraseIS{liggja flatur fyrir e‑u\,/\addthin e‑m} \dicLangCat{přen.} \dicDirectTranslationCS{být z~(čeho\,/\addthin koho) hotový};  \dicPhraseIS{halla undir flatt} \dicDirectTranslationCS{sklonit\,/\addthin sehnout\,/\addthin schýlit hlavu (na stranu)}
\dicEntry[flauel] \dicTerm{flauel} \dicIPA{{f}{l}{\oe i}{\textepsilon}{\textsubring{l}}} \dicPos{n}[2] \dicFlx{(‑s)}[2] \textbf{1.} \dicDirectTranslationCS{samet}  \textbf{2.} \dicDirectTranslationCS{manšestr} \dicExampleIS{flauelsbuxur} \dicExampleCS{manšestrové kalhoty}
\dicEntry[flaug] \dicTerm{flaug} \dicIPA{{f}{l}{\oe i}{\textlengthmark}{x}} \dicPos{v} \dicFlx{ind pf sg 1 pers} \dicLink{fljúga}
\dicEntry[flausturslegur] \dicTerm{flausturs··legur} \dicIPA{{f}{l}{\oe i}{s}{\textsubring{d}}{\textscy}{\textsubring{r}}{s}{l}{\textepsilon}{\textbabygamma}{\textscy}{\textsubring{r}}} \dicPos{adj}[1]\dicFlx{}[-8] \dicDirectTranslationCS{uspěchaný, nedbalý, odfláknutý} \dicExampleIS{flausturslegur frágangur} \dicExampleCS{ukvapený závěr}
\dicEntry[flaut] \dicTerm{flaut\smash{\textsuperscript{1}}} \dicIPA{{f}{l}{\oe i}{\textlengthmark}{\textsubring{d}}} \dicPos{n}[2] \dicFlx{(‑s, ‑)}[5] \dicDirectTranslationCS{pískot, hvízdnutí, písknutí (sudího ap.), (za)troubení (auta ap.)}
\dicEntry[flaut] \dicTerm{flaut\smash{\textsuperscript{2}}} \dicIPA{{f}{l}{\oe i}{\textlengthmark}{\textsubring{d}}} \dicPos{v} \dicFlx{ind pf sg 1 pers} \dicLink{fljóta}
\dicEntry[flauta] \dicTerm{flaut|a\smash{\textsuperscript{1}}} \dicIPA{{f}{l}{\oe i}{\textlengthmark}{\textsubring{d}}{a}} \dicPos{f}[1] \dicFlx{(‑u, ‑ur)}[19] \textbf{1.} \dicFieldCat{hud.} \dicDirectTranslationCS{flétna}  \textbf{2.} \dicDirectTranslationCS{klakson} \dicExampleIS{bílflauta} \dicExampleCS{klakson}  \textbf{3.} \dicDirectTranslationCS{píšťala, píšťalka}  \textbf{4.} \dicFieldCat{nám.} \dicDirectTranslationCS{flauta} \dicIndirectTranslationCS{(loď)}
\dicEntry[flauta] \dicTerm{flaut|a\smash{\textsuperscript{2}}} \dicIPA{{f}{l}{\oe i}{\textlengthmark}{\textsubring{d}}{a}} \dicPos{v}[1] \dicFlx{(‑aði)}[1] \dicFlx{acc} \textbf{1.} \dicSynonym{blístra\smash{\textsuperscript{2}}} \dicDirectTranslationCS{pískat, hvízdat}  \textbf{2.} \dicDirectTranslationCS{hrát na flétnu}  \textbf{3.} \dicDirectTranslationCS{(za)troubit (v~autě ap.)}
\dicEntry[flautuleikari] \dicTerm{flautu··leik·ar|i} \dicIPA{{f}{l}{\oe i}{\textlengthmark}{\textsubring{d}}{\textscy}{l}{ei}{\r{g}}{a}{r}{\textsci}} \dicPos{m}[1] \dicFlx{(‑a, ‑ar)}[13] \dicDirectTranslationCS{flétnista, flétnistka}
\dicEntry[flámæli] \dicTerm{flá··mæli} \dicIPA{{f}{l}{au}{\textlengthmark}{m}{a}{i}{l}{\textsci}} \dicPos{n}[2] \dicFlx{(‑s)}[20] \dicFieldCat{jaz.} \dicIndirectTranslationCS{výslovnostní varianta v~islandském jazyce}
\dicEntry[fláræði] \dicTerm{flá··ræði} \dicIPA{{f}{l}{au}{\textlengthmark}{r}{a}{i}{ð}{\textsci}} \dicPos{n}[2] \dicFlx{(‑s)}[20] \dicSynonym{undirferli} \dicDirectTranslationCS{proradnost, podlost}
\dicEntry[fláttskapur] \dicTerm{flátt··skap|ur} \dicIPA{{f}{l}{au}{h}{\textsubring{d}}{s}{\r{g}}{a}{\textsubring{b}}{\textscy}{\textsubring{r}}} \dicPos{m}[10] \dicFlx{(‑ar)}[15] \dicSynonym{slægð} \dicDirectTranslationCS{pokrytectví, lstivost}
\dicEntry[fleðulegur] \dicTerm{fleðu··legur} \textls[15]{\dicIPA{{f}{l}{\textepsilon}{\textlengthmark}{ð}{\textscy}{l}{\textepsilon}{\textbabygamma}{\textscy}{\textsubring{r}}} \dicPos{adj}[1]\dicFlx{}[-8] \dicDirectTranslationCS{servilní, podlézavý}} \dicExampleIS{fleðulegur svipur} \dicExampleCS{servilní výraz}
\dicEntry[fleipra] \dicTerm{fleipr|a} \dicIPA{{f}{l}{ei}{\textlengthmark}{\textsubring{b}}{r}{a}} \dicPos{v}[1] \dicFlx{(‑aði)}[1] \dicFlx{dat} \dicSynonym*{aula út úr sér} \dicDirectTranslationCS{vybrebtat, vyžvanit, prokecnout};  \dicPhraseIS{fleipra e‑u út úr sér} \dicDirectTranslationCS{vyžvanit (co)}
\dicEntry[fleipur] \dicTerm{fleipur} \dicIPA{{f}{l}{ei}{\textlengthmark}{\textsubring{b}}{\textscy}{\textsubring{r}}} \dicPos{n}[2] \dicFlx{(‑s)}[28] \dicDirectTranslationCS{(plané) řeči, klepy};  \dicPhraseIS{fara með fleipur} \dicDirectTranslationCS{tlachat, drbat}
\dicEntry[fleirtala] \dicTerm{fleir··|tala} \dicIPA{{f}{l}{ei}{\textsubring{r}}{t\smash{\textsuperscript{h}}}{a}{l}{a}} \dicPos{f}[1] \dicFlx{(‑tölu)}[2] \dicFieldCat{jaz.} \dicDirectTranslationCS{množné číslo, plurál}
\dicEntry[fleiryrtur] \dicTerm{fleir··yrtur} \dicIPA{{f}{l}{ei}{\textlengthmark}{r}{\textsci}{\textsubring{r}}{\textsubring{d}}{\textscy}{\textsubring{r}}} \dicPos{adj}[1]\dicFlx{}[-13] \dicFieldCat{jaz.} \dicDirectTranslationCS{víceslovný};  \dicPhraseIS{fleiryrt samtenging} \dicFieldCat{jaz.} \dicDirectTranslationCS{víceslovná spojka}
\dicEntry[fleka] \dicTerm{flek|a} \dicIPA{{f}{l}{\textepsilon}{\textlengthmark}{\r{g}}{a}} \dicPos{v}[1] \dicFlx{(‑aði)}[1] \dicFlx{acc} \dicSynonym{fífla} \dicDirectTranslationCS{svést, svádět} \dicExampleIS{fleka stúlku} \dicExampleCS{svádět dívku}
\dicEntry[fleki] \dicTerm{flek|i} \dicIPA{{f}{l}{\textepsilon}{\textlengthmark}{\r{\textObardotlessj}}{\textsci}} \dicPos{m}[1] \dicFlx{(‑a, ‑ar)}[1] \textbf{1.} \dicDirectTranslationCS{vor, prám} \dicExampleIS{fara yfir vatnið á fleka} \dicExampleCS{přejet vodu na voru}  \textbf{2.} \dicFieldCat{geol.} \dicDirectTranslationCS{tektonická deska}
\dicEntry[flekka] \dicTerm{flekk|a} \dicIPA{{f}{l}{\textepsilon}{h}{\r{g}}{a}} \dicPos{v}[1] \dicFlx{(‑aði)}[1] \dicFlx{acc} \dicSynonym{bletta} \dicDirectTranslationCS{pošpinit, poskvrnit} \dicExampleIS{flekka hendur sínar} \dicExampleCS{pošpinit si ruce}
\dicEntry[flekklaus] \dicTerm{flekk··laus} \dicIPA{{f}{l}{\textepsilon}{h}{\r{g}}{l}{\oe i}{s}} \dicPos{adj}[5]\dicFlx{}[-1] \dicDirectTranslationCS{neposkvrněný, (jsoucí) bez poskvrny}
\dicEntry[flekkóttur] \dicTerm{flekk··óttur} \dicIPA{{f}{l}{\textepsilon}{h}{\r{g}}{ou}{h}{\textsubring{d}}{\textscy}{\textsubring{r}}} \dicPos{adj}[1]\dicFlx{}[-10] \dicDirectTranslationCS{flekatý, skvrnitý}
\dicEntry[flekkur] \dicTerm{flekk|ur} \dicIPA{{f}{l}{\textepsilon}{h}{\r{g}}{\textscy}{\textsubring{r}}} \dicPos{m}[9] \dicFlx{(‑s\,/\addthin ‑jar, ‑ir)}[26] \textbf{1.} \dicDirectTranslationCS{flek, skvrna}  \textbf{2.} \dicLangCat{přen.} \dicDirectTranslationCS{(po)skvrna} \dicExampleIS{flekkur á mannorði} \dicExampleCS{poskvrna na dobrém jménu}  \textbf{3.} \dicDirectTranslationCS{shrabané seno v~řádcích}
\dicEntry[flengja] \dicTerm{fleng|ja} \dicIPA{{f}{l}{ei}{\textltailn}{\r{\textObardotlessj}}{a}} \dicPos{v}[2] \dicFlx{(‑di, ‑t)}[129] \dicFlx{acc} \dicSynonym{hýða} \dicDirectTranslationCS{naplácat, dát na zadek, nařezat} \dicExampleIS{flengja e‑n með vendi} \dicExampleCS{nařezat (komu) rákoskou}
\dicEntry[flenna] \dicTerm{flenn|a\smash{\textsuperscript{1}}} \dicIPA{{f}{l}{\textepsilon}{n}{\textlengthmark}{a}} \dicPos{f}[1] \dicFlx{(‑u, ‑ur)}[7] \dicDirectTranslationCS{rajda, coura, běhna}
\dicEntry[flenna] \dicTerm{flenn|a\smash{\textsuperscript{2}}} \dicIPA{{f}{l}{\textepsilon}{n}{\textlengthmark}{a}} \dicPos{v}[2] \dicFlx{(‑ti, ‑t)}[63] \dicFlx{acc} \dicDirectTranslationCS{roztáhnout, rozevřít} \dicExampleIS{flenna upp gluggann} \dicExampleCS{rozevřít okno}
\dicEntry[flensa] \dicTerm{flens|a} \dicIPA{{f}{l}{\textepsilon}{n}{s}{a}} \dicPos{f}[1] \dicFlx{(‑u)}[5] \dicFieldCat{med.} \dicDirectTranslationCS{chřipka}
\dicEntry[flesk] \dicTerm{flesk} \dicIPA{{f}{l}{\textepsilon}{s}{\r{g}}} \dicPos{n}[2] \dicFlx{(‑s)}[2] \dicSynonym{beikon} \dicDirectTranslationCS{slanina}
\dicEntry[flestallir] \dicTerm{flest··allir} \dicIPA{{f}{l}{\textepsilon}{s}{\textsubring{d}}{a}{\textsubring{d}}{l}{\textsci}{\textsubring{r}}} \dicPos{pron}[14] \dicFlx{pl} \dicDirectTranslationCS{většina, téměř všichni}
\dicEntry[flestur] \dicTerm{flestur} \dicIPA{{f}{l}{\textepsilon}{s}{\textsubring{d}}{\textscy}{\textsubring{r}}} \dicPos{adj} \dicFlx{m sg nom sup} \dicLink{margur}
\dicEntry[flet] \dicTerm{flet} \dicsymFrequent\  \dicIPA{{f}{l}{\textepsilon}{\textlengthmark}{\textsubring{d}}} \dicPos{n}[2] \dicFlx{(‑s, ‑)}[5] \dicDirectTranslationCS{pryčna, palanda} \dicExampleIS{leggjast á flet} \dicExampleCS{lehnout si na palandu}
\dicEntry[fletir] \dicTerm{fletir} \dicIPA{{f}{l}{\textepsilon}{\textlengthmark}{\textsubring{d}}{\textsci}{\textsubring{r}}} \dicPos{m} \dicFlx{pl nom} \dicLink{flötur}
\dicEntry[fletja] \dicTerm{fletja} \dicIPA{{f}{l}{\textepsilon}{\textlengthmark}{\textsubring{d}}{j}{a}} \dicPos{v}[4] \dicFlx{(flet, flatti, flöttum, fletti, flatt)}[48] \dicFlx{acc} \textbf{1.} \dicPhraseIS{fletja út e‑ð} \dicDirectTranslationCS{vyválet (co) (těsto ap.)}  \textbf{2.} \dicDirectTranslationCS{rozporcovat (rybu ap.)};  \dicIdiom{fletjast}[út]{ \dicPhraseIS{fletjast út}} \dicFlx{refl} \dicDirectTranslationCS{zploštit se, stát se plochým}
\dicEntry[fletta] \dicTerm{flett|a} \dicsymFrequent\  \dicIPA{{f}{l}{\textepsilon}{h}{\textsubring{d}}{a}} \dicPos{v}[2] \dicFlx{(‑i, ‑)}[1] \dicFlx{dat} \textbf{1.} \dicSynonym*{taka af} \dicDirectTranslationCS{sundat, svléknout (povlak na peřinu ap.)};  \dicPhraseIS{fletta sig klæðum} \dicDirectTranslationCS{svléknout si šaty}  \textbf{2.} \dicSynonym*{blaða í} \dicDirectTranslationCS{(na)listovat} \dicExampleIS{fletta blaðinu} \dicExampleCS{listovat v~novinách};  \dicIdiom{fletta}[ofan af]{ \dicPhraseIS{fletta ofan af e‑u}} \dicDirectTranslationCS{odhalit (co)};  \dicIdiom{fletta}[upp]{ \dicPhraseIS{fletta e‑ð upp}} \dicDirectTranslationCS{vyhledat (co) (v~knize ap.)}
\dicEntry[flettiorð] \dicTerm{fletti··orð} \dicIPA{{f}{l}{\textepsilon}{h}{\textsubring{d}}{\textsci}{\textopeno}{r}{\texttheta}} \dicPos{n}[2] \dicFlx{(‑s, ‑)}[5] \dicFieldCat{jaz.} \dicDirectTranslationCS{lemma}
\dicEntry[fley] \dicTerm{fley} \dicIPA{{f}{l}{ei}{\textlengthmark}} \dicPos{n}[2] \dicFlx{(‑s, ‑)}[13] \dicLangCat{básn.} \dicDirectTranslationCS{člun, loď}
\dicEntry[fleyga] \dicTerm{fleyg|a} \dicIPA{{f}{l}{ei}{\textlengthmark}{\textbabygamma}{a}} \dicPos{v}[1] \dicFlx{(‑aði)}[1] \dicFlx{acc} \dicDirectTranslationCS{vklínit, zaklínit}
\dicEntry[fleygja] \dicTerm{fleyg|ja} \dicsymFrequent\  \dicIPA{{f}{l}{ei}{j}{\textlengthmark}{a}} \dicPos{v}[2] \dicFlx{(‑ði, ‑t)}[95] \dicFlx{dat} \textbf{1.} \dicDirectTranslationCS{zahodit, vyhodit, zbavit se} \dicExampleIS{fleygja gömlu skónum sínum} \dicExampleCS{zahodit své staré boty}  \textbf{2.} \dicDirectTranslationCS{odhodit (od sebe ap.)};  \dicPhraseIS{fleygja sér e‑s staðar} \dicDirectTranslationCS{vrhnout se (kam)};  \dicPhraseIS{e‑u fleygir fram} \dicFlx{impers} \dicLangCat{přen.} \dicDirectTranslationCS{(co) jde mílovými kroky kupředu}
\dicEntry[fleygur] \dicTerm{fleyg|ur\smash{\textsuperscript{1}}} \dicIPA{{f}{l}{ei}{\textlengthmark}{\textbabygamma}{\textscy}{\textsubring{r}}} \dicPos{m}[6] \dicFlx{(‑s, ‑ar)}[22] \textbf{1.} \dicDirectTranslationCS{klín, klínek}  \textbf{2.} \dicSynonym*{lítil flaska} \dicDirectTranslationCS{kapesní láhev, čutora, placatka}
\dicEntry[fleygur] \dicTerm{fleygur\smash{\textsuperscript{2}}} \dicIPA{{f}{l}{ei}{\textlengthmark}{\textbabygamma}{\textscy}{\textsubring{r}}} \dicPos{adj}[1]\dicFlx{}[-1] \dicDirectTranslationCS{létavý, schopný letu, letuschopný (pták ap.)}
\dicEntry[fleyta] \dicTerm{fleyt|a\smash{\textsuperscript{1}}} \dicIPA{{f}{l}{ei}{\textlengthmark}{\textsubring{d}}{a}} \dicPos{f}[1] \dicFlx{(‑u, ‑ur)}[19] \dicDirectTranslationCS{malý člun, loďka}
\dicEntry[fleyta] \dicTerm{fleyt|a\smash{\textsuperscript{2}}} \dicIPA{{f}{l}{ei}{\textlengthmark}{\textsubring{d}}{a}} \dicPos{v}[2] \dicFlx{(‑ti, ‑t)}[54] \dicFlx{dat} \textbf{1.} \dicDirectTranslationCS{nést, unášet} \dicExampleIS{fleyta trénu niður ána} \dicExampleCS{unášet dřevo po řece}  \textbf{2.} \dicSynonym{bera} \dicDirectTranslationCS{unést, pojmout}  \textbf{3.} \dicDirectTranslationCS{překonat, přenést} \dicExampleIS{fleyta fjölskyldunni yfir erfiðleikana} \dicExampleCS{překonat s~rodinou potíže};  \dicPhraseIS{fleyta rjómann af e‑u} \dicLangCat{přen.} \dicDirectTranslationCS{slíznout smetanu z~(čeho)};  \dicPhraseIS{geta fleytt sér í e‑ð} \dicDirectTranslationCS{poradit si s~(čím)}
\dicEntry[flétta] \dicTerm{flétt|a\smash{\textsuperscript{1}}} \dicIPA{{f}{l}{j}{\textepsilon}{h}{\textsubring{d}}{a}} \dicPos{f}[1] \dicFlx{(‑u, ‑ur)}[19] \textbf{1.} \dicDirectTranslationCS{cop, pletenec} \dicExampleIS{greiða hárið í fléttur} \dicExampleCS{učesat vlasy do copů}  \textbf{2.} \dicFieldCat{lit.} \dicDirectTranslationCS{děj}  \textbf{3.} \dicFieldCat{bot.} \dicDirectTranslationCS{lišejník} \textit{(l.~{\textLA{Lichenes}})} \dicAntonym*{skóf}
\dicEntry[flétta] \dicTerm{flétt|a\smash{\textsuperscript{2}}} \dicIPA{{f}{l}{j}{\textepsilon}{h}{\textsubring{d}}{a}} \dicPos{v}[3] \dicFlx{(‑aði) (‑i, ‑)}[1] \dicFlx{dat} \dicDirectTranslationCS{(za)plést, splétat, zavinout (do copů ap.)} \dicExampleIS{flétta hárið} \dicExampleCS{zaplést vlasy};  \dicIdiom{fléttast}{ \dicPhraseIS{fléttast}} \dicFlx{refl} \dicDirectTranslationCS{proplétat se, prolínat se}
\dicEntry[flibbi] \dicTerm{flibb|i} \dicIPA{{f}{l}{\textsci}{\textsubring{b}}{\textlengthmark}{\textsci}} \dicPos{m}[1] \dicFlx{(‑a, ‑ar)}[1] \dicDirectTranslationCS{límec, límeček}
\dicEntry[flikka] \dicTerm{flikk|a} \dicIPA{{f}{l}{\textsci}{h}{\r{g}}{a}} \dicPos{v}[1] \dicFlx{(‑aði)}[1] \dicFlx{acc} \dicPhraseIS{flikka upp á e‑ð} \dicSynonym{lagfæra} \dicDirectTranslationCS{vyspravit (co), dát (čemu) novou tvář}
\dicEntry[flikki] \dicTerm{flikki} \dicIPA{{f}{l}{\textsci}{h}{\r{\textObardotlessj}}{\textsci}} \dicPos{n}[2] \dicFlx{(‑s, ‑)}[16] \dicDirectTranslationCS{flák, (velký) kus (nábytku ap.)}
\dicEntry[flimting] \dicTerm{flimt··ing} \dicIPA{{f}{l}{\textsci}{\textsubring{m}}{\textsubring{d}}{i}{\ng}{\r{g}}} \dicPos{f}[4] \dicFlx{(‑ar, ‑ar)}[5] \dicPhraseIS{hafa e‑ð í flimtingum} \dicDirectTranslationCS{dělat si z~(čeho) legraci}
\dicEntry[flippa] \dicTerm{flipp|a} \dicIPA{{f}{l}{\textsci}{h}{\textsubring{b}}{a}} \dicPos{v}[1] \dicFlx{(‑aði)}[44] \dicLangCat{hovor.} \dicDirectTranslationCS{blbnout, užívat si (narkotik ap.)};  \dicIdiom{flippa}[út]{ \dicPhraseIS{flippa út}} \dicLangCat{hovor.} \dicDirectTranslationCS{zjančit se, ztratit nad sebou kontrolu}
\dicEntry[fliss] \dicTerm{fliss} \dicIPA{{f}{l}{\textsci}{s}{\textlengthmark}} \dicPos{n}[2] \dicFlx{(‑)}[24] \dicDirectTranslationCS{hihňání, chichotání, chichot}
\dicEntry[flissa] \dicTerm{fliss|a} \dicsymFrequent\  \dicIPA{{f}{l}{\textsci}{s}{\textlengthmark}{a}} \dicPos{v}[1] \dicFlx{(‑aði)}[44] \dicDirectTranslationCS{hihňat se, chichotat se} \dicExampleIS{Þær flissuðu eins og smástelpur.} \dicExampleCS{Chichotaly se jako nějaké malé holky.}
\dicEntry[flík] \dicTerm{flík} \dicIPA{{f}{l}{i}{\textlengthmark}{\r{g}}} \dicPos{f}[10] \dicFlx{(‑ar\,/\addthin ‑ur, ‑ur\,/\addthin ‑ar)}[8] \dicSynonym{fat} \dicDirectTranslationCS{věc, část oblečení\,/\addthin oděvu (kalhoty, sukně ap.)} \dicExampleIS{fá sér nýja flík fyrir jólin} \dicExampleCS{pořídit si na Vánoce novou věc}
\dicEntry[flíka] \dicTerm{flík|a} \dicIPA{{f}{l}{i}{\textlengthmark}{\r{g}}{a}} \dicPos{v}[1] \dicFlx{(‑aði)}[44] \dicFlx{dat} \dicDirectTranslationCS{stavět\,/\addthin vystavovat na odiv, dát\,/\addthin dávat najevo} \dicExampleIS{flíka ekki tilfinningum sínum} \dicExampleCS{nedávat své pocity najevo}
\dicEntry[flís] \dicTerm{flís\smash{\textsuperscript{1}}} \dicIPA{{f}{l}{i}{\textlengthmark}{s}} \dicPos{f}[4] \dicFlx{(‑ar, ‑ar)}[1] \textbf{1.} \dicDirectTranslationCS{tříska, odštěpek, úlomek} \dicExampleIS{ná flísinni úr fingrinum} \dicExampleCS{vytáhnout třísku z~prstu}  \textbf{2.} \dicSynonym*{keramikflís} \dicDirectTranslationCS{dlaždice} \dicExampleIS{steinflís} \dicExampleCS{kamenná dlaždice}  \textbf{3.} \dicSynonym*{lítil sneið} \dicDirectTranslationCS{(tenký) plátek} \dicExampleIS{flís af kjöti} \dicExampleCS{plátek masa};  \dicPhraseIS{e‑að fellur eins og flís við rass} \dicLangCat{přen.} \dicDirectTranslationCS{(co) padne jako ulité}
\dicEntry[flís] \dicTerm{flís\smash{\textsuperscript{2}}} \dicIPA{{f}{l}{i}{\textlengthmark}{s}} \dicPos{n}[2] \dicFlx{(‑s)}[2] \dicDirectTranslationCS{flís, fleece} \dicIndirectTranslationCS{(materiál)}
\dicEntry[flísatöng] \dicTerm{flísa··|töng} \dicIPA{{f}{l}{i}{\textlengthmark}{s}{a}{t\smash{\textsuperscript{h}}}{\oe i}{\ng}{\r{g}}} \dicPos{f}[8] \dicFlx{(‑tangar, ‑tengur\,/\addthin ‑tangir)}[9] \dicDirectTranslationCS{pinzeta}
\dicEntry[flíspeysa] \dicTerm{flís··peys|a} \dicIPA{{f}{l}{i}{\textlengthmark}{s}{p\smash{\textsuperscript{h}}}{ei}{s}{a}} \dicPos{f}[1] \dicFlx{(‑u, ‑ur)}[7] \dicDirectTranslationCS{flísová\,/\addthin fleecová mikina}
\dicEntry[fljót] \dicTerm{fljót} \dicIPA{{f}{l}{j}{ou}{\textlengthmark}{\textsubring{d}}} \dicPos{n}[2] \dicFlx{(‑s, ‑)}[5] \dicDirectTranslationCS{(velká) řeka, veletok}
\dicEntry[fljóta] \dicTerm{fljóta} \dicsymFrequent\  \dicIPA{{f}{l}{j}{ou}{\textlengthmark}{\textsubring{d}}{a}} \dicPos{v}[6] \dicFlx{(flýt, flaut, flutum, flyti, flotið)}[103] \textbf{1.} \dicSynonym*{vera á floti} \dicDirectTranslationCS{plout, plovat, vznášet se} \dicExampleIS{Báturinn flýtur á vatninu.} \dicExampleCS{Člun pluje po vodě.}  \textbf{2.} \dicSynonym{streyma} \dicDirectTranslationCS{téct, plynout} \dicExampleIS{Blóð flaut úr sárinu.} \dicExampleCS{Z~rány tekla krev.};  \dicPhraseIS{augu e‑rs fljóta í tárum} \dicDirectTranslationCS{(čí) oči se topí v~slzách};  \dicIdiom{fljóta}[með]{ \dicPhraseIS{fá að fljóta með}} \dicDirectTranslationCS{moci se přidat}; { \dicPhraseIS{fá að fljóta með e‑m}} \dicDirectTranslationCS{moci se svézt s~(kým)}; { \dicPhraseIS{láta e‑ð fljóta með}} \dicDirectTranslationCS{přidat (co), zahrnout (co)}
\dicEntry[fljótandi] \dicTerm{fljót··andi} \dicIPA{{f}{l}{j}{ou}{\textlengthmark}{\textsubring{d}}{a}{n}{\textsubring{d}}{\textsci}} \dicPos{adj}[13] \dicFlx{indecl}[1] \textbf{1.} \dicDirectTranslationCS{tekutý} \dicExampleIS{fljótandi efni} \dicExampleCS{tekutá látka}  \textbf{2.} \dicDirectTranslationCS{plovoucí}
\dicEntry[fljótaskrift] \dicTerm{fljóta··skrift} \dicIPA{{f}{l}{j}{ou}{\textlengthmark}{\textsubring{d}}{a}{s}{\r{g}}{r}{\textsci}{f}{\textsubring{d}}} \dicPos{f}[7] \dicFlx{(‑ar)}[3] \textbf{1.} \dicSynonym*{sérstök skrift} \dicIndirectTranslationCS{druh německého psacího písma, který se rozšířil ve Skandinávii v~17. století}  \textbf{2.} \dicSynonym*{óvandað verk} \dicDirectTranslationCS{odbytá práce, břídilství}
\dicEntry[fljótfær] \dicTerm{fljót··fær} \dicIPA{{f}{l}{j}{ou}{\textsubring{d}}{f}{a}{i}{\textsubring{r}}} \dicPos{adj}[5] \dicFlx{(f ‑)}[8] \dicSynonym{glannalegur} \dicDirectTranslationCS{zbrklý, ukvapený, unáhlený, neuvážený}
\dicEntry[fljótfærni] \dicTerm{fljót··færn|i} \dicIPA{{f}{l}{j}{ou}{\textsubring{d}}{f}{a}{i}{r}{\textsubring{d}}{n}{\textsci}} \dicPos{f}[3] \dicFlx{(‑i)}[3] \dicDirectTranslationCS{zbrklost, ukvapenost, unáhlenost} \dicExampleIS{í fljótfærni} \dicExampleCS{zbrkle}
\dicEntry[fljótlega] \dicTerm{fljót··lega} \dicsymFrequent\  \dicIPA{{f}{l}{j}{ou}{\textlengthmark}{\textsubring{d}}{l}{\textepsilon}{\textbabygamma}{a}} \dicPos{adv} \dicSynonym{bráðum} \dicDirectTranslationCS{brzy, záhy, rychle} \dicExampleIS{það kom fljótlega í ljós að} \dicExampleCS{rychle vyšlo najevo, že}
\begin{xtolerant}{}{1pt}
\dicEntry[fljótlegur] \dicTerm{fljót··legur}\addthinS\dicIPA{{f}{l}{j}{ou}{\textlengthmark}{\textsubring{d}}{l}{\textepsilon}{\textbabygamma}{\textscy}{\textsubring{r}}}\addthinS\dicPos{adj}[1]\dicFlx{}[-8]\addthinS\dicSynonym{fljótur} \dicDirectTranslationCS{rychlý, kvapný, spěšný} \dicExampleIS{fljótlegt verk} \dicExampleCS{rychlá práce}
\end{xtolerant}
\dicEntry[fljótt] \dicTerm{fljót|t} \dicsymFrequent\  \dicIPA{{f}{l}{j}{ou}{h}{\textsubring{d}}} \dicPos{adv} \dicFlx{(comp ‑ar, sup ‑ast)} \dicSynonym*{skjótlega} \dicDirectTranslationCS{rychle, spěšně, pohotově} \dicExampleIS{byrja eins fljótt aftur á æfingunum og hægt er} \dicExampleCS{vrátit se ke cvičení tak rychle, jak to jen jde}
\dicEntry[fljótur] \dicTerm{fljótur} \dicsymFrequent\  \dicIPA{{f}{l}{j}{ou}{\textlengthmark}{\textsubring{d}}{\textscy}{\textsubring{r}}} \dicPos{adj}[1]\dicFlx{}[-1] \dicSynonym{hraður} \dicDirectTranslationCS{rychlý, spěšný, pohotový} \dicExampleIS{Hann var fljótur að borða.} \dicExampleCS{Rychle jedl.};  \dicPhraseIS{vera fljótur á sér} \dicDirectTranslationCS{ukvapit se, být ukvapený}
\dicEntry[fljótvirkur] \dicTerm{fljót··virkur} \dicIPA{{f}{l}{j}{ou}{\textlengthmark}{\textsubring{d}}{v}{\textsci}{\textsubring{r}}{\r{g}}{\textscy}{\textsubring{r}}} \dicPos{adj}[1]\dicFlx{}[-1] \dicSynonym*{hraðvirkur} \dicDirectTranslationCS{rychlý, výkonný, efektivní}
\dicEntry[fljúga] \dicTerm{fljúga} \dicsymFrequent\  \dicIPA{{f}{l}{j}{u}{\textlengthmark}{a}} \dicPos{v}[6] \dicFlx{(flýg, flaug, flugum, flygi, flogið)}[96] \textbf{1.} \dicDirectTranslationCS{letět, létat, poletovat} \dicExampleIS{Við fljúgum til Stokkhólms.} \dicExampleCS{Letíme do Stockholmu.}  \textbf{2.} \dicDirectTranslationCS{letět, pádit};  \dicPhraseIS{tíminn flýgur} \dicLangCat{přen.} \dicDirectTranslationCS{čas letí}  \textbf{3.} \dicPhraseIS{e‑m flaug e‑að í hug} \dicFlx{impers} \dicDirectTranslationCS{(co koho) napadlo, (co komu) přišlo na mysl}  \textbf{4.} \dicPhraseIS{e‑að hefur flogið fyrir} \dicDirectTranslationCS{proslýchá se (co), říká se (co)}
\dicEntry[flog] \dicTerm{flog} \dicIPA{{f}{l}{\textopeno}{\textlengthmark}{x}} \dicPos{n}[2] \dicFlx{(‑s, ‑)}[5] \dicFieldCat{med.} \dicDirectTranslationCS{záchvat} \dicExampleIS{Ég fékk flog.} \dicExampleCS{Dostal jsem záchvat.}
\dicEntry[flogaveiki] \dicTerm{floga··veik|i} \dicIPA{{f}{l}{\textopeno}{\textlengthmark}{\textbabygamma}{a}{v}{ei}{\r{\textObardotlessj}}{\textsci}} \dicPos{f}[3] \dicFlx{(‑i)}[3] \dicFieldCat{med.} \dicDirectTranslationCS{epilepsie, padoucnice}
\dicEntry[flogið] \dicTerm{flogið} \dicIPA{{f}{l}{\textopeno i}{j}{\textlengthmark}{\textsci}{\texttheta}} \dicPos{v} \dicFlx{supin} \dicLink{fljúga}
\dicEntry[flokka] \dicTerm{flokk|a} \dicIPA{{f}{l}{\textopeno}{h}{\r{g}}{a}} \dicPos{v}[1] \dicFlx{(‑aði)}[1] \dicFlx{acc} \dicDirectTranslationCS{roztřídit, klasifikovat, zařadit} \dicExampleIS{flokka heimildirnar nákvæmlega eftir aldri} \dicExampleCS{roztřídit prameny přesně podle stáří};  \dicIdiom{flokkast}{ \dicPhraseIS{flokkast}} \dicFlx{refl} \dicDirectTranslationCS{klasifikovat se, řadit se}
\dicEntry[flokksmaður] \dicTerm{flokks··|maður} \dicIPA{{f}{l}{\textopeno}{h}{\r{g}}{s}{m}{a}{ð}{\textscy}{\textsubring{r}}} \dicPos{m}[13] \dicFlx{(‑manns, ‑menn)}[2] \dicDirectTranslationCS{člen(ka) strany, straník, stranička}
\dicEntry[flokkun] \dicTerm{flokk|un} \dicIPA{{f}{l}{\textopeno}{h}{\r{g}}{\textscy}{\textsubring{n}}} \dicPos{f}[7] \dicFlx{(‑unar)}[9] \dicDirectTranslationCS{(roz)dělení, (roz)třídění, klasifikace} \dicExampleIS{flokkun eftir stærð} \dicExampleCS{roztřídění podle velikosti}
\dicEntry[flokkur] \dicTerm{flokk|ur} \dicsymFrequent\  \dicIPA{{f}{l}{\textopeno}{h}{\r{g}}{\textscy}{\textsubring{r}}} \dicPos{m}[6] \dicFlx{(‑s, ‑ar)}[8] \textbf{1.} \dicSynonym{hópur} \dicDirectTranslationCS{skupina (lidí ap.)} \dicExampleIS{vera í flokki með þeim} \dicExampleCS{být s~nimi ve skupině}  \textbf{2.} \dicSynonym{stjórnmálaflokkur} \dicDirectTranslationCS{(politická) strana} \dicExampleIS{leiða flokkinn í kosningunum} \dicExampleCS{vést stranu do voleb}  \textbf{3.} \dicSynonym*{undirdeild} \dicDirectTranslationCS{třída, kategorie}  \textbf{4.} \dicSynonym{gæðaflokkur} \dicDirectTranslationCS{třída (kvalita), kategorie} \dicExampleIS{fyrsta flokks vara} \dicExampleCS{prvotřídní zboží}  \textbf{5.} \dicFieldCat{biol.} \dicDirectTranslationCS{třída};  \dicIdiom{flokkur}{ \dicPhraseIS{fylla flokk e‑rs}} \dicDirectTranslationCS{hlásit se ke (komu)}
\dicEntry[flos] \dicTerm{flos} \dicIPA{{f}{l}{\textopeno}{\textlengthmark}{s}} \dicPos{n}[2] \dicFlx{(‑s)}[2] \dicDirectTranslationCS{plyš}
\dicEntry[flosna] \dicTerm{flosn|a} \dicIPA{{f}{l}{\textopeno}{s}{n}{a}} \dicPos{v}[1] \dicFlx{(‑aði)}[34] \dicPhraseIS{flosna upp} \dicDirectTranslationCS{vzdát, skončit (kvůli obtížím ap.)} \dicExampleIS{Hann flosnaði upp frá námi.} \dicExampleCS{Vzdal studium.}
\dicEntry[flot] \dicTerm{flot} \dicIPA{{f}{l}{\textopeno}{\textlengthmark}{\textsubring{d}}} \dicPos{n}[2] \dicFlx{(‑s)}[2] \textbf{1.} \dicSynonym*{það að fljóta} \dicDirectTranslationCS{plavení, plování} \dicExampleIS{Báturinn er á floti.} \dicExampleCS{Člun pluje.}  \textbf{2.} \dicSynonym*{bráðin feiti} \dicDirectTranslationCS{rozpuštěný tuk};  \dicPhraseIS{e‑að er á floti} \dicDirectTranslationCS{(co) se topí ve vodě, (co) je zatopené};  \dicPhraseIS{koma e‑u á flot} \dicDirectTranslationCS{zahájit (co), spustit (co) (změny ap.)}
\dicEntry[flotaforingi] \dicTerm{flota··for·ing|i} \dicIPA{{f}{l}{\textopeno}{\textlengthmark}{\textsubring{d}}{a}{f}{\textopeno}{r}{i}{\textltailn}{\r{\textObardotlessj}}{\textsci}} \dicPos{m}[1] \dicFlx{(‑ja, ‑jar)}[14] \dicFieldCat{nám.} \dicDirectTranslationCS{admirál}
\dicEntry[flotamálaráðherra] \dicTerm{flota·mála··ráð·herr|a} \dicIPA{{f}\-{l}\-{\textopeno}\-{\textlengthmark}\-{\textsubring{d}}\-{a}\-{m}\-{au}\-{l}\-{a}\-{r}\-{au}\-{\texttheta}\-{h}\-{\textepsilon}\-{r}\-{a}\-} \dicPos{m}[1] \dicFlx{(‑a, ‑ar)}[17] \dicDirectTranslationCS{ministr(yně) námořnictví}
\dicEntry[flotastöð] \dicTerm{flota··stöð} \dicIPA{{f}{l}{\textopeno}{\textlengthmark}{\textsubring{d}}{a}{s}{\textsubring{d}}{\oe}{\texttheta}} \dicPos{f}[6] \dicFlx{(‑var, ‑var)}[1] \dicFieldCat{nám.} \dicDirectTranslationCS{základna}
\dicEntry[flotaveldi] \dicTerm{flota··veldi} \dicIPA{{f}{l}{\textopeno}{\textlengthmark}{\textsubring{d}}{a}{v}{\textepsilon}{l}{\textsubring{d}}{\textsci}} \dicPos{n}[2] \dicFlx{(‑s, ‑)}[14] \dicDirectTranslationCS{námořní velmoc}
\dicEntry[flotholt] \dicTerm{flot··holt} \dicIPA{{f}{l}{\textopeno}{\textlengthmark}{t\smash{\textsuperscript{h}}}{\textopeno}{\textsubring{l}}{\textsubring{d}}} \dicPos{n}[2] \dicFlx{(‑s, ‑)}[5] \dicDirectTranslationCS{plovák}
\dicEntry[floti] \dicTerm{flot|i} \dicIPA{{f}{l}{\textopeno}{\textlengthmark}{\textsubring{d}}{\textsci}} \dicPos{m}[1] \dicFlx{(‑a, ‑ar)}[1] \textbf{1.} \dicDirectTranslationCS{loďstvo, flotila}  \textbf{2.} \dicDirectTranslationCS{skupina, družstvo} \dicExampleIS{vera í einum flota} \dicExampleCS{být v~jedné skupině}
\dicEntry[flotið] \dicTerm{flotið} \dicIPA{{f}{l}{\textopeno}{\textlengthmark}{\textsubring{d}}{\textsci}{\texttheta}} \dicPos{v} \dicFlx{supin} \dicLink{fljóta}
\dicEntry[flottur] \dicTerm{flottur} \dicsymFrequent\  \dicIPA{{f}{l}{\textopeno}{h}{\textsubring{d}}{\textscy}{\textsubring{r}}} \dicPos{adj}[1]\dicFlx{}[-10] \textbf{1.} \dicSynonym{fallegur} \dicDirectTranslationCS{skvělý, parádní, senzační} \dicExampleIS{flott peysa} \dicExampleCS{skvělý svetr}  \textbf{2.} \dicDirectTranslationCS{krásný, elegantní} \dicExampleIS{flott kona} \dicExampleCS{elegantní žena}
\dicEntry[fló] \dicTerm{fló} \dicIPA{{f}{l}{ou}{\textlengthmark}} \dicPos{f}[9] \dicFlx{(‑ar, flær)}[1] \dicFieldCat{zool.} \dicDirectTranslationCS{blecha} \textit{(l.~{\textLA{Siphonaptera}})}
\dicEntry[flóa] \dicTerm{fló|a\smash{\textsuperscript{1}}} \dicIPA{{f}{l}{ou}{\textlengthmark}{a}} \dicPos{v}[1] \dicFlx{(‑aði)}[1] \dicFlx{acc} \dicLangCat{zast.} \dicDirectTranslationCS{vařit, ohřívat} \dicExampleIS{flóa mjólk} \dicExampleCS{ohřívat mléko}
\dicEntry[flóa] \dicTerm{fló|a\smash{\textsuperscript{2}}} \dicIPA{{f}{l}{ou}{\textlengthmark}{a}} \dicPos{v}[2] \dicFlx{(‑ði, ‑ð)}[121] \dicSynonym{renna\smash{\textsuperscript{3}}} \dicDirectTranslationCS{téct, plynout, plovat, plavat} \dicExampleIS{Hér flóir allt í vatni.} \dicExampleCS{Všechno tady plave ve vodě.};  \dicPhraseIS{augun flóa í tárum} \dicDirectTranslationCS{oči se topí v~slzách};  \dicPhraseIS{það flóir út úr e‑u} \dicFlx{impers} \dicDirectTranslationCS{teče z~(čeho), vytéká z~(čeho)} \dicExampleIS{Það flóir út úr pottinum.} \dicExampleCS{Z~hrnce teče (voda).}
\dicEntry[flóastelkur] \dicTerm{flóa··stelk|ur} \dicIPA{{f}{l}{ou}{\textlengthmark}{a}{s}{\textsubring{d}}{\textepsilon}{\textsubring{l}}{\r{g}}{\textscy}{\textsubring{r}}} \dicPos{m}[6] \dicFlx{(‑s, ‑ar)}[15] \dicFieldCat{zool.} \dicDirectTranslationCS{vodouš bahenní} \textit{(l.~{\textLA{Tringa glareola}})}  \dicsymPhoto\ 
\dicFigure{21203.jpg}{Flóastelkur}{Flóastelkur - Elison, Glen, Biolib, PD}
\dicEntry[flóð] \dicTerm{flóð} \dicIPA{{f}{l}{ou}{\textlengthmark}{\texttheta}} \dicPos{n}[2] \dicFlx{(‑s, ‑)}[5] \textbf{1.} \dicSynonym*{hásævi} \dicDirectTranslationCS{(vysoký) příliv} \dicAntonym{fjara\smash{\textsuperscript{1}}}  \textbf{2.} \dicSynonym*{vatnavöxtur} \dicDirectTranslationCS{záplava, povodeň, potopa}  \textbf{3.} \dicSynonym{snjóflóð} \dicDirectTranslationCS{lavina}
\dicEntry[flóðapuntur] \dicTerm{flóða··punt|ur} \dicIPA{{f}{l}{ou}{\textlengthmark}{ð}{a}{p\smash{\textsuperscript{h}}}{\textscy}{\textsubring{n}}{\textsubring{d}}{\textscy}{\textsubring{r}}} \dicPos{m}[6] \dicFlx{(‑s)}[7] \dicFieldCat{bot.} \dicDirectTranslationCS{zblochan, zblochan vzplývavý} \textit{(l.~{\textLA{Glyceria fluitans}})}  \dicsymPhoto\ 
\dicFigure{ds_image_flodapuntur_0_2.jpg}{Flóðapuntur}{Flóðapuntur - Christian Fischer, GFDL}
\dicEntry[flóðbylgja] \dicTerm{flóð··bylgj|a} \dicIPA{{f}{l}{ou}{ð}{\textsubring{b}}{\textsci}{l}{\r{\textObardotlessj}}{a}} \dicPos{f}[1] \dicFlx{(‑u, ‑ur)}[25] \dicDirectTranslationCS{tsunami}
\dicEntry[flóðgarður] \dicTerm{flóð··garð|ur} \dicIPA{{f}{l}{ou}{ð}{\r{g}}{a}{r}{ð}{\textscy}{\textsubring{r}}} \dicPos{m}[6] \dicFlx{(‑s, ‑ar)}[5] \dicDirectTranslationCS{povodňová hráz, povodňový násep}
\dicEntry[flóðgátt] \dicTerm{flóð··|gátt} \dicIPA{{f}{l}{ou}{ð}{\r{g}}{au}{h}{\textsubring{d}}} \dicPos{f}[7] \dicFlx{(‑gáttar, ‑gáttir\,/\addthin ‑gættir)}[25] \dicDirectTranslationCS{propust, stavidlo}
\dicEntry[flóðhestur] \dicTerm{flóð··hest|ur} \dicIPA{{f}{l}{ou}{\textlengthmark}{\texttheta}{h}{\textepsilon}{s}{\textsubring{d}}{\textscy}{\textsubring{r}}} \dicPos{m}[6] \dicFlx{(‑s, ‑ar)}[4] \dicFieldCat{zool.} \dicDirectTranslationCS{hroch} \textit{(l.~{\textLA{Hippopotamus}})}
\dicEntry[flói] \dicTerm{fló|i} \dicsymFrequent\  \dicIPA{{f}{l}{ou}{\textlengthmark}{\textsci}} \dicPos{m}[1] \dicFlx{(‑a, ‑ar)}[1] \textbf{1.} \dicSynonym*{breiður fjörður} \dicDirectTranslationCS{(široký) záliv, zátoka} \dicExampleIS{í öllum stærri flóum} \dicExampleCS{ve všech větších zálivech};  \dicPhraseIS{Finnski flói} \dicFieldCat{geog.} \dicDirectTranslationCS{Finský záliv}  \textbf{2.} \dicSynonym*{mýrlendi} \dicDirectTranslationCS{mokřina, mokřad}
\dicEntry[flóki] \dicTerm{flók|i} \dicIPA{{f}{l}{ou}{\textlengthmark}{\r{\textObardotlessj}}{\textsci}} \dicPos{m}[1] \dicFlx{(‑a, ‑ar)}[1] \textbf{1.} \dicSynonym{flækja\smash{\textsuperscript{1}}} \dicDirectTranslationCS{spleť, zapletení}  \textbf{2.} \dicSynonym*{ullarflóki} \dicDirectTranslationCS{plst, plsť}  \textbf{3.} \dicSynonym*{dökkt ský} \dicDirectTranslationCS{černý mrak}
\dicEntry[flókinn] \dicTerm{flókinn} \dicsymFrequent\  \dicIPA{{f}{l}{ou}{\textlengthmark}{\r{\textObardotlessj}}{\textsci}{\textsubring{n}}} \dicPos{adj}[6]\dicFlx{}[-2] \dicDirectTranslationCS{složitý, komplikovaný, komplexní} \dicExampleIS{flókin staða} \dicExampleCS{komplikovaná situace}
\dicEntry[flón] \dicTerm{flón} \dicIPA{{f}{l}{ou}{\textlengthmark}{\textsubring{n}}} \dicPos{n}[2] \dicFlx{(‑s, ‑)}[5] \dicSynonym{bjáni} \dicDirectTranslationCS{hlupák, pitomec, trouba}
\dicEntry[flónel] \dicTerm{flónel} \dicIPA{{f}{l}{ou}{\textlengthmark}{n}{\textepsilon}{\textsubring{l}}} \dicPos{n}[2] \dicFlx{(‑s)}[2] \dicDirectTranslationCS{flanel}
\dicEntry[flónska] \dicTerm{flónsk|a} \dicIPA{{f}{l}{ou}{n}{s}{\r{g}}{a}} \dicPos{f}[1] \dicFlx{(‑u)}[5] \dicSynonym{heimska} \dicDirectTranslationCS{hloupost, hlupáctví}
\dicEntry[flór] \dicTerm{flór} \dicIPA{{f}{l}{ou}{\textlengthmark}{\textsubring{r}}} \dicPos{m}[4] \dicFlx{(‑s, ‑ar)}[14] \dicDirectTranslationCS{koryto (na odtoky v~chlévě)};  \dicPhraseIS{moka flórinn} {\textbf{a.}} \dicDirectTranslationCS{kydat hnůj};  {\textbf{b.}} \dicLangCat{přen.} \dicDirectTranslationCS{kydat hnůj, dělat špinavou práci}
\dicEntry[flóra] \dicTerm{flór|a} \dicIPA{{f}{l}{ou}{\textlengthmark}{r}{a}} \dicPos{f}[1] \dicFlx{(‑u, ‑ur)}[7] \dicDirectTranslationCS{flóra} \dicExampleIS{flóra Íslands} \dicExampleCS{flóra Islandu}
\dicEntry[flórgoði] \dicTerm{flór··goð|i} \dicIPA{{f}{l}{ou}{r}{\r{g}}{\textopeno}{ð}{\textsci}} \dicPos{m}[1] \dicFlx{(‑a, ‑ar)}[1] \dicFieldCat{zool.} \dicDirectTranslationCS{potápka žlutorohá} \textit{(l.~{\textLA{Podiceps auritus}})}  \dicsymPhoto\ 
\dicFigure{20700.jpg}{Flórgoði}{Flórgoði - Dewhurst, Donna, Biolib, PD}
\dicEntry[flórsykur] \dicTerm{flór··syk|ur} \dicIPA{{f}{l}{ou}{\textsubring{r}}{s}{\textsci}{\r{g}}{\textscy}{\textsubring{r}}} \dicPos{m}[5] \dicFlx{(‑urs)}[2] \dicDirectTranslationCS{moučkový cukr}
\dicEntry[flóttalegur] \dicTerm{flótta··legur} \dicIPA{{f}{l}{ou}{h}{\textsubring{d}}{a}{l}{\textepsilon}{\textbabygamma}{\textscy}{\textsubring{r}}} \dicPos{adj}[1]\dicFlx{}[-8] \dicSynonym{órólegur} \dicDirectTranslationCS{neklidný, nepokojný} \dicExampleIS{flóttalegt augnaráð} \dicExampleCS{neklidný pohled}
\dicEntry[flóttamaður] \dicTerm{flótta··|maður} \dicIPA{{f}{l}{ou}{h}{\textsubring{d}}{a}{m}{a}{ð}{\textscy}{\textsubring{r}}} \dicPos{m}[13] \dicFlx{(‑manns, ‑menn)}[2] \dicDirectTranslationCS{uprchlík, uprchlice, utečenec, utečenka} \dicExampleIS{pólitískur flóttamaður} \dicExampleCS{politický uprchlík}
\dicEntry[flótti] \dicTerm{flótt|i} \dicsymFrequent\  \dicIPA{{f}{l}{ou}{h}{\textsubring{d}}{\textsci}} \dicPos{m}[1] \dicFlx{(‑a, ‑ar)}[1] \dicDirectTranslationCS{útěk, uprchnutí, únik};  \dicPhraseIS{leggja á flótta} \dicDirectTranslationCS{začít prchat\,/\addthin utíkat};  \dicPhraseIS{stökkva e‑m á flótta} \dicDirectTranslationCS{přinutit (koho) uprchnout};  \dicPhraseIS{vera á flótta} \dicDirectTranslationCS{být na útěku, prchat} \dicExampleIS{vera á flótta undan lögreglunni} \dicExampleCS{prchat před policií}
\dicEntry[flt.] \dicTerm{flt.} \dicPos{zkr} \dicPhraseIS{fleirtala} \dicFieldCat{jaz.} \dicDirectTranslationCS{množné číslo, plurál}
\dicEntry[flug] \dicTerm{flug} \dicsymFrequent\  \dicIPA{{f}{l}{\textscy}{\textlengthmark}{x}} \dicPos{n}[2] \dicFlx{(‑s, ‑)}[5] \textbf{1.} \dicSynonym*{það að fljúga} \dicDirectTranslationCS{(pře)let} \dicExampleIS{flugið til Íslands} \dicExampleCS{let na Island}  \textbf{2.} \dicSynonym{flugpóstur} \dicDirectTranslationCS{letecká pošta} \dicExampleIS{senda e‑ð með flugi} \dicExampleCS{poslat (co) leteckou poštou};  \dicPhraseIS{vera á ferð og flugi} \dicDirectTranslationCS{být na cestách, cestovat, putovat}
\dicEntry[fluga] \dicTerm{flug|a} \dicsymFrequent\  \dicIPA{{f}{l}{\textscy}{\textlengthmark}{\textbabygamma}{a}} \dicPos{f}[1] \dicFlx{(‑u, ‑ur)}[13] \textbf{1.} \dicFieldCat{zool.} \dicDirectTranslationCS{moucha} \textit{(l.~{\textLA{Brachycera}})};  \dicPhraseIS{slá tvær flugur í einu höggi} \dicLangCat{přen.} \dicDirectTranslationCS{zabít dvě mouchy jednou ranou}  \textbf{2.} \dicSynonym{agn} \dicDirectTranslationCS{muška} \dicIndirectTranslationCS{(návnada)};  \dicPhraseIS{veiða á flugu} \dicDirectTranslationCS{chytat na mušku}  \textbf{3.} \dicSynonym{hugmynd} \dicDirectTranslationCS{brouk} \dicIndirectTranslationCS{(utkvělá myšlenka)};  \dicPhraseIS{fá þá flugu í höfuðið að} \dicLangCat{přen.} \dicDirectTranslationCS{mít brouka v~hlavě, že}
\dicEntry[flugafgreiðsla] \dicTerm{flug··af·greiðsl|a} \dicIPA{{f}{l}{\textscy}{\textlengthmark}{\textbabygamma}{a}{v}{\r{g}}{r}{ei}{ð}{s}{\textsubring{d}}{l}{a}} \dicPos{f}[1] \dicFlx{(‑u, ‑ur)}[13] \dicDirectTranslationCS{letištní obsluha}
\dicEntry[flugáhöfn] \dicTerm{flug··á·|höfn} \dicIPA{{f}{l}{\textscy}{\textlengthmark}{\textbabygamma}{au}{h}{\oe}{\textsubring{b}}{\textsubring{n}}} \dicPos{f}[7] \dicFlx{(‑hafnar, ‑hafnir)}[16] \dicDirectTranslationCS{posádka letadla, letecký personál}
\dicEntry[flugáætlun] \dicTerm{flug··á·ætl|un} \dicIPA{{f}{l}{\textscy}{\textlengthmark}{\textbabygamma}{au}{a}{i}{h}{\textsubring{d}}{l}{\textscy}{\textsubring{n}}} \dicPos{f}[7] \dicFlx{(‑unar, ‑anir)}[8] \dicDirectTranslationCS{letový plán\,/\addthin řád}
\dicEntry[flugbátur] \dicTerm{flug··bát|ur} \dicIPA{{f}{l}{\textscy}{\textbabygamma}{\textsubring{b}}{au}{\textsubring{d}}{\textscy}{\textsubring{r}}} \dicPos{m}[6] \dicFlx{(‑s, ‑ar)}[22] \dicDirectTranslationCS{hydroplán}
\dicEntry[flugbraut] \dicTerm{flug··braut} \dicIPA{{f}{l}{\textscy}{\textbabygamma}{\textsubring{b}}{r}{\oe i}{\textsubring{d}}} \dicPos{f}[7] \dicFlx{(‑ar, ‑ir)}[1] \dicDirectTranslationCS{letištní\,/\addthin přistávací\,/\addthin vzletová dráha, ranvej}
\dicEntry[flugdreki] \dicTerm{flug··drek|i} \dicIPA{{f}{l}{\textscy}{\textbabygamma}{\textsubring{d}}{r}{\textepsilon}{\r{\textObardotlessj}}{\textsci}} \dicPos{m}[1] \dicFlx{(‑a, ‑ar)}[1] \dicDirectTranslationCS{drak (papírový ap.)}
\dicEntry[flugeldasýning] \dicTerm{flug·elda··sýn·ing} \dicIPA{{f}{l}{\textscy}{\textlengthmark}{\textbabygamma}{\textepsilon}{l}{\textsubring{d}}{a}{s}{i}{n}{i}{\ng}{\r{g}}} \dicPos{f}[4] \dicFlx{(‑ar, ‑ar)}[5] \dicDirectTranslationCS{ohňostroj}
\dicEntry[flugeldur] \dicTerm{flug··eld|ur} \dicIPA{{f}{l}{\textscy}{\textlengthmark}{\textbabygamma}{\textepsilon}{l}{\textsubring{d}}{\textscy}{\textsubring{r}}} \dicPos{m}[6] \dicFlx{(‑s, ‑ar)}[4] \dicDirectTranslationCS{raketa, petarda, zábavní pyrotechnika}
\dicEntry[flugfarþegi] \dicTerm{flug··far·þeg|i} \dicIPA{{f}{l}{\textscy}{\textbabygamma}{f}{a}{\textsubring{r}}{\texttheta}{ei}{\textsci}} \dicPos{m}[1] \dicFlx{(‑a, ‑ar)}[1] \dicDirectTranslationCS{cestující v~letecké dopravě}
\dicEntry[flugfélag] \dicTerm{flug··fé·|lag} \dicIPA{{f}{l}{\textscy}{\textbabygamma}{f}{j}{\textepsilon}{l}{a}{x}} \dicPos{n}[2] \dicFlx{(‑lags, ‑lög)}[8] \dicDirectTranslationCS{letecká společnost, aerolinie, aerolinka}
\dicEntry[flugfreyja] \dicTerm{flug··freyj|a} \dicIPA{{f}{l}{\textscy}{\textbabygamma}{f}{r}{ei}{j}{a}} \dicPos{f}[1] \dicFlx{(‑u, ‑ur)}[7] \dicDirectTranslationCS{letuška, stevardka}
\dicEntry[flugher] \dicTerm{flug··her} \dicIPA{{f}{l}{\textscy}{\textlengthmark}{\textbabygamma}{h}{\textepsilon}{\textsubring{r}}} \dicPos{m}[9] \dicFlx{(‑s, ‑ir)}[21] \dicDirectTranslationCS{letectvo, letecké síly}
\dicEntry[flughræddur] \dicTerm{flug··hræddur} \dicIPA{{f}{l}{\textscy}{\textbabygamma}{\textsubring{r}}{a}{i}{\textsubring{d}}{\textscy}{\textsubring{r}}} \dicPos{adj}[2]\dicFlx{}[-18] \dicDirectTranslationCS{mající strach z~létání}
\dicEntry[flughræðsla] \dicTerm{flug··hræðsl|a} \dicIPA{{f}{l}{\textscy}{\textbabygamma}{\textsubring{r}}{a}{i}{ð}{s}{\textsubring{d}}{l}{a}} \dicPos{f}[1] \dicFlx{(‑u)}[5] \dicDirectTranslationCS{strach z~létání}
\dicEntry[flughæð] \dicTerm{flug··hæð} \dicIPA{{f}{l}{\textscy}{\textlengthmark}{\textbabygamma}{h}{a}{i}{\texttheta}} \dicPos{f}[7] \dicFlx{(‑ar)}[3] \dicDirectTranslationCS{výška letu}
\dicEntry[flughæfni] \dicTerm{flug··hæfn|i} \dicIPA{{f}{l}{\textscy}{\textlengthmark}{\textbabygamma}{h}{a}{i}{\textsubring{b}}{n}{\textsci}} \dicPos{f}[3] \dicFlx{(‑i)}[3] \dicDirectTranslationCS{způsobilost k~letu}
\dicEntry[flughöfn] \dicTerm{flug··|höfn} \dicIPA{{f}{l}{\textscy}{\textlengthmark}{\textbabygamma}{h}{\oe}{\textsubring{b}}{\textsubring{n}}} \dicPos{f}[7] \dicFlx{(‑hafnar, ‑hafnir)}[16] \dicDirectTranslationCS{(velké) letiště}
\dicEntry[flugkennari] \dicTerm{flug··kenn·ar|i} \dicIPA{{f}{l}{\textscy}{\textbabygamma}{c\smash{\textsuperscript{h}}}{\textepsilon}{n}{a}{r}{\textsci}} \dicPos{m}[1] \dicFlx{(‑a, ‑ar)}[13] \dicDirectTranslationCS{letecký instruktor, letecká instruktorka}
\dicEntry[flugleið] \dicTerm{flug··leið} \dicIPA{{f}{l}{\textscy}{\textbabygamma}{l}{ei}{\texttheta}} \dicPos{f}[7] \dicFlx{(‑ar, ‑ir)}[1] \textbf{1.} \dicSynonym*{leið flugvélar} \dicDirectTranslationCS{letová dráha}  \textbf{2.} \dicSynonym*{áætlunarleið flugvélar} \dicDirectTranslationCS{letecká linka}
\dicEntry[fluglæs] \dicTerm{flug··læs} \dicIPA{{f}{l}{\textscy}{\textbabygamma}{l}{a}{i}{s}} \dicPos{adj}[5]\dicFlx{}[-1] \dicDirectTranslationCS{plynně čtoucí, schopný plynně číst}
\dicEntry[flugmaður] \dicTerm{flug··|maður} \dicIPA{{f}{l}{\textscy}{\textbabygamma}{m}{a}{ð}{\textscy}{\textsubring{r}}} \dicPos{m}[13] \dicFlx{(‑manns, ‑menn)}[2] \dicDirectTranslationCS{pilot(ka)}
\dicEntry[flugmiði] \dicTerm{flug··mið|i} \dicIPA{{f}{l}{\textscy}{\textbabygamma}{m}{\textsci}{ð}{\textsci}} \dicPos{m}[1] \dicFlx{(‑a, ‑ar)}[1] \dicDirectTranslationCS{letenka}
\dicEntry[flugnager] \dicTerm{flugna··ger} \dicIPA{{f}{l}{\textscy}{\r{g}}{n}{a}{\r{\textObardotlessj}}{\textepsilon}{\textsubring{r}}} \dicPos{n}[2] \dicFlx{(‑s)}[2] \dicDirectTranslationCS{hejno\,/\addthin roj mušek}
\dicEntry[flugnanet] \dicTerm{flugna··net} \dicIPA{{f}{l}{\textscy}{\r{g}}{n}{a}{n}{\textepsilon}{\textsubring{d}}} \dicPos{n}[2] \dicFlx{(‑s, ‑)}[5] \dicDirectTranslationCS{síť proti hmyzu, moskytiéra}
\dicEntry[flugnaspaði] \dicTerm{flugna··spað|i} \dicIPA{{f}{l}{\textscy}{\r{g}}{n}{a}{s}{\textsubring{b}}{a}{ð}{\textsci}} \dicPos{m}[1] \dicFlx{(‑a, ‑ar)}[8] \dicDirectTranslationCS{plácačka na mouchy}
\dicEntry[flugpóstur] \dicTerm{flug··póst|ur} \dicIPA{{f}{l}{\textscy}{\textbabygamma}{p\smash{\textsuperscript{h}}}{ou}{s}{\textsubring{d}}{\textscy}{\textsubring{r}}} \dicPos{m}[6] \dicFlx{(‑s)}[7] \dicDirectTranslationCS{letecká pošta}
\dicEntry[flugrán] \dicTerm{flug··rán} \dicIPA{{f}{l}{\textscy}{\textbabygamma}{r}{au}{\textsubring{n}}} \dicPos{n}[2] \dicFlx{(‑s, ‑)}[5] \dicDirectTranslationCS{únos letadla}
\dicEntry[flugrit] \dicTerm{flug··rit} \dicIPA{{f}{l}{\textscy}{\textbabygamma}{r}{\textsci}{\textsubring{d}}} \dicPos{n}[2] \dicFlx{(‑s, ‑)}[5] \dicDirectTranslationCS{leták, brožura}
\dicEntry[flugræningi] \dicTerm{flug··ræn·ing|i} \dicIPA{{f}{l}{\textscy}{\textbabygamma}{r}{a}{i}{n}{i}{\textltailn}{\r{\textObardotlessj}}{\textsci}} \dicPos{m}[1] \dicFlx{(‑ja, ‑jar)}[14] \dicDirectTranslationCS{únosce\,/\addthin únoskyně letadla}
\dicEntry[flugskeyti] \dicTerm{flug··skeyti} \dicIPA{{f}{l}{\textscy}{x}{s}{\r{\textObardotlessj}}{ei}{\textsubring{d}}{\textsci}} \dicPos{n}[2] \dicFlx{(‑s, ‑)}[14] \dicDirectTranslationCS{raketa, střela}
\dicEntry[flugskóli] \dicTerm{flug··skól|i} \dicIPA{{f}{l}{\textscy}{x}{s}{\r{g}}{ou}{l}{\textsci}} \dicPos{m}[1] \dicFlx{(‑a, ‑ar)}[1] \dicDirectTranslationCS{letecká škola}
\dicEntry[flugskýli] \dicTerm{flug··skýli} \dicIPA{{f}{l}{\textscy}{x}{s}{\r{\textObardotlessj}}{i}{l}{\textsci}} \dicPos{n}[2] \dicFlx{(‑s, ‑)}[14] \dicDirectTranslationCS{hangár}
\dicEntry[flugstjóri] \dicTerm{flug··stjór|i} \dicIPA{{f}{l}{\textscy}{x}{s}{\textsubring{d}}{j}{ou}{r}{\textsci}} \dicPos{m}[1] \dicFlx{(‑a, ‑ar)}[1] \dicDirectTranslationCS{kapitán(ka) (letadla)}
\dicEntry[flugstöð] \dicTerm{flug··stöð} \dicIPA{{f}{l}{\textscy}{x}{s}{\textsubring{d}}{\oe}{\texttheta}} \dicPos{f}[6] \dicFlx{(‑var, ‑var)}[1] \dicDirectTranslationCS{letiště, letištní terminál}
\dicEntry[flugsund] \dicTerm{flug··sund} \dicIPA{{f}{l}{\textscy}{x}{s}{\textscy}{n}{\textsubring{d}}} \dicPos{n}[2] \dicFlx{(‑s)}[2] \dicFieldCat{sport.} \dicDirectTranslationCS{motýlek} \dicIndirectTranslationCS{(plavecký styl)}
\dicEntry[flugsveit] \dicTerm{flug··sveit} \dicIPA{{f}{l}{\textscy}{x}{s}{v}{ei}{\textsubring{d}}} \dicPos{f}[7] \dicFlx{(‑ar, ‑ir)}[1] \dicDirectTranslationCS{letka, eskadra}
\dicEntry[flugtak] \dicTerm{flug··|tak} \dicIPA{{f}{l}{\textscy}{\textbabygamma}{t\smash{\textsuperscript{h}}}{a}{\r{g}}} \dicPos{n}[2] \dicFlx{(‑taks, ‑tök)}[8] \dicDirectTranslationCS{vzlet, vzlétnutí}
\dicEntry[flugtími] \dicTerm{flug··tím|i} \dicIPA{{f}{l}{\textscy}{\textbabygamma}{t\smash{\textsuperscript{h}}}{i}{m}{\textsci}} \dicPos{m}[1] \dicFlx{(‑a, ‑ar)}[1] \dicDirectTranslationCS{doba letu}
\dicEntry[flugturn] \dicTerm{flug··turn} \dicIPA{{f}{l}{\textscy}{\textbabygamma}{t\smash{\textsuperscript{h}}}{\textscy}{r}{\textsubring{d}}{\textsubring{n}}} \dicPos{m}[4] \dicFlx{(‑s, ‑ar)}[3] \dicDirectTranslationCS{kontrolní věž}
\dicEntry[flugum] \dicTerm{flugum} \dicIPA{{f}{l}{\textscy}{\textlengthmark}{\textbabygamma}{\textscy}{\textsubring{m}}} \dicPos{v} \dicFlx{ind pf pl 1 pers} \dicLink{fljúga}
\dicEntry[flugumaður] \dicTerm{flugu··|maður} \dicIPA{{f}{l}{\textscy}{\textlengthmark}{\textbabygamma}{\textscy}{m}{a}{ð}{\textscy}{\textsubring{r}}} \dicPos{m}[13] \dicFlx{(‑manns, ‑menn)}[2] \dicDirectTranslationCS{špión(ka), agent(ka)}
\dicEntry[flugumferðarstjóri] \dicTerm{flug··um·ferðar·stjór|i} \dicIPA{{f}\-{l}\-{\textscy}\-{\textlengthmark}\-{\textbabygamma}\-{\textscy}\-{m}\-{f}\-{\textepsilon}\-{\textsubring{r}}\-{ð}\-{a}\-{\textsubring{r}}\-{s}\-{\textsubring{d}}\-{j}\-{ou}\-{r}\-{\textsci}\-} \dicPos{m}[1] \dicFlx{(‑a, ‑ar)}[1] \dicDirectTranslationCS{letový dispečer, letová dispečerka}
\dicEntry[fluguveiði] \dicTerm{flugu··veið|i} \dicIPA{{f}{l}{\textscy}{\textlengthmark}{\textbabygamma}{\textscy}{v}{ei}{ð}{\textsci}} \dicPos{f}[2] \dicFlx{(‑i\,/\addthin ‑ar, ‑ar)}[3] \dicDirectTranslationCS{lovení ryb na mušku}
\dicEntry[flugveður] \dicTerm{flug··veður} \dicIPA{{f}{l}{\textscy}{\textbabygamma}{v}{\textepsilon}{ð}{\textscy}{\textsubring{r}}} \dicPos{n}[2] \dicFlx{(‑s)}[28] \dicFieldCat{let.} \dicDirectTranslationCS{letecké počasí}
\dicEntry[flugvél] \dicTerm{flug··vél} \dicsymFrequent\  \dicIPA{{f}{l}{\textscy}{\textbabygamma}{v}{j}{\textepsilon}{\textsubring{l}}} \dicPos{f}[4] \dicFlx{(‑ar, ‑ar)}[1] \dicDirectTranslationCS{letadlo, letoun} \dicExampleIS{Þetta var talsvert stærri flugvél en sú sem þið höfðuð flogið.} \dicExampleCS{Bylo to značně větší letadlo než to, kterým jste letěli.}
\dicEntry[flugvélamóðurskip] \dicTerm{flug·véla··móður·skip} \dicIPA{{f}\-{l}\-{\textscy}\-{\textbabygamma}\-{v}\-{j}\-{\textepsilon}\-{l}\-{a}\-{m}\-{ou}\-{ð}\-{\textscy}\-{\textsubring{r}}\-{s}\-{\r{\textObardotlessj}}\-{\textsci}\-{\textsubring{b}}\-} \dicPos{n}[2] \dicFlx{(‑s, ‑)}[5] \dicDirectTranslationCS{letadlová loď}
\dicEntry[flugvélstjóri] \dicTerm{flug·vél··stjór|i} \dicIPA{{f}{l}{\textscy}{\textbabygamma}{v}{j}{\textepsilon}{\textsubring{l}}{s}{\textsubring{d}}{j}{ou}{r}{\textsci}} \dicPos{m}[1] \dicFlx{(‑a, ‑ar)}[1] \dicDirectTranslationCS{letecký inženýr, letecká inženýrka}
\dicEntry[flugvirki] \dicTerm{flug··virk|i} \dicIPA{{f}{l}{\textscy}{\textbabygamma}{v}{\textsci}{\textsubring{r}}{\r{\textObardotlessj}}{\textsci}} \dicPos{m}[1] \dicFlx{(‑ja, ‑jar)}[14] \dicDirectTranslationCS{letecký mechanik, letecká mechanička}
\dicEntry[flugvöllur] \dicTerm{flug··|völlur} \dicsymFrequent\  \dicIPA{{f}{l}{\textscy}{\textbabygamma}{v}{\oe}{\textsubring{d}}{l}{\textscy}{\textsubring{r}}} \dicPos{m}[11] \dicFlx{(‑vallar, ‑vellir)}[5] \dicDirectTranslationCS{letiště} \dicExampleIS{Hann kemur beint frá flugvellinum.} \dicExampleCS{Jede rovnou z~letiště.}
\dicEntry[flugþjónn] \dicTerm{flug··þjón|n} \dicIPA{{f}{l}{\textscy}{x}{\texttheta}{j}{ou}{\textsubring{d}}{\textsubring{n}}} \dicPos{m}[6] \dicFlx{(‑s, ‑ar)}[42] \dicDirectTranslationCS{stevard}
\dicEntry[flugþol] \dicTerm{flug··þol} \dicIPA{{f}{l}{\textscy}{x}{\texttheta}{\textopeno}{\textsubring{l}}} \dicPos{n}[2] \dicFlx{(‑s)}[2] \dicDirectTranslationCS{dolet}
\dicEntry[flunkunýr] \dicTerm{flunku··nýr} \dicIPA{{f}{l}{u}{\r{\ng}}{\r{g}}{\textscy}{n}{i}{\textsubring{r}}} \dicPos{adj}[4]\dicFlx{}[-8] \dicDirectTranslationCS{zbrusu nový}
\dicEntry[flutningabifreið] \dicTerm{flutninga··bif·reið} \dicIPA{{f}{l}{\textscy}{h}{\textsubring{d}}{n}{i}{\ng}{\r{g}}{a}{\textsubring{b}}{\textsci}{v}{r}{ei}{\texttheta}} \dicPos{f}[4] \dicFlx{(‑ar, ‑ar\,/\addthin ‑ir)}[30] \dicLangCat{form.} \dicLink{flutningabíll}
\dicEntry[flutningabíll] \dicTerm{flutninga··bíl|l}\dicTerm{, flutningabifreið} \dicIPA{{f}\-{l}\-{\textscy}\-{h}\-{\textsubring{d}}\-{n}\-{i}\-{\ng}\-{\r{g}}\-{a}\-{\textsubring{b}}\-{i}\-{\textsubring{d}}\-{\textsubring{l}}\-} \dicPos{m}[6] \dicFlx{(‑s, ‑ar)}[48] \dicDirectTranslationCS{nákladní auto, náklaďák, dodávka, kamion}
\dicEntry[flutningalest] \dicTerm{flutninga··lest} \dicIPA{{f}{l}{\textscy}{h}{\textsubring{d}}{n}{i}{\ng}{\r{g}}{a}{l}{\textepsilon}{s}{\textsubring{d}}} \dicPos{f}[7] \dicFlx{(‑ar, ‑ir)}[1] \dicDirectTranslationCS{nákladní vlak}
\dicEntry[flutningamiðlari] \dicTerm{flutninga··miðl·ar|i} \dicIPA{{f}{l}{\textscy}{h}{\textsubring{d}}{n}{i}{\ng}{\r{g}}{a}{m}{\textsci}{ð}{l}{a}{r}{\textsci}} \dicPos{m}[1] \dicFlx{(‑a, ‑ar)}[13] \dicDirectTranslationCS{přepravce, přepravkyně}
\dicEntry[flutningaskip] \dicTerm{flutninga··skip} \dicIPA{{f}{l}{\textscy}{h}{\textsubring{d}}{n}{i}{\ng}{\r{g}}{a}{s}{\r{\textObardotlessj}}{\textsci}{\textsubring{b}}} \dicPos{n}[2] \dicFlx{(‑s, ‑)}[5] \dicFieldCat{nám.} \dicDirectTranslationCS{nákladní\,/\addthin obchodní loď}
\dicEntry[flutningsgjald] \dicTerm{flutnings··|gjald} \dicIPA{{f}{l}{\textscy}{h}{\textsubring{d}}{n}{i}{\ng}{s}{\r{\textObardotlessj}}{a}{l}{\textsubring{d}}} \dicPos{n}[2] \dicFlx{(‑gjalds, ‑gjöld)}[8] \dicDirectTranslationCS{přepravné, poplatek za přepravu}
\dicEntry[flutningskostnaður] \dicTerm{flutnings··kost·nað|ur} \dicIPA{{f}\-{l}\-{\textscy}\-{h}\-{\textsubring{d}}\-{n}\-{i}\-{\ng}\-{s}\-{k\smash{\textsuperscript{h}}}\-{\textopeno}\-{s}\-{\textsubring{d}}\-{n}\-{a}\-{ð}\-{\textscy}\-{\textsubring{r}}\-} \dicPos{m}[10] \dicFlx{(‑ar)}[9] \dicDirectTranslationCS{náklady na přepravu}
\dicEntry[flutningur] \dicTerm{flut··ning|ur} \dicIPA{{f}{l}{\textscy}{h}{\textsubring{d}}{n}{i}{\ng}{\r{g}}{\textscy}{\textsubring{r}}} \dicPos{m}[6] \dicFlx{(‑s, ‑ar)}[8] \textbf{1.} \dicDirectTranslationCS{přeprava, doprava, transport} \dicExampleIS{vöruflutningur} \dicExampleCS{přeprava zboží}  \textbf{2.} \dicSynonym{farangur} \dicDirectTranslationCS{náklad (lodní ap.)}  \textbf{3.} \dicDirectTranslationCS{(pře)stěhování}  \textbf{4.} \dicSynonym{upplestur} \dicDirectTranslationCS{přednes (řeči ap.), recitace} \dicExampleIS{flutningur kvæðis} \dicExampleCS{přednes básně}
\dicEntry[flutt] \dicTerm{flutt} \dicIPA{{f}{l}{\textscy}{h}{\textsubring{d}}} \dicPos{v} \dicFlx{supin} \dicLink{flytja}
\dicEntry[flutti] \dicTerm{flutti} \dicIPA{{f}{l}{\textscy}{h}{\textsubring{d}}{\textsci}} \dicPos{v} \dicFlx{ind pf sg 1 pers} \dicLink{flytja}
\dicEntry[fluttum] \dicTerm{fluttum} \dicIPA{{f}{l}{\textscy}{h}{\textsubring{d}}{\textscy}{\textsubring{m}}} \dicPos{v} \dicFlx{ind pf pl 1 pers} \dicLink{flytja}
\dicEntry[flutum] \dicTerm{flutum} \dicIPA{{f}{l}{\textscy}{\textlengthmark}{\textsubring{d}}{\textscy}{\textsubring{m}}} \dicPos{v} \dicFlx{ind pf pl 1 pers} \dicLink{fljóta}
\dicEntry[flúð] \dicTerm{flúð} \dicIPA{{f}{l}{u}{\textlengthmark}{\texttheta}} \dicPos{f}[7] \dicFlx{(‑ar, ‑ir)}[1] \dicDirectTranslationCS{peřej, slap}
\dicEntry[flúðasigling] \dicTerm{flúða··sigl·ing} \dicIPA{{f}{l}{u}{\textlengthmark}{ð}{a}{s}{\textsci}{\r{g}}{l}{i}{\ng}{\r{g}}} \dicPos{f}[4] \dicFlx{(‑ar, ‑ar)}[5] \dicFieldCat{sport.} \dicDirectTranslationCS{rafting}
\dicEntry[flúði] \dicTerm{flúði} \dicIPA{{f}{l}{u}{\textlengthmark}{ð}{\textsci}} \dicPos{v} \dicFlx{ind pf sg 1 pers} \dicLink{flýja}
\dicEntry[flúðum] \dicTerm{flúðum} \dicIPA{{f}{l}{u}{\textlengthmark}{ð}{\textscy}{\textsubring{m}}} \dicPos{v} \dicFlx{ind pf pl 1 pers} \dicLink{flýja}
\dicEntry[flúið] \dicTerm{flúið} \dicIPA{{f}{l}{u}{\textsci}{\texttheta}} \dicPos{v} \dicFlx{supin} \dicLink{flýja}
\dicEntry[flúor] \dicTerm{flúor}\dicTerm{, flúr\smash{\textsuperscript{1}}} \dicIPA{{f}\-{l}\-{u}\-{\textopeno}\-{\textsubring{r}}\-} \dicPos{n}[2] \dicFlx{(‑s)}[33] \dicFieldCat{chem.} \dicDirectTranslationCS{fluor} \textit{(l.~{\textLA{F, Fluorum}})}
\dicEntry[flúr] \dicTerm{flúr\smash{\textsuperscript{1}}} \dicIPA{{f}{l}{u}{\textlengthmark}{\textsubring{r}}} \dicPos{m}[4] \dicFlx{(‑s)}[27] \dicLink{flúor}
\dicEntry[flúr] \dicTerm{flúr\smash{\textsuperscript{2}}} \dicIPA{{f}{l}{u}{\textlengthmark}{\textsubring{r}}} \dicPos{n}[2] \dicFlx{(‑s)}[2] \dicDirectTranslationCS{ornament, okrasa}
\dicEntry[flyðra] \dicTerm{flyðr|a} \dicIPA{{f}{l}{\textsci}{ð}{r}{a}} \dicPos{f}[1] \dicFlx{(‑u, ‑ur)}[7] \dicFieldCat{zool.} \dicDirectTranslationCS{platýs, platýs obecný} \textit{(l.~{\textLA{Hippoglossus hippoglossus}})}
\dicEntry[flygi] \dicTerm{flygi} \dicIPA{{f}{l}{\textsci}{j}{\textlengthmark}{\textsci}} \dicPos{v} \dicFlx{con pf sg 1 pers} \dicLink{fljúga}
\dicEntry[flygill] \dicTerm{flyg|ill} \dicIPA{{f}{l}{i}{j}{\textlengthmark}{\textsci}{\textsubring{d}}{\textsubring{l}}} \dicPos{m}[6] \dicFlx{(‑ils, ‑lar)}[35] \dicFieldCat{hud.} \dicDirectTranslationCS{(koncertní) křídlo, klavír}
\dicEntry[flygsa] \dicTerm{flygs|a} \dicIPA{{f}{l}{\textsci}{x}{s}{a}} \dicPos{f}[1] \dicFlx{(‑u, ‑ur)}[13] \dicSynonym{tætla} \dicDirectTranslationCS{útržek, ústřižek, cár (látky ap.)}
\dicEntry[flykkjast] \dicTerm{flykk|jast} \dicIPA{{f}{l}{\textsci}{h}{\r{\textObardotlessj}}{a}{s}{\textsubring{d}}} \dicPos{v}[2] \dicFlx{(‑tist, ‑st)}[194] \dicFlx{refl} \dicDirectTranslationCS{shromáždit se, houfovat se, shluknout se, shlukovat se} \dicExampleIS{Fólkið flykktist út á götu.} \dicExampleCS{Lidé se shromáždili na ulici.}
\dicEntry[flyksa] \dicTerm{flyks|a} \dicIPA{{f}{l}{\textsci}{\r{g}}{s}{a}} \dicPos{f}[1] \dicFlx{(‑u, ‑ur)}[13] \dicDirectTranslationCS{vločka (sněhová ap.)}
\dicEntry[flysja] \dicTerm{flysj|a} \dicIPA{{f}{l}{\textsci}{\textlengthmark}{s}{j}{a}} \dicPos{v}[1] \dicFlx{(‑aði)}[1] \dicFlx{acc} \dicSynonym{afhýða} \dicDirectTranslationCS{oloupat, oškrábat} \dicExampleIS{flysja kartöflur} \dicExampleCS{oloupat brambory}
\dicEntry[flysjaður] \dicTerm{flysj··|aður} \dicIPA{{f}{l}{\textsci}{\textlengthmark}{s}{j}{a}{ð}{\textscy}{\textsubring{r}}} \dicPos{adj}[3] \dicFlx{(f ‑uð)}[4] \dicDirectTranslationCS{oloupaný, oškrábaný} \dicExampleIS{flysjaðar kartöflur} \dicExampleCS{oloupané brambory}
\dicEntry[flyt] \dicTerm{flyt} \dicIPA{{f}{l}{\textsci}{\textlengthmark}{\textsubring{d}}} \dicPos{v} \dicFlx{ind praes sg 1 pers} \dicLink{flytja}
\dicEntry[flyti] \dicTerm{flyti} \dicIPA{{f}{l}{\textsci}{\textlengthmark}{\textsubring{d}}{\textsci}} \dicPos{v} \dicFlx{con pf sg 1 pers} \dicLink{fljóta}
\dicEntry[flytja] \dicTerm{flytja} \dicsymFrequent\  \dicIPA{{f}{l}{\textsci}{\textlengthmark}{\textsubring{d}}{j}{a}} \dicPos{v}[4] \dicFlx{(flyt, flutti, fluttum, flytti, flutt)}[62] \dicFlx{acc} \textbf{1.} \dicDirectTranslationCS{přemístit, přesunout, přestěhovat, transportovat} \dicExampleIS{flytja vörur} \dicExampleCS{přemístit zboží}  \textbf{2.} \dicDirectTranslationCS{(pře)stěhovat se} \dicExampleIS{Við flytjum í nýja íbúð.} \dicExampleCS{Stěhujeme se do nového bytu.}  \textbf{3.} \dicSynonym*{færa fram} \dicDirectTranslationCS{přednést, přednášet, pronést, ztvárnit (roli ap.)} \dicExampleIS{flytja kvæði} \dicExampleCS{přednést báseň};  \dicIdiom{flytja}[inn]{ \dicPhraseIS{flytja inn e‑ð}} \dicDirectTranslationCS{dovážet (co) (zboží ap.), importovat (co)};  \dicPhraseIS{flytja út e‑ð} \dicDirectTranslationCS{vyvážet (co) (zboží ap.), exportovat (co)};  \dicIdiom{flytjast}{ \dicPhraseIS{flytjast}} \dicFlx{refl} {\textbf{a.}} \dicDirectTranslationCS{(pře)stěhovat se, přesunout se};  {\textbf{b.}} \dicDirectTranslationCS{dovážet se}
\dicEntry[flytjandi] \dicTerm{flytj··|andi} \dicIPA{{f}{l}{\textsci}{\textlengthmark}{\textsubring{d}}{j}{a}{n}{\textsubring{d}}{\textsci}} \dicPos{m}[2] \dicFlx{(‑anda, ‑endur)}[1] \textbf{1.} \dicDirectTranslationCS{účinkující, vystupující, aktér(ka)}  \textbf{2.} \dicDirectTranslationCS{přednášející, řečník, řečnice}
\dicEntry[flytti] \dicTerm{flytti} \dicIPA{{f}{l}{\textsci}{h}{\textsubring{d}}{\textsci}} \dicPos{v} \dicFlx{con pf sg 1 pers} \dicLink{flytja}
\dicEntry[flý] \dicTerm{flý} \dicIPA{{f}{l}{i}{\textlengthmark}} \dicPos{v} \dicFlx{ind praes sg 1 pers} \dicLink{flýja}
\dicEntry[flýði] \dicTerm{flýði} \dicIPA{{f}{l}{i}{\textlengthmark}{ð}{\textsci}} \dicPos{v} \dicFlx{ind\,/\addthin con pf sg 1 pers} \dicLink{flýja}
\dicEntry[flýðum] \dicTerm{flýðum} \dicIPA{{f}{l}{i}{\textlengthmark}{ð}{\textscy}{\textsubring{m}}} \dicPos{v} \dicFlx{ind pf pl 1 pers} \dicLink{flýja}
\dicEntry[flýg] \dicTerm{flýg} \dicIPA{{f}{l}{i}{\textlengthmark}{x}} \dicPos{v} \dicFlx{ind praes sg 1 pers} \dicLink{fljúga}
\dicEntry[flýja] \dicTerm{flýja} \dicsymFrequent\  \dicIPA{{f}{l}{i}{j}{\textlengthmark}{a}} \dicPos{v}[4] \dicFlx{(flý, flúði\,/\addthin flýði, flúðum\,/\addthin flýðum, flýði, flúið)}[53] \dicFlx{acc} \dicDirectTranslationCS{utéct, utíkat, uniknout, unikat, uprchnout, prchat} \dicExampleIS{flýja burt} \dicExampleCS{utéct pryč}
\dicEntry[flýt] \dicTerm{flýt} \dicIPA{{f}{l}{i}{\textlengthmark}{\textsubring{d}}} \dicPos{v} \dicFlx{ind praes sg 1 pers} \dicLink{fljóta}
\dicEntry[flýta] \dicTerm{flýt|a} \dicsymFrequent\  \dicIPA{{f}{l}{i}{\textlengthmark}{\textsubring{d}}{a}} \dicPos{v}[2] \dicFlx{(‑ti, ‑t)}[57] \dicFlx{dat} \dicDirectTranslationCS{zrychlit, urychlit, popohnat};  \dicPhraseIS{flýta sér} \dicDirectTranslationCS{spěchat, pospíchat, pospíšit si};  \dicPhraseIS{klukkan flýtir sér} \dicDirectTranslationCS{hodiny se předcházejí, hodiny jdou napřed};  \dicIdiom{flýta}[fyrir]{ \dicPhraseIS{flýta fyrir e‑m}} \dicDirectTranslationCS{pomoct (komu) (aby byl dřív hotový ap.)}; { \dicPhraseIS{flýta fyrir e‑u}} \dicDirectTranslationCS{uspíšit (co), pomoci (čemu)}
\dicEntry[flýtir] \dicTerm{flýt|ir} \dicIPA{{f}{l}{i}{\textlengthmark}{\textsubring{d}}{\textsci}{\textsubring{r}}} \dicPos{m}[7] \dicFlx{(‑is)}[2] \dicSynonym{hraði} \dicDirectTranslationCS{spěch, shon};  \dicPhraseIS{í flýti} \dicFlx{adv} \dicSynonym*{í hasti} \dicDirectTranslationCS{ve spěchu}
\dicEntry[flæða] \dicTerm{flæ|ða} \dicIPA{{f}{l}{a}{i}{\textlengthmark}{ð}{a}} \dicPos{v}[2] \dicFlx{(‑ddi, ‑tt)}[179] \textbf{1.} \dicDirectTranslationCS{zaplavit, zatopit};  \dicPhraseIS{það flæðir e‑s staðar} \dicFlx{impers} \dicDirectTranslationCS{voda teče (kam), (co) je zaplavené} \dicExampleIS{Það flæðir inn í kjallarann.} \dicExampleCS{Sklep je zatopený.}  \textbf{2.} \dicDirectTranslationCS{proudit, plynout}  \textbf{3.} \dicPhraseIS{það flæðir að} \dicFlx{impers} \dicDirectTranslationCS{nastává (vysoký) příliv}
\dicEntry[flæðarmál] \dicTerm{flæðar··mál} \dicIPA{{f}{l}{a}{i}{ð}{a}{r}{m}{au}{\textsubring{l}}} \dicPos{n}[2] \dicFlx{(‑s)}[2] \dicDirectTranslationCS{pobřeží, přílivové pobřežní pásmo}
\dicEntry[flæði] \dicTerm{flæði} \dicIPA{{f}{l}{a}{i}{\textlengthmark}{ð}{\textsci}} \dicPos{n}[2] \dicFlx{(‑s)}[20] \textbf{1.} \dicSynonym{flóð} \dicDirectTranslationCS{příliv}  \textbf{2.} \dicDirectTranslationCS{příliv, tok (pracovní síly ap.)}  \textbf{3.} \dicFieldCat{fyz.} \dicDirectTranslationCS{tok}
\dicEntry[flæðisker] \dicTerm{flæði··sker} \dicIPA{{f}{l}{a}{i}{\textlengthmark}{ð}{\textsci}{s}{\r{\textObardotlessj}}{\textepsilon}{\textsubring{r}}} \dicPos{n}[2] \dicFlx{(‑s, ‑)}[13] \dicDirectTranslationCS{úskalí, útes zaplavovaný při přílivu};  \dicPhraseIS{vera ekki á flæðiskeri staddur} \dicLangCat{přen.} \dicDirectTranslationCS{nezůstat na suchu, být za vodou (finančně)}
\dicEntry[flækingsköttur] \dicTerm{flæk·ings··|köttur} \dicIPA{{f}{l}{a}{i}{\textlengthmark}{\r{\textObardotlessj}}{i}{\ng}{s}{k\smash{\textsuperscript{h}}}{\oe}{h}{\textsubring{d}}{\textscy}{\textsubring{r}}} \dicPos{m}[11] \dicFlx{(‑kattar, ‑kettir)}[5] \dicDirectTranslationCS{toulavá kočka}
\dicEntry[flækingur] \dicTerm{flæk··ing|ur} \dicIPA{{f}{l}{a}{i}{\textlengthmark}{\r{\textObardotlessj}}{i}{\ng}{\r{g}}{\textscy}{\textsubring{r}}} \dicPos{m}[6] \dicFlx{(‑s, ‑ar)}[8] \textbf{1.} \dicSynonym{flakkari} \dicDirectTranslationCS{tulák, tulačka, vandrák, vandračka}  \textbf{2.} \dicSynonym{flakk} \dicDirectTranslationCS{tuláctví, potulka};  \dicPhraseIS{það er flækingur á e‑m} \dicFlx{impers} \dicDirectTranslationCS{(kdo) je neustále na cestách};  \dicPhraseIS{vera á flækingi} \dicDirectTranslationCS{vést tulácký život}
\dicEntry[flækja] \dicTerm{flækj|a\smash{\textsuperscript{1}}} \dicIPA{{f}{l}{a}{i}{\textlengthmark}{\r{\textObardotlessj}}{a}} \dicPos{f}[1] \dicFlx{(‑u, ‑ur)}[7] \textbf{1.} \dicSynonym{benda\smash{\textsuperscript{1}}} \dicDirectTranslationCS{zapeklitost, komplikace, spletitost}  \textbf{2.} \dicSynonym*{flókið band} \dicDirectTranslationCS{motanice, propletenec}  \textbf{3.} \dicFieldCat{bot.} \dicSynonym{umfeðmingsgras} \dicDirectTranslationCS{vikev ptačí} \textit{(l.~{\textLA{Vicia cracca}})}  \dicsymPhoto\   \textbf{4.} \dicFieldCat{lit.} \dicDirectTranslationCS{zápletka}
\dicFigure{317.jpg}{Flækja}{Flækja - Zicha Ondřej, Biolib, Copyright/CC-BY-NC}
\dicEntry[flækja] \dicTerm{flæk|ja\smash{\textsuperscript{2}}} \dicsymFrequent\  \dicIPA{{f}{l}{a}{i}{\textlengthmark}{\r{\textObardotlessj}}{a}} \dicPos{v}[2] \dicFlx{(‑ti, ‑t)}[25] \dicFlx{acc} \textbf{1.} \dicDirectTranslationCS{zamotat, zaplést, zašmodrchat} \dicExampleIS{flækja bandið} \dicExampleCS{zamotat provázek}  \textbf{2.} \dicDirectTranslationCS{zamotat, (z)komplikovat} \dicExampleIS{flækja málið} \dicExampleCS{zkomplikovat případ};  \dicPhraseIS{flækja sig} \dicDirectTranslationCS{zamotat se, zaplést se (do maléru ap.)};  \dicIdiom{flækjast}{ \dicPhraseIS{flækjast}} \dicFlx{refl} {\textbf{a.}} \dicDirectTranslationCS{zamotat se, zaplést se (provázek ap.)};  {\textbf{b.}} \dicDirectTranslationCS{zamotat se, zkomplikovat se};  {\textbf{c.}} \dicDirectTranslationCS{toulat se, potulovat se};  \dicIdiom{flækjast}[fyrir]{ \dicPhraseIS{flækjast fyrir}} \dicFlx{refl} \dicDirectTranslationCS{(při)plést se (do cesty)};  \dicIdiom{flækjast}[í]{ \dicPhraseIS{flækjast í e‑ð}} \dicFlx{refl} \dicDirectTranslationCS{zaplést se do (čeho), zamotat se do (čeho)}
\dicEntry[flæma] \dicTerm{flæm|a} \dicIPA{{f}{l}{a}{i}{\textlengthmark}{m}{a}} \dicPos{v}[2] \dicFlx{(‑di, ‑t)}[141] \dicFlx{acc} \dicSynonym{hrekja} \dicDirectTranslationCS{vyhnat, vyhánět, vypudit, vystrnadit} \dicExampleIS{flæma e‑n burt} \dicExampleCS{vypudit (koho) pryč};  \dicIdiom{flæmast}{ \dicPhraseIS{flæmast}} \dicFlx{refl} \dicDirectTranslationCS{být vyhnán\,/\addthin vypuzen}
\dicEntry[flæmingi] \dicTerm{flæm··ing|i\smash{\textsuperscript{1}}} \dicIPA{{f}{l}{a}{i}{\textlengthmark}{m}{i}{\textltailn}{\r{\textObardotlessj}}{\textsci}} \dicPos{m}[1] \dicFlx{(‑ja, ‑jar)}[14] \dicFieldCat{zool.} \dicDirectTranslationCS{plameňák, plameňák růžový} \textit{(l.~{\textLA{Phoenicopterus ruber}})}
\dicEntry[Flæmingi] \dicTerm{Flæm··ing|i\smash{\textsuperscript{2}}} \dicIPA{{f}{l}{a}{i}{\textlengthmark}{m}{i}{\textltailn}{\r{\textObardotlessj}}{\textsci}} \dicPos{m}[1] \dicFlx{(‑ja, ‑jar)}[14] \dicDirectTranslationCS{Vlám(ka)}
\dicEntry[flæmska] \dicTerm{flæmsk|a} \dicIPA{{f}{l}{a}{i}{m}{s}{\r{g}}{a}} \dicPos{f}[1] \dicFlx{(‑u)}[5] \dicDirectTranslationCS{vlámština}
\dicEntry[flæmskur] \dicTerm{flæmskur} \dicIPA{{f}{l}{a}{i}{m}{s}{\r{g}}{\textscy}{\textsubring{r}}} \dicPos{adj}[1]\dicFlx{}[-1] \dicDirectTranslationCS{vlámský}
\dicEntry[flær] \dicTerm{flær} \dicIPA{{f}{l}{a}{i}{\textlengthmark}{\textsubring{r}}} \dicPos{f} \dicFlx{pl nom} \dicLink{fló}
\dicEntry[flærð] \dicTerm{flærð} \dicIPA{{f}{l}{a}{i}{r}{\texttheta}} \dicPos{f}[7] \dicFlx{(‑ar, ‑ir)}[1] \dicDirectTranslationCS{faleš, proradnost}
\dicEntry[flög] \dicTerm{flög} \dicIPA{{f}{l}{\oe}{\textlengthmark}{x}} \dicPos{n} \dicFlx{pl nom} \dicLink{flag}
\dicEntry[flögð] \dicTerm{flögð} \dicIPA{{f}{l}{\oe}{\textbabygamma}{\texttheta}} \dicPos{n} \dicFlx{pl nom} \dicLink{flagð}
\dicEntry[flögg] \dicTerm{flögg} \dicIPA{{f}{l}{\oe}{\r{g}}{\textlengthmark}} \dicPos{n} \dicFlx{pl nom} \dicLink{flagg}
\dicEntry[flögra] \dicTerm{flögr|a} \dicIPA{{f}{l}{\oe}{\textbabygamma}{r}{a}} \dicPos{v}[1] \dicFlx{(‑aði)}[34] \textbf{1.} \dicDirectTranslationCS{poletovat, vznášet se} \dicExampleIS{Fuglinn flögraði lengi yfir húsinu.} \dicExampleCS{Ptáček dlouho poletoval nad domem.}  \textbf{2.} \dicPhraseIS{það flögrar að e‑m að (gera e‑ð)} \dicFlx{impers} \dicDirectTranslationCS{(komu) prolétne hlavou (udělat (co))}
\dicEntry[flögu] \dicTerm{flögu} \dicIPA{{f}{l}{\oe}{\textlengthmark}{\textbabygamma}{\textscy}} \dicPos{f} \dicFlx{sg gen} \dicLink{flaga}
\dicEntry[flök] \dicTerm{flök} \dicIPA{{f}{l}{\oe}{\textlengthmark}{\r{g}}} \dicPos{n} \dicFlx{pl nom} \dicLink{flak}
\dicEntry[flökra] \dicTerm{flökr|a} \dicIPA{{f}{l}{\oe}{\textlengthmark}{\r{g}}{r}{a}} \dicPos{v}[1] \dicFlx{(‑aði)}[85] \dicFlx{impers} \dicPhraseIS{það flökrar að e‑m} \dicDirectTranslationCS{(komu) se mihne hlavou, (koho) napadá (myšlenka ap.)};  \dicPhraseIS{e‑n\,/\addthin e‑m flökrar við e‑u} \dicDirectTranslationCS{(komu) se dělá nevolno z~(čeho), (komu) je špatně z~(čeho)} \dicExampleIS{Mig flökrar við matnum.} \dicExampleCS{Dělá se mi nevolno z~jídla.}
\dicEntry[flökta] \dicTerm{flökt|a} \dicIPA{{f}{l}{\oe}{x}{\textsubring{d}}{a}} \dicPos{v}[2] \dicFlx{(‑i, ‑)}[23] \textbf{1.} \dicDirectTranslationCS{míhat se, mihotat se (plamen ap.)} \dicExampleIS{Ljósið flöktir.} \dicExampleCS{Světlo se mihotá.}  \textbf{2.} \dicDirectTranslationCS{těkat, toulat se, bloudit}
\dicEntry[flökun] \dicTerm{flök|un} \dicIPA{{f}{l}{\oe}{\textlengthmark}{\r{g}}{\textscy}{\textsubring{n}}} \dicPos{f}[7] \dicFlx{(‑unar)}[12] \dicDirectTranslationCS{výroba filet} \dicExampleIS{vinna við flökun} \dicExampleCS{pracovat při výrobě filet}
\dicEntry[flösku] \dicTerm{flösku} \dicIPA{{f}{l}{\oe}{s}{\r{g}}{\textscy}} \dicPos{f} \dicFlx{sg gen} \dicLink{flaska}
\dicEntry[flöt] \dicTerm{flöt} \dicIPA{{f}{l}{\oe}{\textlengthmark}{\textsubring{d}}} \dicPos{f}[7] \dicFlx{(flatar, flatir)}[16] \dicDirectTranslationCS{rovinka, rovná zem, plácek} \dicExampleIS{reisa tjald á flötinni} \dicExampleCS{postavit stan na rovince}
\dicEntry[flötur] \dicTerm{flötur} \dicIPA{{f}{l}{\oe}{\textlengthmark}{\textsubring{d}}{\textscy}{\textsubring{r}}} \dicPos{m}[11] \dicFlx{(flatar, fletir)}[5] \textbf{1.} \dicSynonym{grund} \dicDirectTranslationCS{rovina, planina} \dicExampleIS{Borðið stendur á sléttum fleti.} \dicExampleCS{Stůl stojí na rovném místě.}  \textbf{2.} \dicSynonym{grunn\-ur\smash{\textsuperscript{1}}} \dicDirectTranslationCS{pozadí} \dicExampleIS{á gulum fleti} \dicExampleCS{na žlutém pozadí}  \textbf{3.} \dicDirectTranslationCS{rovina, stránka, aspekt (věci ap.)} \dicExampleIS{nýr flötur á málinu} \dicExampleCS{nová rovina v~případu}
\dicEntry[fn.] \dicTerm{fn.} \dicPos{zkr} \dicPhraseIS{fornafn} \dicFieldCat{jaz.} \dicDirectTranslationCS{zájmeno}
\dicEntry[fok] \dicTerm{fok} \dicIPA{{f}{\textopeno}{\textlengthmark}{\r{g}}} \dicPos{n}[2] \dicFlx{(‑s)}[2] \dicDirectTranslationCS{větrná eroze, denudace}
\dicEntry[fokheldur] \dicTerm{fok··heldur} \dicIPA{{f}{\textopeno}{\textlengthmark}{k\smash{\textsuperscript{h}}}{\textepsilon}{l}{\textsubring{d}}{\textscy}{\textsubring{r}}} \dicPos{adj}[2]\dicFlx{}[-14] \dicDirectTranslationCS{(o~domě ap.) chráněný před větrem a~vodou}
\dicEntry[fokið] \dicTerm{fokið} \dicIPA{{f}{\textopeno}{\textlengthmark}{\r{\textObardotlessj}}{\textsci}{\texttheta}} \dicPos{v} \dicFlx{supin} \dicLink{fjúka}
\dicEntry[fokka] \dicTerm{fokk|a} \dicIPA{{f}{\textopeno}{h}{\r{g}}{a}} \dicPos{f}[1] \dicFlx{(‑u, ‑ur)}[19] \dicFieldCat{nám.} \dicDirectTranslationCS{kosatka, stěhovka} \dicIndirectTranslationCS{(trojúhelníková lodní plachta)}
\dicEntry[fokreiður] \textls[15]{\dicTerm{fok··reiður} \dicIPA{{f}{\textopeno}{\textlengthmark}{\r{g}}{r}{ei}{ð}{\textscy}{\textsubring{r}}} \dicPos{adj}[2]\dicFlx{}[-10] \dicDirectTranslationCS{rozzuřený, rozlícený, vzteklý}}
\dicEntry[folald] \dicTerm{fol··|ald} \dicIPA{{f}{\textopeno}{\textlengthmark}{l}{a}{l}{\textsubring{d}}} \dicPos{n}[2] \dicFlx{(‑alds, ‑öld)}[8] \dicDirectTranslationCS{hříbě}
\dicEntry[fold] \dicTerm{fold} \dicIPA{{f}{\textopeno}{l}{\textsubring{d}}} \dicPos{f}[7] \dicFlx{(‑ar, ‑ir)}[1] \dicLangCat{básn.} \dicSynonym{land} \dicDirectTranslationCS{zem, země}
\dicEntry[foli] \dicTerm{fol|i} \dicIPA{{f}{\textopeno}{\textlengthmark}{l}{\textsci}} \dicPos{m}[1] \dicFlx{(‑a, ‑ar)}[1] \textbf{1.} \dicDirectTranslationCS{mladý kůň}  \textbf{2.} \dicLangCat{slang.} \dicDirectTranslationCS{hřebec} \dicIndirectTranslationCS{(mladý, neprovdaný muž k~sexuálním hrátkám)}
\dicEntry[fontur] \dicTerm{font|ur} \dicIPA{{f}{\textopeno}{\textsubring{n}}{\textsubring{d}}{\textscy}{\textsubring{r}}} \dicPos{m}[6] \dicFlx{(‑s, ‑ar)}[4] \dicFieldCat{polygr.} \dicSynonym{leturgerð} \dicDirectTranslationCS{font}
\dicEntry[for] \dicTerm{for\smash{\textsuperscript{1}}} \dicIPA{{f}{\textopeno}{\textlengthmark}{\textsubring{r}}} \dicPos{f}[7] \dicFlx{(‑ar, ‑ir)}[1] \dicSynonym{bleyta\smash{\textsuperscript{1}}} \dicDirectTranslationCS{bláto, bahno}
\dicEntry[for] \dicTerm{for-\smash{\textsuperscript{2}}} \dicIPA{{f}{\textopeno}{r}} \dicPos{predp} \dicDirectTranslationCS{před-} \dicExampleIS{forsala} \dicExampleCS{předprodej}
\dicEntry[forað] \dicTerm{for|að} \dicIPA{{f}{\textopeno}{\textlengthmark}{r}{a}{\texttheta}} \dicPos{n}[3] \dicFlx{(‑aðs, ‑öð\,/\addthin ‑uð)}[4] \dicSynonym{fen} \dicDirectTranslationCS{močál, bahnisko, bažina};  \dicPhraseIS{etja e‑m á foraðið} \dicLangCat{přen.} \dicDirectTranslationCS{vlákat (koho) na tenký led}
\dicEntry[forátta] \dicTerm{for··átt|a} \dicIPA{{f}{\textopeno}{\textlengthmark}{r}{au}{h}{\textsubring{d}}{a}} \dicPos{f}[1] \dicFlx{(‑u)}[5] \dicDirectTranslationCS{chyba, vada, nedostatek};  \dicPhraseIS{finna e‑u allt til foráttu} \dicDirectTranslationCS{nacházet na (čem) jen samé chyby}
\dicEntry[forboð] \dicTerm{for··boð} \dicIPA{{f}{\textopeno}{r}{\textsubring{b}}{\textopeno}{\texttheta}} \dicPos{n}[2] \dicFlx{(‑s, ‑)}[5] \dicSynonym{bann} \dicDirectTranslationCS{zákaz, zapovězení}
\dicEntry[fordild] \dicTerm{for··dild} \dicIPA{{f}{\textopeno}{r}{\textsubring{d}}{\textsci}{l}{\textsubring{d}}} \dicPos{f}[4] \dicFlx{(‑ar)}[3] \dicSynonym{tilgerð} \dicDirectTranslationCS{marnivost, strojenost, přetvářka}
\dicEntry[fordómur] \dicTerm{for··dóm|ur} \dicIPA{{f}{\textopeno}{r}{\textsubring{d}}{ou}{m}{\textscy}{\textsubring{r}}} \dicPos{m}[6] \dicFlx{(‑s, ‑ar)}[10] \dicDirectTranslationCS{předsudek} \dicExampleIS{hafa fordóma gagnvart útlendingum} \dicExampleCS{mít předsudky vůči cizincům}
\dicEntry[fordæma] \dicTerm{for··dæm|a} \dicIPA{{f}{\textopeno}{r}{\textsubring{d}}{a}{i}{m}{a}} \dicPos{v}[2] \dicFlx{(‑di, ‑t)}[141] \dicFlx{acc} \dicDirectTranslationCS{odsoudit, odsuzovat, zavrhnout, zavrhovat} \dicExampleIS{fordæma árásir hryðjuverkamanna} \dicExampleCS{odsoudit teroristické útoky}
\dicEntry[fordæmi] \dicTerm{for··dæmi} \dicIPA{{f}{\textopeno}{r}{\textsubring{d}}{a}{i}{m}{\textsci}} \dicPos{n}[2] \dicFlx{(‑s, ‑)}[14] \dicDirectTranslationCS{příklad, precedens};  \dicPhraseIS{fylgja fordæmi e‑rs} \dicDirectTranslationCS{následovat příkladu (koho)};  \dicPhraseIS{ganga á undan með góðu fordæmi} \dicDirectTranslationCS{jít dobrým příkladem}
\dicEntry[forða] \dicTerm{forð|a} \dicsymFrequent\  \dicIPA{{f}{\textopeno}{r}{ð}{a}} \dicPos{v}[1] \dicFlx{(‑aði)}[1] \dicFlx{dat} \dicSynonym*{koma undan} \dicDirectTranslationCS{zachránit, zachraňovat} \dicExampleIS{forða verðmætum frá eyðileggingu} \dicExampleCS{zachránit hodnoty před zničením};  \dicPhraseIS{forða sér} \dicDirectTranslationCS{zachránit se};  \dicPhraseIS{guð forði mér frá því} \dicDirectTranslationCS{Bože uchraň mě od toho};  \dicIdiom{forðast}{ \dicPhraseIS{forðast e‑ð}} \dicFlx{refl} \dicDirectTranslationCS{vyhýbat se (čemu), vyvarovat se (čeho)}; { \dicPhraseIS{forðast e‑n}} \dicFlx{refl} \dicDirectTranslationCS{vyhýbat se (komu), snažit se nepotkat (koho)}
\dicEntry[forði] \dicTerm{forð|i} \dicIPA{{f}{\textopeno}{r}{ð}{\textsci}} \dicPos{m}[1] \dicFlx{(‑a, ‑ar)}[1] \dicSynonym{birgðir} \dicDirectTranslationCS{zásoba} \dicExampleIS{safna forða} \dicExampleCS{sbírat\,/\addthin hromadit zásoby}
\dicEntry[forðum] \dicTerm{forðum} \dicsymFrequent\  \dicIPA{{f}{\textopeno}{r}{ð}{\textscy}{\textsubring{m}}} \dicPos{adv} \dicSynonym*{fyrir löngu} \dicDirectTranslationCS{dávno, dříve};  \dicPhraseIS{forðum daga} \dicFlx{adv} \dicDirectTranslationCS{v~dřívějších dobách}
\dicEntry[foreldrafélag] \dicTerm{for·eldra··fé·|lag} \dicIPA{{f}{\textopeno}{\textlengthmark}{r}{\textepsilon}{l}{\textsubring{d}}{r}{a}{f}{j}{\textepsilon}{l}{a}{x}} \dicPos{n}[2] \dicFlx{(‑lags, ‑lög)}[8] \dicFieldCat{škol.} \dicDirectTranslationCS{sdružení rodičů (při škole ap.)}
\dicEntry[foreldrar] \dicTerm{for··eldrar} \dicsymFrequent\  \dicIPA{{f}{\textopeno}{\textlengthmark}{r}{\textepsilon}{l}{\textsubring{d}}{r}{a}{\textsubring{r}}} \dicPos{m}[14] \dicFlx{pl}[1] \dicDirectTranslationCS{rodiče} \dicExampleIS{nýbakaðir foreldrar} \dicExampleCS{novopečení rodiče}
\dicEntry[foreldri] \dicTerm{for··eldri} \dicIPA{{f}{\textopeno}{\textlengthmark}{r}{\textepsilon}{l}{\textsubring{d}}{r}{\textsci}} \dicPos{n}[2] \dicFlx{(‑s, ‑)}[14] \dicDirectTranslationCS{rodič}
\dicEntry[forfaðir] \dicTerm{for··|faðir} \dicIPA{{f}{\textopeno}{\textsubring{r}}{f}{a}{ð}{\textsci}{\textsubring{r}}} \dicPos{m}[13] \dicFlx{(‑föður, ‑feður)}[3] \dicDirectTranslationCS{předek, praotec}
\dicEntry[forfallaður] \dicTerm{for··|fall·aður} \dicIPA{{f}{\textopeno}{\textsubring{r}}{f}{a}{\textsubring{d}}{l}{a}{ð}{\textscy}{\textsubring{r}}} \dicPos{adj}[3] \dicFlx{(f ‑fölluð)}[1] \dicDirectTranslationCS{indisponovaný, neschopný účasti}
\dicEntry[forfallakennari] \dicTerm{for·falla··kenn·ar|i} \dicIPA{{f}{\textopeno}{\textsubring{r}}{f}{a}{\textsubring{d}}{l}{a}{c\smash{\textsuperscript{h}}}{\textepsilon}{n}{a}{r}{\textsci}} \dicPos{m}[1] \dicFlx{(‑a, ‑ar)}[13] \dicFieldCat{škol.} \dicDirectTranslationCS{sup\-lu\-jí\-cí učitel(ka)}
\dicEntry[forfallast] \dicTerm{for··fall|ast} \dicIPA{{f}{\textopeno}{\textsubring{r}}{f}{a}{\textsubring{d}}{l}{a}{s}{\textsubring{d}}} \dicPos{v}[1] \dicFlx{(‑aðist)}[96] \dicFlx{refl} \dicDirectTranslationCS{nedostavit se, nezúčastnit se (z~určitých důvodů ap.)} \dicExampleIS{forfallast frá vinnu} \dicExampleCS{nedostavit se do práce}
\dicEntry[forfallinn] \dicTerm{for··fallinn} \dicIPA{{f}{\textopeno}{\textsubring{r}}{f}{a}{\textsubring{d}}{l}{\textsci}{\textsubring{n}}} \dicPos{adj}[6]\dicFlx{}[-3] \dicSynonym{óbetranlegur} \dicDirectTranslationCS{nenapravitelný, propadlý (alkoholu ap.), skalní (fanoušek ap.)} \dicExampleIS{forfallinn drykkjumaður} \dicExampleCS{nenapravitelný alkoholik}
\dicEntry[forfeður] \dicTerm{for··feður} \dicIPA{{f}{\textopeno}{\textsubring{r}}{f}{\textepsilon}{ð}{\textscy}{\textsubring{r}}} \dicPos{m} \dicFlx{pl nom} \dicLink{forfaðir}
\dicEntry[forföður] \dicTerm{for··föður} \dicIPA{{f}{\textopeno}{\textsubring{r}}{f}{\oe}{ð}{\textscy}{\textsubring{r}}} \dicPos{m} \dicFlx{sg gen} \dicLink{forfaðir}
\dicEntry[forföll] \dicTerm{for··föll} \dicIPA{{f}{\textopeno}{\textsubring{r}}{f}{\oe}{\textsubring{d}}{\textsubring{l}}} \dicPos{n}[2] \dicFlx{pl}[9] \dicDirectTranslationCS{absence, nepřítomnost} \dicExampleIS{tilkynna forföll} \dicExampleCS{oznámit absenci}
\dicEntry[forganga] \dicTerm{for··|ganga} \dicIPA{{f}{\textopeno}{r}{\r{g}}{au}{\ng}{\r{g}}{a}} \dicPos{f}[1] \dicFlx{(‑göngu)}[2] \dicDirectTranslationCS{iniciativa};  \dicPhraseIS{hafa forgöngu um e‑ð} \dicDirectTranslationCS{vyvinout iniciativu v~(čem)};  \dicPhraseIS{fyrir forgöngu e‑rs} \dicDirectTranslationCS{z~(čí) iniciativy}
\dicEntry[forgangshlutabréf] \dicTerm{for·gangs··hluta·bréf} \dicIPA{{f}\-{\textopeno}\-{r}\-{\r{g}}\-{au}\-{\ng}\-{s}\-{\textsubring{l}}\-{\textscy}\-{\textsubring{d}}\-{a}\-{\textsubring{b}}\-{r}\-{j}\-{\textepsilon}\-{f}\-} \dicPos{n}[2] \dicFlx{(‑s, ‑)}[5] \dicFieldCat{ekon.} \dicDirectTranslationCS{prioritní akcie}
\dicEntry[forgangsréttur] \dicTerm{for·gangs··rétt|ur} \dicIPA{{f}{\textopeno}{r}{\r{g}}{au}{\ng}{s}{r}{j}{\textepsilon}{h}{\textsubring{d}}{\textscy}{\textsubring{r}}} \dicPos{m}[10] \dicFlx{(‑ar)}[7] \dicDirectTranslationCS{přednost\-ní právo (na nákup ap.)}
\dicEntry[forgangsröð] \dicTerm{for·gangs··|röð} \dicIPA{{f}{\textopeno}{r}{\r{g}}{au}{\ng}{s}{r}{\oe}{\texttheta}} \dicPos{f}[7] \dicFlx{(‑raðar)}[19] \dicDirectTranslationCS{pořadí priorit}
\dicEntry[forgangur] \dicTerm{for··gang|ur} \dicIPA{{f}{\textopeno}{r}{\r{g}}{au}{\ng}{\r{g}}{\textscy}{\textsubring{r}}} \dicPos{m}[6] \dicFlx{(‑s, ‑ar)}[5] \dicDirectTranslationCS{přednost, priorita};  \dicPhraseIS{e‑að hefur forgang} \dicDirectTranslationCS{(co) má přednost}
\dicEntry[forgjöf] \dicTerm{for··|gjöf} \dicIPA{{f}{\textopeno}{r}{\r{\textObardotlessj}}{\oe}{f}} \dicPos{f}[7] \dicFlx{(‑gjafar, ‑gjafir)}[16] \textbf{1.} \dicFieldCat{sport.} \dicDirectTranslationCS{výhoda (dvou branek ap.), handicap} \dicIndirectTranslationCS{(náskok poskytovaný v~golfu ap. pro vyrovnání šancí hráčů)}  \textbf{2.} \dicSynonym{forréttindi} \dicDirectTranslationCS{privilegium, výsada}
\dicEntry[forgrunnur] \dicTerm{for··grunn|ur} \dicIPA{{f}{\textopeno}{r}{\r{g}}{r}{\textscy}{n}{\textscy}{\textsubring{r}}} \dicPos{m}[6] \dicFlx{(‑s)}[7] \dicDirectTranslationCS{popředí} \dicExampleIS{Í forgrunni myndarinnar eru þrjú tré.} \dicExampleCS{V~popředí obrazu jsou tři stromy.} \dicAntonym{bakgrunnur}
\dicEntry[forgöngumaður] \dicTerm{for·göngu··|maður} \dicIPA{{f}{\textopeno}{r}{\r{g}}{\oe i}{\ng}{\r{g}}{\textscy}{m}{a}{ð}{\textscy}{\textsubring{r}}} \dicPos{m}[13] \dicFlx{(‑manns, ‑menn)}[2] \dicDirectTranslationCS{iniciátor, hnací motor\,/\addthin síla};  \dicPhraseIS{vera forgöngumaður um e‑ð} \dicDirectTranslationCS{být hnací silou (čeho)}
\dicEntry[forheimskun] \dicTerm{for··heimsk|un} \dicIPA{{f}{\textopeno}{\textlengthmark}{\textsubring{r}}{h}{ei}{m}{s}{\r{g}}{\textscy}{\textsubring{n}}} \dicPos{f}[7] \dicFlx{(‑unar)}[9] \dicDirectTranslationCS{ohlupování, oblbování}
\dicEntry[forhertur] \dicTerm{for··hertur} \dicIPA{{f}{\textopeno}{\textlengthmark}{\textsubring{r}}{h}{\textepsilon}{\textsubring{r}}{\textsubring{d}}{\textscy}{\textsubring{r}}} \dicPos{adj}[1]\dicFlx{}[-10] \dicSynonym*{iðrunarlaus} \dicDirectTranslationCS{zarytý, zatvrzelý} \dicExampleIS{forhertur vantrúarmaður} \dicExampleCS{zarytý ateista}
\dicEntry[forhúð] \dicTerm{for··húð} \dicIPA{{f}{\textopeno}{\textlengthmark}{\textsubring{r}}{h}{u}{\texttheta}} \dicPos{f}[7] \dicFlx{(‑ar, ‑ir)}[1] \dicFieldCat{anat.} \dicDirectTranslationCS{předkožka}
\dicEntry[foringi] \dicTerm{for··ing|i} \dicIPA{{f}{\textopeno}{\textlengthmark}{r}{i}{\textltailn}{\r{\textObardotlessj}}{\textsci}} \dicPos{m}[1] \dicFlx{(‑ja, ‑jar)}[14] \dicDirectTranslationCS{předák, předačka, velitel(ka), vůdce, vůdkyně} \dicExampleIS{foringi í flokki} \dicExampleCS{stranický předák}
\dicEntry[forkastanlegur] \dicTerm{for·kastan··legur} \dicIPA{{f}{\textopeno}{\textsubring{r}}{k\smash{\textsuperscript{h}}}{a}{s}{\textsubring{d}}{a}{n}{l}{\textepsilon}{\textbabygamma}{\textscy}{\textsubring{r}}} \dicPos{adj}[1]\dicFlx{}[-8] \dicDirectTranslationCS{nepřijatelný, trestuhodný} \dicExampleIS{forkastanleg vinnubrögð} \dicExampleCS{nepřijatelný přístup k~práci}
\begin{xtolerant}{}{1pt}
\dicEntry[forkaupsréttur] \dicTerm{for·kaups··rétt|ur}\addthinS\dicIPA{{f}{\textopeno}{\textsubring{r}}{k\smash{\textsuperscript{h}}}{\oe i}{\textsubring{b}}{s}{r}{j}{\textepsilon}{h}{\textsubring{d}}{\textscy}{\textsubring{r}}}\addthinS\dicPos{m}[10]\addthinS\dicFlx{(‑ar)}[7] \dicFieldCat{ekon.} \dicDirectTranslationCS{předkupní právo}
\end{xtolerant}
\dicEntry[forkeppni] \dicTerm{for··keppn|i} \dicIPA{{f}{\textopeno}{\textsubring{r}}{c\smash{\textsuperscript{h}}}{\textepsilon}{h}{\textsubring{b}}{n}{\textsci}} \dicPos{f}[3] \dicFlx{(‑i, ‑ir)}[1] \dicFieldCat{sport.} \dicDirectTranslationCS{kvalifikace, předkolo}
\dicEntry[forkosningar] \dicTerm{for··kos·ningar} \dicIPA{{f}{\textopeno}{\textsubring{r}}{k\smash{\textsuperscript{h}}}{\textopeno}{s}{\textsubring{d}}{n}{i}{\ng}{\r{g}}{a}{\textsubring{r}}} \dicPos{f}[4] \dicFlx{pl}[6] \dicFieldCat{pol.} \dicDirectTranslationCS{primární volby, primárky}
\dicEntry[forkólfur] \dicTerm{for··kólf|ur} \dicIPA{{f}{\textopeno}{\textsubring{r}}{k\smash{\textsuperscript{h}}}{ou}{l}{v}{\textscy}{\textsubring{r}}} \dicPos{m}[6] \dicFlx{(‑s, ‑ar)}[4] \dicDirectTranslationCS{průkopník, hnací síla}
\dicEntry[forlag] \dicTerm{for··|lag} \dicIPA{{f}{\textopeno}{r}{l}{a}{x}} \dicPos{n}[2] \dicFlx{(‑lags, ‑lög)}[8] \dicDirectTranslationCS{nakladatelství, vydavatelství}
\dicEntry[forlagatrú] \dicTerm{for·laga··trú} \dicIPA{{f}{\textopeno}{r}{l}{a}{\textbabygamma}{a}{t\smash{\textsuperscript{h}}}{r}{u}} \dicPos{f}[4] \dicFlx{(‑ar)}[27] \dicDirectTranslationCS{fatalismus}
\dicEntry[forlát] \dicTerm{for··lát} \dicIPA{{f}{\textopeno}{r}{l}{au}{\textsubring{d}}} \dicPos{n}[2] \dicFlx{(‑s)}[2] \dicLangCat{zast.} \dicSynonym{afsökun} \dicDirectTranslationCS{omluva};  \dicPhraseIS{biðja  e‑n forláts vegna e‑s} \dicDirectTranslationCS{omlouvat se (komu) kvůli (čemu)}
\dicEntry[forleggjari] \dicTerm{for·leggj··ar|i} \dicIPA{{f}{\textopeno}{r}{l}{\textepsilon}{\r{\textObardotlessj}}{a}{r}{\textsci}} \dicPos{m}[1] \dicFlx{(‑a, ‑ar)}[13] \dicDirectTranslationCS{nakladatel(ka), vydavatel(ka)}
\dicEntry[forleikur] \dicTerm{for··leik|ur} \dicIPA{{f}{\textopeno}{r}{l}{ei}{\r{g}}{\textscy}{\textsubring{r}}} \dicPos{m}[9] \dicFlx{(‑s, ‑ir)}[15] \textbf{1.} \dicDirectTranslationCS{předehra, preludium, prolog};  \dicPhraseIS{forleikur að e‑u} \dicDirectTranslationCS{předehra k~(čemu)}  \textbf{2.} \dicSynonym{undanfari} \dicDirectTranslationCS{předehra} \dicIndirectTranslationCS{(událost, která přímo před něčím předchází)}  \textbf{3.} \dicSynonym{atlot} \dicDirectTranslationCS{(milostná) předehra}
\dicEntry[forliður] \dicTerm{for··lið|ur} \dicIPA{{f}{\textopeno}{r}{l}{\textsci}{ð}{\textscy}{\textsubring{r}}} \dicPos{m}[10] \dicFlx{(‑ar\,/\addthin ‑s, ‑ir)}[33] \textbf{1.} \dicFieldCat{jaz.} \dicDirectTranslationCS{první\,/\addthin předcházející část složeniny}  \textbf{2.} \dicFieldCat{jaz.} \dicDirectTranslationCS{nepřízvučná slabika stojící na počátku rytmického celku}  \textbf{3.} \dicFieldCat{filos.} \dicDirectTranslationCS{antecedent}
\dicEntry[forljótur] \dicTerm{for··ljótur} \dicIPA{{f}{\textopeno}{r}{l}{j}{ou}{\textsubring{d}}{\textscy}{\textsubring{r}}} \dicPos{adj}[1]\dicFlx{}[-6] \dicDirectTranslationCS{ohavný, ohyzdný}
\dicEntry[forlög] \dicTerm{for··lög} \dicIPA{{f}{\textopeno}{r}{l}{\oe}{x}} \dicPos{n}[2] \dicFlx{pl}[9] \dicSynonym{örlög} \dicDirectTranslationCS{osud}
\dicEntry[form] \dicTerm{form} \dicsymFrequent\  \dicIPA{{f}{\textopeno}{r}{\textsubring{m}}} \dicPos{n}[2] \dicFlx{(‑s, ‑)}[5] \textbf{1.} \dicSynonym{mót\smash{\textsuperscript{2}}} \dicDirectTranslationCS{forma (na pečení ap.)}  \textbf{2.} \dicSynonym{háttur\smash{\textsuperscript{1}}} \dicDirectTranslationCS{forma} \dicIndirectTranslationCS{(způsob podání)}  \textbf{3.} \dicSynonym{lögun} \dicDirectTranslationCS{forma, tvar}  \textbf{4.} \dicSynonym*{bygging ritverks} \dicDirectTranslationCS{forma, struktura (literárního díla ap.)}  \textbf{5.} \dicSynonym*{líkamlegt ástand} \dicDirectTranslationCS{forma} \dicIndirectTranslationCS{(dobrý stav)};  \dicPhraseIS{halda sér í formi} \dicDirectTranslationCS{udržovat se ve formě, udržovat si formu};  \dicPhraseIS{vera í góðu formi} \dicDirectTranslationCS{být v~dobré formě}
\dicEntry[forma] \dicTerm{form|a} \dicIPA{{f}{\textopeno}{r}{m}{a}} \dicPos{v}[1] \dicFlx{(‑aði)}[1] \dicFlx{acc} \dicSynonym{móta} \dicDirectTranslationCS{utvářet, vytvářet, formovat} \dicExampleIS{forma brjóstmynd} \dicExampleCS{formovat bustu}
\dicEntry[formaður] \dicTerm{for··|maður} \dicsymFrequent\  \dicIPA{{f}{\textopeno}{r}{m}{a}{ð}{\textscy}{\textsubring{r}}} \dicPos{m}[13] \dicFlx{(‑manns, ‑menn)}[2] \textbf{1.} \dicSynonym{leiðtogi} \dicDirectTranslationCS{předseda, předsedkyně, leader, lídr} \dicExampleIS{formenn stjórnarflokkanna} \dicExampleCS{předsedové vládních stran}  \textbf{2.} \dicSynonym{skipstjóri} \dicDirectTranslationCS{kapitán(ka)} \dicIndirectTranslationCS{(velitel lodi)}
\dicEntry[formáli] \dicTerm{for··mál|i} \dicIPA{{f}{\textopeno}{r}{m}{au}{l}{\textsci}} \dicPos{m}[1] \dicFlx{(‑a, ‑ar)}[1] \dicSynonym{inngangsorð} \dicDirectTranslationCS{předmluva, prolog (knihy ap.)} \dicExampleIS{formáli að e‑u} \dicExampleCS{předmluva k~(čemu)}
\dicEntry[formbylting] \dicTerm{form··bylt·ing} \dicIPA{{f}{\textopeno}{r}{m}{\textsubring{b}}{\textsci}{\textsubring{l}}{\textsubring{d}}{i}{\ng}{\r{g}}} \dicPos{f}[4] \dicFlx{(‑ar, ‑ar)}[5] \dicFieldCat{lit.} \dicIndirectTranslationCS{změna formy v~islandském básnictví v~pol. 20. století}
\dicEntry[formennska] \dicTerm{for··mennsk|a} \dicIPA{{f}{\textopeno}{r}{m}{\textepsilon}{n}{s}{\r{g}}{a}} \dicPos{f}[1] \dicFlx{(‑u)}[5] \textbf{1.} \dicDirectTranslationCS{předsednictví} \dicExampleIS{formennska í félagi} \dicExampleCS{předsednictví ve společnosti}  \textbf{2.} \dicSynonym*{bátstjórn} \dicDirectTranslationCS{velení (na lodi ap.)}
\dicEntry[Formentera] \dicTerm{Formenter|a} \dicIPA{{f}{\textopeno}{r}{m}{\textepsilon}{\textsubring{n}}{\textsubring{d}}{\textepsilon}{r}{a}} \dicPos{f}[1] \dicFlx{(‑u)}[6] \dicFieldCat{geog.} \dicDirectTranslationCS{Formentera} \dicIndirectTranslationCS{(ostrov v~souostroví Baleáry)}
\dicEntry[formerki] \dicTerm{for··merki} \dicIPA{{f}{\textopeno}{r}{m}{\textepsilon}{\textsubring{r}}{\r{\textObardotlessj}}{\textsci}} \dicPos{n}[2] \dicFlx{(‑s, ‑)}[16] \textbf{1.} \dicSynonym{auðkenni} \dicDirectTranslationCS{(charakteristický) znak, znamení}  \textbf{2.} \dicFieldCat{hud.} \dicDirectTranslationCS{posuvka}
\dicEntry[formfastur] \dicTerm{form··|fastur} \dicIPA{{f}{\textopeno}{r}{m}{f}{a}{s}{\textsubring{d}}{\textscy}{\textsubring{r}}} \dicPos{adj}[1] \dicFlx{(f ‑föst)}[11] \dicDirectTranslationCS{klasický, formální} \dicIndirectTranslationCS{(lpící na formě)}
\dicEntry[formfræði] \dicTerm{form··fræð|i} \dicIPA{{f}{\textopeno}{r}{m}{f}{r}{a}{i}{ð}{\textsci}} \dicPos{f}[3] \dicFlx{(‑i)}[3] \dicDirectTranslationCS{morfologie} \dicIndirectTranslationCS{(nauka o~tvarech)}
\dicEntry[formgerð] \dicTerm{form··gerð} \dicIPA{{f}{\textopeno}{r}{m}{\r{\textObardotlessj}}{\textepsilon}{r}{\texttheta}} \dicPos{f}[7] \dicFlx{(‑ar, ‑ir)}[1] \dicDirectTranslationCS{struktura, stavba (literárního díla ap.)}
\dicEntry[formkaka] \dicTerm{form··|kaka} \dicIPA{{f}{\textopeno}{r}{m}{k\smash{\textsuperscript{h}}}{a}{\r{g}}{a}} \dicPos{f}[1] \dicFlx{(‑köku, ‑kökur)}[20] \dicFieldCat{kulin.} \dicIndirectTranslationCS{dort z~formy}
\dicEntry[formlega] \dicTerm{form··lega} \dicIPA{{f}{\textopeno}{r}{m}{l}{\textepsilon}{\textbabygamma}{a}} \dicPos{adv} \dicDirectTranslationCS{formálně}
\dicEntry[formlegur] \dicTerm{form··legur} \dicIPA{{f}{\textopeno}{r}{m}{l}{\textepsilon}{\textbabygamma}{\textscy}{\textsubring{r}}} \dicPos{adj}[1]\dicFlx{}[-8] \dicDirectTranslationCS{formální}
\dicEntry[formsatriði] \dicTerm{forms··at·riði} \dicIPA{{f}{\textopeno}{r}{m}{s}{a}{\textsubring{d}}{r}{\textsci}{ð}{\textsci}} \dicPos{n}[2] \dicFlx{(‑s, ‑)}[14] \dicDirectTranslationCS{formalita, formálnost}
\dicEntry[formúla] \dicTerm{formúl|a} \dicIPA{{f}{\textopeno}{r}{m}{u}{l}{a}} \dicPos{f}[1] \dicFlx{(‑u, ‑ur)}[19] \textbf{1.} \dicSynonym{uppskrift} \dicDirectTranslationCS{recept, předpis, formule}  \textbf{2.} \dicFieldCat{mat.} \dicDirectTranslationCS{formule, vzorec} \dicExampleIS{formúla fyrir flatarmáli þríhyrnings} \dicExampleCS{vzorec pro výpočet obsahu trojúhelníku}
\dicEntry[forn] \dicTerm{forn} \dicsymFrequent\  \dicIPA{{f}{\textopeno}{r}{\textsubring{d}}{\textsubring{n}}} \dicPos{adj}[5]\dicFlx{}[-5] \dicDirectTranslationCS{starověký, starobylý, dávný, archaický} \dicExampleIS{fornar siðvenjur} \dicExampleCS{starobylé tradice};  \dicPhraseIS{frá fornu fari} \dicFlx{adv} \dicDirectTranslationCS{odjakživa, od dávných dob};  \dicPhraseIS{til forna} \dicFlx{adv} \dicDirectTranslationCS{za dávných časů, v~dávných dobách}
\dicEntry[fornafn] \dicTerm{for··|nafn} \dicIPA{{f}{\textopeno}{r}{n}{a}{\textsubring{b}}{\textsubring{n}}} \dicPos{n}[2] \dicFlx{(‑nafns, ‑nöfn)}[8] \textbf{1.} \dicSynonym{skírnarnafn} \dicDirectTranslationCS{křestní jméno}  \textbf{2.} \dicFieldCat{jaz.} \dicDirectTranslationCS{zájmeno} \textit{(l.~{\textLA{pronomen}})}
\dicEntry[fornbókabúð] \dicTerm{forn·bóka··búð} \dicIPA{{f}{\textopeno}{r}{\textsubring{d}}{\textsubring{n}}{\textsubring{b}}{ou}{\r{g}}{a}{\textsubring{b}}{u}{\texttheta}} \dicPos{f}[7] \dicFlx{(‑ar, ‑ir)}[1] \dicDirectTranslationCS{antikvariát}
\dicEntry[fornbókasala] \dicTerm{forn··bóka·|sala} \dicIPA{{f}{\textopeno}{r}{\textsubring{d}}{\textsubring{n}}{\textsubring{b}}{ou}{\r{g}}{a}{s}{a}{l}{a}} \dicPos{f}[1] \dicFlx{(‑sölu, ‑sölur)}[20] \dicDirectTranslationCS{antikvariát}
\dicEntry[fornbókasali] \dicTerm{forn·bóka··sal|i} \dicIPA{{f}{\textopeno}{r}{\textsubring{d}}{\textsubring{n}}{\textsubring{b}}{ou}{\r{g}}{a}{s}{a}{l}{\textsci}} \dicPos{m}[1] \dicFlx{(‑a, ‑ar)}[8] \dicDirectTranslationCS{antikvář(ka)} \dicIndirectTranslationCS{(obchodník se staršími knihami)}
\dicEntry[fornbréf] \dicTerm{forn··bréf} \dicIPA{{f}{\textopeno}{r}{\textsubring{d}}{\textsubring{n}}{\textsubring{b}}{r}{j}{\textepsilon}{f}} \dicPos{n}[2] \dicFlx{(‑s, ‑)}[5] \dicFieldCat{hist.} \dicDirectTranslationCS{smlouva, listina}
\dicEntry[forneskja] \dicTerm{forn··eskj|a} \dicIPA{{f}{\textopeno}{r}{\textsubring{d}}{n}{\textepsilon}{s}{\r{\textObardotlessj}}{a}} \dicPos{f}[1] \dicFlx{(‑u)}[5] \textbf{1.} \dicSynonym*{ævaforn tími} \dicDirectTranslationCS{dávnověk, dávné časy}  \textbf{2.} \dicSynonym{galdur} \dicDirectTranslationCS{čarodějnictví, magie}
\dicEntry[fornfálegur] \dicTerm{forn··fá·legur} \dicIPA{{f}{\textopeno}{r}{\textsubring{d}}{\textsubring{n}}{f}{au}{l}{\textepsilon}{\textbabygamma}{\textscy}{\textsubring{r}}} \dicPos{adj}[1]\dicFlx{}[-8] \dicDirectTranslationCS{staromódní, zastaralý}
\dicEntry[fornfræði] \dicTerm{forn··fræð|i} \dicIPA{{f}{\textopeno}{r}{\textsubring{d}}{\textsubring{n}}{f}{r}{a}{i}{ð}{\textsci}} \dicPos{f}[3] \dicFlx{(‑i)}[3] \textbf{1.} \dicDirectTranslationCS{klasická filologie}  \textbf{2.} \dicDirectTranslationCS{věda (filologie\,/\addthin archeologie\,/\addthin etnologie) zaměřená na období doby vikinské}
\dicEntry[fornfræðingur] \dicTerm{forn·fræð··ing|ur} \dicIPA{{f}{\textopeno}{r}{\textsubring{d}}{\textsubring{n}}{f}{r}{a}{i}{ð}{i}{\ng}{\r{g}}{\textscy}{\textsubring{r}}} \dicPos{m}[6] \dicFlx{(‑s, ‑ar)}[8] \textbf{1.} \dicLangCat{zast.} \dicSynonym*{stundar klassísk fræði} \dicDirectTranslationCS{klasický filolog}  \textbf{2.} \dicLangCat{zast.} \dicSynonym*{stundar forn fræði} \dicIndirectTranslationCS{filolog specializující se na staroseverštinu, archeolog\,/\addthin etnolog specializující se na období doby vikinské}
\dicEntry[forngripur] \dicTerm{forn··grip|ur} \dicIPA{{f}{\textopeno}{r}{\textsubring{d}}{\textsubring{n}}{\r{g}}{r}{\textsci}{\textsubring{b}}{\textscy}{\textsubring{r}}} \dicPos{m}[9] \dicFlx{(‑s, ‑ir)}[8] \dicDirectTranslationCS{starožitnost, starožitný předmět, antikvita}
\dicEntry[forngríska] \dicTerm{forn··grísk|a} \dicIPA{{f}{\textopeno}{r}{\textsubring{d}}{\textsubring{n}}{\r{g}}{r}{i}{s}{\r{g}}{a}} \dicPos{f}[1] \dicFlx{(‑u)}[5] \dicDirectTranslationCS{starořečtina}
\dicEntry[forníslenska] \dicTerm{forn··ís·lensk|a} \dicIPA{{f}{\textopeno}{r}{\textsubring{d}}{n}{i}{s}{\textsubring{d}}{l}{\textepsilon}{n}{s}{\r{g}}{a}} \dicPos{f}[1] \dicFlx{(‑u)}[5] \dicDirectTranslationCS{staroislandština} \dicIndirectTranslationCS{(islandský jazyk do roku 1350)}
\dicEntry[fornkvæði] \dicTerm{forn··kvæði} \dicIPA{{f}{\textopeno}{r}{\textsubring{d}}{\textsubring{n}}{k\smash{\textsuperscript{h}}}{v}{a}{i}{ð}{\textsci}} \dicPos{n}[2] \dicFlx{(‑s, ‑)}[14] \textbf{1.} \dicFieldCat{lit.} \dicSynonym*{forn kveðskapur} \dicDirectTranslationCS{staroseverská báseň}  \textbf{2.} \dicFieldCat{lit.} \dicSynonym*{sagnadansar} \dicDirectTranslationCS{balada}
\dicEntry[fornleifafræði] \dicTerm{forn·leifa··fræð|i} \dicIPA{{f}{\textopeno}{r}{\textsubring{d}}{n}{l}{ei}{v}{a}{f}{r}{a}{i}{ð}{\textsci}} \dicPos{f}[3] \dicFlx{(‑i)}[3] \dicDirectTranslationCS{archeologie}
\dicEntry[fornleifafræðingur] \dicTerm{forn·leifa·fræð··ing|ur} \dicIPA{{f}\-{\textopeno}\-{r}\-{\textsubring{d}}\-{n}\-{l}\-{ei}\-{v}\-{a}\-{f}\-{r}\-{a}\-{i}\-{ð}\-{i}\-{\ng}\-{\r{g}}\-{\textscy}\-{\textsubring{r}}\-} \dicPos{m}[6] \dicFlx{(‑s, ‑ar)}[8] \dicDirectTranslationCS{archeolog, archeoložka}
\dicEntry[fornleifafundur] \dicTerm{forn·leifa··fund|ur} \dicIPA{{f}{\textopeno}{r}{\textsubring{d}}{n}{l}{ei}{v}{a}{f}{\textscy}{n}{\textsubring{d}}{\textscy}{\textsubring{r}}} \dicPos{m}[10] \dicFlx{(‑ar, ‑ir)}[4] \dicDirectTranslationCS{archeologický nález}
\dicEntry[fornleifar] \dicTerm{forn··leifar} \dicIPA{{f}{\textopeno}{r}{\textsubring{d}}{n}{l}{ei}{v}{a}{\textsubring{r}}} \dicPos{f}[4] \dicFlx{pl}[2] \dicDirectTranslationCS{starožitnosti, archeologické památky\,/\addthin nálezy}
\dicEntry[fornlífsöld] \dicTerm{forn·lífs··|öld} \dicIPA{{f}{\textopeno}{r}{\textsubring{d}}{n}{l}{i}{f}{s}{\oe}{l}{\textsubring{d}}} \dicPos{f}[7] \dicFlx{(‑aldar)}[19] \dicFieldCat{geol.} \dicDirectTranslationCS{paleozoikum}
\dicEntry[fornmaður] \dicTerm{forn··|maður} \dicIPA{{f}{\textopeno}{r}{\textsubring{d}}{n}{m}{a}{ð}{\textscy}{\textsubring{r}}} \dicPos{m}[13] \dicFlx{(‑manns, ‑menn)}[2] \dicFieldCat{hist.} \dicIndirectTranslationCS{člověk z~časů osidlování Islandu}
\dicEntry[fornmál] \dicTerm{forn··mál} \dicIPA{{f}{\textopeno}{r}{\textsubring{d}}{n}{m}{au}{\textsubring{l}}} \dicPos{n}[2] \dicFlx{(‑s, ‑)}[5] \dicDirectTranslationCS{starý\,/\addthin dávný jazyk}
\dicEntry[fornmenntastefna] \dicTerm{forn·mennta··stefn|a} \dicIPA{{f}{\textopeno}{r}{\textsubring{d}}{n}{m}{\textepsilon}{\textsubring{n}}{\textsubring{d}}{a}{s}{\textsubring{d}}{\textepsilon}{\textsubring{b}}{n}{a}} \dicPos{f}[1] \dicFlx{(‑u)}[5] \dicFieldCat{filos.} \dicDirectTranslationCS{humanismus}
\dicEntry[fornminjar] \dicTerm{forn··minjar} \dicIPA{{f}{\textopeno}{r}{\textsubring{d}}{n}{m}{\textsci}{n}{j}{a}{\textsubring{r}}} \dicPos{f}[12] \dicFlx{pl}[1] \dicDirectTranslationCS{(starověké) památky, starožitnosti}
\dicEntry[fornminjasafn] \dicTerm{forn·minja··|safn} \dicIPA{{f}{\textopeno}{r}{\textsubring{d}}{n}{m}{\textsci}{n}{j}{a}{s}{a}{\textsubring{b}}{\textsubring{n}}} \dicPos{n}[2] \dicFlx{(‑safns, ‑söfn)}[8] \dicDirectTranslationCS{historické\,/\addthin archeologické muzeum}
\dicEntry[fornnorræna] \dicTerm{forn··nor·ræn|a} \dicIPA{{f}{\textopeno}{r}{\textsubring{d}}{n}{\textopeno}{r}{a}{i}{n}{a}} \dicPos{f}[1] \dicFlx{(‑u)}[5] \dicSynonym{norræna} \dicDirectTranslationCS{staroseverština}
\dicEntry[fornnorrænn] \dicTerm{forn··nor·rænn} \dicIPA{{f}{\textopeno}{r}{\textsubring{d}}{n}{\textopeno}{r}{a}{i}{\textsubring{d}}{\textsubring{n}}} \dicPos{adj}[7]\dicFlx{}[-2] \dicDirectTranslationCS{staroseverský}
\dicEntry[fornrit] \dicTerm{forn··rit} \dicIPA{{f}{\textopeno}{r}{\textsubring{d}}{n}{r}{\textsci}{\textsubring{d}}} \dicPos{n}[2] \dicFlx{(‑s, ‑)}[5] \dicFieldCat{lit.} \dicDirectTranslationCS{starověký spis\,/\addthin rukopis}
\dicEntry[fornsaga] \dicTerm{forn··|saga} \dicIPA{{f}{\textopeno}{r}{\textsubring{d}}{\textsubring{n}}{s}{a}{\textbabygamma}{a}} \dicPos{f}[1] \dicFlx{(‑sögu, ‑sögur)}[14] \textbf{1.} \dicDirectTranslationCS{starý\,/\addthin starodávný příběh}  \textbf{2.} \dicFieldCat{lit.} \dicDirectTranslationCS{sága}
\dicEntry[fornsala] \dicTerm{forn··|sala} \dicIPA{{f}{\textopeno}{r}{\textsubring{d}}{\textsubring{n}}{s}{a}{l}{a}} \dicPos{f}[1] \dicFlx{(‑sölu, ‑sölur)}[20] \dicDirectTranslationCS{starožitnictví, obchod se starožitnostmi}
\dicEntry[fornsali] \dicTerm{forn··sal|i} \dicIPA{{f}{\textopeno}{r}{\textsubring{d}}{\textsubring{n}}{s}{a}{l}{\textsci}} \dicPos{m}[1] \dicFlx{(‑a, ‑ar)}[8] \dicDirectTranslationCS{starožitník, starožitnice, antikvář(ka)}
\dicEntry[fornsteinöld] \dicTerm{forn·stein··|öld} \dicIPA{{f}{\textopeno}{r}{\textsubring{d}}{\textsubring{n}}{s}{\textsubring{d}}{ei}{n}{\oe}{l}{\textsubring{d}}} \dicPos{f}[7] \dicFlx{(‑aldar)}[19] \dicDirectTranslationCS{paleolit}
\dicEntry[fornyrði] \dicTerm{forn··yrði} \dicIPA{{f}{\textopeno}{r}{\textsubring{d}}{n}{\textsci}{r}{ð}{\textsci}} \dicPos{n}[2] \dicFlx{(‑s, ‑)}[14] \dicFieldCat{jaz.} \dicDirectTranslationCS{archaismus, zastaralý jazykový prostředek}
\dicEntry[fornyrðislag] \dicTerm{forn·yrðis··lag} \dicIPA{{f}{\textopeno}{r}{\textsubring{d}}{n}{\textsci}{r}{ð}{\textsci}{s}{l}{a}{x}} \dicPos{n}[2] \dicFlx{(‑s)}[2] \dicFieldCat{lit.} \dicIndirectTranslationCS{druh eddického metra (osmiřádkové, bez rýmu, s~aliterací)}
\dicEntry[fornöld] \dicTerm{forn··|öld} \dicIPA{{f}{\textopeno}{r}{\textsubring{d}}{n}{\oe}{l}{\textsubring{d}}} \dicPos{f}[7] \dicFlx{(‑aldar)}[19] \textbf{1.} \dicFieldCat{hist.} \dicDirectTranslationCS{pravěk}  \textbf{2.} \dicFieldCat{hist.} \dicDirectTranslationCS{středověk na Islandu}
\dicEntry[forpróf] \dicTerm{for··próf} \dicIPA{{f}{\textopeno}{\textsubring{r}}{p\smash{\textsuperscript{h}}}{r}{ou}{f}} \dicPos{n}[2] \dicFlx{(‑s, ‑)}[5] \dicDirectTranslationCS{přijímací zkouška}
\dicEntry[forráð] \dicTerm{for··ráð} \dicIPA{{f}{\textopeno}{r}{\textlengthmark}{au}{\texttheta}} \dicPos{n}[2] \dicFlx{(‑s, ‑)}[5] \dicSynonym{umráð} \dicDirectTranslationCS{správa, administrace, řízení};  \dicPhraseIS{kunna fótum sínum forráð} \dicLangCat{přen.} \dicDirectTranslationCS{umět se o~sebe postarat}
\dicEntry[forráðamaður] \dicTerm{for·ráða··|maður} \dicIPA{{f}{\textopeno}{r}{\textlengthmark}{au}{ð}{a}{m}{a}{ð}{\textscy}{\textsubring{r}}} \dicPos{m}[13] \dicFlx{(‑manns, ‑menn)}[2] \textbf{1.} \dicSynonym{yfirmaður} \dicDirectTranslationCS{správce, správkyně, vedoucí, leader, lídr}  \textbf{2.} \dicDirectTranslationCS{poručník, poručnice, kurátor(ka), opatrovník, opatrovnice} \dicExampleIS{forráðamenn barna og unglinga} \dicExampleCS{poručníci dětí a~mládeže}
\dicEntry[forréttindi] \dicTerm{for··rétt·indi} \dicIPA{{f}{\textopeno}{r}{\textlengthmark}{j}{\textepsilon}{h}{\textsubring{d}}{\textsci}{n}{\textsubring{d}}{\textsci}} \dicPos{n}[2] \dicFlx{pl}[19] \dicDirectTranslationCS{výsada, privilegium} \dicExampleIS{njóta forréttinda} \dicExampleCS{užívat privilegia}
\dicEntry[forréttur] \dicTerm{for··rétt|ur} \dicIPA{{f}{\textopeno}{r}{\textlengthmark}{j}{\textepsilon}{h}{\textsubring{d}}{\textscy}{\textsubring{r}}} \dicPos{m}[10] \dicFlx{(‑ar, ‑ir)}[4] \dicDirectTranslationCS{předkrm} \dicExampleIS{Ég ætla að fá rækjur í forrétt.} \dicExampleCS{Rád bych si dal krevety jako předkrm.}
\dicEntry[forrit] \dicTerm{for··rit} \dicsymFrequent\  \dicIPA{{f}{\textopeno}{r}{\textlengthmark}{\textsci}{\textsubring{d}}} \dicPos{n}[2] \dicFlx{(‑s, ‑)}[5] \textbf{1.} \dicFieldCat{poč.} \dicSynonym{tölvuforrit} \dicDirectTranslationCS{program, aplikace}  \textbf{2.} \dicFieldCat{lit.} \dicDirectTranslationCS{originál, původní rukopis} \dicExampleIS{forrit handrits} \dicExampleCS{originál rukopisu}
\dicEntry[forrita] \dicTerm{for··rit|a} \dicIPA{{f}{\textopeno}{r}{\textlengthmark}{\textsci}{\textsubring{d}}{a}} \dicPos{v}[1] \dicFlx{(‑aði)}[1] \dicFieldCat{poč.} \dicDirectTranslationCS{programovat}
\dicEntry[forritanlegur] \dicTerm{for·ritan··legur} \dicIPA{{f}{\textopeno}{r}{\textlengthmark}{\textsci}{\textsubring{d}}{a}{n}{l}{\textepsilon}{\textbabygamma}{\textscy}{\textsubring{r}}} \dicPos{adj}[1]\dicFlx{}[-8] \dicFieldCat{poč.} \dicDirectTranslationCS{programovatelný}
\dicEntry[forritari] \dicTerm{for··rit·ar|i} \dicIPA{{f}{\textopeno}{r}{\textlengthmark}{\textsci}{\textsubring{d}}{a}{r}{\textsci}} \dicPos{m}[1] \dicFlx{(‑a, ‑ar)}[13] \dicDirectTranslationCS{programátor(ka)}
\dicEntry[forritun] \dicTerm{for··rit|un} \dicIPA{{f}{\textopeno}{r}{\textlengthmark}{\textsci}{\textsubring{d}}{\textscy}{\textsubring{n}}} \dicPos{f}[7] \dicFlx{(‑unar)}[9] \dicFieldCat{poč.} \dicDirectTranslationCS{programování}
\dicEntry[forritunarmál] \dicTerm{for·ritunar··mál} \dicIPA{{f}{\textopeno}{r}{\textlengthmark}{\textsci}{\textsubring{d}}{\textscy}{n}{a}{r}{m}{au}{\textsubring{l}}} \dicPos{n}[2] \dicFlx{(‑s, ‑)}[5] \dicFieldCat{poč.} \dicDirectTranslationCS{programovací jazyk}
\dicEntry[forríkur] \dicTerm{for··ríkur} \dicIPA{{f}{\textopeno}{r}{\textlengthmark}{i}{\r{g}}{\textscy}{\textsubring{r}}} \dicPos{adj}[1]\dicFlx{}[-6] \dicDirectTranslationCS{zámožný, velmi bohatý}
\dicEntry[forræði] \dicTerm{for··ræði} \dicIPA{{f}{\textopeno}{r}{\textlengthmark}{a}{i}{ð}{\textsci}} \dicPos{n}[2] \dicFlx{(‑s)}[20] \dicFieldCat{práv.} \dicDirectTranslationCS{poručnictví} \dicExampleIS{fá forræði yfir barni} \dicExampleCS{získat poručnictví nad dítětem}
\dicEntry[forsala] \dicTerm{for··|sala} \dicIPA{{f}{\textopeno}{\textsubring{r}}{s}{a}{l}{a}} \dicPos{f}[1] \dicFlx{(‑sölu, ‑sölur)}[20] \dicDirectTranslationCS{předprodej}
\dicEntry[forsenda] \dicTerm{for··send|a} \dicsymFrequent\  \dicIPA{{f}{\textopeno}{\textsubring{r}}{s}{\textepsilon}{n}{\textsubring{d}}{a}} \dicPos{f}[1] \dicFlx{(‑u, ‑ur)}[19] \textbf{1.} \dicSynonym{skilyrði} \dicDirectTranslationCS{předpoklad} \dicExampleIS{grandskoða allar forsendur} \dicExampleCS{přezkoumat všechny předpoklady}  \textbf{2.} \dicFieldCat{filos.} \dicDirectTranslationCS{premisa, předpoklad}
\dicEntry[forsetaembætti] \dicTerm{for·seta··embætti} \dicIPA{{f}{\textopeno}{\textsubring{r}}{s}{\textepsilon}{\textsubring{d}}{a}{\textepsilon}{m}{\textsubring{b}}{a}{i}{h}{\textsubring{d}}{\textsci}} \dicPos{n}[2] \dicFlx{(‑s, ‑)}[14] \dicDirectTranslationCS{prezidentský úřad}
\dicEntry[forsetaframbjóðandi] \textls[15]{\dicTerm{for·seta··fram·bjóð·|andi} \dicIPA{{f}\-{\textopeno}\-{\textsubring{r}}\-{s}\-{\textepsilon}\-{\textsubring{d}}\-{a}\-{f}\-{r}\-{a}\-{m}\-{\textsubring{b}}\-{j}\-{ou}\-{ð}\-{a}\-{n}\-{\textsubring{d}}\-{\textsci}} \dicPos{m}[2]} \dicFlx{(‑anda, ‑endur)}[1] \dicDirectTranslationCS{prezidentský kandidát, prezidentská kandidátka}
\dicEntry[forsetaframboð] \dicTerm{for·seta··fram·boð} \dicIPA{{f}{\textopeno}{\textsubring{r}}{s}{\textepsilon}{\textsubring{d}}{a}{f}{r}{a}{m}{\textsubring{b}}{\textopeno}{\texttheta}} \dicPos{n}[2] \dicFlx{(‑s, ‑)}[5] \dicDirectTranslationCS{prezidentská kandidatura}
\dicEntry[forsetafrú] \dicTerm{for·seta··frú} \dicIPA{{f}{\textopeno}{\textsubring{r}}{s}{\textepsilon}{\textsubring{d}}{a}{f}{r}{u}} \dicPos{f}[4] \dicFlx{(‑ar, ‑r)}[26] \dicDirectTranslationCS{první dáma, žena prezidenta}
\dicEntry[forsetakosningar] \dicTerm{for·seta··kos·ningar} \dicIPA{{f}{\textopeno}{\textsubring{r}}{s}{\textepsilon}{\textsubring{d}}{a}{k\smash{\textsuperscript{h}}}{\textopeno}{s}{\textsubring{d}}{n}{i}{\ng}{\r{g}}{a}{\textsubring{r}}} \dicPos{f}[4] \dicFlx{pl}[6] \dicFieldCat{pol.} \dicDirectTranslationCS{prezidentské volby}
\dicEntry[forsetaskrifstofa] \dicTerm{for·seta··skrif·stof|a} \dicIPA{{f}\-{\textopeno}\-{\textsubring{r}}\-{s}\-{\textepsilon}\-{\textsubring{d}}\-{a}\-{s}\-{\r{g}}\-{r}\-{\textsci}\-{f}\-{s}\-{\textsubring{d}}\-{\textopeno}\-{v}\-{a}\-} \dicPos{f}[1] \dicFlx{(‑u, ‑ur)}[7] \dicDirectTranslationCS{prezidentská kancelář}
\dicEntry[forseti] \dicTerm{for··set|i} \dicsymFrequent\  \dicIPA{{f}{\textopeno}{\textsubring{r}}{s}{\textepsilon}{\textsubring{d}}{\textsci}} \dicPos{m}[1] \dicFlx{(‑a, ‑ar)}[1] \textbf{1.} \dicSynonym{formaður} \dicDirectTranslationCS{předseda, předsedkyně} \dicExampleIS{forseti alþingis} \dicExampleCS{předseda parlamentu}  \textbf{2.} \dicFieldCat{pol.} \dicDirectTranslationCS{prezident(ka)} \dicExampleIS{forseti Íslands} \dicExampleCS{prezident Islandu}
\dicEntry[forsetning] \dicTerm{for··set·ning} \dicIPA{{f}{\textopeno}{\textsubring{r}}{s}{\textepsilon}{h}{\textsubring{d}}{n}{i}{\ng}{\r{g}}} \dicPos{f}[4] \dicFlx{(‑ar, ‑ar)}[5] \dicFieldCat{jaz.} \dicDirectTranslationCS{předložka} \textit{(l.~{\textLA{praepositio}})}
\begin{xtolerant}{}{1pt}
\dicEntry[forsetningarliður] \dicTerm{for·setningar··lið|ur}\addthinS\dicIPA{{f}\-{\textopeno}\-{\textsubring{r}}\-{s}\-{\textepsilon}\-{h}\-{\textsubring{d}}\-{n}\-{i}\-{\ng}\-{\r{g}}\-{a}\-{r}\-{l}\-{\textsci}\-{ð}\-{\textscy}\-{\textsubring{r}}\-}\addthinS\dicPos{m}[10]\addthinS\dicFlx{(‑ar\,/\addthinS ‑s, ‑ir)}[33] \dicFieldCat{jaz.} \dicDirectTranslationCS{předložková fráze}
\end{xtolerant}
\dicEntry[forsíða] \dicTerm{for··síð|a} \dicIPA{{f}{\textopeno}{\textsubring{r}}{s}{i}{ð}{a}} \dicPos{f}[1] \dicFlx{(‑u, ‑ur)}[13] \dicDirectTranslationCS{přední\,/\addthin titulní strana (novin ap.)}
\dicEntry[forsjá] \dicTerm{for··sjá} \dicIPA{{f}{\textopeno}{\textsubring{r}}{s}{j}{au}} \dicPos{f}[4] \dicFlx{(‑r)}[19] \textbf{1.} \dicSynonym{umsjón} \dicDirectTranslationCS{dohled, dozor}  \textbf{2.} \dicSynonym{fyrirhyggja} \dicDirectTranslationCS{prozíravost, prozřetelnost}  \textbf{3.} \dicFieldCat{práv.} \dicDirectTranslationCS{opatrovnictví (dítěte)}
\dicEntry[forsjáll] \dicTerm{for··sjáll} \dicIPA{{f}{\textopeno}{\textsubring{r}}{s}{j}{au}{\textsubring{d}}{\textsubring{l}}} \dicPos{adj}[8]\dicFlx{}[-1] \dicSynonym*{ráðspakur} \dicDirectTranslationCS{předvídavý, prozíravý, prozřetelný} \dicAntonym{óforsjáll}
\dicEntry[forsjálni] \dicTerm{for··sjáln|i} \dicIPA{{f}{\textopeno}{\textsubring{r}}{s}{j}{au}{l}{n}{\textsci}} \dicPos{f}[3] \dicFlx{(‑i)}[3] \dicDirectTranslationCS{prozíravost, prozřetelnost} \dicAntonym{óforsjálni}
\dicEntry[forsjón] \dicTerm{for··sjón} \dicIPA{{f}{\textopeno}{\textsubring{r}}{s}{j}{ou}{\textsubring{n}}} \dicPos{f}[7] \dicFlx{(‑ar)}[3] \dicDirectTranslationCS{prozíravost, prozřetelnost, obezřetnost} \dicExampleIS{treysta forsjóninni} \dicExampleCS{důvěřovat prozíravosti}
\dicEntry[forsk.] \dicTerm{forsk.} \dicPos{zkr} \dicPhraseIS{forskeyti} \dicFieldCat{jaz.} \dicDirectTranslationCS{předpona}
\dicEntry[forskaut] \dicTerm{for··skaut} \dicIPA{{f}{\textopeno}{\textsubring{r}}{s}{\r{g}}{\oe i}{\textsubring{d}}} \dicPos{n}[2] \dicFlx{(‑s, ‑)}[5] \dicFieldCat{fyz.} \dicDirectTranslationCS{anoda} \dicAntonym{bakskaut}
\dicEntry[forskeyti] \dicTerm{for··skeyti} \dicIPA{{f}{\textopeno}{\textsubring{r}}{s}{\r{\textObardotlessj}}{ei}{\textsubring{d}}{\textsci}} \dicPos{n}[2] \dicFlx{(‑s, ‑)}[14] \dicFieldCat{jaz.} \dicDirectTranslationCS{předpona, prefix}
\dicEntry[forskot] \dicTerm{for··skot} \dicIPA{{f}{\textopeno}{\textsubring{r}}{s}{\r{g}}{\textopeno}{\textsubring{d}}} \dicPos{n}[2] \dicFlx{(‑s, ‑)}[5] \dicDirectTranslationCS{náskok, výhoda (prvního tahu ap.)} \dicExampleIS{50 metra forskot} \dicExampleCS{náskok 50 metrů}
\dicEntry[forskrift] \dicTerm{for··skrift} \dicIPA{{f}{\textopeno}{\textsubring{r}}{s}{\r{g}}{r}{\textsci}{f}{\textsubring{d}}} \dicPos{f}[7] \dicFlx{(‑ar, ‑ir)}[1] \textbf{1.} \dicDirectTranslationCS{předloha, předepsaný text}  \textbf{2.} \dicSynonym{fyrirmæli} \dicDirectTranslationCS{předpis, pokyn}
\dicEntry[forskrúfaður] \dicTerm{for··skrúf·|aður} \dicIPA{{f}{\textopeno}{\textsubring{r}}{s}{\r{g}}{r}{u}{a}{ð}{\textscy}{\textsubring{r}}} \dicPos{adj}[3] \dicFlx{(f ‑uð)}[3] \textbf{1.} \dicDirectTranslationCS{mající stržený závit}  \textbf{2.} \dicSynonym{truflaður} \dicDirectTranslationCS{vyšinutý, výstřední}
\dicEntry[forsmá] \dicTerm{for··smá} \dicIPA{{f}{\textopeno}{\textsubring{r}}{s}{m}{au}} \dicPos{v}[2] \dicFlx{(‑ði, ‑ð)}[118] \dicFlx{acc} \dicDirectTranslationCS{přehlížet, zlehčovat, brát na lehkou váhu} \dicExampleIS{forsmá reynslu annarra} \dicExampleCS{přehlížet zkušenost druhých}
\dicEntry[forsmán] \dicTerm{for··smán} \dicIPA{{f}{\textopeno}{\textsubring{r}}{s}{m}{au}{\textsubring{n}}} \dicPos{f}[7] \dicFlx{(‑ar, ‑ir)}[1] \dicSynonym{svívirða\smash{\textsuperscript{1}}} \dicDirectTranslationCS{ostuda, pobouření, potupa} \dicExampleIS{þvílík forsmán} \dicExampleCS{taková ostuda}
\dicEntry[forsmekkur] \dicTerm{for··smekk|ur} \dicIPA{{f}{\textopeno}{\textsubring{r}}{s}{m}{\textepsilon}{h}{\r{g}}{\textscy}{\textsubring{r}}} \dicPos{m}[6] \dicFlx{(‑s)}[17] \dicDirectTranslationCS{předzvěst, náznak}
\dicEntry[forspjall] \dicTerm{for··|spjall} \dicIPA{{f}{\textopeno}{\textsubring{r}}{s}{\textsubring{b}}{j}{a}{\textsubring{d}}{\textsubring{l}}} \dicPos{n}[2] \dicFlx{(‑spjalls, ‑spjöll)}[8] \dicDirectTranslationCS{předmluva, prolog, úvodní poznámka}
\dicEntry[forsprakki] \dicTerm{for··sprakk|i} \dicIPA{{f}{\textopeno}{\textsubring{r}}{s}{\textsubring{b}}{r}{a}{h}{\r{\textObardotlessj}}{\textsci}} \dicPos{m}[1] \dicFlx{(‑a, ‑ar)}[8] \textbf{1.} \dicSynonym{upphafsmaður} \dicDirectTranslationCS{podněcovatel(ka), iniciátor(ka)}  \textbf{2.} \dicSynonym{leiðtogi} \dicDirectTranslationCS{vůdce, vůdkyně, vedoucí osoba}
\dicEntry[forspurður] \dicTerm{for··spurður} \dicIPA{{f}{\textopeno}{\textsubring{r}}{s}{\textsubring{b}}{\textscy}{r}{ð}{\textscy}{\textsubring{r}}} \dicPos{adj}[2]\dicFlx{}[-4] \dicPhraseIS{(gera e‑ð) að e‑m forspurðum} \dicDirectTranslationCS{(udělat (co)) bez vědomí (koho)}
\dicEntry[forstaða] \dicTerm{for··|staða} \dicIPA{{f}{\textopeno}{\textsubring{r}}{s}{\textsubring{d}}{a}{ð}{a}} \dicPos{f}[1] \dicFlx{(‑stöðu)}[2] \dicSynonym{forysta} \dicDirectTranslationCS{vedení};  \dicPhraseIS{veita e‑u forstöðu} \dicDirectTranslationCS{řídit (co), vládnout (čemu)} \dicExampleIS{veita fyrirtækinu forstöðu} \dicExampleCS{řídit firmu}
\dicEntry[forstjóri] \dicTerm{for··stjór|i} \dicsymFrequent\  \dicIPA{{f}{\textopeno}{\textsubring{r}}{s}{\textsubring{d}}{j}{ou}{r}{\textsci}} \dicPos{m}[1] \dicFlx{(‑a, ‑ar)}[1] \dicDirectTranslationCS{ředitel(ka)} \dicExampleIS{fjalla um skipun forstjóra} \dicExampleCS{zabývat se jmenováním ředitele}
\dicEntry[forstofa] \dicTerm{for··stof|a} \dicsymFrequent\  \dicIPA{{f}{\textopeno}{\textsubring{r}}{s}{\textsubring{d}}{\textopeno}{v}{a}} \dicPos{f}[1] \dicFlx{(‑u, ‑ur)}[7] \dicDirectTranslationCS{předsíň} \dicExampleIS{Hún sparkaði af sér skónum í forstofunni.} \dicExampleCS{V~předsíni ze sebe setřásla boty.}
\dicEntry[forstöðumaður] \dicTerm{for·stöðu··|maður} \dicIPA{{f}{\textopeno}{\textsubring{r}}{s}{\textsubring{d}}{\oe}{ð}{\textscy}{m}{a}{ð}{\textscy}{\textsubring{r}}} \dicPos{m}[13] \dicFlx{(‑manns, ‑menn)}[2] \dicDirectTranslationCS{ředitel(ka), vedoucí, šéf(ka), představitel(ka)}
\dicEntry[forsvar] \dicTerm{for··svar} \dicIPA{{f}{\textopeno}{\textsubring{r}}{s}{v}{a}{\textsubring{r}}} \dicPos{n}[2] \dicFlx{(‑s)}[2] \dicSynonym{vörn} \dicDirectTranslationCS{reprezentace, zastoupení};  \dicPhraseIS{vera í forsvari fyrir e‑u\,/\addthin e‑m} \dicDirectTranslationCS{reprezentovat (co\,/\addthin koho), zastupovat (co\,/\addthin koho)}
\dicEntry[forsvaranlegur] \dicTerm{for·svaran··legur} \dicIPA{{f}{\textopeno}{\textsubring{r}}{s}{v}{a}{r}{a}{n}{l}{\textepsilon}{\textbabygamma}{\textscy}{\textsubring{r}}} \dicPos{adj}[1]\dicFlx{}[-8] \dicSynonym{afsakanlegur} \dicDirectTranslationCS{omluvitelný, ospravedlnitelný} \dicIndirectTranslationCS{(používané v~otázce nebo záporu)} \dicExampleIS{Þetta er ekki forsvaranleg framkoma.} \dicExampleCS{Toto je neomluvitelné chování.}
\dicEntry[forsvarsmaður] \dicTerm{for·svars··|maður} \dicIPA{{f}{\textopeno}{\textsubring{r}}{s}{v}{a}{\textsubring{r}}{s}{m}{a}{ð}{\textscy}{\textsubring{r}}} \dicPos{m}[13] \dicFlx{(‑manns, ‑menn)}[2] \dicDirectTranslationCS{představitel(ka), reprezentant(ka), zodpovědná osoba}
\dicEntry[forsæla] \dicTerm{for··sæl|a} \dicIPA{{f}{\textopeno}{\textsubring{r}}{s}{a}{i}{l}{a}} \dicPos{f}[1] \dicFlx{(‑u)}[5] \dicSynonym{skuggi} \dicDirectTranslationCS{stín}
\dicEntry[forsæti] \dicTerm{for··sæti} \dicIPA{{f}{\textopeno}{\textsubring{r}}{s}{a}{i}{\textsubring{d}}{\textsci}} \dicPos{n}[2] \dicFlx{(‑s, ‑)}[14] \dicDirectTranslationCS{předsednictví} \dicExampleIS{vera í forsæti á fundinum} \dicExampleCS{předsedat schůzi}
\dicEntry[forsætisráðherra] \dicTerm{for·sætis··ráð·herr|a} \dicIPA{{f}\-{\textopeno}\-{\textsubring{r}}\-{s}\-{a}\-{i}\-{\textsubring{d}}\-{\textsci}\-{s}\-{r}\-{au}\-{\texttheta}\-{h}\-{\textepsilon}\-{r}\-{a}\-} \dicPos{m}[1] \dicFlx{(‑a, ‑ar)}[17] \dicDirectTranslationCS{premiér(ka), předseda\,/\addthin předsedkyně vlády}
\dicEntry[forsætisráðuneyti] \textls[15]{\dicTerm{for·sætis··ráðu·neyti} \dicIPA{{f}\-{\textopeno}\-{\textsubring{r}}\-{s}\-{a}\-{i}\-{\textsubring{d}}\-{\textsci}\-{s}\-{r}\-{au}\-{ð}\-{\textscy}\-{n}\-{ei}\-{\textsubring{d}}\-{\textsci}\-} \dicPos{n}[2] \dicFlx{(‑s, ‑)}[14]} \dicDirectTranslationCS{úřad vlády\,/\addthin premiéra}
\dicEntry[forsögn] \dicTerm{for··|sögn} \dicIPA{{f}{\textopeno}{\textsubring{r}}{s}{\oe}{\r{g}}{\textsubring{n}}} \dicPos{f}[7] \dicFlx{(‑sagnar, ‑sagnir)}[16] \textbf{1.} \dicSynonym{tilvísun} \dicDirectTranslationCS{pokyn, nařízení, instrukce};  \dicPhraseIS{eftir minni forsögn} \dicFlx{adv} \dicDirectTranslationCS{podle mých pokynů}  \textbf{2.} \dicDirectTranslationCS{prognóza, předpověď}
\dicEntry[forsögulegur] \dicTerm{for·sögu··legur} \dicIPA{{f}{\textopeno}{\textsubring{r}}{s}{\oe}{\textbabygamma}{\textscy}{l}{\textepsilon}{\textbabygamma}{\textscy}{\textsubring{r}}} \dicPos{adj}[1]\dicFlx{}[-8] \dicDirectTranslationCS{pravěký, prehistorický}
\dicEntry[forsöngvari] \dicTerm{for··söngv·ar|i} \dicIPA{{f}{\textopeno}{\textsubring{r}}{s}{\oe i}{\ng}{\r{g}}{v}{a}{r}{\textsci}} \dicPos{m}[1] \dicFlx{(‑a, ‑ar)}[13] \dicDirectTranslationCS{předzpěvák, předzpěvačka} \dicExampleIS{forsöngvari í kirkju} \dicExampleCS{předzpěvák v~kostele}
\dicEntry[fortaka] \dicTerm{for··|taka} \dicIPA{{f}{\textopeno}{\textsubring{r}}{t\smash{\textsuperscript{h}}}{a}{\r{g}}{a}} \dicPos{v}[6] \dicFlx{(‑tek, ‑tók, ‑tókum, ‑tæki, ‑tekið)}[128] \dicFlx{acc} \dicSynonym{neita} \dicDirectTranslationCS{popřít, popírat, odmítnout, odmítat} \dicExampleIS{Ég fortek það ekki.} \dicExampleCS{Nepopírám to.}
\dicEntry[fortakslaus] \dicTerm{for·taks··laus} \dicIPA{{f}{\textopeno}{\textsubring{r}}{t\smash{\textsuperscript{h}}}{a}{\r{g}}{s}{l}{\oe i}{s}} \dicPos{adj}[5]\dicFlx{}[-1] \dicDirectTranslationCS{nepopiratelný, jednoznačný}
\dicEntry[fortakslaust] \dicTerm{for·taks··laust} \dicIPA{{f}{\textopeno}{\textsubring{r}}{t\smash{\textsuperscript{h}}}{a}{\r{g}}{s}{l}{\oe i}{s}{\textsubring{d}}} \dicPos{adv} \dicDirectTranslationCS{nepopiratelně, jednoznačně}
\dicEntry[fortek] \dicTerm{for··tek} \dicIPA{{f}{\textopeno}{\textsubring{r}}{t\smash{\textsuperscript{h}}}{\textepsilon}{\r{g}}} \dicPos{v} \dicFlx{ind praes sg 1 pers} \dicLink{fortaka}
\dicEntry[fortekið] \dicTerm{for··tekið} \dicIPA{{f}{\textopeno}{\textsubring{r}}{t\smash{\textsuperscript{h}}}{\textepsilon}{\r{\textObardotlessj}}{\textsci}{\texttheta}} \dicPos{v} \dicFlx{supin} \dicLink{fortaka}
\dicEntry[fortíð] \dicTerm{for··tíð} \dicsymFrequent\  \dicIPA{{f}{\textopeno}{\textsubring{r}}{t\smash{\textsuperscript{h}}}{i}{\texttheta}} \dicPos{f}[7] \dicFlx{(‑ar)}[3] \dicDirectTranslationCS{minulost} \dicExampleIS{Hvað gerðist í fortíðinni?} \dicExampleCS{Co se stalo v~minulosti?}
\dicEntry[fortjald] \dicTerm{for··|tjald} \dicIPA{{f}{\textopeno}{\textsubring{r}}{t\smash{\textsuperscript{h}}}{j}{a}{l}{\textsubring{d}}} \dicPos{n}[2] \dicFlx{(‑tjalds, ‑tjöld)}[8] \textbf{1.} \dicDirectTranslationCS{opona}  \textbf{2.} \dicDirectTranslationCS{závěs}
\dicEntry[fortók] \dicTerm{for··tók} \dicIPA{{f}{\textopeno}{\textsubring{r}}{t\smash{\textsuperscript{h}}}{ou}{\r{g}}} \dicPos{v} \dicFlx{ind pf sg 1 pers} \dicLink{fortaka}
\dicEntry[fortókum] \dicTerm{for··tókum} \dicIPA{{f}{\textopeno}{\textsubring{r}}{t\smash{\textsuperscript{h}}}{ou}{\r{g}}{\textscy}{\textsubring{m}}} \dicPos{v} \dicFlx{ind pf pl 1 pers} \dicLink{fortaka}
\dicEntry[fortæki] \dicTerm{for··tæki} \dicIPA{{f}{\textopeno}{\textsubring{r}}{t\smash{\textsuperscript{h}}}{a}{i}{\r{\textObardotlessj}}{\textsci}} \dicPos{v} \dicFlx{con pf sg 1 pers} \dicLink{fortaka}
\dicEntry[fortölur] \dicTerm{for··tölur} \dicIPA{{f}{\textopeno}{\textsubring{r}}{t\smash{\textsuperscript{h}}}{\oe}{l}{\textscy}{\textsubring{r}}} \dicPos{f}[1] \dicFlx{pl}[15] \dicDirectTranslationCS{přesvědčování, přemlouvání} \dicExampleIS{Fortölur hans urðu árangurslausar.} \dicExampleCS{Jeho přemlouvání bylo neúspěšné.}
\dicEntry[forugur] \dicTerm{for··ugur} \dicIPA{{f}{\textopeno}{\textlengthmark}{r}{\textscy}{\textbabygamma}{\textscy}{\textsubring{r}}} \dicPos{adj}[1]\dicFlx{}[-8] \dicSynonym{skítugur} \dicDirectTranslationCS{zablácený, zabahněný, blátivý}
\dicEntry[forusta] \dicTerm{for··ust|a} \dicIPA{{f}{\textopeno}{\textlengthmark}{r}{\textscy}{s}{\textsubring{d}}{a}} \dicPos{f}[1] \dicFlx{(‑u)}[5] \dicLink{forysta}
\dicEntry[forveri] \dicTerm{for··ver|i} \dicIPA{{f}{\textopeno}{r}{v}{\textepsilon}{r}{\textsci}} \dicPos{m}[1] \dicFlx{(‑a, ‑ar)}[1] \dicDirectTranslationCS{předchůdce, předchůdkyně}
\dicEntry[forvextir] \dicTerm{for··vextir} \dicIPA{{f}{\textopeno}{r}{v}{\textepsilon}{x}{s}{\textsubring{d}}{\textsci}{\textsubring{r}}} \dicPos{m}[11] \dicFlx{pl}[1] \dicFieldCat{ekon.} \dicDirectTranslationCS{předem placený úrok}
\dicEntry[forviða] \dicTerm{for··viða} \dicsymFrequent\  \dicIPA{{f}{\textopeno}{r}{v}{\textsci}{ð}{a}} \dicPos{adj}[13] \dicFlx{indecl}[1] \dicDirectTranslationCS{zaskočený, ohromený, užaslý, překvapený} \dicExampleIS{Hún varð forviða.} \dicExampleCS{Zůstala ohromená.}
\dicEntry[forvitinn] \dicTerm{for··vitinn} \dicsymFrequent\  \dicIPA{{f}{\textopeno}{r}{v}{\textsci}{\textsubring{d}}{\textsci}{\textsubring{n}}} \dicPos{adj}[6]\dicFlx{}[-2] \dicDirectTranslationCS{zvědavý, zvídavý, všetečný} \dicExampleIS{virða hann fyrir sér forvitnum augum} \dicExampleCS{zvědavě si ho prohlédnout}
\dicEntry[forvitnast] \dicTerm{for··vitn|ast} \dicIPA{{f}{\textopeno}{r}{v}{\textsci}{h}{\textsubring{d}}{n}{a}{s}{\textsubring{d}}} \dicPos{v}[1] \dicFlx{(‑aðist)}[91] \dicFlx{refl} \dicDirectTranslationCS{vyzvídat, dotazovat se, informovat se};  \dicPhraseIS{forvitnast um e‑ð} \dicDirectTranslationCS{dotazovat se na (co), informovat se o~(čem)}
\dicEntry[forvitni] \dicTerm{for··vitn|i} \dicsymFrequent\  \dicIPA{{f}{\textopeno}{r}{v}{\textsci}{h}{\textsubring{d}}{n}{\textsci}} \dicPos{f}[3] \dicFlx{(‑i)}[3] \dicDirectTranslationCS{zvědavost, zvídavost, všetečnost};  \dicPhraseIS{e‑m leikur forvitni á e‑u} \dicFlx{impers} \dicDirectTranslationCS{(kdo) je zvědavý na (co)}
\dicEntry[forvitnilegur] \dicTerm{for·vitni··legur} \dicIPA{{f}{\textopeno}{r}{v}{\textsci}{h}{\textsubring{d}}{n}{\textsci}{l}{\textepsilon}{\textbabygamma}{\textscy}{\textsubring{r}}} \dicPos{adj}[1]\dicFlx{}[-8] \dicSynonym{hnýsilegur} \dicDirectTranslationCS{zajímavý, pozoruhodný}
\dicEntry[forvöð] \dicTerm{for··vöð} \dicIPA{{f}{\textopeno}{r}{v}{\oe}{\texttheta}} \dicPos{n}[2] \dicFlx{pl}[9] \dicPhraseIS{það eru síðustu forvöð að (gera e‑ð)} \dicDirectTranslationCS{to je poslední šance (udělat (co))}
\dicEntry[forynja] \dicTerm{for··ynj|a} \dicIPA{{f}{\textopeno}{\textlengthmark}{r}{\textsci}{n}{j}{a}} \dicPos{f}[1] \dicFlx{(‑u, ‑ur)}[7] \dicSynonym{ófreskja} \dicDirectTranslationCS{netvor, příšera, obluda, zrůda, monstrum, nestvůra}
\dicEntry[forysta] \dicTerm{for··yst|a}\dicTerm{, forusta} \dicIPA{{f}\-{\textopeno}\-{\textlengthmark}\-{r}\-{\textsci}\-{s}\-{\textsubring{d}}\-{a}\-} \dicPos{f}[1] \dicFlx{(‑u)}[5] \textbf{1.} \dicDirectTranslationCS{vedení, vůdcovství};  \dicPhraseIS{vera í forystu fyrir e‑u} \dicDirectTranslationCS{vést (co)};  \dicPhraseIS{undir forystu e‑rs} \dicDirectTranslationCS{pod (čím) vedením}  \textbf{2.} \dicDirectTranslationCS{vedení, náskok} \dicExampleIS{hafa forystu} \dicExampleCS{mít náskok}
\dicEntry[forþekking] \dicTerm{for··þekk·ing} \dicIPA{{f}{\textopeno}{\textsubring{r}}{\texttheta}{\textepsilon}{h}{\r{\textObardotlessj}}{i}{\ng}{\r{g}}} \dicPos{f}[4] \dicFlx{(‑ar)}[7] \dicDirectTranslationCS{(předchozí) informovanost, povědomost}
\dicEntry[fosfór] \dicTerm{fosfór} \dicIPA{{f}{\textopeno}{s}{f}{ou}{\textsubring{r}}} \dicPos{m}[4] \dicFlx{(‑s)}[18] \dicFieldCat{chem.} \dicDirectTranslationCS{fosfor} \textit{(l.~{\textLA{P, Phosphorum}})}
\dicEntry[foss] \dicTerm{foss} \dicsymFrequent\  \dicIPA{{f}{\textopeno}{s}{\textlengthmark}} \dicPos{m}[4] \dicFlx{(‑, ‑ar)}[28] \dicDirectTranslationCS{vodopád} \dicExampleIS{Glymur er hæsti foss landsins.} \dicExampleCS{Glymur je nejvyšší vodopád v~zemi.}
\dicEntry[fossa] \dicTerm{foss|a} \dicIPA{{f}{\textopeno}{s}{\textlengthmark}{a}} \dicPos{v}[1] \dicFlx{(‑aði)}[1] \dicDirectTranslationCS{stříkat, vystřikovat} \dicExampleIS{Blóðið fossar úr sárinu.} \dicExampleCS{Krev stříká z~rány.}
\dicEntry[fossasig] \dicTerm{fossa··sig} \dicIPA{{f}{\textopeno}{s}{\textlengthmark}{a}{s}{\textsci}{x}} \dicPos{n}[2] \dicFlx{(‑s)}[2] \dicDirectTranslationCS{slaňování vodopádu}
\dicEntry[fóbía] \dicTerm{fóbí|a} \dicIPA{{f}{ou}{\textlengthmark}{\textsubring{b}}{i}{j}{a}} \dicPos{f}[1] \dicFlx{(‑u)}[5] \dicDirectTranslationCS{fobie, (chorobný) strach, bázeň}
\dicEntry[fóðra] \dicTerm{fóðr|a} \dicIPA{{f}{ou}{ð}{r}{a}} \dicPos{v}[1] \dicFlx{(‑aði)}[1] \dicFlx{acc} \textbf{1.} \dicSynonym*{gefa fæðu} \dicDirectTranslationCS{(na)krmit} \dicExampleIS{fóðra fugla} \dicExampleCS{krmit ptáky}  \textbf{2.} \dicSynonym*{fóðra flík} \dicDirectTranslationCS{podšít, podšívat} \dicExampleIS{fóðra jakkann} \dicExampleCS{podšít bundu}  \textbf{3.} \dicDirectTranslationCS{vyložit} \dicIndirectTranslationCS{(pokrýt vnitřní stranu\,/\addthin povrch něčím)}
\dicEntry[fóðraður] \dicTerm{fóðr··|aður} \dicIPA{{f}{ou}{ð}{r}{a}{ð}{\textscy}{\textsubring{r}}} \dicPos{adj}[3] \dicFlx{(f ‑uð)}[3] \dicDirectTranslationCS{podšitý, mající podšívku\,/\addthin vložku}
\dicEntry[fóður] \dicTerm{fóður} \dicIPA{{f}{ou}{\textlengthmark}{ð}{\textscy}{\textsubring{r}}} \dicPos{n}[2] \dicFlx{(‑s, ‑)}[25] \textbf{1.} \dicDirectTranslationCS{krmení, krmivo} \dicExampleIS{skepnufóður} \dicExampleCS{krmení pro zvířata}  \textbf{2.} \dicDirectTranslationCS{podšívka} \dicExampleIS{jakkafóður} \dicExampleCS{podšívka bundy}  \textbf{3.} \dicSynonym{veggfóður} \dicDirectTranslationCS{tapeta}
\dicEntry[fógeti] \dicTerm{fóget|i} \dicIPA{{f}{ou}{\textlengthmark}{\textbabygamma}{\textepsilon}{\textsubring{d}}{\textsci}} \dicPos{m}[1] \dicFlx{(‑a, ‑ar)}[1] \dicFieldCat{práv.} \dicDirectTranslationCS{městský správce} \dicIndirectTranslationCS{(úředník se soudní pravomocí)}
\dicEntry[fókus] \dicTerm{fókus} \dicIPA{{f}{ou}{\textlengthmark}{\r{g}}{\textscy}{s}} \dicPos{m}[10] \dicFlx{(‑s\,/\addthin ‑ar)}[35] \dicDirectTranslationCS{zaostření}
\dicEntry[fól] \dicTerm{fól\smash{\textsuperscript{1}}} \dicIPA{{f}{ou}{\textlengthmark}{\textsubring{l}}} \dicPos{n}[2] \dicFlx{(‑s, ‑)}[5] \dicSynonym{fauti} \dicDirectTranslationCS{brutální\,/\addthin surový člověk, surovec}
\dicEntry[fól] \dicTerm{fól\smash{\textsuperscript{2}}} \dicIPA{{f}{ou}{\textlengthmark}{\textsubring{l}}} \dicPos{v} \dicFlx{ind pf sg 1 pers} \dicLink{fela\smash{\textsuperscript{2}}}
\dicEntry[fólginn] \dicTerm{fólginn} \dicIPA{{f}{ou}{l}{\r{\textObardotlessj}}{\textsci}{\textsubring{n}}} \dicPos{adj}[6]\dicFlx{}[-2] \textbf{1.} \dicSynonym*{falinn} \dicDirectTranslationCS{schovaný, skrytý}  \textbf{2.} \dicPhraseIS{e‑að er fólgið í e‑u} \dicDirectTranslationCS{(co) spočívá v~(čem), (co) je obsaženo v~(čem)}
\dicEntry[fólk] \dicTerm{fólk} \dicsymFrequent\  \dicIPA{{f}{ou}{\textsubring{l}}{\r{g}}} \dicPos{n}[2] \dicFlx{(‑s)}[2] \textbf{1.} \dicSynonym*{manneskjur} \dicDirectTranslationCS{lidé, lidi};  \dicPhraseIS{fullorðna fólkið} \dicDirectTranslationCS{dospělí};  \dicPhraseIS{unga fólkið} \dicDirectTranslationCS{mladí lidé, mládež}  \textbf{2.} \dicSynonym*{ættingjar} \dicDirectTranslationCS{příbuzní, rodina}
\dicEntry[fólksfjöldi] \dicTerm{fólks··fjöld|i} \dicIPA{{f}{ou}{\textsubring{l}}{\r{g}}{s}{f}{j}{\oe}{l}{\textsubring{d}}{\textsci}} \dicPos{m}[1] \dicFlx{(‑a)}[3] \textbf{1.} \dicSynonym*{fjöldatala} \dicDirectTranslationCS{počet obyvatel, populace}  \textbf{2.} \dicSynonym*{hópur fólks} \dicDirectTranslationCS{dav\,/\addthin zástup lidí}
\dicEntry[fólksfjölgun] \dicTerm{fólks··fjölg|un} \dicIPA{{f}{ou}{\textsubring{l}}{\r{g}}{s}{f}{j}{\oe}{l}{\r{g}}{\textscy}{\textsubring{n}}} \dicPos{f}[7] \dicFlx{(‑unar)}[12] \dicDirectTranslationCS{nárůst počtu obyvatel}
\dicEntry[fólksflutningur] \dicTerm{fólks··flut·ning|ur} \dicIPA{{f}{ou}{\textsubring{l}}{\r{g}}{s}{f}{l}{\textscy}{h}{\textsubring{d}}{n}{i}{\ng}{\r{g}}{\textscy}{\textsubring{r}}} \dicPos{m}[6] \dicFlx{(‑s, ‑ar)}[8] \textbf{1.} \dicDirectTranslationCS{přesídlení, migrace (lidí)}  \textbf{2.} \dicPhraseIS{fólksflutningar} \dicFlx{pl} \dicDirectTranslationCS{osobní doprava, přeprava\,/\addthin převoz osob}
\dicEntry[fólksfækkun] \dicTerm{fólks··fækk|un} \dicIPA{{f}{ou}{\textsubring{l}}{\r{g}}{s}{f}{a}{i}{h}{\r{g}}{\textscy}{\textsubring{n}}} \dicPos{f}[7] \dicFlx{(‑unar)}[9] \dicDirectTranslationCS{pokles\,/\addthin úbytek počtu obyvatel}
\dicEntry[fólska] \dicTerm{fólsk|a} \dicIPA{{f}{ou}{l}{s}{\r{g}}{a}} \dicPos{f}[1] \dicFlx{(‑u)}[5] \dicSynonym{harðleikni} \dicDirectTranslationCS{surovost, brutalita}
\dicEntry[fólum] \dicTerm{fólum} \dicIPA{{f}{ou}{\textlengthmark}{l}{\textscy}{\textsubring{m}}} \dicPos{v} \dicFlx{ind pf pl 1 pers} \dicLink{fela\smash{\textsuperscript{2}}}
\dicEntry[fónem] \dicTerm{fónem} \dicIPA{{f}{ou}{\textlengthmark}{n}{\textepsilon}{\textsubring{m}}} \dicPos{n}[2] \dicFlx{(‑s, ‑)}[5] \dicFieldCat{jaz.} \dicDirectTranslationCS{foném}
\dicEntry[fónn] \dicTerm{fón|n} \dicIPA{{f}{ou}{\textsubring{d}}{\textsubring{n}}} \dicPos{m}[6] \dicFlx{(‑s, ‑ar)}[42] \textbf{1.} \dicSynonym*{grammófónn} \dicDirectTranslationCS{gramofon}  \textbf{2.} \dicSynonym{diktafónn} \dicDirectTranslationCS{diktafon}
\dicEntry[fór] \dicTerm{fór} \dicIPA{{f}{ou}{\textlengthmark}{\textsubring{r}}} \dicPos{v} \dicFlx{ind pf sg 1 pers} \dicLink{fara}
\dicEntry[fórn] \dicTerm{fórn} \dicIPA{{f}{ou}{r}{\textsubring{d}}{\textsubring{n}}} \dicPos{f}[7] \dicFlx{(‑ar, ‑ir)}[1] \textbf{1.} \dicSynonym*{gjöf til guðs} \dicDirectTranslationCS{oběť} \dicIndirectTranslationCS{(dar zaslíbený božstvu)};  \dicPhraseIS{færa e‑m e‑ð að fórn} \dicDirectTranslationCS{obětovat (co) pro (koho), obětovat (co komu)}  \textbf{2.} \dicDirectTranslationCS{oběť, obětování} \dicIndirectTranslationCS{(obřad)}
\dicEntry[fórna] \dicTerm{fórn|a} \dicsymFrequent\  \dicIPA{{f}{ou}{r}{\textsubring{d}}{n}{a}} \dicPos{v}[1] \dicFlx{(‑aði)}[1] \dicFlx{dat} \textbf{1.} \dicDirectTranslationCS{obětovat, přinést oběť (Bohu ap.)} \dicExampleIS{fórna guði dýri} \dicExampleCS{obětovat Bohu zvíře}  \textbf{2.} \dicDirectTranslationCS{obětovat, dát jako oběť};  \dicPhraseIS{fórna lífi sínu, fórna lífinu} \dicDirectTranslationCS{obětovat (svůj) život};  \dicPhraseIS{fórna sér} \dicDirectTranslationCS{obětovat se} \dicExampleIS{Hann fórnaði sjálfum sér fyrir okkur.} \dicExampleCS{Obětoval se pro nás.}
\dicEntry[fórnargjöf] \dicTerm{fórnar··|gjöf} \dicIPA{{f}{ou}{r}{\textsubring{d}}{n}{a}{\textsubring{r}}{\r{\textObardotlessj}}{\oe}{f}} \dicPos{f}[7] \dicFlx{(‑gjafar, ‑gjafir)}[16] \dicSynonym{fórn} \dicDirectTranslationCS{obětina, obětní dar}
\dicEntry[fórnarlamb] \dicTerm{fórnar··|lamb} \dicIPA{{f}{ou}{r}{\textsubring{d}}{n}{a}{r}{l}{a}{m}{\textsubring{b}}} \dicPos{n}[2] \dicFlx{(‑lambs, ‑lömb)}[8] \dicDirectTranslationCS{oběť (postižený), obětní beránek} \dicExampleIS{fórnarlamb nasista} \dicExampleCS{oběť nacistů}
\dicEntry[fórnfús] \dicTerm{fórn··fús} \dicIPA{{f}{ou}{r}{\textsubring{d}}{\textsubring{n}}{f}{u}{s}} \dicPos{adj}[5]\dicFlx{}[-1] \dicSynonym{hjálpsamur} \dicDirectTranslationCS{obětavý, nesobecký, altruistický}
\dicEntry[fórnfýsi] \dicTerm{fórn··fýs|i} \dicIPA{{f}{ou}{r}{\textsubring{d}}{\textsubring{n}}{f}{i}{s}{\textsci}} \dicPos{f}[3] \dicFlx{(‑i)}[3] \dicDirectTranslationCS{obětavost, nesobeckost, altruismus}
\dicEntry[fórum] \dicTerm{fórum} \dicIPA{{f}{ou}{r}{\textlengthmark}{\textscy}{\textsubring{m}}} \dicPos{v} \dicFlx{ind pf pl 1 pers} \dicLink{fara}
\dicEntry[fóstbróðir] \dicTerm{fóst··|bróðir} \dicIPA{{f}{ou}{s}{\textsubring{d}}{\textsubring{b}}{r}{ou}{ð}{\textsci}{\textsubring{r}}} \dicPos{m}[13] \dicFlx{(‑bróður, ‑bræður)}[4] \dicSynonym*{fósturbróðir} \dicDirectTranslationCS{pokrevní bratr}
\dicEntry[fóstra] \dicTerm{fóstr|a\smash{\textsuperscript{1}}} \dicIPA{{f}{ou}{s}{\textsubring{d}}{r}{a}} \dicPos{f}[1] \dicFlx{(‑u, ‑ur)}[7] \textbf{1.} \dicSynonym{fósturmóðir} \dicDirectTranslationCS{pěstounka}  \textbf{2.} \dicFieldCat{škol.} \dicDirectTranslationCS{učitelka v~mateřské školce}
\dicEntry[fóstra] \dicTerm{fóstr|a\smash{\textsuperscript{2}}} \dicIPA{{f}{ou}{s}{\textsubring{d}}{r}{a}} \dicPos{v}[1] \dicFlx{(‑aði)}[1] \dicFlx{acc} \dicDirectTranslationCS{opatrovat, pečovat (o~cizí dítě ap.)}
\dicEntry[fóstri] \dicTerm{fóstr|i} \dicIPA{{f}{ou}{s}{\textsubring{d}}{r}{\textsci}} \dicPos{m}[1] \dicFlx{(‑a, ‑ar)}[1] \textbf{1.} \dicSynonym{fósturfaðir} \dicDirectTranslationCS{pěstoun}  \textbf{2.} \dicSynonym{fóstursonur} \dicDirectTranslationCS{schovanec}  \textbf{3.} \dicFieldCat{škol.} \dicDirectTranslationCS{učitel v~mateřské školce}
\dicEntry[fóstur] \dicTerm{fóstur} \dicIPA{{f}{ou}{s}{\textsubring{d}}{\textscy}{\textsubring{r}}} \dicPos{n}[2] \dicFlx{(‑s, ‑)}[25] \textbf{1.} \dicSynonym{uppeldi} \dicDirectTranslationCS{pěstounství, opatrovnictví} \dicExampleIS{taka e‑n í fóstur} \dicExampleCS{vzít (koho) do pěstounské péče}  \textbf{2.} \dicDirectTranslationCS{embryo, zárodek} \dicExampleIS{Fóstrið þroskast eðlilega.} \dicExampleCS{Embryo se vyvíjí přirozeně.};  \dicPhraseIS{eyða fóstri} \dicDirectTranslationCS{provést potrat};  \dicPhraseIS{missa fóstur} \dicDirectTranslationCS{potratit}
\dicEntry[fósturbarn] \dicTerm{fóstur··|barn} \dicIPA{{f}{ou}{s}{\textsubring{d}}{\textscy}{r}{\textsubring{b}}{a}{r}{\textsubring{d}}{\textsubring{n}}} \dicPos{n}[2] \dicFlx{(‑barns, ‑börn)}[8] \dicDirectTranslationCS{schovanec, schovanka}
\dicEntry[fósturdóttir] \dicTerm{fóstur··|dóttir} \dicIPA{{f}{ou}{s}{\textsubring{d}}{\textscy}{r}{\textsubring{d}}{ou}{h}{\textsubring{d}}{\textsci}{\textsubring{r}}} \dicPos{f}[11] \dicFlx{(‑dóttur, ‑dætur)}[2] \dicDirectTranslationCS{schovanka}
\dicEntry[fóstureyðing] \dicTerm{fóstur··eyð·ing} \dicIPA{{f}{ou}{s}{\textsubring{d}}{\textscy}{r}{ei}{ð}{i}{\ng}{\r{g}}} \dicPos{f}[4] \dicFlx{(‑ar, ‑ar)}[5] \dicFieldCat{med.} \dicDirectTranslationCS{potrat}
\dicEntry[fósturfaðir] \dicTerm{fóstur··|faðir} \dicIPA{{f}{ou}{s}{\textsubring{d}}{\textscy}{\textsubring{r}}{f}{a}{ð}{\textsci}{\textsubring{r}}} \dicPos{m}[13] \dicFlx{(‑föður, ‑feður)}[3] \dicDirectTranslationCS{pěstoun, opatrovník}
\dicEntry[fósturfræði] \dicTerm{fóstur··fræð|i} \dicIPA{{f}{ou}{s}{\textsubring{d}}{\textscy}{\textsubring{r}}{f}{r}{a}{i}{ð}{\textsci}} \dicPos{f}[3] \dicFlx{(‑i)}[3] \dicDirectTranslationCS{embryologie}
\dicEntry[fósturfræðingur] \dicTerm{fóstur·fræð··ing|ur} \dicIPA{{f}{ou}{s}{\textsubring{d}}{\textscy}{\textsubring{r}}{f}{r}{a}{i}{ð}{i}{\ng}{\r{g}}{\textscy}{\textsubring{r}}} \dicPos{m}[6] \dicFlx{(‑s, ‑ar)}[8] \dicDirectTranslationCS{embryolog, embryoložka}
\dicEntry[fósturheimili] \dicTerm{fóstur··heimili} \dicIPA{{f}{ou}{s}{\textsubring{d}}{\textscy}{\textsubring{r}}{h}{ei}{m}{\textsci}{l}{\textsci}} \dicPos{n}[2] \dicFlx{(‑s, ‑)}[14] \dicDirectTranslationCS{pěstounská péče, náhradní rodinná péče}
\dicEntry[fósturjörð] \dicTerm{fóstur··|jörð} \dicIPA{{f}{ou}{s}{\textsubring{d}}{\textscy}{r}{j}{\oe}{r}{\texttheta}} \dicPos{f}[7] \dicFlx{(‑jarðar)}[15] \dicDirectTranslationCS{rodná země, rod\-ný kraj, domovina}
\dicEntry[fósturlát] \dicTerm{fóstur··lát} \dicIPA{{f}{ou}{s}{\textsubring{d}}{\textscy}{r}{l}{au}{\textsubring{d}}} \dicPos{n}[2] \dicFlx{(‑s, ‑)}[5] \dicFieldCat{med. } \dicDirectTranslationCS{(samovolný) potrat}
\dicEntry[fósturmóðir] \dicTerm{fóstur··|móðir} \dicIPA{{f}{ou}{s}{\textsubring{d}}{\textscy}{r}{m}{ou}{ð}{\textsci}{\textsubring{r}}} \dicPos{f}[11] \dicFlx{(‑móður, ‑mæður)}[1] \dicDirectTranslationCS{pěstounka, opatrovnice}
\dicEntry[fóstursonur] \dicTerm{fóstur··|sonur} \dicIPA{{f}{ou}{s}{\textsubring{d}}{\textscy}{\textsubring{r}}{s}{\textopeno}{n}{\textscy}{\textsubring{r}}} \dicPos{m}[12] \dicFlx{(‑sonar, ‑synir)}[9] \dicDirectTranslationCS{schovanec}
\dicEntry[fósturstig] \dicTerm{fóstur··stig} \dicIPA{{f}{ou}{s}{\textsubring{d}}{\textscy}{\textsubring{r}}{s}{\textsubring{d}}{\textsci}{x}} \dicPos{n}[2] \dicFlx{(‑s)}[2] \dicFieldCat{med.} \dicDirectTranslationCS{prenatální stadium}
\dicEntry[fótalaus] \dicTerm{fóta··laus} \dicIPA{{f}{ou}{\textlengthmark}{\textsubring{d}}{a}{l}{\oe i}{s}} \dicPos{adj}[5]\dicFlx{}[-1] \textbf{1.} \dicSynonym*{með engafætur} \dicDirectTranslationCS{beznohý, (jsoucí) bez noh}  \textbf{2.} \dicSynonym{þreyttur} \dicDirectTranslationCS{uchozený, utahaný}
\dicEntry[fótaskortur] \dicTerm{fóta··skort|ur} \dicIPA{{f}{ou}{\textlengthmark}{\textsubring{d}}{a}{s}{\r{g}}{\textopeno}{\textsubring{r}}{\textsubring{d}}{\textscy}{\textsubring{r}}} \dicPos{m}[6] \dicFlx{(‑s)}[7] \dicDirectTranslationCS{zakopnutí};  \dicPhraseIS{e‑m verður fótaskortur} \dicFlx{impers} \dicDirectTranslationCS{(kdo) zakopne}
\dicEntry[fótatak] \dicTerm{fóta··tak} \dicsymFrequent\  \dicIPA{{f}{ou}{\textlengthmark}{\textsubring{d}}{a}{t\smash{\textsuperscript{h}}}{a}{\r{g}}} \dicPos{n}[2] \dicFlx{(‑s)}[2] \dicDirectTranslationCS{kroky} \dicIndirectTranslationCS{(zvuk kroků)} \dicExampleIS{heyra fótatak} \dicExampleCS{slyšet kroky}
\dicEntry[fótboltalið] \dicTerm{fót·bolta··lið} \dicIPA{{f}{ou}{\textlengthmark}{\textsubring{d}}{\textsubring{b}}{\textopeno}{\textsubring{l}}{\textsubring{d}}{a}{l}{\textsci}{\texttheta}} \dicPos{n}[2] \dicFlx{(‑s, ‑)}[5] \dicDirectTranslationCS{fotbalové družstvo, fotbalový tým}
\dicEntry[fótboltamaður] \dicTerm{fót·bolta··|maður} \dicIPA{{f}{ou}{\textlengthmark}{\textsubring{d}}{\textsubring{b}}{\textopeno}{\textsubring{l}}{\textsubring{d}}{a}{m}{a}{ð}{\textscy}{\textsubring{r}}} \dicPos{m}[13] \dicFlx{(‑manns, ‑menn)}[2] \dicDirectTranslationCS{fotbalista, fotbalistka}
\dicEntry[fótboltavöllur] \dicTerm{fót·bolta··|völlur} \dicIPA{{f}{ou}{\textlengthmark}{\textsubring{d}}{\textsubring{b}}{\textopeno}{\textsubring{l}}{\textsubring{d}}{a}{v}{\oe}{\textsubring{d}}{l}{\textscy}{\textsubring{r}}} \dicPos{m}[11] \dicFlx{(‑vallar, ‑vellir)}[5] \dicDirectTranslationCS{fotbalové hřiště, fotbalový stadion}
\dicEntry[fótbolti] \dicTerm{fót··bolt|i} \dicIPA{{f}{ou}{\textlengthmark}{\textsubring{d}}{\textsubring{b}}{\textopeno}{\textsubring{l}}{\textsubring{d}}{\textsci}} \dicPos{m}[1] \dicFlx{(‑a, ‑ar)}[1] \textbf{1.} \dicFieldCat{sport.} \dicSynonym{knattspyrna} \dicDirectTranslationCS{fotbal, kopaná} \dicExampleIS{spila\,/\addthin leika fótbolta} \dicExampleCS{hrát fotbal}  \textbf{2.} \dicFieldCat{sport.} \dicDirectTranslationCS{fotbalový míč}
\dicEntry[fótbremsa] \dicTerm{fót··brems|a} \dicIPA{{f}{ou}{\textlengthmark}{\textsubring{d}}{\textsubring{b}}{r}{\textepsilon}{m}{s}{a}} \dicPos{f}[1] \dicFlx{(‑u, ‑ur)}[7] \dicDirectTranslationCS{nožní brzda}
\dicEntry[fótbrotna] \dicTerm{fót··brotn|a} \dicIPA{{f}{ou}{\textlengthmark}{\textsubring{d}}{\textsubring{b}}{r}{\textopeno}{h}{\textsubring{d}}{n}{a}} \dicPos{v}[1] \dicFlx{(‑aði)}[84] \dicFlx{acc} \dicDirectTranslationCS{zlomit si nohu} \dicExampleIS{fótbrotna sig} \dicExampleCS{zlomit si nohu}
\dicEntry[fótfesta] \dicTerm{fót··fest|a} \dicIPA{{f}{ou}{\textlengthmark}{\textsubring{d}}{f}{\textepsilon}{s}{\textsubring{d}}{a}} \dicPos{f}[1] \dicFlx{(‑u)}[5] \dicDirectTranslationCS{opora\,/\addthin opěra pro nohu};  \dicPhraseIS{missa fótfestuna} {\textbf{a.}} \dicDirectTranslationCS{ztratit oporu (pro nohu)};  {\textbf{b.}} \dicLangCat{přen.} \dicDirectTranslationCS{ztratit pevnou půdu pod nohama};  \dicPhraseIS{ná fótfestu} \dicLangCat{přen.} \dicDirectTranslationCS{získat pevnou půdu pod nohama}
\dicEntry[fótgangandi] \dicTerm{fót··gang·|andi\smash{\textsuperscript{1}}} \dicIPA{{f}{ou}{\textlengthmark}{\textsubring{d}}{\r{g}}{au}{\ng}{\r{g}}{a}{n}{\textsubring{d}}{\textsci}} \dicPos{m}[2] \dicFlx{(‑anda, ‑endur)}[1] \dicDirectTranslationCS{chodec, chodkyně, pěší}
\dicEntry[fótgangandi] \dicTerm{fót··gang·andi\smash{\textsuperscript{2}}} \dicIPA{{f}{ou}{\textlengthmark}{\textsubring{d}}{\r{g}}{au}{\ng}{\r{g}}{a}{n}{\textsubring{d}}{\textsci}} \dicPos{adj}[13] \dicFlx{indecl}[1] \dicDirectTranslationCS{jdoucí pěšky}
\dicEntry[fótgöngulið] \dicTerm{fót·göngu··lið} \dicIPA{{f}{ou}{\textlengthmark}{\textsubring{d}}{\r{g}}{\oe i}{\ng}{\r{g}}{\textscy}{l}{\textsci}{\texttheta}} \dicPos{n}[2] \dicFlx{(‑s, ‑)}[5] \dicFieldCat{voj.} \dicDirectTranslationCS{pěchota, infanterie}
\dicEntry[fótgönguliði] \dicTerm{fót·göngu··lið|i} \dicIPA{{f}{ou}{\textlengthmark}{\textsubring{d}}{\r{g}}{\oe i}{\ng}{\r{g}}{\textscy}{l}{\textsci}{ð}{\textsci}} \dicPos{m}[1] \dicFlx{(‑a, ‑ar)}[1] \dicFieldCat{voj.} \dicDirectTranslationCS{pěšák, voják(yně) pěchoty}
\dicEntry[fótleggur] \dicTerm{fót··legg|ur} \dicsymFrequent\  \dicIPA{{f}{ou}{\textlengthmark}{\textsubring{d}}{l}{\textepsilon}{\r{g}}{\textscy}{\textsubring{r}}} \dicPos{m}[9] \dicFlx{(‑jar\,/\addthin ‑s, ‑ir)}[26] \dicFieldCat{anat.} \dicDirectTranslationCS{noha, bérec} \dicIndirectTranslationCS{(dolní část nohy)} \dicExampleIS{hreyfingar fótleggja} \dicExampleCS{pohyby nohou}
\dicEntry[fótsnyrting] \dicTerm{fót··snyrt·ing} \dicIPA{{f}{ou}{\textlengthmark}{\textsubring{d}}{s}{\textsubring{d}}{n}{\textsci}{\textsubring{r}}{\textsubring{d}}{i}{\ng}{\r{g}}} \dicPos{f}[4] \dicFlx{(‑ar)}[7] \dicDirectTranslationCS{pedikúra} \dicIndirectTranslationCS{(odborné ošetření nohou a~nehtů na nohou)}
\dicEntry[fótspor] \dicTerm{fót··spor} \dicIPA{{f}{ou}{\textlengthmark}{\textsubring{d}}{s}{\textsubring{b}}{\textopeno}{\textsubring{r}}} \dicPos{n}[2] \dicFlx{(‑s, ‑)}[5] \dicDirectTranslationCS{stopa, šlápota, šlépěj} \dicIndirectTranslationCS{(otisk chodidla)}
\dicEntry[fótstig] \dicTerm{fót··stig} \dicIPA{{f}{ou}{\textlengthmark}{\textsubring{d}}{s}{\textsubring{d}}{\textsci}{x}} \dicPos{n}[2] \dicFlx{(‑s, ‑)}[5] \dicDirectTranslationCS{pedál, šlapka}
\dicEntry[fótur] \dicTerm{fótur} \dicsymFrequent\  \dicIPA{{f}{ou}{\textlengthmark}{\textsubring{d}}{\textscy}{\textsubring{r}}} \dicPos{m}[13] \dicFlx{(fótar, fætur)}[5] \textbf{1.} \dicFieldCat{anat.} \dicDirectTranslationCS{noha}  \textbf{2.} \dicDirectTranslationCS{noha (u~židle ap.)} \dicExampleIS{Borðið stóð á fjórum fótum.} \dicExampleCS{Stůl stál na čtyřech nohách.}  \textbf{3.} \dicDirectTranslationCS{sokl, podstavec};  \dicPhraseIS{bregða fæti fyrir e‑n} \dicLangCat{přen.} \dicDirectTranslationCS{podrážet (komu) nohy};  \dicPhraseIS{e‑að stendur höllum fæti} \dicLangCat{přen.} \dicDirectTranslationCS{(co) stagnuje, (co) přešlapuje na místě};  \dicPhraseIS{e‑að stendur traustum fótum} \dicLangCat{přen.} \dicDirectTranslationCS{(co) stojí na pevných základech};  \dicPhraseIS{fara á fætur} \dicDirectTranslationCS{vstát (z~postele)};  \dicPhraseIS{gefa e‑m undir fótinn} \dicLangCat{přen.} \dicSynonym*{reyna við e‑m} \dicDirectTranslationCS{flirtovat s~(kým)} \dicExampleIS{Strákurinn gefur stúlkunni undir fótinn.} \dicExampleCS{Chlapec flirtuje s~dívkou.};  \dicPhraseIS{setja á fót e‑u} \dicDirectTranslationCS{založit (co), postavit (co) na nohy};  \dicPhraseIS{taka til fótanna, taka undir sig fætur, kasta fótum undir sig} \dicDirectTranslationCS{rozběhnout se, začít běžet};  \dicPhraseIS{vita ekki í hvorn fótinn maður á að stíga} \dicLangCat{přen.} \dicDirectTranslationCS{nevědět, kam šlápnout} \dicIndirectTranslationCS{(nevědět, pro kterou možnost se rozhodnout)}
\dicEntry[fr.] \dicTerm{fr.} \dicPos{zkr} \dicPhraseIS{franska, franskur} \dicDirectTranslationCS{francouzština, francouzský}
\dicEntry[fragt] \dicTerm{fragt} \dicIPA{{f}{r}{a}{x}{\textsubring{d}}} \dicPos{f}[4] \dicFlx{(‑ar)}[3] \dicLink{frakt}
\dicEntry[fragtskip] \dicTerm{fragt··skip} \dicIPA{{f}{r}{a}{x}{\textsubring{d}}{s}{\r{\textObardotlessj}}{\textsci}{\textsubring{b}}} \dicPos{n}[2] \dicFlx{(‑s, ‑)}[5] \dicLink{fraktskip}
\dicEntry[frakki] \dicTerm{frakk|i\smash{\textsuperscript{1}}} \dicsymFrequent\  \dicIPA{{f}{r}{a}{h}{\r{\textObardotlessj}}{\textsci}} \dicPos{m}[1] \dicFlx{(‑a, ‑ar)}[8] \dicDirectTranslationCS{(svrchní) kabát} \dicExampleIS{Hann var í nýjum frakka.} \dicExampleCS{Byl v~novém kabátě.}
\dicEntry[Frakki] \dicTerm{Frakk|i\smash{\textsuperscript{2}}} \dicIPA{{f}{r}{a}{h}{\r{\textObardotlessj}}{\textsci}} \dicPos{m}[1] \dicFlx{(‑a, ‑ar)}[8] \dicDirectTranslationCS{Francouz(ka)}
\dicEntry[Frakkland] \dicTerm{Frakk··land} \dicsymFrequent\  \dicIPA{{f}{r}{a}{h}{\r{g}}{l}{a}{n}{\textsubring{d}}} \dicPos{n}[2] \dicFlx{(‑s)}[4] \dicFieldCat{geog.} \dicDirectTranslationCS{Francie} \dicExampleIS{á landamærum Frakklands og Spánar} \dicExampleCS{na hranicích Francie a~Španělska}
\dicEntry[frakkur] \dicTerm{frakkur} \dicIPA{{f}{r}{a}{h}{\r{g}}{\textscy}{\textsubring{r}}} \dicPos{adj}[1] \dicFlx{(f frökk)}[2] \dicSynonym{frekur} \dicDirectTranslationCS{troufalý, smělý, drzý}
\dicEntry[frakt] \dicTerm{frakt}\dicTerm{, fragt} \dicIPA{{f}\-{r}\-{a}\-{x}\-{\textsubring{d}}\-} \dicPos{f}[4] \dicFlx{(‑ar)}[3] \textbf{1.} \dicSynonym{farmgjald} \dicDirectTranslationCS{přepravné}  \textbf{2.} \dicSynonym*{vöruflutningur} \dicDirectTranslationCS{náklad, přepravované zboží}
\dicEntry[fraktskip] \dicTerm{frakt··skip}\dicTerm{, fragtskip} \dicIPA{{f}\-{r}\-{a}\-{x}\-{\textsubring{d}}\-{s}\-{\r{\textObardotlessj}}\-{\textsci}\-{\textsubring{b}}\-} \dicPos{n}[2] \dicFlx{(‑s, ‑)}[5] \dicDirectTranslationCS{nákladní loď}
\dicEntry[fram] \dicTerm{fram} \dicsymFrequent\  \dicIPA{{f}{r}{a}{m}{\textlengthmark}} \dicPos{prep\,/\addthin adv} \dicSynonym{áfram} \dicDirectTranslationCS{vpřed, dopředu, kupředu} \dicExampleIS{stökkva fram} \dicExampleCS{skočit dopředu};  \dicIdiom{fram}[að]{ \dicPhraseIS{fram að e‑u}} \dicFlx{prep} \dicDirectTranslationCS{až do (čeho)} \dicExampleIS{fram að jólum} \dicExampleCS{až do Vánoc};  \dicPhraseIS{fram að þessu} \dicFlx{adv} \dicDirectTranslationCS{dosud, doposud, doposavad};  \dicIdiom{fram}[af]{ \dicPhraseIS{fram af e‑u}} \dicFlx{prep} {\textbf{a.}} \dicDirectTranslationCS{ven z~(čeho)};  {\textbf{b.}} \dicDirectTranslationCS{vpředu (čeho), v~přední části (čeho)};  \dicIdiom{fram}[á]{ \dicPhraseIS{fram á e‑ð}} \dicFlx{prep} \dicDirectTranslationCS{až do (čeho)} \dicIndirectTranslationCS{(časově)} \dicExampleIS{fram á morgun} \dicExampleCS{až do rána};  \dicIdiom{fram}[á við]{ \dicPhraseIS{fram á við}} \dicFlx{adv} \dicDirectTranslationCS{dopředu};  \dicIdiom{fram}[eftir]{ \dicPhraseIS{fram eftir}} \dicFlx{adv} \dicDirectTranslationCS{dlouho (spát, ponocovat ap.)} \dicExampleIS{vaka lengi fram eftir} \dicExampleCS{být dlouho vzhůru}; { \dicPhraseIS{fram eftir e‑u}} \dicFlx{prep} {\textbf{a.}} \dicDirectTranslationCS{podél (čeho)};  {\textbf{b.}} \dicDirectTranslationCS{dlouho do (čeho) (noci ap.)};  \dicIdiom{fram}[fyrir]{ \dicPhraseIS{fram fyrir}} \dicFlx{adv} \dicDirectTranslationCS{dopředu} \dicExampleIS{troða sér fram fyrir} \dicExampleCS{tlačit se dopředu}; { \dicPhraseIS{fram fyrir e‑n}} \dicFlx{prep} \dicDirectTranslationCS{(dopředu) před (koho)} \dicExampleIS{troða sér fram fyrir e‑n} \dicExampleCS{cpát se před (koho)};  \dicIdiom{fram}[hjá]{ \dicPhraseIS{fram hjá e‑u}} \dicFlx{prep} \dicDirectTranslationCS{kolem (čeho), okolo (čeho)} \dicExampleIS{Við fórum fram hjá mörgum skógum.} \dicExampleCS{Projížděli jsme kolem mnoha lesů.}; { \dicPhraseIS{fram hjá}} \dicFlx{adv} \dicDirectTranslationCS{kolem, podél};  \dicIdiom{fram}[í]{ \dicPhraseIS{fram í}} \dicFlx{adv} \dicDirectTranslationCS{vepředu, zepředu} \dicExampleIS{sitja fram í} \dicExampleCS{sedět vepředu};  \dicIdiom{fram}[undan]{ \dicPhraseIS{fram undan}} \dicFlx{adv} \dicDirectTranslationCS{před} \dicIndirectTranslationCS{(v~budoucnu)} \dicExampleIS{mikil vinna fram undan hjá okkur} \dicExampleCS{hodně práce před námi};  \dicPhraseIS{fram undan e‑u} \dicFlx{prep} \dicDirectTranslationCS{ven z~(čeho)};  \dicIdiom{fram}[undir]{ \dicPhraseIS{fram undir e‑ð}} \dicFlx{prep} \dicDirectTranslationCS{až do (čeho)} \dicIndirectTranslationCS{(časově)} \dicExampleIS{fram undir jól} \dicExampleCS{až do Vánoc};  \dicPhraseIS{fram undir þetta} \dicFlx{adv} \dicDirectTranslationCS{(téměř) až do teď};  \dicIdiom{fram}[úr]{ \dicPhraseIS{fram úr e‑u}} \dicFlx{prep} {\textbf{a.}} \dicDirectTranslationCS{(ven) z~(čeho)} \dicExampleIS{detta fram úr rúminu} \dicExampleCS{spadnout z~postele};  {\textbf{b.}} \dicDirectTranslationCS{kolem (čeho), podél (čeho)};  \dicIdiom{fram}[yfir]{ \dicPhraseIS{fram yfir e‑ð}} \dicFlx{prep} {\textbf{a.}} \dicDirectTranslationCS{až do (čeho)} \dicIndirectTranslationCS{(časově)} \dicExampleIS{fram yfir páska} \dicExampleCS{až do Velikonoc};  {\textbf{b.}} \dicDirectTranslationCS{před (čím), nad (co) (rámec ap.)} \dicIndirectTranslationCS{(o~preferencích)}
\dicEntry[framan] \dicTerm{framan} \dicsymFrequent\  \dicIPA{{f}{r}{a}{\textlengthmark}{m}{a}{\textsubring{n}}} \dicPos{prep\,/\addthin adv} \dicDirectTranslationCS{vepředu, zepředu} \dicIndirectTranslationCS{(o~směru zepředu k~něčemu)} \dicExampleIS{koma framan að e‑m} \dicExampleCS{přijít ke (komu) zepředu};  \dicIdiom{að}[framan]{ \dicPhraseIS{að framan}} \dicFlx{adv} {\textbf{a.}} \dicDirectTranslationCS{zepředu, z~přední strany};  {\textbf{b.}} \dicDirectTranslationCS{výše (v~textu ap.), nahoře} \dicExampleIS{umfjöllun hér að framan} \dicExampleCS{debata (diskutovaná) výše};  \dicIdiom{framan}[að]{ \dicPhraseIS{framan að e‑u\,/\addthin e‑m}} \dicFlx{prep} \dicDirectTranslationCS{(zepředu) k(e) (čemu\,/\addthin komu)};  \dicIdiom{framan}[af]{ \dicPhraseIS{framan af e‑u}} \dicFlx{prep} \dicDirectTranslationCS{ze začátku (čeho), z~počátku (čeho)}; { \dicPhraseIS{framan af}} \dicFlx{adv} \dicDirectTranslationCS{na počátku, ze začátku};  \dicIdiom{framan}[á]{ \dicPhraseIS{framan á e‑ð}} \dicFlx{prep} \dicDirectTranslationCS{na přední část (čeho), na předek (čeho)} \dicIndirectTranslationCS{(o~pohybu)}; { \dicPhraseIS{framan á e‑u}} \dicFlx{prep} \dicDirectTranslationCS{vpředu (čeho)} \dicIndirectTranslationCS{(o~pozici)}; { \dicPhraseIS{framan á}} \dicFlx{adv} \dicDirectTranslationCS{vpředu};  \dicIdiom{framan}[frá]{ \dicPhraseIS{framan frá}} \dicFlx{adv} \dicDirectTranslationCS{zepředu, z~přední strany};  \dicIdiom{framan}[í]{ \dicPhraseIS{framan í e‑ð\,/\addthin e‑u}} \dicFlx{prep} \dicDirectTranslationCS{v~přední části (čeho)};  \dicIdiom{framan}[til]{ \dicPhraseIS{framan til}} \dicFlx{adv} \dicDirectTranslationCS{vepředu};  \dicIdiom{framan}[úr]{ \dicPhraseIS{framan úr e‑u}} \dicFlx{prep} \dicDirectTranslationCS{z~(čeho)} \dicExampleIS{hróp framan úr baðstofunni} \dicExampleCS{výkřik z~koupelny};  \dicIdiom{framan}[við]{ \dicPhraseIS{framan við e‑ð}} \dicFlx{prep} \dicDirectTranslationCS{před (čím)} \dicExampleIS{garður framan við húsið} \dicExampleCS{park před domem};  \dicIdiom{í}[framan]{ \dicPhraseIS{í framan}} \dicFlx{adv} \dicDirectTranslationCS{vpředu, vepředu (ve tváři)}
\dicEntry[framandi] \dicTerm{fram··andi} \dicsymFrequent\  \dicIPA{{f}{r}{a}{\textlengthmark}{m}{a}{n}{\textsubring{d}}{\textsci}} \dicPos{adj}[13] \dicFlx{indecl}[1] \textbf{1.} \dicSynonym{erlendur} \dicDirectTranslationCS{cizí, exotický} \dicExampleIS{fara til framandi landa} \dicExampleCS{jezdit do exotických zemí}  \textbf{2.} \dicSynonym{annarlegur} \dicDirectTranslationCS{podivný, zvláštní, tajuplný}
\dicEntry[framandlegur] \dicTerm{framand··legur} \dicIPA{{f}{r}{a}{\textlengthmark}{m}{a}{n}{\textsubring{d}}{l}{\textepsilon}{\textbabygamma}{\textscy}{\textsubring{r}}} \dicPos{adj}[1]\dicFlx{}[-8] \dicDirectTranslationCS{cizokrajný, exotický}
\dicEntry[framandleiki] \dicTerm{framand··leik|i} \dicIPA{{f}{r}{a}{\textlengthmark}{m}{a}{n}{\textsubring{d}}{l}{ei}{\r{\textObardotlessj}}{\textsci}} \dicPos{m}[1] \dicFlx{(‑a)}[3] \dicDirectTranslationCS{cizokrajnost, exotičnost}
\dicEntry[framangreindur] \dicTerm{framan··greindur} \dicIPA{{f}{r}{a}{\textlengthmark}{m}{a}{n}{\r{g}}{r}{ei}{n}{\textsubring{d}}{\textscy}{\textsubring{r}}} \dicPos{adj}[2]\dicFlx{}[-17] \dicDirectTranslationCS{dříve zmíněný, předchozí,  předcházející}
\dicEntry[framanverður] \dicTerm{framan··verður} \dicIPA{{f}{r}{a}{\textlengthmark}{m}{a}{n}{v}{\textepsilon}{r}{ð}{\textscy}{\textsubring{r}}} \dicPos{adj}[2]\dicFlx{}[-4] \dicDirectTranslationCS{přední} \dicIndirectTranslationCS{(o~poloze)} \dicExampleIS{í framanverðum dalnum} \dicExampleCS{vpředu v~údolí};  \dicPhraseIS{að framanverðu} \dicFlx{adv} \dicDirectTranslationCS{zepředu, z~přední strany}
\dicEntry[framar] \dicTerm{framar} \dicsymFrequent\  \dicIPA{{f}{r}{a}{\textlengthmark}{m}{a}{\textsubring{r}}} \dicPos{adv} \dicFlx{comp (pos frammi, sup fremst)} \textbf{1.} \dicDirectTranslationCS{více} \dicIndirectTranslationCS{(často se záporem)};  \dicPhraseIS{aldrei framar} \dicFlx{adv} \dicDirectTranslationCS{nikdy více};  \dicPhraseIS{ekki framar} \dicFlx{adv} \dicDirectTranslationCS{ne více, ne déle}  \textbf{2.} \dicPhraseIS{framar öðru} \dicFlx{adv} \dicDirectTranslationCS{především};  \dicPhraseIS{framar öllu} \dicFlx{adv} \dicDirectTranslationCS{ze všeho nejvíce; víc než cokoliv jiného}  \textbf{3.} \dicSynonym*{á undan öðrum} \dicDirectTranslationCS{dále, více vepředu}
\dicEntry[framarlega] \dicTerm{framar··lega} \dicIPA{{f}{r}{a}{\textlengthmark}{m}{a}{r}{l}{\textepsilon}{\textbabygamma}{a}} \dicPos{adv} \textbf{1.} \dicDirectTranslationCS{vpředu, vepředu}  \textbf{2.} \dicDirectTranslationCS{hluboko, daleko uvnitř (údolí ap.)} \dicExampleIS{framarlega í dalnum} \dicExampleCS{daleko uvnitř údolí};  \dicPhraseIS{svo framarlega sem} \dicFlx{conj} \dicSynonym*{með því skilyrði að} \dicDirectTranslationCS{pokud; dokud; tak dlouho, dokud}
\begin{xtolerant}{}{1pt}
\dicEntry[framámaður] \dicTerm{framá··|maður}\dicTerm{, frammámaður}\addthinS\textls[-10]{\dicIPA{{f}\-{r}\-{a}\-{m}\-{au}\-{m}\-{a}\-{ð}\-{\textscy}\-{\textsubring{r}}}} \dicPos{m}[13]\addthinS\dicFlx{(‑manns, ‑menn)}[2] \dicSynonym*{forystumaður} \dicDirectTranslationCS{vlivná osoba, prominentní člověk}
\end{xtolerant}
\dicEntry[frambjóðandi] \dicTerm{fram··bjóð·|andi} \dicIPA{{f}{r}{a}{m}{\textsubring{b}}{j}{ou}{ð}{a}{n}{\textsubring{d}}{\textsci}} \dicPos{m}[2] \dicFlx{(‑anda, ‑endur)}[1] \dicDirectTranslationCS{kandidát(ka), uchazeč(ka)} \dicExampleIS{frambjóðandi til þings} \dicExampleCS{kandidátka do parlamentu}
\dicEntry[framboð] \dicTerm{fram··boð} \dicsymFrequent\  \dicIPA{{f}{r}{a}{m}{\textsubring{b}}{\textopeno}{\texttheta}} \dicPos{n}[2] \dicFlx{(‑s, ‑)}[5] \textbf{1.} \dicDirectTranslationCS{kandidatura} \dicExampleIS{framboð til þings} \dicExampleCS{kandidatura do parlamentu}  \textbf{2.} \dicDirectTranslationCS{nabídka} \dicExampleIS{framboð og eftirspurn} \dicExampleCS{nabídka a~poptávka}
\dicEntry[framburðargalli] \dicTerm{fram·burðar··gall|i} \dicIPA{{f}{r}{a}{m}{\textsubring{b}}{\textscy}{r}{ð}{a}{\textsubring{r}}{\r{g}}{a}{\textsubring{d}}{l}{\textsci}} \dicPos{m}[1] \dicFlx{(‑a, ‑ar)}[8] \dicDirectTranslationCS{výslovnostní chyba, chyba ve výslovnosti}
\dicEntry[framburður] \dicTerm{fram··burð|ur} \dicIPA{{f}{r}{a}{m}{\textsubring{b}}{\textscy}{r}{ð}{\textscy}{\textsubring{r}}} \dicPos{m}[10] \dicFlx{(‑ar)}[7] \textbf{1.} \dicFieldCat{jaz.} \dicDirectTranslationCS{výslovnost} \dicExampleIS{framburður orða} \dicExampleCS{výslovnost slov}  \textbf{2.} \dicFieldCat{práv.} \dicDirectTranslationCS{svědectví, svědecká výpověď}  \textbf{3.} \dicDirectTranslationCS{usazenina, sedlina, sediment}
\dicEntry[frambúð] \dicTerm{fram··búð} \dicIPA{{f}{r}{a}{m}{\textsubring{b}}{u}{\texttheta}} \dicPos{f}[7] \dicFlx{(‑ar)}[3] \dicPhraseIS{til frambúðar} \dicFlx{adv} \dicDirectTranslationCS{trvale, natrvalo, nastálo} \dicExampleIS{gegna starfinu til frambúðar} \dicExampleCS{zastávat pozici natrvalo}
\dicEntry[frambærilegur] \dicTerm{fram·bæri··legur} \dicIPA{{f}{r}{a}{m}{\textsubring{b}}{a}{i}{r}{\textsci}{l}{\textepsilon}{\textbabygamma}{\textscy}{\textsubring{r}}} \dicPos{adj}[1]\dicFlx{}[-8] \dicSynonym{boðlegur} \dicDirectTranslationCS{reprezentativní, reprezentační}
\dicEntry[framdi] \dicTerm{framdi} \dicIPA{{f}{r}{a}{m}{\textsubring{d}}{\textsci}} \dicPos{v} \dicFlx{ind pf sg 1 pers} \dicLink{fremja}
\dicEntry[framdráttur] \dicTerm{fram··drátt|ur} \dicIPA{{f}{r}{a}{m}{\textsubring{d}}{r}{au}{h}{\textsubring{d}}{\textscy}{\textsubring{r}}} \dicPos{m}[12] \dicFlx{(‑ar)}[8] \dicPhraseIS{e‑að er e‑u til framdráttar} \dicDirectTranslationCS{(co) je (čemu) ku prospěchu, (co) je (čemu) k~dobru}
\dicEntry[framdrif] \dicTerm{fram··drif} \dicIPA{{f}{r}{a}{m}{\textsubring{d}}{r}{\textsci}{f}} \dicPos{n}[2] \dicFlx{(‑s, ‑)}[5] \dicDirectTranslationCS{přední pohon, pohon předních kol, náhon na přední kola}
\dicEntry[framfara] \dicTerm{fram··fara-} \dicIPA{{f}{r}{a}{m}{f}{a}{r}{a}} \dicPos{predp} \dicDirectTranslationCS{pokrokový} \dicExampleIS{framfaraspor} \dicExampleCS{pokrok}
\dicEntry[framfleyta] \dicTerm{fram··fleyt|a} \dicIPA{{f}{r}{a}{m}{f}{l}{ei}{\textsubring{d}}{a}} \dicPos{v}[2] \dicFlx{(‑ti, ‑t)}[59] \dicFlx{dat} \dicDirectTranslationCS{(u)živit, vydržovat} \dicExampleIS{framfleyta fjölskyldu sinni} \dicExampleCS{živit svou rodinu}
\dicEntry[framfylgja] \dicTerm{fram··fylg|ja} \dicIPA{{f}{r}{a}{m}{f}{\textsci}{l}{\r{\textObardotlessj}}{a}} \dicPos{v}[2] \dicFlx{(‑di, ‑t)}[134] \dicFlx{dat} \dicDirectTranslationCS{následovat, provést, provádět} \dicExampleIS{framfylgja reglum} \dicExampleCS{následovat pravidla}
\dicEntry[framfæra] \dicTerm{fram··fær|a} \dicIPA{{f}{r}{a}{m}{f}{a}{i}{r}{a}} \dicPos{v}[2] \dicFlx{(‑ði, ‑t)}[99] \dicSynonym{framfleyta} \dicDirectTranslationCS{živit, vydržovat}
\dicEntry[framfærandi] \dicTerm{fram··fær·|andi} \dicIPA{{f}{r}{a}{m}{f}{a}{i}{r}{a}{n}{\textsubring{d}}{\textsci}} \dicPos{m}[2] \dicFlx{(‑anda, ‑endur)}[1] \dicLangCat{form.} \dicDirectTranslationCS{živitel(ka), chlebodárce, chlebodárkyně} \dicExampleIS{framfærandi barns} \dicExampleCS{živitelka dítěte}
\dicEntry[framfæri] \dicTerm{fram··færi} \dicsymFrequent\  \dicIPA{{f}{r}{a}{m}{f}{a}{i}{r}{\textsci}} \dicPos{n}[2] \dicFlx{(‑s)}[20] \textbf{1.} \dicSynonym{afkoma} \dicDirectTranslationCS{(hmotná) podpora, podporování, vydržování};  \dicPhraseIS{hafa e‑n á framfæri} \dicDirectTranslationCS{(hmotně) podporovat (koho)};  \dicPhraseIS{vera á framfæri e‑rs} \dicDirectTranslationCS{být (hmotně) podporován (kým)}  \textbf{2.} \dicPhraseIS{koma sér á framfæri} \dicDirectTranslationCS{prosadit se, sklidit úspěch};  \dicPhraseIS{koma e‑u á framfæri} \dicDirectTranslationCS{prosadit (co), realizovat (co)}
\dicEntry[framfærsla] \dicTerm{fram··færsl|a} \dicIPA{{f}{r}{a}{m}{f}{a}{i}{\textsubring{r}}{s}{\textsubring{d}}{l}{a}} \dicPos{f}[1] \dicFlx{(‑u)}[5] \dicDirectTranslationCS{(hmotná) podpora, živobytí}
\begin{xtolerant}{}{1pt}
\dicEntry[framfærslukostnaður] \dicTerm{fram·færslu··kost·nað|ur} \textls[-10]{\dicIPA{{f}\-{r}\-{a}\-{m}\-{f}\-{a}\-{i}\-{\textsubring{r}}\-{s}\-{\textsubring{d}}\-{l}\-{\textscy}\-{k\smash{\textsuperscript{h}}}\-{\textopeno}\-{s}\-{\textsubring{d}}\-{n}\-{a}\-{ð}\-{\textscy}\-{\textsubring{r}}}} \dicPos{m}[10] \dicFlx{(‑ar)}[9] \dicDirectTranslationCS{životní náklady, náklady na živobytí}
\end{xtolerant}
\dicEntry[framfærsluvísitala] \dicTerm{fram·færslu··vísi·|tala} \textls[-10]{\dicIPA{{f}\-{r}\-{a}\-{m}\-{f}\-{a}\-{i}\-{\textsubring{r}}\-{s}\-{\textsubring{d}}\-{l}\-{\textscy}\-{v}\-{i}\-{s}\-{\textsci}\-{t\smash{\textsuperscript{h}}}\-{a}\-{l}\-{a}}} \dicPos{f}[1] \dicFlx{(‑tö\-lu, ‑töl\-ur)}[14] \dicFieldCat{ekon.} \dicDirectTranslationCS{index životní úrovně}
\dicEntry[framför] \dicTerm{fram··|för} \dicsymFrequent\  \dicIPA{{f}{r}{a}{m}{f}{\oe}{\textsubring{r}}} \dicPos{f}[7] \dicFlx{(‑farar, ‑farir)}[16] \textbf{1.} \dicSynonym{þroski} \dicDirectTranslationCS{vzestup, rozmach, rozvoj}  \textbf{2.} \dicPhraseIS{framfarir} \dicFlx{pl} \dicDirectTranslationCS{pokrok, změna k~lepšímu};  \dicPhraseIS{taka framförum} \dicDirectTranslationCS{udělat pokrok, dělat pokroky}
\dicEntry[framgangur] \dicTerm{fram··gang|ur} \dicIPA{{f}{r}{a}{m}{\r{g}}{au}{\ng}{\r{g}}{\textscy}{\textsubring{r}}} \dicPos{m}[6] \dicFlx{(‑s, ‑ar)}[5] \textbf{1.} \dicDirectTranslationCS{rozvoj, postup, posun, vývoj}  \textbf{2.} \dicSynonym{stöðuhækkun} \dicDirectTranslationCS{postup} \dicIndirectTranslationCS{(přechod do vyššího stupně)}
\dicEntry[framgjarn] \dicTerm{fram··|gjarn} \dicIPA{{f}{r}{a}{m}{\r{\textObardotlessj}}{a}{r}{\textsubring{d}}{\textsubring{n}}} \dicPos{adj}[5] \dicFlx{(f ‑gjörn)}[6] \textbf{1.} \dicSynonym{metnaðargjarn} \dicDirectTranslationCS{ctižádostivý, ambiciózní}  \textbf{2.} \dicSynonym{frakkur} \dicDirectTranslationCS{troufalý, drzý}
\dicEntry[framgómhljóð] \dicTerm{fram·góm··hljóð} \dicIPA{{f}{r}{a}{m}{\r{g}}{ou}{m}{\textsubring{l}}{j}{ou}{\texttheta}} \dicPos{n}[2] \dicFlx{(‑s, ‑)}[5] \dicFieldCat{jaz.} \dicDirectTranslationCS{předopatrová souhláska, palatála}
\dicEntry[framgómmæltur] \dicTerm{fram·góm··mæltur} \dicIPA{{f}{r}{a}{m}{\r{g}}{ou}{m}{m}{a}{i}{\textsubring{l}}{\textsubring{d}}{\textscy}{\textsubring{r}}} \dicPos{adj}[1]\dicFlx{}[-13] \dicFieldCat{jaz.} \dicDirectTranslationCS{předopatrový, palatální}
\dicEntry[framgómur] \dicTerm{fram··góm|ur} \dicIPA{{f}{r}{a}{m}{\r{g}}{ou}{m}{\textscy}{\textsubring{r}}} \dicPos{m}[6] \dicFlx{(‑s, ‑ar)}[10] \dicFieldCat{anat.} \dicDirectTranslationCS{patro}
\dicEntry[framhald] \dicTerm{fram··hald} \dicsymFrequent\  \dicIPA{{f}{r}{a}{m}{h}{a}{l}{\textsubring{d}}} \dicPos{n}[2] \dicFlx{(‑s)}[2] \dicDirectTranslationCS{pokračování};  \dicPhraseIS{framhald á e‑u} \dicDirectTranslationCS{pokračování (čeho)};  \dicPhraseIS{í framhaldi af e‑u} \dicDirectTranslationCS{jako pokračování (čeho), v~návaznosti na (co)}
\dicEntry[framhaldsflug] \dicTerm{fram·halds··flug} \dicIPA{{f}{r}{a}{m}{h}{a}{l}{\textsubring{d}}{s}{f}{l}{\textscy}{x}} \dicPos{n}[2] \dicFlx{(‑s, ‑)}[5] \dicDirectTranslationCS{navazující\,/\addthin přípojný let}
\dicEntry[framhaldsnám] \dicTerm{fram·halds··nám} \dicIPA{{f}{r}{a}{m}{h}{a}{l}{\textsubring{d}}{s}{n}{au}{\textsubring{m}}} \dicPos{n}[2] \dicFlx{(‑s)}[2] \dicDirectTranslationCS{doktorské\,/\addthin postgraduální studium} \dicExampleIS{stunda framhaldsnám erlendis} \dicExampleCS{být na doktorském studiu v~zahraničí}
\dicEntry[framhaldsnámskeið] \dicTerm{fram·halds··nám·skeið} \dicIPA{{f}\-{r}\-{a}\-{m}\-{h}\-{a}\-{l}\-{\textsubring{d}}\-{s}\-{n}\-{au}\-{m}\-{s}\-{\r{\textObardotlessj}}\-{ei}\-{\texttheta}\-} \dicPos{n}[2] \dicFlx{(‑s, ‑)}[5] \dicFieldCat{škol.} \dicDirectTranslationCS{pokračovací kurz}
\dicEntry[framhaldsnemi] \dicTerm{fram·halds··nem|i} \dicIPA{{f}{r}{a}{m}{h}{a}{l}{\textsubring{d}}{s}{n}{\textepsilon}{m}{\textsci}} \dicPos{m}[1] \dicFlx{(‑a, ‑ar)}[1] \dicDirectTranslationCS{doktorand(ka)}
\dicEntry[framhaldssaga] \dicTerm{fram·halds··|saga} \dicIPA{{f}{r}{a}{m}{h}{a}{l}{\textsubring{d}}{s}{a}{\textbabygamma}{a}} \dicPos{f}[1] \dicFlx{(‑sögu, ‑sögur)}[14] \dicDirectTranslationCS{seriál, příběh na pokračování}
\dicEntry[framhaldsskólastig] \dicTerm{fram·halds··skóla·stig} \dicIPA{{f}\-{r}\-{a}\-{m}\-{h}\-{a}\-{l}\-{\textsubring{d}}\-{s}\-{\r{g}}\-{ou}\-{l}\-{a}\-{s}\-{\textsubring{d}}\-{\textsci}\-{x}\-} \dicPos{n}[2] \dicFlx{(‑s, ‑)}[5] \dicFieldCat{škol.} \dicDirectTranslationCS{středoškolský stupeň}
\begin{xtolerant}{}{1pt}
\dicEntry[framhaldsskóli] \dicTerm{fram·halds··skól|i}\addthinS\dicIPA{{f}{r}{a}{m}{h}{a}{l}{\textsubring{d}}{s}{\r{g}}{ou}{l}{\textsci}}\addthinS\dicPos{m}[1]\addthinS\dicFlx{(‑a, ‑ar)}[1]\addthinS\dicFieldCat{škol.}\addthinS\dicDirectTranslationCS{střední škola}
\end{xtolerant}
\dicEntry[framhandleggur] \dicTerm{fram··hand·legg|ur} \dicIPA{{f}{r}{a}{m}{h}{a}{n}{\textsubring{d}}{l}{\textepsilon}{\r{g}}{\textscy}{\textsubring{r}}} \dicPos{m}[9] \dicFlx{(‑jar\,/\addthin ‑s, ‑ir)}[26] \dicFieldCat{anat.} \dicDirectTranslationCS{předloktí}
\dicEntry[framherji] \dicTerm{fram··herj|i} \dicIPA{{f}{r}{a}{m}{h}{\textepsilon}{r}{j}{\textsci}} \dicPos{m}[1] \dicFlx{(‑a, ‑ar)}[1] \dicFieldCat{sport.} \dicDirectTranslationCS{(hrotový) útočník}
\dicEntry[framhjá] \dicTerm{fram··hjá} \dicsymFrequent\  \dicIPA{{f}{r}{a}{m}{\c{c}}{au}} \dicPos{prep\,/\addthin adv} \dicFlx{dat} \dicFlx{adv} \dicDirectTranslationCS{kolem, okolo} \dicExampleIS{ganga framhjá} \dicExampleCS{jít kolem};  \dicPhraseIS{framhjá e‑u} \dicFlx{prep} \dicDirectTranslationCS{kolem (čeho), podél (čeho)}
\dicEntry[framhjáhald] \dicTerm{fram·hjá··hald} \dicIPA{{f}{r}{a}{m}{\c{c}}{au}{h}{a}{l}{\textsubring{d}}} \dicPos{n}[2] \dicFlx{(‑s)}[2] \dicDirectTranslationCS{nevěra, cizoložství}
\dicEntry[framhjól] \dicTerm{fram··hjól} \dicIPA{{f}{r}{a}{m}{\c{c}}{ou}{\textsubring{l}}} \dicPos{n}[2] \dicFlx{(‑s, ‑)}[5] \dicDirectTranslationCS{přední kolo}
\dicEntry[framhjóladrif] \dicTerm{fram··hjóla·drif} \dicIPA{{f}{r}{a}{m}{\c{c}}{ou}{l}{a}{\textsubring{d}}{r}{\textsci}{f}} \dicPos{n}[2] \dicFlx{(‑s, ‑)}[5] \dicDirectTranslationCS{přední pohon, pohon předních kol, náhon na přední kola} \dicAntonym{afturhjóladrif}
\dicEntry[framhleypinn] \dicTerm{fram··hleypinn} \dicIPA{{f}{r}{a}{m}{\textsubring{l}}{ei}{\textsubring{b}}{\textsci}{\textsubring{n}}} \dicPos{adj}[6]\dicFlx{}[-2] \dicDirectTranslationCS{opovážlivý, drzý}
\dicEntry[framhleypni] \dicTerm{fram··hleypn|i} \dicIPA{{f}{r}{a}{m}{\textsubring{l}}{ei}{h}{\textsubring{b}}{n}{\textsci}} \dicPos{f}[3] \dicFlx{(‑i)}[3] \dicDirectTranslationCS{opovážlivost, drzost}
\dicEntry[framhlið] \dicTerm{fram··hlið} \dicIPA{{f}{r}{a}{m}{\textsubring{l}}{\textsci}{\texttheta}} \dicPos{f}[4] \dicFlx{(‑ar, ‑ar)}[1] \textbf{1.} \dicDirectTranslationCS{přední strana, předek}  \textbf{2.} \dicDirectTranslationCS{fasáda, průčelí (domu ap.)}
\dicEntry[frami] \dicTerm{fram|i} \dicIPA{{f}{r}{a}{\textlengthmark}{m}{\textsci}} \dicPos{m}[1] \dicFlx{(‑a)}[3] \dicSynonym{frægð} \dicDirectTranslationCS{věhlas, proslulost, renomé}
\dicEntry[framið] \dicTerm{framið} \dicIPA{{f}{r}{a}{\textlengthmark}{m}{\textsci}{\texttheta}} \dicPos{v} \dicFlx{supin} \dicLink{fremja}
\dicEntry[framjaxl] \dicTerm{fram··jaxl} \dicIPA{{f}{r}{a}{m}{j}{a}{x}{s}{\textsubring{d}}{\textsubring{l}}} \dicPos{m}[4] \dicFlx{(‑s, ‑ar)}[4] \dicFieldCat{anat.} \dicDirectTranslationCS{třenový zub}
\dicEntry[framkalla] \dicTerm{fram··kall|a} \dicIPA{{f}{r}{a}{m}{k\smash{\textsuperscript{h}}}{a}{\textsubring{d}}{l}{a}} \dicPos{v}[1] \dicFlx{(‑aði)}[13] \dicFlx{acc} \textbf{1.} \dicSynonym{orsaka} \dicDirectTranslationCS{vyvolat, způsobit, zapříčinit, přivodit} \dicExampleIS{framkalla gleði} \dicExampleCS{vyvolat radost}  \textbf{2.} \dicDirectTranslationCS{vyvolat (film ap.)} \dicExampleIS{framkalla filmu} \dicExampleCS{vyvolat film}
\dicEntry[framkoma] \dicTerm{fram··kom|a} \dicsymFrequent\  \dicIPA{{f}{r}{a}{m}{k\smash{\textsuperscript{h}}}{\textopeno}{m}{a}} \dicPos{f}[1] \dicFlx{(‑u)}[5] \textbf{1.} \dicSynonym{hegðun} \dicDirectTranslationCS{(za)chování, jednání} \dicExampleIS{vera hissa á framkomu hennar} \dicExampleCS{být překvapený jejím chováním}  \textbf{2.} \dicSynonym*{það að koma fram} \dicDirectTranslationCS{výstup, vystoupení} \dicExampleIS{fyrsta framkoma e‑rs á sviði} \dicExampleCS{první výstup (koho) na jevišti}
\dicEntry[framkrókar] \dicTerm{fram··krókar} \dicIPA{{f}{r}{a}{m}{k\smash{\textsuperscript{h}}}{r}{ou}{\r{g}}{a}{\textsubring{r}}} \dicPos{m}[6] \dicFlx{pl}[2] \dicPhraseIS{leggja sig í framkróka um\,/\addthin við að (gera e‑ð)} \dicDirectTranslationCS{vynasnažit se (udělat (co)), snažit se ze všech sil (udělat (co))}
\dicEntry[framkvæma] \dicTerm{fram··kvæm|a} \dicsymFrequent\  \dicIPA{{f}{r}{a}{m}{k\smash{\textsuperscript{h}}}{v}{a}{i}{m}{a}} \dicPos{v}[2] \dicFlx{(‑di, ‑t)}[141] \dicFlx{acc} \dicDirectTranslationCS{provést, provádět, uskutečnit, vykonat, realizovat} \dicExampleIS{framkvæma aðgerð} \dicExampleCS{provést operaci}
\dicEntry[framkvæmanlegur] \dicTerm{fram·kvæman··legur} \dicIPA{{f}{r}{a}{m}{k\smash{\textsuperscript{h}}}{v}{a}{i}{m}{a}{n}{l}{\textepsilon}{\textbabygamma}{\textscy}{\textsubring{r}}} \dicPos{adj}[1]\dicFlx{}[-8] \dicSynonym{gerlegur} \dicDirectTranslationCS{proveditelný, uskutečnitelný, realizovatelný} \dicAntonym{óframkvæmanlegur}
\dicEntry[framkvæmd] \textls[15]{\dicTerm{fram··kvæmd} \dicsymFrequent\  \dicIPA{{f}{r}{a}{m}{k\smash{\textsuperscript{h}}}{v}{a}{i}{m}{\textsubring{d}}} \dicPos{f}[7] \dicFlx{(‑ar, ‑ir)}[1] \textbf{1.} \dicDirectTranslationCS{realizace, provedení, uskutečnění};  \dicPhraseIS{koma e‑u í framkvæmd} \dicSynonym{framkvæma} \dicDirectTranslationCS{uskutečnit (co), realizovat (co)}  \textbf{2.} \dicSynonym{verk} \dicDirectTranslationCS{renovace, stavební práce}}
\dicEntry[framkvæmdastjóri] \dicTerm{fram·kvæmda··stjór|i} \dicsymFrequent\  \dicIPA{{f}{r}{a}{m}{k\smash{\textsuperscript{h}}}{v}{a}{i}{m}{\textsubring{d}}{a}{s}{\textsubring{d}}{j}{ou}{r}{\textsci}} \dicPos{m}[1] \dicFlx{(‑a, ‑ar)}[1] \dicDirectTranslationCS{výkonný ředitel, výkonná ředitelka} \dicExampleIS{styrkja stöðu framkvæmdastjórans} \dicExampleCS{upevnit pozici výkonné ředitelky}
\dicEntry[framkvæmdastjórn] \dicTerm{fram·kvæmda··stjórn} \dicIPA{{f}{r}{a}{m}{k\smash{\textsuperscript{h}}}{v}{a}{i}{m}{\textsubring{d}}{a}{s}{\textsubring{d}}{j}{ou}{r}{\textsubring{d}}{\textsubring{n}}} \dicPos{f}[7] \dicFlx{(‑ar, ‑ir)}[1] \dicDirectTranslationCS{výkonná rada, exekutiva}
\dicEntry[framköllun] \dicTerm{fram··köll|un} \dicIPA{{f}{r}{a}{m}{k\smash{\textsuperscript{h}}}{\oe}{\textsubring{d}}{l}{\textscy}{\textsubring{n}}} \dicPos{f}[7] \dicFlx{(‑unar)}[12] \dicDirectTranslationCS{vyvolání (filmu ap.)}
\dicEntry[framlag] \dicTerm{fram··|lag} \dicsymFrequent\  \dicIPA{{f}{r}{a}{m}{l}{a}{x}} \dicPos{n}[2] \dicFlx{(‑lags, ‑lög)}[8] \dicDirectTranslationCS{(finanční) příspěvek, dotace, grant} \dicExampleIS{framlag til verksins} \dicExampleCS{příspěvek na práci}
\dicEntry[framleiða] \dicTerm{fram··lei|ða} \dicsymFrequent\  \dicIPA{{f}{r}{a}{m}{l}{ei}{ð}{a}} \dicPos{v}[2] \dicFlx{(‑ddi, ‑tt)}[167] \dicFlx{acc} \textbf{1.} \dicDirectTranslationCS{(vy)produkovat, vyrobit, vyrábět} \dicExampleIS{framleiða vöru úr hráefninu} \dicExampleCS{vyrábět zboží ze suroviny}  \textbf{2.} \dicDirectTranslationCS{produkovat (film ap.)}
\dicEntry[framleiðandi] \dicTerm{fram··leið·|andi} \dicIPA{{f}{r}{a}{m}{l}{ei}{ð}{a}{n}{\textsubring{d}}{\textsci}} \dicPos{m}[2] \dicFlx{(‑anda, ‑endur)}[1] \textbf{1.} \dicSynonym*{vöruframleiðandi} \dicDirectTranslationCS{výrobce, producent} \dicExampleIS{framleiðandi vörunnar} \dicExampleCS{výrobce zboží}  \textbf{2.} \dicSynonym*{kvikmyndaframleiðandi} \dicDirectTranslationCS{producent(ka) (filmu ap.)}
\dicEntry[framleiðni] \dicTerm{fram··leiðn|i} \dicIPA{{f}{r}{a}{m}{l}{ei}{ð}{n}{\textsci}} \dicPos{f}[3] \dicFlx{(‑i)}[3] \dicDirectTranslationCS{produktivita}
\dicEntry[framleiðsla] \dicTerm{fram··leiðsl|a} \dicsymFrequent\  \dicIPA{{f}{r}{a}{m}{l}{ei}{ð}{s}{\textsubring{d}}{l}{a}} \dicPos{f}[1] \dicFlx{(‑u)}[5] \textbf{1.} \dicDirectTranslationCS{výroba, vyrábění, produkce} \dicExampleIS{Framleiðslan er í fullum gangi.} \dicExampleCS{Výroba je v~plném proudu.};  \dicPhraseIS{framleiðsla á e‑u} \dicDirectTranslationCS{produkce (čeho)}  \textbf{2.} \dicDirectTranslationCS{produkt, výrobek}
\dicEntry[framleiðslugeta] \dicTerm{fram·leiðslu··get|a} \dicIPA{{f}{r}{a}{m}{l}{ei}{ð}{s}{\textsubring{d}}{l}{\textscy}{\r{\textObardotlessj}}{\textepsilon}{\textsubring{d}}{a}} \dicPos{f}[1] \dicFlx{(‑u)}[5] \dicDirectTranslationCS{výrobní kapacita}
\begin{xtolerant}{}{1pt}
\dicEntry[framleiðslukostnaður] \dicTerm{fram·leiðslu··kost·nað|ur} \textls[-10]{\dicIPA{{f}\-{r}\-{a}\-{m}\-{l}\-{ei}\-{ð}\-{s}\-{\textsubring{d}}\-{l}\-{\textscy}\-{k\smash{\textsuperscript{h}}}\-{\textopeno}\-{s}\-{\textsubring{d}}\-{n}\-{a}\-{ð}\-{\textscy}\-{\textsubring{r}}}} \dicPos{m}[10] \dicFlx{(‑ar)}[9] \dicDirectTranslationCS{výrobní náklady, náklady na výrobu}
\end{xtolerant}
\dicEntry[framleigja] \dicTerm{fram··leig|ja} \dicIPA{{f}{r}{a}{m}{l}{ei}{j}{a}} \dicPos{v}[2] \dicFlx{(‑ði, ‑t)}[92] \dicFlx{acc} \dicDirectTranslationCS{dále pronajmout}
\dicEntry[framlenging] \dicTerm{fram··leng·ing} \dicIPA{{f}{r}{a}{m}{l}{ei}{\textltailn}{\r{\textObardotlessj}}{i}{\ng}{\r{g}}} \dicPos{f}[4] \dicFlx{(‑ar, ‑ar)}[5] \textbf{1.} \dicDirectTranslationCS{prodloužení, prodlužování} \dicExampleIS{framlenging á leyfinu} \dicExampleCS{prodloužení povolení}  \textbf{2.} \dicFieldCat{sport.} \dicDirectTranslationCS{prodloužení}
\dicEntry[framlengja] \dicTerm{fram··leng|ja} \dicIPA{{f}{r}{a}{m}{l}{ei}{\textltailn}{\r{\textObardotlessj}}{a}} \dicPos{v}[2] \dicFlx{(‑di, ‑t)}[129] \dicFlx{acc} \textbf{1.} \dicDirectTranslationCS{prodloužit (lhůtu ap.), prodlužovat, prolongovat} \dicExampleIS{framlengja námskeiðið um eina viku} \dicExampleCS{prodloužit kurz o~jeden týden}  \textbf{2.} \dicDirectTranslationCS{prodloužit, protáhnout (stůl ap.)} \dicExampleIS{framlengja borðið um 10 sm} \dicExampleCS{prodloužit stůl o~10 cm}
\dicEntry[framliðinn] \dicTerm{fram··liðinn} \dicIPA{{f}{r}{a}{m}{l}{\textsci}{ð}{\textsci}{\textsubring{n}}} \dicPos{adj}[6]\dicFlx{}[-6] \dicSynonym{dáinn} \dicDirectTranslationCS{zesnulý, zemřelý}
\dicEntry[framljós] \dicTerm{fram··ljós} \dicIPA{{f}{r}{a}{m}{l}{j}{ou}{s}} \dicPos{n}[2] \dicFlx{(‑s, ‑)}[5] \dicDirectTranslationCS{přední světlo (u~auta ap.)}
\dicEntry[frammámaður] \textls[15]{\dicTerm{frammá··|maður} \dicIPA{{f}{r}{a}{m}{\textlengthmark}{au}{m}{a}{ð}{\textscy}{\textsubring{r}}} \dicPos{m}[13] \dicFlx{(‑manns, ‑menn)}[2]} \dicLink{framámaður}
\dicEntry[frammi] \dicTerm{frammi} \dicsymFrequent\  \dicIPA{{f}{r}{a}{m}{\textlengthmark}{\textsci}} \dicPos{prep\,/\addthin adv} \dicFlx{(comp framar, sup fremst)} \dicFlx{adv} \dicDirectTranslationCS{vpředu, vepředu} \dicExampleIS{bíddu frammi} \dicExampleCS{počkej vepředu};  \dicPhraseIS{frammi fyrir e‑u} \dicFlx{prep} \dicSynonym{andspænis} \dicDirectTranslationCS{naproti (čemu)}
\dicEntry[frammistaða] \dicTerm{frammi··|staða} \dicIPA{{f}{r}{a}{m}{\textlengthmark}{\textsci}{s}{\textsubring{d}}{a}{ð}{a}} \dicPos{f}[1] \dicFlx{(‑stöðu)}[2] \dicDirectTranslationCS{(předvedený) výkon, vykonání} \dicExampleIS{glæsileg frammistaða} \dicExampleCS{vynikající výkon}
\dicEntry[frammí] \dicTerm{framm··í} \dicIPA{{f}{r}{a}{m}{\textlengthmark}{i}} \dicPos{adv} \dicPhraseIS{fram í} \dicLink{fram}
\dicEntry[frammæltur] \dicTerm{fram··mæltur} \dicIPA{{f}{r}{a}{m}{\textlengthmark}{a}{i}{\textsubring{l}}{\textsubring{d}}{\textscy}{\textsubring{r}}} \dicPos{adj}[1]\dicFlx{}[-13] \dicFieldCat{jaz.} \dicDirectTranslationCS{přední (samo\-hláska ap.)} \dicAntonym{uppmæltur}
\dicEntry[framreiða] \dicTerm{fram··rei|ða} \dicIPA{{f}{r}{a}{m}{r}{ei}{ð}{a}} \dicPos{v}[2] \dicFlx{(‑ddi, ‑tt)}[167] \dicDirectTranslationCS{podávat (večeři ap.), servírovat} \dicExampleIS{framreiða kvöldverð} \dicExampleCS{podávat večeři}
\dicEntry[framreiðsla] \dicTerm{fram··reiðsl|a} \dicIPA{{f}{r}{a}{m}{r}{ei}{ð}{s}{\textsubring{d}}{l}{a}} \dicPos{f}[1] \dicFlx{(‑u)}[5] \dicDirectTranslationCS{podávání, servírování}
\dicEntry[framreiðslumaður] \dicTerm{fram·reiðslu··|maður} \dicIPA{{f}{r}{a}{m}{r}{ei}{ð}{s}{\textsubring{d}}{l}{\textscy}{m}{a}{ð}{\textscy}{\textsubring{r}}} \dicPos{m}[13] \dicFlx{(‑manns, ‑menn)}[2] \dicDirectTranslationCS{číšník}
\dicEntry[framreiðslustúlka] \dicTerm{fram·reiðslu··stúlk|a} \dicIPA{{f}\-{r}\-{a}\-{m}\-{r}\-{ei}\-{ð}\-{s}\-{\textsubring{d}}\-{l}\-{\textscy}\-{s}\-{\textsubring{d}}\-{u}\-{\textsubring{l}}\-{\r{g}}\-{a}\-} \dicPos{f}[1] \dicFlx{(‑u, ‑ur)}[13] \dicDirectTranslationCS{číšnice, servírka}
\dicEntry[framrúða] \dicTerm{fram··rúð|a} \dicIPA{{f}{r}{a}{m}{r}{u}{ð}{a}} \dicPos{f}[1] \dicFlx{(‑u, ‑ur)}[19] \dicDirectTranslationCS{čelní\,/\addthin přední sklo}
\dicEntry[framsal] \dicTerm{fram··|sal} \dicIPA{{f}{r}{a}{m}{s}{a}{\textsubring{l}}} \dicPos{n}[2] \dicFlx{(‑sals, ‑söl)}[8] \textbf{1.} \dicFieldCat{práv.} \dicSynonym*{framsal verðmæta} \dicDirectTranslationCS{postoupení (pohledávky ap.), převedení};  \dicPhraseIS{framsal á e‑u} \dicDirectTranslationCS{postoupení (čeho)} \dicExampleIS{framsal á kröfu} \dicExampleCS{postoupení pohledávky}  \textbf{2.} \dicSynonym*{framsal fanga} \dicDirectTranslationCS{vydání (vězně do jiné země ap.)}
\dicEntry[framselja] \dicTerm{fram··sel|ja} \dicIPA{{f}{r}{a}{m}{s}{\textepsilon}{l}{j}{a}} \dicPos{v}[4] \dicFlx{(‑di, ‑t)}[23] \dicFlx{acc} \textbf{1.} \dicDirectTranslationCS{vydat (k~potrestání ap.)} \dicExampleIS{framselja e‑n til Bandaríkjanna} \dicExampleCS{vydat (koho) do Spojených států}  \textbf{2.} \dicFieldCat{ekon.} \dicDirectTranslationCS{postoupit, převést (nárok ap.)} \dicExampleIS{framselja víxil} \dicExampleCS{převést směnku}
\dicEntry[framseljanlegur] \dicTerm{fram··seljan·legur}\addthinS\dicIPA{{f}{r}{a}{m}{s}{\textepsilon}{l}{j}{a}{n}{l}{\textepsilon}{\textbabygamma}{\textscy}{\textsubring{r}}}\addthinS\dicPos{adj}[1]\dicFlx{}[-8] \dicDirectTranslationCS{převoditelný, směnitelný}
\dicEntry[framsenda] \dicTerm{fram··sen|da} \dicIPA{{f}{r}{a}{m}{s}{\textepsilon}{n}{\textsubring{d}}{a}} \dicPos{v}[2] \dicFlx{(‑di, ‑t)}[141] \dicFlx{acc} \textbf{1.} \dicDirectTranslationCS{poslat dále, přeposlat (dopis ap.)}  \textbf{2.} \dicFieldCat{poč.} \dicDirectTranslationCS{přeposlat (e‑mail ap.)}
\dicEntry[framsetning] \dicTerm{fram··set·ning} \dicIPA{{f}{r}{a}{m}{s}{\textepsilon}{h}{\textsubring{d}}{n}{i}{\ng}{\r{g}}} \dicPos{f}[4] \dicFlx{(‑ar)}[7] \dicSynonym{túlkun} \dicDirectTranslationCS{výklad, interpretace} \dicExampleIS{skýr framsetning} \dicExampleCS{jasná interpretace}
\dicEntry[framsigla] \dicTerm{fram··sigl|a} \dicIPA{{f}{r}{a}{m}{s}{\textsci}{\r{g}}{l}{a}} \dicPos{f}[1] \dicFlx{(‑u, ‑ur)}[13] \dicFieldCat{nám.} \dicDirectTranslationCS{přední stěžeň}
\begin{xtolerant}{}{1pt}
\dicEntry[framsókn] \dicTerm{fram··sókn} \dicIPA{{f}{r}{a}{m}{s}{ou}{h}{\r{g}}{\textsubring{n}}}\addthinS\dicPos{f}[7]\addthinS\dicFlx{(‑ar)}[3] \dicDirectTranslationCS{postup, pokrok} \dicExampleIS{stöðva framsókn hersins} \dicExampleCS{zastavit postup armády}
\end{xtolerant}
\dicEntry[Framsóknarflokkurinn] \dicTerm{Fram·sóknar··flokk|urinn} \dicIPA{{f}\-{r}\-{a}\-{m}\-{s}\-{ou}\-{h}\-{\r{g}}\-{n}\-{a}\-{\textsubring{r}}\-{f}\-{l}\-{\textopeno}\-{h}\-{\r{g}}\-{\textscy}\-{r}\-{\textsci}\-{\textsubring{n}}\-} \dicPos{m}[4] \dicFlx{(‑sins)}[1] \dicFlx{def} \dicFieldCat{pol.} \dicDirectTranslationCS{Strana pokroku} \dicIndirectTranslationCS{(založena roku 1916)}
\dicEntry[framsóknarmaður] \dicTerm{fram·sóknar··|maður} \dicIPA{{f}{r}{a}{m}{s}{ou}{h}{\r{g}}{n}{a}{r}{m}{a}{ð}{\textscy}{\textsubring{r}}} \dicPos{m}[13] \dicFlx{(‑manns, ‑menn)}[2] \dicFieldCat{pol.} \dicDirectTranslationCS{člen(ka) Strany pokroku}
\dicEntry[framstaða] \dicTerm{fram··|staða} \dicIPA{{f}{r}{a}{m}{s}{\textsubring{d}}{a}{ð}{a}} \dicPos{f}[1] \dicFlx{(‑stöðu)}[2] \dicFieldCat{jaz.} \dicDirectTranslationCS{přední pozice}
\dicEntry[framstag] \dicTerm{fram··|stag} \dicIPA{{f}{r}{a}{m}{s}{\textsubring{d}}{a}{x}} \dicPos{n}[2] \dicFlx{(‑stags, ‑stög)}[8] \dicFieldCat{nám.} \dicDirectTranslationCS{přední stěh}
\dicEntry[framsýni] \dicTerm{fram··sýn|i} \dicIPA{{f}{r}{a}{m}{s}{i}{n}{\textsci}} \dicPos{f}[3] \dicFlx{(‑i)}[3] \dicDirectTranslationCS{prozíravost, předvídavost}
\dicEntry[framsýnn] \dicTerm{fram··sýnn} \textls[15]{\dicIPA{{f}{r}{a}{m}{s}{i}{\textsubring{d}}{\textsubring{n}}} \dicPos{adj}[7]\dicFlx{}[-1] \dicSynonym*{forsjálegur}} \dicDirectTranslationCS{prozíravý, před\-vídavý, vizionářský}
\dicEntry[framsækinn] \dicTerm{fram··sækinn} \dicIPA{{f}{r}{a}{m}{s}{a}{i}{\r{\textObardotlessj}}{\textsci}{\textsubring{n}}} \dicPos{adj}[6]\dicFlx{}[-2] \dicDirectTranslationCS{pokrokový, progresivní}
\dicEntry[framsæti] \dicTerm{fram··sæti} \dicIPA{{f}{r}{a}{m}{s}{a}{i}{\textsubring{d}}{\textsci}} \dicPos{n}[2] \dicFlx{(‑s, ‑)}[14] \dicDirectTranslationCS{přední sedadlo (auta ap.)} \dicAntonym{aftursæti}
\dicEntry[framsögn] \dicTerm{fram··|sögn} \dicIPA{{f}{r}{a}{m}{s}{\oe}{\r{g}}{\textsubring{n}}} \dicPos{f}[7] \dicFlx{(‑sagnar)}[19] \dicDirectTranslationCS{výslovnost, dikce, styl\,/\addthin způsob vyjadřování}
\begin{xtolerant}{}{1pt}
\dicEntry[framsöguháttur] \dicTerm{fram·sögu··|háttur} \dicIPA{{f}{r}{a}{m}{s}{\oe}{\textbabygamma}{\textscy}{h}{au}{h}{\textsubring{d}}{\textscy}{\textsubring{r}}}\addthinS\dicPos{m}[12]\addthinS\dicFlx{(‑hátt\-ar, ‑hætt\-ir)}[7]\addthinS\dicFieldCat{jaz.} \dicDirectTranslationCS{oznamovací způsob, indikativ}
\end{xtolerant}
\begin{xtolerant}{}{1pt}
\dicEntry[framsögumaður] \dicTerm{fram·sögu··|maður} \dicIPA{{f}{r}{a}{m}{s}{\oe}{\textbabygamma}{\textscy}{m}{a}{ð}{\textscy}{\textsubring{r}}}\addthinS\dicPos{m}[13]\addthinS\dicFlx{(‑manns, ‑menn)}[2] \dicDirectTranslationCS{mluvčí}
\end{xtolerant}
\dicEntry[framtak] \dicTerm{fram··tak} \dicIPA{{f}{r}{a}{m}{t\smash{\textsuperscript{h}}}{a}{\r{g}}} \dicPos{n}[2] \dicFlx{(‑s)}[2] \dicDirectTranslationCS{iniciativa, podnět}
\dicEntry[framtakslaus] \dicTerm{fram·taks··laus} \dicIPA{{f}{r}{a}{m}{t\smash{\textsuperscript{h}}}{a}{\r{g}}{s}{l}{\oe i}{s}} \dicPos{adj}[5]\dicFlx{}[-1] \dicSynonym{óduglegur} \dicDirectTranslationCS{nepodnikavý}
\dicEntry[framtakssamur] \dicTerm{fram·taks··|samur} \dicIPA{{f}{r}{a}{m}{t\smash{\textsuperscript{h}}}{a}{\r{g}}{s}{a}{m}{\textscy}{\textsubring{r}}} \dicPos{adj}[1] \dicFlx{(f ‑söm)}[2] \dicDirectTranslationCS{podnikavý, iniciativní}
\dicEntry[framtakssemi] \dicTerm{fram·taks··sem|i} \dicIPA{{f}{r}{a}{m}{t\smash{\textsuperscript{h}}}{a}{\r{g}}{s}{\textepsilon}{m}{\textsci}} \dicPos{f}[2] \dicFlx{(‑i)}[2] \dicSynonym{atorka} \dicDirectTranslationCS{podnikavost, píle}
\dicEntry[framtal] \dicTerm{fram··|tal} \dicIPA{{f}{r}{a}{m}{t\smash{\textsuperscript{h}}}{a}{\textsubring{l}}} \dicPos{n}[2] \dicFlx{(‑tals, ‑töl)}[8] \dicFieldCat{ekon.} \dicDirectTranslationCS{daňové přiznání}
\dicEntry[framtíð] \dicTerm{fram··tíð} \dicsymFrequent\  \dicIPA{{f}{r}{a}{m}{t\smash{\textsuperscript{h}}}{i}{\texttheta}} \dicPos{f}[7] \dicFlx{(‑ar)}[3] \textbf{1.} \dicDirectTranslationCS{budoucnost} \dicExampleIS{í framtíðinni} \dicExampleCS{v~budoucnosti}  \textbf{2.} \dicFieldCat{jaz.} \dicDirectTranslationCS{budoucí čas}
\dicEntry[framtönn] \dicTerm{fram··|tönn} \dicIPA{{f}{r}{a}{m}{t\smash{\textsuperscript{h}}}{\oe}{\textsubring{n}}} \dicPos{f}[8] \dicFlx{(‑tannar, ‑tennur)}[2] \dicFieldCat{anat.} \dicDirectTranslationCS{řezák}
\dicEntry[framundan] \dicTerm{fram··undan} \dicsymFrequent\  \dicIPA{{f}{r}{a}{m}{\textlengthmark}{\textscy}{n}{\textsubring{d}}{a}{\textsubring{n}}} \dicPos{adv} \dicPhraseIS{fram undan} \dicLink{fram}
\dicEntry[framúrskarandi] \dicTerm{fram·úr··skar·andi} \dicIPA{{f}{r}{a}{m}{\textlengthmark}{u}{\textsubring{r}}{s}{\r{g}}{a}{r}{a}{n}{\textsubring{d}}{\textsci}} \dicPos{adj}[13] \dicFlx{indecl}[1] \dicSynonym{mikilhæfur} \dicDirectTranslationCS{vynikající, prvotřídní, excelentní} \dicExampleIS{framúrskarandi hæfileikar} \dicExampleCS{prvotřídní schopnosti}
\dicEntry[framúrstefna] \dicTerm{fram·úr··stefn|a} \dicIPA{{f}{r}{a}{m}{\textlengthmark}{u}{\textsubring{r}}{s}{\textsubring{d}}{\textepsilon}{\textsubring{b}}{n}{a}} \dicPos{f}[1] \dicFlx{(‑u)}[5] \dicDirectTranslationCS{avantgarda}
\dicEntry[framvegis] \dicTerm{fram··vegis} \dicsymFrequent\  \dicIPA{{f}{r}{a}{m}{v}{ei}{\textsci}{s}} \dicPos{adv} \dicDirectTranslationCS{od této chvíle, nadále, v~budoucnu} \dicExampleIS{búa hér framvegis} \dicExampleCS{bydlet zde nadále};  \dicPhraseIS{og svo framvegis} \dicDirectTranslationCS{a~tak dále}
\dicEntry[framvinda] \dicTerm{fram··vind|a} \dicIPA{{f}{r}{a}{m}{v}{\textsci}{n}{\textsubring{d}}{a}} \dicPos{f}[1] \dicFlx{(‑u)}[5] \dicSynonym{ferli} \dicDirectTranslationCS{vývoj, průběh} \dicExampleIS{fylgjast með framvindu leiksins} \dicExampleCS{sledovat průběh hry}
\dicEntry[framvísa] \dicTerm{fram··vís|a} \dicIPA{{f}{r}{a}{m}{v}{i}{s}{a}} \dicPos{v}[1] \dicFlx{(‑aði)}[1] \dicFlx{dat} \dicDirectTranslationCS{předložit, předkládat (k~nahlédnutí ap.)} \dicExampleIS{framvísa vegabréfum} \dicExampleCS{předložit pasy}
\dicEntry[framþróun] \dicTerm{fram··þró|un} \dicIPA{{f}{r}{a}{m}{\texttheta}{r}{ou}{\textscy}{\textsubring{n}}} \dicPos{f}[7] \dicFlx{(‑unar)}[9] \dicDirectTranslationCS{rozvoj, vývoj, rozvinutí}
\dicEntry[franki] \dicTerm{frank|i} \dicIPA{{f}{r}{au}{\r{\textltailn}}{\r{\textObardotlessj}}{\textsci}} \dicPos{m}[1] \dicFlx{(‑a, ‑ar)}[8] \dicFieldCat{ekon.} \dicDirectTranslationCS{frank}
\dicEntry[fransín] \dicTerm{fransín} \dicIPA{{f}{r}{a}{n}{s}{i}{\textsubring{n}}} \dicPos{n}[2] \dicFlx{(‑s)}[2] \dicFieldCat{chem.} \dicDirectTranslationCS{francium} \textit{(l.~{\textLA{Fr, Francium}})}
\dicEntry[franska] \dicTerm{franska} \dicIPA{{f}{r}{a}{n}{s}{\r{g}}{a}} \dicPos{f}[1] \dicFlx{(frönsku, frönskur)}[20] \textbf{1.} \dicDirectTranslationCS{francouzština}  \textbf{2.} \dicPhraseIS{frönskur} \dicFlx{pl} \dicSynonym{franskar} \dicDirectTranslationCS{hranolky}
\dicEntry[Franska-Gvæjana] \dicTerm{Franska-Gvæjana} \dicIPA{{f}{r}{a}{n}{s}{\r{g}}{a}{\textbabygamma}{v}{a}{i}{j}{a}{n}{a}} \dicPos{n}[4] \dicFlx{indecl}[2] \dicFieldCat{geog.} \dicDirectTranslationCS{Francouzská Guyana}
\dicEntry[franskar] \dicTerm{franskar} \dicIPA{{f}{r}{a}{n}{s}{\r{g}}{a}{\textsubring{r}}} \dicPos{adj} \dicFlx{pl nom f} \dicLink{franskur}
\dicEntry[franskbrauð] \dicTerm{fransk··brauð} \dicIPA{{f}{r}{a}{n}{s}{\r{g}}{\textsubring{b}}{r}{\oe i}{\texttheta}} \dicPos{n}[2] \dicFlx{(‑s, ‑)}[5] \dicDirectTranslationCS{bílý chléb}
\dicEntry[franskur] \dicTerm{franskur} \dicsymFrequent\  \dicIPA{{f}{r}{a}{n}{s}{\r{g}}{\textscy}{\textsubring{r}}} \dicPos{adj}[1] \dicFlx{(f frönsk)}[3] \dicDirectTranslationCS{francouzský} \dicExampleIS{lesa frönsk ljóð} \dicExampleCS{číst francouzskou báseň};  \dicPhraseIS{franskar (kartöflur)} \dicDirectTranslationCS{hranolky}
\dicEntry[Fransmaður] \dicTerm{Frans··|maður} \dicIPA{{f}{r}{a}{n}{s}{m}{a}{ð}{\textscy}{\textsubring{r}}} \dicPos{m}[13] \dicFlx{(‑manns, ‑menn)}[2] \dicLangCat{hovor.} \dicSynonym{Frakki\smash{\textsuperscript{2}}} \dicDirectTranslationCS{Francouz(ka)}
\dicEntry[frasi] \dicTerm{fras|i} \dicIPA{{f}{r}{a}{\textlengthmark}{s}{\textsci}} \dicPos{m}[1] \dicFlx{(‑a, ‑ar)}[8] \dicSynonym{orðtak} \dicDirectTranslationCS{fráze, klišé}
\dicEntry[frat] \dicTerm{frat} \dicIPA{{f}{r}{a}{\textlengthmark}{\textsubring{d}}} \dicPos{n}[2] \dicFlx{(‑s)}[2] \dicLangCat{hovor.} \dicDirectTranslationCS{harampádí, šunt};  \dicPhraseIS{gefa frat í e‑ð} \dicLangCat{přen.} \dicDirectTranslationCS{kašlat na (co)}
\dicEntry[frauðplast] \dicTerm{frauð··plast} \dicIPA{{f}{r}{\oe i}{ð}{p\smash{\textsuperscript{h}}}{l}{a}{s}{\textsubring{d}}} \dicPos{n}[2] \dicFlx{(‑s)}[2] \dicDirectTranslationCS{polystyren}
\dicEntry[fraus] \dicTerm{fraus} \dicIPA{{f}{r}{\oe i}{\textlengthmark}{s}} \dicPos{v} \dicFlx{ind pf sg 1 pers} \dicLink{frjósa}
\dicEntry[frá] \dicTerm{frá} \dicsymFrequent\  \dicIPA{{f}{r}{au}{\textlengthmark}} \dicPos{prep\,/\addthin adv} \dicFlx{dat} \textbf{1.} \dicDirectTranslationCS{z(e), od(e)} \dicIndirectTranslationCS{(o~pohybu nebo směru pryč)} \dicExampleIS{frá landi} \dicExampleCS{ze země};  \dicPhraseIS{farðu frá} \dicFlx{imper} \dicDirectTranslationCS{uhni}  \textbf{2.} \dicDirectTranslationCS{z(e)} \dicIndirectTranslationCS{(o~původu něčeho\,/\addthin někoho)} \dicExampleIS{frá Tékklandi} \dicExampleCS{z~České republiky}  \textbf{3.} \dicDirectTranslationCS{od(e)} \dicIndirectTranslationCS{(o~vzdálenosti)} \dicExampleIS{langt frá öðrum} \dicExampleCS{daleko od ostatních}  \textbf{4.} \dicDirectTranslationCS{od(e)} \dicIndirectTranslationCS{(o~čase)} \dicExampleIS{Hann vann frá klukkan fjögur til klukkan fimm.} \dicExampleCS{Pracoval od čtyř do pěti.};  \dicPhraseIS{frá upphafi} \dicFlx{adv} \dicDirectTranslationCS{od začátku}  \textbf{5.} \dicDirectTranslationCS{za} \dicIndirectTranslationCS{(o~pořadí)} \dicExampleIS{dag frá degi} \dicExampleCS{den za dnem};  \dicIdiom{frá}{ \dicPhraseIS{frá sér}} \dicFlx{adj} \dicDirectTranslationCS{(jsoucí) mimo (ze strachu ap.)}; { \dicPhraseIS{til og frá}} \dicFlx{adv} \dicDirectTranslationCS{tam a~zpátky\,/\addthin zpět};  \dicPhraseIS{vera frá} {\textbf{a.}} \dicDirectTranslationCS{být nepřítomný, chybět};  {\textbf{b.}} \dicDirectTranslationCS{mít velké bolesti};  {\textbf{c.}} \dicDirectTranslationCS{být hotový\,/\addthin skončený (úkol ap.)}
\dicEntry[frábitinn] \dicTerm{frá··bitinn} \dicIPA{{f}{r}{au}{\textlengthmark}{\textsubring{b}}{\textsci}{\textsubring{d}}{\textsci}{\textsubring{n}}} \dicPos{adj}[6]\dicFlx{}[-2] \dicPhraseIS{vera frábitinn e‑u} \dicDirectTranslationCS{mít nechuť\,/\addthin odpor k~(čemu), nemít se k~(čemu)} \dicExampleIS{vera frábitinn námi} \dicExampleCS{mít nechuť ke studiu}
\dicEntry[fráblásinn] \dicTerm{frá··blásinn} \dicIPA{{f}{r}{au}{\textlengthmark}{\textsubring{b}}{l}{au}{s}{\textsci}{\textsubring{n}}} \dicPos{adj}[6]\dicFlx{}[-6] \dicFieldCat{jaz.} \dicDirectTranslationCS{aspirovaný}
\dicEntry[fráblástur] \dicTerm{frá··blást|ur} \dicIPA{{f}{r}{au}{\textlengthmark}{\textsubring{b}}{l}{au}{s}{\textsubring{d}}{\textscy}{\textsubring{r}}} \dicPos{m}[5] \dicFlx{(‑urs)}[12] \dicFieldCat{jaz.} \dicDirectTranslationCS{aspirace, přídech}
\dicEntry[frábrugðinn] \dicTerm{frá··brugðinn} \dicsymFrequent\  \dicIPA{{f}{r}{au}{\textlengthmark}{\textsubring{b}}{r}{\textscy}{\textbabygamma}{ð}{\textsci}{\textsubring{n}}} \dicPos{adj}[6]\dicFlx{}[-2] \dicDirectTranslationCS{odlišný, rozdílný, atypický};  \dicPhraseIS{e‑að er frábrugðið e‑u} \dicDirectTranslationCS{(co) je odlišné od (čeho)}
\dicEntry[frábær] \dicTerm{frá··bær} \dicsymFrequent\  \dicIPA{{f}{r}{au}{\textlengthmark}{\textsubring{b}}{a}{i}{\textsubring{r}}} \dicPos{adj}[5] \dicFlx{(f ‑)}[8] \dicSynonym{ágætur} \dicDirectTranslationCS{vynikající, výborný, fantastický, skvělý, senzační, úžasný, super} \dicExampleIS{frábær tónlist} \dicExampleCS{vynikající hudba}
\dicEntry[frábærlega] \dicTerm{frá·bær··lega} \dicIPA{{f}{r}{au}{\textlengthmark}{\textsubring{b}}{a}{i}{r}{l}{\textepsilon}{\textbabygamma}{a}} \dicPos{adv} \dicSynonym{ákaflega} \dicDirectTranslationCS{výborně, skvěle, senzačně, úžasně, bezvadně}
\dicEntry[frádráttarbær] \dicTerm{frá·dráttar··bær} \dicIPA{{f}{r}{au}{\textlengthmark}{\textsubring{d}}{r}{au}{h}{\textsubring{d}}{a}{r}{\textsubring{b}}{a}{i}{\textsubring{r}}} \dicPos{adj}[5] \dicFlx{(f ‑)}[10] \dicDirectTranslationCS{odečitatelný, odpočitatelný};  \dicPhraseIS{frádráttarbær útgjöld} \dicFieldCat{ekon.} \dicDirectTranslationCS{odečitatelné výdaje}
\dicEntry[frádráttarmerki] \dicTerm{frá·dráttar··merki} \dicIPA{{f}{r}{au}{\textlengthmark}{\textsubring{d}}{r}{au}{h}{\textsubring{d}}{a}{r}{m}{\textepsilon}{\textsubring{r}}{\r{\textObardotlessj}}{\textsci}} \dicPos{n}[2] \dicFlx{(‑s, ‑)}[16] \dicFieldCat{mat.} \dicDirectTranslationCS{znaménko odčítání, minus}
\dicEntry[frádráttur] \dicTerm{frá··drátt|ur} \dicIPA{{f}{r}{au}{\textlengthmark}{\textsubring{d}}{r}{au}{h}{\textsubring{d}}{\textscy}{\textsubring{r}}} \dicPos{m}[12] \dicFlx{(‑ar)}[8] \textbf{1.} \dicFieldCat{mat.} \dicDirectTranslationCS{odčítání, odečítání}  \textbf{2.} \dicFieldCat{ekon.} \dicDirectTranslationCS{odpočet, odečítání, odečtení, srážka (z~daní)}
\dicEntry[fráfall] \dicTerm{frá··|fall} \dicIPA{{f}{r}{au}{\textlengthmark}{f}{a}{\textsubring{d}}{\textsubring{l}}} \dicPos{n}[2] \dicFlx{(‑falls, ‑föll)}[8] \dicSynonym{dauði} \dicDirectTranslationCS{úmrtí, skon}
\dicEntry[frágangssök] \dicTerm{frá·gangs··|sök} \dicIPA{{f}{r}{au}{\textlengthmark}{\r{g}}{au}{\ng}{s}{\oe}{\r{g}}} \dicPos{f}[7] \dicFlx{(‑sakar)}[19] \dicPhraseIS{það er ekki\,/\addthin engin frágangssök að (gera e‑ð)} \dicDirectTranslationCS{není nemožné (dělat (co))}
\dicEntry[frágangur] \dicTerm{frá··gang|ur} \dicIPA{{f}{r}{au}{\textlengthmark}{\r{g}}{au}{\ng}{\r{g}}{\textscy}{\textsubring{r}}} \dicPos{m}[6] \dicFlx{(‑s)}[7] \dicDirectTranslationCS{zpracování, dokončení, kompletace} \dicExampleIS{góður frágangur á ritinu} \dicExampleCS{dobré zpracování písemné práce}
\dicEntry[frágenginn] \dicTerm{frá··genginn} \dicIPA{{f}{r}{au}{\textlengthmark}{\r{\textObardotlessj}}{ei}{\textltailn}{\r{\textObardotlessj}}{\textsci}{\textsubring{n}}} \dicPos{adj}[6]\dicFlx{}[-6] \dicDirectTranslationCS{provedený, dokončený} \dicExampleIS{Málið er endanlega frágengið.} \dicExampleCS{Věc je definitivně vyřízená.}
\dicEntry[fráhnepptur] \dicTerm{frá··hnepptur} \dicIPA{{f}{r}{au}{\textlengthmark}{\textsubring{n}}{\textepsilon}{f}{\textsubring{d}}{\textscy}{\textsubring{r}}} \dicPos{adj}[1]\dicFlx{}[-13] \dicDirectTranslationCS{rozepnutý (kabát z~knoflíků ap.)} \dicExampleIS{Frakkinn er fráhnepptur.} \dicExampleCS{Kabát je rozepnutý.}
\dicEntry[fráhrindandi] \dicTerm{frá··hrind·andi} \dicIPA{{f}{r}{au}{\textlengthmark}{\textsubring{r}}{\textsci}{n}{\textsubring{d}}{a}{n}{\textsubring{d}}{\textsci}} \dicPos{adj}[13] \dicFlx{indecl}[1] \dicSynonym{þurrlegur} \dicDirectTranslationCS{odpudivý, nepřívětivý} \dicExampleIS{vera fráhrindandi í útliti} \dicExampleCS{mít odpudivý vzhled}
\dicEntry[fráhverfur] \dicTerm{frá··hverfur} \dicIPA{{f}{r}{au}{\textlengthmark}{k\smash{\textsuperscript{h}}}{v}{\textepsilon}{r}{v}{\textscy}{\textsubring{r}}} \dicPos{adj}[1]\dicFlx{}[-1] \dicDirectTranslationCS{mající nechuť\,/\addthin odpor\,/\addthin averzi};  \dicPhraseIS{vera fráhverfur e‑u} \dicDirectTranslationCS{mít odpor k~(čemu)};  \dicPhraseIS{gerast e‑u fráhverfur} \dicFlx{refl} \dicDirectTranslationCS{odvrátit se od (čeho), upustit od (čeho)}
\dicEntry[frákast] \dicTerm{frá··|kast} \dicIPA{{f}{r}{au}{\textlengthmark}{k\smash{\textsuperscript{h}}}{a}{s}{\textsubring{d}}} \dicPos{n}[2] \dicFlx{(‑kasts, ‑köst)}[8] \dicFieldCat{sport.} \dicDirectTranslationCS{doskok}
\dicEntry[frálag] \dicTerm{frá··lag} \dicIPA{{f}{r}{au}{\textlengthmark}{l}{a}{x}} \dicPos{n}[2] \dicFlx{(‑s)}[2] \dicFieldCat{poč.} \dicDirectTranslationCS{výstup, output} \dicAntonym{ílag}
\dicEntry[fráleitur] \dicTerm{frá··leitur} \dicIPA{{f}{r}{au}{\textlengthmark}{l}{ei}{\textsubring{d}}{\textscy}{\textsubring{r}}} \dicPos{adj}[1]\dicFlx{}[-1] \dicSynonym{afleitur} \dicDirectTranslationCS{absurdní, nesmyslný} \dicExampleIS{fráleit hugmynd} \dicExampleCS{nesmyslný nápad}
\dicEntry[frár] \dicTerm{frár} \dicIPA{{f}{r}{au}{\textlengthmark}{\textsubring{r}}} \dicPos{adj}[4]\dicFlx{}[-1] \dicSynonym{fljótur} \dicDirectTranslationCS{hbitý} \dicExampleIS{frár á fæti} \dicExampleCS{rychlonohý}
\dicEntry[frárennsli] \dicTerm{frá··rennsli} \dicIPA{{f}{r}{au}{\textlengthmark}{r}{\textepsilon}{n}{s}{\textsubring{d}}{l}{\textsci}} \dicPos{n}[2] \dicFlx{(‑s, ‑)}[14] \dicDirectTranslationCS{odtok, odpad} \dicAntonym{aðrennsli}
\dicEntry[frásaga] \dicTerm{frá··|saga} \dicIPA{{f}{r}{au}{\textlengthmark}{s}{a}{\textbabygamma}{a}} \dicPos{f}[1] \dicFlx{(‑sögu, ‑sögur)}[14] \dicSynonym{umtal} \dicDirectTranslationCS{vyprávění, povídání};  \dicPhraseIS{þetta er ekki í frásögur færandi} \dicDirectTranslationCS{to nestojí za řeč}
\dicEntry[frásagnarfræði] \dicTerm{frá·sagnar··fræð|i} \dicIPA{{f}{r}{au}{\textlengthmark}{s}{a}{\r{g}}{n}{a}{\textsubring{r}}{f}{r}{a}{i}{ð}{\textsci}} \dicPos{f}[3] \dicFlx{(‑i)}[3] \dicFieldCat{lit.} \dicDirectTranslationCS{techniky vyprávění}
\dicEntry[frásagnarlist] \dicTerm{frá·sagnar··list} \dicIPA{{f}{r}{au}{\textlengthmark}{s}{a}{\r{g}}{n}{a}{r}{l}{\textsci}{s}{\textsubring{d}}} \dicPos{f}[7] \dicFlx{(‑ar)}[3] \dicDirectTranslationCS{vyprávěčství, vypravěčství}
\dicEntry[fráskilinn] \dicTerm{frá··skilinn} \dicIPA{{f}{r}{au}{\textlengthmark}{s}{\r{\textObardotlessj}}{\textsci}{l}{\textsci}{\textsubring{n}}} \dicPos{adj}[6]\dicFlx{}[-10] \textbf{1.} \dicSynonym{skilinn} \dicDirectTranslationCS{rozvedený}  \textbf{2.} \dicSynonym{undanskilinn} \dicDirectTranslationCS{nezapočítaný, vynechaný}
\dicEntry[frásögn] \dicTerm{frá··|sögn} \dicsymFrequent\  \dicIPA{{f}{r}{au}{\textlengthmark}{s}{\oe}{\r{g}}{\textsubring{n}}} \dicPos{f}[7] \dicFlx{(‑sagnar, ‑sagnir)}[16] \dicDirectTranslationCS{vyprávění, vypravování} \dicExampleIS{frásögn af ferðinni} \dicExampleCS{vyprávění z~cesty}
\dicEntry[frátalinn] \dicTerm{frá··talinn} \dicIPA{{f}{r}{au}{\textlengthmark}{t\smash{\textsuperscript{h}}}{a}{l}{\textsci}{\textsubring{n}}} \dicPos{adj}[6]\dicFlx{}[-9] \dicSynonym{undanskilinn} \dicDirectTranslationCS{nezapočítaný, vynechaný};  \dicPhraseIS{að frátöldum e‑u} \dicFlx{prep} \dicDirectTranslationCS{s~výjimkou (čeho), vyjma (čeho)}
\dicEntry[frátekinn] \dicTerm{frá··tekinn} \dicIPA{{f}{r}{au}{\textlengthmark}{t\smash{\textsuperscript{h}}}{\textepsilon}{\r{\textObardotlessj}}{\textsci}{\textsubring{n}}} \dicPos{adj}[6]\dicFlx{}[-6] \textbf{1.} \dicSynonym*{pantaður} \dicDirectTranslationCS{rezervovaný, zamluvený, vyhrazený} \dicExampleIS{Sætið er frátekið fyrir hana.} \dicExampleCS{Má rezervované tohle místo.}  \textbf{2.} \dicSynonym{undanskilinn} \dicDirectTranslationCS{nezapočítaný, nepočítaný} \dicExampleIS{að fráteknum þremur dögum} \dicExampleCS{s~výjimkou tří dnů}
\dicEntry[frávera] \dicTerm{frá··ver|a} \dicIPA{{f}{r}{au}{\textlengthmark}{v}{\textepsilon}{r}{a}} \dicPos{f}[1] \dicFlx{(‑u)}[5] \dicSynonym{fjarvist} \dicDirectTranslationCS{nepřítomnost, absence} \dicExampleIS{í fráveru e‑rs} \dicExampleCS{v~(čí) nepřítomnosti}
\dicEntry[frávik] \dicTerm{frá··vik} \dicIPA{{f}{r}{au}{\textlengthmark}{v}{\textsci}{\r{g}}} \dicPos{n}[2] \dicFlx{(‑s, ‑)}[5] \textbf{1.} \dicDirectTranslationCS{odchylka, odchýlení} \dicExampleIS{frávik frá réttri stefnu} \dicExampleCS{odchylka od správného směru}  \textbf{2.} \dicSynonym{undantekning} \dicDirectTranslationCS{výjimka} \dicExampleIS{frávik frá almennum reglum} \dicExampleCS{výjimka od všeobecných pravidel}
\dicEntry[freðfiskur] \dicTerm{freð··fisk|ur} \dicIPA{{f}{r}{\textepsilon}{ð}{f}{\textsci}{s}{\r{g}}{\textscy}{\textsubring{r}}} \dicPos{m}[6] \dicFlx{(‑s, ‑ar)}[10] \dicDirectTranslationCS{mražená ryba}
\dicEntry[freðinn] \dicTerm{freðinn} \dicIPA{{f}{r}{\textepsilon}{\textlengthmark}{ð}{\textsci}{\textsubring{n}}} \dicPos{adj}[6]\dicFlx{}[-6] \textbf{1.} \dicSynonym{frosinn} \dicDirectTranslationCS{(z)mražený, zmrzlý, zamrzlý}  \textbf{2.} \dicDirectTranslationCS{studený, mrazivý (pohled ap.)}
\dicEntry[freðmýri] \dicTerm{freð··mýr|i} \dicIPA{{f}{r}{\textepsilon}{ð}{m}{i}{r}{\textsci}} \dicPos{f}[5] \dicFlx{(‑ar, ‑ar)}[3] \dicDirectTranslationCS{tundra}
\dicEntry[fregn] \dicTerm{fregn} \dicIPA{{f}{r}{\textepsilon}{\r{g}}{\textsubring{n}}} \dicPos{f}[7] \dicFlx{(‑ar, ‑ir)}[1] \dicSynonym{frétt} \dicDirectTranslationCS{zpráva, novina} \dicExampleIS{Fregnir hafa borist til Íslands.} \dicExampleCS{Zprávy se dostaly na Island.}
\dicEntry[freigáta] \dicTerm{freigát|a} \dicIPA{{f}{r}{ei}{\textlengthmark}{\textbabygamma}{au}{\textsubring{d}}{a}} \dicPos{f}[1] \dicFlx{(‑u, ‑ur)}[19] \dicFieldCat{nám.} \dicDirectTranslationCS{fregata}
\dicEntry[freista] \dicTerm{freist|a} \dicsymFrequent\  \dicIPA{{f}{r}{ei}{s}{\textsubring{d}}{a}} \dicPos{v}[1] \dicFlx{(‑aði)}[22] \dicFlx{gen} \textbf{1.} \dicPhraseIS{e‑að freistar e‑rs} \dicDirectTranslationCS{(co koho) láká, (co koho) uvádí do pokušení} \dicExampleIS{Hluturinn freistar mín ekki.} \dicExampleCS{Ta věc mě neláká.}  \textbf{2.} \dicPhraseIS{e‑r freistar e‑s} \dicSynonym*{e‑r reynir} \dicDirectTranslationCS{(co kdo) zkouší};  \dicPhraseIS{freista gæfunnar} \dicDirectTranslationCS{pokoušet štěstí};  \dicIdiom{freistast}{ \dicPhraseIS{freistast}} \dicFlx{refl} \dicDirectTranslationCS{(u)padnout v~pokušení} \dicExampleIS{Hún freistaðist.} \dicExampleCS{Upadla v~pokušení.}
\dicEntry[freistandi] \dicTerm{freist··andi} \dicIPA{{f}{r}{ei}{s}{\textsubring{d}}{a}{n}{\textsubring{d}}{\textsci}} \dicPos{adj}[13] \dicFlx{indecl}[1] \dicSynonym{teygilegur} \dicDirectTranslationCS{lákající, lákavý, svůdný, vábný, vábivý};  \dicPhraseIS{það er freistandi að (gera e‑ð)} \dicDirectTranslationCS{je lákavé (udělat (co))}
\dicEntry[freistari] \dicTerm{freist··ar|i} \dicIPA{{f}{r}{ei}{s}{\textsubring{d}}{a}{r}{\textsci}} \dicPos{m}[1] \dicFlx{(‑a, ‑ar)}[13] \dicDirectTranslationCS{pokušitel(ka), svůdce, svůdkyně};  \dicPhraseIS{freistarinn} \dicFlx{def} \dicFieldCat{náb.} \dicDirectTranslationCS{ďábel}
\dicEntry[freisting] \dicTerm{freist··ing} \dicIPA{{f}{r}{ei}{s}{\textsubring{d}}{i}{\ng}{\r{g}}} \dicPos{f}[4] \dicFlx{(‑ar, ‑ar)}[5] \dicSynonym{raun} \dicDirectTranslationCS{pokušení, lákání, vábení} \dicExampleIS{Þetta er honum mikil freisting.} \dicExampleCS{Je to pro něj velké pokušení.};  \dicPhraseIS{falla fyrir freistingunni} \dicDirectTranslationCS{podlehnout pokušení};  \dicPhraseIS{sigrast á freistingunni} \dicFlx{refl} \dicDirectTranslationCS{odolat pokušení}
\dicEntry[frekar] \dicTerm{frekar} \dicsymFrequent\  \dicIPA{{f}{r}{\textepsilon}{\textlengthmark}{\r{g}}{a}{\textsubring{r}}} \dicPos{adv} \dicFlx{comp (sup frekast)} \textbf{1.} \dicDirectTranslationCS{spíš(e)} \dicExampleIS{Það er frekar regla en undantekning.} \dicExampleCS{Je to spíš pravidlo než výjimka.}  \textbf{2.} \dicDirectTranslationCS{poměrně, docela, celkem} \dicExampleIS{Hún var frekar lagleg.} \dicExampleCS{Byla docela\,/\addthin celkem pohledná.}  \textbf{3.} \dicDirectTranslationCS{blíže, podrobněji}
\dicEntry[frekari] \dicTerm{frek··ari} \dicIPA{{f}{r}{\textepsilon}{\textlengthmark}{\r{g}}{a}{r}{\textsci}} \dicPos{adj}[1] \dicFlx{comp}[5] \dicDirectTranslationCS{další} \dicExampleIS{frekari upplýsingar} \dicExampleCS{další informace}
\dicEntry[frekast] \dicTerm{frekast} \dicIPA{{f}{r}{\textepsilon}{\textlengthmark}{\r{g}}{a}{s}{\textsubring{d}}} \dicPos{adv} \dicFlx{sup} \dicLink{frekar}
\dicEntry[frekja] \dicTerm{frekj|a} \dicIPA{{f}{r}{\textepsilon}{\textlengthmark}{\r{\textObardotlessj}}{a}} \dicPos{f}[1] \dicFlx{(‑u, ‑ur)}[7] \textbf{1.} \dicSynonym{ágengni} \dicDirectTranslationCS{drzost, opovážlivost, troufalost}  \textbf{2.} \dicSynonym*{frek manneskja} \dicDirectTranslationCS{drzoun, troufalec, dotěrný\,/\addthin opovážlivý člověk}
\dicEntry[frekjulegur] \dicTerm{frekju··legur} \dicIPA{{f}{r}{\textepsilon}{\textlengthmark}{\r{\textObardotlessj}}{\textscy}{l}{\textepsilon}{\textbabygamma}{\textscy}{\textsubring{r}}} \dicPos{adj}[1]\dicFlx{}[-8] \dicDirectTranslationCS{drzý, opovážlivý, troufalý}
\dicEntry[frekna] \dicTerm{frekn|a} \dicIPA{{f}{r}{\textepsilon}{h}{\r{g}}{n}{a}} \dicPos{f}[1] \dicFlx{(‑u, ‑ur)}[7] \dicDirectTranslationCS{piha}
\dicEntry[freknóttur] \dicTerm{frekn··óttur} \dicIPA{{f}{r}{\textepsilon}{h}{\r{g}}{\textsubring{n}}{ou}{h}{\textsubring{d}}{\textscy}{\textsubring{r}}} \dicPos{adj}[1]\dicFlx{}[-10] \dicDirectTranslationCS{pihovatý, pihatý}
\dicEntry[frekur] \dicTerm{frekur} \dicsymFrequent\  \dicIPA{{f}{r}{\textepsilon}{\textlengthmark}{\r{g}}{\textscy}{\textsubring{r}}} \dicPos{adj}[1]\dicFlx{}[-1] \textbf{1.} \dicSynonym{ráðríkur} \dicDirectTranslationCS{panovačný, pánovitý}  \textbf{2.} \dicSynonym{gráðugur} \dicDirectTranslationCS{chamtivý, nenasytný, náročný};  \dicPhraseIS{vera frekur á e‑ð} \dicDirectTranslationCS{lačnit po (čem), lakotit s~(čím)}
\dicEntry[frelsa] \dicTerm{frels|a} \dicIPA{{f}{r}{\textepsilon}{l}{s}{a}} \dicPos{v}[1] \dicFlx{(‑aði)}[1] \dicFlx{acc} \dicDirectTranslationCS{osvobodit, osvobozovat, vysvobodit, vysvobozovat} \dicExampleIS{frelsa þjóðina af ófriði} \dicExampleCS{osvobodit národ od války};  \dicIdiom{frelsast}{ \dicPhraseIS{frelsast}} \dicFlx{refl} \dicFieldCat{náb.} \dicDirectTranslationCS{být očištěn}
\dicEntry[frelsari] \dicTerm{frels··ar|i} \dicIPA{{f}{r}{\textepsilon}{l}{s}{a}{r}{\textsci}} \dicPos{m}[1] \dicFlx{(‑a, ‑ar)}[13] \textbf{1.} \dicLangCat{han.} \dicDirectTranslationCS{spasitel, zachránce}  \textbf{2.} \dicPhraseIS{frelsarinn} \dicFlx{def} \dicFieldCat{náb.} \dicDirectTranslationCS{Spasitel (Kristus)}
\dicEntry[frelsi] \dicTerm{frelsi} \dicsymFrequent\  \dicIPA{{f}{r}{\textepsilon}{l}{s}{\textsci}} \dicPos{n}[2] \dicFlx{(‑s)}[20] \dicSynonym{sjálfstæði} \dicDirectTranslationCS{svoboda, volnost} \dicExampleIS{frelsi einstaklingsins} \dicExampleCS{svoboda jednotlivce}
\dicEntry[frelsingi] \dicTerm{frels··ing|i} \dicIPA{{f}{r}{\textepsilon}{l}{s}{i}{\textltailn}{\r{\textObardotlessj}}{\textsci}} \dicPos{m}[1] \dicFlx{(‑ja, ‑jar)}[14] \textbf{1.} \dicLangCat{zast.} \dicSynonym{leysingi} \dicDirectTranslationCS{osvobozený otrok}  \textbf{2.} \dicDirectTranslationCS{volný pěšec} \dicIndirectTranslationCS{(v~šachách)}
\dicEntry[frelsisbarátta] \dicTerm{frelsis··bar·átt|a} \dicIPA{{f}{r}{\textepsilon}{l}{s}{\textsci}{s}{\textsubring{b}}{a}{r}{au}{h}{\textsubring{d}}{a}} \dicPos{f}[1] \dicFlx{(‑u)}[5] \dicDirectTranslationCS{boj za svobodu}
\dicEntry[frelsishetja] \dicTerm{frelsis··hetj|a} \dicIPA{{f}{r}{\textepsilon}{l}{s}{\textsci}{s}{h}{\textepsilon}{\textsubring{d}}{j}{a}} \dicPos{f}[1] \dicFlx{(‑u, ‑ur)}[7] \dicDirectTranslationCS{osvoboditel(ka), bojovník\,/\addthin bojovnice za svobodu}
\dicEntry[frelsistilfinning] \dicTerm{frelsis··til·finn·ing} \dicIPA{{f}\-{r}\-{\textepsilon}\-{l}\-{s}\-{\textsci}\-{s}\-{t\smash{\textsuperscript{h}}}\-{\textsci}\-{l}\-{f}\-{\textsci}\-{n}\-{i}\-{\ng}\-{\r{g}}\-} \dicPos{f}[4] \dicFlx{(‑ar, ‑ar)}[5] \dicDirectTranslationCS{pocit svobody}
\dicEntry[frelsun] \dicTerm{frels|un} \dicIPA{{f}{r}{\textepsilon}{l}{s}{\textscy}{\textsubring{n}}} \dicPos{f}[7] \dicFlx{(‑unar)}[9] \dicSynonym{lausn} \dicDirectTranslationCS{osvobození, vysvobození}
\dicEntry[frem] \dicTerm{frem} \dicIPA{{f}{r}{\textepsilon}{\textlengthmark}{\textsubring{m}}} \dicPos{v} \dicFlx{ind praes sg 1 pers} \dicLink{fremja}
\dicEntry[fremdardýr] \dicTerm{fremdar··dýr} \dicIPA{{f}{r}{\textepsilon}{m}{\textsubring{d}}{a}{r}{\textsubring{d}}{i}{\textsubring{r}}} \dicPos{n}[2] \dicFlx{(‑s, ‑)}[5] \dicFieldCat{zool.} \dicDirectTranslationCS{primát} \textit{(l.~{\textLA{Primates}})}
\dicEntry[fremdi] \dicTerm{fremdi} \dicIPA{{f}{r}{\textepsilon}{m}{\textsubring{d}}{\textsci}} \dicPos{v} \dicFlx{con pf sg 1 pers} \dicLink{fremja}
\dicEntry[fremi] \dicTerm{fremi} \dicIPA{{f}{r}{\textepsilon}{\textlengthmark}{m}{\textsci}} \dicPos{adv} \dicPhraseIS{svo fremi sem} \dicFlx{conj} \dicDirectTranslationCS{pokud, dokud}
\dicEntry[fremja] \dicTerm{fremja} \dicsymFrequent\  \dicIPA{{f}{r}{\textepsilon}{m}{j}{a}} \dicPos{v}[4] \dicFlx{(frem, framdi, frömdum, fremdi, framið)}[6] \dicFlx{acc} \dicDirectTranslationCS{(s)páchat, dopustit se, dopouštět se} \dicExampleIS{fremja lögbrot} \dicExampleCS{spáchat trestný čin}
\dicEntry[fremri] \dicTerm{fremri} \dicsymFrequent\  \dicIPA{{f}{r}{\textepsilon}{m}{r}{\textsci}} \dicPos{adj}[12] \dicFlx{comp (sup fremstur)}[6] \textbf{1.} \dicSynonym*{sem er framar} \dicDirectTranslationCS{první (ze dvou), přední};  \dicPhraseIS{fremri röðin} \dicDirectTranslationCS{přední řada}  \textbf{2.} \dicSynonym{betri} \dicDirectTranslationCS{lepší} \dicExampleIS{Hún er mér fremri í ensku.} \dicExampleCS{Je lepší na angličtinu než já.}
\dicEntry[fremst] \dicTerm{fremst} \dicIPA{{f}{r}{\textepsilon}{m}{s}{\textsubring{d}}} \dicPos{adv} \dicFlx{sup (pos frammi\,/\addthin harla, comp framar\,/\addthin fremur)} \dicDirectTranslationCS{nejdále, nejvíce vpřed}
\dicEntry[fremstur] \dicTerm{fremstur} \dicIPA{{f}{r}{\textepsilon}{m}{s}{\textsubring{d}}{\textscy}{\textsubring{r}}} \dicPos{adj}[12] \dicFlx{sup (comp fremri)}[6] \textbf{1.} \dicSynonym*{sem er fremst} \dicDirectTranslationCS{nejpřednější, (jsoucí) nejvíce vepředu};  \dicPhraseIS{vera fremstur í flokki} \dicDirectTranslationCS{stát v~čele}  \textbf{2.} \dicSynonym{bestur} \dicDirectTranslationCS{nejpřednější, stojící v~čele}
\dicEntry[fremur] \dicTerm{fremur} \dicsymFrequent\  \dicIPA{{f}{r}{\textepsilon}{\textlengthmark}{m}{\textscy}{\textsubring{r}}} \dicPos{adv} \dicFlx{comp (pos harla, sup fremst)} \textbf{1.} \dicSynonym{held\-ur\smash{\textsuperscript{2}}} \dicDirectTranslationCS{spíše, více};  \dicPhraseIS{fremur en} \dicFlx{conj} \dicDirectTranslationCS{spíše než}  \textbf{2.} \dicSynonym*{þó nokkuð} \dicDirectTranslationCS{celkem, docela, relativně} \dicExampleIS{fremur einfaldur} \dicExampleCS{relativně jednoduchý}
\dicEntry[freska] \dicTerm{fresk|a} \dicIPA{{f}{r}{\textepsilon}{s}{\r{g}}{a}} \dicPos{f}[1] \dicFlx{(‑u, ‑ur)}[13] \dicDirectTranslationCS{freska}
\dicEntry[fress] \dicTerm{fress} \dicIPA{{f}{r}{\textepsilon}{s}{\textlengthmark}} \dicPos{m}[4] \dicFlx{(‑, ‑ar)}[31] \dicSynonym{högni} \dicDirectTranslationCS{kocour}
\dicEntry[fressköttur] \dicTerm{fress··|köttur} \dicIPA{{f}{r}{\textepsilon}{s}{\textlengthmark}{k\smash{\textsuperscript{h}}}{\oe}{h}{\textsubring{d}}{\textscy}{\textsubring{r}}} \dicPos{m}[11] \dicFlx{(‑kattar, ‑kettir)}[5] \dicDirectTranslationCS{kocour}
\dicEntry[fresta] \dicTerm{frest|a} \dicIPA{{f}{r}{\textepsilon}{s}{\textsubring{d}}{a}} \dicPos{v}[1] \dicFlx{(‑aði)}[1] \dicFlx{dat} \dicDirectTranslationCS{odložit, odkládat, odsunout, odsouvat, odročit, odročovat} \dicExampleIS{Fundinum verður frestað.} \dicExampleCS{Schůze bude odložena.};  \dicIdiom{frestast}{ \dicPhraseIS{frestast}} \dicFlx{refl} \dicDirectTranslationCS{být odložen\,/\addthin odročen, odložit se}
\dicEntry[frestun] \dicTerm{frest|un} \dicIPA{{f}{r}{\textepsilon}{s}{\textsubring{d}}{\textscy}{\textsubring{n}}} \dicPos{f}[7] \dicFlx{(‑unar, ‑anir)}[8] \dicDirectTranslationCS{odložení, odklad} \dicExampleIS{frestun málsins} \dicExampleCS{odklad věci}
\dicEntry[frestur] \dicTerm{frest|ur} \dicsymFrequent\  \dicIPA{{f}{r}{\textepsilon}{s}{\textsubring{d}}{\textscy}{\textsubring{r}}} \dicPos{m}[9] \dicFlx{(‑s, ‑ir)}[4] \textbf{1.} \dicDirectTranslationCS{lhůta, uzávěrka, termín} \dicExampleIS{Fresturinn rennur út 13. maí.} \dicExampleCS{Lhůta vyprší třináctého května.}  \textbf{2.} \dicDirectTranslationCS{odklad};  \dicPhraseIS{á hálftíma fresti} \dicFlx{adv} \dicDirectTranslationCS{co půl hodiny, každou půlhodinu};  \dicPhraseIS{slá e‑u á frest} \dicDirectTranslationCS{odložit (co), odročit (co)}
\dicEntry[freta] \dicTerm{fret|a} \dicIPA{{f}{r}{\textepsilon}{\textlengthmark}{\textsubring{d}}{a}} \dicPos{v}[1] \dicFlx{(‑aði)}[22] \dicDirectTranslationCS{(u)prdnout si, prdět, pšouknout si}
\dicEntry[freyða] \dicTerm{frey|ða} \dicIPA{{f}{r}{ei}{\textlengthmark}{ð}{a}} \dicPos{v}[2] \dicFlx{(‑ddi, ‑tt)}[180] \dicDirectTranslationCS{(z)pěnit, perlit, bublat} \dicExampleIS{Sápan freyðir.} \dicExampleCS{Mýdlo pění.}
\dicEntry[freyðandi] \dicTerm{freyð··andi} \dicIPA{{f}{r}{ei}{\textlengthmark}{ð}{a}{n}{\textsubring{d}}{\textsci}} \dicPos{adj}[13] \dicFlx{indecl}[1] \dicDirectTranslationCS{pěnivý, šumivý, zpěněný (nápoj ap.)}
\dicEntry[Freyja] \dicTerm{Freyj|a} \dicIPA{{f}{r}{ei}{j}{\textlengthmark}{a}} \dicPos{f}[1] \dicFlx{(‑u)}[6] \dicFlx{prop} \dicFieldCat{myt.} \dicDirectTranslationCS{Freya} \dicIndirectTranslationCS{(bohyně lásky a~krásy)}
\dicEntry[freyjubrá] \dicTerm{freyju··brá} \dicIPA{{f}{r}{ei}{j}{\textlengthmark}{\textscy}{\textsubring{b}}{r}{au}} \dicPos{f}[4] \dicFlx{(‑r, ‑r)}[18] \dicFieldCat{bot.} \dicDirectTranslationCS{kopretina, kopretina bílá} \textit{(l.~{\textLA{Leucanthemum vulgare}})}  \dicsymPhoto\ 
\dicFigure{ds_image_freyjubra_0_1.jpg}{Freyjubrá}{Freyjubrá - Kaldari, PD}
\dicEntry[frétt] \dicTerm{frétt} \dicsymFrequent\  \dicIPA{{f}{r}{j}{\textepsilon}{h}{\textsubring{d}}} \dicPos{f}[7] \dicFlx{(‑ar, ‑ir)}[1] \textbf{1.} \dicSynonym{fregn} \dicDirectTranslationCS{zpráva, novina} \dicExampleIS{Fyrsta frétt var viðtal við konu.} \dicExampleCS{První zpráva byl rozhovor s~ženou.};  \dicPhraseIS{frétt \,/\addthin fréttir um e‑ð} \dicDirectTranslationCS{zpráva\,/\addthin zprávy o~(čem)};  \dicPhraseIS{Hvað er í fréttum?} \dicDirectTranslationCS{Co je nového?};  \dicPhraseIS{leita frétta} \dicDirectTranslationCS{vyhledávat informace}  \textbf{2.} \dicPhraseIS{fréttir} \dicFlx{pl} \dicDirectTranslationCS{zprávy, zpravodajství}
\dicEntry[frétta] \dicTerm{frétt|a} \dicsymFrequent\  \dicIPA{{f}{r}{j}{\textepsilon}{h}{\textsubring{d}}{a}} \dicPos{v}[2] \dicFlx{(‑i, ‑)}[11] \dicFlx{acc} \dicDirectTranslationCS{dozvědět se, zaslechnout (novinu ap.)} \dicExampleIS{Ég var að frétta að þú værir skilinn.} \dicExampleCS{Dozvěděl jsem se, že jsi rozvedený.};  \dicIdiom{fréttast}{ \dicPhraseIS{fréttast}} \dicFlx{refl} \dicDirectTranslationCS{roznést se, rozšířit se (zpráva ap.)}; { \dicPhraseIS{það fréttist af e‑m}} \dicFlx{refl impers} \dicDirectTranslationCS{je slyšet o~(kom)}
\dicEntry[Fréttablaðið] \dicTerm{Frétta··blað|ið} \dicIPA{{f}{r}{j}{\textepsilon}{h}{\textsubring{d}}{a}{\textsubring{b}}{l}{a}{ð}{\textsci}{\texttheta}} \dicPos{n}[2] \dicFlx{(‑sins)}[3] \dicFlx{def} \dicDirectTranslationCS{Fréttablaðið} \dicIndirectTranslationCS{(islandský deník)}
\dicEntry[fréttabréf] \dicTerm{frétta··bréf} \dicIPA{{f}{r}{j}{\textepsilon}{h}{\textsubring{d}}{a}{\textsubring{b}}{r}{j}{\textepsilon}{f}} \dicPos{n}[2] \dicFlx{(‑s, ‑)}[5] \dicDirectTranslationCS{věstník, informační bulletin, newsletter}
\dicEntry[fréttakvikmynd] \dicTerm{frétta··kvik·mynd} \dicIPA{{f}{r}{j}{\textepsilon}{h}{\textsubring{d}}{a}{k\smash{\textsuperscript{h}}}{v}{\textsci}{\r{g}}{m}{\textsci}{n}{\textsubring{d}}} \dicPos{f}[7] \dicFlx{(‑ar, ‑ir)}[1] \dicDirectTranslationCS{zpravodajský\,/\addthin dokumentární film}
\dicEntry[fréttaljósmynd] \dicTerm{frétta··ljós·mynd} \dicIPA{{f}{r}{j}{\textepsilon}{h}{\textsubring{d}}{a}{l}{j}{ou}{s}{m}{\textsci}{n}{\textsubring{d}}} \dicPos{f}[7] \dicFlx{(‑ar, ‑ir)}[1] \dicDirectTranslationCS{žurnalistická fotografie}
\dicEntry[fréttamaður] \dicTerm{frétta··|maður} \dicIPA{{f}{r}{j}{\textepsilon}{h}{\textsubring{d}}{a}{m}{a}{ð}{\textscy}{\textsubring{r}}} \dicPos{m}[13] \dicFlx{(‑manns, ‑menn)}[2] \dicDirectTranslationCS{novinář(ka), reportér(ka)}
\dicEntry[fréttamannafundur] \dicTerm{frétta·manna··fund|ur} \dicIPA{{f}\-{r}\-{j}\-{\textepsilon}\-{h}\-{\textsubring{d}}\-{a}\-{m}\-{a}\-{n}\-{a}\-{f}\-{\textscy}\-{n}\-{\textsubring{d}}\-{\textscy}\-{\textsubring{r}}\-} \dicPos{m}[10] \dicFlx{(‑ar, ‑ir)}[4] \dicDirectTranslationCS{tisková konference}
\dicEntry[fréttamennska] \dicTerm{frétta··mennsk|a} \dicIPA{{f}{r}{j}{\textepsilon}{h}{\textsubring{d}}{a}{m}{\textepsilon}{n}{s}{\r{g}}{a}} \dicPos{f}[1] \dicFlx{(‑u)}[5] \dicDirectTranslationCS{novinářství, žurnalistika}
\dicEntry[fréttaritari] \dicTerm{frétta··rit·ar|i} \dicIPA{{f}{r}{j}{\textepsilon}{h}{\textsubring{d}}{a}{r}{\textsci}{\textsubring{d}}{a}{r}{\textsci}} \dicPos{m}[1] \dicFlx{(‑a, ‑ar)}[13] \dicDirectTranslationCS{zpravodaj(ka), dopisovatel(ka), korespondent(ka)}
\dicEntry[fréttastofa] \dicTerm{frétta··stof|a} \dicIPA{{f}{r}{j}{\textepsilon}{h}{\textsubring{d}}{a}{s}{\textsubring{d}}{\textopeno}{v}{a}} \dicPos{f}[1] \dicFlx{(‑u, ‑ur)}[7] \dicDirectTranslationCS{redakce, tisková kancelář, zpravodajská místnost}
\dicEntry[fréttatilkynning] \dicTerm{frétta··til·kynn·ing} \dicIPA{{f}{r}{j}{\textepsilon}{h}{\textsubring{d}}{a}{t\smash{\textsuperscript{h}}}{\textsci}{\textsubring{l}}{c\smash{\textsuperscript{h}}}{\textsci}{n}{i}{\ng}{\r{g}}} \dicPos{f}[4] \dicFlx{(‑ar, ‑ar)}[5] \dicDirectTranslationCS{tisková zpráva, tiskové oznámení, komuniké}
\dicEntry[fréttaþáttur] \dicTerm{frétta··|þáttur} \dicIPA{{f}{r}{j}{\textepsilon}{h}{\textsubring{d}}{a}{\texttheta}{au}{h}{\textsubring{d}}{\textscy}{\textsubring{r}}} \dicPos{m}[12] \dicFlx{(‑þáttar, ‑þættir)}[7] \dicDirectTranslationCS{zpravodajský pořad}
\dicEntry[friða] \dicTerm{frið|a} \dicIPA{{f}{r}{\textsci}{\textlengthmark}{ð}{a}} \dicPos{v}[1] \dicFlx{(‑aði)}[1] \dicFlx{acc} \textbf{1.} \dicSynonym{róa\smash{\textsuperscript{2}}} \dicDirectTranslationCS{uklidnit, uklidňovat, uspokojit, uspokojovat, utišit, utišovat}  \textbf{2.} \dicSynonym*{friðlýsa} \dicDirectTranslationCS{(o)chránit, ochraňovat} \dicExampleIS{friða hvali} \dicExampleCS{chránit velryby}
\dicEntry[friðaður] \dicTerm{frið··|aður} \dicIPA{{f}{r}{\textsci}{\textlengthmark}{ð}{a}{ð}{\textscy}{\textsubring{r}}} \dicPos{adj}[3] \dicFlx{(f ‑uð)}[3] \dicDirectTranslationCS{chráněný}
\dicEntry[friðargæslulið] \dicTerm{friðar··gæslu·lið} \dicIPA{{f}{r}{\textsci}{ð}{a}{r}{\r{\textObardotlessj}}{a}{i}{s}{\textsubring{d}}{l}{\textscy}{l}{\textsci}{\texttheta}} \dicPos{n}[2] \dicFlx{(‑s, ‑)}[5] \dicFieldCat{voj.} \dicDirectTranslationCS{mírová jednotka}
\dicEntry[friðarsamningur] \dicTerm{friðar··samn·ing|ur} \dicIPA{{f}{r}{\textsci}{ð}{a}{\textsubring{r}}{s}{a}{m}{n}{i}{\ng}{\r{g}}{\textscy}{\textsubring{r}}} \dicPos{m}[6] \dicFlx{(‑s, ‑ar)}[8] \dicDirectTranslationCS{mírová smlouva\,/\addthin dohoda}
\dicEntry[friðarsinni] \dicTerm{friðar··sinn|i} \dicIPA{{f}{r}{\textsci}{ð}{a}{\textsubring{r}}{s}{\textsci}{n}{\textsci}} \dicPos{m}[1] \dicFlx{(‑a, ‑ar)}[1] \dicDirectTranslationCS{pacifista, pacifistka, mírový aktivista, mírová aktivistka}
\dicEntry[friðarverðlaun] \dicTerm{friðar··verð·laun} \dicIPA{{f}{r}{\textsci}{ð}{a}{r}{v}{\textepsilon}{r}{ð}{l}{\oe i}{\textsubring{n}}} \dicPos{n}[2] \dicFlx{pl}[1] \dicDirectTranslationCS{cena míru, cena za mír};  \dicPhraseIS{friðarverðlaun Nóbels} \dicDirectTranslationCS{Nobelova cena míru}
\dicEntry[friðhelgi] \dicTerm{frið··helg|i} \dicIPA{{f}{r}{\textsci}{\textlengthmark}{\texttheta}{h}{\textepsilon}{l}{\r{\textObardotlessj}}{\textsci}} \dicPos{f}[3] \dicFlx{(‑i)}[3] \dicDirectTranslationCS{integrita, nedotknutelnost}
\dicEntry[friðhelgur] \dicTerm{frið··helgur} \dicIPA{{f}{r}{\textsci}{\textlengthmark}{\texttheta}{h}{\textepsilon}{l}{\r{g}}{\textscy}{\textsubring{r}}} \dicPos{adj}[1]\dicFlx{}[-1] \dicDirectTranslationCS{nedotknutelný, neporušitelný}
\dicEntry[friðill] \dicTerm{frið|ill} \dicIPA{{f}{r}{\textsci}{\textlengthmark}{ð}{\textsci}{\textsubring{d}}{\textsubring{l}}} \dicPos{m}[6] \dicFlx{(‑ils, ‑lar)}[35] \dicDirectTranslationCS{milenec (vdané ženy)}
\dicEntry[friðland] \dicTerm{frið··|land} \dicIPA{{f}{r}{\textsci}{ð}{l}{a}{n}{\textsubring{d}}} \dicPos{n}[2] \dicFlx{(‑lands, ‑lönd)}[8] \dicDirectTranslationCS{(přírodní) rezervace, chráněné území}
\dicEntry[friðlaus] \dicTerm{frið··laus} \dicIPA{{f}{r}{\textsci}{ð}{l}{\oe i}{s}} \dicPos{adj}[5]\dicFlx{}[-1] \dicDirectTranslationCS{neklidný, nepokojný};  \dicPhraseIS{vera friðlaus af e‑u} \dicDirectTranslationCS{být plný (čeho) (očekávání ap.)} \dicExampleIS{vera friðlaus af spenningi} \dicExampleCS{být plný napětí}
\dicEntry[friðsamur] \dicTerm{frið··|samur} \dicIPA{{f}{r}{\textsci}{ð}{s}{a}{m}{\textscy}{\textsubring{r}}} \dicPos{adj}[1] \dicFlx{(f ‑söm)}[2] \dicSynonym{spakur} \dicDirectTranslationCS{(po)klidný, pokojný}
\dicEntry[friðsemd] \dicTerm{frið··semd} \dicIPA{{f}{r}{\textsci}{ð}{s}{\textepsilon}{m}{\textsubring{d}}} \dicPos{f}[7] \dicFlx{(‑ar)}[3] \dicDirectTranslationCS{poklidnost, pokojnost}
\dicEntry[friðsæld] \dicTerm{frið··sæld} \dicIPA{{f}{r}{\textsci}{ð}{s}{a}{i}{l}{\textsubring{d}}} \dicPos{f}[7] \dicFlx{(‑ar)}[3] \dicDirectTranslationCS{pokoj, poklid}
\dicEntry[friðsæll] \dicTerm{frið··sæll} \dicIPA{{f}{r}{\textsci}{ð}{s}{a}{i}{\textsubring{d}}{\textsubring{l}}} \dicPos{adj}[8]\dicFlx{}[-1] \dicSynonym{kyrr} \dicDirectTranslationCS{(po)klidný, pokojný, tichý} \dicExampleIS{friðsælt umhverfi} \dicExampleCS{poklidné prostředí}
\dicEntry[friðun] \dicTerm{frið|un} \dicIPA{{f}{r}{\textsci}{\textlengthmark}{ð}{\textscy}{\textsubring{n}}} \dicPos{f}[7] \dicFlx{(‑unar)}[9] \dicDirectTranslationCS{ochrana (životního prostředí ap.)}
\dicEntry[friður] \dicTerm{frið|ur} \dicsymFrequent\  \dicIPA{{f}{r}{\textsci}{\textlengthmark}{ð}{\textscy}{\textsubring{r}}} \dicPos{m}[10] \dicFlx{(‑ar)}[7] \textbf{1.} \dicSynonym*{friðarástand} \dicDirectTranslationCS{mír};  \dicPhraseIS{semja frið við e‑n} \dicDirectTranslationCS{uzavřít mír s~(kým)};  \dicPhraseIS{stilla til friðar} \dicDirectTranslationCS{nastolit mír}  \textbf{2.} \dicSynonym{ró\smash{\textsuperscript{2}}} \dicDirectTranslationCS{klid, pokoj};  \dicPhraseIS{láta e‑n í friði} \dicDirectTranslationCS{nechat (koho) na pokoji}
\dicEntry[friðþæging] \dicTerm{frið··þæg·ing} \dicIPA{{f}{r}{\textsci}{ð}{\texttheta}{a}{i}{j}{i}{\ng}{\r{g}}} \dicPos{f}[4] \dicFlx{(‑ar)}[7] \dicDirectTranslationCS{(u)smíření}
\dicEntry[friggjargras] \dicTerm{friggjar··|gras} \dicIPA{{f}{r}{\textsci}{\r{\textObardotlessj}}{a}{r}{\r{g}}{r}{a}{s}} \dicPos{n}[2] \dicFlx{(‑grass, ‑grös)}[8] \dicFieldCat{bot.} \dicDirectTranslationCS{vemeník, vemeník severní} \textit{(l.~{\textLA{Platanthera hyperborea}})}  \dicsymPhoto\ 
\dicFigure{ds_image_friggjargras_0_2.jpg}{Friggjargras}{Friggjargras - Stickpen, CC BY 3.0 US}
\dicEntry[frilla] \dicTerm{frill|a} \dicIPA{{f}{r}{\textsci}{\textsubring{d}}{l}{a}} \dicPos{f}[1] \dicFlx{(‑u, ‑ur)}[19] \dicLangCat{han.} \dicDirectTranslationCS{milenka, milá (ženatého muže)}
\dicEntry[frí] \dicTerm{frí} \dicsymFrequent\  \dicIPA{{f}{r}{i}{\textlengthmark}} \dicPos{n}[2] \dicFlx{(‑s, ‑)}[5] \dicDirectTranslationCS{prázdniny, volno, dovolená};  \dicPhraseIS{eiga frí} \dicDirectTranslationCS{mít prázdniny\,/\addthin dovolenou};  \dicPhraseIS{fara í frí} \dicDirectTranslationCS{jet na prázdniny\,/\addthin dovolenou};  \dicPhraseIS{taka sér frí} \dicDirectTranslationCS{vzít si volno\,/\addthin dovolenou};  \dicPhraseIS{vera í fríi} \dicDirectTranslationCS{být na prázdninách\,/\addthin dovolené}
\dicEntry[frídagur] \dicTerm{frí··dag|ur} \dicIPA{{f}{r}{i}{\textlengthmark}{\textsubring{d}}{a}{\textbabygamma}{\textscy}{\textsubring{r}}} \dicPos{m}[6] \dicFlx{(‑s, ‑ar)}[62] \dicDirectTranslationCS{(státní) svátek, volný den, volno}
\dicEntry[fríðindi] \dicTerm{fríð··indi} \dicIPA{{f}{r}{i}{\textlengthmark}{ð}{\textsci}{n}{\textsubring{d}}{\textsci}} \dicPos{n}[2] \dicFlx{pl}[19] \dicDirectTranslationCS{výhody, bonusy, benefity (zaměstnanecké ap.)}
\dicEntry[fríðleiki] \dicTerm{fríð··leik|i} \dicIPA{{f}{r}{i}{ð}{l}{ei}{\r{\textObardotlessj}}{\textsci}} \dicPos{m}[1] \dicFlx{(‑a)}[3] \dicDirectTranslationCS{krása, půvab, líbeznost}
\dicEntry[fríður] \dicTerm{fríður} \dicsymFrequent\  \dicIPA{{f}{r}{i}{\textlengthmark}{ð}{\textscy}{\textsubring{r}}} \dicPos{adj}[2]\dicFlx{}[-6] \dicDirectTranslationCS{pěkný, půvabný, pohledný} \dicExampleIS{fríð stúlka} \dicExampleCS{půvabná dívka};  \dicPhraseIS{vera fríður sýnum} \dicDirectTranslationCS{být pohledný}
\dicEntry[fríhendis] \dicTerm{frí··hendis} \dicIPA{{f}{r}{i}{\textlengthmark}{h}{\textepsilon}{n}{\textsubring{d}}{\textsci}{s}} \dicPos{adv} \dicDirectTranslationCS{od ruky} \dicExampleIS{teikna myndina fríhendis} \dicExampleCS{nakreslit obrázek od ruky}
\dicEntry[fríhyggja] \dicTerm{frí··hyggj|a} \dicIPA{{f}{r}{i}{\textlengthmark}{h}{\textsci}{\r{\textObardotlessj}}{a}} \dicPos{f}[1] \dicFlx{(‑u)}[5] \dicFieldCat{náb.} \dicDirectTranslationCS{ateismus, volnomyšlenkářství}
\dicEntry[fríhöfn] \dicTerm{frí··|höfn} \dicIPA{{f}{r}{i}{\textlengthmark}{h}{\oe}{\textsubring{b}}{\textsubring{n}}} \dicPos{f}[7] \dicFlx{(‑hafnar, ‑hafnir)}[16] \textbf{1.} \dicDirectTranslationCS{svobodný přístav}  \textbf{2.} \dicDirectTranslationCS{bezcelní obchod} \dicIndirectTranslationCS{(obchod osvobozený od cla)}
\dicEntry[frík] \dicTerm{frík} \dicIPA{{f}{r}{i}{\textlengthmark}{\r{g}}} \dicPos{n}[2] \dicFlx{(‑s, ‑)}[5] \dicLangCat{hovor.} \dicDirectTranslationCS{podivín(ka), exot(ka), zjev}
\dicEntry[fríkirkja] \dicTerm{frí··kirkj|a} \dicIPA{{f}{r}{i}{\textlengthmark}{c\smash{\textsuperscript{h}}}{\textsci}{\textsubring{r}}{\r{\textObardotlessj}}{a}} \dicPos{f}[1] \dicFlx{(‑u, ‑ur)}[17] \dicFieldCat{náb.} \dicDirectTranslationCS{svobodná církev} \dicIndirectTranslationCS{(církev na Islandu, která se oddělila od státní luteránské církve)}
\dicEntry[fríkka] \dicTerm{fríkk|a} \dicIPA{{f}{r}{i}{h}{\r{g}}{a}} \dicPos{v}[1] \dicFlx{(‑aði)}[34] \dicDirectTranslationCS{zkrásnět, stát se hezčím\,/\addthin krásnějším}
\dicEntry[frílisti] \dicTerm{frí··list|i} \dicIPA{{f}{r}{i}{\textlengthmark}{l}{\textsci}{s}{\textsubring{d}}{\textsci}} \dicPos{m}[1] \dicFlx{(‑a, ‑ar)}[1] \dicFieldCat{hist.} \dicDirectTranslationCS{bezcelní listina}
\dicEntry[frímerki] \dicTerm{frí··merki} \dicIPA{{f}{r}{i}{\textlengthmark}{m}{\textepsilon}{\textsubring{r}}{\r{\textObardotlessj}}{\textsci}} \dicPos{n}[2] \dicFlx{(‑s, ‑)}[16] \dicDirectTranslationCS{(poštovní) známka} \dicExampleIS{frímerkjabók} \dicExampleCS{album na známky}
\dicEntry[frímerkja] \dicTerm{frí··merk|ja} \dicIPA{{f}{r}{i}{\textlengthmark}{m}{\textepsilon}{\textsubring{r}}{\r{\textObardotlessj}}{a}} \dicPos{v}[2] \dicFlx{(‑ti, ‑t)}[25] \dicFlx{acc} \dicDirectTranslationCS{(o)razítkovat, (o)frankovat}
\dicEntry[frímerkjasafn] \dicTerm{frí·merkja··|safn} \dicIPA{{f}{r}{i}{\textlengthmark}{m}{\textepsilon}{\textsubring{r}}{\r{\textObardotlessj}}{a}{s}{a}{\textsubring{b}}{\textsubring{n}}} \dicPos{n}[2] \dicFlx{(‑safns, ‑söfn)}[8] \dicDirectTranslationCS{sbírka známek}
\dicEntry[frímerkjasafnari] \dicTerm{frí·merkja··safn·ar|i} \dicIPA{{f}\-{r}\-{i}\-{\textlengthmark}\-{m}\-{\textepsilon}\-{\textsubring{r}}\-{\r{\textObardotlessj}}\-{a}\-{s}\-{a}\-{\textsubring{b}}\-{n}\-{a}\-{r}\-{\textsci}\-} \dicPos{m}[1] \dicFlx{(‑a, ‑ar)}[10] \dicDirectTranslationCS{filatelista, filatelistka, sběratel(ka) známek}
\dicEntry[frímerkjasöfnun] \dicTerm{frí·merkja··söfn|un} \dicIPA{{f}{r}{i}{\textlengthmark}{m}{\textepsilon}{\textsubring{r}}{\r{\textObardotlessj}}{a}{s}{\oe}{\textsubring{b}}{n}{\textscy}{\textsubring{n}}} \dicPos{f}[7] \dicFlx{(‑unar)}[12] \dicDirectTranslationCS{filatelie}
\dicEntry[frímínútur] \dicTerm{frí··mínútur} \dicIPA{{f}{r}{i}{\textlengthmark}{m}{i}{n}{u}{\textsubring{d}}{\textscy}{\textsubring{r}}} \dicPos{f}[1] \dicFlx{pl}[16] \dicFieldCat{škol.} \dicDirectTranslationCS{přestávka (ve škole ap.)}
\dicEntry[frímúrari] \dicTerm{frí··múr·ar|i} \dicIPA{{f}{r}{i}{\textlengthmark}{m}{u}{r}{a}{r}{\textsci}} \dicPos{m}[1] \dicFlx{(‑a, ‑ar)}[13] \dicDirectTranslationCS{svobodný zednář}
\dicEntry[frír] \dicTerm{frír} \dicIPA{{f}{r}{i}{\textlengthmark}{\textsubring{r}}} \dicPos{adj}[4]\dicFlx{}[-1] \textbf{1.} \dicSynonym{ókeypis\smash{\textsuperscript{1}}} \dicDirectTranslationCS{(jsoucí) zdarma\,/\addthin gratis}  \textbf{2.} \dicSynonym{frjáls} \dicDirectTranslationCS{volný, svobodný}
\dicEntry[frískur] \dicTerm{frískur} \dicIPA{{f}{r}{i}{s}{\r{g}}{\textscy}{\textsubring{r}}} \dicPos{adj}[1]\dicFlx{}[-1] \textbf{1.} \dicSynonym{fjörugur} \dicDirectTranslationCS{živý, čilý}  \textbf{2.} \dicSynonym{heilbrigður} \dicDirectTranslationCS{zdravý, (jsoucí) fit}
\dicEntry[frísneska] \dicTerm{frís··nesk|a} \dicIPA{{f}{r}{i}{s}{\textsubring{d}}{n}{\textepsilon}{s}{\r{g}}{a}} \dicPos{f}[1] \dicFlx{(‑u)}[5] \dicDirectTranslationCS{fríština}
\dicEntry[frísneskur] \dicTerm{frís··neskur} \dicIPA{{f}{r}{i}{s}{\textsubring{d}}{n}{\textepsilon}{s}{\r{g}}{\textscy}{\textsubring{r}}} \dicPos{adj}[1]\dicFlx{}[-1] \dicDirectTranslationCS{fríský}
\dicEntry[fríspark] \dicTerm{frí··|spark} \dicIPA{{f}{r}{i}{\textlengthmark}{s}{\textsubring{b}}{a}{\textsubring{r}}{\r{g}}} \dicPos{n}[2] \dicFlx{(‑sparks, ‑spörk)}[8] \dicFieldCat{sport.} \dicDirectTranslationCS{volný kop (ve fotbale)}
\dicEntry[frístund] \dicTerm{frí··stund} \dicIPA{{f}{r}{i}{\textlengthmark}{s}{\textsubring{d}}{\textscy}{n}{\textsubring{d}}} \dicPos{f}[7] \dicFlx{(‑ar, ‑ir)}[1] \dicPhraseIS{frístundir} \dicFlx{pl} \dicDirectTranslationCS{volný čas, volno} \dicExampleIS{í frístundum} \dicExampleCS{ve volném čase}
\dicEntry[fríverslun] \dicTerm{frí··versl|un} \dicIPA{{f}{r}{i}{\textlengthmark}{v}{\textepsilon}{\textsubring{r}}{s}{\textsubring{d}}{l}{\textscy}{\textsubring{n}}} \dicPos{f}[7] \dicFlx{(‑unar)}[9] \dicFieldCat{ekon.} \dicDirectTranslationCS{volný obchod}
\dicEntry[frjáls] \dicTerm{frjáls} \dicsymFrequent\  \dicIPA{{f}{r}{j}{au}{l}{s}} \dicPos{adj}[5] \dicFlx{(f frjáls)}[1] \textbf{1.} \dicSynonym{laus} \dicDirectTranslationCS{svobodný, volný} \dicExampleIS{Ég var frjáls eins og fuglinn.} \dicExampleCS{Byl jsem volný jako pták.};  \dicPhraseIS{vera frjáls að því að (gera e‑ð)} \dicDirectTranslationCS{smět (dělat (co)) bez obav}  \textbf{2.} \dicDirectTranslationCS{volný, (jsoucí) bez zábran};  \dicPhraseIS{frjálsar ástir} \dicDirectTranslationCS{volná láska};  \dicPhraseIS{frjáls samkeppni} \dicDirectTranslationCS{volná soutěž};  \dicPhraseIS{frjálsar (íþróttir)} \dicFieldCat{sport.} \dicDirectTranslationCS{atletika};  \dicPhraseIS{hafa frjálsar hendur} \dicLangCat{přen.} \dicDirectTranslationCS{mít volné ruce}
\dicEntry[frjálshyggja] \dicTerm{frjáls··hyggj|a} \dicIPA{{f}{r}{j}{au}{l}{s}{h}{\textsci}{\r{\textObardotlessj}}{a}} \dicPos{f}[1] \dicFlx{(‑u)}[5] \textbf{1.} \dicFieldCat{ekon.} \dicDirectTranslationCS{liberalismus}  \textbf{2.} \dicFieldCat{náb.} \dicSynonym{fríhyggja} \dicDirectTranslationCS{svobodomyslnost, volnomyšlenkářství}
\begin{xtolerant}{}{1pt}
\dicEntry[frjálsíþróttamaður] \dicTerm{frjáls·í·þrótta··|maður}\addthinSS\textls[-13]{\dicIPA{{f}\-{r}\-{j}\-{au}\-{l}\-{s}\-{i}\-{\texttheta}\-{r}\-{ou}\-{h}\-{\textsubring{d}}\-{a}\-{m}\-{a}\-{ð}\-{\textscy}\-{\textsubring{r}}}} \dicPos{m}[1\hspace{-0.06667em}3]\addthinSS\dicFlx{(‑manns, ‑menn)}[2] \dicDirectTranslationCS{atlet(ka)}
\end{xtolerant}
\dicEntry[frjálsíþróttamót] \dicTerm{frjáls·í·þrótta··mót} \dicIPA{{f}{r}{j}{au}{l}{s}{i}{\texttheta}{r}{ou}{h}{\textsubring{d}}{a}{m}{ou}{\textsubring{d}}} \dicPos{n}[2] \dicFlx{(‑s, ‑)}[5] \dicDirectTranslationCS{atletické závody}
\dicEntry[frjálsíþróttir] \dicTerm{frjáls··í·þróttir} \dicIPA{{f}{r}{j}{au}{l}{s}{i}{\texttheta}{r}{ou}{h}{\textsubring{d}}{\textsci}{\textsubring{r}}} \dicPos{f}[7] \dicFlx{pl}[2] \dicFieldCat{sport.} \dicDirectTranslationCS{atletika}
\dicEntry[frjálslegur] \dicTerm{frjáls··legur} \dicIPA{{f}{r}{j}{au}{l}{s}{l}{\textepsilon}{\textbabygamma}{\textscy}{\textsubring{r}}} \dicPos{adj}[1]\dicFlx{}[-8] \textbf{1.} \dicSynonym{óþvingaður} \dicDirectTranslationCS{volný, přirozený, uvolněný}  \textbf{2.} \dicSynonym*{laus í forminu} \dicDirectTranslationCS{volný, nekonvenční}
\dicEntry[frjálslyndi] \dicTerm{frjáls··lyndi} \dicIPA{{f}{r}{j}{au}{l}{s}{l}{\textsci}{n}{\textsubring{d}}{\textsci}} \dicPos{n}[2] \dicFlx{(‑s)}[20] \dicDirectTranslationCS{liberalismus, svobodomyslnost}
\dicEntry[frjálslyndur] \dicTerm{frjáls··lyndur} \dicsymFrequent\  \dicIPA{{f}{r}{j}{au}{l}{s}{l}{\textsci}{n}{\textsubring{d}}{\textscy}{\textsubring{r}}} \dicPos{adj}[2]\dicFlx{}[-14] \textbf{1.} \dicSynonym*{fordómalaus} \dicDirectTranslationCS{tolerantní, nepředpojatý, liberální} \dicExampleIS{vera frjálslyndur í skoðunum} \dicExampleCS{mít liberální názory}  \textbf{2.} \dicDirectTranslationCS{liberální (politická strana ap.)}
\dicEntry[frjálsræði] \dicTerm{frjáls··ræði} \dicIPA{{f}{r}{j}{au}{l}{s}{r}{a}{i}{ð}{\textsci}} \dicPos{n}[2] \dicFlx{(‑s)}[20] \dicSynonym{frelsi} \dicDirectTranslationCS{svoboda (rozhodování ap.)};  \dicPhraseIS{búa við (mikið) frjálsræði} \dicDirectTranslationCS{mít (velkou) svobodu}
\dicEntry[frjó] \dicTerm{frjó} \dicIPA{{f}{r}{j}{ou}{\textlengthmark}} \dicPos{n}[2] \dicFlx{(‑s)}[2] \textbf{1.} \dicFieldCat{biol.} \dicSynonym{sæði} \dicDirectTranslationCS{spermie, semeno}  \textbf{2.} \dicSynonym{fræ} \dicDirectTranslationCS{pyl}
\dicEntry[frjónæmi] \dicTerm{frjó··næm|i\smash{\textsuperscript{1}}} \dicIPA{{f}{r}{j}{ou}{\textlengthmark}{n}{a}{i}{m}{\textsci}} \dicPos{f}[3] \dicFlx{(‑i)}[3] \dicLink{frjónæmi\smash{\textsuperscript{3}}}
\dicEntry[frjónæmi] \dicTerm{frjó··næmi\smash{\textsuperscript{2}}}\dicTerm{, frjónæmi\smash{\textsuperscript{1}}} \dicIPA{{f}\-{r}\-{j}\-{ou}\-{\textlengthmark}\-{n}\-{a}\-{i}\-{m}\-{\textsci}\-} \dicPos{n}[2] \dicFlx{(‑s)}[20] \dicFieldCat{med.} \dicDirectTranslationCS{alergie na pyl}
\dicEntry[frjór] \dicTerm{frjór} \dicIPA{{f}{r}{j}{ou}{\textlengthmark}{\textsubring{r}}} \dicPos{adj}[4]\dicFlx{}[-1] \textbf{1.} \dicDirectTranslationCS{plodný (žena ap.)}  \textbf{2.} \dicDirectTranslationCS{úrodný}  \textbf{3.} \dicSynonym*{hugmyndaríkur} \dicDirectTranslationCS{plodný, vynalézavý (autor ap.)}
\dicEntry[frjósa] \dicTerm{frjósa} \dicIPA{{f}{r}{j}{ou}{\textlengthmark}{s}{a}} \dicPos{v}[6] \dicFlx{(frýs, fraus, frusum, frysi, frosið)}[103] \textbf{1.} \dicDirectTranslationCS{zmrznout, zamrznout, zamrzat} \dicExampleIS{Hvers vegna frýs vatn?} \dicExampleCS{Proč zamrzá voda?};  \dicPhraseIS{það er farið að frjósa} \dicDirectTranslationCS{začíná mrznout}  \textbf{2.} \dicDirectTranslationCS{zamrznout, zaseknout se (o~počítači ap.)}
\dicEntry[frjósamur] \dicTerm{frjó··|samur} \dicIPA{{f}{r}{j}{ou}{\textlengthmark}{s}{a}{m}{\textscy}{\textsubring{r}}} \dicPos{adj}[1] \dicFlx{(f ‑söm)}[2] \textbf{1.} \dicDirectTranslationCS{plodný}  \textbf{2.} \dicDirectTranslationCS{úrodný}
\dicEntry[frjósemi] \dicTerm{frjó··sem|i} \dicIPA{{f}{r}{j}{ou}{\textlengthmark}{s}{\textepsilon}{m}{\textsci}} \dicPos{f}[2] \dicFlx{(‑i)}[2] \dicDirectTranslationCS{plodnost, fertilita}
\dicEntry[frjósemisguð] \dicTerm{frjó·semis··guð} \dicIPA{{f}{r}{j}{ou}{\textlengthmark}{s}{\textepsilon}{m}{\textsci}{s}{\r{g}}{v}{\textscy}{\texttheta}} \dicPos{m}[9] \dicFlx{(‑s, ‑ir)}[18] \dicDirectTranslationCS{bůh plodnosti}
\dicEntry[frjósemisgyðja] \dicTerm{frjó·semis··gyðj|a} \dicIPA{{f}{r}{j}{ou}{\textlengthmark}{s}{\textepsilon}{m}{\textsci}{s}{\r{\textObardotlessj}}{\textsci}{ð}{j}{a}} \dicPos{f}[1] \dicFlx{(‑u, ‑ur)}[7] \dicDirectTranslationCS{bohyně plodnosti}
\dicEntry[frjótæknir] \dicTerm{frjó··tækn|ir} \dicIPA{{f}{r}{j}{ou}{\textlengthmark}{t\smash{\textsuperscript{h}}}{a}{i}{h}{\r{g}}{n}{\textsci}{\textsubring{r}}} \dicPos{m}[7] \dicFlx{(‑is, ‑ar)}[1] \dicDirectTranslationCS{inseminátor}
\dicEntry[frjóvga] \dicTerm{frjóvg|a} \dicIPA{{f}{r}{j}{ou}{v}{\textbabygamma}{a}} \dicPos{v}[1] \dicFlx{(‑aði)}[1] \dicFlx{acc} \dicFieldCat{biol.} \dicDirectTranslationCS{oplodnit};  \dicIdiom{frjóvgast}{ \dicPhraseIS{frjóvgast}} \dicFlx{refl} \dicDirectTranslationCS{oplodnit se}
\dicEntry[frjóvgun] \dicTerm{frjóvg|un} \dicIPA{{f}{r}{j}{ou}{v}{\textbabygamma}{\textscy}{\textsubring{n}}} \dicPos{f}[7] \dicFlx{(‑unar, ‑anir)}[8] \dicFieldCat{biol.} \dicDirectTranslationCS{oplodnění}
\dicEntry[frl.] \dicTerm{frl.} \dicPos{zkr} \dicPhraseIS{frumlag} \dicFieldCat{jaz.} \dicDirectTranslationCS{podmět}
\dicEntry[froða] \dicTerm{froð|a} \dicIPA{{f}{r}{\textopeno}{\textlengthmark}{ð}{a}} \dicPos{f}[1] \dicFlx{(‑u)}[5] \textbf{1.} \dicSynonym{löður} \dicDirectTranslationCS{pěna} \dicExampleIS{bjórfroða} \dicExampleCS{pivní pěna}  \textbf{2.} \dicSynonym{þvaður} \dicDirectTranslationCS{nesmysl, blábol}
\dicEntry[froðufella] \dicTerm{froðu··fell|a} \dicIPA{{f}{r}{\textopeno}{\textlengthmark}{ð}{\textscy}{f}{\textepsilon}{\textsubring{d}}{l}{a}} \dicPos{v}[2] \dicFlx{(‑di, ‑t)}[160] \dicDirectTranslationCS{soptit s~pěnou u~úst}
\dicEntry[frosið] \dicTerm{frosið} \dicIPA{{f}{r}{\textopeno}{\textlengthmark}{s}{\textsci}{\texttheta}} \dicPos{v} \dicFlx{supin} \dicLink{frjósa}
\dicEntry[frosinn] \dicTerm{frosinn} \dicIPA{{f}{r}{\textopeno}{\textlengthmark}{s}{\textsci}{\textsubring{n}}} \dicPos{adj}[6]\dicFlx{}[-6] \textbf{1.} \dicDirectTranslationCS{zmrzlý, zamrzlý} \dicExampleIS{frosin jörð} \dicExampleCS{zmrzlá půda}  \textbf{2.} \dicSynonym{freðinn} \dicDirectTranslationCS{(z)mražený} \dicExampleIS{frosið grænmeti} \dicExampleCS{mražená zelenina}  \textbf{3.} \dicDirectTranslationCS{zkamenělý, (jsoucí) v~šoku}
\begin{xtolerant}{}{1pt}
\dicEntry[froskalappir] \dicTerm{froska··lappir}\addthinS\dicIPA{{f}{r}{\textopeno}{s}{\r{g}}{a}{l}{a}{h}{\textsubring{b}}{\textsci}{\textsubring{r}}}\addthinS\dicPos{f}[7]\addthinS\dicFlx{pl}[18]\addthinS\dicFieldCat{sport.} \dicDirectTranslationCS{plavecké ploutve}
\end{xtolerant}
\dicEntry[froskdýr] \dicTerm{frosk··dýr} \dicIPA{{f}{r}{\textopeno}{s}{\r{g}}{\textsubring{d}}{i}{\textsubring{r}}} \dicPos{n}[2] \dicFlx{(‑s, ‑)}[5] \dicFieldCat{zool.} \dicDirectTranslationCS{obojživelník} \textit{(l.~{\textLA{Amphibia}})}
\dicEntry[froskur] \dicTerm{frosk|ur} \dicIPA{{f}{r}{\textopeno}{s}{\r{g}}{\textscy}{\textsubring{r}}} \dicPos{m}[6] \dicFlx{(‑s, ‑ar)}[10] \dicFieldCat{zool.} \dicDirectTranslationCS{žába} \textit{(l.~{\textLA{Anura}})}
\dicEntry[frost] \dicTerm{frost} \dicsymFrequent\  \dicIPA{{f}{r}{\textopeno}{s}{\textsubring{d}}} \dicPos{n}[2] \dicFlx{(‑s, ‑)}[5] \dicDirectTranslationCS{mráz} \dicExampleIS{Það var 20 stiga frost í gær.} \dicExampleCS{Včera bylo --20 stupňů.}
\dicEntry[frostlögur] \dicTerm{frost··|lögur} \dicIPA{{f}{r}{\textopeno}{s}{\textsubring{d}}{l}{\oe}{\textbabygamma}{\textscy}{\textsubring{r}}} \dicPos{m}[11] \dicFlx{(‑lagar)}[6] \dicDirectTranslationCS{nemrznoucí směs}
\dicEntry[frostmark] \dicTerm{frost··mark} \dicIPA{{f}{r}{\textopeno}{s}{\textsubring{d}}{m}{a}{\textsubring{r}}{\r{g}}} \dicPos{n}[2] \dicFlx{(‑s)}[2] \dicDirectTranslationCS{bod mrazu}
\dicEntry[frostþurrka] \dicTerm{frost··þurrk|a} \dicIPA{{f}{r}{\textopeno}{s}{\textsubring{d}}{\texttheta}{\textscy}{\textsubring{r}}{\r{g}}{a}} \dicPos{v}[1] \dicFlx{(‑aði)}[1] \dicFlx{acc} \dicDirectTranslationCS{sušit na mrazu}
\dicEntry[fró] \dicTerm{fró} \dicIPA{{f}{r}{ou}{\textlengthmark}} \dicPos{f}[4] \dicFlx{(‑ar)}[27] \dicSynonym{linun} \dicDirectTranslationCS{úleva, útěcha}
\dicEntry[fróa] \dicTerm{fró|a} \dicIPA{{f}{r}{ou}{\textlengthmark}{a}} \dicPos{v}[1] \dicFlx{(‑aði)}[1] \dicFlx{dat} \textbf{1.} \dicPhraseIS{fróa sér} \dicSynonym{runka} \dicDirectTranslationCS{onanovat, masturbovat, ukájet se}  \textbf{2.} \dicSynonym{hugga} \dicDirectTranslationCS{utěšit, utěšovat, (u)konejšit} \dicExampleIS{Orð hans fróuðu mér.} \dicExampleCS{Jeho slova mě utěšila.}
\dicEntry[fróðlegur] \dicTerm{fróð··legur} \dicsymFrequent\  \dicIPA{{f}{r}{ou}{ð}{l}{\textepsilon}{\textbabygamma}{\textscy}{\textsubring{r}}} \dicPos{adj}[1]\dicFlx{}[-8] \dicDirectTranslationCS{vzdělávací, osvětový, poučný, informativní} \dicExampleIS{fróðlegt erindi} \dicExampleCS{osvětová přednáška}
\dicEntry[fróðleikur] \textls[15]{\dicTerm{fróð··leik|ur} \dicIPA{{f}{r}{ou}{ð}{l}{ei}{\r{g}}{\textscy}{\textsubring{r}}} \dicPos{m}[6] \dicFlx{(‑s)}[17] \dicDirectTranslationCS{znalost, poučení, obe\-zná\-me\-ní, informace} \dicExampleIS{Bókin inniheldur mikinn fróð\-leik.} \dicExampleCS{Kniha obsahuje spoustu informací.};  \dicPhraseIS{til fróðleiks} \dicFlx{adv} \dicDirectTranslationCS{k~obeznámení, pro informaci}}
\dicEntry[fróður] \dicTerm{fróður} \dicIPA{{f}{r}{ou}{\textlengthmark}{ð}{\textscy}{\textsubring{r}}} \dicPos{adj}[2]\dicFlx{}[-6] \dicDirectTranslationCS{znalý, poučený, obeznámený}
\dicEntry[frómur] \textls[15]{\dicTerm{frómur} \dicIPA{{f}{r}{ou}{\textlengthmark}{m}{\textscy}{\textsubring{r}}} \dicPos{adj}[1]\dicFlx{}[-1] \dicDirectTranslationCS{čestný, poctivý, upřímný};}  \dicPhraseIS{frómt frá sagt} \dicFlx{adv} \dicDirectTranslationCS{upřímně řečeno}
\dicEntry[frón] \dicTerm{frón} \dicIPA{{f}{r}{ou}{\textlengthmark}{\textsubring{n}}} \dicPos{n}[2] \dicFlx{(‑s, ‑)}[5] \textbf{1.} \dicSynonym{Ísland} \dicDirectTranslationCS{Island} \dicExampleIS{heima á fróni\,/\addthin Fróni} \dicExampleCS{doma na Islandu}  \textbf{2.} \dicLangCat{básn.} \dicSynonym{land} \dicDirectTranslationCS{zem, kraj}
\dicEntry[frónskur] \dicTerm{frónskur} \dicIPA{{f}{r}{ou}{n}{s}{\r{g}}{\textscy}{\textsubring{r}}} \dicPos{adj}[1]\dicFlx{}[-1] \dicSynonym{íslenskur} \dicDirectTranslationCS{islandský}
\dicEntry[fróun] \dicTerm{fró|un} \dicIPA{{f}{r}{ou}{\textscy}{\textsubring{n}}} \dicPos{f}[7] \dicFlx{(‑unar)}[9] \textbf{1.} \dicDirectTranslationCS{útěcha, utěšení, ukonejšení}  \textbf{2.} \dicDirectTranslationCS{masturbace, onanování}
\dicEntry[frt.] \dicTerm{frt.} \dicPos{zkr} \dicPhraseIS{framtíð} \dicFieldCat{jaz.} \dicDirectTranslationCS{budoucí čas}
\dicEntry[frum] \dicTerm{frum-} \dicIPA{{f}{r}{\textscy}{\textsubring{m}}} \dicPos{predp} \dicDirectTranslationCS{prvo-, původní}
\dicEntry[fruma] \dicTerm{frum|a} \dicIPA{{f}{r}{\textscy}{\textlengthmark}{m}{a}} \dicPos{f}[1] \dicFlx{(‑u, ‑ur)}[19] \dicFieldCat{biol.} \dicDirectTranslationCS{buňka}
\dicEntry[frumatriði] \dicTerm{frum··at·riði} \dicIPA{{f}{r}{\textscy}{\textlengthmark}{m}{a}{\textsubring{d}}{r}{\textsci}{ð}{\textsci}} \dicPos{n}[2] \dicFlx{(‑s, ‑)}[14] \dicDirectTranslationCS{základ, základní bod}
\dicEntry[frumbernska] \dicTerm{frum··bernsk|a} \dicIPA{{f}{r}{\textscy}{m}{\textsubring{b}}{\textepsilon}{r}{n}{s}{\r{g}}{a}} \dicPos{f}[1] \dicFlx{(‑u)}[5] \dicDirectTranslationCS{rané dětství}
\dicEntry[frumbjarga] \dicTerm{frum··bjarga} \dicIPA{{f}{r}{\textscy}{m}{\textsubring{b}}{j}{a}{r}{\r{g}}{a}} \dicPos{adj}[13] \dicFlx{indecl}[1] \dicFieldCat{biol.} \dicDirectTranslationCS{autotrofní} \dicAntonym{ófrumbjarga}
\dicEntry[frumburður] \dicTerm{frum··burð|ur} \dicIPA{{f}{r}{\textscy}{m}{\textsubring{b}}{\textscy}{r}{ð}{\textscy}{\textsubring{r}}} \dicPos{m}[10] \dicFlx{(‑ar, ‑ir)}[4] \dicDirectTranslationCS{prvorozený, prvorozená (dítě ap.)}
\dicEntry[frumbyggi] \dicTerm{frum··bygg|i} \dicIPA{{f}{r}{\textscy}{m}{\textsubring{b}}{\textsci}{\r{\textObardotlessj}}{\textsci}} \dicPos{m}[1] \dicFlx{(‑ja, ‑jar)}[14] \dicDirectTranslationCS{domorodec, domorodka, původní obyvatel(ka)}
\begin{xtolerant}{}{1pt}
\dicEntry[frumdráttur] \dicTerm{frum··|dráttur}\addthinS\dicIPA{{f}{r}{\textscy}{m}{\textsubring{d}}{r}{au}{h}{\textsubring{d}}{\textscy}{\textsubring{r}}}\addthinS\dicPos{m}[12]\addthinS\dicFlx{(‑dráttar, ‑drættir)}[7]\addthinS\dicSynonym{uppkast} \dicDirectTranslationCS{náčrt, návrh, skica};  \dicPhraseIS{frumdrættir að e‑u} \dicDirectTranslationCS{návrhy na (co)}
\end{xtolerant}
\dicEntry[frumdýr] \dicTerm{frum··dýr} \dicIPA{{f}{r}{\textscy}{m}{\textsubring{d}}{i}{\textsubring{r}}} \dicPos{n}[2] \dicFlx{(‑s, ‑)}[5] \dicFieldCat{zool.} \dicDirectTranslationCS{prvok} \textit{(l.~{\textLA{Protozoa}})}
\dicEntry[frumefni] \dicTerm{frum··efni} \dicIPA{{f}{r}{\textscy}{\textlengthmark}{m}{\textepsilon}{\textsubring{b}}{n}{\textsci}} \dicPos{n}[2] \dicFlx{(‑s, ‑)}[14] \dicFieldCat{chem.} \dicDirectTranslationCS{(chemický) prvek}
\dicEntry[frumeind] \dicTerm{frum··eind} \dicIPA{{f}{r}{\textscy}{\textlengthmark}{m}{ei}{n}{\textsubring{d}}} \dicPos{f}[7] \dicFlx{(‑ar, ‑ir)}[1] \dicFieldCat{fyz.} \dicSynonym{atóm} \dicDirectTranslationCS{atom}
\dicEntry[frumeindakjarni] \dicTerm{frum·einda··kjarn|i} \dicIPA{{f}{r}{\textscy}{\textlengthmark}{m}{ei}{n}{\textsubring{d}}{a}{c\smash{\textsuperscript{h}}}{a}{r}{\textsubring{d}}{n}{\textsci}} \dicPos{m}[1] \dicFlx{(‑a, ‑ar)}[8] \dicFieldCat{fyz.} \dicDirectTranslationCS{atomové jádro}
\dicEntry[frumeining] \dicTerm{frum··ein·ing} \dicIPA{{f}{r}{\textscy}{\textlengthmark}{m}{ei}{n}{i}{\ng}{\r{g}}} \dicPos{f}[4] \dicFlx{(‑ar, ‑ar)}[5] \dicDirectTranslationCS{základní jednotka\,/\addthin veličina}
\dicEntry[frumframleiðsla] \dicTerm{frum··fram·leiðsl|a} \dicIPA{{f}{r}{\textscy}{m}{f}{r}{a}{m}{l}{ei}{ð}{s}{\textsubring{d}}{l}{a}} \dicPos{f}[1] \dicFlx{(‑u)}[5] \textbf{1.} \dicFieldCat{ekon.} \dicDirectTranslationCS{prvovýroba}  \textbf{2.} \dicFieldCat{biol.} \dicDirectTranslationCS{primární produkce}
\dicEntry[frumgerð] \dicTerm{frum··gerð} \dicIPA{{f}{r}{\textscy}{m}{\r{\textObardotlessj}}{\textepsilon}{r}{\texttheta}} \dicPos{f}[7] \dicFlx{(‑ar, ‑ir)}[1] \textbf{1.} \dicFieldCat{lit.} \dicDirectTranslationCS{původní text, originál}  \textbf{2.} \dicDirectTranslationCS{prototyp (auta ap.)}
\dicEntry[frumhandrit] \dicTerm{frum··hand·rit} \dicIPA{{f}{r}{\textscy}{\textlengthmark}{m}{h}{a}{n}{\textsubring{d}}{r}{\textsci}{\textsubring{d}}} \dicPos{n}[2] \dicFlx{(‑s, ‑)}[5] \dicDirectTranslationCS{původní rukopis}
\dicEntry[frumheimild] \dicTerm{frum··heimild} \dicIPA{{f}{r}{\textscy}{\textlengthmark}{m}{h}{ei}{m}{\textsci}{l}{\textsubring{d}}} \dicPos{f}[7] \dicFlx{(‑ar, ‑ir)}[1] \dicDirectTranslationCS{původní zdroj\,/\addthin pramen}
\dicEntry[frumherji] \dicTerm{frum··herj|i} \dicIPA{{f}{r}{\textscy}{\textlengthmark}{m}{h}{\textepsilon}{r}{j}{\textsci}} \dicPos{m}[1] \dicFlx{(‑a, ‑ar)}[1] \dicDirectTranslationCS{průkopník, průkopnice, pionýr(ka)}
\dicEntry[frumkvæði] \dicTerm{frum··kvæði} \dicsymFrequent\  \dicIPA{{f}{r}{\textscy}{m}{k\smash{\textsuperscript{h}}}{v}{a}{i}{ð}{\textsci}} \dicPos{n}[2] \dicFlx{(‑s)}[20] \dicDirectTranslationCS{iniciativa};  \dicPhraseIS{eiga frumkvæði\,/\addthin frumkvæðið að e‑u} \dicDirectTranslationCS{iniciovat (co)};  \dicPhraseIS{að eigin frumkvæði} \dicFlx{adv} \dicDirectTranslationCS{z~vlastní iniciativy}
\dicEntry[frumkvöðull] \dicTerm{frum··kvöð|ull} \dicIPA{{f}{r}{\textscy}{m}{k\smash{\textsuperscript{h}}}{v}{\oe}{ð}{\textscy}{\textsubring{d}}{\textsubring{l}}} \dicPos{m}[6] \dicFlx{(‑uls, ‑lar)}[31] \dicDirectTranslationCS{iniciátor(ka), průkopník, průkopnice}
\dicEntry[frumlag] \dicTerm{frum··|lag} \dicIPA{{f}{r}{\textscy}{m}{l}{a}{x}} \dicPos{n}[2] \dicFlx{(‑lags, ‑lög)}[8] \dicFieldCat{jaz.} \dicDirectTranslationCS{podmět} \dicExampleIS{frumlag -- umsögn -- andlag} \dicExampleCS{podmět -- přísudek -- předmět}
\dicEntry[frumlagsígildi] \dicTerm{frum·lags··í·gildi} \dicIPA{{f}{r}{\textscy}{m}{l}{a}{x}{s}{i}{\r{\textObardotlessj}}{\textsci}{l}{\textsubring{d}}{\textsci}} \dicPos{n}[2] \dicFlx{(‑s, ‑)}[14] \dicFieldCat{jaz.} \dicDirectTranslationCS{logický podmět}
\dicEntry[frumlegur] \dicTerm{frum··legur} \dicIPA{{f}{r}{\textscy}{m}{l}{\textepsilon}{\textbabygamma}{\textscy}{\textsubring{r}}} \dicPos{adj}[1]\dicFlx{}[-8] \dicSynonym{nýstárlegur} \dicDirectTranslationCS{původní, originální, prvotní} \dicExampleIS{frumlegt listaverk} \dicExampleCS{původní umělecké dílo}
\dicEntry[frumlyf] \dicTerm{frum··lyf} \dicIPA{{f}{r}{\textscy}{m}{l}{\textsci}{f}} \dicPos{n}[2] \dicFlx{(‑s, ‑)}[13] \dicFieldCat{med.} \dicDirectTranslationCS{originální lék} \dicAntonym{samheitalyf}
\dicEntry[frummál] \dicTerm{frum··mál} \dicIPA{{f}{r}{\textscy}{m}{\textlengthmark}{au}{\textsubring{l}}} \dicPos{n}[2] \dicFlx{(‑s, ‑)}[5] \dicDirectTranslationCS{výchozí\,/\addthin zdrojový jazyk}
\dicEntry[frumminni] \dicTerm{frum··minni} \dicIPA{{f}{r}{\textscy}{m}{\textlengthmark}{\textsci}{n}{\textsci}} \dicPos{n}[2] \dicFlx{(‑s, ‑)}[14] \dicFieldCat{lit.\,/\addthin psych.} \dicDirectTranslationCS{archetyp}
\dicEntry[frummynd] \dicTerm{frum··mynd} \dicIPA{{f}{r}{\textscy}{m}{\textlengthmark}{\textsci}{n}{\textsubring{d}}} \dicPos{f}[7] \dicFlx{(‑ar, ‑ir)}[1] \textbf{1.} \dicDirectTranslationCS{originál} \dicIndirectTranslationCS{(původní dílo)}  \textbf{2.} \dicFieldCat{filos.} \dicDirectTranslationCS{idea}  \textbf{3.} \dicDirectTranslationCS{prototyp}
\dicEntry[frummælandi] \dicTerm{frum··mæl·|andi} \dicIPA{{f}{r}{\textscy}{m}{\textlengthmark}{a}{i}{l}{a}{n}{\textsubring{d}}{\textsci}} \dicPos{m}[2] \dicFlx{(‑anda, ‑endur)}[1] \dicDirectTranslationCS{úvodní řečník\,/\addthin řečnice}
\dicEntry[frumnorræna] \dicTerm{frum··nor·ræn|a} \dicIPA{{f}{r}{\textscy}{m}{n}{\textopeno}{r}{a}{i}{n}{a}} \dicPos{f}[1] \dicFlx{(‑u)}[5] \dicFieldCat{jaz.} \dicDirectTranslationCS{praseverština}
\dicEntry[frumnorrænn] \dicTerm{frum··nor·rænn} \dicIPA{{f}{r}{\textscy}{m}{n}{\textopeno}{r}{a}{i}{\textsubring{d}}{\textsubring{n}}} \dicPos{adj}[7]\dicFlx{}[-2] \dicDirectTranslationCS{praseverský}
\dicEntry[frumrit] \dicTerm{frum··rit} \dicIPA{{f}{r}{\textscy}{m}{r}{\textsci}{\textsubring{d}}} \dicPos{n}[2] \dicFlx{(‑s, ‑)}[5] \dicDirectTranslationCS{originál, původní text}
\dicEntry[frumsamdi] \dicTerm{frum··samdi} \dicIPA{{f}{r}{\textscy}{m}{s}{a}{m}{\textsubring{d}}{\textsci}} \dicPos{v} \dicFlx{ind pf sg 1 pers} \dicLink{frumsemja}
\dicEntry[frumsamið] \dicTerm{frum··samið} \dicIPA{{f}{r}{\textscy}{m}{s}{a}{m}{\textsci}{\texttheta}} \dicPos{v} \dicFlx{supin} \dicLink{frumsemja}
\dicEntry[frumsem] \dicTerm{frum··sem} \dicIPA{{f}{r}{\textscy}{m}{s}{\textepsilon}{\textsubring{m}}} \dicPos{v} \dicFlx{ind praes sg 1 pers} \dicLink{frumsemja}
\dicEntry[frumsemdi] \dicTerm{frum··semdi} \dicIPA{{f}{r}{\textscy}{m}{s}{\textepsilon}{m}{\textsubring{d}}{\textsci}} \dicPos{v} \dicFlx{con pf sg 1 pers} \dicLink{frumsemja}
\dicEntry[frumsemja] \dicTerm{frum··|semja} \dicIPA{{f}{r}{\textscy}{m}{s}{\textepsilon}{m}{j}{a}} \dicPos{v}[4] \dicFlx{(‑sem, ‑samdi, ‑sömdum, ‑semdi, ‑samið)}[6] \dicFlx{acc} \dicDirectTranslationCS{vytvořit původní\,/\addthin originální dílo}
\dicEntry[frumsetning] \dicTerm{frum··set·ning} \dicIPA{{f}{r}{\textscy}{m}{s}{\textepsilon}{h}{\textsubring{d}}{n}{i}{\ng}{\r{g}}} \dicPos{f}[4] \dicFlx{(‑ar, ‑ar)}[5] \dicSynonym{prinsipp} \dicDirectTranslationCS{postulát, axiom}
\dicEntry[frumskilyrði] \dicTerm{frum··skil·yrði} \dicIPA{{f}{r}{\textscy}{m}{s}{\r{\textObardotlessj}}{\textsci}{l}{\textsci}{r}{ð}{\textsci}} \dicPos{n}[2] \dicFlx{(‑s, ‑)}[14] \dicDirectTranslationCS{základní předpoklad\,/\addthin podmínka}
\dicEntry[frumskógur] \dicTerm{frum··skóg|ur} \dicIPA{{f}{r}{\textscy}{m}{s}{\r{g}}{ou}{\textscy}{\textsubring{r}}} \dicPos{m}[6] \dicFlx{(‑ar, ‑ar)}[56] \dicDirectTranslationCS{džungle, prales}
\dicEntry[frumspeki] \dicTerm{frum··spek|i} \dicIPA{{f}{r}{\textscy}{m}{s}{\textsubring{b}}{\textepsilon}{\r{\textObardotlessj}}{\textsci}} \dicPos{f}[3] \dicFlx{(‑i)}[3] \dicFieldCat{filos.} \dicDirectTranslationCS{metafyzika}
\dicEntry[frumstig] \dicTerm{frum··stig} \dicIPA{{f}{r}{\textscy}{m}{s}{\textsubring{d}}{\textsci}{x}} \dicPos{n}[2] \dicFlx{(‑s, ‑)}[5] \textbf{1.} \dicDirectTranslationCS{počáteční fáze, počátek} \dicExampleIS{e‑að er á frumstigi} \dicExampleCS{(co) je v~počáteční fázi}  \textbf{2.} \dicFieldCat{jaz.} \dicDirectTranslationCS{pozitiv, 1. stupeň stupňování}
\dicEntry[frumstæður] \dicTerm{frum··stæður} \dicsymFrequent\  \dicIPA{{f}{r}{\textscy}{m}{s}{\textsubring{d}}{a}{i}{ð}{\textscy}{\textsubring{r}}} \dicPos{adj}[2]\dicFlx{}[-6] \textbf{1.} \dicSynonym{einfaldur} \dicDirectTranslationCS{primitivní, velmi prostý}  \textbf{2.} \dicDirectTranslationCS{primitivní} \dicIndirectTranslationCS{(jsoucí v~prvotním stupni vývoje)} \dicExampleIS{Það er dæmigert fyrir frumstæð þjóðfélög.} \dicExampleCS{To je typické pro primitivní společnosti.}
\dicEntry[frumsýna] \dicTerm{frum··sýn|a} \dicIPA{{f}{r}{\textscy}{m}{s}{i}{n}{a}} \dicPos{v}[2] \dicFlx{(‑di, ‑t)}[155] \dicFlx{acc} \dicDirectTranslationCS{uvést premiéru} \dicExampleIS{frumsýna óperuna} \dicExampleCS{uvést premiéru opery}
\dicEntry[frumsýning] \dicTerm{frum··sýn·ing} \dicIPA{{f}{r}{\textscy}{m}{s}{i}{n}{i}{\ng}{\r{g}}} \dicPos{f}[4] \dicFlx{(‑ar, ‑ar)}[5] \dicDirectTranslationCS{premiéra} \dicExampleIS{frumsýning á e‑u} \dicExampleCS{premiéra (čeho)}
\dicEntry[frumsömdum] \dicTerm{frum··sömdum} \dicIPA{{f}{r}{\textscy}{m}{s}{\oe}{m}{\textsubring{d}}{\textscy}{\textsubring{m}}} \dicPos{v} \dicFlx{ind pf pl 1 pers} \dicLink{frumsemja}
\dicEntry[frumtala] \dicTerm{frum··|tala} \dicIPA{{f}{r}{\textscy}{m}{t\smash{\textsuperscript{h}}}{a}{l}{a}} \dicPos{f}[1] \dicFlx{(‑tölu, ‑tölur)}[14] \textbf{1.} \dicFieldCat{jaz.} \dicSynonym*{grunntala} \dicDirectTranslationCS{základní číslovka}  \textbf{2.} \dicFieldCat{mat.} \dicSynonym{prímtala} \dicDirectTranslationCS{prvočíslo}
\dicEntry[frumtexti] \dicTerm{frum··text|i} \dicIPA{{f}{r}{\textscy}{m}{t\smash{\textsuperscript{h}}}{\textepsilon}{x}{s}{\textsubring{d}}{\textsci}} \dicPos{m}[1] \dicFlx{(‑a, ‑ar)}[1] \dicDirectTranslationCS{původní\,/\addthin originální text, originál}
\dicEntry[frumuskipting] \dicTerm{frumu··skipt·ing} \dicIPA{{f}{r}{\textscy}{\textlengthmark}{m}{\textscy}{s}{\r{\textObardotlessj}}{\textsci}{f}{\textsubring{d}}{i}{\ng}{\r{g}}} \dicPos{f}[4] \dicFlx{(‑ar, ‑ar)}[5] \dicFieldCat{biol.} \dicDirectTranslationCS{buněčné dělení}
\dicEntry[frumútgáfa] \dicTerm{frum··út·gáf|a} \dicIPA{{f}{r}{\textscy}{\textlengthmark}{m}{u}{\textsubring{d}}{\r{g}}{au}{v}{a}} \dicPos{f}[1] \dicFlx{(‑u, ‑ur)}[13] \dicDirectTranslationCS{původní vydání}
\dicEntry[frumvarp] \dicTerm{frum··|varp} \dicsymFrequent\  \dicIPA{{f}{r}{\textscy}{m}{v}{a}{\textsubring{r}}{\textsubring{b}}} \dicPos{n}[2] \dicFlx{(‑varps, ‑vörp)}[8] \dicFieldCat{pol.} \dicDirectTranslationCS{návrh (zákona)} \dicExampleIS{frumvarp að fjárhagsáætlun} \dicExampleCS{návrh na rozpočet}
\dicEntry[frunsa] \dicTerm{fruns|a} \dicIPA{{f}{r}{\textscy}{n}{s}{a}} \dicPos{f}[1] \dicFlx{(‑u, ‑ur)}[7] \dicFieldCat{med.} \dicDirectTranslationCS{opar, herpes}
\dicEntry[frusum] \dicTerm{frusum} \dicIPA{{f}{r}{\textscy}{\textlengthmark}{s}{\textscy}{\textsubring{m}}} \dicPos{v} \dicFlx{ind pf pl 1 pers} \dicLink{frjósa}
\dicEntry[frú] \dicTerm{frú} \dicsymFrequent\  \dicIPA{{f}{r}{u}{\textlengthmark}} \dicPos{f}[4] \dicFlx{(‑ar, ‑r)}[26] \dicSynonym{kona} \dicDirectTranslationCS{paní, madam, dáma} \dicExampleIS{frú forseti} \dicExampleCS{paní prezidentka}
\dicEntry[frúktósi] \dicTerm{frúktós|i} \dicIPA{{f}{r}{u}{x}{\textsubring{d}}{ou}{s}{\textsci}} \dicPos{m}[1] \dicFlx{(‑a)}[3] \dicFieldCat{biol.} \dicDirectTranslationCS{fruktóza}
\dicEntry[frygð] \dicTerm{frygð} \dicIPA{{f}{r}{\textsci}{\textbabygamma}{\texttheta}} \dicPos{f}[4] \dicFlx{(‑ar)}[3] \dicSynonym{losti} \dicDirectTranslationCS{chtíč, choutka}
\dicEntry[frymi] \dicTerm{frymi} \dicIPA{{f}{r}{\textsci}{\textlengthmark}{m}{\textsci}} \dicPos{n}[2] \dicFlx{(‑s, ‑)}[14] \dicFieldCat{biol.} \dicDirectTranslationCS{protoplazma}
\dicEntry[frysi] \dicTerm{frysi} \dicIPA{{f}{r}{\textsci}{\textlengthmark}{s}{\textsci}} \dicPos{v} \dicFlx{con pf sg 1 pers} \dicLink{frjósa}
\dicEntry[frysta] \dicTerm{fryst|a} \dicIPA{{f}{r}{\textsci}{s}{\textsubring{d}}{a}} \dicPos{v}[2] \dicFlx{(‑i, ‑)}[16] \dicFlx{acc} \textbf{1.} \dicDirectTranslationCS{(z)mrazit, zamrazit} \dicExampleIS{frysta fiskinn} \dicExampleCS{zamrazit rybu}  \textbf{2.} \dicDirectTranslationCS{zmrazit} \dicIndirectTranslationCS{(zastavit oběh peněz ap.)}
\dicEntry[frystihólf] \dicTerm{frysti··hólf} \dicIPA{{f}{r}{\textsci}{s}{\textsubring{d}}{\textsci}{h}{ou}{l}{f}} \dicPos{n}[2] \dicFlx{(‑s, ‑)}[5] \dicDirectTranslationCS{mrazák, mrazicí box}
\dicEntry[frystihús] \dicTerm{frysti··hús} \dicsymFrequent\  \dicIPA{{f}{r}{\textsci}{s}{\textsubring{d}}{\textsci}{h}{u}{s}} \dicPos{n}[2] \dicFlx{(‑s, ‑)}[5] \dicDirectTranslationCS{mrazírna} \dicExampleIS{reka eigið frystihús} \dicExampleCS{provozovat vlastní mrazírnu}
\dicEntry[frystikista] \dicTerm{frysti··kist|a} \dicIPA{{f}{r}{\textsci}{s}{\textsubring{d}}{\textsci}{c\smash{\textsuperscript{h}}}{\textsci}{s}{\textsubring{d}}{a}} \dicPos{f}[1] \dicFlx{(‑u, ‑ur)}[19] \dicDirectTranslationCS{mrazicí pult}
\dicEntry[frysting] \dicTerm{fryst··ing} \dicIPA{{f}{r}{\textsci}{s}{\textsubring{d}}{i}{\ng}{\r{g}}} \dicPos{f}[4] \dicFlx{(‑ar, ‑ar)}[5] \textbf{1.} \dicDirectTranslationCS{(z)mrazení (ryb ap.)}  \textbf{2.} \dicDirectTranslationCS{zmrazení (cen ap.)}
\dicEntry[frystir] \dicTerm{fryst|ir} \dicIPA{{f}{r}{\textsci}{s}{\textsubring{d}}{\textsci}{\textsubring{r}}} \dicPos{m}[7] \dicFlx{(‑is, ‑ar)}[1] \dicDirectTranslationCS{mrazák, mrazicí zařízení}
\dicEntry[frystiskápur] \dicTerm{frysti··skáp|ur} \dicIPA{{f}{r}{\textsci}{s}{\textsubring{d}}{\textsci}{s}{\r{g}}{au}{\textsubring{b}}{\textscy}{\textsubring{r}}} \dicPos{m}[6] \dicFlx{(‑s, ‑ar)}[15] \dicDirectTranslationCS{mrazák, mrazicí skříň}
\dicEntry[frýja] \dicTerm{frý|ja} \dicIPA{{f}{r}{i}{j}{\textlengthmark}{a}} \dicPos{v}[3] \dicFlx{(‑jaði)\,/\addthin (‑, ‑ði, ‑juðum, ‑ði, ‑ð)}[10] \dicFlx{dat} \dicPhraseIS{frýja e‑m vits} \dicDirectTranslationCS{pochybovat o~(čí) inteligenci}
\dicEntry[frýnilegur] \dicTerm{frýni··legur} \dicIPA{{f}{r}{i}{\textlengthmark}{n}{\textsci}{l}{\textepsilon}{\textbabygamma}{\textscy}{\textsubring{r}}} \dicPos{adj}[1]\dicFlx{}[-8] \dicSynonym{elskulegur} \dicDirectTranslationCS{přátelský, vlídný} \dicIndirectTranslationCS{(používané zvláště v~záporu)}
\dicEntry[frýs] \dicTerm{frýs} \dicIPA{{f}{r}{i}{\textlengthmark}{s}} \dicPos{v} \dicFlx{ind praes sg 1 pers} \dicLink{frjósa}
\dicEntry[fræ] \dicTerm{fræ} \dicIPA{{f}{r}{a}{i}{\textlengthmark}} \dicPos{n}[2] \dicFlx{(‑s, ‑)}[13] \dicDirectTranslationCS{semeno, zrno} \dicIndirectTranslationCS{(plod rostliny)}
\dicEntry[fræða] \dicTerm{fræ|ða} \dicIPA{{f}{r}{a}{i}{\textlengthmark}{ð}{a}} \dicPos{v}[2] \dicFlx{(‑ddi, ‑tt)}[167] \dicFlx{acc} \dicDirectTranslationCS{vzdělat, vzdělávat, informovat, poučit} \dicExampleIS{fræða e‑n um e‑ð} \dicExampleCS{poučit (koho) o~(čem)};  \dicIdiom{fræðast}{ \dicPhraseIS{fræðast um e‑ð}} \dicFlx{refl} \dicDirectTranslationCS{vzdělat se v~(čem), dozvědět se více o~(čem)}
\dicEntry[fræði] \dicTerm{fræð|i\smash{\textsuperscript{1}}} \dicIPA{{f}{r}{a}{i}{\textlengthmark}{ð}{\textsci}} \dicPos{f}[3] \dicFlx{(‑i)}[3] \dicLink{fræði\smash{\textsuperscript{2}}}
\dicEntry[fræði] \dicTerm{fræði\smash{\textsuperscript{2}}}\dicTerm{, fræði\smash{\textsuperscript{1}}} \dicIPA{{f}\-{r}\-{a}\-{i}\-{\textlengthmark}\-{ð}\-{\textsci}\-} \dicPos{n}[2] \dicFlx{pl}[19] \textbf{1.} \dicFlx{f} \dicDirectTranslationCS{věda, nauka, vědní disciplína} \dicExampleIS{málfræði} \dicExampleCS{mluvnice}  \textbf{2.} \dicFlx{n pl} \dicSynonym{vísindi} \dicDirectTranslationCS{věda, vědění}
\dicEntry[fræðigrein] \dicTerm{fræði··grein} \dicIPA{{f}{r}{a}{i}{\textlengthmark}{ð}{\textsci}{\r{g}}{r}{ei}{\textsubring{n}}} \dicPos{f}[4] \dicFlx{(‑ar, ‑ar)}[1] \dicDirectTranslationCS{vědní obor\,/\addthin disciplína}
\dicEntry[fræðiheiti] \dicTerm{fræði··heiti} \dicIPA{{f}{r}{a}{i}{\textlengthmark}{ð}{\textsci}{h}{ei}{\textsubring{d}}{\textsci}} \dicPos{n}[2] \dicFlx{(‑s, ‑)}[14] \dicDirectTranslationCS{odborný\,/\addthin vědecký název}
\dicEntry[fræðilegur] \dicTerm{fræði··legur} \dicIPA{{f}{r}{a}{i}{\textlengthmark}{ð}{\textsci}{l}{\textepsilon}{\textbabygamma}{\textscy}{\textsubring{r}}} \dicPos{adj}[1]\dicFlx{}[-8] \dicDirectTranslationCS{akademický, teoretický, vzdělávací}
\dicEntry[fræðimaður] \dicTerm{fræði··|maður} \dicsymFrequent\  \dicIPA{{f}{r}{a}{i}{\textlengthmark}{ð}{\textsci}{m}{a}{ð}{\textscy}{\textsubring{r}}} \dicPos{m}[13] \dicFlx{(‑manns, ‑menn)}[2] \dicDirectTranslationCS{učenec, učenkyně, vědecký pracovník, vědecká pracovnice} \dicExampleIS{samkvæmt fræðimönnum} \dicExampleCS{podle vědeckých pracovníků}
\dicEntry[fræðimennska] \dicTerm{fræði··mennsk|a} \dicIPA{{f}{r}{a}{i}{\textlengthmark}{ð}{\textsci}{m}{\textepsilon}{n}{s}{\r{g}}{a}} \dicPos{f}[1] \dicFlx{(‑u)}[5] \dicDirectTranslationCS{učenost}
\dicEntry[fræðiorð] \dicTerm{fræði··orð} \dicIPA{{f}{r}{a}{i}{\textlengthmark}{ð}{\textsci}{\textopeno}{r}{\texttheta}} \dicPos{n}[2] \dicFlx{(‑s, ‑)}[5] \dicDirectTranslationCS{odborný název, termín}
\dicEntry[fræðiorðakerfi] \dicTerm{fræði·orða··kerfi} \dicIPA{{f}{r}{a}{i}{\textlengthmark}{ð}{\textsci}{\textopeno}{r}{ð}{a}{c\smash{\textsuperscript{h}}}{\textepsilon}{r}{v}{\textsci}} \dicPos{n}[2] \dicFlx{(‑s, ‑)}[14] \dicDirectTranslationCS{názvosloví, terminologie, nomenklatura}
\dicEntry[fræðirit] \dicTerm{fræði··rit} \dicIPA{{f}{r}{a}{i}{\textlengthmark}{ð}{\textsci}{r}{\textsci}{\textsubring{d}}} \dicPos{n}[2] \dicFlx{(‑s, ‑)}[5] \dicDirectTranslationCS{odborná práce (kniha ap.), vědecké pojednání}
\dicEntry[fræðsla] \dicTerm{fræðsl|a} \dicIPA{{f}{r}{a}{i}{ð}{s}{\textsubring{d}}{l}{a}} \dicPos{f}[1] \dicFlx{(‑u)}[5] \dicSynonym{kennsla} \dicDirectTranslationCS{vzdělávání, výuka, osvěta}
\dicEntry[fræðsluefni] \dicTerm{fræðslu··efni} \dicIPA{{f}{r}{a}{i}{ð}{s}{\textsubring{d}}{l}{\textscy}{\textepsilon}{\textsubring{b}}{n}{\textsci}} \dicPos{n}[2] \dicFlx{(‑s, ‑)}[14] \dicDirectTranslationCS{vzdělávací\,/\addthin výukový materiál}
\dicEntry[fræðslumiðstöð] \dicTerm{fræðslu··mið·stöð} \dicIPA{{f}{r}{a}{i}{ð}{s}{\textsubring{d}}{l}{\textscy}{m}{\textsci}{ð}{s}{\textsubring{d}}{\oe}{\texttheta}} \dicPos{f}[6] \dicFlx{(‑var, ‑var)}[1] \dicFieldCat{škol.} \dicDirectTranslationCS{vzdělávací středisko}
\dicEntry[fræðslumynd] \dicTerm{fræðslu··mynd} \dicIPA{{f}{r}{a}{i}{ð}{s}{\textsubring{d}}{l}{\textscy}{m}{\textsci}{n}{\textsubring{d}}} \dicPos{f}[7] \dicFlx{(‑ar, ‑ir)}[1] \dicDirectTranslationCS{vzdělávací film}
\dicEntry[fræðsluþáttur] \dicTerm{fræðslu··|þáttur} \dicIPA{{f}{r}{a}{i}{ð}{s}{\textsubring{d}}{l}{\textscy}{\texttheta}{au}{h}{\textsubring{d}}{\textscy}{\textsubring{r}}} \dicPos{m}[12] \dicFlx{(‑þáttar, ‑þættir)}[7] \dicDirectTranslationCS{vzdělávací pořad (v~TV ap.)}
\dicEntry[frægð] \dicTerm{frægð} \dicIPA{{f}{r}{a}{i}{\textbabygamma}{\texttheta}} \dicPos{f}[4] \dicFlx{(‑ar)}[3] \dicDirectTranslationCS{sláva, proslulost}
\dicEntry[frægja] \dicTerm{fræg|ja} \dicIPA{{f}{r}{a}{i}{j}{\textlengthmark}{a}} \dicPos{v}[2] \dicFlx{(‑ði, ‑t)}[89] \dicFlx{acc} \dicDirectTranslationCS{chválit, vychvalovat, velebit}
\dicEntry[frægur] \dicTerm{frægur} \dicsymFrequent\  \dicIPA{{f}{r}{a}{i}{\textlengthmark}{\textbabygamma}{\textscy}{\textsubring{r}}} \dicPos{adj}[1]\dicFlx{}[-1] \dicDirectTranslationCS{slavný, věhlasný, proslulý};  \dicPhraseIS{gera garðinn frægan} \dicLangCat{přen.} \dicDirectTranslationCS{stát se slavným};  \dicPhraseIS{vera frægur að endemum} \dicDirectTranslationCS{být notoricky známý}
\dicEntry[frækorn] \dicTerm{fræ··korn} \dicIPA{{f}{r}{a}{i}{\textlengthmark}{k\smash{\textsuperscript{h}}}{\textopeno}{r}{\textsubring{d}}{\textsubring{n}}} \dicPos{n}[2] \dicFlx{(‑s, ‑)}[5] \dicSynonym{fræ} \dicDirectTranslationCS{semeno} \dicIndirectTranslationCS{(plod rostliny)}
\dicEntry[frændfólk] \dicTerm{frænd··fólk} \dicIPA{{f}{r}{a}{i}{n}{\textsubring{d}}{f}{ou}{\textsubring{l}}{\r{g}}} \dicPos{n}[2] \dicFlx{(‑s)}[2] \dicDirectTranslationCS{příbuzní, příbuzenstvo, rodina}
\dicEntry[frændi] \dicTerm{frænd|i} \dicsymFrequent\  \dicIPA{{f}{r}{a}{i}{n}{\textsubring{d}}{\textsci}} \dicPos{m}[2] \dicFlx{(‑a, ‑ur)}[3] \dicDirectTranslationCS{příbuzný (strýc, bratranec, synovec ap.)} \dicExampleIS{Ég á marga frændur.} \dicExampleCS{Mám mnoho příbuzných.}
\dicEntry[frændsemi] \dicTerm{frænd··sem|i} \dicIPA{{f}{r}{a}{i}{n}{\textsubring{d}}{s}{\textepsilon}{m}{\textsci}} \dicPos{f}[2] \dicFlx{(‑i)}[2] \dicDirectTranslationCS{příbuzenství}
\dicEntry[frændsystkin] \dicTerm{frænd··syst·kin} \dicIPA{{f}{r}{a}{i}{n}{\textsubring{d}}{s}{\textsci}{s}{c\smash{\textsuperscript{h}}}{\textsci}{\textsubring{n}}} \dicPos{n}[2] \dicFlx{pl}[1] \dicDirectTranslationCS{bratranci, sestřenice}
\dicEntry[frænka] \dicTerm{frænk|a} \dicsymFrequent\  \dicIPA{{f}{r}{a}{i}{\r{\ng}}{\r{g}}{a}} \dicPos{f}[1] \dicFlx{(‑u, ‑ur)}[19] \dicDirectTranslationCS{příbuzná (teta, sestřenice, neteř ap.)} \dicExampleIS{leita að frænku sinni} \dicExampleCS{hledat svoji příbuznou}
\dicEntry[fræva] \dicTerm{fræv|a} \dicIPA{{f}{r}{a}{i}{\textlengthmark}{v}{a}} \dicPos{f}[1] \dicFlx{(‑u, ‑ur)}[7] \dicFieldCat{bot.} \dicDirectTranslationCS{pestík}
\dicEntry[frævun] \dicTerm{fræv|un} \dicIPA{{f}{r}{a}{i}{\textlengthmark}{v}{\textscy}{\textsubring{n}}} \dicPos{f}[7] \dicFlx{(‑unar)}[9] \dicDirectTranslationCS{opylení, opylování}
\dicEntry[fröken] \dicTerm{fröken} \dicIPA{{f}{r}{\oe}{\textlengthmark}{\r{\textObardotlessj}}{\textepsilon}{\textsubring{n}}} \dicPos{f}[4] \dicFlx{(‑ar, ‑ar)}[1] \dicLangCat{zast.} \dicDirectTranslationCS{slečna}
\dicEntry[frömdum] \dicTerm{frömdum} \dicIPA{{f}{r}{\oe}{m}{\textsubring{d}}{\textscy}{\textsubring{m}}} \dicPos{v} \dicFlx{ind pf pl 1 pers} \dicLink{fremja}
\dicEntry[frömuður] \dicTerm{fröm··uð|ur} \dicIPA{{f}{r}{\oe}{\textlengthmark}{m}{\textscy}{ð}{\textscy}{\textsubring{r}}} \dicPos{m}[10] \dicFlx{(‑ar, ‑ir)}[4] \dicSynonym{forgöngumaður} \dicDirectTranslationCS{propagátor(ka), průkopník, průkopnice}
\dicEntry[frönsku] \dicTerm{frönsku} \dicIPA{{f}{r}{\oe}{n}{s}{\r{g}}{\textscy}} \dicPos{f} \dicFlx{sg gen} \dicLink{franska}
\dicEntry[frönskumælandi] \dicTerm{frönsku··mæl·andi} \dicIPA{{f}{r}{\oe}{n}{s}{\r{g}}{\textscy}{m}{a}{i}{l}{a}{n}{\textsubring{d}}{\textsci}} \dicPos{adj}[13] \dicFlx{indecl}[1] \dicDirectTranslationCS{francouzsky mluvící, frankofonní}
\dicEntry[fs.] \dicTerm{fs.} \dicPos{zkr} \dicPhraseIS{forsetning} \dicFieldCat{jaz.} \dicDirectTranslationCS{předložka}
\dicEntry[fsh.] \dicTerm{fsh.} \dicPos{zkr} \dicPhraseIS{framsöguháttur} \dicFieldCat{jaz.} \dicDirectTranslationCS{oznamovací způsob, indikativ}
\dicEntry[fsl.] \dicTerm{fsl.} \dicPos{zkr} \dicPhraseIS{forsetningarliður} \dicFieldCat{jaz.} \dicDirectTranslationCS{předložková fráze}
\dicEntry[ft.] \dicTerm{ft.} \dicPos{zkr} \dicPhraseIS{fleirtala} \dicFieldCat{jaz.} \dicDirectTranslationCS{množné číslo, plurál}
\dicEntry[fugl] \dicTerm{fugl} \dicsymFrequent\  \dicIPA{{f}{\textscy}{\r{g}}{\textsubring{l}}} \dicPos{m}[4] \dicFlx{(‑s, ‑ar)}[3] \dicDirectTranslationCS{pták} \dicExampleIS{vængjaþytur fugla} \dicExampleCS{šum ptačích křídel};  \dicPhraseIS{hvorki fugl né fiskur} \dicLangCat{přen.} \dicDirectTranslationCS{ani ryba ani rak}
\dicEntry[fuglabjarg] \dicTerm{fugla··|bjarg} \dicIPA{{f}{\textscy}{\r{g}}{l}{a}{\textsubring{b}}{j}{a}{r}{\r{g}}} \dicPos{n}[2] \dicFlx{(‑bjargs, ‑björg)}[8] \dicDirectTranslationCS{ptačí útes}
\dicEntry[fuglabúr] \dicTerm{fugla··búr} \dicIPA{{f}{\textscy}{\r{g}}{l}{a}{\textsubring{b}}{u}{\textsubring{r}}} \dicPos{n}[2] \dicFlx{(‑s, ‑)}[5] \dicDirectTranslationCS{ptačí klec}
\dicEntry[fuglabæli] \dicTerm{fugla··bæli} \dicIPA{{f}{\textscy}{\r{g}}{l}{a}{\textsubring{b}}{a}{i}{l}{\textsci}} \dicPos{n}[2] \dicFlx{(‑s, ‑)}[14] \dicSynonym{bæli} \dicDirectTranslationCS{ptačí hnízdo}
\dicEntry[fugladrit] \dicTerm{fugla··drit} \dicIPA{{f}{\textscy}{\r{g}}{l}{a}{\textsubring{d}}{r}{\textsci}{\textsubring{d}}} \dicPos{n}[2] \dicFlx{(‑s)}[2] \dicDirectTranslationCS{ptačí trus}
\dicEntry[fuglaflensa] \dicTerm{fugla··flens|a} \dicIPA{{f}{\textscy}{\r{g}}{l}{a}{f}{l}{\textepsilon}{n}{s}{a}} \dicPos{f}[1] \dicFlx{(‑u)}[5] \dicFieldCat{med.} \dicDirectTranslationCS{ptačí chřipka}
\dicEntry[fuglafræði] \dicTerm{fugla··fræð|i} \dicIPA{{f}{\textscy}{\r{g}}{l}{a}{f}{r}{a}{i}{ð}{\textsci}} \dicPos{f}[3] \dicFlx{(‑i)}[3] \dicDirectTranslationCS{ornitologie}
\dicEntry[fuglafræðingur] \dicTerm{fugla·fræð··ing|ur} \dicIPA{{f}{\textscy}{\r{g}}{l}{a}{f}{r}{a}{i}{ð}{i}{\ng}{\r{g}}{\textscy}{\textsubring{r}}} \dicPos{m}[6] \dicFlx{(‑s, ‑ar)}[8] \dicDirectTranslationCS{ornitolog, ornitoložka}
\dicEntry[fuglahræða] \dicTerm{fugla··hræð|a} \dicIPA{{f}{\textscy}{\r{g}}{l}{a}{\textsubring{r}}{a}{i}{ð}{a}} \dicPos{f}[1] \dicFlx{(‑u, ‑ur)}[19] \dicDirectTranslationCS{strašák, bubák} \dicIndirectTranslationCS{(panák z~hadrů ke strašení ptáků)}
\dicEntry[fuglaskoðun] \dicTerm{fugla··skoð|un} \dicIPA{{f}{\textscy}{\r{g}}{l}{a}{s}{\r{g}}{\textopeno}{ð}{\textscy}{\textsubring{n}}} \dicPos{f}[7] \dicFlx{(‑unar)}[9] \dicDirectTranslationCS{pozorování\,/\addthin observace ptáků}
\dicEntry[fuglasöngur] \dicTerm{fugla··söng|ur} \dicIPA{{f}{\textscy}{\r{g}}{l}{a}{s}{\oe i}{\ng}{\r{g}}{\textscy}{\textsubring{r}}} \dicPos{m}[8] \dicFlx{(‑s, ‑var)}[2] \dicDirectTranslationCS{ptačí zpěv}
\dicEntry[fuglatré] \dicTerm{fugla··tré} \dicIPA{{f}{\textscy}{\r{g}}{l}{a}{t\smash{\textsuperscript{h}}}{r}{j}{\textepsilon}} \dicPos{n}[2] \dicFlx{(‑s, ‑)}[36] \dicFieldCat{bot.} \dicDirectTranslationCS{trnovník akát} \textit{(l.~{\textLA{Robinia pseudoacacia}})}  \dicsymPhoto\ 
\dicFigure{ds_image_fuglatre_0_1.jpg}{Fuglatré}{Fuglatré - Jean-Pol GRANDMONT, CC BY-SA 3.0}
\dicEntry[fukum] \dicTerm{fukum} \dicIPA{{f}{\textscy}{\textlengthmark}{\r{g}}{\textscy}{\textsubring{m}}} \dicPos{v} \dicFlx{ind pf pl 1 pers} \dicLink{fjúka}
\dicEntry[full] \dicTerm{full-} \dicIPA{{f}{\textscy}{\textsubring{d}}{l}} \dicPos{predp} \dicDirectTranslationCS{plně, zcela}
\dicEntry[fullbúinn] \dicTerm{full··búinn} \dicIPA{{f}{\textscy}{\textsubring{d}}{\textsubring{l}}{\textsubring{b}}{u}{\textsci}{\textsubring{n}}} \dicPos{adj}[6]\dicFlx{}[-6] \textbf{1.} \dicSynonym{tilbúinn} \dicDirectTranslationCS{hotový, připravený}  \textbf{2.} \dicSynonym*{vel útbúinn} \dicDirectTranslationCS{plně vybavený\,/\addthin vystrojený}
\dicEntry[fullgera] \dicTerm{full··ger|a} \dicIPA{{f}{\textscy}{\textsubring{d}}{\textsubring{l}}{\r{\textObardotlessj}}{\textepsilon}{r}{a}} \dicPos{v}[2] \dicFlx{(‑ði, ‑t)}[105] \dicFlx{acc} \dicSynonym*{ljúka við} \dicDirectTranslationCS{dokončit, dokončovat, dodělat, dodělávat} \dicExampleIS{fullgera verkið} \dicExampleCS{dokončit práci}
\dicEntry[fullgerður] \dicTerm{full··gerður} \dicIPA{{f}{\textscy}{\textsubring{d}}{\textsubring{l}}{\r{\textObardotlessj}}{\textepsilon}{r}{ð}{\textscy}{\textsubring{r}}} \dicPos{adj}[2]\dicFlx{}[-1] \dicDirectTranslationCS{dokončený, dodělaný, hotový} \dicExampleIS{fullgerðar íbúðir til sölu} \dicExampleCS{dokončené byty na prodej}
\dicEntry[fullgilda] \dicTerm{full··gil|da} \dicIPA{{f}{\textscy}{\textsubring{d}}{\textsubring{l}}{\r{\textObardotlessj}}{\textsci}{l}{\textsubring{d}}{a}} \dicPos{v}[2] \dicFlx{(‑ti, ‑t)}[43] \dicFlx{acc} \dicFieldCat{práv.} \dicDirectTranslationCS{schválit, ratifikovat (smlouvu ap.)}
\dicEntry[fullkominn] \dicTerm{full··kominn} \dicsymFrequent\  \dicIPA{{f}{\textscy}{\textsubring{d}}{\textsubring{l}}{k\smash{\textsuperscript{h}}}{\textopeno}{m}{\textsci}{\textsubring{n}}} \dicPos{adj}[6]\dicFlx{}[-2] \textbf{1.} \dicDirectTranslationCS{kompletní, perfektní, dokonalý} \dicExampleIS{fullkomin áætlun} \dicExampleCS{dokonalý plán}  \textbf{2.} \dicDirectTranslationCS{naprostý} \dicIndirectTranslationCS{(k~zdůraznění)} \dicExampleIS{fullkominn sigur} \dicExampleCS{naprosté vítězství}
\dicEntry[fullkomlega] \dicTerm{full·kom··lega} \dicsymFrequent\  \dicIPA{{f}{\textscy}{\textsubring{d}}{\textsubring{l}}{k\smash{\textsuperscript{h}}}{\textopeno}{m}{l}{\textepsilon}{\textbabygamma}{a}} \dicPos{adv} \textbf{1.} \dicSynonym*{af fullkomnun} \dicDirectTranslationCS{perfektně, dokonale}  \textbf{2.} \dicSynonym{alveg} \dicDirectTranslationCS{naprosto, úplně} \dicExampleIS{Það er fullkomlega eðlilegt.} \dicExampleCS{Je to úplně normální.}
\dicEntry[fullkomleiki] \dicTerm{full·kom··leik|i} \dicIPA{{f}{\textscy}{\textsubring{d}}{\textsubring{l}}{k\smash{\textsuperscript{h}}}{\textopeno}{m}{l}{ei}{\r{\textObardotlessj}}{\textsci}} \dicPos{m}[1] \dicFlx{(‑a)}[3] \dicDirectTranslationCS{dokonalost}
\dicEntry[fullkomna] \dicTerm{full··komn|a} \dicIPA{{f}{\textscy}{\textsubring{d}}{\textsubring{l}}{k\smash{\textsuperscript{h}}}{\textopeno}{m}{n}{a}} \dicPos{v}[1] \dicFlx{(‑aði)}[1] \dicFlx{acc} \dicDirectTranslationCS{zdokonalit, dovršit, završit} \dicExampleIS{fullkomna kerfi} \dicExampleCS{zdokonalit systém}
\dicEntry[fullkomnun] \dicTerm{full··komn|un} \dicIPA{{f}{\textscy}{\textsubring{d}}{\textsubring{l}}{k\smash{\textsuperscript{h}}}{\textopeno}{m}{n}{\textscy}{\textsubring{n}}} \dicPos{f}[7] \dicFlx{(‑unar)}[9] \dicDirectTranslationCS{dokonalost, dovršení, završení}
\dicEntry[fullljós] \dicTerm{full··ljós} \dicIPA{{f}{\textscy}{\textsubring{d}}{l}{j}{ou}{s}} \dicPos{adj}[5]\dicFlx{}[-4] \dicDirectTranslationCS{zcela zřejmý\,/\addthin jasný}
\dicEntry[fullnuma] \dicTerm{full··numa\smash{\textsuperscript{1}}} \dicIPA{{f}{\textscy}{\textsubring{d}}{l}{n}{\textscy}{m}{a}} \dicPos{adj}[13] \dicFlx{indecl}[1] \dicDirectTranslationCS{vyučený, vystudovaný} \dicExampleIS{fullnuma í e‑u} \dicExampleCS{vyučený v~(čem)}
\dicEntry[fullnuma] \dicTerm{full··num|a\smash{\textsuperscript{2}}} \dicIPA{{f}{\textscy}{\textsubring{d}}{l}{n}{\textscy}{m}{a}} \dicPos{v}[1] \dicFlx{(‑aði)}[1] \dicFlx{acc} \dicPhraseIS{fullnuma sig í e‑u} \dicDirectTranslationCS{vyučit se (čemu), dokončit vzdělání v~(čem)}
\dicEntry[fullnusta] \dicTerm{fulln··ust|a} \dicIPA{{f}{\textscy}{\textsubring{d}}{l}{n}{\textscy}{s}{\textsubring{d}}{a}} \dicPos{f}[1] \dicFlx{(‑u)}[5] \dicFieldCat{práv.} \dicDirectTranslationCS{vykonání, exekuce};  \dicPhraseIS{til fullnustu} \dicFlx{adv} \dicDirectTranslationCS{(ú)plně, důkladně}
\dicEntry[fullnýta] \dicTerm{full··nýt|a} \dicIPA{{f}{\textscy}{\textsubring{d}}{l}{n}{i}{\textsubring{d}}{a}} \dicPos{v}[2] \dicFlx{(‑ti, ‑t)}[57] \dicFlx{acc} \dicDirectTranslationCS{plně využít\,/\addthin zužitkovat}
\dicEntry[fullnýting] \dicTerm{full··nýt·ing} \dicIPA{{f}{\textscy}{\textsubring{d}}{l}{n}{i}{\textsubring{d}}{i}{\ng}{\r{g}}} \dicPos{f}[4] \dicFlx{(‑ar)}[7] \dicDirectTranslationCS{plné využití\,/\addthin zužitkování}
\dicEntry[fullnæging] \dicTerm{full··næg·ing} \dicIPA{{f}{\textscy}{\textsubring{d}}{l}{n}{a}{i}{j}{i}{\ng}{\r{g}}} \dicPos{f}[4] \dicFlx{(‑ar, ‑ar)}[5] \textbf{1.} \dicDirectTranslationCS{naplnění, uspokojení}  \textbf{2.} \dicDirectTranslationCS{orgasmus} \dicExampleIS{kynferðisleg fullnæging} \dicExampleCS{sexuální uspokojení}
\dicEntry[fullnægja] \dicTerm{full··næg|ja} \dicsymFrequent\  \dicIPA{{f}{\textscy}{\textsubring{d}}{l}{n}{a}{i}{j}{a}} \dicPos{v}[2] \dicFlx{(‑ði, ‑t)}[94] \dicFlx{dat} \textbf{1.} \dicDirectTranslationCS{uspokojit, naplnit, splnit} \dicExampleIS{fullnægja ósk sinni} \dicExampleCS{splnit své přání}  \textbf{2.} \dicDirectTranslationCS{(sexuálně) uspokojit}
\dicEntry[fullnægjandi] \dicTerm{full··nægj·andi} \dicIPA{{f}{\textscy}{\textsubring{d}}{l}{n}{a}{i}{j}{a}{n}{\textsubring{d}}{\textsci}} \dicPos{adj}[13] \dicFlx{indecl}[1] \dicDirectTranslationCS{uspokojivý, dostatečný, dostačující} \dicExampleIS{Svarið er varla fullnægjandi.} \dicExampleCS{Odpověď není příliš uspokojivá.}
\dicEntry[fullorðinn] \dicTerm{full··orðinn} \dicsymFrequent\  \dicIPA{{f}{\textscy}{\textsubring{d}}{l}{\textopeno}{r}{ð}{\textsci}{\textsubring{n}}} \dicPos{adj}[6]\dicFlx{}[-1] \textbf{1.} \dicSynonym*{fullvaxinn} \dicDirectTranslationCS{dospělý} \dicExampleIS{fullorðinn einstaklingur} \dicExampleCS{dospělý jedinec}  \textbf{2.} \dicSynonym{roskinn} \dicDirectTranslationCS{postarší, (jsoucí) v~letech}
\dicEntry[fullorðinsaldur] \dicTerm{full·orðins··ald|ur} \dicIPA{{f}{\textscy}{\textsubring{d}}{l}{\textopeno}{r}{ð}{\textsci}{n}{s}{a}{l}{\textsubring{d}}{\textscy}{\textsubring{r}}} \dicPos{m}[5] \dicFlx{(‑urs, ‑rar)}[3] \dicDirectTranslationCS{dospělost}
\dicEntry[fullorðinsfræðsla] \textls[15]{\dicTerm{full·orðins··fræðsl|a} \dicIPA{{f}\-{\textscy}\-{\textsubring{d}}\-{l}\-{\textopeno}\-{r}\-{ð}\-{\textsci}\-{n}\-{s}\-{f}\-{r}\-{a}\-{i}\-{ð}\-{s}\-{\textsubring{d}}\-{l}\-{a}} \dicPos{f}[1] \dicFlx{(‑u)}[5] \dicFieldCat{škol.}} \dicDirectTranslationCS{vzdělávání dospělých}
\dicEntry[fullreyna] \dicTerm{full··reyn|a} \dicIPA{{f}{\textscy}{\textsubring{d}}{l}{r}{ei}{n}{a}} \dicPos{v}[2] \dicFlx{(‑di, ‑t)}[145] \dicFlx{acc} \dicDirectTranslationCS{plně\,/\addthin maximálně vyzkoušet}
\dicEntry[fullsaddur] \dicTerm{full··|saddur} \dicIPA{{f}{\textscy}{\textsubring{d}}{\textsubring{l}}{s}{a}{\textsubring{d}}{\textscy}{\textsubring{r}}} \dicPos{adj}[2] \dicFlx{(f ‑södd)}[20] \dicPhraseIS{fá sig fullsaddan á e‑u} \dicLangCat{přen.} \dicDirectTranslationCS{mít (čeho) dost, mít (čeho) plné zuby}
\dicEntry[fullskipaður] \dicTerm{full··skip·|aður} \dicIPA{{f}{\textscy}{\textsubring{d}}{\textsubring{l}}{s}{\r{\textObardotlessj}}{\textsci}{\textsubring{b}}{a}{ð}{\textscy}{\textsubring{r}}} \dicPos{adj}[3] \dicFlx{(f ‑uð)}[4] \dicDirectTranslationCS{zcela obsazený}
\dicEntry[fulltingi] \dicTerm{full··tingi} \dicIPA{{f}{\textscy}{\textsubring{d}}{\textsubring{l}}{t\smash{\textsuperscript{h}}}{i}{\textltailn}{\r{\textObardotlessj}}{\textsci}} \dicPos{n}[2] \dicFlx{(‑s)}[20] \dicSynonym{atbeini} \dicDirectTranslationCS{pomoc, podpora, asistence} \dicExampleIS{veita e‑m fulltingi} \dicExampleCS{poskytnout (komu) pomoc}
\dicEntry[fulltíða] \dicTerm{full··tíða} \dicIPA{{f}{\textscy}{\textsubring{d}}{\textsubring{l}}{t\smash{\textsuperscript{h}}}{i}{ð}{a}} \dicPos{adj}[13] \dicFlx{indecl}[1] \dicSynonym{fullorðinn} \dicDirectTranslationCS{dospělý} \dicExampleIS{fulltíða maður} \dicExampleCS{dospělý člověk}
\dicEntry[fulltr.] \dicTerm{fulltr.} \dicPos{zkr} \dicPhraseIS{fulltrúi} \dicLink{fulltrúi}
\dicEntry[fulltrúaráð] \dicTerm{full·trúa··ráð} \dicIPA{{f}{\textscy}{\textsubring{d}}{\textsubring{l}}{t\smash{\textsuperscript{h}}}{r}{u}{a}{r}{au}{\texttheta}} \dicPos{n}[2] \dicFlx{(‑s, ‑)}[5] \dicDirectTranslationCS{zastupitelstvo}
\dicEntry[fulltrúi] \dicTerm{full··trú|i} \dicsymFrequent\  \dicIPA{{f}{\textscy}{\textsubring{d}}{\textsubring{l}}{t\smash{\textsuperscript{h}}}{r}{u}{\textsci}} \dicPos{m}[1] \dicFlx{(‑a, ‑ar)}[1] \textbf{1.} \dicDirectTranslationCS{zástupce, zástupkyně, zastupitel(ka), představitel(ka), reprezentant(ka)} \dicExampleIS{eiga fulltrúa í nefndinni} \dicExampleCS{mít zástupce v~komisi}  \textbf{2.} \dicDirectTranslationCS{zmocněnec, zmocněnkyně, delegát(ka), důvěrník, důvěrnice}
\dicEntry[fullur] \dicTerm{fullur} \dicsymFrequent\  \dicIPA{{f}{\textscy}{\textsubring{d}}{l}{\textscy}{\textsubring{r}}} \dicPos{adj}[10] \dicFlx{(comp fyllri, sup fyllstur)}[7] \textbf{1.} \dicSynonym{troðinn} \dicDirectTranslationCS{plný, naplněný, zaplněný} \dicExampleIS{fullur af vatni} \dicExampleCS{plný vody}  \textbf{2.} \dicSynonym{drukkinn} \dicDirectTranslationCS{opilý}  \textbf{3.} \dicSynonym{óskertur} \dicDirectTranslationCS{plný, neohraničený} \dicExampleIS{fullur aðgangur} \dicExampleCS{plný přístup};  \dicIdiom{fullur}{ \dicPhraseIS{að fullu}} \dicFlx{adv} \dicDirectTranslationCS{(ú)plně, zcela}; { \dicPhraseIS{á fullu}} \dicFlx{adv} \dicDirectTranslationCS{naplno, v~plném proudu}; { \dicPhraseIS{til fulls}} \dicFlx{adv} \dicDirectTranslationCS{úplně, zcela, naplno}
\dicEntry[fullvalda] \dicTerm{full··valda} \dicIPA{{f}{\textscy}{\textsubring{d}}{l}{v}{a}{l}{\textsubring{d}}{a}} \dicPos{adj}[13] \dicFlx{indecl}[1] \dicSynonym{sjálfstæður} \dicDirectTranslationCS{svrchovaný, suverénní};  \dicPhraseIS{fullvalda ríki} \dicDirectTranslationCS{svrchovaný stát}
\dicEntry[fullveðja] \dicTerm{full··veðja} \dicIPA{{f}{\textscy}{\textsubring{d}}{l}{v}{\textepsilon}{ð}{j}{a}} \dicPos{adj}[13] \dicFlx{indecl}[1] \textbf{1.} \dicSynonym*{fjárráða} \dicDirectTranslationCS{nezávislý, soběstačný}  \textbf{2.} \dicSynonym{fullvalda} \dicDirectTranslationCS{svrchovaný, suverénní} \dicExampleIS{fullveðja þjóð} \dicExampleCS{svrchovaný národ}
\dicEntry[fullveldi] \dicTerm{full··veldi} \dicIPA{{f}{\textscy}{\textsubring{d}}{l}{v}{\textepsilon}{l}{\textsubring{d}}{\textsci}} \dicPos{n}[2] \dicFlx{(‑s, ‑)}[14] \dicFieldCat{pol.} \dicDirectTranslationCS{suverenita, svrchovanost}
\dicEntry[fullveldisdagur] \dicTerm{full·veldis··dag|ur} \dicIPA{{f}{\textscy}{\textsubring{d}}{l}{v}{\textepsilon}{l}{\textsubring{d}}{\textsci}{s}{\textsubring{d}}{a}{\textbabygamma}{\textscy}{\textsubring{r}}} \dicPos{m}[6] \dicFlx{(‑s, ‑ar)}[62] \dicDirectTranslationCS{den nezávislosti} \dicIndirectTranslationCS{(den vyhlášení nezávislosti Islandu roku 1918, islandský státní svátek -- 1. prosince)}
\dicEntry[fullvirðisréttur] \dicTerm{full·virðis··rétt|ur} \dicIPA{{f}{\textscy}{\textsubring{d}}{l}{v}{\textsci}{r}{ð}{\textsci}{s}{r}{j}{\textepsilon}{h}{\textsubring{d}}{\textscy}{\textsubring{r}}} \dicPos{m}[10] \dicFlx{(‑ar)}[7] \dicFieldCat{hist.} \dicDirectTranslationCS{státní garance na zemědělské produkty}
\dicEntry[fullvissa] \dicTerm{full··viss|a\smash{\textsuperscript{1}}} \dicIPA{{f}{\textscy}{\textsubring{d}}{l}{v}{\textsci}{s}{a}} \dicPos{f}[1] \dicFlx{(‑u)}[5] \dicSynonym{sannfæring} \dicDirectTranslationCS{jistota, přesvědčení, ujištění}
\dicEntry[fullvissa] \dicTerm{full··viss|a\smash{\textsuperscript{2}}} \dicIPA{{f}{\textscy}{\textsubring{d}}{l}{v}{\textsci}{s}{a}} \dicPos{v}[1] \dicFlx{(‑aði)}[1] \dicFlx{acc} \textbf{1.} \dicDirectTranslationCS{ujistit, ubezpečit};  \dicPhraseIS{fullvissa e‑n um e‑ð} \dicDirectTranslationCS{ujistit (koho) o~(čem)}  \textbf{2.} \dicPhraseIS{fullvissa sig} \dicDirectTranslationCS{ujistit se, ubezpečit se};  \dicPhraseIS{fullvissa sig um e‑ð} \dicDirectTranslationCS{ujistit se o~(čem)}
\dicEntry[fullyrða] \dicTerm{full··yr|ða} \dicsymFrequent\  \dicIPA{{f}{\textscy}{\textsubring{d}}{l}{\textsci}{r}{ð}{a}} \dicPos{v}[2] \dicFlx{(‑ti, ‑t)}[52] \dicFlx{acc} \dicSynonym{staðhæfa} \dicDirectTranslationCS{(po)tvrdit, prohlásit} \dicExampleIS{Hann fullyrti að þetta væri rangt.} \dicExampleCS{Potvrdil, že je to špatně.}
\dicEntry[fullyrðing] \dicTerm{full··yrð·ing} \dicIPA{{f}{\textscy}{\textsubring{d}}{l}{\textsci}{r}{ð}{i}{\ng}{\r{g}}} \dicPos{f}[4] \dicFlx{(‑ar, ‑ar)}[5] \dicDirectTranslationCS{(po)tvrzení, prohlášení}
\dicEntry[fum] \dicTerm{fum} \dicIPA{{f}{\textscy}{\textlengthmark}{\textsubring{m}}} \dicPos{n}[2] \dicFlx{(‑s)}[2] \dicSynonym{óðagot} \dicDirectTranslationCS{zmatek, panika, rychlé neuvážené jednání};  \dicPhraseIS{það er fum á e‑m} \dicDirectTranslationCS{(kdo) zmatkuje\,/\addthin panikaří}
\dicEntry[funda] \dicTerm{fund|a} \dicIPA{{f}{\textscy}{n}{\textsubring{d}}{a}} \dicPos{v}[1] \dicFlx{(‑aði)}[44] \dicDirectTranslationCS{pořádat setkání\,/\addthin schůzi, zasedat}
\dicEntry[fundargerð] \dicTerm{fundar··gerð} \dicIPA{{f}{\textscy}{n}{\textsubring{d}}{a}{r}{\r{\textObardotlessj}}{\textepsilon}{r}{\texttheta}} \dicPos{f}[7] \dicFlx{(‑ar, ‑ir)}[1] \dicDirectTranslationCS{protokol, zápis (ze schůze ap.)}
\dicEntry[fundarhald] \dicTerm{fundar··|hald} \dicIPA{{f}{\textscy}{n}{\textsubring{d}}{a}{\textsubring{r}}{h}{a}{l}{\textsubring{d}}} \dicPos{n}[2] \dicFlx{(‑halds, ‑höld)}[8] \dicDirectTranslationCS{schůze, pořádání schůze}
\dicEntry[fundarlaun] \dicTerm{fundar··laun} \dicIPA{{f}{\textscy}{n}{\textsubring{d}}{a}{r}{l}{\oe i}{\textsubring{n}}} \dicPos{n}[2] \dicFlx{pl}[1] \dicDirectTranslationCS{nálezné}
\dicEntry[fundarmaður] \dicTerm{fundar··|maður} \dicIPA{{f}{\textscy}{n}{\textsubring{d}}{a}{r}{m}{a}{ð}{\textscy}{\textsubring{r}}} \dicPos{m}[13] \dicFlx{(‑manns, ‑menn)}[2] \dicDirectTranslationCS{účastník\,/\addthin účastnice schůze, zasedající}
\dicEntry[fundarstjóri] \dicTerm{fundar··stjór|i} \dicIPA{{f}{\textscy}{n}{\textsubring{d}}{a}{\textsubring{r}}{s}{\textsubring{d}}{j}{ou}{r}{\textsci}} \dicPos{m}[1] \dicFlx{(‑a, ‑ar)}[1] \dicDirectTranslationCS{moderátor(ka) (schůze ap.), předsedající}
\dicEntry[fundið] \dicTerm{fundið} \dicIPA{{f}{\textscy}{n}{\textsubring{d}}{\textsci}{\texttheta}} \dicPos{v} \dicFlx{supin} \dicLink{finna}
\dicEntry[fundum] \dicTerm{fundum} \dicIPA{{f}{\textscy}{n}{\textsubring{d}}{\textscy}{\textsubring{m}}} \dicPos{v} \dicFlx{ind pf pl 1 pers} \dicLink{finna}
\dicEntry[fundur] \dicTerm{fund|ur} \dicsymFrequent\  \dicIPA{{f}{\textscy}{n}{\textsubring{d}}{\textscy}{\textsubring{r}}} \dicPos{m}[10] \dicFlx{(‑ar, ‑ir)}[4] \textbf{1.} \dicDirectTranslationCS{schůzka, schůze, jednání, zasedání} \dicExampleIS{árlegur fundur} \dicExampleCS{každoroční schůze};  \dicPhraseIS{e‑ir eiga fund (með e‑m)} \dicDirectTranslationCS{(kdo) má schůzku\,/\addthin schůzi (s~(kým))}  \textbf{2.} \dicDirectTranslationCS{nález, nalezení, objevení}
\dicEntry[fundvís] \dicTerm{fund··vís} \dicIPA{{f}{\textscy}{n}{\textsubring{d}}{v}{i}{s}} \dicPos{adj}[5]\dicFlx{}[-1] \dicDirectTranslationCS{dobrý v~hledání (věcí, východiska ap.)}
\dicEntry[funkisstefna] \dicTerm{funkis··stefn|a} \dicIPA{{f}{u}{\r{\textltailn}}{\r{\textObardotlessj}}{\textsci}{s}{\textsubring{d}}{\textepsilon}{\textsubring{b}}{n}{a}} \dicPos{f}[1] \dicFlx{(‑u)}[5] \dicDirectTranslationCS{funkcionalismus}
\dicEntry[fura] \dicTerm{fur|a} \dicIPA{{f}{\textscy}{\textlengthmark}{r}{a}} \dicPos{f}[1] \dicFlx{(‑u, ‑ur)}[7] \dicFieldCat{bot.} \dicDirectTranslationCS{borovice} \textit{(l.~{\textLA{Pinus}})}  \dicsymPhoto\ 
\dicFigure{ds_image_fura_0_1.jpg}{Fura}{Fura - Zp at cs.wikipedia, PD}
\dicEntry[furða] \dicTerm{furð|a\smash{\textsuperscript{1}}} \dicsymFrequent\  \dicIPA{{f}{\textscy}{r}{ð}{a}} \dicPos{f}[1] \dicFlx{(‑u, ‑ur)}[19] \dicSynonym{undrun} \dicDirectTranslationCS{údiv, úžas};  \dicPhraseIS{það er ekki\,/\addthin engin furða að} \dicDirectTranslationCS{není se co divit, že};  \dicPhraseIS{það er mesta furða (hvað e‑r gerir)} \dicDirectTranslationCS{to je úžasné, (co (kdo) dělá)}
\dicEntry[furða] \dicTerm{furð|a\smash{\textsuperscript{2}}} \dicIPA{{f}{\textscy}{r}{ð}{a}} \dicPos{v}[1] \dicFlx{(‑aði)}[57] \dicPhraseIS{e‑n furðar} \dicFlx{impers} \dicDirectTranslationCS{(kdo) se diví, (koho) překvapuje} \dicExampleIS{mig furðar að} \dicExampleCS{divím se, že};  \dicPhraseIS{furða sig á e‑u} \dicDirectTranslationCS{divit se (čemu), být (čím) překvapený} \dicExampleIS{furða sig á framkomu hans} \dicExampleCS{být překvapený jeho chováním}
\dicEntry[furðanlegur] \dicTerm{furðan··legur} \dicIPA{{f}{\textscy}{r}{ð}{a}{n}{l}{\textepsilon}{\textbabygamma}{\textscy}{\textsubring{r}}} \dicPos{adj}[1]\dicFlx{}[-8] \dicDirectTranslationCS{překvapivý, překvapující}
\dicEntry[furðulegur] \dicTerm{furðu··legur} \dicsymFrequent\  \dicIPA{{f}{\textscy}{r}{ð}{\textscy}{l}{\textepsilon}{\textbabygamma}{\textscy}{\textsubring{r}}} \dicPos{adj}[1]\dicFlx{}[-8] \dicDirectTranslationCS{zvláštní, (po)divný, bizarní} \dicExampleIS{furðulegur staður} \dicExampleCS{podivné místo}
\dicEntry[furðuverk] \dicTerm{furðu··verk} \dicIPA{{f}{\textscy}{r}{ð}{\textscy}{v}{\textepsilon}{\textsubring{r}}{\r{g}}} \dicPos{n}[2] \dicFlx{(‑s, ‑)}[5] \dicSynonym{undur\smash{\textsuperscript{1}}} \dicDirectTranslationCS{zvláštnost, podivnost}
\dicEntry[furstadæmi] \dicTerm{fursta··dæmi} \dicIPA{{f}{\textscy}{\textsubring{r}}{s}{\textsubring{d}}{a}{\textsubring{d}}{a}{i}{m}{\textsci}} \dicPos{n}[2] \dicFlx{(‑s, ‑)}[14] \dicDirectTranslationCS{knížectví};  \dicPhraseIS{Sameinuðu arabísku furstadæmin} \dicFieldCat{geog.} \dicDirectTranslationCS{Spojené arabské emiráty}
\dicEntry[fursti] \dicTerm{furst|i} \dicIPA{{f}{\textscy}{\textsubring{r}}{s}{\textsubring{d}}{\textsci}} \dicPos{m}[1] \dicFlx{(‑a, ‑ar)}[1] \dicDirectTranslationCS{kníže}
\dicEntry[furstynja] \dicTerm{furst··ynj|a} \dicIPA{{f}{\textscy}{\textsubring{r}}{s}{\textsubring{d}}{\textsci}{n}{j}{a}} \dicPos{f}[1] \dicFlx{(‑u, ‑ur)}[7] \dicDirectTranslationCS{kněžna}
\dicEntry[furunál] \dicTerm{furu··nál} \dicIPA{{f}{\textscy}{\textlengthmark}{r}{\textscy}{n}{au}{\textsubring{l}}} \dicPos{f}[4] \dicFlx{(‑ar, ‑ar)}[1] \dicDirectTranslationCS{borová jehlice}
\dicEntry[fussa] \dicTerm{fuss|a} \dicIPA{{f}{\textscy}{s}{\textlengthmark}{a}} \dicPos{v}[1] \dicFlx{(‑aði)}[44] \dicDirectTranslationCS{ohrnovat nos}
\dicEntry[fúi] \dicTerm{fú|i} \dicIPA{{f}{u}{\textlengthmark}{\textsci}} \dicPos{m}[1] \dicFlx{(‑a)}[3] \dicDirectTranslationCS{hnití, hniloba}
\dicEntry[fúinn] \dicTerm{fúinn} \dicIPA{{f}{u}{\textlengthmark}{\textsci}{\textsubring{n}}} \dicPos{adj}[6]\dicFlx{}[-2] \dicSynonym{meyr} \dicDirectTranslationCS{shnilý, prohnilý} \dicExampleIS{Timbrið er fúið.} \dicExampleCS{Dřevo je shnilé.}
\dicEntry[fúkalyf] \dicTerm{fúka··lyf}\dicTerm{, fúkkalyf} \dicIPA{{f}\-{u}\-{\textlengthmark}\-{\r{g}}\-{a}\-{l}\-{\textsci}\-{f}\-} \dicPos{n}[2] \dicFlx{(‑s, ‑)}[13] \dicFieldCat{med.} \dicDirectTranslationCS{antibiotikum}
\dicEntry[fúkkalyf] \dicTerm{fúkka··lyf} \dicIPA{{f}{u}{h}{\r{g}}{a}{l}{\textsci}{f}} \dicPos{n}[2] \dicFlx{(‑s, ‑)}[13] \dicLink{fúkalyf}
\dicEntry[fúksía] \dicTerm{fúksí|a} \dicIPA{{f}{u}{\r{g}}{s}{i}{j}{a}} \dicPos{f}[1] \dicFlx{(‑u, ‑ur)}[7] \dicFieldCat{bot.} \dicDirectTranslationCS{fuchsie} \textit{(l.~{\textLA{Fuchsia}})}  \dicsymPhoto\ 
\dicFigure{ds_image_fuksia_0_1.jpg}{Fúksía}{Fúksía - André Karwath aka Aka, CC BY-SA 2.5}
\dicEntry[fúkyrði] \dicTerm{fúk··yrði} \dicIPA{{f}{u}{\textlengthmark}{\r{g}}{\textsci}{r}{ð}{\textsci}} \dicPos{n}[2] \dicFlx{(‑s, ‑)}[14] \dicSynonym*{illyrði} \dicDirectTranslationCS{nadávka}
\dicEntry[fúlegg] \dicTerm{fúl··egg} \dicIPA{{f}{u}{\textlengthmark}{l}{\textepsilon}{\r{g}}} \dicPos{n}[2] \dicFlx{(‑s, ‑)}[13] \dicDirectTranslationCS{záprtek, zkažené vejce}
\dicEntry[fúlga] \dicTerm{fúlg|a} \dicIPA{{f}{u}{l}{\r{g}}{a}} \dicPos{f}[1] \dicFlx{(‑u, ‑ur)}[19] \dicDirectTranslationCS{hromada (peněz), balík (peněz)}
\dicEntry[fúll] \dicTerm{fúll} \dicsymFrequent\  \dicIPA{{f}{u}{\textsubring{d}}{\textsubring{l}}} \dicPos{adj}[8]\dicFlx{}[-1] \textbf{1.} \dicSynonym{skemmdur} \dicDirectTranslationCS{zatuchlý, smrdící} \dicExampleIS{fúlt loft} \dicExampleCS{zatuchlé povětří}  \textbf{2.} \dicSynonym*{í vondu skapi} \dicDirectTranslationCS{mrzutý, nevrlý, otrávený, naštvaný};  \dicPhraseIS{vera fúll út í e‑n} \dicDirectTranslationCS{být na (koho) nevrlý};  \dicPhraseIS{e‑m er fúlasta alvara} \dicFlx{impers} \dicDirectTranslationCS{(kdo) to myslí naprosto vážně}
\dicEntry[fúlmenni] \dicTerm{fúl··menni} \dicIPA{{f}{u}{l}{m}{\textepsilon}{n}{\textsci}} \dicPos{n}[2] \dicFlx{(‑s, ‑)}[14] \dicDirectTranslationCS{ničema, neřád, padouch, padouška}
\dicEntry[fúlmennska] \dicTerm{fúl··mennsk|a} \dicIPA{{f}{u}{l}{m}{\textepsilon}{n}{s}{\r{g}}{a}} \dicPos{f}[1] \dicFlx{(‑u)}[5] \dicDirectTranslationCS{ničemnost, podlost, darebáctví}
\dicEntry[fúlna] \dicTerm{fúln|a} \dicIPA{{f}{u}{l}{n}{a}} \dicPos{v}[1] \dicFlx{(‑aði)}[34] \dicDirectTranslationCS{zatuchnout, zkazit se (vejce ap.)} \dicExampleIS{Vatnið fúlnaði.} \dicExampleCS{Voda zatuchla.}
\dicEntry[fúlsa] \dicTerm{fúls|a} \dicIPA{{f}{u}{l}{s}{a}} \dicPos{v}[1] \dicFlx{(‑aði)}[44] \dicPhraseIS{fúlsa við e‑u} \dicDirectTranslationCS{ohrnovat nad (čím) nos} \dicExampleIS{fúlsa við matnum} \dicExampleCS{ohrnovat nos nad jídlem}
\dicEntry[fúna] \dicTerm{fún|a} \dicIPA{{f}{u}{\textlengthmark}{n}{a}} \dicPos{v}[1] \dicFlx{(‑aði)}[34] \dicDirectTranslationCS{(s)hnít, tlít, zahnívat}
\dicEntry[fús] \dicTerm{fús} \dicIPA{{f}{u}{\textlengthmark}{s}} \dicPos{adj}[5]\dicFlx{}[-1] \dicSynonym{gjarn} \dicDirectTranslationCS{ochotný, ochoten (k~nějaké činnosti ap.)};  \dicPhraseIS{af fúsum og frjálsum vilja} \dicFlx{adv} \dicDirectTranslationCS{dobrovolně, ochotně}
\dicEntry[fúska] \dicTerm{fúsk|a} \dicIPA{{f}{u}{s}{\r{g}}{a}} \dicPos{v}[1] \dicFlx{(‑aði)}[34] \dicPhraseIS{fúska við e‑ð} \dicDirectTranslationCS{fušovat do (čeho), plést se do (čeho)}
\dicEntry[fúskari] \dicTerm{fúsk··ar|i} \dicIPA{{f}{u}{s}{\r{g}}{a}{r}{\textsci}} \dicPos{m}[1] \dicFlx{(‑a, ‑ar)}[13] \dicDirectTranslationCS{fušer(ka), břídil(ka)}
\dicEntry[fúslega] \dicTerm{fús··lega} \dicIPA{{f}{u}{\textlengthmark}{s}{l}{\textepsilon}{\textbabygamma}{a}} \dicPos{adv} \dicSynonym{gjarnan} \dicDirectTranslationCS{ochotně, rád} \dicExampleIS{Kennari svarar fúslega spurningum í tíma.} \dicExampleCS{Během vyučování odpovídá učitel ochotně na otázky.}
\dicEntry[fútúrismi] \dicTerm{fútúr··ism|i} \dicIPA{{f}{u}{\textlengthmark}{\textsubring{d}}{u}{r}{\textsci}{s}{m}{\textsci}} \dicPos{m}[1] \dicFlx{(‑a)}[3] \dicDirectTranslationCS{futurismus}
\dicEntry[fyki] \dicTerm{fyki} \dicIPA{{f}{\textsci}{\textlengthmark}{\r{\textObardotlessj}}{\textsci}} \dicPos{v} \dicFlx{con pf sg 1 pers} \dicLink{fjúka}
\dicEntry[fylgd] \dicTerm{fylgd} \dicIPA{{f}{\textsci}{l}{\textsubring{d}}} \dicPos{f}[7] \dicFlx{(‑ar, ‑ir)}[1] \dicDirectTranslationCS{společnost, doprovod, doprovázení};  \dicPhraseIS{í fylgd e‑rs} \dicDirectTranslationCS{v~doprovodu (koho)}
\dicEntry[fylgdarlið] \dicTerm{fylgdar··lið} \dicIPA{{f}{\textsci}{l}{\textsubring{d}}{a}{r}{l}{\textsci}{\texttheta}} \dicPos{n}[2] \dicFlx{(‑s, ‑)}[5] \dicDirectTranslationCS{doprovod, eskorta}
\dicEntry[fylgdarmaður] \dicTerm{fylgdar··|maður} \dicIPA{{f}{\textsci}{l}{\textsubring{d}}{a}{r}{m}{a}{ð}{\textscy}{\textsubring{r}}} \dicPos{m}[13] \dicFlx{(‑manns, ‑menn)}[2] \dicDirectTranslationCS{průvodce, průvodkyně, doprovod, doprovázející}
\dicEntry[fylgi] \dicTerm{fylgi} \dicIPA{{f}{\textsci}{l}{\r{\textObardotlessj}}{\textsci}} \dicPos{n}[2] \dicFlx{(‑s)}[20] \dicDirectTranslationCS{podpora, podporování} \dicExampleIS{veita e‑m fylgi} \dicExampleCS{podporovat (koho)}
\dicEntry[fylgibréf] \dicTerm{fylgi··bréf} \dicIPA{{f}{\textsci}{l}{\r{\textObardotlessj}}{\textsci}{\textsubring{b}}{r}{j}{\textepsilon}{f}} \dicPos{n}[2] \dicFlx{(‑s, ‑)}[5] \dicDirectTranslationCS{průvodka, průvodní list}
\dicEntry[fylgifiskur] \dicTerm{fylgi··fisk|ur} \dicIPA{{f}{\textsci}{l}{\r{\textObardotlessj}}{\textsci}{f}{\textsci}{s}{\r{g}}{\textscy}{\textsubring{r}}} \dicPos{m}[6] \dicFlx{(‑s, ‑ar)}[10] \dicDirectTranslationCS{průvodní jev} \dicExampleIS{e‑að er fylgifiskur e‑s} \dicExampleCS{(co) je průvodní jev (čeho)}
\dicEntry[fylgihlutur] \dicTerm{fylgi··hlut|ur} \dicIPA{{f}{\textsci}{l}{\r{\textObardotlessj}}{\textsci}{\textsubring{l}}{\textscy}{\textsubring{d}}{\textscy}{\textsubring{r}}} \dicPos{m}[10] \dicFlx{(‑ar, ‑ir)}[13] \dicDirectTranslationCS{doplněk, příslušenství}
\dicEntry[fylgihnöttur] \dicTerm{fylgi··|hnöttur} \dicIPA{{f}{\textsci}{l}{\r{\textObardotlessj}}{\textsci}{\textsubring{n}}{\oe}{h}{\textsubring{d}}{\textscy}{\textsubring{r}}} \dicPos{m}[11] \dicFlx{(‑hnattar, ‑hnettir)}[5] \dicFieldCat{astro.} \dicDirectTranslationCS{oběžnice, satelit}
\dicEntry[fylginautur] \dicTerm{fylgi··naut|ur} \dicIPA{{f}{\textsci}{l}{\r{\textObardotlessj}}{\textsci}{n}{\oe i}{\textsubring{d}}{\textscy}{\textsubring{r}}} \dicPos{m}[6] \dicFlx{(‑s\,/\addthin ‑ar, ‑ar)}[65] \dicDirectTranslationCS{společník, společnice, průvodce, průvodkyně}
\dicEntry[fylgirit] \dicTerm{fylgi··rit} \dicIPA{{f}{\textsci}{l}{\r{\textObardotlessj}}{\textsci}{r}{\textsci}{\textsubring{d}}} \dicPos{n}[2] \dicFlx{(‑s, ‑)}[5] \dicDirectTranslationCS{příloha (v~novinách ap.)}
\dicEntry[fylgiskjal] \dicTerm{fylgi··|skjal} \dicIPA{{f}{\textsci}{l}{\r{\textObardotlessj}}{\textsci}{s}{\r{\textObardotlessj}}{a}{\textsubring{l}}} \dicPos{n}[2] \dicFlx{(‑skjals, ‑skjöl)}[8] \dicDirectTranslationCS{příloha (k~dokumentu ap.)} \dicExampleIS{fylgiskjal samnings} \dicExampleCS{příloha smlouvy}
\dicEntry[fylgismaður] \dicTerm{fylgis··|maður} \dicIPA{{f}{\textsci}{l}{\r{\textObardotlessj}}{\textsci}{s}{m}{a}{ð}{\textscy}{\textsubring{r}}} \dicPos{m}[13] \dicFlx{(‑manns, ‑menn)}[2] \dicDirectTranslationCS{přívrženec, přívrženkyně, stoupenec, stoupenkyně, následovník, následovnice}
\dicEntry[fylgja] \dicTerm{fylgj|a\smash{\textsuperscript{1}}} \dicIPA{{f}{\textsci}{l}{\r{\textObardotlessj}}{a}} \dicPos{f}[1] \dicFlx{(‑u, ‑ur)}[17] \textbf{1.} \dicFieldCat{anat.} \dicDirectTranslationCS{placenta, plodové lůžko}  \textbf{2.} \dicFieldCat{pov.} \dicSynonym{vofa\smash{\textsuperscript{1}}} \dicDirectTranslationCS{(ochranný) duch, přízrak}
\dicEntry[fylgja] \dicTerm{fylg|ja\smash{\textsuperscript{2}}} \dicsymFrequent\  \dicIPA{{f}{\textsci}{l}{\r{\textObardotlessj}}{a}} \dicPos{v}[2] \dicFlx{(‑di, ‑t)}[129] \dicFlx{dat} \textbf{1.} \dicSynonym*{fara með} \dicDirectTranslationCS{doprovodit, (do)provázet, vést} \dicExampleIS{fylgja e‑m til dyra} \dicExampleCS{doprovodit (koho) ke dveřím}  \textbf{2.} \dicSynonym{styðja} \dicDirectTranslationCS{podporovat, podpořit}  \textbf{3.} \dicDirectTranslationCS{řídit se, plnit (rady ap.)} \dicExampleIS{fylgja fyrirmælum} \dicExampleCS{řídit se pokyny}  \textbf{4.} \dicDirectTranslationCS{doprovázet, být součástí};  \dicIdiom{fylgja}[eftir]{ \dicPhraseIS{fylgja e‑m eftir}} {\textbf{a.}} \dicDirectTranslationCS{sledovat (koho), jít za (kým)};  {\textbf{b.}} \dicDirectTranslationCS{sledovat (koho), stíhat (koho) (při projevu ap.)};  \dicIdiom{fylgja}[með]{ \dicPhraseIS{fylgja með e‑u}} \dicDirectTranslationCS{být součástí (čeho), doprovázet (co)};  \dicIdiom{fylgjast}[með]{ \dicPhraseIS{fylgjast með e‑u}} \dicFlx{refl} \dicDirectTranslationCS{sledovat (co), pozorovat (co)} \dicExampleIS{Ég fylgist með fréttum.} \dicExampleCS{Sleduji zprávy.}
\dicEntry[fylgjandi] \dicTerm{fylgj··andi} \dicIPA{{f}{\textsci}{l}{\r{\textObardotlessj}}{a}{n}{\textsubring{d}}{\textsci}} \dicPos{adj}[13] \dicFlx{indecl}[1] \dicDirectTranslationCS{(jsoucí) pro, podporující};  \dicPhraseIS{vera fylgjandi e‑u} \dicDirectTranslationCS{být pro (co)}
\dicEntry[fylgni] \dicTerm{fylgn|i} \dicIPA{{f}{\textsci}{l}{\r{g}}{n}{\textsci}} \dicPos{f}[3] \dicFlx{(‑i)}[3] \dicDirectTranslationCS{vzájemná závislost, korelace}
\dicEntry[fylgsni] \dicTerm{fylgsni} \dicIPA{{f}{\textsci}{l}{s}{n}{\textsci}} \dicPos{n}[2] \dicFlx{(‑s, ‑)}[14] \dicSynonym{felustaður} \dicDirectTranslationCS{skrýš, úkryt} \dicExampleIS{eiga fylgsni undir klettunum} \dicExampleCS{mít skrýš pod skálou}
\dicEntry[fylki] \dicTerm{fylki} \dicIPA{{f}{\textsci}{\textsubring{l}}{\r{\textObardotlessj}}{\textsci}} \dicPos{n}[2] \dicFlx{(‑s, ‑)}[16] \textbf{1.} \dicSynonym{ríki} \dicDirectTranslationCS{stát (americký ap.), kraj, kanton}  \textbf{2.} \dicFieldCat{mat.} \dicDirectTranslationCS{matice}
\dicEntry[fylking] \dicTerm{fylk··ing} \dicIPA{{f}{\textsci}{\textsubring{l}}{\r{\textObardotlessj}}{i}{\ng}{\r{g}}} \dicPos{f}[4] \dicFlx{(‑ar, ‑ar)}[5] \textbf{1.} \dicSynonym*{skoðanahópur} \dicDirectTranslationCS{skupina (názorová ap.), frakce}  \textbf{2.} \dicSynonym{hópur} \dicDirectTranslationCS{houf, skupina, zástup}  \textbf{3.} \dicFieldCat{biol.} \dicDirectTranslationCS{kmen, phylum}
\dicEntry[fylkja] \dicTerm{fylk|ja} \dicIPA{{f}{\textsci}{\textsubring{l}}{\r{\textObardotlessj}}{a}} \dicPos{v}[2] \dicFlx{(‑ti, ‑t)}[25] \dicFlx{dat} \dicPhraseIS{fylkja liði} \dicSynonym{fjölmenna} \dicDirectTranslationCS{shromáždit se, vytvořit zástup}
\dicEntry[fylla] \dicTerm{fyll|a} \dicsymFrequent\  \dicIPA{{f}{\textsci}{\textsubring{d}}{l}{a}} \dicPos{v}[2] \dicFlx{(‑ti, ‑t)}[71] \dicFlx{acc} \dicDirectTranslationCS{(na)plnit, naplňovat, vyplnit, vyplňovat, zaplnit, zaplňovat} \dicExampleIS{fylla ílátið af vatni} \dicExampleCS{naplnit nádobu vodou};  \dicPhraseIS{fylla e‑n} \dicDirectTranslationCS{opít (koho)};  \dicIdiom{fylla}[á]{ \dicPhraseIS{fylla á e‑u}} \dicDirectTranslationCS{naplnit (co) (láhev vodou ap.)}; { \dicPhraseIS{fylla á bílinn\,/\addthin tankinn}} \dicDirectTranslationCS{natankovat};  \dicIdiom{fylla}[upp]{ \dicPhraseIS{fylla upp í e‑ð}} \dicDirectTranslationCS{zaplnit (co), vyplnit (co) (mezeru ap.)};  \dicIdiom{fylla}[út]{ \dicPhraseIS{fylla út e‑ð}} \dicDirectTranslationCS{vyplnit (co), vyplňovat (co)} \dicExampleIS{fylla út eyðublað} \dicExampleCS{vyplnit dotazník};  \dicIdiom{fyllast}{ \dicPhraseIS{fyllast}} \dicFlx{refl} \dicDirectTranslationCS{naplnit se, zaplnit se} \dicExampleIS{Stofan fylltist af fólki.} \dicExampleCS{Místnost se zaplnila lidmi.}; { \dicPhraseIS{fyllast reiði}} \dicFlx{refl} \dicDirectTranslationCS{rozčílit se, rozzlobit se}
\dicEntry[fyllerí] \dicTerm{fyllerí} \dicIPA{{f}{\textsci}{\textsubring{d}}{l}{\textepsilon}{r}{i}} \dicPos{n}[2] \dicFlx{(‑s, ‑)}[5] \dicLink{fyllirí}
\dicEntry[fylli] \dicTerm{fyll|i} \dicIPA{{f}{\textsci}{\textsubring{d}}{l}{\textsci}} \dicPos{f}[5] \dicFlx{(‑i\,/\addthin ‑ar, ‑ir\,/\addthin ‑ar)}[9] \dicDirectTranslationCS{sytost, nasycenost, nasycení};  \dicPhraseIS{éta fylli sína} \dicDirectTranslationCS{nasytit se}
\dicEntry[fyllibytta] \dicTerm{fylli··bytt|a} \dicIPA{{f}{\textsci}{\textsubring{d}}{l}{\textsci}{\textsubring{b}}{\textsci}{h}{\textsubring{d}}{a}} \dicPos{f}[1] \dicFlx{(‑u, ‑ur)}[19] \dicLangCat{han.} \dicDirectTranslationCS{opilec, opilkyně, ochlasta, ochmelka}
\dicEntry[fyllilega] \dicTerm{fylli··lega} \dicIPA{{f}{\textsci}{\textsubring{d}}{l}{\textsci}{l}{\textepsilon}{\textbabygamma}{a}} \dicPos{adv} \dicSynonym{fullkomlega} \dicDirectTranslationCS{plně, zcela}
\dicEntry[fylliliður] \dicTerm{fylli··lið|ur} \dicIPA{{f}{\textsci}{\textsubring{d}}{l}{\textsci}{l}{\textsci}{ð}{\textscy}{\textsubring{r}}} \dicPos{m}[10] \dicFlx{(‑ar\,/\addthin ‑s, ‑ir)}[33] \dicFieldCat{jaz.} \dicDirectTranslationCS{doplněk}
\dicEntry[fylling] \dicTerm{fyll··ing} \dicIPA{{f}{\textsci}{\textsubring{d}}{l}{i}{\ng}{\r{g}}} \dicPos{f}[4] \dicFlx{(‑ar, ‑ar)}[5] \textbf{1.} \dicSynonym{tannfylling} \dicDirectTranslationCS{výplň, (zubní) plomba} \dicExampleIS{fylling í tönn} \dicExampleCS{zubní plomba}  \textbf{2.} \dicDirectTranslationCS{náplň (do tiskárny ap.)} \dicExampleIS{fylling í kúlupenna} \dicExampleCS{náplň do kuličkového pera}  \textbf{3.} \dicFieldCat{kulin.} \dicDirectTranslationCS{nádivka};  \dicPhraseIS{tunglið er í fyllingu} \dicDirectTranslationCS{měsíc je v~úplňku}
\dicEntry[fyllingarefni] \dicTerm{fyllingar··efni} \dicIPA{{f}{\textsci}{\textsubring{d}}{l}{i}{\ng}{\r{g}}{a}{r}{\textepsilon}{\textsubring{b}}{n}{\textsci}} \dicPos{n}[2] \dicFlx{(‑s, ‑)}[14] \dicDirectTranslationCS{výplň, plnivo, plnicí materiál}
\dicEntry[fyllirí] \dicTerm{fyllirí}\dicTerm{, fyllerí} \dicIPA{{f}\-{\textsci}\-{\textsubring{d}}\-{l}\-{\textsci}\-{r}\-{i}\-} \dicPos{n}[2] \dicFlx{(‑s)}[2] \dicLangCat{hovor.} \dicDirectTranslationCS{flám, tah, mejdan};  \dicPhraseIS{fara á fyllirí} \dicDirectTranslationCS{jít flámovat}
\dicEntry[fyllstur] \dicTerm{fyllstur} \dicIPA{{f}{\textsci}{l}{s}{\textsubring{d}}{\textscy}{\textsubring{r}}} \dicPos{adj} \dicFlx{m sg nom sup} \dicLink{fullur}
\dicEntry[fyndi] \dicTerm{fyndi} \dicIPA{{f}{\textsci}{n}{\textsubring{d}}{\textsci}} \dicPos{v} \dicFlx{con pf sg 1 pers} \dicLink{finna}
\dicEntry[fyndinn] \dicTerm{fyndinn} \dicsymFrequent\  \dicIPA{{f}{\textsci}{n}{\textsubring{d}}{\textsci}{\textsubring{n}}} \dicPos{adj}[6]\dicFlx{}[-2] \dicSynonym{hlægilegur} \dicDirectTranslationCS{vtipný, legrační} \dicExampleIS{reyna að vera fyndinn} \dicExampleCS{snažit se být vtipný}
\dicEntry[fyndni] \dicTerm{fyndn|i} \dicIPA{{f}{\textsci}{n}{\textsubring{d}}{n}{\textsci}} \dicPos{f}[3] \dicFlx{(‑i)}[3] \textbf{1.} \dicDirectTranslationCS{humor, vtipnost}  \textbf{2.} \dicSynonym*{skrýtla} \dicDirectTranslationCS{vtip, žert, šprým}
\dicEntry[fyrir] \dicTerm{fyrir} \dicsymFrequent\  \dicIPA{{f}{\textsci}{\textlengthmark}{r}{\textsci}{\textsubring{r}}} \dicPos{prep\,/\addthin adv} \dicFlx{acc\,/\addthin dat} \textbf{1.} \dicFlx{dat} \dicDirectTranslationCS{před} \dicIndirectTranslationCS{(o~něčem, co stojí v~cestě někomu)} \dicExampleIS{Stóllinn er fyrir mér.} \dicExampleCS{Židle je přede mnou.\,/\addthin Židle mi stojí v~cestě.}  \textbf{2.} \dicFlx{dat} \dicDirectTranslationCS{před} \dicIndirectTranslationCS{(o~směru nebo pohybu)} \dicExampleIS{ganga fyrir konung} \dicExampleCS{předstoupit před krále}  \textbf{3.} \dicFlx{dat} \dicDirectTranslationCS{před} \dicIndirectTranslationCS{(o~čase, který již nastal)} \dicExampleIS{fyrir ári síðan} \dicExampleCS{před jedním rokem}  \textbf{4.} \dicFlx{acc} \dicDirectTranslationCS{před} \dicIndirectTranslationCS{(o~čase, který ještě nenastal)} \dicExampleIS{fyrir jól} \dicExampleCS{před Vánocemi}  \textbf{5.} \dicFlx{dat} \dicDirectTranslationCS{o} \dicIndirectTranslationCS{(s~určitými podstatnými jmény)} \dicExampleIS{sönnun fyrir e‑u} \dicExampleCS{důkaz o~(čem)}  \textbf{6.} \dicFlx{acc} \dicDirectTranslationCS{pro} \dicIndirectTranslationCS{(v~zájmu někoho, pro někoho)} \dicExampleIS{bók fyrir almenning} \dicExampleCS{kniha pro veřejnost}  \textbf{7.} \dicFlx{acc\,/\addthin dat} \dicLangCat{hovor.} \dicDirectTranslationCS{pro} \dicIndirectTranslationCS{(o~pozici ve hře)} \dicExampleIS{3-1 fyrir Reykjavík} \dicExampleCS{3:1 pro Reykjavík}  \textbf{8.} \dicFlx{dat} \dicDirectTranslationCS{kvůli, díky} \dicExampleIS{heyra ekkert fyrir hávaða} \dicExampleCS{neslyšet nic kvůli hluku}  \textbf{9.} \dicFlx{acc} \dicSynonym{sökum} \dicDirectTranslationCS{díky, kvůli} \dicIndirectTranslationCS{(o~příčině)} \dicExampleIS{fyrir tilviljun} \dicExampleCS{díky náhodě}  \textbf{10.} \dicFlx{acc} \dicDirectTranslationCS{za} \dicIndirectTranslationCS{(o~prodeji nebo nákupu)} \dicExampleIS{gott verð fyrir e‑ð} \dicExampleCS{dobrá cena za (co)}  \textbf{11.} \dicFlx{acc} \dicDirectTranslationCS{za, místo} \dicIndirectTranslationCS{(v~náhradě za)}  \textbf{12.} \dicFlx{acc} \dicDirectTranslationCS{po, za} \dicIndirectTranslationCS{(o~pořadí)} \dicExampleIS{lið fyrir lið} \dicExampleCS{bod po bodu};  \dicIdiom{fyrir}[aftan]{ \dicPhraseIS{fyrir aftan e‑ð}} \dicFlx{prep} \dicDirectTranslationCS{za (čím)} \dicIndirectTranslationCS{(místně)};  \dicIdiom{fyrir}[austan]{ \dicPhraseIS{fyrir austan e‑ð}} \dicFlx{prep} \dicDirectTranslationCS{na východ od (čeho)};  \dicPhraseIS{fyrir austan} \dicFlx{adv} {\textbf{a.}} \dicDirectTranslationCS{na východě};  {\textbf{b.}} \dicDirectTranslationCS{na východě (Islandu)};  {\textbf{c.}} \dicDirectTranslationCS{na východě, ve východní Evropě};  \dicIdiom{fyrir}[fram]{ \dicPhraseIS{fyrir aldur fram}} \dicFlx{adv} \dicDirectTranslationCS{předčasně, příliš brzy}; { \dicPhraseIS{fyrir fram}} \dicFlx{adv} \dicDirectTranslationCS{předem, dopředu (zaplatit ap.)} \dicIndirectTranslationCS{(o~čase)};  \dicIdiom{fyrir}[framan]{ \dicPhraseIS{fyrir framan e‑ð}} \dicFlx{prep} \dicDirectTranslationCS{před (čím)} \dicIndirectTranslationCS{(místně)} \dicExampleIS{fyrir framan apótekið} \dicExampleCS{před lékárnou};  \dicIdiom{fyrir}[handan]{ \dicPhraseIS{fyrir handan e‑ð}} \dicFlx{prep} \dicDirectTranslationCS{na druhé straně (čeho)};  \dicPhraseIS{fyrir handan} \dicFlx{adv} \dicDirectTranslationCS{na onom světě, na druhé straně} \dicIndirectTranslationCS{(v~ráji)};  \dicIdiom{fyrir}[innan]{ \dicPhraseIS{fyrir innan e‑ð}} \dicFlx{prep} \dicDirectTranslationCS{(uvnitř) v~(čem)};  \dicPhraseIS{fyrir innan} \dicFlx{adv} \dicDirectTranslationCS{(dále) uvnitř};  \dicIdiom{fyrir}[neðan]{ \dicPhraseIS{fyrir neðan e‑ð}} \dicFlx{prep} \dicDirectTranslationCS{pod (čím)} \dicIndirectTranslationCS{(místně)};  \dicIdiom{fyrir}[norðan]{ \dicPhraseIS{fyrir norðan e‑ð}} \dicFlx{prep} \dicDirectTranslationCS{na sever od (čeho)}; { \dicPhraseIS{fyrir norðan}} \dicFlx{adv} {\textbf{a.}} \dicDirectTranslationCS{na severu};  {\textbf{b.}} \dicDirectTranslationCS{na severu (Islandu)};  \dicIdiom{fyrir}[ofan]{ \dicPhraseIS{fyrir ofan e‑ð}} \dicFlx{prep} \dicDirectTranslationCS{nad (čím)} \dicIndirectTranslationCS{(místně)};  \dicIdiom{fyrir}[sunnan]{ \dicPhraseIS{fyrir sunnan e‑ð}} \dicFlx{prep} \dicDirectTranslationCS{na jih od (čeho)};  \dicPhraseIS{fyrir sunnan} \dicFlx{adv} {\textbf{a.}} \dicDirectTranslationCS{na jihu};  {\textbf{b.}} \dicDirectTranslationCS{na jihu (v~oblasti hlavního města Islandu)};  \dicIdiom{fyrir}[utan]{ \dicPhraseIS{fyrir utan e‑ð}} \dicFlx{prep} {\textbf{a.}} \dicDirectTranslationCS{před (čím)} \dicIndirectTranslationCS{(místně)} \dicExampleIS{fyrir utan dyrnar} \dicExampleCS{přede dveřmi};  {\textbf{b.}} \dicDirectTranslationCS{kromě (čeho), mimo (čeho)} \dicExampleIS{fyrir utan þig} \dicExampleCS{kromě tebe};  \dicPhraseIS{fyrir utan} \dicFlx{adv} \dicDirectTranslationCS{(venku) před};  \dicIdiom{fyrir}[vestan]{ \dicPhraseIS{fyrir vestan e‑ð}} \dicFlx{prep} \dicDirectTranslationCS{na západ od (čeho)};  \dicPhraseIS{fyrir vestan} \dicFlx{adv} {\textbf{a.}} \dicDirectTranslationCS{na západě};  {\textbf{b.}} \dicDirectTranslationCS{na západě (Islandu)};  {\textbf{c.}} \dicDirectTranslationCS{na západě (v~západní Evropě)};  {\textbf{d.}} \dicDirectTranslationCS{v~Americe (v~USA)}
\dicEntry[fyrirbauð] \dicTerm{fyrir··bauð} \dicIPA{{f}{\textsci}{\textlengthmark}{r}{\textsci}{r}{\textsubring{b}}{\oe i}{\texttheta}} \dicPos{v} \dicFlx{ind pf sg 1 pers} \dicLink{fyrirbjóða}
\dicEntry[fyrirbjóða] \dicTerm{fyrir··|bjóða} \dicIPA{{f}{\textsci}{\textlengthmark}{r}{\textsci}{r}{\textsubring{b}}{j}{ou}{ð}{a}} \dicPos{v}[6] \dicFlx{(‑býð, ‑bauð, ‑buðum, ‑byði, ‑boð\-ið)}[104] \dicFlx{dat} \dicLangCat{zast.} \dicSynonym{banna} \dicDirectTranslationCS{zakázat, zakazovat, nedovolit, nedovolovat} \dicExampleIS{fyrirbjóða e‑m að gera e‑ð} \dicExampleCS{nedovolit (co komu) udělat}
\dicEntry[fyrirboði] \dicTerm{fyrir··boð|i} \dicIPA{{f}{\textsci}{\textlengthmark}{r}{\textsci}{r}{\textsubring{b}}{\textopeno}{ð}{\textsci}} \dicPos{m}[1] \dicFlx{(‑a, ‑ar)}[1] \dicDirectTranslationCS{předzvěst, znamení}
\dicEntry[fyrirboðið] \dicTerm{fyrir··boðið} \dicIPA{{f}{\textsci}{\textlengthmark}{r}{\textsci}{r}{\textsubring{b}}{\textopeno}{ð}{\textsci}{\texttheta}} \dicPos{v} \dicFlx{supin} \dicLink{fyrirbjóða}
\dicEntry[fyrirbrigði] \dicTerm{fyrir··brigði} \dicIPA{{f}{\textsci}{\textlengthmark}{r}{\textsci}{r}{\textsubring{b}}{r}{\textsci}{\textbabygamma}{ð}{\textsci}} \dicPos{n}[2] \dicFlx{(‑s, ‑)}[14] \dicDirectTranslationCS{jev, úkaz, fenomén} \dicExampleIS{náttúrufyrirbrigði} \dicExampleCS{přírodní úkaz}
\dicEntry[fyrirbuðum] \dicTerm{fyrir··buðum} \dicIPA{{f}{\textsci}{\textlengthmark}{r}{\textsci}{r}{\textsubring{b}}{\textscy}{ð}{\textscy}{\textsubring{m}}} \dicPos{v} \dicFlx{ind pf pl 1 pers} \dicLink{fyrirbjóða}
\dicEntry[fyrirbyði] \dicTerm{fyrir··byði} \dicIPA{{f}{\textsci}{\textlengthmark}{r}{\textsci}{r}{\textsubring{b}}{\textsci}{ð}{\textsci}} \dicPos{v} \dicFlx{con pf sg 1 pers} \dicLink{fyrirbjóða}
\dicEntry[fyrirbýð] \dicTerm{fyrir··býð} \dicIPA{{f}{\textsci}{\textlengthmark}{r}{\textsci}{r}{\textsubring{b}}{i}{\texttheta}} \dicPos{v} \dicFlx{ind praes sg 1 pers} \dicLink{fyrirbjóða}
\dicEntry[fyrirbæri] \dicTerm{fyrir··bæri} \dicsymFrequent\  \dicIPA{{f}{\textsci}{\textlengthmark}{r}{\textsci}{r}{\textsubring{b}}{a}{i}{r}{\textsci}} \dicPos{n}[2] \dicFlx{(‑s, ‑)}[14] \textbf{1.} \dicSynonym{fyrirbrigði} \dicDirectTranslationCS{jev, úkaz, fenomén} \dicExampleIS{hversdagslegt fyrirbæri} \dicExampleCS{každodenní jev}  \textbf{2.} \dicDirectTranslationCS{fenomén} \dicIndirectTranslationCS{(člověk v~některém oboru mimořádně vynikající)}
\dicEntry[fyrirfara] \dicTerm{fyrir··|fara} \dicIPA{{f}{\textsci}{\textlengthmark}{r}{\textsci}{\textsubring{r}}{f}{a}{r}{a}} \dicPos{v}[8] \dicFlx{(‑fer, ‑fór, ‑fórum, ‑færi, ‑farið)}[4] \dicPhraseIS{fyrirfara sér} \dicDirectTranslationCS{zabít se, spáchat sebevraždu}
\dicEntry[fyrirfarið] \dicTerm{fyrir··farið} \dicIPA{{f}{\textsci}{\textlengthmark}{r}{\textsci}{\textsubring{r}}{f}{a}{r}{\textsci}{\texttheta}} \dicPos{v} \dicFlx{supin} \dicLink{fyrirfara}
\dicEntry[fyrirfer] \dicTerm{fyrir··fer} \dicIPA{{f}{\textsci}{\textlengthmark}{r}{\textsci}{\textsubring{r}}{f}{\textepsilon}{\textsubring{r}}} \dicPos{v} \dicFlx{ind praes sg 1 pers} \dicLink{fyrirfara}
\dicEntry[fyrirferð] \dicTerm{fyrir··ferð} \dicIPA{{f}{\textsci}{\textlengthmark}{r}{\textsci}{\textsubring{r}}{f}{\textepsilon}{r}{\texttheta}} \dicPos{f}[7] \dicFlx{(‑ar)}[3] \dicSynonym{stærð} \dicDirectTranslationCS{rozsah, velikost, objem}
\dicEntry[fyrirferðarmikill] \dicTerm{fyrir·ferðar··|mikill} \dicIPA{{f}\-{\textsci}\-{\textlengthmark}\-{r}\-{\textsci}\-{\textsubring{r}}\-{f}\-{\textepsilon}\-{r}\-{ð}\-{a}\-{r}\-{m}\-{\textsci}\-{\r{\textObardotlessj}}\-{\textsci}\-{\textsubring{d}}\-{\textsubring{l}}\-} \dicPos{adj}[11] \dicFlx{(comp ‑meiri, sup ‑mestur)}[12] \textbf{1.} \dicDirectTranslationCS{objemný, rozměrný} \dicExampleIS{Skreiðin er fyrirferðarmikil vara.} \dicExampleCS{Sušená ryba je objemné zboží.}  \textbf{2.} \dicDirectTranslationCS{náročný, nekontrolovatelný (dítě ap.)}
\dicEntry[fyrirfór] \dicTerm{fyrir··fór} \dicIPA{{f}{\textsci}{\textlengthmark}{r}{\textsci}{\textsubring{r}}{f}{ou}{\textsubring{r}}} \dicPos{v} \dicFlx{ind pf sg 1 pers} \dicLink{fyrirfara}
\dicEntry[fyrirfórum] \dicTerm{fyrir··fórum} \dicIPA{{f}{\textsci}{\textlengthmark}{r}{\textsci}{\textsubring{r}}{f}{ou}{r}{\textscy}{\textsubring{m}}} \dicPos{v} \dicFlx{ind pf pl 1 pers} \dicLink{fyrirfara}
\dicEntry[fyrirfram] \dicTerm{fyrir··fram} \dicsymFrequent\  \dicIPA{{f}{\textsci}{\textlengthmark}{r}{\textsci}{\textsubring{r}}{f}{r}{a}{\textsubring{m}}} \dicPos{adv} \dicDirectTranslationCS{předem, dopředu} \dicIndirectTranslationCS{(o~čase)} \dicExampleIS{greiða fyrirfram} \dicExampleCS{zaplatit předem}
\begin{xtolerant}{}{1pt}
\dicEntry[fyrirframgreiðsla] \dicTerm{fyrir·fram··greiðsl|a}\addthinS\dicIPA{{f}\-{\textsci}\-{\textlengthmark}\-{r}\-{\textsci}\-{\textsubring{r}}\-{f}\-{r}\-{a}\-{m}\-{\r{g}}\-{r}\-{ei}\-{ð}\-{s}\-{\textsubring{d}}\-{l}\-{a}}\addthinS\dicPos{f}[1]\addthinS\dicFlx{(‑u,\addthinS‑ur)}[13] \dicDirectTranslationCS{plat\-ba předem}
\end{xtolerant}
\dicEntry[fyrirfæri] \dicTerm{fyrir··færi} \dicIPA{{f}{\textsci}{\textlengthmark}{r}{\textsci}{\textsubring{r}}{f}{a}{i}{r}{\textsci}} \dicPos{v} \dicFlx{con pf sg 1 pers} \dicLink{fyrirfara}
\dicEntry[fyrirgaf] \dicTerm{fyrir··gaf} \dicIPA{{f}{\textsci}{\textlengthmark}{r}{\textsci}{r}{\r{g}}{a}{f}} \dicPos{v} \dicFlx{ind pf sg 1 pers} \dicLink{fyrirgefa}
\dicEntry[fyrirgangur] \dicTerm{fyrir··gang|ur} \dicIPA{{f}{\textsci}{\textlengthmark}{r}{\textsci}{r}{\r{g}}{au}{\ng}{\r{g}}{\textscy}{\textsubring{r}}} \dicPos{m}[6] \dicFlx{(‑s)}[7] \dicSynonym{bægslagangur} \dicDirectTranslationCS{povyk, pozdvižení}
\dicEntry[fyrirgáfum] \dicTerm{fyrir··gáfum} \dicIPA{{f}{\textsci}{\textlengthmark}{r}{\textsci}{r}{\r{g}}{au}{v}{\textscy}{\textsubring{m}}} \dicPos{v} \dicFlx{ind pf pl 1 pers} \dicLink{fyrirgefa}
\dicEntry[fyrirgef] \dicTerm{fyrir··gef} \dicIPA{{f}{\textsci}{\textlengthmark}{r}{\textsci}{r}{\r{\textObardotlessj}}{\textepsilon}{f}} \dicPos{v} \dicFlx{ind praes sg 1 pers} \dicLink{fyrirgefa}
\dicEntry[fyrirgefa] \dicTerm{fyrir··|gefa} \dicsymFrequent\  \dicIPA{{f}{\textsci}{\textlengthmark}{r}{\textsci}{r}{\r{\textObardotlessj}}{\textepsilon}{v}{a}} \dicPos{v}[6] \dicFlx{(‑gef, ‑gaf, ‑gáfum, ‑gæfi, ‑gefið)}[8] \dicFlx{dat + acc} \dicDirectTranslationCS{odpustit, odpouštět, prominout, promíjet} \dicExampleIS{fyrirgefa e‑m e‑ð} \dicExampleCS{prominout (komu co)};  \dicPhraseIS{fyrirgefðu} \dicDirectTranslationCS{promiň, odpusť, pardon};  \dicIdiom{fyrirgefast}{ \dicPhraseIS{e‑m fyrirgefst}} \dicFlx{refl impers} \dicDirectTranslationCS{(komu) je odpuštěno\,/\addthin prominuto}
\dicEntry[fyrirgefanlegur] \dicTerm{fyrir·gefan··legur} \dicIPA{{f}{\textsci}{\textlengthmark}{r}{\textsci}{r}{\r{\textObardotlessj}}{\textepsilon}{v}{a}{n}{l}{\textepsilon}{\textbabygamma}{\textscy}{\textsubring{r}}} \dicPos{adj}[1]\dicFlx{}[-8] \dicDirectTranslationCS{omluvitelný, prominutelný} \dicExampleIS{fyrirgefanleg mistök} \dicExampleCS{omluvitelná chyba}
\dicEntry[fyrirgefið] \dicTerm{fyrir··gefið} \dicIPA{{f}{\textsci}{\textlengthmark}{r}{\textsci}{r}{\r{\textObardotlessj}}{\textepsilon}{v}{\textsci}{\texttheta}} \dicPos{v} \dicFlx{supin} \dicLink{fyrirgefa}
\dicEntry[fyrirgefning] \dicTerm{fyrir··gef·ning} \dicIPA{{f}{\textsci}{\textlengthmark}{r}{\textsci}{r}{\r{\textObardotlessj}}{\textepsilon}{\textsubring{b}}{n}{i}{\ng}{\r{g}}} \dicPos{f}[4] \dicFlx{(‑ar, ‑ar)}[5] \dicDirectTranslationCS{odpuštění, prominutí} \dicExampleIS{biðja e‑n fyrirgefningar á e‑u} \dicExampleCS{požádat (koho) o~prominutí (čeho)}
\dicEntry[fyrirgengilegur] \dicTerm{fyrir·gengi··legur} \dicIPA{{f}{\textsci}{\textlengthmark}{r}{\textsci}{r}{\r{\textObardotlessj}}{ei}{\textltailn}{\r{\textObardotlessj}}{\textsci}{l}{\textepsilon}{\textbabygamma}{\textscy}{\textsubring{r}}} \dicPos{adj}[1]\dicFlx{}[-8] \dicDirectTranslationCS{vetchý, chatrný}
\dicEntry[fyrirgreiðsla] \dicTerm{fyrir··greiðsl|a} \dicIPA{{f}{\textsci}{\textlengthmark}{r}{\textsci}{r}{\r{g}}{r}{ei}{ð}{s}{\textsubring{d}}{l}{a}} \dicPos{f}[1] \dicFlx{(‑u, ‑ur)}[13] \dicDirectTranslationCS{pomoc, podpora, služba}
\dicEntry[fyrirgæfi] \dicTerm{fyrir··gæfi} \dicIPA{{f}{\textsci}{\textlengthmark}{r}{\textsci}{r}{\r{\textObardotlessj}}{a}{i}{v}{\textsci}} \dicPos{v} \dicFlx{con pf sg 1 pers} \dicLink{fyrirgefa}
\dicEntry[fyrirhafnarlaus] \dicTerm{fyrir·hafnar··laus} \dicIPA{{f}{\textsci}{\textlengthmark}{r}{\textsci}{\textsubring{r}}{h}{a}{\textsubring{b}}{n}{a}{r}{l}{\oe i}{s}} \dicPos{adj}[5]\dicFlx{}[-1] \dicDirectTranslationCS{snadný, nenamáhavý} \dicExampleIS{fyrirhafnarlaus akstur} \dicExampleCS{snadná jízda}
\dicEntry[fyrirheit] \dicTerm{fyrir··heit} \dicsymFrequent\  \dicIPA{{f}{\textsci}{\textlengthmark}{r}{\textsci}{\textsubring{r}}{h}{ei}{\textsubring{d}}} \dicPos{n}[2] \dicFlx{(‑s, ‑)}[5] \dicDirectTranslationCS{příslib};  \dicPhraseIS{fyrirheit um e‑ð} \dicDirectTranslationCS{příslib (čeho)}
\dicEntry[fyrirhuga] \dicTerm{fyrir··hug|a} \dicIPA{{f}{\textsci}{\textlengthmark}{r}{\textsci}{\textsubring{r}}{h}{\textscy}{\textbabygamma}{a}} \dicPos{v}[1] \dicFlx{(‑aði)}[1] \dicFlx{acc} \dicDirectTranslationCS{zamýšlet, (na)plánovat}
\dicEntry[fyrirhugaður] \dicTerm{fyrir··hug·|aður} \dicIPA{{f}{\textsci}{\textlengthmark}{r}{\textsci}{\textsubring{r}}{h}{\textscy}{\textbabygamma}{a}{ð}{\textscy}{\textsubring{r}}} \dicPos{adj}[3] \dicFlx{(f ‑uð)}[4] \dicDirectTranslationCS{zamýšlený, plánovaný}
\dicEntry[fyrirhyggja] \dicTerm{fyrir··hyggj|a} \dicIPA{{f}{\textsci}{\textlengthmark}{r}{\textsci}{\textsubring{r}}{h}{\textsci}{\r{\textObardotlessj}}{a}} \dicPos{f}[1] \dicFlx{(‑u)}[5] \dicDirectTranslationCS{prozíravost, předvídavost} \dicExampleIS{hafa litla fyrirhyggju} \dicExampleCS{nebýt příliš prozíravý}
\dicEntry[fyrirhyggjulaus] \dicTerm{fyrir·hyggju··laus} \dicIPA{{f}{\textsci}{\textlengthmark}{r}{\textsci}{\textsubring{r}}{h}{\textsci}{\r{\textObardotlessj}}{\textscy}{l}{\oe i}{s}} \dicPos{adj}[5]\dicFlx{}[-1] \dicDirectTranslationCS{neprozíravý, nepředvídavý}
\dicEntry[fyrirhyggjuleysi] \dicTerm{fyrir·hyggju··leysi} \dicIPA{{f}{\textsci}{\textlengthmark}{r}{\textsci}{\textsubring{r}}{h}{\textsci}{\r{\textObardotlessj}}{\textscy}{l}{ei}{s}{\textsci}} \dicPos{n}[2] \dicFlx{(‑s)}[20] \dicDirectTranslationCS{neprozíravost, nepředvídavost}
\dicEntry[fyrirhyggjumaður] \dicTerm{fyrir·hyggju··|maður} \dicIPA{{f}{\textsci}{\textlengthmark}{r}{\textsci}{\textsubring{r}}{h}{\textsci}{\r{\textObardotlessj}}{\textscy}{m}{a}{ð}{\textscy}{\textsubring{r}}} \dicPos{m}[13] \dicFlx{(‑manns, ‑menn)}[2] \dicDirectTranslationCS{prozíravý\,/\addthin předvídavý člověk}
\dicEntry[fyrirhöfn] \dicTerm{fyrir··|höfn} \dicIPA{{f}{\textsci}{\textlengthmark}{r}{\textsci}{\textsubring{r}}{h}{\oe}{\textsubring{b}}{\textsubring{n}}} \dicPos{f}[7] \dicFlx{(‑hafnar)}[19] \dicSynonym{áreynsla} \dicDirectTranslationCS{úsilí, snaha} \dicExampleIS{Þetta kostar hann mikla fyrirhöfn.} \dicExampleCS{Stojí ho to velké úsilí.}
\dicEntry[fyrirkallaður] \dicTerm{fyrir··|kall·aður} \dicIPA{{f}{\textlengthmark}{\textsci}{r}{\textsci}{\textsubring{r}}{k\smash{\textsuperscript{h}}}{a}{\textsubring{d}}{l}{a}{ð}{\textscy}{\textsubring{r}}} \dicPos{adj}[3] \dicFlx{(f ‑kölluð)}[2] \dicPhraseIS{vera illa fyrirkallaður} \dicDirectTranslationCS{být ve špatném stavu, být indisponovaný}
\dicEntry[fyrirkomulag] \textls[15]{\dicTerm{fyrir·komu··lag} \dicIPA{{f}{\textlengthmark}{\textsci}{r}{\textsci}{\textsubring{r}}{k\smash{\textsuperscript{h}}}{\textopeno}{m}{\textscy}{l}{a}{x}} \dicPos{n}[2] \dicFlx{(‑s)}[2] \dicSynonym{tilhögun}} \dicDirectTranslationCS{uspo\-řádání, organizace, režim}
\dicEntry[fyrirleit] \dicTerm{fyrir··leit} \dicIPA{{f}{\textsci}{\textlengthmark}{r}{\textsci}{r}{l}{ei}{\textsubring{d}}} \dicPos{v} \dicFlx{ind pf sg 1 pers} \dicLink{fyrirlíta}
\dicEntry[fyrirlesari] \dicTerm{fyrir··les·ar|i} \dicIPA{{f}{\textsci}{\textlengthmark}{r}{\textsci}{r}{l}{\textepsilon}{s}{a}{r}{\textsci}} \dicPos{m}[1] \dicFlx{(‑a, ‑ar)}[13] \dicDirectTranslationCS{přednášející (na vysoké škole ap.), lektor(ka)}
\dicEntry[fyrirlestrasalur] \dicTerm{fyrir·lestra··sal|ur} \dicIPA{{f}{\textsci}{\textlengthmark}{r}{\textsci}{r}{l}{\textepsilon}{s}{\textsubring{d}}{r}{a}{s}{a}{l}{\textscy}{\textsubring{r}}} \dicPos{m}[10] \dicFlx{(‑ar, ‑ir)}[14] \dicDirectTranslationCS{posluchárna, auditorium, přednáškový sál}
\dicEntry[fyrirlestur] \dicTerm{fyrir··lest|ur} \dicsymFrequent\  \dicIPA{{f}{\textsci}{\textlengthmark}{r}{\textsci}{r}{l}{\textepsilon}{s}{\textsubring{d}}{\textscy}{\textsubring{r}}} \dicPos{m}[5] \dicFlx{(‑rar\,/\addthin ‑urs, ‑rar)}[9] \dicDirectTranslationCS{přednáška, výklad} \dicExampleIS{sækja fyrirlestra hjá prófessornum} \dicExampleCS{docházet na přednášky k~profesorovi}
\dicEntry[fyrirliði] \dicTerm{fyrir··lið|i} \dicIPA{{f}{\textsci}{\textlengthmark}{r}{\textsci}{r}{l}{\textsci}{ð}{\textsci}} \dicPos{m}[1] \dicFlx{(‑a, ‑ar)}[1] \dicFieldCat{sport.} \dicSynonym{foringi} \dicDirectTranslationCS{kapitán(ka) (týmu ap.)}
\dicEntry[fyrirliggjandi] \textls[15]{\dicTerm{fyrir··liggj·andi} \dicIPA{{f}{\textsci}{\textlengthmark}{r}{\textsci}{r}{l}{\textsci}{\r{\textObardotlessj}}{a}{n}{\textsubring{d}}{\textsci}} \dicPos{adj}[13] \dicFlx{indecl}[1] \dicDirectTranslationCS{(jsoucí) na skladě, (jsoucí) skladem}}
\dicEntry[fyrirliti] \dicTerm{fyrir··liti} \dicIPA{{f}{\textsci}{\textlengthmark}{r}{\textsci}{r}{l}{\textsci}{\textsubring{d}}{\textsci}} \dicPos{v} \dicFlx{con pf sg 1 pers} \dicLink{fyrirlíta}
\dicEntry[fyrirlitið] \dicTerm{fyrir··litið} \dicIPA{{f}{\textsci}{\textlengthmark}{r}{\textsci}{r}{l}{\textsci}{\textsubring{d}}{\textsci}{\texttheta}} \dicPos{v} \dicFlx{supin} \dicLink{fyrirlíta}
\dicEntry[fyrirlitlegur] \textls[15]{\dicTerm{fyrir·lit··legur} \dicIPA{{f}{\textsci}{\textlengthmark}{r}{\textsci}{r}{l}{\textsci}{\textsubring{d}}{l}{\textepsilon}{\textbabygamma}{\textscy}{\textsubring{r}}} \dicPos{adj}[1]\dicFlx{}[-8]} \dicSynonym{andstyggilegur} \dicDirectTranslationCS{opovrženíhodný, odporný, ohavný}
\dicEntry[fyrirlitning] \dicTerm{fyrir··lit·ning} \dicIPA{{f}{\textsci}{\textlengthmark}{r}{\textsci}{r}{l}{\textsci}{h}{\textsubring{d}}{n}{i}{\ng}{\r{g}}} \dicPos{f}[4] \dicFlx{(‑ar)}[7] \dicDirectTranslationCS{opovržení, opovrhování, pohrdání};  \dicPhraseIS{hafa fyrirlitningu á e‑m} \dicDirectTranslationCS{opovrhovat (kým)}
\dicEntry[fyrirlitum] \dicTerm{fyrir··litum} \dicIPA{{f}{\textsci}{\textlengthmark}{r}{\textsci}{r}{l}{\textsci}{\textsubring{d}}{\textscy}{\textsubring{m}}} \dicPos{v} \dicFlx{ind pf pl 1 pers} \dicLink{fyrirlíta}
\dicEntry[fyrirlít] \dicTerm{fyrir··lít} \dicIPA{{f}{\textsci}{\textlengthmark}{r}{\textsci}{r}{l}{i}{\textsubring{d}}} \dicPos{v} \dicFlx{ind praes sg 1 pers} \dicLink{fyrirlíta}
\dicEntry[fyrirlíta] \dicTerm{fyrir··|líta} \dicIPA{{f}{\textsci}{\textlengthmark}{r}{\textsci}{r}{l}{i}{\textsubring{d}}{a}} \dicPos{v}[6] \dicFlx{(‑lít, ‑leit, ‑litum, ‑liti, ‑litið)}[77] \dicFlx{acc} \dicDirectTranslationCS{opovrhovat, pohrdat} \dicExampleIS{fyrirlíta e‑n fyrir e‑ð} \dicExampleCS{opovrhovat (kým) kvůli (čemu)}
\dicEntry[fyrirmuna] \dicTerm{fyrir··mun|a} \dicIPA{{f}{\textsci}{\textlengthmark}{r}{\textsci}{r}{m}{\textscy}{n}{a}} \dicPos{v}[1] \dicFlx{(‑aði)}[44] \dicFlx{dat} \dicPhraseIS{e‑m er fyrirmunað að (gera e‑ð)} \dicFlx{impers} \dicDirectTranslationCS{pro (koho) není možné (udělat (co))}
\dicEntry[fyrirmynd] \dicTerm{fyrir··mynd} \dicsymFrequent\  \dicIPA{{f}{\textsci}{\textlengthmark}{r}{\textsci}{r}{m}{\textsci}{n}{\textsubring{d}}} \dicPos{f}[7] \dicFlx{(‑ar, ‑ir)}[1] \dicDirectTranslationCS{předloha, vzor, model, příklad};  \dicPhraseIS{fyrirmynd að e‑u} \dicDirectTranslationCS{předloha pro (co)};  \dicPhraseIS{taka e‑n sér til fyrirmyndar} \dicDirectTranslationCS{vzít si (koho) za vzor, mít (koho) za vzor};  \dicPhraseIS{e‑að er til fyrirmyndar} \dicDirectTranslationCS{(co) je ukázkové (chování ap.)}
\dicEntry[fyrirmæli] \dicTerm{fyrir··mæli} \dicIPA{{f}{\textsci}{\textlengthmark}{r}{\textsci}{r}{m}{a}{i}{l}{\textsci}} \dicPos{n}[2] \dicFlx{pl}[19] \dicSynonym{skipun} \dicDirectTranslationCS{instrukce, pokyn, nařízení};  \dicPhraseIS{fara að fyrirmælunum} \dicDirectTranslationCS{řídit se instrukcemi};  \dicPhraseIS{að fyrirmælum e‑rs} \dicDirectTranslationCS{podle (čích) pokynů}
\dicEntry[fyrirrennari] \dicTerm{fyrir··renn·ar|i} \dicIPA{{f}{\textsci}{\textlengthmark}{r}{\textsci}{r}{\textepsilon}{n}{a}{r}{\textsci}} \dicPos{m}[1] \dicFlx{(‑a, ‑ar)}[13] \dicDirectTranslationCS{předchůdce, předchůdkyně}
\dicEntry[fyrirsát] \dicTerm{fyrir··sát} \dicIPA{{f}{\textsci}{\textlengthmark}{r}{\textsci}{\textsubring{r}}{s}{au}{\textsubring{d}}} \dicPos{f}[7] \dicFlx{(‑ar, ‑ir)}[1] \dicDirectTranslationCS{léčka, útok ze zálohy} \dicExampleIS{gera\,/\addthin veita herflokknum fyrirsát} \dicExampleCS{napadnout armádní jednotku ze zálohy}
\dicEntry[fyrirsjáanlegur] \dicTerm{fyrir··sjáan·legur} \dicIPA{{f}{\textsci}{\textlengthmark}{r}{\textsci}{\textsubring{r}}{s}{j}{au}{a}{n}{l}{\textepsilon}{\textbabygamma}{\textscy}{\textsubring{r}}} \dicPos{adj}[1]\dicFlx{}[-8] \dicDirectTranslationCS{předvídatel\-ný, očekávatelný}
\dicEntry[fyrirskipun] \dicTerm{fyrir··skip|un} \dicIPA{{f}{\textsci}{\textlengthmark}{r}{\textsci}{\textsubring{r}}{s}{\r{\textObardotlessj}}{\textsci}{\textsubring{b}}{\textscy}{\textsubring{n}}} \dicPos{f}[7] \dicFlx{(‑unar, ‑anir)}[8] \dicSynonym{skipun} \dicDirectTranslationCS{příkaz, pokyn};  \dicPhraseIS{fyrirskipun um e‑ð} \dicDirectTranslationCS{příkaz (čeho)} \dicExampleIS{fyrirskipun um uppgjöf} \dicExampleCS{příkaz vzdání se};  \dicPhraseIS{að fyrirskipun e‑rs} \dicDirectTranslationCS{na (čí) příkaz}
\dicEntry[fyrirsláttur] \dicTerm{fyrir··slátt|ur} \dicIPA{{f}{\textsci}{\textlengthmark}{r}{\textsci}{\textsubring{r}}{s}{\textsubring{d}}{l}{au}{h}{\textsubring{d}}{\textscy}{\textsubring{r}}} \dicPos{m}[12] \dicFlx{(‑ar)}[8] \dicDirectTranslationCS{záminka, vytáčka} \dicExampleIS{finna fyrirslátt af e‑u} \dicExampleCS{nalézt záminku v~(čem)}
\dicEntry[fyrirspurn] \dicTerm{fyrir··spurn} \dicIPA{{f}{\textsci}{\textlengthmark}{r}{\textsci}{\textsubring{r}}{s}{\textsubring{b}}{\textscy}{r}{\textsubring{d}}{\textsubring{n}}} \dicPos{f}[7] \dicFlx{(‑ar, ‑ir)}[1] \dicDirectTranslationCS{(oficiální) dotaz, interpelace};  \dicPhraseIS{leggja fram fyrirspurn} \dicDirectTranslationCS{položit dotaz}
\dicEntry[fyrirstaða] \dicTerm{fyrir··|staða} \dicIPA{{f}{\textsci}{\textlengthmark}{r}{\textsci}{\textsubring{r}}{s}{\textsubring{d}}{a}{ð}{a}} \dicPos{f}[1] \dicFlx{(‑stöðu, ‑stöður)}[8] \dicDirectTranslationCS{překážka, zábrana, obstrukce} \dicIndirectTranslationCS{(okolnost něčemu\,/\addthin někomu překážející)};  \dicPhraseIS{það er ekkert því til fyrirstöðu} \dicDirectTranslationCS{nic tomu nebrání, nic tomu nestojí v~cestě}
\dicEntry[fyrirsæta] \dicTerm{fyrir··sæt|a} \dicIPA{{f}{\textsci}{\textlengthmark}{r}{\textsci}{\textsubring{r}}{s}{a}{i}{\textsubring{d}}{a}} \dicPos{f}[1] \dicFlx{(‑u, ‑ur)}[19] \textbf{1.} \dicDirectTranslationCS{model(ka)} \dicIndirectTranslationCS{(pózující malířům ap.)}  \textbf{2.} \dicDirectTranslationCS{model(ka)} \dicIndirectTranslationCS{(v~módním návrhářství)}
\dicEntry[fyrirsögn] \dicTerm{fyrir··|sögn} \dicIPA{{f}{\textsci}{\textlengthmark}{r}{\textsci}{\textsubring{r}}{s}{\oe}{\r{g}}{\textsubring{n}}} \dicPos{f}[7] \dicFlx{(‑sagnar, ‑sagnir)}[16] \textbf{1.} \dicSynonym{skipun} \dicDirectTranslationCS{pokyn, instrukce} \dicExampleIS{fara að fyrirsögn kennarans} \dicExampleCS{řídit se učitelovým pokynem}  \textbf{2.} \dicSynonym{yfirskrift} \dicDirectTranslationCS{titul, nadpis (kapitoly ap.)}
\dicEntry[fyrirtekt] \dicTerm{fyrir··tekt} \dicIPA{{f}{\textsci}{\textlengthmark}{r}{\textsci}{\textsubring{r}}{t\smash{\textsuperscript{h}}}{\textepsilon}{x}{\textsubring{d}}} \dicPos{f}[7] \dicFlx{(‑ar)}[3] \dicDirectTranslationCS{vrtoch, rozmar, chvilkový nápad}
\dicEntry[fyrirtæki] \dicTerm{fyrir··tæki} \dicsymFrequent\  \dicIPA{{f}{\textsci}{\textlengthmark}{r}{\textsci}{\textsubring{r}}{t\smash{\textsuperscript{h}}}{a}{i}{\r{\textObardotlessj}}{\textsci}} \dicPos{n}[2] \dicFlx{(‑s, ‑)}[16] \textbf{1.} \dicDirectTranslationCS{firma, podnik, společnost} \dicExampleIS{reka fyrirtæki} \dicExampleCS{provozovat firmu}  \textbf{2.} \dicDirectTranslationCS{aktivita, předsevzetí, projekt}
\dicEntry[fyrirtækjasamsteypa] \textls[15]{\dicTerm{fyrir·tækja··sam·steyp|a} \dicIPA{{f}\-{\textsci}\-{\textlengthmark}\-{r}\-{\textsci}\-{\textsubring{r}}\-{t\smash{\textsuperscript{h}}}\-{a}\-{i}\-{\r{\textObardotlessj}}\-{a}\-{s}\-{a}\-{m}\-{s}\-{\textsubring{d}}\-{ei}\-{\textsubring{b}}\-{a}}} \dicPos{f}[1] \dicFlx{(‑u, ‑ur)}[19] \dicFieldCat{ekon.} \dicDirectTranslationCS{konsorcium podniků}
\dicEntry[fyrirtækjaskrá] \dicTerm{fyrir·tækja··skrá} \dicIPA{{f}{\textsci}{\textlengthmark}{r}{\textsci}{\textsubring{r}}{t\smash{\textsuperscript{h}}}{a}{i}{\r{\textObardotlessj}}{a}{s}{\r{g}}{r}{au}} \dicPos{f}[4] \dicFlx{(‑r\,/\addthin ‑ar, ‑r)}[21] \dicDirectTranslationCS{obchod\-ní rejstřík, adresář firem}
\dicEntry[fyrirvari] \dicTerm{fyrir··var|i} \dicsymFrequent\  \dicIPA{{f}{\textsci}{\textlengthmark}{r}{\textsci}{r}{v}{a}{r}{\textsci}} \dicPos{m}[1] \dicFlx{(‑a, ‑ar)}[8] \textbf{1.} \dicSynonym{frestur} \dicDirectTranslationCS{oznámení, avízo}  \textbf{2.} \dicSynonym{skilyrði} \dicDirectTranslationCS{výhrada};  \dicPhraseIS{með fyrirvara} \dicFlx{adv} \dicDirectTranslationCS{s~výhradami} \dicExampleIS{samþykkja tillöguna með fyrirvara} \dicExampleCS{přijmout návrh s~výhradami}
\dicEntry[fyrirverða] \dicTerm{fyrir··verð|a} \dicIPA{{f}{\textsci}{\textlengthmark}{r}{\textsci}{r}{v}{\textepsilon}{r}{ð}{a}} \dicPos{v}[1] \dicFlx{(‑aði)}[1] \dicPhraseIS{fyrirverða sig} \dicDirectTranslationCS{stydět se}
\dicEntry[fyrirvinna] \dicTerm{fyrir··vinn|a} \dicIPA{{f}{\textsci}{\textlengthmark}{r}{\textsci}{r}{v}{\textsci}{n}{a}} \dicPos{f}[1] \dicFlx{(‑u)}[5] \dicDirectTranslationCS{živitel(ka)}
\dicEntry[fyrirætlun] \dicTerm{fyrir··ætl|un} \dicIPA{{f}{\textsci}{\textlengthmark}{r}{\textsci}{r}{a}{i}{h}{\textsubring{d}}{l}{\textscy}{\textsubring{n}}} \dicPos{f}[7] \dicFlx{(‑unar, ‑anir)}[8] \dicSynonym{áform} \dicDirectTranslationCS{záměr, plán} \dicExampleIS{segja e‑m frá fyrirætlunum sínum} \dicExampleCS{říci (komu) o~svých záměrech}
\dicEntry[fyrna] \dicTerm{fyrn|a} \dicIPA{{f}{\textsci}{r}{\textsubring{d}}{n}{a}} \dicPos{v}[2] \dicFlx{(‑di, ‑t)}[141] \dicFlx{acc} \dicDirectTranslationCS{archaizovat};  \dicIdiom{fyrnast}{ \dicPhraseIS{fyrnast}} \dicFlx{refl} \dicDirectTranslationCS{vypršet, promlčet se} \dicExampleIS{Skuldin fyrnist fljótt.} \dicExampleCS{Dluh bude brzy promlčen.}
\dicEntry[fyrnd] \dicTerm{fyrnd} \dicIPA{{f}{\textsci}{r}{n}{\textsubring{d}}} \dicPos{f}[4] \dicFlx{(‑ar)}[3] \dicDirectTranslationCS{dávná doba, dávnověk};  \dicPhraseIS{í fyrndinni} \dicFlx{adv} \dicDirectTranslationCS{v~dávných dobách}
\dicEntry[fyrning] \dicTerm{fyrn··ing} \dicIPA{{f}{\textsci}{r}{\textsubring{d}}{n}{i}{\ng}{\r{g}}} \dicPos{f}[4] \dicFlx{(‑ar, ‑ar)}[5] \textbf{1.} \dicFieldCat{práv.} \dicSynonym{úrelding} \dicDirectTranslationCS{promlčení}  \textbf{2.} \dicFieldCat{ekon.} \dicSynonym{afskrift} \dicDirectTranslationCS{devalvace, znehodnocení}  \textbf{3.} \dicSynonym*{gamalt hey} \dicDirectTranslationCS{seno (nebo jiné věci) nechané z~předchozího roku}  \textbf{4.} \dicDirectTranslationCS{archaizace (řeči ap.)}
\dicEntry[fyrr] \dicTerm{fyrr} \dicsymFrequent\  \dicIPA{{f}{\textsci}{r}{\textlengthmark}} \dicPos{adv} \dicFlx{comp (pos snemma, sup fyrst)} \dicSynonym{áður} \dicDirectTranslationCS{dříve};  \dicPhraseIS{fyrr á tímum} \dicFlx{adv} \dicDirectTranslationCS{dříve, v~dřívějších dobách};  \dicPhraseIS{fyrr á öldum} \dicFlx{adv} \dicDirectTranslationCS{v~dřívějších dobách};  \dicPhraseIS{fyrr eða síðar} \dicFlx{adv} \dicDirectTranslationCS{dříve nebo později};  \dicPhraseIS{fyrr en} \dicFlx{conj} \dicDirectTranslationCS{dříve než} \dicIndirectTranslationCS{(používané v~záporu)}
\dicEntry[fyrradagur] \dicTerm{fyrra··dag|ur} \dicIPA{{f}{\textsci}{r}{\textlengthmark}{a}{\textsubring{d}}{a}{\textbabygamma}{\textscy}{\textsubring{r}}} \dicPos{m}[6] \dicFlx{(‑s)}[61] \dicPhraseIS{í fyrradag} \dicFlx{adv} \dicDirectTranslationCS{předevčírem}
\dicEntry[fyrramál] \dicTerm{fyrra··mál} \dicIPA{{f}{\textsci}{r}{\textlengthmark}{a}{m}{au}{\textsubring{l}}} \dicPos{n}[2] \dicFlx{(‑s)}[2] \dicPhraseIS{í fyrramálið} \dicFlx{adv} \dicDirectTranslationCS{zítra ráno}
\dicEntry[fyrrasumar] \dicTerm{fyrra··sum|ar} \dicIPA{{f}{\textsci}{r}{\textlengthmark}{a}{s}{\textscy}{m}{a}{\textsubring{r}}} \dicPos{n}[3] \dicFlx{(‑ars)}[2] \dicPhraseIS{í fyrrasumar} \dicFlx{adv} \dicDirectTranslationCS{předešlé léto, minulé léto}
\dicEntry[fyrri] \dicTerm{fyrri} \dicsymFrequent\  \dicIPA{{f}{\textsci}{r}{\textlengthmark}{\textsci}} \dicPos{adj}[12] \dicFlx{comp}[5] \textbf{1.} \dicDirectTranslationCS{první (ze dvou)} \dicExampleIS{sá fyrri} \dicExampleCS{dříve jmenovaný}  \textbf{2.} \dicDirectTranslationCS{předchozí, dřívější} \dicIndirectTranslationCS{(o~čase)};  \dicPhraseIS{í fyrra} \dicFlx{adv} \dicDirectTranslationCS{(v)loni, minulý rok};  \dicPhraseIS{í fyrra mánuði} \dicFlx{adv} \dicDirectTranslationCS{minulý měsíc, v~předchozím měsíci}
\dicEntry[fyrrn.] \dicTerm{fyrrn.} \dicPos{zkr} \dicPhraseIS{fyrrnefndur} \dicDirectTranslationCS{dříve uvedený\,/\addthin jmenovaný}
\dicEntry[fyrrnefndur] \dicTerm{fyrr··nefndur} \dicsymFrequent\  \dicIPA{{f}{\textsci}{r}{n}{\textepsilon}{m}{\textsubring{d}}{\textscy}{\textsubring{r}}} \dicPos{adj}[2]\dicFlx{}[-17] \dicDirectTranslationCS{prvně jmenovaný, dříve uvedený\,/\addthin zmíněný} \dicExampleIS{fyrrnefndur þáttur} \dicExampleCS{(dříve) zmíněný pořad} \dicAntonym{síðarnefndur}
\dicEntry[fyrrum] \dicTerm{fyrrum} \dicsymFrequent\  \dicIPA{{f}{\textsci}{r}{\textlengthmark}{\textscy}{\textsubring{m}}} \dicPos{adv} \dicSynonym{áður} \dicDirectTranslationCS{dříve, kdysi} \dicExampleIS{Hann var fyrrum tóbaksbóndi.} \dicExampleCS{Kdysi byl pěstitelem tabáku.}
\dicEntry[fyrrv.] \dicTerm{fyrrv.} \dicPos{zkr} \dicPhraseIS{fyrrverandi} \dicDirectTranslationCS{dřívější, bývalý, minulý}
\dicEntry[fyrrverandi] \dicTerm{fyrr··ver·andi} \dicsymFrequent\  \dicIPA{{f}{\textsci}{r}{v}{\textepsilon}{r}{a}{n}{\textsubring{d}}{\textsci}} \dicPos{adj}[13] \dicFlx{indecl}[1] \dicDirectTranslationCS{bývalý, dřívější, minulý} \dicExampleIS{fyrrverandi eiginmaður hennar} \dicExampleCS{její dřívější manžel}
\dicEntry[fyrst] \dicTerm{fyrst\smash{\textsuperscript{1}}} \dicsymFrequent\  \dicIPA{{f}{\textsci}{\textsubring{r}}{s}{\textsubring{d}}} \dicPos{adv} \dicFlx{sup (pos snemma, comp fyrr)} \dicDirectTranslationCS{nejdříve, nejprve, prvně} \dicExampleIS{Þegar ég heyrði þetta fyrst var ég á móti því.} \dicExampleCS{Když jsem to uslyšel, byl jsem nejdříve proti tomu.};  \dicPhraseIS{fyrst og fremst} \dicFlx{adv} \dicSynonym{aðallega} \dicDirectTranslationCS{především, hlavně};  \dicPhraseIS{sem fyrst} \dicFlx{adv} \dicDirectTranslationCS{co nejdříve}
\dicEntry[fyrst] \dicTerm{fyrst\smash{\textsuperscript{2}}} \dicIPA{{f}{\textsci}{\textsubring{r}}{s}{\textsubring{d}}} \dicPos{conj} \dicDirectTranslationCS{neboť, když}
\dicEntry[fyrsti] \dicTerm{fyrsti} \dicIPA{{f}{\textsci}{\textsubring{r}}{s}{\textsubring{d}}{\textsci}} \dicPos{num}[5] \dicFlx{ord} \dicDirectTranslationCS{první}
\dicEntry[fyrstur] \dicTerm{fyrstur} \dicsymFrequent\  \dicIPA{{f}{\textsci}{\textsubring{r}}{s}{\textsubring{d}}{\textscy}{\textsubring{r}}} \dicPos{adj}[1]\dicFlx{}[-13] \dicDirectTranslationCS{první} \dicExampleIS{Hann er fyrstur í röðinni.} \dicExampleCS{Je první na řadě.};  \dicPhraseIS{í fyrsta lagi} \dicFlx{adv} \dicDirectTranslationCS{za prvé};  \dicPhraseIS{í fyrstu} \dicFlx{adv} \dicDirectTranslationCS{z~počátku, nejdříve}
\dicEntry[fyrtast] \dicTerm{fyrtast} \dicIPA{{f}{\textsci}{\textsubring{r}}{\textsubring{d}}{a}{s}{\textsubring{d}}} \dicPos{v}[2] \dicFlx{(fyrtist, fyrst)}[193] \dicFlx{refl} \dicSynonym*{móðgast} \dicDirectTranslationCS{urazit se};  \dicPhraseIS{fyrtast við} \dicDirectTranslationCS{urazit se}
\dicEntry[fýk] \dicTerm{fýk} \dicIPA{{f}{i}{\textlengthmark}{\r{g}}} \dicPos{v} \dicFlx{ind praes sg 1 pers} \dicLink{fjúka}
\dicEntry[fýla] \dicTerm{fýl|a} \dicsymFrequent\  \dicIPA{{f}{i}{\textlengthmark}{l}{a}} \dicPos{f}[1] \dicFlx{(‑u, ‑ur)}[19] \textbf{1.} \dicSynonym*{vond lykt} \dicDirectTranslationCS{(odporný) puch, zápach, smrad} \dicExampleIS{Það er fýla af ruslinu.} \dicExampleCS{Páchnou tu odpadky.}  \textbf{2.} \dicSynonym{ólund} \dicDirectTranslationCS{špatná nálada, mrzutost, naštvanost};  \dicPhraseIS{fara í fýlu} \dicDirectTranslationCS{naštvat se};  \dicPhraseIS{vera í fýlu} \dicDirectTranslationCS{mít špatnou náladu, být naštvaný}  \textbf{3.} \dicSynonym*{þokuveður} \dicDirectTranslationCS{špatné\,/\addthin mizerné počasí}
\dicEntry[fýll] \dicTerm{fýl|l} \dicIPA{{f}{i}{\textsubring{d}}{\textsubring{l}}} \dicPos{m}[6] \dicFlx{(‑s, ‑ar)}[48] \dicFieldCat{zool.} \dicDirectTranslationCS{buřňák lední} \textit{(l.~{\textLA{Fulmarus glacialis}})}  \dicsymPhoto\ 
\dicFigure{20972.jpg}{Fýll}{Fýll - Morkill, Anne, Biolib, PD}
\dicEntry[fýlulegur] \dicTerm{fýlu··legur} \dicIPA{{f}{i}{\textlengthmark}{l}{\textscy}{l}{\textepsilon}{\textbabygamma}{\textscy}{\textsubring{r}}} \dicPos{adj}[1]\dicFlx{}[-8] \dicDirectTranslationCS{podrážděný, mrzutý, nevrlý}
\dicEntry[fýlupoki] \dicTerm{fýlu··pok|i} \dicIPA{{f}{i}{\textlengthmark}{l}{\textscy}{p\smash{\textsuperscript{h}}}{\textopeno}{\r{\textObardotlessj}}{\textsci}} \dicPos{m}[1] \dicFlx{(‑a, ‑ar)}[1] \dicDirectTranslationCS{škarohlíd(ka), bručoun(ka)}
\dicEntry[fýsa] \dicTerm{fýs|a} \dicIPA{{f}{i}{\textlengthmark}{s}{a}} \dicPos{v}[2] \dicFlx{(‑ti, ‑t)}[82] \dicFlx{impers} \dicPhraseIS{e‑n fýsir} \dicDirectTranslationCS{(kdo) touží, (kdo) (moc) chce} \dicExampleIS{Mig fýsir að fara.} \dicExampleCS{Chce se mi jít.}
\dicEntry[fýsilegur] \dicTerm{fýsi··legur} \dicIPA{{f}{i}{\textlengthmark}{s}{\textsci}{l}{\textepsilon}{\textbabygamma}{\textscy}{\textsubring{r}}} \dicPos{adj}[1]\dicFlx{}[-8] \dicSynonym{æskilegur} \dicDirectTranslationCS{žádoucí, kýžený};  \dicPhraseIS{fýsilegur kostur} \dicDirectTranslationCS{atraktivní možnost, optimální řešení}
\dicEntry[fýsn] \dicTerm{fýsn} \dicIPA{{f}{i}{s}{\textsubring{d}}{\textsubring{n}}} \dicPos{f}[7] \dicFlx{(‑ar, ‑ir)}[1] \textbf{1.} \dicSynonym*{holdleg girnd} \dicDirectTranslationCS{chtíč, (tělesná) žádostivost};  \dicPhraseIS{holdleg fýsn} \dicDirectTranslationCS{tělesná žádostivost}  \textbf{2.} \dicSynonym*{áköf löngun} \dicDirectTranslationCS{tužba, chtění}
\dicEntry[fæ] \dicTerm{fæ} \dicIPA{{f}{a}{i}{\textlengthmark}} \dicPos{v} \dicFlx{ind praes sg 1 pers} \dicLink{fá}
\dicEntry[fæddur] \dicTerm{fæddur} \dicIPA{{f}{a}{i}{\textsubring{d}}{\textlengthmark}{\textscy}{\textsubring{r}}} \dicPos{adj}[2]\dicFlx{}[-21] \dicDirectTranslationCS{narozený} \dicExampleIS{Hún er fædd árið 1979.} \dicExampleCS{Narodila se roku 1979.}
\dicEntry[fæð] \dicTerm{fæð} \dicIPA{{f}{a}{i}{\textlengthmark}{\texttheta}} \dicPos{f}[7] \dicFlx{(‑ar, ‑ir)}[1] \textbf{1.} \dicPhraseIS{leggja fæð á e‑n} \dicDirectTranslationCS{nenávidět (koho)}  \textbf{2.} \dicDirectTranslationCS{trocha, málo (malé množství)}
\dicEntry[fæða] \dicTerm{fæð|a\smash{\textsuperscript{1}}} \dicIPA{{f}{a}{i}{\textlengthmark}{ð}{a}} \dicPos{f}[1] \dicFlx{(‑u)}[5] \dicDirectTranslationCS{výživa, potrava}
\dicEntry[fæða] \dicTerm{fæ|ða\smash{\textsuperscript{2}}} \dicsymFrequent\  \dicIPA{{f}{a}{i}{\textlengthmark}{ð}{a}} \dicPos{v}[2] \dicFlx{(‑ddi, ‑tt)}[167] \dicFlx{acc} \textbf{1.} \dicSynonym{fóðra} \dicDirectTranslationCS{(u)živit, (na)krmit} \dicExampleIS{fæða (e‑n) og klæða} \dicExampleCS{živit (koho) a~oblékat}  \textbf{2.} \dicSynonym{ala} \dicDirectTranslationCS{(po)rodit} \dicExampleIS{fæða barn} \dicExampleCS{porodit dítě};  \dicIdiom{fæðast}{ \dicPhraseIS{fæðast}} \dicFlx{refl} \dicDirectTranslationCS{narodit se} \dicExampleIS{Ég fæddist árið 1989.} \dicExampleCS{Narodil jsem se roku 1989.}
\dicEntry[fæði] \dicTerm{fæði} \dicIPA{{f}{a}{i}{\textlengthmark}{ð}{\textsci}} \dicPos{n}[2] \dicFlx{(‑s)}[20] \dicSynonym{matur} \dicDirectTranslationCS{strava, potrava};  \dicPhraseIS{fullt fæði} \dicDirectTranslationCS{plná penze};  \dicPhraseIS{fæði og húsnæði} \dicDirectTranslationCS{strava a~ubytování};  \dicPhraseIS{hálft fæði} \dicDirectTranslationCS{polopenze}
\dicEntry[fæðing] \dicTerm{fæð··ing} \dicsymFrequent\  \dicIPA{{f}{a}{i}{\textlengthmark}{ð}{i}{\ng}{\r{g}}} \dicPos{f}[4] \dicFlx{(‑ar, ‑ar)}[5] \dicSynonym{burður} \dicDirectTranslationCS{narození, porod};  \dicPhraseIS{frá fæðingu} \dicFlx{adv} \dicDirectTranslationCS{od narození}
\dicEntry[fæðingarblettur] \dicTerm{fæðingar··blett|ur} \dicIPA{{f}{a}{i}{ð}{i}{\ng}{\r{g}}{a}{r}{\textsubring{b}}{l}{\textepsilon}{h}{\textsubring{d}}{\textscy}{\textsubring{r}}} \dicPos{m}[9] \dicFlx{(‑s, ‑ir)}[4] \dicDirectTranslationCS{mateřské znaménko}
\dicEntry[fæðingardagur] \dicTerm{fæðingar··dag|ur} \dicIPA{{f}{a}{i}{ð}{i}{\ng}{\r{g}}{a}{r}{\textsubring{d}}{a}{\textbabygamma}{\textscy}{\textsubring{r}}} \dicPos{m}[6] \dicFlx{(‑s, ‑ar)}[62] \dicSynonym{afmæli} \dicDirectTranslationCS{den narození}
\dicEntry[fæðingardeild] \dicTerm{fæðingar··deild} \dicIPA{{f}{a}{i}{ð}{i}{\ng}{\r{g}}{a}{r}{\textsubring{d}}{ei}{l}{\textsubring{d}}} \dicPos{f}[7] \dicFlx{(‑ar, ‑ir)}[1] \dicDirectTranslationCS{porodní oddělení}
\dicEntry[fæðingarheimili] \dicTerm{fæðingar··heimili} \dicIPA{{f}{a}{i}{\textlengthmark}{ð}{i}{\ng}{\r{g}}{a}{\textsubring{r}}{h}{ei}{m}{\textsci}{l}{\textsci}} \dicPos{n}[2] \dicFlx{(‑s, ‑)}[14] \dicDirectTranslationCS{porodnice} \dicIndirectTranslationCS{(zdravotnické zařízení)}
\dicEntry[fæðingarlæknir] \dicTerm{fæðingar··lækn|ir} \dicIPA{{f}{a}{i}{ð}{i}{\ng}{\r{g}}{a}{r}{l}{a}{i}{h}{\r{g}}{n}{\textsci}{\textsubring{r}}} \dicPos{m}[7] \dicFlx{(‑is, ‑ar)}[1] \dicDirectTranslationCS{porodník, porodnice, porodní lékař(ka)}
\dicEntry[fæðingarorlof] \dicTerm{fæðingar··or·lof} \dicIPA{{f}{a}{i}{\textlengthmark}{ð}{i}{\ng}{\r{g}}{a}{r}{\textopeno}{r}{l}{\textopeno}{f}} \dicPos{n}[2] \dicFlx{(‑s, ‑)}[5] \dicDirectTranslationCS{mateřská dovolená}
\dicEntry[fæðingarorlofssjóður] \dicTerm{fæðingar·or·lofs··sjóð|ur} \dicIPA{{f}\-{a}\-{i}\-{\textlengthmark}\-{ð}\-{i}\-{\ng}\-{\r{g}}\-{a}\-{r}\-{\textopeno}\-{r}\-{l}\-{\textopeno}\-{f}\-{s}\-{j}\-{ou}\-{ð}\-{\textscy}\-{\textsubring{r}}\-} \dicPos{m}[9] \dicFlx{(‑s, ‑ir)}[6] \dicDirectTranslationCS{fond, ze kterého je placena mateřská dovolená}
\dicEntry[fæðingarstyrkur] \textls[15]{\dicTerm{fæðingar··styrk|ur} \dicIPA{{f}{a}{i}{ð}{i}{\ng}{\r{g}}{a}{\textsubring{r}}{s}{\textsubring{d}}{\textsci}{\textsubring{r}}{\r{g}}{\textscy}{\textsubring{r}}} \dicPos{m}[9] \dicFlx{(‑s, ‑ir)}[15] \dicDirectTranslationCS{porodné}}
\dicEntry[fæðingarvottorð] \dicTerm{fæðingar··vott·orð} \dicIPA{{f}{a}{i}{ð}{i}{\ng}{\r{g}}{a}{r}{v}{\textopeno}{h}{\textsubring{d}}{\textopeno}{r}{\texttheta}} \dicPos{n}[2] \dicFlx{(‑s, ‑)}[5] \dicDirectTranslationCS{rodný list}
\dicEntry[fægilögur] \dicTerm{fægi··|lögur} \dicIPA{{f}{a}{i}{j}{\textlengthmark}{\textsci}{l}{\oe}{\textbabygamma}{\textscy}{\textsubring{r}}} \dicPos{m}[11] \dicFlx{(‑lagar)}[6] \dicDirectTranslationCS{leštidlo}
\dicEntry[fægja] \dicTerm{fæg|ja} \dicIPA{{f}{a}{i}{j}{\textlengthmark}{a}} \dicPos{v}[2] \dicFlx{(‑ði, ‑t)}[89] \dicFlx{acc} \dicDirectTranslationCS{(vy)leštit, (vy)cídit} \dicExampleIS{fægja silfrið} \dicExampleCS{vycídit stříbro}
\dicEntry[fækka] \dicTerm{fækk|a} \dicsymFrequent\  \dicIPA{{f}{a}{i}{h}{\r{g}}{a}} \dicPos{v}[1] \dicFlx{(‑aði)}[47] \dicFlx{dat} \textbf{1.} \dicDirectTranslationCS{snížit, snižovat, zmenšit, zmenšovat  (počet ap.), (z)redukovat} \dicExampleIS{fækka ferðum} \dicExampleCS{snížit počet cest}  \textbf{2.} \dicPhraseIS{e‑m\,/\addthin e‑u fækkar} \dicFlx{impers} \dicDirectTranslationCS{(koho\,/\addthin čeho) ubývá, snižuje se počet (koho\,/\addthin čeho)} \dicExampleIS{Verslunum hefur fækkað.} \dicExampleCS{Ubylo obchodů.}
\dicEntry[fækkun] \dicTerm{fækk|un} \dicIPA{{f}{a}{i}{h}{\r{g}}{\textscy}{\textsubring{n}}} \dicPos{f}[7] \dicFlx{(‑unar)}[9] \dicDirectTranslationCS{snížení, úbytek, redukce}
\dicEntry[fæla] \dicTerm{fæl|a} \dicIPA{{f}{a}{i}{\textlengthmark}{l}{a}} \dicPos{v}[2] \dicFlx{(‑di, ‑t)}[141] \dicDirectTranslationCS{vystrašit, vylekat, polekat, vyplašit} \dicExampleIS{fæla e‑n burt með hávaða} \dicExampleCS{odstrašit (koho) hlukem};  \dicIdiom{fælast}{ \dicPhraseIS{fælast}} \dicFlx{refl} \dicSynonym*{hræðast} \dicDirectTranslationCS{vystrašit se, vyplašit se} \dicExampleIS{Hesturinn fældist.} \dicExampleCS{Kůň se vyplašil.}
\dicEntry[fæli] \dicTerm{fæli} \dicIPA{{f}{a}{i}{\textlengthmark}{l}{\textsci}} \dicPos{v} \dicFlx{con pf sg 1 pers} \dicLink{fela\smash{\textsuperscript{2}}}
\dicEntry[fælinn] \dicTerm{fælinn} \dicIPA{{f}{a}{i}{\textlengthmark}{l}{\textsci}{\textsubring{n}}} \dicPos{adj}[6]\dicFlx{}[-2] \dicDirectTranslationCS{plachý (kůň ap.)}
\dicEntry[fælni] \dicTerm{fæln|i} \dicIPA{{f}{a}{i}{l}{n}{\textsci}} \dicPos{f}[3] \dicFlx{(‑i)}[3] \dicDirectTranslationCS{bázeň, fobie}
\dicEntry[fær] \dicTerm{fær} \dicsymFrequent\  \dicIPA{{f}{a}{i}{\textlengthmark}{\textsubring{r}}} \dicPos{adj}[5] \dicFlx{(f ‑)}[8] \textbf{1.} \dicSynonym{duglegur} \dicDirectTranslationCS{schopný, způsobilý}  \textbf{2.} \dicSynonym*{farandi} \dicDirectTranslationCS{sjízdný, schůdný} \dicExampleIS{Vegurinn er orðinn fær á ný.} \dicExampleCS{Cesta je už znovu sjízdná.}  \textbf{3.} \dicSynonym{mögulegur} \dicDirectTranslationCS{možný, proveditelný};  \dicPhraseIS{vera fær um að (gera e‑ð)} \dicDirectTranslationCS{být schopný ((co) udělat)};  \dicPhraseIS{vera fær í flestan sjó} \dicLangCat{přen.} \dicDirectTranslationCS{umět si poradit, mít hodně zkušeností}
\dicEntry[færa] \dicTerm{fær|a} \dicsymFrequent\  \dicIPA{{f}{a}{i}{\textlengthmark}{r}{a}} \dicPos{v}[2] \dicFlx{(‑ði, ‑t)}[99] \dicFlx{acc} \textbf{1.} \dicSynonym{flytja} \dicDirectTranslationCS{posunout, posouvat, přesunout, přesouvat, přemístit, přemísťovat, přenést, přenášet};  \dicPhraseIS{færa sig} \dicDirectTranslationCS{přesunout se, přemístit se} \dicExampleIS{Hann færir sig í annað sæti.} \dicExampleCS{Přesouvá se na druhé místo.}  \textbf{2.} \dicSynonym*{koma með} \dicDirectTranslationCS{přinést, přinášet, donést};  \dicPhraseIS{færa e‑m e‑ð} \dicDirectTranslationCS{přinést (komu co)};  \dicPhraseIS{færa e‑m e‑ð að gjöf} \dicDirectTranslationCS{dát (komu co) jako dárek};  \dicIdiom{færa}[fram]{ \dicPhraseIS{færa fram e‑ð}} \dicDirectTranslationCS{předložit (co), představit (co)};  \dicIdiom{færa}[í]{ \dicPhraseIS{færa e‑n í e‑ð}} \dicDirectTranslationCS{obléknout (koho) do (čeho), ustrojit (koho) do (čeho)} \dicExampleIS{færa hana í buxur} \dicExampleCS{obléknout jí kalhoty};  \dicIdiom{færa}[upp]{ \dicPhraseIS{færa upp e‑ð}} \dicDirectTranslationCS{inscenovat (co)};  \dicIdiom{færa}[úr]{ \dicPhraseIS{færa e‑n úr e‑u}} \dicDirectTranslationCS{svléci (koho) z~(čeho), svléknout (koho) z~(čeho)} \dicExampleIS{færa e‑n úr fötunum} \dicExampleCS{svléknout (koho) z~šatů};  \dicIdiom{færa}[út]{ \dicPhraseIS{færa út e‑ð}} \dicDirectTranslationCS{rozšířit (co)}; { \dicPhraseIS{færa út kvíarnar}} \dicDirectTranslationCS{expandovat};  \dicIdiom{færast}{ \dicPhraseIS{færast}} \dicFlx{refl} \dicDirectTranslationCS{posunout se, přesunout se, přemístit se}; { \dicPhraseIS{færast í aukana}} \dicFlx{refl} \dicDirectTranslationCS{nabývat na síle, zvětšovat se}; { \dicPhraseIS{e‑að færist í vöxt}} \dicFlx{refl} \dicDirectTranslationCS{(co) se stává častějším};  \dicIdiom{færast}[undan]{ \dicPhraseIS{færast undan e‑u}} \dicFlx{refl} \dicDirectTranslationCS{vyhýbat se (čemu)}
\dicEntry[færanlegur] \dicTerm{færan··legur} \dicIPA{{f}{a}{i}{\textlengthmark}{r}{a}{n}{l}{\textepsilon}{\textbabygamma}{\textscy}{\textsubring{r}}} \dicPos{adj}[1]\dicFlx{}[-8] \dicDirectTranslationCS{mobilní, přenosný}
\dicEntry[færð] \dicTerm{færð} \dicIPA{{f}{a}{i}{r}{\texttheta}} \dicPos{f}[7] \dicFlx{(‑ar, ‑ir)}[1] \dicDirectTranslationCS{sjízdnost\,/\addthin stav vozovky}
\dicEntry[Færeyingur] \dicTerm{Fær·ey··ing|ur} \dicIPA{{f}{a}{i}{\textlengthmark}{r}{ei}{\textlengthmark}{\ng}{\r{g}}{\textscy}{\textsubring{r}}} \dicPos{m}[6] \dicFlx{(‑s, ‑ar)}[8] \dicDirectTranslationCS{Faeřan(ka)}
\dicEntry[Færeyjar] \dicTerm{Fær··eyjar} \dicIPA{{f}{a}{i}{\textlengthmark}{r}{ei}{j}{a}{\textsubring{r}}} \dicPos{f}[4] \dicFlx{pl}[15] \dicFieldCat{geog.} \dicDirectTranslationCS{Faerské ostrovy}
\dicEntry[færeyska] \dicTerm{fær··eysk|a} \dicIPA{{f}{a}{i}{\textlengthmark}{r}{ei}{s}{\r{g}}{a}} \dicPos{f}[1] \dicFlx{(‑u)}[5] \dicDirectTranslationCS{faerština}
\dicEntry[færeyskur] \dicTerm{fær··ey·skur} \dicIPA{{f}{a}{i}{\textlengthmark}{r}{ei}{s}{\r{g}}{\textscy}{\textsubring{r}}} \dicPos{adj}[1]\dicFlx{}[-1] \dicDirectTranslationCS{faerský}
\dicEntry[færi] \dicTerm{færi\smash{\textsuperscript{1}}} \dicsymFrequent\  \dicIPA{{f}{a}{i}{\textlengthmark}{r}{\textsci}} \dicPos{n}[2] \dicFlx{(‑s, ‑)}[14] \textbf{1.} \dicSynonym{tækifæri} \dicDirectTranslationCS{příležitost, šance};  \dicPhraseIS{fá færi á að (gera e‑ð)} \dicDirectTranslationCS{dostat příležitost k~(udělání (čeho))};  \dicPhraseIS{gefa e‑m færi á að (gera e‑ð)} \dicDirectTranslationCS{dát (komu) příležitost (udělat (co))};  \dicPhraseIS{setja\,/\addthin sitja sig ekki úr færi að (gera e‑ð))} \dicDirectTranslationCS{nenechat si ujít příležitost (udělat (co))}  \textbf{2.} \dicSynonym*{vegalengd skots} \dicDirectTranslationCS{dosah, dostřel}  \textbf{3.} \dicSynonym{veiðarfæri} \dicDirectTranslationCS{rybářské náčiní, rybářská výstroj}  \textbf{4.} \dicSynonym*{færð á leið} \dicDirectTranslationCS{sjízdnost}
\dicEntry[færi] \dicTerm{færi\smash{\textsuperscript{2}}} \dicIPA{{f}{a}{i}{\textlengthmark}{r}{\textsci}} \dicPos{v} \dicFlx{con pf sg 1 pers} \dicLink{fara}
\dicEntry[færiband] \dicTerm{færi··|band} \dicIPA{{f}{a}{i}{\textlengthmark}{r}{\textsci}{\textsubring{b}}{a}{n}{\textsubring{d}}} \dicPos{n}[2] \dicFlx{(‑bands, ‑bönd)}[8] \dicDirectTranslationCS{výrobní linka}
\dicEntry[færni] \dicTerm{færn|i} \dicIPA{{f}{a}{i}{r}{\textsubring{d}}{n}{\textsci}} \dicPos{f}[3] \dicFlx{(‑i)}[3] \dicDirectTranslationCS{dovednost, schopnost} \dicExampleIS{færni í dansi} \dicExampleCS{dovednost tancování}
\dicEntry[færri] \dicTerm{færri} \dicIPA{{f}{a}{i}{r}{\textlengthmark}{\textsci}} \dicPos{adj} \dicFlx{comp m} \dicLink{fár\smash{\textsuperscript{2}}}
\dicEntry[færsla] \dicTerm{færsl|a} \dicIPA{{f}{a}{i}{\textsubring{r}}{s}{\textsubring{d}}{l}{a}} \dicPos{f}[1] \dicFlx{(‑u, ‑ur)}[13] \textbf{1.} \dicSynonym*{það að færa} \dicDirectTranslationCS{přesunutí, převedení}  \textbf{2.} \dicSynonym*{skráning upphæðar} \dicDirectTranslationCS{transakce}  \textbf{3.} \dicFieldCat{poč.} \dicDirectTranslationCS{záznam, položka}
\dicEntry[færslugjald] \dicTerm{færslu··|gjald} \dicIPA{{f}{a}{i}{\textsubring{r}}{s}{\textsubring{d}}{l}{\textscy}{\r{\textObardotlessj}}{a}{l}{\textsubring{d}}} \dicPos{n}[2] \dicFlx{(‑gjalds, ‑gjöld)}[8] \dicFieldCat{ekon.} \dicDirectTranslationCS{transakční poplatek}
\dicEntry[færsluhnappur] \dicTerm{færslu··hnapp|ur} \dicIPA{{f}{a}{i}{\textsubring{r}}{s}{\textsubring{d}}{l}{\textscy}{\textsubring{n}}{a}{h}{\textsubring{b}}{\textscy}{\textsubring{r}}} \dicPos{m}[6] \dicFlx{(‑s, ‑ar)}[25] \dicFieldCat{poč.} \dicDirectTranslationCS{klávesa Enter}
\dicEntry[fæstur] \dicTerm{fæstur} \dicIPA{{f}{a}{i}{s}{\textsubring{d}}{\textscy}{\textsubring{r}}} \dicPos{adj} \dicFlx{m sg nom sup} \dicLink{fár\smash{\textsuperscript{2}}}
\dicEntry[fætur] \dicTerm{fætur} \dicIPA{{f}{a}{i}{\textlengthmark}{\textsubring{d}}{\textscy}{\textsubring{r}}} \dicPos{m} \dicFlx{pl nom} \dicLink{fótur}
\dicEntry[föd.] \dicTerm{föd.} \dicPos{zkr} \dicPhraseIS{föstudagur} \dicDirectTranslationCS{pátek}
\dicEntry[föður] \dicTerm{föður} \dicIPA{{f}{\oe}{\textlengthmark}{ð}{\textscy}{\textsubring{r}}} \dicPos{m} \dicFlx{sg gen} \dicLink{faðir}
\dicEntry[föðurbróðir] \dicTerm{föður··|bróðir} \dicIPA{{f}{\oe}{\textlengthmark}{ð}{\textscy}{r}{\textsubring{b}}{r}{ou}{ð}{\textsci}{\textsubring{r}}} \dicPos{m}[13] \dicFlx{(‑bróður, ‑bræður)}[4] \dicDirectTranslationCS{strýc} \dicIndirectTranslationCS{(bratr otce)} \dicAntonym{móðurbróðir}
\dicEntry[föðurland] \dicTerm{föður··|land} \dicIPA{{f}{\oe}{\textlengthmark}{ð}{\textscy}{r}{l}{a}{n}{\textsubring{d}}} \dicPos{n}[2] \dicFlx{(‑lands, ‑lönd)}[8] \textbf{1.} \dicSynonym{upprunaland} \dicDirectTranslationCS{vlast, rodná země}  \textbf{2.} \dicSynonym*{síðar ullarnærbuxur} \dicDirectTranslationCS{vlněné spodky}
\dicEntry[föðurlandsást] \dicTerm{föður·lands··ást} \dicIPA{{f}{\oe}{\textlengthmark}{ð}{\textscy}{r}{l}{a}{n}{\textsubring{d}}{s}{au}{s}{\textsubring{d}}} \dicPos{f}[7] \dicFlx{(‑ar)}[3] \dicDirectTranslationCS{vlastenectví, patriotismus}
\dicEntry[föðurlandssvik] \dicTerm{föður·lands··svik} \dicIPA{{f}{\oe}{\textlengthmark}{ð}{\textscy}{r}{l}{a}{n}{\textsubring{d}}{s}{v}{\textsci}{\r{g}}} \dicPos{n}[2] \dicFlx{pl}[1] \dicDirectTranslationCS{vlastizrada}
\dicEntry[föðurlandssvikari] \dicTerm{föður·lands··svik·ar|i} \dicIPA{{f}\-{\oe}\-{\textlengthmark}\-{ð}\-{\textscy}\-{r}\-{l}\-{a}\-{n}\-{\textsubring{d}}\-{s}\-{v}\-{\textsci}\-{\r{g}}\-{a}\-{r}\-{\textsci}\-} \dicPos{m}[1] \dicFlx{(‑a, ‑ar)}[13] \dicDirectTranslationCS{vlastizrádce, vlastizrádkyně}
\dicEntry[föðurlandsvinur] \dicTerm{föður·lands··vin|ur} \dicIPA{{f}{\oe}{\textlengthmark}{ð}{\textscy}{r}{l}{a}{n}{\textsubring{d}}{s}{v}{\textsci}{n}{\textscy}{\textsubring{r}}} \dicPos{m}[10] \dicFlx{(‑ar, ‑ir)}[17] \dicDirectTranslationCS{vlastenec, vlastenka, patriot(ka)}
\dicEntry[föðurlegur] \dicTerm{föður··legur} \dicIPA{{f}{\oe}{\textlengthmark}{ð}{\textscy}{r}{l}{\textepsilon}{\textbabygamma}{\textscy}{\textsubring{r}}} \dicPos{adj}[1]\dicFlx{}[-8] \dicDirectTranslationCS{otcovský}
\dicEntry[föðursystir] \dicTerm{föður··syst|ir} \dicIPA{{f}{\oe}{\textlengthmark}{ð}{\textscy}{\textsubring{r}}{s}{\textsci}{s}{\textsubring{d}}{\textsci}{\textsubring{r}}} \dicPos{f}[11] \dicFlx{(‑ur, ‑ur)}[3] \dicDirectTranslationCS{teta} \dicIndirectTranslationCS{(sestra otce)} \dicAntonym{móðursystir}
\dicEntry[föflu] \dicTerm{föflu} \dicIPA{{f}{\oe}{\textsubring{b}}{l}{\textscy}} \dicPos{f} \dicFlx{sg gen} \dicLink{fafla}
\dicEntry[fög] \dicTerm{fög} \dicIPA{{f}{\oe}{\textlengthmark}{x}} \dicPos{n} \dicFlx{pl nom} \dicLink{fag}
\dicEntry[föggur] \dicTerm{föggur} \dicIPA{{f}{\oe}{\r{g}}{\textlengthmark}{\textscy}{\textsubring{r}}} \dicPos{f}[12] \dicFlx{pl}[3] \dicDirectTranslationCS{věci, zavazadla}
\dicEntry[fögnuður] \dicTerm{fögn··uður} \dicsymFrequent\  \dicIPA{{f}{\oe}{\r{g}}{n}{\textscy}{ð}{\textscy}{\textsubring{r}}} \dicPos{m}[10] \dicFlx{(fögnuðar\,/\addthin fagnaðar, fögnuðir)}[44] \dicSynonym{ánægja} \dicDirectTranslationCS{radost, radování se, slavení} \dicExampleIS{fögnuður yfir e‑u} \dicExampleCS{radost z~(čeho)}
\dicEntry[fögur] \dicTerm{fögur} \dicIPA{{f}{\oe}{\textlengthmark}{\textbabygamma}{\textscy}{\textsubring{r}}} \dicPos{adj} \dicFlx{f sg nom pos} \dicLink{fagur}
\dicEntry[föl] \dicTerm{föl\smash{\textsuperscript{1}}} \dicIPA{{f}{\oe}{\textlengthmark}{\textsubring{l}}} \dicPos{f}[4] \dicFlx{(‑ar)}[3] \dicLink{föl\smash{\textsuperscript{2}}}
\dicEntry[föl] \dicTerm{föl\smash{\textsuperscript{2}}} \dicIPA{{f}{\oe}{\textlengthmark}{\textsubring{l}}} \dicPos{n}[2] \dicFlx{(‑s)}[2] \dicDirectTranslationCS{sněhový poprašek}
\dicEntry[föl] \dicTerm{föl\smash{\textsuperscript{3}}} \dicIPA{{f}{\oe}{\textlengthmark}{\textsubring{l}}} \dicPos{adj} \dicFlx{f sg nom pos} \dicLink{falur}
\dicEntry[földum] \dicTerm{földum} \dicIPA{{f}{\oe}{l}{\textsubring{d}}{\textscy}{\textsubring{m}}} \dicPos{v} \dicFlx{ind pf pl 1 pers} \dicLink{fela\smash{\textsuperscript{1}}}
\dicEntry[föll] \dicTerm{föll} \dicIPA{{f}{\oe}{\textsubring{d}}{\textsubring{l}}} \dicPos{n} \dicFlx{pl nom} \dicLink{fall}
\dicEntry[fölleitur] \dicTerm{föl··leitur} \dicIPA{{f}{\oe}{l}{\textlengthmark}{ei}{\textsubring{d}}{\textscy}{\textsubring{r}}} \dicPos{adj}[1]\dicFlx{}[-1] \dicDirectTranslationCS{bledý, pobledlý}
\dicEntry[fölna] \dicTerm{föln|a} \dicIPA{{f}{\oe}{l}{n}{a}} \dicPos{v}[1] \dicFlx{(‑aði)}[34] \textbf{1.} \dicDirectTranslationCS{(z)blednout}  \textbf{2.} \dicSynonym{sölna} \dicDirectTranslationCS{(z)vadnout, uvadnout, uvadat} \dicExampleIS{Gróðurinn er farinn að fölna.} \dicExampleCS{Vegetace začíná vadnout.}
\dicEntry[fölsk] \dicTerm{fölsk} \dicIPA{{f}{\oe}{l}{s}{\r{g}}} \dicPos{adj} \dicFlx{f sg nom pos} \dicLink{falskur}
\dicEntry[fölsun] \dicTerm{fölsun} \dicIPA{{f}{\oe}{l}{s}{\textscy}{\textsubring{n}}} \dicPos{f}[7] \dicFlx{(‑ar, falsanir)}[16] \dicDirectTranslationCS{padělání, falšování, falzifikace} \dicExampleIS{fölsun á peningum} \dicExampleCS{padělání peněz}
\dicEntry[fölur] \dicTerm{fölur} \dicsymFrequent\  \dicIPA{{f}{\oe}{\textlengthmark}{l}{\textscy}{\textsubring{r}}} \dicPos{adj}[1]\dicFlx{}[-1] \dicDirectTranslationCS{bledý, pobledlý, sinalý, popelavý} \dicExampleIS{fölur á vangann} \dicExampleCS{bledý v~tváři}
\dicEntry[föndra] \dicTerm{föndr|a} \dicIPA{{f}{\oe}{n}{\textsubring{d}}{r}{a}} \dicPos{v}[1] \dicFlx{(‑aði)}[34] \dicDirectTranslationCS{tvořit, vytvářet (rukodělné výrobky ap.)};  \dicPhraseIS{föndra við e‑ð} \dicDirectTranslationCS{hrát si s~(čím) (usilovně se něčím zabývat)}
\dicEntry[föndur] \dicTerm{föndur} \dicIPA{{f}{\oe}{n}{\textsubring{d}}{\textscy}{\textsubring{r}}} \dicPos{n}[2] \dicFlx{(‑s)}[28] \textbf{1.} \dicSynonym*{það að föndra} \dicDirectTranslationCS{tvoření, kutilství}  \textbf{2.} \dicDirectTranslationCS{výtvor, vytvořená věc (ze zájmu)}
\dicEntry[föng] \dicTerm{föng} \dicIPA{{f}{\oe i}{\ng}{\r{g}}} \dicPos{n} \dicFlx{pl nom} \dicLink{fang}
\dicEntry[föngulegur] \dicTerm{föngu··legur} \dicIPA{{f}{\oe i}{\ng}{\r{g}}{\textscy}{l}{\textepsilon}{\textbabygamma}{\textscy}{\textsubring{r}}} \dicPos{adj}[1]\dicFlx{}[-8] \dicDirectTranslationCS{atraktivní, přitažlivý}
\dicEntry[fönn] \dicTerm{fönn} \dicIPA{{f}{\oe}{n}{\textlengthmark}} \dicPos{f}[7] \dicFlx{(fannar, fannir)}[16] \dicDirectTranslationCS{(sněhová) závěj}
\dicEntry[för] \dicTerm{för\smash{\textsuperscript{1}}} \dicsymFrequent\  \dicIPA{{f}{\oe}{\textlengthmark}{\textsubring{r}}} \dicPos{f}[7] \dicFlx{(farar, farir)}[16] \dicDirectTranslationCS{cesta, výlet, výprava, plavba} \dicExampleIS{för hennar um Asíu} \dicExampleCS{její cesta po Asii};  \dicPhraseIS{e‑að hefur e‑ð í för með sér} \dicDirectTranslationCS{(co) sebou nese (co), (co) má za následek (co)};  \dicPhraseIS{slást í för með e‑m} \dicFlx{refl} \dicDirectTranslationCS{přidat se ke (komu), připojit se ke (komu)};  \dicPhraseIS{vera á förum} \dicDirectTranslationCS{být na cestě};  \dicPhraseIS{vera í för með e‑m} \dicDirectTranslationCS{cestovat spolu s~(kým)};  \dicPhraseIS{vera til fara} \dicDirectTranslationCS{být dobře oblečen}
\dicEntry[för] \dicTerm{för\smash{\textsuperscript{2}}} \dicIPA{{f}{\oe}{\textlengthmark}{\textsubring{r}}} \dicPos{n} \dicFlx{pl nom} \dicLink{far}
\dicEntry[förðun] \dicTerm{förð|un} \dicIPA{{f}{\oe}{r}{ð}{\textscy}{\textsubring{n}}} \dicPos{f}[7] \dicFlx{(‑unar)}[12] \dicDirectTranslationCS{(na)líčení, make‑up}
\dicEntry[förg] \dicTerm{förg} \dicIPA{{f}{\oe}{r}{\r{g}}} \dicPos{n} \dicFlx{pl nom} \dicLink{farg}
\dicEntry[förufálki] \dicTerm{föru··fálk|i} \dicIPA{{f}{\oe}{\textlengthmark}{r}{\textscy}{f}{au}{\textsubring{l}}{\r{\textObardotlessj}}{\textsci}} \dicPos{m}[1] \dicFlx{(‑a, ‑ar)}[1] \dicFieldCat{zool.} \dicDirectTranslationCS{sokol stěhovavý} \textit{(l.~{\textLA{Falco peregrinus}})}  \dicsymPhoto\ 
\dicFigure{6174.jpg}{Förufálki}{Förufálki - Deml Miroslav, Biolib, Copyright/CC-BY}
\dicEntry[förumunkur] \dicTerm{föru··munk|ur} \dicIPA{{f}{\oe}{\textlengthmark}{r}{\textscy}{m}{u}{\r{\ng}}{\r{g}}{\textscy}{\textsubring{r}}} \dicPos{m}[6] \dicFlx{(‑s, ‑ar)}[10] \dicDirectTranslationCS{žebravý mnich}
\dicEntry[förunautur] \dicTerm{föru··naut|ur} \dicIPA{{f}{\oe}{\textlengthmark}{r}{\textscy}{n}{\oe i}{\textsubring{d}}{\textscy}{\textsubring{r}}} \dicPos{m}[6] \dicFlx{(‑s\,/\addthin ‑ar, ‑ar)}[65] \dicSynonym{samferðamaður} \dicDirectTranslationCS{společník, společnice, druh\,/\addthin družka (na cestě ap.)}
\dicEntry[föruneyti] \dicTerm{föru··neyti} \dicIPA{{f}{\oe}{\textlengthmark}{r}{\textscy}{n}{ei}{\textsubring{d}}{\textsci}} \dicPos{n}[2] \dicFlx{(‑s, ‑)}[14] \dicDirectTranslationCS{doprovod, svita, suita}
\dicEntry[föst] \dicTerm{föst} \dicIPA{{f}{\oe}{s}{\textsubring{d}}} \dicPos{adj} \dicFlx{f sg nom pos} \dicLink{fastur}
\dicEntry[föstu] \dicTerm{föstu} \dicIPA{{f}{\oe}{s}{\textsubring{d}}{\textscy}} \dicPos{f} \dicFlx{sg gen} \dicLink{fasta\smash{\textsuperscript{1}}}
\dicEntry[föstudagskvöld] \dicTerm{föstu·dags··kvöld} \dicIPA{{f}{\oe}{s}{\textsubring{d}}{\textscy}{\textsubring{d}}{a}{x}{s}{k\smash{\textsuperscript{h}}}{v}{\oe}{l}{\textsubring{d}}} \dicPos{n}[2] \dicFlx{(‑s, ‑)}[5] \dicDirectTranslationCS{páteční večer, pátek večer}
\dicEntry[föstudagur] \dicTerm{föstu··dag|ur} \dicsymFrequent\  \dicIPA{{f}{\oe}{s}{\textsubring{d}}{\textscy}{\textsubring{d}}{a}{\textbabygamma}{\textscy}{\textsubring{r}}} \dicPos{m}[6] \dicFlx{(‑s, ‑ar)}[62] \dicDirectTranslationCS{pátek};  \dicPhraseIS{á föstudaginn} \dicFlx{adv} \dicDirectTranslationCS{v~pátek};  \dicPhraseIS{á föstudaginn kemur} \dicFlx{adv} \dicDirectTranslationCS{v~(nejbližší) pátek};  \dicPhraseIS{á föstudaginn var} \dicFlx{adv} \dicDirectTranslationCS{minulý pátek};  \dicPhraseIS{á föstudögum} \dicFlx{adv} \dicDirectTranslationCS{každý pátek};  \dicPhraseIS{föstudagurinn langi} \dicFieldCat{náb.} \dicDirectTranslationCS{Velký pátek}
\dicEntry[föstuinngangur] \dicTerm{föstu··inn·gang|ur} \dicIPA{{f}{\oe}{s}{\textsubring{d}}{\textscy}{\textsci}{n}{\r{g}}{au}{\ng}{\r{g}}{\textscy}{\textsubring{r}}} \dicPos{m}[6] \dicFlx{(‑s)}[7] \dicFieldCat{náb.} \dicDirectTranslationCS{počátek půstu}
\dicEntry[föstur] \dicTerm{föstur} \dicIPA{{f}{\oe}{s}{\textsubring{d}}{\textscy}{\textsubring{r}}} \dicPos{f} \dicFlx{pl nom} \dicLink{fasta\smash{\textsuperscript{1}}}
\dicEntry[föt] \dicTerm{föt} \dicIPA{{f}{\oe}{\textlengthmark}{\textsubring{d}}} \dicPos{n}[2] \dicFlx{pl}[9] \dicDirectTranslationCS{oblečení, šaty, oděv}
\dicEntry[fötluð] \dicTerm{fötluð} \dicIPA{{f}{\oe}{h}{\textsubring{d}}{l}{\textscy}{\texttheta}} \dicPos{adj} \dicFlx{f sg nom pos} \dicLink{fatlaður}
\dicEntry[fötlun] \dicTerm{fötlun} \dicIPA{{f}{\oe}{h}{\textsubring{d}}{l}{\textscy}{\textsubring{n}}} \dicPos{f}[7] \dicFlx{(‑ar, fatlanir)}[16] \dicDirectTranslationCS{invalidita, handicap}
\dicEntry[fötu] \dicTerm{fötu} \dicIPA{{f}{\oe}{\textlengthmark}{\textsubring{d}}{\textscy}} \dicPos{f} \dicFlx{sg gen} \dicLink{fata\smash{\textsuperscript{1}}}
\dicEntry[föx] \dicTerm{föx} \dicIPA{{f}{\oe}{x}{s}} \dicPos{n} \dicFlx{pl nom} \dicLink{fax\smash{\textsuperscript{1}}}

\else
\makeatletter
\@for\next:=\alphabetF\do{\input{letters/\next.tex}}
\makeatother 
\fi
\restoregeometry
\pagestyle{plain}

\fi

% ================================================================ PHONETICS =

\ifinputphon

\cleardoublepage

\chapter{Seznam islandských fonémů}              \label{sec:phon_phonems}

\begin{table}[h]
\begin{tabular}{lllll} \toprule

\textbf{foném} & \textbf{vysvětlivka} & \textbf{příklad} & \textbf{fonetický přepis} & \textbf{překlad} \\
\midrule
{\textipa{[{a}]}} & otevřená přední nezaokrouhlená samohl. & rata & {\textipa{[{r}{a}{\textlengthmark}{\textsubring{d}}{a}]}} & trefit \\ 
{\textipa{[{\textepsilon}]}} & polootevřená přední nezaokrouhlená samohl. & meta & {\textipa{[{m}{\textepsilon}{\textlengthmark}{\textsubring{d}}{a}]}} & ocenit \\ 
{\textipa{[{e}]}} & polozavřená přední nezaokrouhlená samohl. & heiður & {\textipa{[{h}{ei}{\textlengthmark}{ð}{\textscy}{\textsubring{r}}]}} & čest \\
{\textipa{[{\textsci}]}} & téměř zavřená téměř přední nezaokrouhlená samohl. & vinna & {\textipa{[{v}{\textsci}{n}{\textlengthmark}{a}]}} & pracovat \\ 
{\textipa{[{i}]}} & zavřená přední nezaokrouhlená samohl. & fínn & {\textipa{[{f}{i}{\textsubring{d}}{\textsubring{n}}]}} & jemný \\ 
{\textipa{[{\textopeno}]}} & polootevřená zadní zaokrouhlená samohl. & lofa & {\textipa{[{l}{\textopeno}{\textlengthmark}{v}{a}]}} & slibovat \\ 
{\textipa{[{o}]}} & polozavřená zadní zaokrouhlená samohl. & bóndi & {\textipa{[{\textsubring{b}}{ou}{n}{\textsubring{d}}{\textsci}]}} & farmář \\ 
{\textipa{[{\textscy}]}} & téměř zavřená téměř přední zaokrouhlená samohl. & hundur & {\textipa{[{h}{\textscy}{n}{\textsubring{d}}{\textscy}{\textsubring{r}}]}} & pes \\ 
{\textipa{[{u}]}} & zavřená zadní zaokrouhlená samohl. & súkkulaði & {\textipa{[{s}{u}{h}{\r{g}}{\textscy}{l}{a}{ð}{\textsci}]}} & čokoláda \\ 
{\textipa{[{\oe}]}} & téměř otevřená přední nezaokrouhlená samohl. & lönd & {\textipa{[{l}{\oe}{n}{\textsubring{d}}]}} & země \\ 
{\textipa{[{m}]}} & bilabiální nazála & mýri & {\textipa{[{m}{i}{\textlengthmark}{r}{\textsci}]}} & mokřina \\ 
{\textipa{[{\textsubring{m}}]}} & neznělá bilabiální nazála & þyrmt & {\textipa{[{\texttheta}{\textsci}{\textsubring{r}}{\textsubring{m}}{\textsubring{d}}]}} & {\textit{$\shortrightarrow$ þyrma}} \\  
{\textipa{[{n}]}} & alveolární nazála & nemandi & {\textipa{[{n}{\textepsilon}{\textlengthmark}{m}{a}{n}{\textsubring{d}}{\textsci}]}} & žák \\ 
{\textipa{[{\textsubring{n}}]}} & neznělá alveolární nazála & mennt & {\textipa{[{m}{\textepsilon}{\textsubring{n}}{\textsubring{d}}]}} & vzdělání \\ 
{\textipa{[{\textltailn}]}} & palatální nazála & angi & {\textipa{[{au}{\textltailn}{\r{\textObardotlessj}}{\textsci}]}} & drobek \\ 
{\textipa{[{\r{\textltailn}}]}} & neznělá palatální nazála & banki & {\textipa{[{\textsubring{b}}{au}{\r{\textltailn}}{\r{\textObardotlessj}}{\textsci}]}} & banka \\
{\textipa{[{\ng}]}} & velární nazála & skyggndur & {\textipa{[{s}{\r{\textObardotlessj}}{\textsci}{\ng}{\textsubring{d}}{\textscy}{\textsubring{r}}]}} & {\textit{$\shortrightarrow$ skyggna}} \\ 
{\textipa{[{\r{\ng}}]}} & neznělá velární nazála & skyggnt & {\textipa{[{s}{\r{\textObardotlessj}}{\textsci}{\r{\ng}}{\textsubring{d}}]}} & {\textit{$\shortrightarrow$ skyggna}} \\ 
{\textipa{[{\textsubring{b}}]}} & neznělá bilabiální ploziva & bara & {\textipa{[{\textsubring{b}}{a}{\textlengthmark}{r}{a}]}} & jenom \\ 
{\textipa{[{p\textsuperscript{h}}]}} & neznělá bilabiální ploziva s přídechem & prestur & {\textipa{[{p\textsuperscript{h}}{r}{\textepsilon}{s}{\textsubring{d}}{\textscy}{\textsubring{r}}]}} & kněz \\ 
{\textipa{[{\textsubring{d}}]}} & neznělá alveolární ploziva & dæmi & {\textipa{[{\textsubring{d}}{a}{i}{\textlengthmark}{m}{\textsci}]}} & příklad \\ 
{\textipa{[{t\textsuperscript{h}}]}} & neznělá alveolární ploziva s přídechem & tómur & {\textipa{[{t\textsuperscript{h}}{ou}{\textlengthmark}{m}{\textscy}{\textsubring{r}}]}} & prázdný \\ 
{\textipa{[{\r{\textObardotlessj}}]}} & neznělá palatální ploziva & gefa & {\textipa{[{\r{\textObardotlessj}}{\textepsilon}{\textlengthmark}{v}{a}]}} & dát \\ 
{\textipa{[{c\textsuperscript{h}}]}} & neznělá palatální ploziva s přídechem & kenna & {\textipa{[{c\textsuperscript{h}}{\textepsilon}{n}{\textlengthmark}{a}]}} & učit \\
{\textipa{[{\r{g}}]}} & neznělá velární ploziva & gata & {\textipa{[{\r{g}}{a}{\textlengthmark}{\textsubring{d}}{a}]}} & ulice \\ 
{\textipa{[{k\textsuperscript{h}}]}} & neznělá velární ploziva s přídechem & koma & {\textipa{[{k\textsuperscript{h}}{\textopeno}{\textlengthmark}{m}{a}]}} & přijít \\
{\textipa{[{v}]}} & znělá labiodentální frikativa & hafa & {\textipa{[{h}{a}{\textlengthmark}{v}{a}]}} & mít \\ 
{\textipa{[{f}]}} & neznělá labiodentální frikativa & finna & {\textipa{[{f}{\textsci}{n}{\textlengthmark}{a}]}} & najít \\ 
{\textipa{[{ð}]}} & znělá dentální frikativa & borða & {\textipa{[{\textsubring{b}}{\textopeno}{r}{ð}{a}]}} & jíst \\ 
{\textipa{[{\texttheta}]}} & neznělá dentální frikativa & þola & {\textipa{[{\texttheta}{\textopeno}{\textlengthmark}{l}{a}]}} & snést \\
{\textipa{[{s}]}} & neznělá alveolární frikativa & systir & {\textipa{[{s}{\textsci}{s}{\textsubring{d}}{\textsci}{\textsubring{r}}]}} & sestra \\ 
{\textipa{[{\c{c}}]}} & neznělá palatální frikativa & hjóla & {\textipa{[{\c{c}}{ou}{\textlengthmark}{l}{a}]}} & jet na kole \\ 
{\textipa{[{\textbabygamma}]}} & znělá velární frikativa & sagði & {\textipa{[{s}{a}{\textbabygamma}{ð}{\textsci}]}} & řekl \\ 
{\textipa{[{x}]}} & neznělá velární frikativa & lag & {\textipa{[{l}{a}{\textlengthmark}{x}]}} & píseň \\
{\textipa{[{h}]}} & neznělá glotální frikativa & heima & {\textipa{[{h}{ei}{\textlengthmark}{m}{a}]}} & doma \\ 
{\textipa{[{l}]}} & alveolární laterální aproximanta & loforð & {\textipa{[{l}{\textopeno}{\textlengthmark}{v}{\textopeno}{r}{\texttheta}]}} & slib \\ 
{\textipa{[{\textsubring{l}}]}} & neznělá alveolární laterální aproximanta & hlutur & {\textipa{[{\textsubring{l}}{\textscy}{\textlengthmark}{\textsubring{d}}{\textscy}{\textsubring{r}}]}} & věc \\  
{\textipa{[{j}]}} & palatální aproximanta & jól & {\textipa{[{j}{ou}{\textlengthmark}{\textsubring{l}}]}} & Vánoce \\ 
{\textipa{[{r}]}} & alveolární vibranta & reynsla & {\textipa{[{r}{ei}{n}{s}{\textsubring{d}}{l}{a}]}} & zkušenost \\ 
{\textipa{[{\textsubring{r}}]}} & neznělá alveolární vibranta & hraun & {\textipa{[{\textsubring{r}}{\oe i}{\textlengthmark}{\textsubring{n}}]}} & láva \\ 
\bottomrule
\end{tabular}
\end{table}

\clearpage

%\section{Výslovnost islandských samohlásek}      \label{sec:phon_vowels}
%\LTXtable{\columnwidth}{phon/vowels}
%\clearpage

%\section{Výslovnost islandských souhlásek}       \label{sec:phon_consonants}
%\LTXtable{\columnwidth}{phon/consonants}
%\clearpage
\fi

% =============================================================== MORPHOLOGY =

\ifinputmorpho

\cleardoublepage

\chapter{Morfologický klíč}                      \label{sec:morpho}

\section{Podstatné jméno, rod mužský}            \label{sec:morpho_m}
{\small\LTXfw{morpho/m}}
\clearpage

\section{Podstatné jméno, rod ženský}            \label{sec:morpho_f}
{\small\LTXfw{morpho/f}}
\clearpage

\section{Podstatné jméno, rod střední}           \label{sec:morpho_n}
{\small\LTXfw{morpho/n}}
\clearpage

\section{Přídavné jméno}                         \label{sec:morpho_adj}
{\small\LTXfw{morpho/adj}}
\clearpage

\section{Zájmeno}                                \label{sec:morpho_pron}
{\small\LTXfw{morpho/pron}}
 \begin{center}
\begin{minipage}[t]{.45\textwidth}
\begin{tabular}{l>{\footnotesize\itshape}lll}
 {\textbf{\textit{třída}}} & {\textit{pád}}   & \textit{sg} & \textit{pl}  \\ 
 \hline
\multirow{3}{*}{{{\textbf{pron} \Large{\textbf{17}}}}}  &  nom & ég & við   \\*
 & acc &  mig  & okkur  \\*
 & dat & mér & okkur   \\*
 & gen & mín  & okkar  \\
\hline

\multirow{3}{*}{{{\textbf{pron} \Large{\textbf{17}}}}}  &  nom & hann & þeir   \\*
 & acc &  hann  & þá  \\*
 & dat & honum & þeim   \\*
 & gen & hans  & þeirra  \\
\hline

\multirow{3}{*}{{{\textbf{pron} \Large{\textbf{17}}}}}  &  nom & hún & þær   \\*
 & acc &  hana  & þær  \\*
 & dat & henni & þeim   \\*
 & gen & hennar  & þeirra  \\
\hline

\multirow{3}{*}{{{\textbf{pron} \Large{\textbf{17}}}}}  &  nom &  & vér   \\*
 & acc &    & oss  \\*
 & dat &  & oss   \\*
 & gen &   & vor  \\
\hline

\multirow{3}{*}{{{\textbf{pron} \Large{\textbf{17}}}}}  &  nom &  & þér   \\*
 & acc &    & yður  \\*
 & dat &  & yður   \\*
 & gen &   & yðar  \\
\hline

\multirow{3}{*}{{{\textbf{pron} \Large{\textbf{17}}}}}  &  nom & þú & þið   \\*
 & acc &  þig  & ykkur  \\*
 & dat & þér & ykkur   \\*
 & gen & þín  & ykkar  \\
\hline

\multirow{3}{*}{{{\textbf{pron} \Large{\textbf{17}}}}}  &  nom & það & þau   \\*
 & acc &  það  & þau  \\*
 & dat & því & þeim   \\*
 & gen & þess  & þeirra  \\
\hline
\end{tabular}
\end{minipage}\hfil
\begin{minipage}[t]{.45\textwidth}
\begin{tabular}{l>{\footnotesize\itshape}lll}
\hline
\multirow{3}{*}{{{\textbf{pron} \Large{\textbf{18}}}}}  &  nom &  &    \\*
 & acc &  sig  &   \\*
 & dat & sér &    \\*
 & gen & sín  &   \\
 \hline
\end{tabular}
\end{minipage}
\end{center}

\clearpage

\section{Číslovka}                               \label{sec:morpho_num}
 {\small\LTXfw{morpho/num}}
\clearpage

\begin{landscape}
\section{Sloveso}                                \label{sec:morpho_v}
{\small\LTXfw{morpho/v}}
\end{landscape}

\clearpage

\fi

% ============================================================== BACK MATTER =

\ifmakebackmatter

% ------------------------------------------------ Photographs authors index -

\cleardoublepage
%\pagestyle{empty}
%\pagestyle{indexstyle}
%\phantomsection
%\addcontentsline{toc}{chapter}
%  {Seznam autorů fotografií a ilustrací}        
\label{sec:photo}
\printindex[figures]

% ------------------------------------------------------------- Bibliography -
\cleardoublepage
\chapter{Bibliografie}
%\twocolumn
%\phantomsection
%\addcontentsline{toc}{chapter}{Bibliografie}
\begin{multicols}{2}
\nocite{*}\printbibliography[heading=none]
\end{multicols}

% ------------------------------------------------------------------ Apendix -
\cleardoublepage
\onecolumn
\ifPDF
\else
\chapter{Apendix: Islandská gramatika}
\clearpage
\covergeometry  
\makeatletter
\includepdf[
  pages={1,2,3,4,5,6,7,8,9,10,11,12,13,14,15,16,17,18,19,20,21,22,23,24,25,26,27,28,29,30,31,32,33,34,35,36,37,38,39,40,41,42,43,44,45,46,47,48,49,50,51,52,53,54,55,56,57,58,59,60,61,62,63,64,65,66,67,68,69,70,71,72,73},
  turn=false,
  scale=0.79,
  keepaspectratio,
  offset={\dimexpr(\paperwidth-\csname Gm@layoutwidth\endcsname)/2-\csname Gm@layouthoffset\endcsname\relax}
  {\dimexpr(\paperheight-\csname  Gm@layoutheight\endcsname)/2-\csname Gm@layoutvoffset\endcsname\relax}
   ]{gram.pdf}
\makeatother
\restoregeometry  
\fi

% ------------------------------------------------------------------ Barcode -
\pagenumbering{gobble}
\ifPDF
\else
\cleardoublepage\null\clearpage % clear to verso
\vspace*\fill
\EANisbn[SC3]
\fi
\fi

% --------------------------------------------------------------- Back cover -

\ifmakecovers
  \cleardoublepage\null\clearpage % inner side of the cover
\covergeometry  
  \makecoverwith\backcoverimages
\restoregeometry  
\fi

\end{document}

